\AtBeginDocument{%
  \providecommand\BibTeX{{Bib\TeX}}
}

\usepackage{mathtools}
\usepackage{mathrsfs}
\usepackage{listings}
\usepackage{multirow}
\usepackage{graphicx}
%% \iflatexml is predefined as true by LaTeXML (the engine arXiv uses to
%% render the HTML version). Under pdfLaTeX it is undefined, so we declare it
%% here defaulting to false; LaTeXML preserves its own true value and ignores
%% this reset. This lets us ship a PDF-only path (absolute page positions via
%% zref-savepos) alongside an HTML fallback that needs no page coordinates.
\newif\iflatexml
\latexmlfalse
\iflatexml\else
  \usepackage{zref-savepos}
\fi
\usepackage{tikz}
\usetikzlibrary{positioning,arrows.meta,calc}

%% Allow multi-line display math (align, etc.) to break across pages, so
%% the long BNF grammars are not trapped as single unbreakable blocks.
\allowdisplaybreaks

%% Right-aligned annotations for a multi-line math display.
%% Nested \left\{...\right\} groups cannot share a common alignment column,
%% so each annotation is pinned to a single right edge via absolute page
%% positions recorded by zref-savepos (resolved on the second compiler pass).
%% \figureAnnotationAnchor records the display's left edge (call once just
%% inside the math, before the outermost \begin{aligned}); \figureAnnotation
%% lays out a zero-width box whose right end sits at (left edge + \linewidth),
%% i.e. the enclosing minipage's right margin, regardless of nesting depth.
\newcounter{figureAnnotationCounter}
\iflatexml
  %% HTML fallback (LaTeXML / arXiv HTML): page-absolute positioning via
  %% zref-savepos is unavailable, and the MathML output flows rather than being
  %% laid out at fixed page coordinates, so a pixel-accurate right edge cannot
  %% be reproduced. Instead, render each annotation as a small inline label that
  %% trails the line it describes; the browser wraps it naturally. The anchor
  %% becomes a no-op since no left edge needs to be recorded.
  \newcommand{\figureAnnotationAnchor}{}
  \newcommand{\figureAnnotation}[1]{%
    \quad\normalfont\text{\scriptsize(#1)}%
  }
\else
  \newcommand{\figureAnnotationAnchor}{\zsavepos{figureAnnotationLeft}}
  \newcommand{\figureAnnotation}[1]{%
    \stepcounter{figureAnnotationCounter}%
    \zsavepos{figureAnnotation\the\value{figureAnnotationCounter}}%
    \rlap{\hbox to
      \dimexpr\zposx{figureAnnotationLeft}sp+\linewidth-\zposx{figureAnnotation\the\value{figureAnnotationCounter}}sp\relax
    {\hfill\normalfont\text{\scriptsize #1}}}%
  }
\fi

\AtBeginDocument{%
  \theoremstyle{acmdefinition}%
  \newtheorem{remark}[theorem]{Remark}%
}

\lstdefinestyle{scala}{
  language=Scala,
  basicstyle=\ttfamily\small,
  breaklines=true,
  showstringspaces=false,
  columns=flexible,
}

\lstdefinestyle{haskell}{
  language=Haskell,
  basicstyle=\ttfamily\small,
  breaklines=true,
  showstringspaces=false,
  columns=flexible,
}

\lstdefinestyle{python}{
  language=Python,
  basicstyle=\ttfamily\small,
  breaklines=true,
  showstringspaces=false,
  columns=flexible,
  commentstyle=\upshape\color[gray]{0.5},
  keywordstyle=\upshape,
  stringstyle=\upshape,
}

\lstdefinestyle{nix}{
  language=,
  basicstyle=\ttfamily\small,
  breaklines=true,
  showstringspaces=false,
  columns=flexible,
  commentstyle=\color[gray]{0.5},
  morecomment=[l]{\#},
}

\newcommand{\Path}{\mathtt{Path}}
\newcommand{\Label}{\mathtt{Label}}
\newcommand{\dom}{\mathtt{dom}}
\newcommand{\properties}{\mathtt{properties}}
\newcommand{\overrides}{\mathtt{overrides}}
\newcommand{\supers}{\mathtt{supers}}
\newcommand{\inherits}{\mathtt{inherits}}
\newcommand{\defines}{\mathtt{defines}}
\newcommand{\Node}{\mathtt{Node}}
\newcommand{\defs}{\mathtt{defs}}
\newcommand{\refs}{\mathtt{refs}}
\newcommand{\ASTroot}{\mathtt{root}}
\newcommand{\bases}{\mathtt{bases}}
\newcommand{\resolve}{\mathtt{resolve}}
\newcommand{\this}{\mathtt{this}}
\newcommand{\snoc}{\mathbin{\triangleright}}
\newcommand{\init}{\mathtt{init}}
\newcommand{\last}{\mathtt{last}}
\newcommand{\qualifiedthis}[1]{#1.\mathbf{this}}
\newcommand{\dbi}[1]{\uparrow^{\!#1}}


% --- conditional-compilation switches -------------------------------
% Entry points choose which parts of this shared file to compile by
% setting these booleans before \AtBeginDocument{%
  \providecommand\BibTeX{{Bib\TeX}}
}

\usepackage{mathtools}
\usepackage{mathrsfs}
\usepackage{listings}
\usepackage{multirow}
\usepackage{graphicx}
%% \iflatexml is predefined as true by LaTeXML (the engine arXiv uses to
%% render the HTML version). Under pdfLaTeX it is undefined, so we declare it
%% here defaulting to false; LaTeXML preserves its own true value and ignores
%% this reset. This lets us ship a PDF-only path (absolute page positions via
%% zref-savepos) alongside an HTML fallback that needs no page coordinates.
\newif\iflatexml
\latexmlfalse
\iflatexml\else
  \usepackage{zref-savepos}
\fi
\usepackage{tikz}
\usetikzlibrary{positioning,arrows.meta,calc}

%% Allow multi-line display math (align, etc.) to break across pages, so
%% the long BNF grammars are not trapped as single unbreakable blocks.
\allowdisplaybreaks

%% Right-aligned annotations for a multi-line math display.
%% Nested \left\{...\right\} groups cannot share a common alignment column,
%% so each annotation is pinned to a single right edge via absolute page
%% positions recorded by zref-savepos (resolved on the second compiler pass).
%% \figureAnnotationAnchor records the display's left edge (call once just
%% inside the math, before the outermost \begin{aligned}); \figureAnnotation
%% lays out a zero-width box whose right end sits at (left edge + \linewidth),
%% i.e. the enclosing minipage's right margin, regardless of nesting depth.
\newcounter{figureAnnotationCounter}
\iflatexml
  %% HTML fallback (LaTeXML / arXiv HTML): page-absolute positioning via
  %% zref-savepos is unavailable, and the MathML output flows rather than being
  %% laid out at fixed page coordinates, so a pixel-accurate right edge cannot
  %% be reproduced. Instead, render each annotation as a small inline label that
  %% trails the line it describes; the browser wraps it naturally. The anchor
  %% becomes a no-op since no left edge needs to be recorded.
  \newcommand{\figureAnnotationAnchor}{}
  \newcommand{\figureAnnotation}[1]{%
    \quad\normalfont\text{\scriptsize(#1)}%
  }
\else
  \newcommand{\figureAnnotationAnchor}{\zsavepos{figureAnnotationLeft}}
  \newcommand{\figureAnnotation}[1]{%
    \stepcounter{figureAnnotationCounter}%
    \zsavepos{figureAnnotation\the\value{figureAnnotationCounter}}%
    \rlap{\hbox to
      \dimexpr\zposx{figureAnnotationLeft}sp+\linewidth-\zposx{figureAnnotation\the\value{figureAnnotationCounter}}sp\relax
    {\hfill\normalfont\text{\scriptsize #1}}}%
  }
\fi

\AtBeginDocument{%
  \theoremstyle{acmdefinition}%
  \newtheorem{remark}[theorem]{Remark}%
}

\lstdefinestyle{scala}{
  language=Scala,
  basicstyle=\ttfamily\small,
  breaklines=true,
  showstringspaces=false,
  columns=flexible,
}

\lstdefinestyle{haskell}{
  language=Haskell,
  basicstyle=\ttfamily\small,
  breaklines=true,
  showstringspaces=false,
  columns=flexible,
}

\lstdefinestyle{python}{
  language=Python,
  basicstyle=\ttfamily\small,
  breaklines=true,
  showstringspaces=false,
  columns=flexible,
  commentstyle=\upshape\color[gray]{0.5},
  keywordstyle=\upshape,
  stringstyle=\upshape,
}

\lstdefinestyle{nix}{
  language=,
  basicstyle=\ttfamily\small,
  breaklines=true,
  showstringspaces=false,
  columns=flexible,
  commentstyle=\color[gray]{0.5},
  morecomment=[l]{\#},
}

\newcommand{\Path}{\mathtt{Path}}
\newcommand{\Label}{\mathtt{Label}}
\newcommand{\dom}{\mathtt{dom}}
\newcommand{\properties}{\mathtt{properties}}
\newcommand{\overrides}{\mathtt{overrides}}
\newcommand{\supers}{\mathtt{supers}}
\newcommand{\inherits}{\mathtt{inherits}}
\newcommand{\defines}{\mathtt{defines}}
\newcommand{\Node}{\mathtt{Node}}
\newcommand{\defs}{\mathtt{defs}}
\newcommand{\refs}{\mathtt{refs}}
\newcommand{\ASTroot}{\mathtt{root}}
\newcommand{\bases}{\mathtt{bases}}
\newcommand{\resolve}{\mathtt{resolve}}
\newcommand{\this}{\mathtt{this}}
\newcommand{\snoc}{\mathbin{\triangleright}}
\newcommand{\init}{\mathtt{init}}
\newcommand{\last}{\mathtt{last}}
\newcommand{\qualifiedthis}[1]{#1.\mathbf{this}}
\newcommand{\dbi}[1]{\uparrow^{\!#1}}


% --- conditional-compilation switches -------------------------------
% Entry points choose which parts of this shared file to compile by
% setting these booleans before \AtBeginDocument{%
  \providecommand\BibTeX{{Bib\TeX}}
}

\usepackage{mathtools}
\usepackage{mathrsfs}
\usepackage{listings}
\usepackage{multirow}
\usepackage{graphicx}
%% \iflatexml is predefined as true by LaTeXML (the engine arXiv uses to
%% render the HTML version). Under pdfLaTeX it is undefined, so we declare it
%% here defaulting to false; LaTeXML preserves its own true value and ignores
%% this reset. This lets us ship a PDF-only path (absolute page positions via
%% zref-savepos) alongside an HTML fallback that needs no page coordinates.
\newif\iflatexml
\latexmlfalse
\iflatexml\else
  \usepackage{zref-savepos}
\fi
\usepackage{tikz}
\usetikzlibrary{positioning,arrows.meta,calc}

%% Allow multi-line display math (align, etc.) to break across pages, so
%% the long BNF grammars are not trapped as single unbreakable blocks.
\allowdisplaybreaks

%% Right-aligned annotations for a multi-line math display.
%% Nested \left\{...\right\} groups cannot share a common alignment column,
%% so each annotation is pinned to a single right edge via absolute page
%% positions recorded by zref-savepos (resolved on the second compiler pass).
%% \figureAnnotationAnchor records the display's left edge (call once just
%% inside the math, before the outermost \begin{aligned}); \figureAnnotation
%% lays out a zero-width box whose right end sits at (left edge + \linewidth),
%% i.e. the enclosing minipage's right margin, regardless of nesting depth.
\newcounter{figureAnnotationCounter}
\iflatexml
  %% HTML fallback (LaTeXML / arXiv HTML): page-absolute positioning via
  %% zref-savepos is unavailable, and the MathML output flows rather than being
  %% laid out at fixed page coordinates, so a pixel-accurate right edge cannot
  %% be reproduced. Instead, render each annotation as a small inline label that
  %% trails the line it describes; the browser wraps it naturally. The anchor
  %% becomes a no-op since no left edge needs to be recorded.
  \newcommand{\figureAnnotationAnchor}{}
  \newcommand{\figureAnnotation}[1]{%
    \quad\normalfont\text{\scriptsize(#1)}%
  }
\else
  \newcommand{\figureAnnotationAnchor}{\zsavepos{figureAnnotationLeft}}
  \newcommand{\figureAnnotation}[1]{%
    \stepcounter{figureAnnotationCounter}%
    \zsavepos{figureAnnotation\the\value{figureAnnotationCounter}}%
    \rlap{\hbox to
      \dimexpr\zposx{figureAnnotationLeft}sp+\linewidth-\zposx{figureAnnotation\the\value{figureAnnotationCounter}}sp\relax
    {\hfill\normalfont\text{\scriptsize #1}}}%
  }
\fi

\AtBeginDocument{%
  \theoremstyle{acmdefinition}%
  \newtheorem{remark}[theorem]{Remark}%
}

\lstdefinestyle{scala}{
  language=Scala,
  basicstyle=\ttfamily\small,
  breaklines=true,
  showstringspaces=false,
  columns=flexible,
}

\lstdefinestyle{haskell}{
  language=Haskell,
  basicstyle=\ttfamily\small,
  breaklines=true,
  showstringspaces=false,
  columns=flexible,
}

\lstdefinestyle{python}{
  language=Python,
  basicstyle=\ttfamily\small,
  breaklines=true,
  showstringspaces=false,
  columns=flexible,
  commentstyle=\upshape\color[gray]{0.5},
  keywordstyle=\upshape,
  stringstyle=\upshape,
}

\lstdefinestyle{nix}{
  language=,
  basicstyle=\ttfamily\small,
  breaklines=true,
  showstringspaces=false,
  columns=flexible,
  commentstyle=\color[gray]{0.5},
  morecomment=[l]{\#},
}

\newcommand{\Path}{\mathtt{Path}}
\newcommand{\Label}{\mathtt{Label}}
\newcommand{\dom}{\mathtt{dom}}
\newcommand{\properties}{\mathtt{properties}}
\newcommand{\overrides}{\mathtt{overrides}}
\newcommand{\supers}{\mathtt{supers}}
\newcommand{\inherits}{\mathtt{inherits}}
\newcommand{\defines}{\mathtt{defines}}
\newcommand{\Node}{\mathtt{Node}}
\newcommand{\defs}{\mathtt{defs}}
\newcommand{\refs}{\mathtt{refs}}
\newcommand{\ASTroot}{\mathtt{root}}
\newcommand{\bases}{\mathtt{bases}}
\newcommand{\resolve}{\mathtt{resolve}}
\newcommand{\this}{\mathtt{this}}
\newcommand{\snoc}{\mathbin{\triangleright}}
\newcommand{\init}{\mathtt{init}}
\newcommand{\last}{\mathtt{last}}
\newcommand{\qualifiedthis}[1]{#1.\mathbf{this}}
\newcommand{\dbi}[1]{\uparrow^{\!#1}}


% --- conditional-compilation switches -------------------------------
% Entry points choose which parts of this shared file to compile by
% setting these booleans before \AtBeginDocument{%
  \providecommand\BibTeX{{Bib\TeX}}
}

\usepackage{mathtools}
\usepackage{mathrsfs}
\usepackage{listings}
\usepackage{multirow}
\usepackage{graphicx}
%% \iflatexml is predefined as true by LaTeXML (the engine arXiv uses to
%% render the HTML version). Under pdfLaTeX it is undefined, so we declare it
%% here defaulting to false; LaTeXML preserves its own true value and ignores
%% this reset. This lets us ship a PDF-only path (absolute page positions via
%% zref-savepos) alongside an HTML fallback that needs no page coordinates.
\newif\iflatexml
\latexmlfalse
\iflatexml\else
  \usepackage{zref-savepos}
\fi
\usepackage{tikz}
\usetikzlibrary{positioning,arrows.meta,calc}

%% Allow multi-line display math (align, etc.) to break across pages, so
%% the long BNF grammars are not trapped as single unbreakable blocks.
\allowdisplaybreaks

%% Right-aligned annotations for a multi-line math display.
%% Nested \left\{...\right\} groups cannot share a common alignment column,
%% so each annotation is pinned to a single right edge via absolute page
%% positions recorded by zref-savepos (resolved on the second compiler pass).
%% \figureAnnotationAnchor records the display's left edge (call once just
%% inside the math, before the outermost \begin{aligned}); \figureAnnotation
%% lays out a zero-width box whose right end sits at (left edge + \linewidth),
%% i.e. the enclosing minipage's right margin, regardless of nesting depth.
\newcounter{figureAnnotationCounter}
\iflatexml
  %% HTML fallback (LaTeXML / arXiv HTML): page-absolute positioning via
  %% zref-savepos is unavailable, and the MathML output flows rather than being
  %% laid out at fixed page coordinates, so a pixel-accurate right edge cannot
  %% be reproduced. Instead, render each annotation as a small inline label that
  %% trails the line it describes; the browser wraps it naturally. The anchor
  %% becomes a no-op since no left edge needs to be recorded.
  \newcommand{\figureAnnotationAnchor}{}
  \newcommand{\figureAnnotation}[1]{%
    \quad\normalfont\text{\scriptsize(#1)}%
  }
\else
  \newcommand{\figureAnnotationAnchor}{\zsavepos{figureAnnotationLeft}}
  \newcommand{\figureAnnotation}[1]{%
    \stepcounter{figureAnnotationCounter}%
    \zsavepos{figureAnnotation\the\value{figureAnnotationCounter}}%
    \rlap{\hbox to
      \dimexpr\zposx{figureAnnotationLeft}sp+\linewidth-\zposx{figureAnnotation\the\value{figureAnnotationCounter}}sp\relax
    {\hfill\normalfont\text{\scriptsize #1}}}%
  }
\fi

\AtBeginDocument{%
  \theoremstyle{acmdefinition}%
  \newtheorem{remark}[theorem]{Remark}%
}

\lstdefinestyle{scala}{
  language=Scala,
  basicstyle=\ttfamily\small,
  breaklines=true,
  showstringspaces=false,
  columns=flexible,
}

\lstdefinestyle{haskell}{
  language=Haskell,
  basicstyle=\ttfamily\small,
  breaklines=true,
  showstringspaces=false,
  columns=flexible,
}

\lstdefinestyle{python}{
  language=Python,
  basicstyle=\ttfamily\small,
  breaklines=true,
  showstringspaces=false,
  columns=flexible,
  commentstyle=\upshape\color[gray]{0.5},
  keywordstyle=\upshape,
  stringstyle=\upshape,
}

\lstdefinestyle{nix}{
  language=,
  basicstyle=\ttfamily\small,
  breaklines=true,
  showstringspaces=false,
  columns=flexible,
  commentstyle=\color[gray]{0.5},
  morecomment=[l]{\#},
}

\newcommand{\Path}{\mathtt{Path}}
\newcommand{\Label}{\mathtt{Label}}
\newcommand{\dom}{\mathtt{dom}}
\newcommand{\properties}{\mathtt{properties}}
\newcommand{\overrides}{\mathtt{overrides}}
\newcommand{\supers}{\mathtt{supers}}
\newcommand{\inherits}{\mathtt{inherits}}
\newcommand{\defines}{\mathtt{defines}}
\newcommand{\Node}{\mathtt{Node}}
\newcommand{\defs}{\mathtt{defs}}
\newcommand{\refs}{\mathtt{refs}}
\newcommand{\ASTroot}{\mathtt{root}}
\newcommand{\bases}{\mathtt{bases}}
\newcommand{\resolve}{\mathtt{resolve}}
\newcommand{\this}{\mathtt{this}}
\newcommand{\snoc}{\mathbin{\triangleright}}
\newcommand{\init}{\mathtt{init}}
\newcommand{\last}{\mathtt{last}}
\newcommand{\qualifiedthis}[1]{#1.\mathbf{this}}
\newcommand{\dbi}[1]{\uparrow^{\!#1}}


% --- conditional-compilation switches -------------------------------
% Entry points choose which parts of this shared file to compile by
% setting these booleans before \input{inheritance-calculus}.  The
% defaults include everything, so the full build (preprint.tex) sets
% nothing; submission.tex clears the appendices (POPL 2027 requires them
% as separate supplemental material), supplement.tex clears the body.
% The single bibliography is emitted exactly once, after the main matter and
% before the appendix, since both cite it. In the appendix-only build (no main
% matter) it instead follows the appendix; see the two guarded copies below.
\makeatletter
\@ifundefined{ifInheritanceBody}{\newif\ifInheritanceBody\InheritanceBodytrue}{}
\@ifundefined{ifInheritanceAppendix}{\newif\ifInheritanceAppendix\InheritanceAppendixtrue}{}
% Language switch for the reader-facing prose. Defaults to false (English), so the
% existing pdfLaTeX entries are unaffected; preprint-zh.tex sets it true and adds
% CJK support, still under pdfLaTeX. See the paper-writing skill (Bilingual
% Translation) for the convention.
\@ifundefined{ifInheritanceChinese}{\newif\ifInheritanceChinese\InheritanceChinesefalse}{}
% acmart writes \newlabel{tocindent<n>}{<dimen>} into the .aux to align the
% table of contents. Those bare-dimen labels break hyperref's \newlabel when
% another build imports this .aux via xr-hyper; the importing entries
% (submission.tex, supplement.tex) drop them at import time through
% xr-tocindent-filter.tex, so nothing is suppressed here at the source.
\makeatother

% --- bilingual reader-facing prose ----------------------------------
% The reader-facing body exists in two languages in this one file, selected by
% \ifInheritanceChinese: English (default) and a Chinese translation. Both build
% with pdfLaTeX, so the math renders identically. Block prose is guarded by
% \ifInheritanceChinese <chinese> \else <english> \fi; inline switches (title,
% captions, theorem notes interleaved with math) use \bilingual{<en>}{<zh>}.
% Author-side notes stay Chinese % comments in both builds.
\newcommand{\bilingual}[2]{#1}
\ifInheritanceChinese
  % CJKutf8 + the Arphic gbsn font render Chinese under pdfLaTeX (see
  % modules/texlive.nix). The document is wrapped in a CJK* environment opened
  % after \begin{document}; ASCII passes through, so English and math render as in
  % the English build.
  \usepackage{CJKutf8}
  \setlength{\emergencystretch}{2em}
  % Display the Chinese, but feed the English to hyperref's PDF strings, which
  % cannot encode the CJK active characters.
  \renewcommand{\bilingual}[2]{\texorpdfstring{#2}{#1}}
\fi

\begin{document}
% Open CJK* before \title so the Chinese in the title, abstract, and body is read
% inside it; \proofname and the theorem-environment names are localized here
% (their bytes must sit inside a CJK environment, and translating \thmname's
% argument leaves the shared counter and numbering identical to English). Closed
% before \end{document}.
\ifInheritanceChinese
  \begin{CJK*}{UTF8}{gbsn}
  \renewcommand{\proofname}{证明}
  \makeatletter
  \newcommand{\inheritanceThmName}[1]{\@ifundefined{inheritance@thm@#1}{#1}{\csname inheritance@thm@#1\endcsname}}
  \renewcommand{\thmname}[1]{\inheritanceThmName{#1}}
  \@namedef{inheritance@thm@Theorem}{定理}
  \@namedef{inheritance@thm@Lemma}{引理}
  \@namedef{inheritance@thm@Corollary}{推论}
  \@namedef{inheritance@thm@Proposition}{命题}
  \@namedef{inheritance@thm@Conjecture}{猜想}
  \@namedef{inheritance@thm@Definition}{定义}
  \@namedef{inheritance@thm@Example}{例}
  \@namedef{inheritance@thm@Remark}{注}
  \@namedef{inheritance@thm@Notation}{记法}
  \makeatother
\fi

% The optional short title feeds the running head and the PDF metadata; keep it
% ASCII English in both builds, since those are typeset outside CJK*.
\title[A Calculus of Inheritance]{\bilingual{A Calculus of Inheritance}{一种继承的演算}}
% The appendix-only build (supplement.tex) is distributed as a standalone
% artifact; mark it so it is not mistaken for the main paper, which carries
% no subtitle.
\ifInheritanceBody\else
  \subtitle{Supplementary Material}
\fi

\author{Bo Yang}
\affiliation{%
  \institution{Figure AI Inc.}
  \city{San Jose}
  \state{California}
  \country{USA}
}
\email{yang-bo@yang-bo.com}
\thanks{This work was conducted independently prior to the author's employment at Figure AI.}

% The abstract, CCS concepts, and keywords are main-matter front matter: they
% belong to the paper, not to the standalone appendix PDF (supplement.tex),
% which sets \ifInheritanceBody false. Without this guard \maketitle typesets
% them in the appendix-only build.
\ifInheritanceBody
\begin{abstract}
  Just as the $\lambda$-calculus uses three primitives (abstraction, application, variable) as the foundation of functional programming, inheritance-calculus uses three primitives (record, definition, inheritance) as the foundation of declarative programming. A record is a set of definitions; by unifying modules, classes, objects, methods, fields, and locals under this single record abstraction, the calculus models inheritance simply as set union of definitions, two definitions of the same name merging recursively rather than one overriding the other. A merge is always meaningful because every value is a record and observation is purely structural: the merged record simply answers queries with the union of what its parts provide, so there is no scalar leaf at which two values could conflict. Consequently, composition is inherently commutative, idempotent, and associative, structurally eliminating the multiple-inheritance linearization problem. Its semantics is first-order, denotational, and computable by tabling (memoized fixpoint iteration as in tabled logic programming), even for cyclic inheritance hierarchies. This first-order, denotational, tabled semantics extends to the $\lambda$-calculus: $\lambda$-terms translate into a sublanguage of a lazy approximation of inheritance-calculus, the approximation that treats cyclic self-reference as divergent instead of iterating it to a fixpoint, and the translation is fully abstract, equating exactly the $\lambda$-terms with equal B\"ohm trees.

  Inheritance-calculus is distilled from MIXINv2, a practical implementation in which the same code can be invoked under different calling conventions (function colors such as synchronous and asynchronous); ordinary arithmetic, encoded within the calculus itself, behaves as multi-directional relations in the sense of logic programming; $\mathtt{this}$ resolves to multiple enclosing targets at once rather than a single receiver; and programs are extensible in both dimensions of the Expression Problem, gaining new variants and new operations by adding code. This makes inheritance-calculus strictly more expressive than the $\lambda$-calculus in Felleisen's sense of macro-expressibility.
\end{abstract}

\begin{CCSXML}
  <ccs2012>
  <concept>
  <concept_id>10003752.10010124.10010131</concept_id>
  <concept_desc>Theory of computation~Program semantics</concept_desc>
  <concept_significance>500</concept_significance>
  </concept>
  <concept>
  <concept_id>10003752.10010124.10010125.10010128</concept_id>
  <concept_desc>Theory of computation~Object oriented constructs</concept_desc>
  <concept_significance>500</concept_significance>
  </concept>
  <concept>
  <concept_id>10003752.10010124.10010125.10010127</concept_id>
  <concept_desc>Theory of computation~Functional constructs</concept_desc>
  <concept_significance>300</concept_significance>
  </concept>
  <concept>
  <concept_id>10003752.10010124.10010125.10010129</concept_id>
  <concept_desc>Theory of computation~Program schemes</concept_desc>
  <concept_significance>100</concept_significance>
  </concept>
  <concept>
  <concept_id>10011007.10011006.10011039</concept_id>
  <concept_desc>Software and its engineering~Formal language definitions</concept_desc>
  <concept_significance>500</concept_significance>
  </concept>
  <concept>
  <concept_id>10011007.10011006.10011008.10011009.10011019</concept_id>
  <concept_desc>Software and its engineering~Extensible languages</concept_desc>
  <concept_significance>500</concept_significance>
  </concept>
  <concept>
  <concept_id>10011007.10011006.10011008.10011009.10011011</concept_id>
  <concept_desc>Software and its engineering~Object oriented languages</concept_desc>
  <concept_significance>100</concept_significance>
  </concept>
  </ccs2012>
\end{CCSXML}

\ccsdesc[500]{Theory of computation~Program semantics}
\ccsdesc[500]{Theory of computation~Object oriented constructs}
\ccsdesc[300]{Theory of computation~Functional constructs}
\ccsdesc[100]{Theory of computation~Program schemes}
\ccsdesc[500]{Software and its engineering~Formal language definitions}
\ccsdesc[500]{Software and its engineering~Extensible languages}
\ccsdesc[100]{Software and its engineering~Object oriented languages}

\keywords{declarative programming, inheritance, fixpoint semantics,
  self-referential records, expression problem,
  first-order semantics,
  \texorpdfstring{$\lambda$-calculus}{lambda-calculus},
\texorpdfstring{B\"ohm trees}{Bohm trees}, configuration languages}
\fi % end \ifInheritanceBody (abstract, CCS concepts, keywords are main-matter only)

\maketitle

\ifInheritanceBody
\section{\bilingual{Introduction}{引言}}
\label{sec:introduction}

\ifInheritanceChinese
诸如 Jsonnet~\cite{jsonnet}、Hydra~\cite{hydra2023}、CUE~\cite{cue2019}、Dhall~\cite{dhall2017}、Kustomize~\cite{kustomize} 与 JSON Patch~\cite{rfc6902-json-patch} 这样的配置语言,都提供了通过继承或合并来组合结构化数据的机制。

在这些之中,Nix 生态系统~\cite{dolstra2010-nixos} 格外突出,无论是就其规模还是其设计的走向而言。它的软件包集合 \texttt{nixpkgs} 收录了超过 113{,}000 个软件包,%
\footnote{
  据 Repology~\cite{repology2025},这接近 Debian 或 Ubuntu 软件包数量的三倍,而贡献者数量相近。并非每个软件包都用到这里所描述的可扩展性机制,但同样的告诫也适用于其他发行版。
}
是现存最大的软件包集合之一。软件包的组合最初依赖一条线性化的 mixin 链,其中每一层都能看见它的前一层,%
\footnote{
  Nix overlay 通过一条线性化的 \texttt{self}/\texttt{super} 链来组合软件包集合,遵循 Bracha--Cook 的 mixin 模式~\cite{bracha1990-mixin-inheritance}:\texttt{self} 晚绑定到不动点结果,而 \texttt{super} 指向前一层。该机制是不对称且依赖顺序的。
}
但该生态系统已日益汇聚到一种对称的替代方案上:NixOS 模块系统~\cite{dolstra2010-nixos,nixosmodules},它通过对嵌套的属性集进行递归合并、配以自引用求值与延迟模块~\cite{nixosmodules} 来组合,而这种组合是可交换、幂等且可结合的。模块系统最初为系统配置而设计,如今已被用于用户环境~\cite{homemanager}、macOS 配置~\cite{nixdarwin}、磁盘分区~\cite{disko}、flake 结构~\cite{flakeparts}、开发者环境~\cite{devenv}、Kubernetes 集群~\cite{kubenix,nixidy},以及多语言软件包构建~\cite{dream2nix},而最后这一项正是 \texttt{nixpkgs} 那套不对称机制原本的领域。Nix 是一门带头等函数的函数式语言;然而一份典型的模块配置几乎完全由嵌套的属性集构成,函数只是偶尔作为模块系统的 DSL 语法出现,例如 \texttt{mkIf} 与 \texttt{mkMerge}。没有任何计算理论刻画了这一机制所展现的可扩展性。
\else
Configuration languages such as Jsonnet~\cite{jsonnet}, Hydra~\cite{hydra2023},
CUE~\cite{cue2019}, Dhall~\cite{dhall2017},
Kustomize~\cite{kustomize}, and JSON
Patch~\cite{rfc6902-json-patch} all provide mechanisms for
composing structured data through inheritance or merging.

Among these, the Nix ecosystem~\cite{dolstra2010-nixos} stands out,
both for scale and for the trajectory of its design.
Its package collection, \texttt{nixpkgs}, contains over
113{,}000 packages,%
\footnote{
  According to Repology~\cite{repology2025},
  nearly three times the package count of Debian or Ubuntu,
  achieved with a similar number of contributors.
  Not every package exercises the extensibility mechanisms
  described here, but the same caveat applies to other distributions.
}
one of the largest in existence.
Package composition originally relied on a linearized mixin
chain in which each layer sees its predecessor,%
\footnote{
  Nix overlays compose package sets via a linearized
  \texttt{self}/\texttt{super} chain following the
  Bracha--Cook mixin pattern~\cite{bracha1990-mixin-inheritance}:
  \texttt{self} is late-bound to the fixed-point result
  and \texttt{super} refers to the previous layer.
  The mechanism is asymmetric and order-dependent.
}
but the ecosystem has increasingly converged on a symmetric
alternative: the NixOS module
system~\cite{dolstra2010-nixos,nixosmodules}, which composes by
recursively merging nested attribute sets with self-referential
evaluation and deferred modules~\cite{nixosmodules}---commutative,
idempotent,
and associative.
Originally designed for system configuration, the module system
has been adopted for user environments~\cite{homemanager},
macOS configuration~\cite{nixdarwin},
disk partitioning~\cite{disko},
flake structure~\cite{flakeparts},
developer environments~\cite{devenv},
Kubernetes clusters~\cite{kubenix,nixidy},
and multi-language package building~\cite{dream2nix}---the
original domain of \texttt{nixpkgs}'s asymmetric mechanism.
Nix is a functional language with first-class functions;
yet a typical module configuration consists almost entirely of
nested attribute sets, with functions appearing only
occasionally as DSL syntax of the module system
such as \texttt{mkIf} and \texttt{mkMerge}.
No computational theory captures the extensibility that this
mechanism exhibits.
\fi

The $\lambda$-calculus has no analogue for declarative
programming. This gap matters because the two paradigms differ in
fundamental ways, which raises a question:

\begin{quote}
  Can configuration alone be practical general-purpose programming?
\end{quote}

\noindent
To make the question precise, we operationalize it with four
\emph{opinionated} proxies, grouped under the question's two halves:
\emph{configuration} and \emph{practical general-purpose} programming.

\begin{itemize}
  \item Immutable and first-order%
    \footnote{
      Here \emph{first-order} holds in three senses at once:
      the semantic domain contains only data, no function spaces,
      closures, or
      strategies~\cite{vanemden1976-predicate-logic-semantics};
      the defining equations map data to
      data~\cite{reynolds1972-definitional-interpreters};
      and the semantics is a first-order inductive
      definition~\cite{aczel1977-inductive-definitions} whose
      quantifiers range only over data.
    }
    is our proxy for
    \emph{configuration}: a configuration is read, not a function that
    reads an argument.
  \item Turing-completeness together with extensibility is our proxy for
    \emph{practical general-purpose} programming: a language must be able
    to compute anything and to grow by adding code rather than rewriting
    it, the Expression
    Problem~\cite{wadler1998-expression-problem}.
\end{itemize}

\noindent
Each classical model satisfies some of these proxies but fails at least
one:

\begin{table}[h]
  \centering
  \small
  \begin{tabular}{lcccl}
    \textbf{Model} & \textbf{Immutable} & \textbf{First-order} & \textbf{Turing complete} & \textbf{Extensibility} \\
    \hline
    Turing Machine        & $\times$     & $\checkmark$ & $\checkmark$ & $\times$ \\
    $\lambda$-calculus    & $\checkmark$ & $\times$     & $\checkmark$ & opt-in \\
    RAM Machine           & $\times$     & $\checkmark$ & $\checkmark$ & $\times$ \\
    Prolog / Horn clauses & $\checkmark$ & $\checkmark$ & $\checkmark$ & disjunct-only \\
  \end{tabular}
  \caption{The question operationalized as four proxies: immutable and
    first-order (configuration) versus Turing-complete and extensible
  (practical general-purpose). No classical model satisfies all four.}
  \label{tab:computational-models}
\end{table}

\noindent
The Turing Machine and RAM Machine require mutable state and have no native
extension mechanism;
the $\lambda$-calculus is higher-order rather than first-order, and its
extensibility is opt-in: a function is extensible only where its author
anticipated extension, for example by threading open-recursion
parameters, while a closed function body cannot be extended without
rewriting it;
Prolog meets the other three proxies but is extensible only by adding
disjuncts (alternative clauses) to existing predicates, not by attaching
new fields or methods to existing values.
No known computational model satisfies all four proxies simultaneously.

The Nix ecosystem already suggests that such a model may exist:
its inheritance-based composition over recursive records, without
explicit functions, is expressive enough to configure entire
operating systems and manage one of the largest package
collections in existence.
Most configuration languages operate on finite, well-founded
structures (initial algebras in the sense of universal algebra).
Nix is an exception: its values are lazily observable
and possibly infinite, allowing any single package or
configuration option to be evaluated without materializing the
entire set of hundreds of thousands of packages and options.
Guided by this observation, we set out to reduce the
NixOS module system to a minimal set of primitives.
The reduction produced MIXIN, an early implementation in Nix,
and subsequently
MIXINv2\anon[~(included as supplementary material)]{~\cite{mixin2025}}, now implemented in Python: a declarative programming
language with only three constructs (record, definition,
and inheritance) and no functions or let-bindings.
MIXINv2 obtains scalars such as floating-point numbers through a
foreign-function interface; inheritance-calculus is what remains after
removing that interface, so names like $\mathrm{PI}$ and
$\mathrm{Float}$ in Figure~\ref{fig:cheatsheet} denote imported
records whose internals the examples never inspect: scalar
computation such as the figure's additions happens behind the
interface, while the calculus itself only ever observes structure;
the case study of Section~\ref{sec:case-study} instead encodes
naturals and their arithmetic purely in the calculus, with no FFI.
When two merged records define the same name with different contents,
the result is simply a record carrying both structures; the calculus
never compares or rejects them
(Section~\ref{sec:mixin-trees} makes the merge precise).
No such collision occurs in Figure~\ref{fig:cheatsheet}: each
instance defines each scalar field exactly once, and the type record
it merges with contributes structure, not a competing value.

\begin{figure*}
  \centering
  \begin{minipage}[t]{0.34\textwidth}
    \textbf{Python}\par
\begin{lstlisting}[style=python]
# complex_module.py
from builtins import *
from math import *
from operator import *
from dataclasses import *
@dataclass(kw_only=True)
class Complex:
    re: float
    im: float
    def plus(self, *,
             other: "Complex"):
        sumRe = add(
            self.re,
            other.re)
        sumIm = add(
            self.im,
            other.im)
        return Complex(
            re=sumRe,
            im=sumIm)
a = Complex(
    re=pi,
    im=e)
b = Complex(
    re=tau,
    im=pi)
total = a.plus(
    other=b)
\end{lstlisting}
  \end{minipage}
  \hfill
  \begin{minipage}[t]{0.65\textwidth}
    \textbf{Inheritance-calculus}\par\smallskip
    {\small\(
        \figureAnnotationAnchor
        \begin{aligned}
          &\mathrm{ComplexModule} \mapsto {} \figureAnnotation{module} \\
          &\left\{
            \begin{aligned}
              &\mathrm{BuiltIns}, \figureAnnotation{import} \\
              &\mathrm{Math}, \figureAnnotation{import} \\
              &\mathrm{Operator}, \figureAnnotation{import} \\
              &\mathrm{Complex} \mapsto {} \figureAnnotation{class} \\
              &\left\{
                \begin{aligned}
                  &\mathrm{re} \mapsto {} \figureAnnotation{field} \\
                  &\left\{ \qualifiedthis{\mathrm{ComplexModule}}.\mathrm{Float} \right\}, \figureAnnotation{imported type} \\
                  &\mathrm{im} \mapsto {} \figureAnnotation{field} \\
                  &\left\{ \qualifiedthis{\mathrm{ComplexModule}}.\mathrm{Float} \right\}, \figureAnnotation{imported type} \\
                  &\mathrm{Plus} \mapsto {} \figureAnnotation{method} \\
                  &\left\{
                    \begin{aligned}
                      &\mathrm{other} \mapsto \left\{ \mathrm{Complex} \right\}, \figureAnnotation{parameter} \\
                      &\mathrm{\_sumReCall} \mapsto {} \figureAnnotation{function application} \\
                      &\left\{
                        \begin{aligned}
                          &\qualifiedthis{\mathrm{ComplexModule}}.\mathrm{Add}, \figureAnnotation{imported function} \\
                          &\mathrm{operand0} \mapsto \left\{ \qualifiedthis{\mathrm{Complex}}.\mathrm{re} \right\}, \figureAnnotation{keyword argument} \\
                          &\mathrm{operand1} \mapsto \left\{ \qualifiedthis{\mathrm{Plus}}.\mathrm{other}.\mathrm{re} \right\} \figureAnnotation{keyword argument} \\
                      \end{aligned}\right\}, \\
                      &\mathrm{\_sumRe} \mapsto \left\{ \mathrm{\_sumReCall}.\mathrm{result} \right\}, \figureAnnotation{local} \\
                      &\mathrm{\_sumImCall} \mapsto {} \figureAnnotation{function application} \\
                      &\left\{
                        \begin{aligned}
                          &\qualifiedthis{\mathrm{ComplexModule}}.\mathrm{Add}, \figureAnnotation{imported function} \\
                          &\mathrm{operand0} \mapsto \left\{ \qualifiedthis{\mathrm{Complex}}.\mathrm{im} \right\}, \figureAnnotation{keyword argument} \\
                          &\mathrm{operand1} \mapsto \left\{ \qualifiedthis{\mathrm{Plus}}.\mathrm{other}.\mathrm{im} \right\} \figureAnnotation{keyword argument} \\
                      \end{aligned}\right\}, \\
                      &\mathrm{\_sumIm} \mapsto \left\{ \mathrm{\_sumImCall}.\mathrm{result} \right\}, \figureAnnotation{local} \\
                      &\mathrm{result} \mapsto {} \figureAnnotation{return} \\
                      &\left\{
                        \begin{aligned}
                          &\mathrm{Complex}, \figureAnnotation{new instance} \\
                          &\mathrm{re} \mapsto \left\{ \mathrm{\_sumRe} \right\}, \figureAnnotation{keyword argument} \\
                          &\mathrm{im} \mapsto \left\{ \mathrm{\_sumIm} \right\} \figureAnnotation{keyword argument} \\
                      \end{aligned}\right\} \\
                  \end{aligned}\right\} \\
              \end{aligned}\right\}, \\
              &\mathrm{a} \mapsto {} \figureAnnotation{module-level variable} \\
              &\left\{
                \begin{aligned}
                  &\mathrm{Complex}, \figureAnnotation{new instance} \\
                  &\mathrm{re} \mapsto \left\{ \qualifiedthis{\mathrm{ComplexModule}}.\mathrm{PI} \right\}, \figureAnnotation{imported constant} \\
                  &\mathrm{im} \mapsto \left\{ \qualifiedthis{\mathrm{ComplexModule}}.\mathrm{E} \right\} \figureAnnotation{imported constant} \\
              \end{aligned}\right\}, \\
              &\mathrm{b} \mapsto {} \figureAnnotation{module-level variable} \\
              &\left\{
                \begin{aligned}
                  &\mathrm{Complex}, \figureAnnotation{new instance} \\
                  &\mathrm{re} \mapsto \left\{ \qualifiedthis{\mathrm{ComplexModule}}.\mathrm{Tau} \right\}, \figureAnnotation{imported constant} \\
                  &\mathrm{im} \mapsto \left\{ \qualifiedthis{\mathrm{ComplexModule}}.\mathrm{PI} \right\} \figureAnnotation{imported constant} \\
              \end{aligned}\right\}, \\
              &\mathrm{\_totalCall} \mapsto {} \figureAnnotation{method application} \\
              &\left\{
                \begin{aligned}
                  &\mathrm{a}.\mathrm{Plus}, \figureAnnotation{method reference} \\
                  &\mathrm{other} \mapsto \left\{ \mathrm{b} \right\} \figureAnnotation{keyword argument} \\
              \end{aligned}\right\}, \\
              &\mathrm{total} \mapsto \left\{ \mathrm{\_totalCall}.\mathrm{result} \right\} \figureAnnotation{module-level variable} \\
          \end{aligned}\right\}
        \end{aligned}
    \)}
  \end{minipage}
  \caption{A first look at inheritance-calculus (right), beside Python
    (left): the same complex-number class with an addition method.
    An entry $name \mapsto \{\dots\}$ is a definition; a bare name
    inside braces (such as $\mathrm{BuiltIns}$ or $\mathrm{Complex}$) is
    an inheritance clause that copies that record's definitions into the
    enclosing record.
    $\mathrm{BuiltIns}$, $\mathrm{Math}$, and $\mathrm{Operator}$ are
    defined in an enclosing scope, elided here.
    A bare name resolves lexically to the innermost enclosing record
    that defines it directly; qualified $\this$ reaches members that
    arrive by inheritance, such as the imported $\mathrm{Float}$ and
    $\mathrm{Add}$.
    Types and values share the same record syntax: a field such as
    $\mathrm{re}$ records its type by inheriting the type's record,
    with no checking force in the untyped calculus.
    Both sides, typing their components with \texttt{builtins.float},
    drawing their constants from the \texttt{math} module, and adding with
    \texttt{operator.add}, compute the complex
    number $(\pi+\tau) + (e+\pi)i$.
    $\mathrm{\_sumReCall}$ and $\mathrm{\_sumImCall}$ are function
    applications, and $\mathrm{\_totalCall}$ a method application, each
    encoded as a command pattern: a record that inherits the callee and
    defines its arguments, the call's result read back from a
    $\mathrm{result}$ field that the callee defines.
    Idempotence only means that listing the same inheritance source
    twice in one record is a no-op; distinct records inheriting the
    same source, like the two call sites here or the two instances
    $\mathrm{a}$ and $\mathrm{b}$, remain distinct.}
  \label{fig:cheatsheet}
\end{figure*}



In developing real programs in MIXINv2, including a database-backed web server whose Python foreign-function interface only faithfully wraps individual host-library calls, we found that module, class, object, method, field, and local binding can all be expressed as the same construct: a record of definitions composed by inheritance. This paper distills this minimal computational model into \textbf{inheritance-calculus}, shown beside Python in Figure~\ref{fig:cheatsheet}.
The key to reading the figure is that a reference
$\qualifiedthis{X}.w$ behaves like a late-bound Java-style
\texttt{Outer.this}: it denotes the enclosing record labeled $X$ as
seen from wherever the reference itself ends up in the fully merged
tree.
Inheriting a record does not duplicate it; it makes the record's
definitions, references included, visible at the inheriting position
as well, and a reference resolves relative to the position from which
it is viewed.
The set-union picture and this late-binding picture are the same
mechanism seen globally and locally: the union says which definitions
a position carries, and late binding says how a definition behaves at
each position that carries it.
Qualified references therefore see members that arrive by merging,
such as the imported $\mathrm{Float}$, while bare names such as
$\mathrm{Complex}$ resolve lexically against directly written
definitions only.
This late binding is what makes inheritance act as function
application: the call record $\mathrm{\_totalCall}$ inherits
$\mathrm{Plus}$, so the merged tree presents
$\mathrm{Plus}$'s body inside $\mathrm{\_totalCall}$ as well, and
viewed from there the reference
$\qualifiedthis{\mathrm{Plus}}.\mathrm{other}$
denotes that copy's enclosing record, $\mathrm{\_totalCall}$ itself,
where the argument $\mathrm{other}$ is defined.
Receivers bind the same way: the dotted bare name
$\mathrm{a}.\mathrm{Plus}$ resolves its first label lexically and
projects the rest, denoting the copy of $\mathrm{Plus}$ inside the
instance $\mathrm{a}$, in which
$\qualifiedthis{\mathrm{Complex}}.\mathrm{re}$ denotes
$\mathrm{a}$'s own field.
Section~\ref{sec:syntax} gives its formal syntax, and
Section~\ref{sec:mixin-trees} expands this first look into the full
semantics.
The linearization problem of mixin-based
systems~\cite{bracha1990-mixin-inheritance,barrett1996-c3-linearization}
did not arise: inheritance is inherently commutative, idempotent,
and associative, and $\this$ naturally resolves to
\emph{multiple targets} rather than one: when inheritance places the
same definition inside several enclosing records, a reference through
$\this$ inherits from all of them at once, their contents merged,%
\footnote{
  Resolution is anchored at the position from which the reference is
  observed. When both $\mathrm{A}$ and $\mathrm{B}$ inherit
  $\mathrm{Base}$, a reference in $\mathrm{Base}$ observed inside
  $\mathrm{A}$ sees only $\mathrm{A}$; the two call sites of
  Figure~\ref{fig:cheatsheet} therefore stay independent. Multiple
  targets arise when a single observation position inherits the same
  definition through two routes, say $\mathrm{C}$ inheriting
  $\mathrm{Base}$ via both $\mathrm{A}$ and $\mathrm{B}$: observed at
  $\mathrm{C}$, a $\this$ reference in $\mathrm{Base}$ resolves to
  $\mathrm{A}$ and $\mathrm{B}$ at once, and projection through it
  yields the union of what they provide.
}
where
prior systems
such as Scala and the NixOS module system reject the multi-target
situation.
Section~\ref{sec:case-study} traces how boolean logic and
arithmetic, encoded purely in the calculus as separate files, extend in both
dimensions of the Expression
Problem~\cite{wadler1998-expression-problem} without any dedicated
mechanism, and how ordinary arithmetic yields the relational
semantics of logic
programming~\cite{vanemden1976-predicate-logic-semantics}.
Section~\ref{sec:emergent-phenomena} discusses these and further
emergent phenomena, including function color
blindness~\cite{nystrom2015-function-color}.
Together these sections answer the four proxies: inheritance-calculus
is immutable and first-order by construction
(Section~\ref{sec:definitions}), Turing-complete because the
$\lambda$-calculus embeds into it
(Section~\ref{sec:forward-translation}), and extensible in both
dimensions of the Expression Problem (Section~\ref{sec:case-study}).

\begin{figure*}
  \centering
  \begin{tikzpicture}[
      calc/.style={align=center, inner sep=2pt, font=\small},
      more/.style={-{Stealth[length=2mm]}, semithick},
      iso/.style={{Stealth[length=2mm]}-{Stealth[length=2mm]}, semithick},
      elbl/.style={font=\scriptsize\itshape, midway, fill=white, inner sep=1pt},
      x=1mm, y=1mm,
    ]
    % rows: bottom = lambda calculi, middle = lambda sublanguages, top = inheritance-calculi
    % columns: eager (x=0), lazy (x=42), fixpoint (x=84)
    \node[calc] (lam)  at (0,0)   {$\lambda$};
    \node[calc] (llam) at (42,0)  {lazy $\lambda$};
    \node[calc] (flam) at (84,0)  {fixpoint $\lambda$~\cite{yang2026-co-lambda}};

    \node[calc] (lsub) at (42,19)  {$\lambda$-lazy-IC (ours)};
    \node[calc] (sub)  at (84,19)  {$\lambda$-IC};

    \node[calc] (lic)  at (42,38)  {lazy IC (ours)};
    \node[calc] (ic)   at (84,38)  {IC (ours)};

    % horizontal edges (convergence): "more defined"
    \draw[more] (lam)  -- node[elbl,below]{more defined (non-strict)} (llam);
    \draw[more] (llam) -- node[elbl,below]{more defined (unproductive $\bot$)} (flam);
    \draw[more] (lsub) -- node[elbl,above]{more defined (powerset)} (sub);
    \draw[more] (lic)  -- node[elbl,above]{more defined (powerset)} (ic);

    % vertical edges (expressiveness): "strictly more expressive"
    \draw[more] (lsub) -- node[elbl,left]{strictly more expressive} (lic);
    \draw[more] (sub)  -- node[elbl,right]{strictly more expressive} (ic);

    % full abstraction: vertical, in the lazy column
    \draw[iso] (llam) -- node[elbl,left]{fully abstract} (lsub);
  \end{tikzpicture}
  \caption{The calculi of this paper; nodes marked
    (ours) are contributions of this paper.
    Abbreviations:
    $\lambda$ is the $\lambda$-calculus;
    the lazy prefix denotes the approximation that stops after the first
    iteration of the fixpoint of this paper's semantic equations
    (written $T_P\!\uparrow\!1$ in
    logic-programming notation): lookups that re-enter themselves
    through cyclic self-reference return divergence instead of being
    iterated further.
    Continuing:
    fixpoint $\lambda$ is the tabled evaluation of the
    $\lambda$-calculus of Yang~\cite{yang2026-co-lambda};
    IC the inheritance-calculus;
    and $\lambda$-IC and $\lambda$-lazy-IC the $\lambda$ sublanguages of
    IC and of lazy IC, respectively.
    Edge labels:
    \emph{more defined} means every term that converges in the source
    calculus also converges in the target calculus, but not conversely.
    The parenthesized qualifiers distinguish the mechanism:
    \emph{non-strict} more-definedness gains convergence by deeming
    abstractions values without evaluating their bodies;
    \emph{unproductive $\bot$} more-definedness gains convergence by
    deciding unproductive loops such as $\Omega$ as the definite answer
    $\bot$ in finite time, preserving the lazy tree semantics;
    and \emph{powerset} more-definedness gains convergence by
    resolving cyclic self-references to definite sets, as least fixed
    points over the powerset lattice.
    Continuing the edge labels:
    \emph{strictly more expressive} is meant in Felleisen's sense of
    macro-expressibility;
    and \emph{fully abstract} marks full abstraction.}
  \label{fig:calculi-matrix}
\end{figure*}

Figure~\ref{fig:calculi-matrix} situates inheritance-calculus among
several related calculi developed below.
Section~\ref{sec:definitions} defines the semantics via six
mutually recursive first-order functions on paths and sets of paths,
with no function spaces, closures, or higher-order constructs;
the functions are computable by
tabling~\cite{tamaki1986-tabled-resolution}.
Section~\ref{sec:forward-translation} shows that
these functions extend to the $\lambda$-calculus, embedding it.
Section~\ref{sec:asymmetry} establishes that the lazy
$\lambda$-calculus cannot macro-express the inheritance constructors,
so inheritance-calculus is strictly more expressive in
Felleisen's sense~\cite{felleisen1991-expressive-power}.

\section{\bilingual{Syntax}{语法}}
\label{sec:syntax}

\ifInheritanceChinese
设 $e$ 表示一个表达式,$s$ 表示一个元素列表,$c$ 表示一个元素,
$q$ 表示一个限定符,$w$ 表示一个投影列表,$\ell$ 表示一个充当
属性名的标签,$n$ 表示一个 de~Bruijn 索引。
空串 $\varepsilon$ 是空元素列表。
\else
Let $e$ denote an expression, $s$ an element list, $c$ an element,
$q$ a qualifier, $w$ a projection list, $\ell$ a label serving as a
property name, and $n$ a de~Bruijn index.
The empty string $\varepsilon$ is the empty element list.
\fi

\begin{align*}
  e \quad ::= \quad & \{\, s \,\}
  && \text{\bilingual{(record)}{(记录)}} \\[6pt]
  s \quad ::= \quad & \varepsilon
  && \text{\bilingual{(empty)}{(空)}} \\*
  \mid\quad & c
  && \text{\bilingual{(one element)}{(一个元素)}} \\*
  \mid\quad & c,\; s
  && \text{\bilingual{(element, then more)}{(一个元素,然后更多)}} \\[6pt]
  c \quad ::= \quad & \ell \mapsto e
  && \text{\bilingual{(definition)}{(定义)}} \\*
  \mid\quad & q.\, w
  && \text{\bilingual{(inheritance: qualifier and projections)}{(继承:限定符与投影)}} \\*
  \mid\quad & q
  && \text{\bilingual{(inheritance: qualifier only)}{(继承:仅限定符)}} \\*
  \mid\quad & w
  && \text{\bilingual{(inheritance: projections only)}{(继承:仅投影)}} \\[6pt]
  q \quad ::= \quad & \qualifiedthis{\ell_{\mathrm{qualifier}}}
  && \text{\bilingual{(named qualifier)}{(具名限定符)}} \\*
  \mid\quad & \dbi{n}
  && \text{\bilingual{(indexed qualifier)}{(索引限定符)}} \\[6pt]
  w \quad ::= \quad & \ell_{\mathrm{projection}}
  && \text{\bilingual{(one projection)}{(一个投影)}} \\*
  \mid\quad & \ell_{\mathrm{projection}}.\, w
  && \text{\bilingual{(projection, then more)}{(一个投影,然后更多)}}
\end{align*}

\ifInheritanceChinese
一个表达式将一个元素列表 $s$ 包裹在花括号中,而每个元素
$c$ 要么是一个定义 $\ell \mapsto e$,要么是一个继承源。
语法把 $s$ 写成一个递归列表,但该列表被读作一个
\emph{集合}:元素的顺序无关紧要,重复元素也没有任何作用,
因此两个仅在重排或重复上有所不同的元素列表表示
同一个表达式。由于表达式是集合,继承本身就是
可交换、幂等且可结合的。
图~\ref{fig:cheatsheet} 中的记录 $\mathrm{Complex} \mapsto \{ \mathrm{re} \mapsto \ldots,\;
\mathrm{im} \mapsto \ldots,\; \mathrm{Plus} \mapsto \ldots \}$
就是一个这样的多元素列表。

$\ell \mapsto e$ 中的 $\mapsto$ 定义了一个名为 $\ell$、其值为
$e$ 的属性;例如,图~\ref{fig:cheatsheet} 中的 $\mathrm{re} \mapsto \{\ldots\}$ 与
$\mathrm{Complex} \mapsto \{\ldots\}$。
MIXINv2 还通过一个
外部函数接口(FFI)提供标量类型;继承演算就是
去掉 FFI 之后所剩下的部分。
本文中没有任何例子用到 FFI 或任何标量值。

\paragraph{\bilingual{Inheritance}{继承}}
一个继承通过一个限定符 $q$ 引用一个作用域;该限定符
沿作用域层级向上导航,然后通过一个投影列表 $w$
向下投影属性。限定符 $q$ 容许两种等价的
记法:
\begin{description}
  \item[\bilingual{Named form}{具名形式}]
    $\qualifiedthis{\ell_{\mathrm{qualifier}}}.\, w$,
    类比于 Java 的 \texttt{Outer.this.field},其中 $w$ 是一个
    投影列表。
    每个 $\{\ldots\}$ 记录在它的外围作用域中都有一个标签;
    $\ell_{\mathrm{qualifier}}$ 命名目标外围作用域。
    例如,图~\ref{fig:cheatsheet} 中的
    $\qualifiedthis{\mathrm{ComplexModule}}.\mathrm{Float}$
    投影单个标签,而
    $\qualifiedthis{\mathrm{Plus}}.\mathrm{other}.\mathrm{re}$ 展示了一个
    多标签的投影列表 $w$。
  \item[\bilingual{Indexed form}{索引形式}]
    $\dbi{n}.\, w$,
    其中 $n \ge 0$ 是一个 de~Bruijn 索引~\cite{debruijn1972-nameless-dummies}:
    $n = 0$ 指向该引用的外围作用域,
    而 $n$ 每加一,就向外移动一个作用域层级。
\end{description}
两种形式都先取出目标作用域,即所有继承都应用完毕之后
那个完全继承后的记录,然后从中投影由投影列表 $w$
所命名的属性。
具名形式在解析期间脱糖为索引形式,
如第~\ref{sec:mixin-trees} 节所述。

每个 $\{\ldots\}$ 记录字面量都会创建一个新的作用域层级。%
\footnote{
  我们用\emph{记录}指那个语法构造,用\emph{作用域}指
  它所创建的名字解析语境,用\emph{mixin} 指它
  作为可组合单元的角色;在继承演算中这些是同一个实体的三种
  视角。
}
该作用域包含那个 mixin 的所有属性,包括那些
通过继承源继承而来的属性。
例如,$\{a \mapsto \{\},\; b \mapsto \{\}\}$ 是单个 mixin,其作用域
同时包含 $a$ 与 $b$;$\{r,\; a \mapsto \{\}\}$ 也是如此,其中 $r$
是对一个包含 $b$ 的记录的引用。
继承源不会创建额外的作用域层级。
在图~\ref{fig:cheatsheet} 中,绑定到 $\mathrm{Plus}$ 的记录
就是一个这样的作用域,它包含 $\mathrm{other}$、$\mathrm{\_sumReCall}$、
$\mathrm{\_sumRe}$ 以及在其内部直接定义的其他属性。

\begin{itemize}
  \item $\qualifiedthis{\mathrm{Foo}}$(等价地,$\dbi{n}$,其中
    $\mathrm{Foo}$ 在外围作用域之上 $n$ 步)
    取出名为 $\mathrm{Foo}$ 的记录。
  \item $\qualifiedthis{\mathrm{Foo}}.\ell$,或等价地,$\dbi{n}.\ell$,
    从 $\mathrm{Foo}$ 投影属性 $\ell$。
  \item 一个仅限定符的源 $q$(没有投影列表)继承
    整个外围作用域,这可能产生一棵无限深的
    树;图~\ref{fig:cheatsheet} 中裸的 $\mathrm{BuiltIns}$、$\mathrm{Math}$ 与
    $\mathrm{Operator}$ 源就是此类
    仅限定符的导入。
\end{itemize}

当限定符无歧义时,它可以省略:这种简写是
一个裸的投影列表 $w$。它的首个投影针对
那个在其记录字面量中直接定义首个投影标签的
最内层外围作用域来解析,而该源脱糖为
对应的索引形式。
例如,图~\ref{fig:cheatsheet} 中的 $\mathrm{\_sumReCall}.\mathrm{result}$、
$\mathrm{a}.\mathrm{Plus}$ 与 $\mathrm{\_totalCall}.\mathrm{result}$
就是此类省略限定符的简写。
由于这一搜索只检查记录字面量的定义,该
简写只能访问在记录字面量中定义的属性,而不能访问
那些通过继承而继承来的属性;对于继承来的属性或要绕过变量遮蔽,
都需要两种限定形式之一。
\else
An expression encloses an element list $s$ in braces, and each element
$c$ is either a definition $\ell \mapsto e$ or an inheritance source.
The grammar writes $s$ as a recursive list, but this list is read as a
\emph{set}: element order is irrelevant and duplicates have no effect,
so two element lists differing only by reordering or repetition denote
the same expression. Because expressions are sets, inheritance is
inherently commutative, idempotent, and associative.
The record $\mathrm{Complex} \mapsto \{ \mathrm{re} \mapsto \ldots,\;
\mathrm{im} \mapsto \ldots,\; \mathrm{Plus} \mapsto \ldots \}$ in
Figure~\ref{fig:cheatsheet} is one such multi-element list.

The $\mapsto$ in $\ell \mapsto e$ defines a property named $\ell$ with
value $e$; for example, $\mathrm{re} \mapsto \{\ldots\}$ and
$\mathrm{Complex} \mapsto \{\ldots\}$ in Figure~\ref{fig:cheatsheet}.
MIXINv2 also provides scalar types through a
foreign-function interface (FFI); inheritance-calculus is what remains
after removing the FFI.
None of the examples in this paper uses the FFI or any scalar values.

\paragraph{Inheritance}
An inheritance references a scope through a qualifier $q$ that
navigates the scope hierarchy upward, then projects properties downward
through a projection list $w$. The qualifier $q$ admits two equivalent
notations:
\begin{description}
  \item[Named form]
    $\qualifiedthis{\ell_{\mathrm{qualifier}}}.\, w$,
    analogous to Java's \texttt{Outer.this.field}, where $w$ is a
    projection list.
    Each $\{\ldots\}$ record has a label in its enclosing scope;
    $\ell_{\mathrm{qualifier}}$ names the target enclosing scope.
    For example, $\qualifiedthis{\mathrm{ComplexModule}}.\mathrm{Float}$
    in Figure~\ref{fig:cheatsheet} projects a single label, while
    $\qualifiedthis{\mathrm{Plus}}.\mathrm{other}.\mathrm{re}$ shows a
    multi-label projection list $w$.
  \item[Indexed form]
    $\dbi{n}.\, w$,
    where $n \ge 0$ is a de~Bruijn index~\cite{debruijn1972-nameless-dummies}:
    $n = 0$ refers to the reference's enclosing scope,
    and each increment of $n$ moves one scope level further out.
\end{description}
Both forms retrieve the target scope, that is, the fully inherited record
after all inheritance has been applied, and then project the properties
named by the projection list $w$ from it.
The named form desugars to the indexed form during parsing,
as described in Section~\ref{sec:mixin-trees}.

A new scope level is created at each $\{\ldots\}$ record literal.%
\footnote{
  We use \emph{record} for the syntactic construct, \emph{scope}
  for the name-resolution context it creates, and \emph{mixin} for its
  role as a composable unit; in inheritance-calculus these are three views
  of the same entity.
}
The scope contains all properties of that mixin, including those
inherited via inheritance sources.
For example, $\{a \mapsto \{\},\; b \mapsto \{\}\}$ is a single mixin whose scope
contains both $a$ and $b$; so is $\{r,\; a \mapsto \{\}\}$ where $r$
is a reference to a record containing $b$.
Inheritance sources do not create additional scope levels.
In Figure~\ref{fig:cheatsheet}, the record bound to $\mathrm{Plus}$ is
one such scope, containing $\mathrm{other}$, $\mathrm{\_sumReCall}$,
$\mathrm{\_sumRe}$, and the other properties defined directly within it.

\begin{itemize}
  \item $\qualifiedthis{\mathrm{Foo}}$ (equivalently, $\dbi{n}$ where
    $\mathrm{Foo}$ is $n$ steps above the enclosing scope)
    retrieves the record named $\mathrm{Foo}$.
  \item $\qualifiedthis{\mathrm{Foo}}.\ell$, or equivalently, $\dbi{n}.\ell$,
    projects property $\ell$ from $\mathrm{Foo}$.
  \item A qualifier-only source $q$ (with no projection list) inherits
    the entire enclosing scope, which may produce an infinitely deep
    tree; the bare $\mathrm{BuiltIns}$, $\mathrm{Math}$, and
    $\mathrm{Operator}$ sources in Figure~\ref{fig:cheatsheet} are
    qualifier-only imports of this kind.
\end{itemize}

When the qualifier is unambiguous, it may be omitted: the shorthand is
a bare projection list $w$. Its leading projection is resolved against
the innermost enclosing scope that defines the leading projection's
label directly in its record literal, and the source desugars to the
corresponding indexed form.
For example, $\mathrm{\_sumReCall}.\mathrm{result}$,
$\mathrm{a}.\mathrm{Plus}$, and $\mathrm{\_totalCall}.\mathrm{result}$
in Figure~\ref{fig:cheatsheet} are such qualifier-omitted shorthands.
Because this search inspects only record-literal definitions, the
shorthand can only access properties defined in record literals, not
those inherited through inheritance; one of the two qualified forms is
required for inherited properties or to bypass variable shadowing.
\fi

\section{\bilingual{Mixin Trees}{Mixin 树}}
\label{sec:mixin-trees}

\ifInheritanceChinese
在面向对象编程中,一个
\emph{mixin}~\cite{bracha1990-mixin-inheritance} 是一段行为片段;
它可以通过继承组合到一个类上。
继承演算只有一个抽象,即
\emph{深可合并 mixin}(deep-mergeable mixin),其成员是
\emph{深可合并属性}(deep-mergeable properties)。%
\footnote{
  在不引起歧义之处,我们简写为``mixin''与``property''。
  正如图~\ref{fig:cheatsheet} 所示,二者各自推广了它们在先前
  工作中的同名物:Bracha--Cook 的 mixin~\cite{bracha1990-mixin-inheritance} 与
  面向对象语言的属性。
}
它可以同时充当模块、类、对象、方法、字段以及
局部绑定,正如图~\ref{fig:cheatsheet} 所展示的那样,因为
没有不透明的方法:每个值都是
一个透明的记录,而继承总是对子树执行一次深度
合并。
我们把所得到的结构称为一棵\emph{mixin 树}。
它的语义是纯\emph{观察式}(observational)的:给定一个
继承演算程序,我们查询某个
标签 $\ell$ 是否是某个给定位置 $p$ 处的一个属性。

本节通过六个相互递归的
路径上的一阶函数来定义语义。
这六个函数把面向对象的概念映射为
集合论的运算:$\properties$ 与 $\supers$
推广成员查找与继承链,
$\overrides$ 与 $\bases$ 捕捉方法覆盖与
基类,而 $\resolve$ 与 $\this$ 处理名字
解析与限定 $\this$ 解析。
\else
In object-oriented programming, a
\emph{mixin}~\cite{bracha1990-mixin-inheritance} is a fragment of behavior
that can be composed onto a class via inheritance.
Inheritance-calculus has a single abstraction, the
\emph{deep-mergeable mixin}, whose members are
\emph{deep-mergeable properties}.%
\footnote{
  Where no ambiguity arises, we write simply ``mixin'' and ``property.''
  As Figure~\ref{fig:cheatsheet} shows, each generalizes its namesake in prior
  work: the Bracha--Cook mixin~\cite{bracha1990-mixin-inheritance} and the
  property of object-oriented languages.
}
which can simultaneously serve as module, class, object, method, field, and
local binding, as Figure~\ref{fig:cheatsheet} illustrated, because there
are no opaque methods: every value is
a transparent record, and inheritance always performs a deep
merge of subtrees.
We call the resulting structure a \emph{mixin tree}.
Its semantics is purely \emph{observational}: given an
inheritance-calculus program, one queries whether a
label $\ell$ is a property at a given position $p$.

This section defines the semantics via six mutually recursive
first-order functions on paths.
The six functions map object-oriented concepts to
set-theoretic operations: $\properties$ and $\supers$
generalize member lookup and inheritance chains,
$\overrides$ and $\bases$ capture method override and
base classes, and $\resolve$ and $\this$ handle name
resolution and qualified $\this$ resolution.
\fi

\subsection{\bilingual{Definitions}{定义}}
\label{sec:definitions}

\paragraph{\bilingual{Path}{路径}}
\ifInheritanceChinese
一个 $\Path$ 是一个标签序列
$(\ell_1, \ell_2, \ldots, \ell_n)$;它标识树中的一个
位置。
我们用 $p$ 表示一条路径。
\emph{根路径}是空序列 $()$。
给定一条路径 $p$ 与一个标签 $\ell$,
$p \snoc \ell$ 是序列 $p$ 用 $\ell$ 扩展后的结果。
对于非根路径,父路径 $\init(p)$ 是 $p$ 去掉
其最后一个元素后的结果,
而末尾标签 $\last(p)$ 是 $p$ 的最后一个元素。
为方便起见,当路径在上下文中清楚时,我们有时
用 $\ell$ 表示 $\last(p)$。
有了路径,我们就能陈述 AST 在
每条路径处提供什么。
\else
A $\Path$ is a sequence of labels
$(\ell_1, \ell_2, \ldots, \ell_n)$ that identifies a position in
the tree.
We write $p$ for a path.
The \emph{root path} is the empty sequence $()$.
Given a path $p$ and a label $\ell$,
$p \snoc \ell$ is the sequence $p$ extended with $\ell$.
For a nonroot path, the parent path $\init(p)$ is $p$ with
its last element removed,
and the final label $\last(p)$ is the last element of $p$.
For convenience, we sometimes write $\ell$ for $\last(p)$
when the path is clear from context.
With paths in hand, we can state what the AST provides
at each path.
\fi

\paragraph{AST}
\ifInheritanceChinese
来自第~\ref{sec:syntax} 节的一个表达式 $e$ 被解析成一棵 AST。
我们以两种等价的风格给出 AST:首先是一个固定其形状的
代数数据类型,然后是两个派生的投影;本节其余部分的
语义函数消费它们。

\emph{代数数据类型风格}把规范化后的 AST 给出为一个由节点项
构成的纯语法。一个节点 $N$ 把一个定义列表 $d$ 与一个
引用列表 $b$ 配成对;每个定义把一个标签绑定到一个子节点,而
每个引用是一个 de~Bruijn 索引 $n$ 连同一个投影列表
$w$:
\else
An expression $e$ from Section~\ref{sec:syntax} is parsed into an AST.
We give the AST in two equivalent styles: first an algebraic
datatype that fixes its shape, then two derived projections that
the semantic functions in the rest of this section consume.

The \emph{algebraic-datatype style} gives the normalized AST as a pure
grammar of node terms. A node $N$ pairs a definition list $d$ with a
reference list $b$; each definition binds a label to a child node, and
each reference is a de~Bruijn index $n$ together with a projection list
$w$:
\fi
\begin{align*}
  N \quad ::= \quad & \langle\, d \,;\; b \,\rangle
  && \text{\bilingual{(node: definitions and references)}{(节点:定义与引用)}} \\[6pt]
  d \quad ::= \quad & \varepsilon
  && \text{\bilingual{(no definitions)}{(没有定义)}} \\*
  \mid\quad & (\ell \mapsto N),\; d
  && \text{\bilingual{(definition, then more)}{(一个定义,然后更多)}} \\[6pt]
  b \quad ::= \quad & \varepsilon
  && \text{\bilingual{(no references)}{(没有引用)}} \\*
  \mid\quad & (n,\; w),\; b
  && \text{\bilingual{(reference, then more)}{(一个引用,然后更多)}} \\[6pt]
  w \quad ::= \quad & \varepsilon
  && \text{\bilingual{(no projections)}{(没有投影)}} \\*
  \mid\quad & \ell.\, w
  && \text{\bilingual{(projection, then more)}{(一个投影,然后更多)}}
\end{align*}
\ifInheritanceChinese
与表达式的元素列表一样,定义列表 $d$ 是
递归写成的,但读取时不计重排与重复:它
表示从其各绑定收集而来的有限偏映射 $\defs : \Label \rightharpoonup \Node$;该映射
对每个被定义的标签恰有一条目,而
$N.\defs(\ell)$ 是绑定到 $\ell$ 的子节点(当 $\ell$
未绑定时无定义)。同样地,引用列表 $b$ 表示由继承源
$(n,\; w)$ 构成的集合
$\refs \subseteq \mathbb{N} \times \Label^{*}$,读取时不计重排与重复。由于继承
是幂等的且记录是集合,这些商不丢失任何
信息。
该语法只固定形状;一个节点项是良构的,当它
满足:
\else
Like the element list of an expression, the definition list $d$ is
written recursively but read up to reordering and duplication: it
denotes the finite partial map $\defs : \Label \rightharpoonup \Node$
collected from its bindings, with one entry per defined label, and
$N.\defs(\ell)$ is the child bound to $\ell$ (undefined when $\ell$ is
unbound). Likewise the reference list $b$ denotes the set
$\refs \subseteq \mathbb{N} \times \Label^{*}$ of inheritance sources
$(n,\; w)$, read up to reordering and duplication. Because inheritance
is idempotent and a record is a set, these quotients lose no
information.
The grammar fixes only the shape; a node term is well-formed when it
satisfies:
\fi
\begin{enumerate}
  \item \ifInheritanceChinese 路径 $p$ 处的每个源 $(n,\; w)$ 满足 $n < |p|$。\else Each source $(n,\; w)$ at path $p$ satisfies $n < |p|$.\fi
  \item \ifInheritanceChinese 一个 $w$ 为空的源 $(n,\; \varepsilon)$ 是
    一个仅限定符源的规范化;第~\ref{sec:syntax} 节的表达式语法
    没有 $w$ 为空的仅投影源。\else A source $(n,\; \varepsilon)$ with empty $w$ is the
    normalization of a qualifier-only source; the expression grammar of
    Section~\ref{sec:syntax} has no projection-only source with empty
    $w$.\fi
\end{enumerate}
\ifInheritanceChinese
该节点语法被\emph{余归纳地}(coinductively)读取:通过
继承,一个子节点可能展开成一棵无界乃至
无限深的树,所以这些产生式描述的是最大不动点,
而非有限生成的最小不动点。整个程序是单个
全局节点 $\ASTroot$,即顶层记录的规范化 AST。

\emph{余代数风格}(coalgebraic style)用两个由路径 $p$ 索引的
原始函数替换全局树,每次观察一个节点。它们通过沿 $p$
导航从 $\ASTroot$ 派生而来:
从 $\ASTroot$ 出发,在 $p$ 的每个标签处跟随 $\defs$ 映射
到达 $p$ 处的节点,而这两个投影读出其
字段。
\else
The node grammar is read \emph{coinductively}: through
inheritance a child node may expand to a tree of unbounded or even
infinite depth, so the productions describe the greatest fixed point,
not a finitely generated least one. The whole program is a single
global node $\ASTroot$, the normalized AST of the top-level record.

The \emph{coalgebraic style} replaces the global tree by two
primitive functions indexed by a path $p$, observing one node at a
time. They are derived from $\ASTroot$ by navigating along $p$:
starting at $\ASTroot$ and following the $\defs$ map at each label
of $p$ reaches the node at $p$, and the two projections read off its
fields.
\fi
\begin{align}
  \defines(p) &= \operatorname{dom}\bigl(
    \ASTroot.\defs(\ell_1).\defs(\ell_2)\cdots\defs(\ell_n).\defs
  \bigr)
  \label{eq:defines}\\
  \inherits(p) &= \ASTroot.\defs(\ell_1).\defs(\ell_2)\cdots\defs(\ell_n).\refs
  \label{eq:inherits}
\end{align}
\ifInheritanceChinese
其中 $p = (\ell_1, \ldots, \ell_n)$,从而 $\defines(p)$ 是
那些使得 $p$ 包含一个定义 $\ell \mapsto e$ 的标签 $\ell$ 的集合,
而 $\inherits(p)$ 是 $p$ 处的引用对
$(n,\; w)$ 的集合。
在根路径 $()$ 处导航为空,给出
$\defines(()) = \operatorname{dom}(\ASTroot.\defs)$ 以及
$\inherits(()) = \ASTroot.\refs$。
下文的语义函数只使用 $\defines$ 与 $\inherits$,
所以任一风格皆可取作原始。

在解析期间,三种语法形式全都被解析为
de~Bruijn 索引对 $(n,\; w)$。
\emph{索引形式} $\dbi{n}.\, w$
已经直接携带 de~Bruijn 索引 $n$;
不需要解析。
路径 $p$ 处的一个\emph{具名引用} $\qualifiedthis{\ell_{\mathrm{qualifier}}}.\, w$
通过在 $p$ 的标签中找到
$\ell_{\mathrm{qualifier}}$ 的最后一次出现来解析。
设 $p_{\mathrm{target}}$ 为 $p$ 的前缀,直到并包括
那次出现。
de~Bruijn 索引为 $n = |p| - |p_{\mathrm{target}}| - 1$,
而投影列表 $w$ 原样保留。
路径 $p$ 处一个首投影为 $\ell$ 的\emph{词法引用} $w$
通过找到 $p$ 的最近前缀 $p'$
使得 $\ell \in \defines(p')$ 来解析。
de~Bruijn 索引为 $n = |p| - |p'| - 1$,而投影列表
$w$ 原样保留。
三种形式都产生同一种表示;下文的语义函数
只在 de~Bruijn 索引对上操作。

AST 是纯数据,由程序文本固定,而不是
运行期状态的函数。%
\footnote{
  一个实现可以非严格地解析 AST:$\ASTroot$ 不必
  预先被完全物化。每个 $\Node$ 把它的源
  文本保持为未解析,直到某个查询观察它,从而读取
  $\defines(p)$ 或 $\inherits(p)$ 时按需强制
  沿 $p$ 的子树。这只影响解析何时发生,而不影响其结果,
  因为 AST 是引用透明的。
}
有了两个原始函数 $\defines$ 与 $\inherits$,我们现在可以回答
本节开头提出的那个问题:
某个标签 $\ell$ 是否是某个给定位置 $p$ 处的一个属性。
\else
for $p = (\ell_1, \ldots, \ell_n)$, so that $\defines(p)$ is the set
of labels $\ell$ for which $p$ contains a definition $\ell \mapsto e$,
and $\inherits(p)$ is the set of reference pairs
$(n,\; w)$ at $p$.
At the root path $()$ the navigation is empty, giving
$\defines(()) = \operatorname{dom}(\ASTroot.\defs)$ and
$\inherits(()) = \ASTroot.\refs$.
The semantic functions below use only $\defines$ and $\inherits$,
so either style may be taken as primitive.

During parsing, all three syntactic forms are resolved to
de~Bruijn index pairs $(n,\; w)$.
The \emph{indexed form} $\dbi{n}.\, w$
already carries the de~Bruijn index $n$ directly;
no resolution is needed.
A \emph{named reference} $\qualifiedthis{\ell_{\mathrm{qualifier}}}.\, w$
at path $p$ is resolved by finding the last occurrence of
$\ell_{\mathrm{qualifier}}$ among the labels of $p$.
Let $p_{\mathrm{target}}$ be the prefix of $p$ up to and including
that occurrence.
The de~Bruijn index is $n = |p| - |p_{\mathrm{target}}| - 1$
and the projection list $w$ is carried over unchanged.
A \emph{lexical reference} $w$ at path $p$, with leading
projection $\ell$, is resolved by finding the nearest prefix $p'$ of $p$
such that $\ell \in \defines(p')$.
The de~Bruijn index is $n = |p| - |p'| - 1$ and the projection list
$w$ is carried over unchanged.
All three forms produce the same representation; the semantic functions
below operate only on de~Bruijn index pairs.

The AST is pure data, fixed by the program text and not a function
of runtime state.%
\footnote{
  An implementation may parse the AST non-strictly: $\ASTroot$ need
  not be materialized in full up front. Each $\Node$ keeps its source
  text unparsed until a query observes it, so that reading
  $\defines(p)$ or $\inherits(p)$ forces the subtree along $p$ on
  demand. This affects only when parsing happens, not its result,
  since the AST is referentially transparent.
}
Given the two primitives $\defines$ and $\inherits$, we can now answer
the question posed at the beginning of this section:
whether a label $\ell$ is a property at a given position $p$.
\fi

\paragraph{\bilingual{Properties}{属性}}
\ifInheritanceChinese
一条路径的属性是它自身的属性,连同
从所有 supers 继承而来的属性:
\else
The properties of a path are its own properties together with
those inherited from all supers:
\fi
\begin{equation}\label{eq:properties}
  \properties(p) =
  \bigl\{\; \ell \;\big|\;
    (\_,\; p_{\mathrm{override}}) \in \supers(p),\;
    \ell \in \defines(p_{\mathrm{override}})
  \;\bigr\}
\end{equation}
\ifInheritanceChinese
本节其余部分定义 $\supers$,
即传递的继承闭包,
及其依赖项。
\else
The remainder of this section defines $\supers$,
the transitive inheritance closure,
and its dependencies.
\fi

\paragraph{Supers}
\ifInheritanceChinese
直观上,$\supers(p)$ 收集 $p$ 所继承的每一条路径:
$p$ 自身的 $\overrides$,
加上 $p$ 的每个
直接基 $\bases$ 的 $\overrides$,如此传递下去。
每个结果与到达它所经由的继承点语境
$\init(p_{\mathrm{base}})$ 配成对。
我们下文定义的限定 $\this$ 解析需要这一来源,
以便把一个定义点 mixin 映射回
那些纳入它的继承点路径:
\else
Intuitively, $\supers(p)$ collects every path that $p$ inherits from:
the $\overrides$ of $p$ itself,
plus the $\overrides$ of each
direct base $\bases$ of $p$, and so on transitively.
Each result is paired with the inheritance-site context
$\init(p_{\mathrm{base}})$ through which it is reached.
This provenance is needed by qualified $\this$ resolution,
which we define below, to map a definition-site mixin back to
the inheritance-site paths that incorporate it:
\fi
\begin{equation}\label{eq:supers}
  \supers(p) =
  \bigl\{\; (\init(p_{\mathrm{base}}),\; p_{\mathrm{override}}) \;\big|\;
    p_{\mathrm{base}} \in \bases^*(p),\;
    p_{\mathrm{override}} \in \overrides(p_{\mathrm{base}})
  \;\bigr\}
\end{equation}
\ifInheritanceChinese
这里 $\bases^*$ 表示 $\bases$ 的自反传递闭包。
$\supers$ 公式依赖两个函数:
$\overrides$ 与单跳引用目标 $\bases$。
我们先定义 $\overrides$。
\else
Here $\bases^*$ denotes the reflexive-transitive closure of $\bases$.
The $\supers$ formula depends on two functions:
$\overrides$ and the one-hop reference targets $\bases$.
We define $\overrides$ first.
\fi

\paragraph{Overrides}
\ifInheritanceChinese
由于一个 mixin 可以扮演一个方法的角色,
一个子作用域应当能够
通过父作用域的标签找到它,即便该子作用域
并不重新定义那个标签;这就是方法覆盖。
具体地说,继承可以在同一个作用域层级引入同一标签的
多个定义;
$\overrides(p)$ 收集所有与 $p$ 共享
\emph{相同身份}(same identity)的这类路径,从而使它们的定义
相互合并而非相互遮蔽,以实现深度合并。
$p$ 的 overrides 包含 $p$ 自身
以及对于 $\init(p)$ 的每个
也定义 $\last(p)$ 的分支 $p_{\mathrm{branch}}$,$p_{\mathrm{branch}} \snoc \last(p)$:
\else
Since a mixin can play the role of a method,
a child scope should be
able to find it through the parent's label even if the child
does not redefine that label; this is method override.
Concretely, inheritance can introduce multiple
definitions of the same label at the same scope level;
$\overrides(p)$ collects all such paths that share the
\emph{same identity} as $p$, so that their definitions
merge rather than shadow each other, enabling deep merge.
The overrides of $p$ include $p$ itself
and $p_{\mathrm{branch}} \snoc \last(p)$ for every
branch $p_{\mathrm{branch}}$ of $\init(p)$ that also defines $\last(p)$:
\fi
\begin{equation}\label{eq:overrides}
  \overrides(p) =
  \begin{cases}
    \{p\} & \text{if } p = ()  \\[6pt]
    \{p\}
    \;\cup\;
    \left\{\; p_{\mathrm{branch}} \snoc \last(p) \;\middle|\;
      \begin{aligned}
        &(\_,\; p_{\mathrm{branch}}) \in \supers(\init(p)), \\
        &\text{s.t.}\; \last(p) \in \defines(p_{\mathrm{branch}})
      \end{aligned}
    \right\}
    & \text{if } p \neq ()
  \end{cases}
\end{equation}
\ifInheritanceChinese
还剩下要定义 $\bases$,即 $\supers$ 的另一个依赖项。
\else
It remains to define $\bases$, the other dependency of $\supers$.
\fi

\paragraph{Bases}
\ifInheritanceChinese
直观上,$\bases(p)$ 是 $p$ 通过引用直接
继承的那些路径,类比于面向对象语言中的直接
基类。
具体地说,$\bases$ 把 $p$ 的 $\overrides$ 中的每个引用
解析一步:
\else
Intuitively, $\bases(p)$ are the paths that $p$ directly
inherits from via references, analogous to the direct base
classes in object-oriented languages.
Concretely, $\bases$ resolves every reference in
$p$'s $\overrides$ one step:
\fi
\begin{equation}\label{eq:bases}
  \bases(p) =
  \left\{\; p_{\mathrm{target}} \;\middle|\;
    \begin{aligned}
      &p_{\mathrm{override}} \in \overrides(p), \\
      &(n,\; w) \in \inherits(p_{\mathrm{override}}), \\
      &p_{\mathrm{target}} \in \resolve(\init(p),\; p_{\mathrm{override}},\; n,\; w)
    \end{aligned}
  \right\}
\end{equation}
\ifInheritanceChinese
$\bases$ 公式调用引用解析函数 $\resolve$,我们接下来定义它。
\else
The $\bases$ formula calls the reference resolution function $\resolve$, which we define next.
\fi

\paragraph{\bilingual{Reference resolution}{引用解析}}
\ifInheritanceChinese
直观上,$\resolve$ 把一个语法引用转化为它在
完全继承后的树中所指向的路径集合。
它取一条继承点路径 $p_{\mathrm{site}}$、
一条定义点路径 $p_{\mathrm{def}}$、
一个 de~Bruijn 索引 $n$,
以及投影列表 $w$。
解析分两个阶段进行:
$\this$ 从外围作用域 $\init(p_{\mathrm{def}})$ 出发
执行 $n$ 次向上步进,
把定义点路径映射为继承点路径;
然后追加
$w = \ell_{\mathrm{projection},1}\ldots\ell_{\mathrm{projection},k}$
的向下投影:
\else
Intuitively, $\resolve$ turns a syntactic reference into the
set of paths it points to in the fully inherited tree.
It takes an inheritance-site path $p_{\mathrm{site}}$,
a definition-site path $p_{\mathrm{def}}$,
a de~Bruijn index $n$,
and the projection list $w$.
Resolution proceeds in two phases:
$\this$ performs $n$ upward steps
starting from the enclosing scope $\init(p_{\mathrm{def}})$,
mapping the definition-site path to inheritance-site paths;
then the downward projections of
$w = \ell_{\mathrm{projection},1}\ldots\ell_{\mathrm{projection},k}$
are appended:
\fi
\begin{multline}\label{eq:resolve}
  \resolve(p_{\mathrm{site}},\; p_{\mathrm{def}},\; n,\; w)
  = \\
  \bigl\{\;
    p_{\mathrm{current}} \snoc \ell_{\mathrm{projection},1} \snoc \cdots \snoc \ell_{\mathrm{projection},k}
    \;\big|
    p_{\mathrm{current}} \in
    \this(\{p_{\mathrm{site}}\},\; \init(p_{\mathrm{def}}),\; n)
  \;\bigr\}
\end{multline}
\ifInheritanceChinese
$\resolve$ 返回一个集合而非单条路径,因为
$\this$ 可能解析到多个目标。
我们现在定义 $\this$。
\else
$\resolve$ returns a set rather than a single path because
$\this$ may resolve to multiple targets.
We now define $\this$.
\fi

\paragraph{\bilingual{Qualified $\this$ resolution}{限定 $\this$ 解析}}
\ifInheritanceChinese
$\this$ 回答这个问题:
``在完全继承后的树中,定义点作用域
$p_{\mathrm{def}}$ 究竟住在哪里?''
每一步在前沿 $S$ 中每条路径的 supers 当中,
找出那些其 override 分量匹配
$p_{\mathrm{def}}$ 的,并把对应的
继承点路径收集为新的前沿。
经过 $n$ 步后,前沿就包含答案:
\else
$\this$ answers the question:
``in the fully inherited tree, where does
the definition-site scope $p_{\mathrm{def}}$ actually live?''
Each step finds, among the supers of every path in the
frontier $S$, those whose override component matches
$p_{\mathrm{def}}$, and collects the corresponding
inheritance-site paths as the new frontier.
After $n$ steps the frontier contains the answer:
\fi
{\small
  \begin{equation}\label{eq:this}
    \this(S,\; p_{\mathrm{def}},\; n) =
    \begin{cases}
      S & \text{if } n = 0 \\[6pt]
      \this\!\left(
        \left\{\; p_{\mathrm{site}} \;\middle|\;
          \begin{aligned}
            &p_{\mathrm{current}} \in S, \\
            &(p_{\mathrm{site}},\; p_{\mathrm{override}})
            \in \supers(p_{\mathrm{current}}), \\
            &\text{s.t.}\; p_{\mathrm{override}} = p_{\mathrm{def}}
          \end{aligned}
        \right\},\;
        \init(p_{\mathrm{def}}),\;
        n - 1
      \right)
      & \text{if } n > 0
    \end{cases}
\end{equation}}%
\ifInheritanceChinese
每一步 $\init$ 把 $p_{\mathrm{def}}$ 缩短一个标签,
而 $n$ 减少一。
由于 $n$ 是一个非负整数,该递归终止。%
\footnote{
  若在某一步不存在匹配对
  $(p_{\mathrm{site}},\; p_{\mathrm{override}})$,
  则前沿变为空集,而
  $\resolve$ 返回 $\varnothing$:该引用没有
  目标;不继承任何属性。
  这是无类型演算最为自洽的
  语义。
  MIXINv2\anon[~(补充材料)]{~\cite{mixin2025}} 更为严格:它的解析器拒绝
  前沿变为空的引用,在编译期报告一个
  解析错误。
}

这就完成了计算 $\properties(p)$ 所需的整条定义链。

附录~\ref{app:evaluation-trace} 把这六个方程逐步
应用到一个突出 $\resolve$ 集合值本质的例子上;
在那里前沿 $S$ 携带多个目标,而
$\this$ 同时解析到它们,而非如 Scala~3 与 NixOS 模块
这类面向对象语言中所见的单个目标
(附录~\ref{app:scala-multi-target})。
\else
At each step $\init$ shortens $p_{\mathrm{def}}$ by one label
and $n$ decreases by one.
Since $n$ is a nonnegative integer, the recursion terminates.%
\footnote{
  If no matching pair
  $(p_{\mathrm{site}},\; p_{\mathrm{override}})$ exists
  at some step, the frontier becomes the empty set, and
  $\resolve$ returns $\varnothing$: the reference has no
  target; no properties are inherited.
  This is the most self-consistent semantics for the
  untyped calculus.
  MIXINv2\anon[~(supplementary material)]{~\cite{mixin2025}} is stricter: its resolver rejects
  references whose frontier becomes empty, reporting a
  resolution error at compile time.
}

This completes the chain of definitions needed to compute $\properties(p)$.

Appendix~\ref{app:evaluation-trace} applies the six equations step by
step to an example highlighting the set-valued nature of $\resolve$,
where the frontier $S$ carries multiple targets that
$\this$ resolves to simultaneously rather than the single target
seen in object-oriented languages such as Scala~3 and NixOS modules
(Appendix~\ref{app:scala-multi-target}).
\fi

\subsection{\bilingual{Recursive Evaluation}{递归求值}}
\ifInheritanceChinese
六个方程~(\ref{eq:properties})--(\ref{eq:this})
通过相互递归定义 $\properties(p)$:
$\ell \in \properties(p)$ 成立,当且仅当它能
由这些函数的有限次应用链
确立。
这些函数\emph{就是}解释器:观察者
通过选择检查哪条路径来驱动计算,
而每个方程按需展开。
由于这些函数是纯的且与顺序无关,
(\ref{eq:supers}) 中的自反传递闭包 $\bases^*$
可以按广度优先或
深度优先来探索而不影响结果;
路径可以被 intern,从而结构相等
退化为指针相等;
而这六个函数构成一个以路径为键的可记忆化
动态规划递推式。

在操作上,求值使用\emph{tabled
求值}~\cite{tamaki1986-tabled-resolution}:
每个查询在首次遇到时被记忆化,而
重入调用,即对应于依赖图中的环的调用,
返回当前的累加器而非发散。
在指称上,六个方程定义了一个单调的
直接后承
算子~$T_P$~\cite{vanemden1976-predicate-logic-semantics};
它的右端只使用 $\cup$、$\in$、相等以及
正向集合概括,没有否定、集合
差或补。
由 Knaster--Tarski
定理~\cite{tarski1955-lattice-fixpoint-theorem},$T_P$ 有一个唯一的
最小不动点 $\mathrm{lfp}(T_P)$;它充当
指称语义。
这一指称刻画把上文引入的观察式
语义形式化:观察者驱动的查询
由同一组不动点方程计算。
求值是否终止取决于程序。
回忆 $\supers$(方程~\ref{eq:supers})取
$\bases$(方程~\ref{eq:bases})的自反传递
闭包;所得的图就是
\emph{继承层级}(inheritance hierarchy)。
由于每个继承源在 mixin 树中占据一个不同的节点,
这一层级可从
树结构中直接读出。
继承层级无环的程序在
一趟内终止。
当层级包含环时,
当每个查询有有限的答案域时求值
终止:有限格上的升链条件保证
收敛~\cite{bancilhon1986-recursive-query-strategies},
如同在 Datalog~\cite{vanemden1976-predicate-logic-semantics} 中那样。
当层级无限时,求值可能
发散;这是图灵完备性所固有的。
附录~\ref{app:well-definedness} 证明 tabled
求值绝不返回错误答案,并把
这些终止条件形式化。
\else
The six equations~(\ref{eq:properties})--(\ref{eq:this})
define $\properties(p)$ by mutual recursion:
$\ell \in \properties(p)$ holds if and only if it can
be established by a finite chain of applications of
these functions.
The functions \emph{are} the interpreter: the observer
drives the computation by choosing which path to
inspect, and each equation unfolds on demand.
Because the functions are pure and order-independent,
the reflexive-transitive closure $\bases^*$
in~(\ref{eq:supers}) may be explored breadth-first or
depth-first without affecting the result;
paths may be interned so that structural equality
reduces to pointer equality;
and the six functions form a memoizable
dynamic-programming recurrence keyed by paths.

Operationally, evaluation uses \emph{tabled
evaluation}~\cite{tamaki1986-tabled-resolution}:
each query is memoised on first encounter, and
re-entrant calls, which correspond to cycles in the dependency graph,
return the current accumulator rather than diverging.
Denotationally, the six equations define a monotone
immediate consequence
operator~$T_P$~\cite{vanemden1976-predicate-logic-semantics}
whose right-hand sides use only $\cup$, $\in$, equality, and
positive set comprehension, with no negation, set
difference, or complement.
By the Knaster--Tarski
theorem~\cite{tarski1955-lattice-fixpoint-theorem}, $T_P$ has a unique
least fixed point $\mathrm{lfp}(T_P)$, which serves as
the denotational semantics.
This denotational characterization formalizes the observational
semantics introduced above: the observer-driven queries are
computed by the same fixed-point equations.
Whether evaluation terminates depends on the program.
Recall that $\supers$ (equation~\ref{eq:supers}) takes the
reflexive-transitive closure of $\bases$
(equation~\ref{eq:bases}); the resulting graph is the
\emph{inheritance hierarchy}.
Since each inheritance source occupies a distinct node in
the mixin tree, this hierarchy is directly readable from
the tree structure.
Programs whose inheritance hierarchy is acyclic terminate in
one pass.
When the hierarchy contains cycles,
evaluation terminates when each query has a finite
answer domain: the ascending chain condition on the finite
lattice guarantees
convergence~\cite{bancilhon1986-recursive-query-strategies},
as in Datalog~\cite{vanemden1976-predicate-logic-semantics}.
When the hierarchy is infinite, evaluation may
diverge; this is inherent to Turing completeness.
Appendix~\ref{app:well-definedness} proves that tabled
evaluation never returns a wrong answer and formalizes
these termination conditions.
\fi

%% === 审稿人备忘：六个方程的渐近复杂度分析 ===
%%
%% 论文声称语义是一阶的，审稿人可能追问：一阶性是否带来额外的
%% 计算开销？下面按场景分析。
%%
%% 关键观察：set comprehension 自动去重（集合语义），加上路径
%% intern（结构相等退化为指针相等）和 memoization，相同的
%% (path, query) 对只计算一次。这不是可选的优化，而是集合语义
%% 的固有性质——集合不会包含重复元素。
%%
%% --- 场景 1：λ-calculus 翻译 ---
%% Single-Path Lemma (Lemma 3) 保证 this 的 frontier 集合 S
%% 始终只有一个元素。每个路径的 overrides 也只有一个元素。
%% bases 每个节点最多一个继承源。supers 的传递闭包是线性链。
%% 复杂度：O(n)，其中 n 是 ANF term 的大小。
%% 与传统环境机（environment machine）一致。
%%
%% --- 场景 2：Object algebra（k 个 constructor，t 个 operation 文件，m 个字段） ---
%% 多个 operation 文件各自继承同一个 data 文件，然后在最终组合时
%% merge。每个 constructor 在每个 operation 文件里各有一个定义，
%% 通过 overrides 合并，所以 overrides 大小为 O(t)。
%%
%% 关键点：this 的 frontier 是 1，不是"去重后 O(1)"，就是 1。
%% 原因：qualified this 解析的是外层 scope（如 NatFactory），
%% 而这个外层 scope 在继承链上只有一条路径到达——因为是从一个
%% 具体的 operation 定义点出发往上走的。multi-target this 的前提
%% 是同一个定义点被继承到多个不同的 site，这只在 trie union 时
%% 发生。Object algebra 的继承是把多个 operation merge 到同一个
%% factory 上，不是把多个 factory merge 成一个值。方向不同。
%%
%% 这也解释了与 Scala 的关系：Scala 拒绝 multi-target self-reference resolution
%%（Appendix B 的 "conflicting base types"），但在 object algebra
%% 场景下 this 本来就是 single-target 的，所以 Scala 的拒绝是
%% 过度保守——它拒绝了一类去重后（实际上根本不需要去重，因为
%% frontier 就是 1）完全合法的组合。
%%
%% 复杂度：O(k · m · t)，与 Scala trait linearization 一致。
%%
%% --- 场景 3：Nat/BinNat/Boolean 算术 ---
%% 数值的递归深度为 n（Nat 值的大小），每次递归触发一次 this
%% 解析（single-target），遍历 O(t) 个 override。
%% Nat 算术复杂度：O(n · t)，即 Peano 算术 + 常数因子。
%% BinNat 算术复杂度：O(log n · t)，即二进制算术 + 常数因子。
%%
%% --- 场景 4：Cartesian product / 逻辑编程 ---
%% 这是唯一出现真正 multi-target this 的场景。
%% trie union 的 bases 包含多条路径，每条路径的 addend 通过
%% this 解析到另一个 trie union，路径数相乘。
%% 对于 |R₁| × |R₂| 的 join，复杂度为 O(|R₁| · |R₂|)
%% （memoization 下）。
%%
%% 这与 B-Tree 引擎的 Datalog nested-loop join 复杂度一致。
%% trie 天然有索引——路径本身是前缀树，共享前缀的值自动共享
%% 计算（memoization 命中）。这与 Datalog 引擎对 B-Tree 做
%% prefix scan 在精神上一致。
%%
%% --- 总结 ---
%%
%% | 场景               | 复杂度              | 对比               |
%% |--------------------|---------------------|--------------------|
%% | λ-calculus         | O(n)                | = 环境机           |
%% | Object algebra     | O(k·m·t)            | = Scala线性化      |
%% | Nat 算术           | O(n·t)              | = Peano 算术       |
%% | BinNat 算术        | O(log n · t)        | = 二进制算术       |
%% | Cartesian product  | O(|R₁|·|R₂|)       | = nested-loop join |
%%
%% 每个场景都与对应的专用引擎复杂度一致。一阶性没有带来额外开销。
%%
%% Single-Path Lemma 的适用范围比论文明示的更大：不仅 λ-calculus
%% 翻译是 single-path，所有不涉及 trie union 的正常 OOP/object
%% algebra 代码都是 single-path。multi-target this 是 trie union
%% 的专属现象，而 trie union 对应的恰好是关系代数的 join，其
%% 复杂度是固有的，与引擎无关。
%% ================================================================

\section{\bilingual{Embedding the \texorpdfstring{$\lambda$}{λ}-Calculus}{嵌入 \texorpdfstring{$\lambda$}{λ}-演算}}
\label{sec:translation}
\label{sec:forward-translation}

\ifInheritanceChinese
以下翻译将 A-正规形式（ANF）~\cite{flanagan1993-essence-compiling-continuations}的 $\lambda$-演算映射到继承演算。
在 ANF 中,每个中间结果都绑定到一个名字,且每个表达式至多包含一次应用。
我们用一个\emph{值绑定} $\mathbf{let}\; x = V \;\mathbf{in}\; M$ 来扩展标准 ANF 语法,该值绑定将一个并非应用结果的值命名:%
\footnote{
  标准 ANF 允许将值内联用作操作数;值绑定产生式为其命名,使得翻译后程序的每个子表达式都可通过词法引用访问(附录~\ref{app:expressiveness})。
}
\[
  \begin{array}{r@{\;::=\;}l}
    M & \mathbf{let}\; x = V_1\; V_2 \;\mathbf{in}\; M
    \mid \mathbf{let}\; x = V \;\mathbf{in}\; M
    \mid V_1\; V_2
    \mid V \\
    V & x \mid \lambda x.\, M
  \end{array}
\]
每个 $\lambda$-项都可以通过对中间结果命名的方式机械地转换为 ANF。\footnote{%
  这可以选择性地包括对值的命名;由于每个 ANF 项也是一个 $\lambda$-项,ANF 既不是 $\lambda$-演算的限制,也不是其扩展。
}
ANF 结构与继承演算天然吻合:每个 $\mathbf{let}$-绑定成为一个命名的记录属性,每次应用则通过 $.\mathrm{result}$ 包装并投影。

设 $\mathcal{T}$ 为翻译函数。每个 $\lambda$-抽象在翻译后的 mixin 树中引入一个作用域层级。
\else
The following translation maps the $\lambda$-calculus in
A-normal form (ANF)~\cite{flanagan1993-essence-compiling-continuations} to inheritance-calculus.
In ANF, every intermediate result is bound to a name and each
expression contains at most one application.
We extend the standard ANF grammar with a \emph{value binding}
$\mathbf{let}\; x = V \;\mathbf{in}\; M$, which gives a name to a
value that is not an application result:%
\footnote{
  Standard ANF allows values inline as operands;
  the value-binding production names them so that every
  subexpression of the translated program is accessible
  via a lexical reference (Appendix~\ref{app:expressiveness}).
}
\[
  \begin{array}{r@{\;::=\;}l}
    M & \mathbf{let}\; x = V_1\; V_2 \;\mathbf{in}\; M
    \mid \mathbf{let}\; x = V \;\mathbf{in}\; M
    \mid V_1\; V_2
    \mid V \\
    V & x \mid \lambda x.\, M
  \end{array}
\]
Every $\lambda$-term can be mechanically converted to ANF by
naming intermediate results.\footnote{%
  This can optionally include naming values as well;
  since every ANF term is also a $\lambda$-term,
  ANF is neither a restriction nor an extension of the $\lambda$-calculus.
}
The ANF structure matches inheritance-calculus naturally: each
$\mathbf{let}$-binding becomes a named record property, and
each application is wrapped and projected via
$.\mathrm{result}$.

Let $\mathcal{T}$ denote the translation function.
Each $\lambda$-abstraction introduces one scope level
in the translated mixin tree.
\fi

\medskip
\begin{center}
  \small
  \begin{tabular}{l@{\quad$\longrightarrow$\quad}l}
    $x$ \text{\scriptsize(\bilingual{$\lambda$-bound, index $n'$}{$\lambda$-绑定,索引 $n'$})} &
    $\dbi{n'}.\mathrm{argument}$ \\[4pt]
    $x$ \text{\scriptsize(\bilingual{let-bound}{let-绑定})} &
    $x.\mathrm{result}$ \\[4pt]
    $\lambda x.\, M$ &
    $\{\mathrm{argument} \mapsto \{\},\;
    \mathrm{result} \mapsto \mathcal{T}(M)\}$ \\[4pt]
    $\mathbf{let}\; x = V_1\; V_2
    \;\mathbf{in}\; M$ &
    $\{x \mapsto \{\mathcal{T}(V_1),\;
      \mathrm{argument} \mapsto \mathcal{T}(V_2)\},$ \\
    & $\;\mathrm{result} \mapsto \mathcal{T}(M)\}$ \\[4pt]
    $\mathbf{let}\; x = V
    \;\mathbf{in}\; M$ &
    $\{x \mapsto \{\mathrm{result} \mapsto \mathcal{T}(V)\},\;
    \mathrm{result} \mapsto \mathcal{T}(M)\}$ \\[4pt]
    $V_1\; V_2$ \text{\scriptsize(\bilingual{tail call}{尾调用})} &
    $\{\mathrm{tailCall} \mapsto \{\mathcal{T}(V_1),\;
      \mathrm{argument} \mapsto \mathcal{T}(V_2)\},$ \\
    & $\;\mathrm{result} \mapsto \{\mathrm{tailCall}.\mathrm{result}\}\}$
  \end{tabular}
\end{center}
\medskip

\ifInheritanceChinese
带 de~Bruijn 索引 $n$ 的 $\lambda$-绑定变量 $x$(即从引用处到其绑定者之间包围的 $\lambda$ 数目)翻译为 $\dbi{n'}.\mathrm{argument}$,该引用到达绑定 $x$ 的抽象并访问其 $\mathrm{argument}$ 槽。%
\footnote{\label{fn:debruijn-adjust}在 $\lambda$-演算中,de~Bruijn 索引只计数 $\lambda$-抽象。在继承演算中,de~Bruijn 索引计数所有包围的作用域层级,包括由 $\mathbf{let}$-绑定和尾调用引入的层级。因此翻译会调整该索引:若一个 $\lambda$-绑定变量跨越 $k$ 个中间的 $\mathbf{let}$-绑定或尾调用作用域(即 $k$ 计数从变量出现处到其绑定的 $\lambda$ 在 ANF 语法树中路径上的 $\mathbf{let}$- 和尾调用构造数),则其继承演算索引为 $n' = n + k$,而非 $n$。let-绑定变量不受影响,因为它们使用词法引用(如 $x.\mathrm{result}$)而非 de~Bruijn 索引。我们假设翻译引入的合成标签($\mathrm{argument}$、$\mathrm{result}$、$\mathrm{tailCall}$)与用户定义的标签(如 let-绑定名 $x$)来自不相交的字母表,因此不会发生冲突。}
let-绑定变量 $x$ 翻译为词法引用 $x.\mathrm{result}$,从封闭记录中的兄弟属性 $x$ 投影应用结果。

抽象 $\lambda x.\, M$ 翻译为具有自身属性 $\{\mathrm{argument}, \mathrm{result}\}$ 的记录;我们称之为\emph{抽象形状}。在应用 $\mathbf{let}$-绑定规则中,名字 $x$ 绑定到继承记录 $\{\mathcal{T}(V_1),\;\mathrm{argument} \mapsto \mathcal{T}(V_2)\}$,而 $\mathcal{T}(M)$ 可引用 $x.\mathrm{result}$ 以获取应用的值。在值 $\mathbf{let}$-绑定规则中,$x$ 将 $\mathcal{T}(V)$ 包装在一个 $\mathrm{result}$ 子节点之下,使得统一访问器 $x.\mathrm{result}$ 产生 $\mathcal{T}(V)$ 本身,即翻译后的值。这个额外的包装在语义上是惰性的,但与证明相关:它使得用于应用的同一 $\mathrm{result}$-追踪论证在充分性证明中能均匀地扩展到值 $\mathbf{let}$-绑定(附录~\ref{app:bohm-tree-proofs})。在尾调用规则中,新标签 $\mathrm{tailCall}$ 起到相同作用,而 $\mathrm{result}$ 投影出答案。在应用和尾调用两种情形中,执行应用的继承被封装在一个命名属性之后,使其内部结构(被调用者的 $\mathrm{argument}$ 和 $\mathrm{result}$ 标签)不会泄漏到封闭作用域。

以上六条规则构成了从 $\lambda$-演算(经由 ANF)到继承演算的完整翻译:每个 ANF 项都有一个像,且构造是复合的。由于 $\lambda$-演算是图灵完备的,继承演算也是。%
\footnote{
  MIXINv2 将符号解析(解析词法引用和限定 $\this$ 引用)与运行时 mixin 求值分离,但符号解析是逐符号即时执行的,而非提前执行,因此它不提供传统意义上的静态类型检查。尽管如此,MIXINv2 的解析阶段比无类型的继承演算更严格。本文中的所有示例均从 MIXINv2 移植而来,并同时符合两种语义。为满足解析器而添加的元素对此处呈现的算法无害,且可能有助于可读性。
}
然而,图灵完备性并未说明翻译保留了什么;下一小节将精确刻画语义对应关系。
\else
A $\lambda$-bound variable $x$ with de~Bruijn index $n$
(the number of enclosing $\lambda$s between
the reference and its binder) becomes
$\dbi{n'}.\mathrm{argument}$,
which reaches the abstraction that binds $x$ and
accesses its $\mathrm{argument}$ slot.%
\footnote{\label{fn:debruijn-adjust}In the $\lambda$-calculus, de~Bruijn
  indices count only $\lambda$-abstractions.
  In inheritance-calculus, de~Bruijn indices count all
  enclosing scope levels, including those introduced by
  $\mathbf{let}$-bindings and tail calls.
  The translation therefore adjusts the index:
  if a $\lambda$-bound variable crosses $k$ intervening
  $\mathbf{let}$-binding or tail-call scopes
  (that is, $k$ counts the $\mathbf{let}$- and tail-call
    constructs on the path from the variable occurrence
  to its binding $\lambda$ in the ANF syntax tree),
  its inheritance-calculus
  index is $n' = n + k$,
  not~$n$.
  Let-bound variables are unaffected, as they use
  lexical references (such as $x.\mathrm{result}$) rather than
  de~Bruijn indices.
  We assume that the synthetic labels introduced by the
  translation ($\mathrm{argument}$, $\mathrm{result}$,
  $\mathrm{tailCall}$) and the user-defined labels
  (let-bound names such as~$x$) are drawn from disjoint
  alphabets, so no collision can arise.
}
A let-bound variable $x$ becomes the lexical reference
$x.\mathrm{result}$, projecting the application result
from the sibling property $x$ in the enclosing record.

An abstraction $\lambda x.\, M$ translates to a record with own
properties $\{\mathrm{argument}, \mathrm{result}\}$; we call this
the \emph{abstraction shape}.
In the application $\mathbf{let}$-binding rule, the name $x$
binds to the inherited record
$\{\mathcal{T}(V_1),\;\mathrm{argument} \mapsto \mathcal{T}(V_2)\}$,
and $\mathcal{T}(M)$ may reference $x.\mathrm{result}$ to obtain
the value of the application.
In the value $\mathbf{let}$-binding rule, $x$ wraps $\mathcal{T}(V)$
under a $\mathrm{result}$ child so that the uniform accessor
$x.\mathrm{result}$ yields $\mathcal{T}(V)$, the translated value
itself.
This extra wrapper is semantically inert but proof-relevant:
it lets the same $\mathrm{result}$-following argument used for
applications extend uniformly to value
$\mathbf{let}$-bindings in the adequacy proof
(Appendix~\ref{app:bohm-tree-proofs}).
In the tail-call rule, the fresh label $\mathrm{tailCall}$ serves
the same purpose, and $\mathrm{result}$ projects the answer.
In the application and tail-call cases, the inheritance that
performs the application is
encapsulated behind a named property, so its internal
structure (the $\mathrm{argument}$ and $\mathrm{result}$ labels
of the callee) does not leak to the enclosing scope.

The six rules above are a complete translation from the
$\lambda$-calculus (via ANF) to inheritance-calculus: every ANF
term has an image, and the construction is compositional.
Since the $\lambda$-calculus is Turing complete, so is
inheritance-calculus.%
\footnote{
  MIXINv2 separates symbol resolution (which resolves
  lexical references and qualified $\this$ references) from runtime
  mixin evaluation, but symbol resolution is performed
  just-in-time per symbol rather than ahead of time, so it does
  not provide static type checking in the traditional sense.
  Nevertheless, MIXINv2's resolution phase is stricter than the
  untyped inheritance-calculus.
  All examples in this paper are ported from MIXINv2
  and conform to both semantics.
  Elements added to satisfy the resolver are harmless
  for the algorithms presented here and may aid readability.
}
Turing completeness, however, says nothing about what the
translation preserves; the next subsection pins down
the semantic correspondence.
\fi

\subsection{\bilingual{B\"ohm Tree Correspondence}{Böhm 树对应}}
\label{sec:bohm-tree}

\ifInheritanceChinese
翻译 $\mathcal{T}$ 将 $\lambda$-演算嵌入继承演算的一个子语言中。直觉上,mixin 树以附加的语义结构增广了 AST;对于翻译后的 $\lambda$-项,这一附加结构的投影与 Böhm 树同构。若从根出发追踪有限次 $\mathrm{result}$ 投影后到达一个同时拥有 $\mathrm{argument}$ 和 $\mathrm{result}$ 的节点(即 $\mathcal{T}(\lambda x.\, M)$ 产生的\emph{抽象形状}),则称 mixin 树 $T$ \emph{收敛},记为 $T{\Downarrow}$。翻译保持并反映收敛性(充分性)。记 $\mathcal{T}(M) \approx_\lambda \mathcal{T}(N)$ 当对所有封闭的 $\lambda$-演算上下文 $C[\cdot]$ 均有 $\mathcal{T}(C[M]){\Downarrow} \Leftrightarrow \mathcal{T}(C[N]){\Downarrow}$;则 $\approx_\lambda$ 恰好与 Böhm 树等价~\cite{barendregt1984-lambda-calculus} 重合。形式化陈述和证明见附录~\ref{app:bohm-tree-proofs}。注意 $\approx_\lambda$ 仅在 $\lambda$-演算上下文上量化;继承演算的上下文区分能力严格更强,如第~\ref{sec:asymmetry} 节所证。以上对应关系在首次迭代近似 $T_P\!\uparrow\!1$ 处成立;我们现在考察在完整最小不动点处会发生什么。

在 $T_P\!\uparrow\!1$ 处,每个方程被求值一次,可重入调用返回 $\bot$;这对应于 $\beta$-规约。在完整的 $\mathrm{lfp}(T_P)$ 处,可重入引用改而解析为幂集格上的最小不动点:迭代在\emph{容器}层级计算不动点,求出满足 $S = F(S)$ 的集合 $S$。这正是递归 Datalog 在循环图上收敛的机制(附录~\ref{app:datalog-encoding},第~\ref{sec:recursive-rules} 节);既非不动点组合子(如 $Y$),也非 Haskell 的 \texttt{mfix}(\texttt{RecursiveDo})能够表达它:两者都在项层级打结,而非在语义所累积的集合上打结。破坏 Böhm 树对应关系的这种\emph{幂集}更多确定性(图~\ref{fig:calculi-matrix})的最小见证例是 $\mathbf{let}\; x = x \;\mathbf{in}\; x$:它在 $\beta$-规约下发散,但其翻译 $\{x \mapsto \{\mathrm{result} \mapsto \{x.\mathrm{result}\}\},\; \mathrm{result} \mapsto \{x.\mathrm{result}\}\}$ 在 $\mathrm{lfp}(T_P)$ 下稳定,tabling 求值检测到循环并返回确定结果 $\properties(\mathrm{root}) = \{x, \mathrm{result}\}$。因此本小节的对应关系是 $T_P\!\uparrow\!1$ 独有的性质。下一小节考察继承演算是否也严格更具表达力。
\else
The translation $\mathcal{T}$ embeds the $\lambda$-calculus
into a sublanguage of inheritance-calculus.
Intuitively, a mixin tree augments the AST with
additional semantic structure; for translated
$\lambda$-terms, the projection of this additional
structure is isomorphic to the B\"ohm tree.
A mixin tree $T$ \emph{converges}, written $T{\Downarrow}$,
if following finitely many $\mathrm{result}$ projections from
the root reaches a node with both $\mathrm{argument}$ and
$\mathrm{result}$ among its properties, that is, the \emph{abstraction shape}
produced by $\mathcal{T}(\lambda x.\, M)$.
The translation preserves and reflects convergence
(adequacy).
Write $\mathcal{T}(M) \approx_\lambda \mathcal{T}(N)$
when $\mathcal{T}(C[M]){\Downarrow} \Leftrightarrow
\mathcal{T}(C[N]){\Downarrow}$ for every closing
$\lambda$-calculus context $C[\cdot]$;
then $\approx_\lambda$ coincides with B\"ohm tree
equivalence~\cite{barendregt1984-lambda-calculus}.
Formal statements and proofs are in
Appendix~\ref{app:bohm-tree-proofs}.
Note that $\approx_\lambda$ quantifies over $\lambda$-calculus
contexts only; inheritance-calculus contexts are strictly more
discriminating, as Section~\ref{sec:asymmetry} establishes.
The correspondence above holds at the first-iteration
approximation $T_P\!\uparrow\!1$; we now ask what happens
at the full least fixed point.

At $T_P\!\uparrow\!1$, each equation is evaluated once and
re-entrant calls return~$\bot$; this corresponds to
$\beta$-reduction.
At the full $\mathrm{lfp}(T_P)$, re-entrant references instead
resolve to least fixed points over the powerset lattice: the
iteration computes fixed points at the \emph{container} level,
finding a set~$S$ such that $S = F(S)$.
This is the mechanism by which recursive Datalog converges on
cyclic graphs (Appendix~\ref{app:datalog-encoding},
Section~\ref{sec:recursive-rules}), and neither a fixed-point
combinator such as~$Y$ nor Haskell's \texttt{mfix}
(\texttt{RecursiveDo}) expresses it: both tie the knot at the
term level, not over the sets the semantics accumulates.
The smallest witness that this \emph{powerset} more-definedness
(Figure~\ref{fig:calculi-matrix}) breaks the B\"ohm-tree
correspondence is
$\mathbf{let}\; x = x \;\mathbf{in}\; x$:
it diverges under $\beta$-reduction, but its translation
$\{x \mapsto \{\mathrm{result} \mapsto \{x.\mathrm{result}\}\},\;
\mathrm{result} \mapsto \{x.\mathrm{result}\}\}$
stabilizes under $\mathrm{lfp}(T_P)$, where tabled evaluation
detects the cycle and returns the definite result
$\properties(\mathrm{root}) = \{x, \mathrm{result}\}$.
The correspondence of this subsection is therefore a property
of $T_P\!\uparrow\!1$ alone.
The next subsection examines whether inheritance-calculus is
also strictly more expressive.
\fi

\subsection{\bilingual{Expressive Asymmetry}{表达力不对称}}
\label{sec:asymmetry}
\label{sec:formal-separation}
\ifInheritanceChinese
两个 Böhm 树等价的 $\lambda$-项,在 $\lambda$-演算中因此无法区分,可以被一个混合了 $\lambda$-演算与继承演算构造子的上下文分离;该上下文通过开放扩展在已有定义上继承新的可观测投影。根据 Felleisen 的表达力判据~\cite{felleisen1991-expressive-power},继承演算严格比惰性 $\lambda$-演算更具表达力(定理~\ref{thm:expressive-asymmetry}):$\lambda$-演算可以通过 ANF 转换和 $\mathcal{T}$ 宏表达于继承演算中(引理~\ref{lem:anf-equiv},定理~\ref{thm:forward-macro}),但惰性 $\lambda$-演算无法宏表达继承(定理~\ref{thm:cartesian-nonexpressibility})。形式化定义、构造和证明见附录~\ref{app:expressiveness};该分离构造本身是表达式问题的一个微型实例,它在不修改已有定义的情况下添加了一个新的可观测投影。
\else
Two $\lambda$-terms that are B\"ohm-tree equivalent,
and therefore indistinguishable in the $\lambda$-calculus,
can be
separated by a context that mixes $\lambda$-calculus and
inheritance-calculus constructors, inheriting new observable
projections onto an existing definition via open extension.
By Felleisen's expressiveness
criterion~\cite{felleisen1991-expressive-power}, inheritance-calculus is
strictly more expressive than the lazy $\lambda$-calculus
(Theorem~\ref{thm:expressive-asymmetry}):
the $\lambda$-calculus can be macro-expressed in
inheritance-calculus via ANF conversion
and~$\mathcal{T}$
(Lemma~\ref{lem:anf-equiv}, Theorem~\ref{thm:forward-macro}),
but the lazy $\lambda$-calculus cannot macro-express inheritance
(Theorem~\ref{thm:cartesian-nonexpressibility}).
The formal definitions, constructions, and proofs are given in
Appendix~\ref{app:expressiveness};
the separating construction is itself a miniature instance of
the Expression Problem, adding a new observable projection to
an existing definition without modifying it.
\fi

\section{\bilingual{Case Study: Nat Arithmetic}{案例研究：Nat(自然数)算术}}
\label{sec:case-study}
\label{sec:expression-problem}

\ifInheritanceChinese
我们在继承演算中实现了布尔逻辑与自然数算术,包括一元与二进制两种形式。
本节详细追踪一元 Nat 的情形:其数据表示、加法与相等判断。
该实现遵循普通的 Gang-of-Four 设计模式~\cite{gamma1994-design-patterns},采用声明式面向对象风格,将每项关注点沿两个轴分解为 mixin:操作轴与情形轴。由于演算本身没有方法,所有名称均为名词或形容词:UpperCamelCase 名称扮演模块、类型、数据构造子与方法的角色,lowerCamelCase 名称扮演字段与参数的角色。
\else
We implemented boolean logic and natural number arithmetic, both unary
and binary, in inheritance-calculus.
This section traces the unary Nat case in detail: its data representation,
addition, and equality.
The implementation follows ordinary Gang-of-Four design
patterns~\cite{gamma1994-design-patterns} in a declarative object-oriented
style, with each concern split into mixins along two axes: operations and
cases. Because the calculus has no methods of its own, every name is a noun or
adjective: UpperCamelCase names play the roles of modules, types, data
constructors, and methods, and lowerCamelCase names the roles of fields and
parameters.
\fi

\paragraph{NatData}
\ifInheritanceChinese
mixin $\mathrm{NatData}$ 是自然数接口及其工厂。$\mathrm{NatFactory}$ 是一个工厂,其抽象产品类型被别名为 $\mathrm{Nat}$:
\else
The mixin $\mathrm{NatData}$ is the natural-number interface and its factory.
$\mathrm{NatFactory}$ is a factory whose abstract product type is aliased as $\mathrm{Nat}$:
\fi
\[
  \mathrm{NatData} \mapsto \left\{
    \begin{aligned}
      &\mathrm{NatFactory} \mapsto \{\mathrm{Product} \mapsto \{\}\},\\
      &\mathrm{Nat} \mapsto \{\mathrm{NatFactory}.\mathrm{Product}\}
  \end{aligned}\right\}
\]

\paragraph{NatDataZero}
\ifInheritanceChinese
每个数据构造子都是一个独立的实现 mixin,继承 $\mathrm{NatData}$。一个 $\mathrm{Zero}$ 值继承 $\mathrm{Product}$:
\else
Each data constructor is a separate implementation mixin that inherits
$\mathrm{NatData}$.
A $\mathrm{Zero}$ value inherits $\mathrm{Product}$:
\fi
\[
  \mathrm{NatDataZero} \mapsto \left\{
    \begin{aligned}
      &\mathrm{NatData},\\
      &\mathrm{NatFactory} \mapsto \{\mathrm{Zero} \mapsto \{\qualifiedthis{\mathrm{NatFactory}}.\mathrm{Product}\}\}
  \end{aligned}\right\}
\]

\paragraph{NatDataSuccessor}
\ifInheritanceChinese
一个 $\mathrm{Successor}$ 值继承 $\mathrm{Product}$,并暴露一个类型为 $\mathrm{Product}$ 的 $\mathrm{predecessor}$ 属性:
\else
A $\mathrm{Successor}$ value inherits $\mathrm{Product}$ and
exposes a $\mathrm{predecessor}$ property whose type is $\mathrm{Product}$:
\fi
\[
  \begin{aligned}
    &\mathrm{NatDataSuccessor} \mapsto \\
    &\left\{
    \begin{aligned}
      &\mathrm{NatData},\\
      &\mathrm{NatFactory} \mapsto \\
      &\left\{\mathrm{Successor} \mapsto \left\{
        \begin{aligned}
          &\qualifiedthis{\mathrm{NatDataSuccessor}}.\mathrm{NatFactory}.\mathrm{Product},\\
          &\mathrm{predecessor} \mapsto \{\qualifiedthis{\mathrm{NatDataSuccessor}}.\mathrm{NatFactory}.\mathrm{Product}\}
      \end{aligned}\right\}\right\}
  \end{aligned}\right\}
  \end{aligned}
\]
\ifInheritanceChinese
与 $\mathrm{NatDataZero}$ 合在一起,这以开放的方式完成了自然数的 Peano 编码,即函子 $F(X) = 1 + X$ 的初始代数:两个构造子作为独立的 mixin 加入,而非固定在单个封闭的数据类型定义中。
\else
Together with $\mathrm{NatDataZero}$, this completes the Peano encoding of
the natural numbers, the initial algebra of the functor $F(X) = 1 + X$,
in an open way: the two constructors are added as independent mixins
rather than fixed in a single closed datatype definition.
\fi

\paragraph{NatPlus}
\ifInheritanceChinese
加法操作沿情形轴分解为一个\emph{接口} mixin $\mathrm{NatPlus}$,该接口 mixin 保留关注点的名称并在抽象 $\mathrm{Product}$ 上声明命令,以及每个情形对应一个\emph{实现} mixin,实现 mixin 在名称后附加构造子名称(即下文的 $\mathrm{NatPlusZero}$ 与 $\mathrm{NatPlusSuccessor}$)。$\mathrm{NatPlus}$ 继承 $\mathrm{NatData}$。面向对象方法会用动词命名的地方,命令模式~\cite{gamma1994-design-patterns} 将其具化为以名词命名的命令,因此加法方法被命名为 $\mathrm{Plus}$ 而非 $\mathrm{Add}$,并通过投影属性 $\mathrm{sum}$ 暴露其结果;这里已指定类型,但尚未计算:
\else
The addition operation is split along the case axis into an \emph{interface}
mixin $\mathrm{NatPlus}$, which keeps the concern's name and declares the command
on the abstract $\mathrm{Product}$, and one \emph{implementation} mixin per case,
which appends the constructor name ($\mathrm{NatPlusZero}$ and
$\mathrm{NatPlusSuccessor}$ below). $\mathrm{NatPlus}$ inherits
$\mathrm{NatData}$. Where an
object-oriented method would be named with a verb, the command
pattern~\cite{gamma1994-design-patterns} reifies it as a noun-named command, so
the addition method is named $\mathrm{Plus}$ rather than $\mathrm{Add}$, and it
exposes its result through the projection property $\mathrm{sum}$, here typed but
not yet computed:
\fi
\[
  \mathrm{NatPlus} \mapsto \left\{
    \begin{aligned}
      &\mathrm{NatData},\\
      &\mathrm{NatFactory} \mapsto \left\{\mathrm{Product} \mapsto \left\{\mathrm{Plus} \mapsto \left\{\mathrm{sum} \mapsto \{\mathrm{Product}\}\right\}\right\}\right\}
  \end{aligned}\right\}
\]

\paragraph{NatPlusZero}
\ifInheritanceChinese
基础情形为 $0 + m = m$。这里 $\mathrm{NatPlusZero}$ 充当一个导入其他模块的模块:它继承 $\mathrm{Zero}$ 数据与 $\mathrm{Plus}$ 接口,并直接返回被加数:
\else
The base case is $0 + m = m$. Here $\mathrm{NatPlusZero}$ acts as a module that
imports other modules: it inherits the $\mathrm{Zero}$ data and the $\mathrm{Plus}$
interface, and returns the addend directly:
\fi
\[
  \begin{aligned}
    &\mathrm{NatPlusZero} \mapsto \\
    &\left\{
    \begin{aligned}
      &\mathrm{NatDataZero},\; \mathrm{NatPlus},\\
      &\mathrm{NatFactory} \mapsto \left\{\mathrm{Zero} \mapsto \left\{\mathrm{Plus} \mapsto \left\{
          \begin{aligned}
            &\mathrm{addend} \mapsto \{\qualifiedthis{\mathrm{NatPlusZero}}.\mathrm{NatFactory}.\mathrm{Product}\},\\
            &\mathrm{sum} \mapsto \{\mathrm{addend}\}
      \end{aligned}\right\}\right\}\right\}
  \end{aligned}\right\}
  \end{aligned}
\]

\paragraph{NatPlusSuccessor}
\ifInheritanceChinese
递归情形将 $\mathrm{S}(n_0) + m$ 规约为 $n_0 + \mathrm{S}(m)$,方法是以递增后的被加数委托给 $n_0$ 的 $\mathrm{Plus}$:
\else
The recursive case reduces $\mathrm{S}(n_0) + m$ to $n_0 + \mathrm{S}(m)$
by delegating to $n_0$'s $\mathrm{Plus}$ with an incremented addend:
\fi
\[
  \begin{aligned}
    &\mathrm{NatPlusSuccessor} \mapsto \\
    &\left\{
    \begin{aligned}
      &\mathrm{NatDataSuccessor},\; \mathrm{NatPlus},\\
      &\mathrm{NatFactory} \mapsto \\
      &\left\{\mathrm{Successor} \mapsto \left\{\mathrm{Plus} \mapsto \left\{
              \begin{aligned}
                &\mathrm{addend} \mapsto \{\qualifiedthis{\mathrm{NatFactory}}.\mathrm{Product}\},\\
                &\mathrm{\_increasedAddend} \mapsto  \\
                &\left\{
                  \begin{aligned}
                    &\mathrm{Successor},\\
                    &\mathrm{predecessor} \mapsto \{\mathrm{addend}\}
                \end{aligned}\right\},\\
                &\mathrm{\_recursiveAddition} \mapsto \\
                &\left\{
                  \begin{aligned}
                    &\qualifiedthis{\mathrm{Successor}}.\mathrm{predecessor}.\mathrm{Plus},\\
                    &\mathrm{addend} \mapsto \{\mathrm{\_increasedAddend}\}
                \end{aligned}\right\},\\
                &\mathrm{sum} \mapsto \{\mathrm{\_recursiveAddition}.\mathrm{sum}\}
          \end{aligned}\right\}\right\}\right\}
  \end{aligned}\right\}
  \end{aligned}
\]

\paragraph{NatVisitor}
\ifInheritanceChinese
相等判断需要对数据构造子进行情形分析;我们使用访问者模式~\cite{gamma1994-design-patterns}。$\mathrm{NatVisitor}$ 向 $\mathrm{Product}$ 添加一个 $\mathrm{Acceptance}$ 命令,即访问者模式的 $\mathrm{accept}$ 方法的名词形式。mixin $\mathrm{NatVisitor}$ 将其声明为空,留下 $\mathrm{VisitorMap}$ 与 $\mathrm{Accepted}$ 属性由各构造子的实现来填充:
\else
Equality requires case analysis on the data constructor; we use the
visitor pattern~\cite{gamma1994-design-patterns}.
$\mathrm{NatVisitor}$ adds an $\mathrm{Acceptance}$ command to $\mathrm{Product}$,
the noun form of the visitor pattern's $\mathrm{accept}$ method.
The mixin $\mathrm{NatVisitor}$ declares it empty, leaving a $\mathrm{VisitorMap}$ and an
$\mathrm{Accepted}$ property for the per-constructor implementations to fill in:
\fi
\[
  \begin{aligned}
    &\mathrm{NatVisitor} \mapsto \\
    &\left\{
    \begin{aligned}
      &\mathrm{NatData},\\
      &\mathrm{NatFactory} \mapsto \left\{\mathrm{Product} \mapsto \{\mathrm{Acceptance} \mapsto \{\mathrm{VisitorMap} \mapsto \{\},\; \mathrm{Accepted} \mapsto \{\}\}\}\right\}
  \end{aligned}\right\}
  \end{aligned}
\]

\paragraph{NatVisitorZero}
\ifInheritanceChinese
$\mathrm{Zero}$ 的实现选择 $\mathrm{ZeroVisitor}$ 分支。这里的访问者模式与~\cite{gamma1994-design-patterns} 中的表述不同,那里数据构造子由 $\mathrm{accept}$ 接口固定;此处通过组合允许添加新的构造子,因为每个构造子作为独立的 mixin 贡献其自身的分支:
\else
The $\mathrm{Zero}$ implementation selects the $\mathrm{ZeroVisitor}$ branch.
The visitor pattern here, unlike its formulation in
\cite{gamma1994-design-patterns} where the data constructors are fixed by the
$\mathrm{accept}$ interface, admits new constructors by composition, since each
constructor contributes its own branch as a separate mixin:
\fi
\[
  \begin{aligned}
    &\mathrm{NatVisitorZero} \mapsto \\
    &\left\{
    \begin{aligned}
      &\mathrm{NatDataZero},\; \mathrm{NatVisitor},\\
      &\mathrm{NatFactory} \mapsto \left\{\mathrm{Zero} \mapsto \left\{\mathrm{Acceptance} \mapsto \left\{
                \begin{aligned}
                  &\mathrm{VisitorMap} \mapsto \{\mathrm{ZeroVisitor} \mapsto \{\}\},\\
                  &\mathrm{Accepted} \mapsto \{\mathrm{VisitorMap}.\mathrm{ZeroVisitor}\}
              \end{aligned}\right\}\right\}\right\}
  \end{aligned}\right\}
  \end{aligned}
\]

\paragraph{NatVisitorSuccessor}
\ifInheritanceChinese
$\mathrm{Successor}$ 的实现选择 $\mathrm{SuccessorVisitor}$ 分支:
\else
The $\mathrm{Successor}$ implementation selects the $\mathrm{SuccessorVisitor}$ branch:
\fi
\[
  \begin{aligned}
    &\mathrm{NatVisitorSuccessor} \mapsto \\
    &\left\{
    \begin{aligned}
      &\mathrm{NatDataSuccessor},\; \mathrm{NatVisitor},\\
      &\mathrm{NatFactory} \mapsto \\
      &\left\{\mathrm{Successor} \mapsto \left\{
            \begin{aligned}
              &\mathrm{predecessor} \mapsto \{\qualifiedthis{\mathrm{NatFactory}}.\mathrm{Product}\},\\
              &\mathrm{Acceptance} \mapsto \left\{
                \begin{aligned}
                  &\mathrm{VisitorMap} \mapsto \{\mathrm{SuccessorVisitor} \mapsto \{\}\},\\
                  &\mathrm{Accepted} \mapsto \{\mathrm{VisitorMap}.\mathrm{SuccessorVisitor}\}
              \end{aligned}\right\}
          \end{aligned}\right\}\right\}
  \end{aligned}\right\}
  \end{aligned}
\]

\paragraph{BooleanData}
\ifInheritanceChinese
$\mathrm{BooleanData}$ 定义输出的域:
\else
$\mathrm{BooleanData}$ defines the output sort:
\fi
\[
  \mathrm{BooleanData} \mapsto \left\{
    \begin{aligned}
      &\mathrm{BooleanFactory} \mapsto \left\{
        \begin{aligned}
          &\mathrm{Product} \mapsto \{\},\\
          &\mathrm{True} \mapsto \{\mathrm{Product}\},\\
          &\mathrm{False} \mapsto \{\mathrm{Product}\}
      \end{aligned}\right\},\\
      &\mathrm{Boolean} \mapsto \{\mathrm{BooleanFactory}.\mathrm{Product}\}
  \end{aligned}\right\}
\]
\paragraph{NatEquality}
\label{sec:nat-equality}
\ifInheritanceChinese
为了验证 $2 + 3 = 5$,我们需要对 Nat 值进行相等判断。mixin $\mathrm{NatEquality}$ 导入 $\mathrm{NatVisitor}$ 与 $\mathrm{BooleanData}$,并在 $\mathrm{Product}$ 上声明一个 $\mathrm{Equal}$ 命令,由各构造子的实现提供。由于 $\mathrm{NatEquality}$ 继承 $\mathrm{BooleanData}$,限定 $\this$ $\qualifiedthis{\mathrm{NatEquality}}$ 可以到达 $\mathrm{Boolean}$,因此 $\mathrm{Equal}$ 以类型 $\qualifiedthis{\mathrm{NatEquality}}.\mathrm{Boolean}$ 声明其结果 $\mathrm{equal}$:
\else
To test that $2 + 3 = 5$, we needed equality on Nat values.
The mixin $\mathrm{NatEquality}$ imports $\mathrm{NatVisitor}$ and $\mathrm{BooleanData}$,
and declares an $\mathrm{Equal}$ command on $\mathrm{Product}$ that the per-constructor implementations provide.
Because $\mathrm{NatEquality}$ inherits $\mathrm{BooleanData}$,
the qualified this $\qualifiedthis{\mathrm{NatEquality}}$ can reach $\mathrm{Boolean}$,
so $\mathrm{Equal}$ declares its result
$\mathrm{equal}$ with the type $\qualifiedthis{\mathrm{NatEquality}}.\mathrm{Boolean}$:
\fi
\[
  \begin{aligned}
    &\mathrm{NatEquality} \mapsto \\
    &\left\{
    \begin{aligned}
      &\mathrm{NatVisitor},\; \mathrm{BooleanData},\\
      &\mathrm{NatFactory} \mapsto \left\{\mathrm{Product} \mapsto \left\{\mathrm{Equal} \mapsto \left\{
          \begin{aligned}
            &\mathrm{other} \mapsto \{\mathrm{Product}\},\\
            &\mathrm{equal} \mapsto \{\qualifiedthis{\mathrm{NatEquality}}.\mathrm{Boolean}\}
      \end{aligned}\right\}\right\}\right\}
  \end{aligned}\right\}
  \end{aligned}
\]

\paragraph{NatEqualityZero}
\ifInheritanceChinese
$\mathrm{Zero}$ 只与另一个 $\mathrm{Zero}$ 相等。这里的继承执行情形分析:继承 $\mathrm{other}.\mathrm{Acceptance}$ 对 $\mathrm{other}$ 是由 $\mathrm{Zero}$ 还是 $\mathrm{Successor}$ 构建进行分派,在 $\mathrm{ZeroVisitor}$ 分支返回 $\mathrm{True}$,否则返回 $\mathrm{False}$,即 $\mathrm{equal} = \{ \mathrm{True} \text{ 若 } \mathrm{other} = \mathrm{Zero},\; \mathrm{False} \text{ 若 } \mathrm{other} = \mathrm{Successor} \}$。$\mathrm{True}$ 与 $\mathrm{False}$ 均为 $\mathrm{BooleanFactory}$ 的数据构造子;该工厂来自 $\mathrm{BooleanData}$,通过继承的 $\mathrm{NatEquality}$ 间接导入,并以 $\qualifiedthis{\mathrm{NatEqualityZero}}.\mathrm{BooleanFactory}$ 到达。同一继承机制也提供了结果类型:在 $\mathrm{Accepted}$ 中,$\mathrm{equal}$ 继承 $\qualifiedthis{\mathrm{NatEqualityZero}}.\mathrm{Boolean}$,从而声明其类型为 $\mathrm{Boolean}$:
\else
$\mathrm{Zero}$ is equal only to another $\mathrm{Zero}$.
Inheriting here performs case analysis:
inheriting $\mathrm{other}.\mathrm{Acceptance}$ dispatches on whether $\mathrm{other}$
was built by $\mathrm{Zero}$ or $\mathrm{Successor}$, returning $\mathrm{True}$
on the $\mathrm{ZeroVisitor}$ branch and $\mathrm{False}$ otherwise, that is,
$\mathrm{equal} = \{ \mathrm{True} \text{ if } \mathrm{other} = \mathrm{Zero},\; \mathrm{False} \text{ if } \mathrm{other} = \mathrm{Successor} \}$.
Both $\mathrm{True}$ and $\mathrm{False}$ are data constructors of
$\mathrm{BooleanFactory}$, which comes from $\mathrm{BooleanData}$, imported
indirectly through the inherited $\mathrm{NatEquality}$, and is reached as
$\qualifiedthis{\mathrm{NatEqualityZero}}.\mathrm{BooleanFactory}$.
The same inheritance mechanism also supplies the result type: in
$\mathrm{Accepted}$, $\mathrm{equal}$ inherits
$\qualifiedthis{\mathrm{NatEqualityZero}}.\mathrm{Boolean}$, declaring its type to
be $\mathrm{Boolean}$:
\fi
{\small
  \[
    \begin{aligned}
      &\mathrm{NatEqualityZero} \mapsto \\
      &\left\{
      \begin{aligned}
        &\mathrm{NatVisitorZero},\; \mathrm{NatEquality},\\
        &\mathrm{NatFactory} \mapsto \\
        &\left\{
        \begin{aligned}
        &\mathrm{Zero} \mapsto \\
        &\left\{
        \begin{aligned}
        &\mathrm{Equal} \mapsto \\
        &\left\{
                  \begin{aligned}
                    &\mathrm{other} \mapsto \{\qualifiedthis{\mathrm{NatEqualityZero}}.\mathrm{NatFactory}.\mathrm{Product}\},\\
                    &\mathrm{\_OtherAcceptance} \mapsto \\
                    &\left\{
                      \begin{aligned}
                        &\mathrm{other}.\mathrm{Acceptance},\\
                        &\mathrm{VisitorMap} \mapsto \\
                        &\left\{
                          \begin{aligned}
                            &\mathrm{ZeroVisitor} \mapsto \{\mathrm{equal} \mapsto \{\qualifiedthis{\mathrm{NatEqualityZero}}.\mathrm{BooleanFactory}.\mathrm{True}\}\},\\
                            &\mathrm{SuccessorVisitor} \mapsto \{\mathrm{equal} \mapsto \{\qualifiedthis{\mathrm{NatEqualityZero}}.\mathrm{BooleanFactory}.\mathrm{False}\}\}
                        \end{aligned}\right\},\\
                        &\mathrm{Accepted} \mapsto \{\mathrm{equal} \mapsto \{\qualifiedthis{\mathrm{NatEqualityZero}}.\mathrm{Boolean}\}\}
                    \end{aligned}\right\},\\
                    &\mathrm{equal} \mapsto \{\mathrm{\_OtherAcceptance}.\mathrm{Accepted}.\mathrm{equal}\}
                \end{aligned}\right\}
        \end{aligned}\right\}
        \end{aligned}\right\}
    \end{aligned}\right\}
    \end{aligned}
\]}

\paragraph{NatEqualitySuccessor}
\ifInheritanceChinese
$\mathrm{Successor}$ 只与另一个具有相等前驱的 $\mathrm{Successor}$ 相等,因此其实现通过 $\qualifiedthis{\mathrm{Successor}}.\mathrm{predecessor}.\mathrm{Equal}$ 进行递归:
\else
$\mathrm{Successor}$ is equal only to another $\mathrm{Successor}$ with an equal
predecessor, so its implementation recurses through
$\qualifiedthis{\mathrm{Successor}}.\mathrm{predecessor}.\mathrm{Equal}$:
\fi
{\small
  \[
    \begin{aligned}
      &\mathrm{NatEqualitySuccessor} \mapsto \\
      &\left\{
      \begin{aligned}
        &\mathrm{NatVisitorSuccessor},\; \mathrm{NatEquality},\\
        &\mathrm{NatFactory} \mapsto \\
        &\left\{
        \begin{aligned}
        &\mathrm{Successor} \mapsto \\
        &\left\{
        \begin{aligned}
        &\mathrm{Equal} \mapsto \\
        &\left\{
                  \begin{aligned}
                    &\mathrm{other} \mapsto \left\{
                      \begin{aligned}
                        &\qualifiedthis{\mathrm{NatEqualitySuccessor}}.\mathrm{NatFactory}.\mathrm{Product},\\
                        &\mathrm{predecessor} \mapsto \{\qualifiedthis{\mathrm{NatEqualitySuccessor}}.\mathrm{NatFactory}.\mathrm{Product}\}
                    \end{aligned}\right\},\\
                    &\mathrm{\_recursiveEquality} \mapsto \left\{
                      \begin{aligned}
                        &\qualifiedthis{\mathrm{Successor}}.\mathrm{predecessor}.\mathrm{Equal},\\
                        &\mathrm{other} \mapsto \{\qualifiedthis{\mathrm{Equal}}.\mathrm{other}.\mathrm{predecessor}\}
                    \end{aligned}\right\},\\
                    &\mathrm{\_OtherAcceptance} \mapsto \\
                    &\left\{
                      \begin{aligned}
                        &\mathrm{other}.\mathrm{Acceptance},\\
                        &\mathrm{VisitorMap} \mapsto \\
                        &\left\{
                          \begin{aligned}
                            &\mathrm{ZeroVisitor} \mapsto \{\mathrm{equal} \mapsto \{\qualifiedthis{\mathrm{NatEqualitySuccessor}}.\mathrm{BooleanFactory}.\mathrm{False}\}\},\\
                            &\mathrm{SuccessorVisitor} \mapsto \{\mathrm{equal} \mapsto \{\mathrm{\_recursiveEquality}.\mathrm{equal}\}\}
                        \end{aligned}\right\},\\
                        &\mathrm{Accepted} \mapsto \{\mathrm{equal} \mapsto \{\qualifiedthis{\mathrm{NatEqualitySuccessor}}.\mathrm{Boolean}\}\}
                    \end{aligned}\right\},\\
                    &\mathrm{equal} \mapsto \{\mathrm{\_OtherAcceptance}.\mathrm{Accepted}.\mathrm{equal}\}
                \end{aligned}\right\}
        \end{aligned}\right\}
        \end{aligned}\right\}
    \end{aligned}\right\}
    \end{aligned}
\]}
%% === 审稿人备忘：关于代码"冗长"的回应 ===
%%
%% NatEquality 的定义看起来比等价的 Haskell 或 ML 模式匹配更长，
%% 但这种冗长与 point-free / 组合子风格的"简洁"形成的对比是表面的。
%% 需要区分两种"长"：
%%
%% 1. 自描述式的长：每个中间步骤都有名字（如 _recursiveEquality、
%%    OtherAcceptance），步骤写得详细。名字本身是零成本的语义锚点。
%% 2. 组合子膨胀式的长：如 SKI bracket abstraction 的指数级膨胀，
%%    在玩具例子里看起来精炼，但在真实程序中复杂度是 λ-calculus 的
%%    多项式甚至指数倍，而且没有涌现出新的计算模式（论文 Related Work
%%    已指出这一点）。
%%
%% Inheritance-calculus 的"冗长"属于第一种，而且是结构上的必要代价：
%% 每个有名字的中间步骤既是代码，也是 mixin 树上可路径寻址、可扩展的
%% 节点。这正是 open extension 的前提——无法扩展匿名的子表达式。
%% 如果加语法糖把名字隐藏（如同 M-表达式之于 S-表达式），反而会
%% 破坏这一性质。Lisp 社区的历史经验表明，S-表达式的"冗长"最终
%% 被证明是优势（代码即数据，宏系统直接操作统一结构），M-表达式
%% 的语法糖从未被真正需要。
%%
%% 在 AI 辅助编程的语境下，这种显式命名的风格尤其有利：
%% _recursiveEquality 这样的名字对 LLM 而言是自带的 chain-of-thought，
%% 比起让 AI 猜测 point-free 管道中第 n 个组合子的意图，
%% 每个自描述的步骤都降低了生成和理解的认知负荷。
%% 考虑到未来越来越多的代码将由 AI 生成和维护，
%% 多写几个 token 的显式名称不是成本，而是投资。
%%
%% 当然，如果确实需要更紧凑的表面语法，为 inheritance-calculus 加一个
%% ANF 预处理器（自动为匿名子表达式生成名字）并不困难，
%% 正如 ANF 变换本身是机械化的。但这是表面语法的选择，
%% 不影响演算本身的语义。
%% ================================================================
\ifInheritanceChinese
通过继承来组合 $\mathrm{Plus}$ 与 $\mathrm{Equal}$ 的实现,无需修改任何 mixin,即可使所得的 Nat 值自动同时具备 $\mathrm{Plus}$ 与 $\mathrm{Equal}$。具体的数字常量存放在一个只继承两个数据构造子的共享 mixin 中:
\else
Composing the $\mathrm{Plus}$ and $\mathrm{Equal}$ implementations by
inheriting them, with no modifications to any mixin, gives the
resulting Nat values both $\mathrm{Plus}$ and $\mathrm{Equal}$
automatically.
The concrete numerals live in a shared mixin that inherits only the
two data constructors:
\fi
\[
  \mathrm{NatConstants} \mapsto \left\{
    \begin{aligned}
      &\mathrm{NatDataZero},\; \mathrm{NatDataSuccessor},\\
      &\mathrm{One} \mapsto \left\{
        \begin{aligned}
          &\qualifiedthis{\mathrm{NatConstants}}.\mathrm{NatFactory}.\mathrm{Successor},\\
          &\mathrm{predecessor} \mapsto \{\qualifiedthis{\mathrm{NatConstants}}.\mathrm{NatFactory}.\mathrm{Zero}\}
      \end{aligned}\right\},\\
      &\mathrm{Two} \mapsto \left\{
        \begin{aligned}
          &\qualifiedthis{\mathrm{NatConstants}}.\mathrm{NatFactory}.\mathrm{Successor},\\
          &\mathrm{predecessor} \mapsto \{\mathrm{One}\}
      \end{aligned}\right\},\\
      &\mathrm{Three} \mapsto \left\{
        \begin{aligned}
          &\qualifiedthis{\mathrm{NatConstants}}.\mathrm{NatFactory}.\mathrm{Successor},\\
          &\mathrm{predecessor} \mapsto \{\mathrm{Two}\}
      \end{aligned}\right\},\\
      &\mathrm{Four} \mapsto \left\{
        \begin{aligned}
          &\qualifiedthis{\mathrm{NatConstants}}.\mathrm{NatFactory}.\mathrm{Successor},\\
          &\mathrm{predecessor} \mapsto \{\mathrm{Three}\}
      \end{aligned}\right\},\\
      &\mathrm{Five} \mapsto \left\{
        \begin{aligned}
          &\qualifiedthis{\mathrm{NatConstants}}.\mathrm{NatFactory}.\mathrm{Successor},\\
          &\mathrm{predecessor} \mapsto \{\mathrm{Four}\}
      \end{aligned}\right\}
  \end{aligned}\right\}
\]
\ifInheritanceChinese
算术测试继承 $\mathrm{NatConstants}$ 以及它所需的操作:
\else
The arithmetic test inherits $\mathrm{NatConstants}$ together with the operations it needs:
\fi
\[
  \mathrm{Test} \mapsto \left\{
    \begin{aligned}
      &\mathrm{NatConstants},\; \mathrm{NatPlusZero},\; \mathrm{NatPlusSuccessor},\\
      &\mathrm{NatEqualityZero},\; \mathrm{NatEqualitySuccessor},\\
      &\mathrm{TwoPlusThree} \mapsto \left\{
        \begin{aligned}
          &\qualifiedthis{\mathrm{Test}}.\mathrm{Two}.\mathrm{Plus},\\
          &\mathrm{addend} \mapsto \{\qualifiedthis{\mathrm{Test}}.\mathrm{Three}\}
      \end{aligned}\right\},\\
      &\mathrm{FiveEqualTwoPlusThree} \mapsto \left\{
        \begin{aligned}
          &\qualifiedthis{\mathrm{Test}}.\mathrm{Five}.\mathrm{Equal},\\
          &\mathrm{other} \mapsto \{\mathrm{TwoPlusThree}.\mathrm{sum}\}
      \end{aligned}\right\}
  \end{aligned}\right\}
\]
\ifInheritanceChinese
路径 $\mathrm{Test}.\mathrm{FiveEqualTwoPlusThree}.\mathrm{equal}$ 继承 $\mathrm{True}$。%
\footnote{
  读取一个结果,需要借助 FFI 来打印,或者借助元语言观察以在该路径上调用 $\supers$。
}
没有任何现有 mixin 被修改。$\mathrm{NatEquality}$(一项新操作)与 $\mathrm{Boolean}$(一种新数据类型)各自只需要一个新 mixin:沿操作轴与情形轴均可扩展。
\else
The path $\mathrm{Test}.\mathrm{FiveEqualTwoPlusThree}.\mathrm{equal}$
inherits $\mathrm{True}$.%
\footnote{
  Reading a result requires either the FFI to print it or a
  metalanguage observation that invokes $\supers$ on the path.
}
No existing mixin was modified.
$\mathrm{NatEquality}$, a new operation, and $\mathrm{Boolean}$,
a new data type, each required only a new mixin: extensibility along both
the operation axis and the case axis.
\fi

\paragraph{\bilingual{Tree-level inheritance and return types}{树层次的继承与返回类型}}
\label{sec:discussion-expression-problem}
\ifInheritanceChinese
每个操作都是一个独立的 mixin,扩展工厂的子树。继承递归地合并共享同一标签的子树,因此工厂中的每个标签都从所有被继承的 mixin 获得全部操作。关键在于,这同样适用于返回类型:$\mathrm{NatPlus}$ 从 $\mathrm{NatFactory}$ 返回值,而 $\mathrm{NatEquality}$ 独立地向同一工厂添加 $\mathrm{Equal}$。同时继承两者,使得 $\mathrm{Plus}$ 返回的值自动携带 $\mathrm{Equal}$,无需修改任何一个 mixin。在对象代数~\cite{oliveira2012-object-algebras} 与 finally tagless 解释器~\cite{carette2009-finally-tagless} 中,独立定义的操作不会自动丰富彼此的返回类型;需要额外的样板代码或类型类机制。

多目标限定 $\this$(方程~\ref{eq:this})至关重要。在上述测试 mixin 中,$\mathrm{NatPlusSuccessor}$ 与 $\mathrm{NatEqualitySuccessor}$ 均继承 $\mathrm{NatDataSuccessor}$。当两个操作被组合时,所得的 $\mathrm{Successor}$ 通过两条独立路由继承 $\mathrm{NatDataSuccessor}$ 的模式。每个操作内部的引用 $\qualifiedthis{\mathrm{Successor}}.\mathrm{predecessor}$ 通过两条路由解析,由于继承是幂等的,所有路由贡献相同的属性。在 Scala 中,同样的模式,即两个独立定义的内部类扩展同一外部类的特质,会触发``冲突基类型''拒绝(附录~\ref{app:scala-multi-target}),因为 DOT 的 self 变量是单值的,无法同时通过多条继承路由解析。表达式问题在 Scala 中仍可使用对象代数或类似模式求解,但组合独立定义的操作本质上会产生上述对共享模式的多目标 $\this$ 解析。
\else
Each operation is an independent mixin that extends a factory's
subtree.
Inheritance recursively merges subtrees that share a label,
so every label in the factory acquires all operations from all
inherited mixins.
Crucially, this applies to return types: $\mathrm{NatPlus}$ returns
values from $\mathrm{NatFactory}$, and $\mathrm{NatEquality}$
independently adds $\mathrm{Equal}$ to the same factory.
Inheriting both causes the values returned by $\mathrm{Plus}$ to
carry $\mathrm{Equal}$ automatically, without modifying either mixin.
In object algebras~\cite{oliveira2012-object-algebras} and finally
tagless interpreters~\cite{carette2009-finally-tagless}, independently defined
operations do not automatically enrich each other's return types;
additional boilerplate or type-class machinery is required.

Multi-target qualified $\this$ (equation~\ref{eq:this}) is essential.
In the test mixin above, $\mathrm{NatPlusSuccessor}$ and
$\mathrm{NatEqualitySuccessor}$ both inherit
$\mathrm{NatDataSuccessor}$.
When the two operations are composed, the resulting
$\mathrm{Successor}$ inherits the $\mathrm{NatDataSuccessor}$ schema
through two independent routes.
The reference
$\qualifiedthis{\mathrm{Successor}}.\mathrm{predecessor}$
inside each operation resolves through both routes, and since
inheritance is idempotent, all routes contribute the same
properties.
In Scala, the same pattern, where two independently defined inner
classes extend the same outer class's trait, triggers a
``conflicting base types'' rejection
(Appendix~\ref{app:scala-multi-target}), because DOT's self variable
is single-valued and cannot resolve through multiple inheritance
routes simultaneously.
The Expression Problem can still be solved in Scala using
object algebras or similar patterns, but composing independently
defined operations inherently produces the multi-target $\this$
resolution to shared schemas just described.
\fi

\paragraph{\bilingual{Church-encoded Nats are tries}{Church 编码的 Nat 是字典树}}
\label{sec:case-study-cartesian}
\ifInheritanceChinese
Church 编码的 Nat 是一个记录,其结构映射了构建它所用的数据构造子路径:$\mathrm{Zero}$ 是一个平坦记录;$\mathrm{S}(\mathrm{Zero})$ 是一个带有 $\mathrm{predecessor}$ 属性、指向一个 $\mathrm{Zero}$ 记录的记录;依此类推。这恰好是字典树的形状。%
\footnote{
  一元 Nat 的字典树可以被视为一个稀疏链表,其深度等于它所代表的数值:在字典树并中,链表中的某些 $\mathrm{Successor}$ 节点指向 $\mathrm{Zero}$(叶节点),因此链表中并非每个位置都被占用。\anon[{} 补充材料]{{} MIXINv2 源码~\cite{mixin2025}}{} 中包含一个字典树深度为对数级的二进制自然数(BinNat)算术。
}
一个同时继承两个 Nat 值的作用域是一个字典树并。下面的测试将 $\mathrm{Plus}$ 应用于 $\{1,2\}$ 的字典树并,被加数为 $\{3,4\}$ 的字典树并,然后用 $\mathrm{Equal}$ 将结果与自身比较:
\else
A Church-encoded Nat is a record whose structure mirrors the
data-constructor path used to build it: $\mathrm{Zero}$ is a flat record;
$\mathrm{S}(\mathrm{Zero})$ is a record with a $\mathrm{predecessor}$
property pointing to a $\mathrm{Zero}$ record; and so on.
This is exactly the shape of a trie.%
\footnote{
  The unary Nat trie can be viewed as a sparse linked list
  whose depth equals the number it represents:
  in a trie union, some $\mathrm{Successor}$ nodes along
  the list point to $\mathrm{Zero}$ (leaf nodes),
  so not every position in the list is occupied.
  The\anon[{} supplementary material]{{} MIXINv2 source~\cite{mixin2025}}{} includes a binary natural number
  (BinNat) arithmetic whose trie depth is logarithmic.
}
A scope that simultaneously inherits two Nat values is a trie union.
The following test applies $\mathrm{Plus}$ to a trie union of $\{1,2\}$
with addend a trie union of $\{3,4\}$, then compares the result to itself
with $\mathrm{Equal}$:
\fi
\[
  \mathrm{CartesianTest} \mapsto \left\{
    \begin{aligned}
      &\mathrm{NatConstants},\; \mathrm{NatPlusZero},\; \mathrm{NatPlusSuccessor},\\
      &\mathrm{NatEqualityZero},\; \mathrm{NatEqualitySuccessor},\\
      &\mathrm{OneOrTwo} \mapsto \left\{\qualifiedthis{\mathrm{CartesianTest}}.\mathrm{One},\; \qualifiedthis{\mathrm{CartesianTest}}.\mathrm{Two}\right\},\\
      &\mathrm{ThreeOrFour} \mapsto \left\{\qualifiedthis{\mathrm{CartesianTest}}.\mathrm{Three},\; \qualifiedthis{\mathrm{CartesianTest}}.\mathrm{Four}\right\},\\
      &\mathrm{Result} \mapsto \left\{\mathrm{OneOrTwo}.\mathrm{Plus},\; \mathrm{addend} \mapsto \{\mathrm{ThreeOrFour}\}\right\},\\
      &\mathrm{Check} \mapsto \left\{\mathrm{Result}.\mathrm{sum}.\mathrm{Equal},\; \mathrm{other} \mapsto \{\mathrm{Result}.\mathrm{sum}\}\right\}
  \end{aligned}\right\}
\]
\ifInheritanceChinese
$\mathrm{Result}.\mathrm{sum}$ 求值为字典树 $\{4,5,6\}$,即所有两两之和的集合。$\mathrm{Equal}$ 随后对从 $\{4,5,6\} \times \{4,5,6\}$ 中取出的每一对进行应用;路径 $\mathrm{CartesianTest}.\mathrm{Check}.\mathrm{equal}$ 同时继承 $\mathrm{True}$ 与 $\mathrm{False}$。笛卡尔积直接来自语义方程。$\mathrm{OneOrTwo}$ 同时继承 $\mathrm{One}$ 与 $\mathrm{Two}$,因此 $\bases(\mathrm{OneOrTwo})$ 包含两条路径。当 $\mathrm{Result}$ 继承 $\mathrm{OneOrTwo}.\mathrm{Plus}$ 时,$\mathrm{Result}$ 的 $\bases$ 包含 $\mathrm{One}$ 与 $\mathrm{Two}$ 各自的 $\mathrm{Plus}$。每个 $\mathrm{Plus}$ 通过 $\this$ 解析其 $\mathrm{addend}$,到达 $\mathrm{ThreeOrFour}$;该值本身是 $\mathrm{Three}$ 与 $\mathrm{Four}$ 的字典树并,因此 $\mathrm{addend}$ 路径的 $\supers$ 包含两个值。由于 $\mathrm{Successor}.\mathrm{Plus}$ 以递增后的被加数递归委托给 $\mathrm{predecessor}.\mathrm{Plus}$,而 $\mathrm{predecessor}$ 也通过字典树并解析,递归在每个字典树节点处分支。因此结果 $\mathrm{sum}$ 收集了通过两个操作数的任意组合所能到达的全部路径,即笛卡尔积。

同样操作单个值的 $\mathrm{Plus}$ 与 $\mathrm{Equal}$ 文件,无需修改,即产生逻辑编程的关系语义:$\overrides$ 与 $\supers$ 收集所有继承路径,每个操作都在其上分布。
\else
$\mathrm{Result}.\mathrm{sum}$ evaluates to the trie $\{4,5,6\}$, the set of all pairwise sums.
$\mathrm{Equal}$ then applies to every pair drawn from $\{4,5,6\} \times \{4,5,6\}$;
the path $\mathrm{CartesianTest}.\mathrm{Check}.\mathrm{equal}$
inherits both $\mathrm{True}$ and $\mathrm{False}$.
The Cartesian product arises directly from the semantic equations.
$\mathrm{OneOrTwo}$ inherits both $\mathrm{One}$ and $\mathrm{Two}$,
so $\bases(\mathrm{OneOrTwo})$ contains two paths.
When $\mathrm{Result}$ inherits
$\mathrm{OneOrTwo}.\mathrm{Plus}$, the $\bases$ of
$\mathrm{Result}$ include the $\mathrm{Plus}$ of both
$\mathrm{One}$ and $\mathrm{Two}$.
Each $\mathrm{Plus}$ resolves its $\mathrm{addend}$ via
$\this$, reaching $\mathrm{ThreeOrFour}$, which is itself a
trie union of $\mathrm{Three}$ and $\mathrm{Four}$:
the $\supers$ of the $\mathrm{addend}$ path contain both values.
Since $\mathrm{Successor}.\mathrm{Plus}$ recursively delegates to
$\mathrm{predecessor}.\mathrm{Plus}$ with an incremented addend,
and $\mathrm{predecessor}$ also resolves through the trie union,
the recursion branches at every trie node.
The result $\mathrm{sum}$ therefore collects all paths reachable
through any combination of the two operands, which is the
Cartesian product.

The same $\mathrm{Plus}$ and $\mathrm{Equal}$ files that operate on
single values thus produce, without modification, the relational
semantics of logic programming:
$\overrides$ and $\supers$ collect all inheritance paths, and
every operation distributes over them.
\fi

\section{\bilingual{Discussion}{讨论}}

\subsection{\bilingual{Emergent Phenomena}{涌现现象}}
\label{sec:emergent-phenomena}

\ifInheritanceChinese
在 MIXINv2 中开发真实程序的过程中，我们发现了三种现象，每一种都是对某种已有设计模式的非预期用法，超出了该模式的原始能力。

\begin{description}
  \item[关系语义]
    第~\ref{sec:case-study}~节案例研究中的笛卡尔积之所以出现，是因为字典树是一种集合结构，因此对字典树的深度合并即为集合并，而对合并后字典树的操作则遍历其操作数的笛卡尔积。
    这并不局限于 $\mathrm{Nat}$：任何 Datalog 元组都可编码为链表，而深度合并在这些元组上计算集合操作，我们认为这正是 Datalog 语义。
    附录~\ref{app:datalog-encoding}~将 Datalog 编码进继承演算，仅用第~\ref{sec:syntax}~节的三个基本构造即可表达变量连接、常量匹配、多体规则和递归规则。
    这种逻辑编程语义是对编码案例研究算术的工厂模式的非预期用法中涌现出来的。

  \item[对不可扩展性的免疫]\label{item:immunity}
    在第~\ref{sec:case-study}~节的案例研究中，$\mathrm{Nat}$ 编码被证明是无类型表达问题~\cite{wadler1998-expression-problem}的一个实例：%
    \footnote{
      表达问题有两个要求：沿操作轴和数据构造子轴均可扩展，以及静态类型安全。无类型演算无法满足类型安全要求，因此我们只考虑双轴可扩展性要求。
    }
    一个新操作和一个新数据类型各自作为独立的 mixin 添加，沿操作轴和情形轴均进行了扩展，且无需修改任何已有 mixin。
    在其他场景中，访问者模式允许添加新操作但不允许添加新情形；而这里两者都允许，因为第~\ref{sec:syntax}~节的每个构造都可通过深度合并（公式~\ref{eq:overrides}）进行扩展，所以不可扩展的那一半根本无法写出；第~\ref{sec:asymmetry}~节（定理~\ref{thm:cartesian-nonexpressibility}）则形式化了使这一结论无条件成立的不对称性。
    在其他地方恢复第二个扩展维度需要专门的机制~\cite{liang1995-monad-transformers, lammel2003-scrap-your-boilerplate, loh2006-open-data-types,
      swierstra2008-data-types-a-la-carte, carette2009-finally-tagless, oliveira2012-object-algebras,
      kiselyov2013-extensible-effects, wu2014-effect-handlers-in-scope, kiselyov2015-freer-monads,
    wang2016-expression-problem-trivially, leijen2017-row-typed-algebraic-effects, poulsen2023-hefty-algebras}。
    第~\ref{sec:semantic-variants}~节分析了为何深度合并是唯一能实现这一点的组合机制。
    这第二个可扩展性维度是对访问者模式的非预期用法中涌现出来的。

\end{description}

\label{sec:practical-function-color}%
更多现象依赖于 FFI，而本文中的示例排除了 FFI。其中一种是函数颜色盲~\cite{nystrom2015-function-color}：同一个 MIXINv2 文件可以在同步和异步运行时下运行，这是通过继承进行依赖注入的非预期用法。补充材料\anon[]{~\cite{mixin2025}}给出了相关示例。
\else
Developing real programs in MIXINv2 surfaced three phenomena, each an
off-label use of an adopted design pattern beyond its original abilities.

\begin{description}
  \item[Relational semantics]
    The Cartesian product in the case study of
    Section~\ref{sec:case-study} arose because a trie is a set
    structure, so deep merge over tries is set union and an operation
    over the merged trie ranges over the product of its operands.
    This is not specific to $\mathrm{Nat}$: any Datalog tuple encodes
    as a linked list, and deep merge computes set operations over such
    tuples, which we believe is Datalog semantics.
    Appendix~\ref{app:datalog-encoding} encodes Datalog into
    inheritance-calculus, expressing variable joins,
    constant matches, multi-body rules, and recursive rules with
    the three primitives of Section~\ref{sec:syntax} alone.
    This logic-programming semantics emerged from an off-label use of
    the factory pattern that encodes the case study's arithmetic.

  \item[Immunity to nonextensibility]\label{item:immunity}
    In the case study of Section~\ref{sec:case-study}, the
    $\mathrm{Nat}$ encoding turned out to be an instance of the
    untyped Expression
    Problem~\cite{wadler1998-expression-problem}:%
    \footnote{
      The Expression Problem has two requirements: extensibility along
      both the operation and the data-constructor axis, and static
      type safety. An untyped calculus cannot meet the type-safety
      requirement, so we consider only the two-axis extensibility
      requirement.
    }
    a new operation and a new data type were each added as a separate
    mixin, extending along both the operation axis and the case axis,
    with no existing mixin modified.
    Elsewhere the visitor pattern admits new operations but not new
    cases; here it admits both, because every construct of
    Section~\ref{sec:syntax} is extensible via deep merge
    (equation~\ref{eq:overrides}), so the nonextensible half cannot be
    written, and Section~\ref{sec:asymmetry}
    (Theorem~\ref{thm:cartesian-nonexpressibility}) formalizes the
    asymmetry that makes this unconditional.
    Recovering that second dimension elsewhere takes dedicated
    mechanisms~\cite{liang1995-monad-transformers, lammel2003-scrap-your-boilerplate, loh2006-open-data-types,
      swierstra2008-data-types-a-la-carte, carette2009-finally-tagless, oliveira2012-object-algebras,
      kiselyov2013-extensible-effects, wu2014-effect-handlers-in-scope, kiselyov2015-freer-monads,
    wang2016-expression-problem-trivially, leijen2017-row-typed-algebraic-effects, poulsen2023-hefty-algebras}.
    Section~\ref{sec:semantic-variants} analyzes why deep merge
    is the only composition mechanism that achieves this.
    This second extensibility dimension emerged from an off-label use
    of the visitor pattern.

\end{description}

\label{sec:practical-function-color}%
Further phenomena rely on the FFI, which the examples in this paper
exclude. One is function color
blindness~\cite{nystrom2015-function-color}: the same MIXINv2 file runs
under both synchronous and asynchronous runtimes, an off-label use of
dependency injection via inheritance. The
supplementary material\anon[]{~\cite{mixin2025}} gives the example.
\fi

\subsection{\bilingual{Semantic Alternatives}{语义替代方案}}
\label{sec:semantic-variants}

\ifInheritanceChinese
上述\hyperref[item:immunity]{对不可扩展性的免疫}条目确立了继承演算对不可扩展性的免疫性。我们将此作为设计目标，并通过调查我们已知的所有候选方案，询问是否有替代设计能共享这一性质。

\paragraph{\bilingual{Degrees of freedom}{自由度}}
以下分析假设 AST 是一棵\emph{命名树}：从父节点到子节点的每条边都携带一个标签，因此每个节点都由从根节点出发的路径唯一标识。对于配置语言，这自然成立：JSON、YAML 或 Nix 属性集中的每个键都是一个标签，而嵌套产生路径。%
\footnote{
  Nix 是一门具有一等函数的函数式语言；仅凭其属性集并不构成命名树 DSL。然而，NixOS 模块系统~\cite{nixosmodules}将用户界面限制为通过递归合并组合的嵌套属性集，而正是这种 DSL 层面的结构自然地构成了命名树。
}
对于主流编程语言，AST 并不直接是命名树，因为匿名表达式位置和隐式作用域缺乏标签。然而，经过 ANF 变换~\cite{flanagan1993-essence-compiling-continuations}后，每个中间结果都绑定到一个名字，而剩余的匿名位置%
\footnote{
  例如，$\mathbf{let}$ 的主体。
}
可以被分配合成的新鲜标签，例如第~\ref{sec:forward-translation}~节中的 $\mathrm{result}$ 和 $\mathrm{tailCall}$，从而得到一棵命名树。

在此前提下，每种程序是带引用的命名树的语言，都可以由两个有限集来刻画：一个由\emph{节点类型}构成的集合 $\mathcal{N}$，每种节点类型决定子树如何组织其子节点；以及一个由\emph{引用类型}构成的集合 $\mathcal{R}$，每种引用类型决定树中一个位置如何引用另一个位置。语法由 $\mathcal{N}$ 和 $\mathcal{R}$ 的选择固定；语义则由每种节点类型如何构造组合以及每种引用类型如何解析其目标来决定。

{\small
  \begin{center}
    \begin{tabular}{lp{0.38\columnwidth}p{0.33\columnwidth}}
      \bilingual{\textbf{Language}}{\textbf{语言}}
      & \bilingual{\textbf{Node types $\mathcal{N}$}}{\textbf{节点类型 $\mathcal{N}$}}
      & \bilingual{\textbf{Reference types $\mathcal{R}$}}{\textbf{引用类型 $\mathcal{R}$}} \\
      \hline
      Java (OOP)
      & \bilingual{class, interface, method,
      block, conditional, loop,
      \texttt{throw}, lambda, \ldots}{类、接口、方法、块、条件、循环、\texttt{throw}、lambda、\ldots}
      & \bilingual{variable, field access, method call,
      \texttt{extends}/\texttt{implements},
      \texttt{import}, \ldots}{变量、字段访问、方法调用、\texttt{extends}/\texttt{implements}、\texttt{import}、\ldots} \\[3pt]
      Haskell (FP)
      & \bilingual{$\lambda$, application, \texttt{let},
      \texttt{case}, \texttt{data},
      type class, instance, \ldots}{$\lambda$、应用、\texttt{let}、\texttt{case}、\texttt{data}、类型类、实例、\ldots}
      & \bilingual{variable, constructor,
      type class dispatch,
      \texttt{import}, \ldots}{变量、构造子、类型类分派、\texttt{import}、\ldots} \\[3pt]
      Scala\footnote{
        \bilingual{Scala is a multi-paradigm language.}{Scala 是一门多范式语言。}
      }
      & \bilingual{class, trait, object, method,
      block, conditional, \texttt{match},
      \texttt{while}, \texttt{throw}, \ldots}{类、trait、对象、方法、块、条件、\texttt{match}、\texttt{while}、\texttt{throw}、\ldots}
      & \bilingual{variable, member access,
      \texttt{extends}, \texttt{with},
      implicits, \texttt{import}, \ldots}{变量、成员访问、\texttt{extends}、\texttt{with}、隐式参数、\texttt{import}、\ldots} \\
      \hline
      $\lambda$-\bilingual{calculus}{演算}
      & \bilingual{abstraction, application}{抽象、应用}
      & \bilingual{variable}{变量} \\
      \bilingual{Inheritance-calculus}{继承演算}
      & \bilingual{record}{记录}
      & \bilingual{inheritance}{继承} \\
    \end{tabular}
  \end{center}
}

\paragraph{\bilingual{Immunity to nonextensibility}{对不可扩展性的免疫}}
程序 $P$ 中的路径 $p$ 是\emph{可扩展的}，若存在程序 $Q$，使得通过引用将 $P$ 与 $Q$ 组合后，在不修改 $P$ 的前提下，$p$ 处能够观察到新的属性。一种语言\emph{对不可扩展性免疫}，若其每个良构程序的每条路径都是可扩展的。

我们现在对设计空间进行分类，以解释为何我们选择第~\ref{sec:mixin-trees}~节的语义。两个维度保持自由：\emph{组合机制}，决定单一引用类型如何合并定义；以及\emph{引用解析}，决定单一引用类型如何找到其目标。

\subsubsection*{\bilingual{The composition mechanism}{组合机制}}

每种组合机制决定了一个抽象的定义如何被合并进另一个抽象。下表按各候选方案是否满足可交换性~(C)、幂等性~(I) 和可结合性~(A)，以及其可扩展性模型对已知候选方案进行分类。

\begin{center}
  \begin{tabular}{lcccll}
    \bilingual{\textbf{Candidate}}{\textbf{候选方案}} & \textbf{C} & \textbf{I} & \textbf{A}
    & \bilingual{\textbf{Extensibility}}{\textbf{可扩展性}} & \bilingual{\textbf{Representative}}{\textbf{代表}} \\
    \hline
    \bilingual{Deep merge}{深度合并}
    & $\checkmark$ & $\checkmark$ & $\checkmark$
    & \bilingual{universal}{通用} & \bilingual{this paper}{本文} \\
    \bilingual{Shallow merge}{浅层合并}
    & $\times$ & $\checkmark$ & $\checkmark$
    & \bilingual{top-level only}{仅顶层} & JS spread \\
    \bilingual{Conflict rejection}{冲突拒绝}
    & $\checkmark$ & $\checkmark$ & $\checkmark$
    & \bilingual{none on overlap}{重叠时无扩展}
    & Harper--Pierce~\cite{harper1991-symmetric-record-concatenation} \\
    \bilingual{Priority merge}{优先级合并}
    & $\times$ & $\times$ & $\checkmark$
    & \bilingual{universal}{通用} & Jsonnet~\texttt{+}, Nix~\texttt{//} \\
    $\beta$-\bilingual{reduction}{归约}
    & $\times$ & $\times$ & $\times$
    & \bilingual{opt-in}{按需选择} & Java, Haskell, Scala, \ldots \\
    \bilingual{Conditional}{条件并入}
    & \multicolumn{3}{c}{\bilingual{causes divergence}{导致发散}}
    & --- & --- \\
  \end{tabular}
\end{center}

\begin{description}
  \item[\bilingual{Deep merge}{深度合并}]
    对同名定义的合并递归进入其子树（$\overrides$，公式~\ref{eq:overrides}）。每个标签都是一个扩展点，无论原作者是否预期了扩展。

  \item[\bilingual{Shallow merge}{浅层合并}]
    嵌套标签无法被独立扩展。

  \item[\bilingual{Conflict rejection}{冲突拒绝}]
    同名定义被作为错误拒绝。三个等式性质均成立，但为空洞成立：这些等式得到满足，是因为用于验证它们的组合被拒绝了。冲突拒绝阻止了对已有标签的扩展~\cite{harper1991-symmetric-record-concatenation}，而这恰恰是表达问题~\cite{wadler1998-expression-problem}所要求的可组合性。

  \item[\bilingual{Priority merge}{优先级合并}]
    一个来源具有优先权~\cite{jsonnet,dolstra2010-nixos}，因此 $A + B \neq B + A$，引入了顺序依赖。

  \item[$\beta$-\bilingual{reduction}{归约}]
    只有作者放置的参数才是扩展点；函数体是封闭的。要使一个新的方面可扩展，原函数必须被重写以接受一个额外参数。任何基于函数应用的组合机制都继承了这一限制，因为 $\beta$-归约使函数体封闭。表达问题~\cite{wadler1998-expression-problem}在此类语言中之所以困难，原因相同：可扩展性是按需选择的。模拟对不可扩展性免疫的框架（访问者模式、对象代数~\cite{oliveira2012-object-algebras}、finally tagless 解释器~\cite{carette2009-finally-tagless}）正是为绕过这一限制而生的机制。

  \item[\bilingual{Conditional incorporation}{条件并入}]
    根据其他定义的存在与否来移除或有条件地包含定义，会产生依赖循环：某定义是否存在，取决于另一定义是否存在，而后者又取决于前者。当添加和移除交替出现时，递归求值（附录~\ref{app:well-definedness}）无法到达不动点，从而导致发散。这与图灵完备性固有的发散不同：图灵完备程序发散是因为计算无界，而条件并入发散是因为直接后果算子 $T_P$ 振荡而非单调递增。
\end{description}

\noindent
在上述候选方案中，深度合并是我们找到的唯一一个可交换、幂等、可结合且对不可扩展性免疫的方案。其他候选方案各自引入了对不可扩展性的脆弱性，或导致发散。

\subsubsection*{\bilingual{The reference resolution}{引用解析}}

一旦将记录上的深度合并确定为组合机制，六个语义方程中的四个便由语法决定：$\properties$~(\ref{eq:properties}) 读取记录的标签，$\bases$~(\ref{eq:bases}) 读取其继承来源，$\supers$~(\ref{eq:supers}) 取传递闭包，$\overrides$~(\ref{eq:overrides}) 实现合并。唯一剩余的自由度是引用的解析方式：$\resolve$~(\ref{eq:resolve}) 和 $\this$~(\ref{eq:this}) 方程。我们考察三个已知候选方案。

\begin{description}
  \item[\bilingual{Early binding}{早期绑定}]
    使用早期绑定时，引用在定义时即解析到一个固定路径，与记录后来如何被继承无关。CUE~\cite{cue2019}是这方面的典型：字段引用被静态解析，嵌套字段无法引用其封闭上下文中由后续组合贡献的属性。这使得 mixin 无法充当 $\lambda$-演算的调用栈：$\beta$-归约要求被调用方能观察到调用点提供的参数，而那是一次后续组合。使用早期绑定时，引用在此组合发生前就已被冻结。在方程中，$\resolve$~(\ref{eq:resolve}) 调用 $\this$~(\ref{eq:this})，后者遍历继承的树结构；迟绑定正是使 $\this$ 能够观察到后续组合的 mixin 所贡献属性的原因。若没有图灵完备的 mixin，函数必须被重新引入为组合机制，从而再次暴露上述按需选择的可扩展性问题。

  \item[\bilingual{Dynamic scope}{动态作用域}]
    使用动态作用域时，引用在使用点而非定义点处解析。Jsonnet~\cite{jsonnet}采用此方式：\texttt{self} 始终指向最终合并后的对象，而非定义点处的对象。联合文件系统在文件系统层面表现出相同的性质：符号链接的相对路径在解引用时相对于挂载上下文解析，而非相对于定义该链接的层。动态作用域破坏了第~\ref{sec:forward-translation}~节的 $\lambda$-演算嵌入。替换引理（引理~\ref{lem:substitution}）依赖于基本情况 $M = y$（$y \neq x$）：$y$ 的 de~Bruijn 指标 $m$ 从定义点路径 $\init(p_{\mathrm{def}})$ 向上走到一个与 $x$ 的绑定器不同的作用域层级，因此 $x$ 作用域处的替换 $\mathrm{argument} \mapsto \mathcal{T}(V)$ 不产生任何效果。动态作用域下，这种分离不再成立。%
    \footnote{
      反例：$(\lambda x.\, \lambda y.\, x)\; V_1\; V_2$。$\lambda y.\, x$ 内部对 $x$ 的引用翻译为 $\dbi{1}.\mathrm{argument}$（de~Bruijn 指标 $1$）。使用迟绑定时，$\this$ 从定义点作用域（内层 $\lambda$）出发，向上走一层到外层 $\lambda$，到达槽 $\mathrm{argument} \mapsto \mathcal{T}(V_1)$。使用动态作用域时，引用从使用点上下文解析，其中最内层封闭作用域是应用 $\{C_1.\mathrm{result},\; \mathrm{argument} \mapsto \mathcal{T}(V_2)\}$；向上走一层到达该作用域的 $\mathrm{argument}$ 槽，返回 $\mathcal{T}(V_2)$ 而非 $\mathcal{T}(V_1)$。这是动态作用域的经典变量捕获失败。
    }
    与早期绑定一样，mixin 无法完全替代函数。

  \item[\bilingual{Single-target $\this$}{单目标 $\this$}]
    单目标 $\this$ 是一个有前途的候选方案：单路径引理~(\ref{lem:single-path})表明，对于翻译后的 $\lambda$-项，$\this$~(\ref{eq:this}) 中的前沿集 $S$ 始终恰好包含一条路径，因此 $\lambda$-演算嵌入不需要多目标解析。然而，单目标 $\this$ 因另一原因而对不可扩展性脆弱。当多条继承路径到达同一作用域时，单目标解析必须拒绝该情况或选择其中一条路径。NixOS 模块系统~\cite{nixosmodules}会引发重复声明错误（附录~\ref{app:scala-multi-target}）。这拒绝了两个独立编写的模块定义同一选项的嵌套扩展，而这种组合在求解表达问题~\cite{wadler1998-expression-problem}时会出现，如第~\ref{sec:expression-problem}~节所示。

    我们对 MIXINv2 的前三个实现使用了 $\this$ 解析的标准技术：遵循 Cook~\cite{cook1989-denotational-inheritance}的基于闭包的不动点，以及使用 de~Bruijn 指标~\cite{debruijn1972-nameless-dummies}的基于栈的环境查找。当继承引入到同一封闭作用域的多条路径时，三种实现都产生了隐晦的错误。困难在于结构上。闭包捕获单一环境，因此不动点组合子只能为单一 $\this$ 目标求解。基于栈的环境是一条线性链，每个作用域层级只有一个父节点。这两种数据结构都内嵌了一个与多重继承自然产生的多目标情况不相容的单值假设（附录~\ref{app:scala-multi-target}）。公式~(\ref{eq:this})是在放弃这一假设后涌现出来的：它跟踪一个\emph{前沿集} $S$，而非单一当前作用域。
\end{description}

\noindent
下表总结了两个层面上所有替代方案的脆弱性。

\begin{center}
  \begin{tabular}{llll}
    & \bilingual{\textbf{Candidate}}{\textbf{候选方案}} & \bilingual{\textbf{Vulnerability}}{\textbf{脆弱性}}
    & \bilingual{\textbf{Representative}}{\textbf{代表}} \\
    \hline
    \multirow{5}{*}{\rotatebox[origin=c]{90}{\scriptsize \bilingual{composition}{组合}}}
    & \bilingual{Shallow merge}{浅层合并}
    & \bilingual{nested extensions not composable}{嵌套扩展不可组合}
    & JS spread \\
    & \bilingual{Conflict rejection}{冲突拒绝}
    & \bilingual{same-label extensions rejected}{同名扩展被拒绝}
    & Harper--Pierce~\cite{harper1991-symmetric-record-concatenation} \\
    & \bilingual{Priority merge}{优先级合并}
    & \bilingual{ordering dependency}{顺序依赖}
    & Jsonnet~\texttt{+}~\cite{jsonnet},
    Nix~\texttt{//}~\cite{dolstra2010-nixos} \\
    & $\beta$-\bilingual{reduction}{归约}
    & \bilingual{opt-in extensibility}{按需选择的可扩展性}
    & Java, Haskell, Scala, \ldots \\
    & \bilingual{Conditional}{条件并入}
    & \bilingual{causes divergence}{导致发散}
    & --- \\
    \hline
    \multirow{3}{*}{\rotatebox[origin=c]{90}{\scriptsize \bilingual{resolution}{解析}}}
    & \bilingual{Early binding}{早期绑定}
    & \bilingual{mixin cannot replace function}{mixin 无法替代函数}
    & CUE~\cite{cue2019} \\
    & \bilingual{Dynamic scope}{动态作用域}
    & \bilingual{mixin cannot replace function}{mixin 无法替代函数}
    & Jsonnet~\cite{jsonnet}, union FS \\
    & \bilingual{Single-target $\this$}{单目标 $\this$}
    & \bilingual{nested extensions rejected}{嵌套扩展被拒绝}
    & NixOS modules~\cite{nixosmodules} \\
  \end{tabular}
\end{center}

\noindent
我们在两个层面考察的所有替代方案，均对不可扩展性不免疫。

我们通过尝试本节中的每个替代方案并发现其失败，从而得到了这一语义：浅层合并无法组合嵌套扩展；优先级合并引入了顺序依赖；早期绑定在组合能够提供引用目标之前就将引用冻结；动态作用域破坏了 $\lambda$-演算嵌入所依赖的替换；单目标 $\this$ 拒绝了表达问题所要求的多目标组合。深度合并加上多目标迟绑定 $\this$，正是最终剩下的选择。
\else
The \hyperref[item:immunity]{immunity-to-nonextensibility item} above established that inheritance-calculus is immune
to nonextensibility.
We adopt this as a design goal and ask whether an alternative
design could share this property
by surveying every candidate known to us.

\paragraph{\bilingual{Degrees of freedom}{自由度}}
The analysis below assumes that the AST is a \emph{named tree}:%
every edge from parent to child carries a label, so that every
node is uniquely identified by the path from the root.
For configuration languages this is naturally the case: every
key in a JSON, YAML, or Nix attribute set is a label, and nesting
produces paths.%
\footnote{
  Nix is a functional language with first-class functions;
  its attribute sets alone do not form a named-tree DSL.
  The NixOS module system~\cite{nixosmodules}, however,
  restricts the user-facing interface to nested attribute sets
  composed by recursive merging, and it is this DSL-level
  structure that naturally forms a named tree.
}
For mainstream programming languages the AST is not directly a
named tree, since anonymous expression positions and implicit
scopes lack labels.
After ANF
transformation~\cite{flanagan1993-essence-compiling-continuations}, however, every intermediate
result is bound to a name, and the remaining anonymous positions%
\footnote{
  For example, the body of a $\mathbf{let}$.
}
can be assigned synthetic
fresh labels such as $\mathrm{result}$ and $\mathrm{tailCall}$
in Section~\ref{sec:forward-translation}, yielding a named
tree.

Under this premise, every language whose programs are such
named trees with references can be characterized by two finite sets:
a set~$\mathcal{N}$ of \emph{node types}, each determining how a
subtree organizes its children, and a set~$\mathcal{R}$ of
\emph{reference types}, each determining how one position in the
tree refers to another.
The syntax is fixed by the choice of $\mathcal{N}$ and
$\mathcal{R}$; the semantics is determined by how each node type
structures composition and how each reference type resolves its
target.

{\small
  \begin{center}
    \begin{tabular}{lp{0.38\columnwidth}p{0.33\columnwidth}}
      \textbf{Language}
      & \textbf{Node types $\mathcal{N}$}
      & \textbf{Reference types $\mathcal{R}$} \\
      \hline
      Java (OOP)
      & class, interface, method,
      block, conditional, loop,
      \texttt{throw}, lambda, \ldots
      & variable, field access, method call,
      \texttt{extends}/\texttt{implements},
      \texttt{import}, \ldots \\[3pt]
      Haskell (FP)
      & $\lambda$, application, \texttt{let},
      \texttt{case}, \texttt{data},
      type class, instance, \ldots
      & variable, constructor,
      type class dispatch,
      \texttt{import}, \ldots \\[3pt]
      Scala\footnote{
        Scala is a multi-paradigm language.
      }
      & class, trait, object, method,
      block, conditional, \texttt{match},
      \texttt{while}, \texttt{throw}, \ldots
      & variable, member access,
      \texttt{extends}, \texttt{with},
      implicits, \texttt{import}, \ldots \\
      \hline
      $\lambda$-calculus
      & abstraction, application
      & variable \\
      Inheritance-calculus
      & record
      & inheritance \\
    \end{tabular}
  \end{center}
}

\paragraph{\bilingual{Immunity to nonextensibility}{对不可扩展性的免疫}}
A path~$p$ in a program~$P$ is \emph{extensible} if there exists
a program~$Q$ such that composing $P$ with~$Q$ through a reference
makes new properties observable at~$p$ without modifying~$P$.
A language is \emph{immune to nonextensibility} if every path in
every well-formed program is extensible.

We now classify the design space to explain why we chose the
semantics of Section~\ref{sec:mixin-trees}.
Two axes remain free:
the \emph{composition mechanism}, which determines how the single reference
type incorporates definitions;
and the \emph{reference resolution}, which determines how the single reference
type finds its target.

\subsubsection*{\bilingual{The composition mechanism}{组合机制}}

Every composition mechanism determines how the definitions of one
abstraction are incorporated into another.
The following table classifies the known candidates by whether
they satisfy commutativity~(C), idempotence~(I), and
associativity~(A), and by their extensibility model.

\begin{center}
  \begin{tabular}{lcccll}
    \bilingual{\textbf{Candidate}}{\textbf{候选方案}} & \textbf{C} & \textbf{I} & \textbf{A}
    & \bilingual{\textbf{Extensibility}}{\textbf{可扩展性}} & \bilingual{\textbf{Representative}}{\textbf{代表}} \\
    \hline
    \bilingual{Deep merge}{深度合并}
    & $\checkmark$ & $\checkmark$ & $\checkmark$
    & \bilingual{universal}{通用} & \bilingual{this paper}{本文} \\
    \bilingual{Shallow merge}{浅层合并}
    & $\times$ & $\checkmark$ & $\checkmark$
    & \bilingual{top-level only}{仅顶层} & JS spread \\
    \bilingual{Conflict rejection}{冲突拒绝}
    & $\checkmark$ & $\checkmark$ & $\checkmark$
    & \bilingual{none on overlap}{重叠时无扩展}
    & Harper--Pierce~\cite{harper1991-symmetric-record-concatenation} \\
    \bilingual{Priority merge}{优先级合并}
    & $\times$ & $\times$ & $\checkmark$
    & \bilingual{universal}{通用} & Jsonnet~\texttt{+}, Nix~\texttt{//} \\
    $\beta$-\bilingual{reduction}{归约}
    & $\times$ & $\times$ & $\times$
    & \bilingual{opt-in}{按需选择} & Java, Haskell, Scala, \ldots \\
    \bilingual{Conditional}{条件并入}
    & \multicolumn{3}{c}{\bilingual{causes divergence}{导致发散}}
    & --- & --- \\
  \end{tabular}
\end{center}

\begin{description}
  \item[Deep merge]
    Merging same-label definitions recurses into their subtrees
    ($\overrides$, equation~\ref{eq:overrides}).
    Every label is an extension point, whether or not the
    original author anticipated extension.

  \item[Shallow merge]
    Nested labels cannot be independently extended.

  \item[Conflict rejection]
    Same-label definitions are rejected as errors.
    All three identities hold, but vacuously:
    the identities are satisfied because the compositions that
    would test them are refused.
    Conflict rejection prevents extending existing
    labels~\cite{harper1991-symmetric-record-concatenation}, precisely the composability that the
    Expression Problem~\cite{wadler1998-expression-problem} requires.

  \item[Priority merge]
    One source takes
    precedence~\cite{jsonnet,dolstra2010-nixos},
    so $A + B \neq B + A$, introducing
    an ordering dependency.

  \item[$\beta$-reduction]
    Only the parameters placed by the author are extension
    points; the function body is closed.
    To make a new aspect extensible, the original function must
    be rewritten to accept an additional parameter.
    Any composition mechanism based on function application
    inherits this limitation, because $\beta$-reduction
    closes the function body.
    The Expression Problem~\cite{wadler1998-expression-problem} is hard
    in such languages for the same reason:
    extensibility is opt-in.
    The frameworks that simulate immunity to nonextensibility
    (visitor pattern,
      object algebras~\cite{oliveira2012-object-algebras},
    finally tagless interpreters~\cite{carette2009-finally-tagless})
    are the machinery needed to work around this limitation.

  \item[Conditional incorporation]
    Removing or conditionally including definitions
    based on the presence of other definitions
    creates a dependency cycle:
    whether a definition exists depends on
    whether another definition exists, which in turn
    depends on the first.
    The recursive evaluation
    (Appendix~\ref{app:well-definedness}) cannot
    reach a fixed point when addition and removal
    alternate, causing divergence.
    This differs from the divergence inherent to
    Turing completeness: Turing-complete programs
    diverge because computation is unbounded, whereas
    conditional incorporation diverges because the
    immediate consequence operator $T_P$ oscillates
    rather than ascending monotonically.
\end{description}

\noindent
Among the candidates above, deep merge is the only one we found
that is commutative, idempotent, associative, and
immune to nonextensibility.
The other candidates each introduce a vulnerability to
nonextensibility or cause divergence.

\subsubsection*{\bilingual{The reference resolution}{引用解析}}

Once deep merge over records is fixed as the composition
mechanism, four of the six semantic equations are determined
by the syntax:
$\properties$~(\ref{eq:properties}) reads the labels of a
record, $\bases$~(\ref{eq:bases}) reads its inheritance
sources, $\supers$~(\ref{eq:supers}) takes the transitive
closure, and $\overrides$~(\ref{eq:overrides}) implements
the merge.
The only remaining degree of freedom is how references
resolve: the $\resolve$~(\ref{eq:resolve}) and
$\this$~(\ref{eq:this}) equations.
We examine the three known candidates.

\begin{description}
  \item[Early binding]
    With early binding, references resolve at definition time to a
    fixed path, independently of how the record is later inherited.
    CUE~\cite{cue2019} exemplifies this: field references are
    statically resolved, and a nested field cannot refer to
    properties of its enclosing context contributed by later
    composition.
    This prevents the mixin from serving as a call stack for the
    $\lambda$-calculus: $\beta$-reduction requires the callee to
    observe the argument supplied at the call site, which is a
    later composition.
    With early binding, references are frozen before this
    composition occurs.
    In the equations,
    $\resolve$~(\ref{eq:resolve}) calls
    $\this$~(\ref{eq:this}), which walks the inherited
    tree structure; late binding is what allows
    $\this$ to see properties contributed by later-composed
    mixins.
    Without a Turing-complete mixin, functions must be
    reintroduced as the composition mechanism, re-exposing
    the opt-in extensibility problem above.

  \item[Dynamic scope]
    With dynamic scope, references resolve at the point of use
    rather than the definition site.
    Jsonnet~\cite{jsonnet} follows this approach: \texttt{self}
    always refers to the final merged object, not the object at
    the definition site.
    Union file systems exhibit the same property at the
    file-system level: a symbolic link's relative path resolves
    against the mount context at dereference time, not against the
    layer that defined the link.
    Dynamic scope breaks the $\lambda$-calculus embedding
    of Section~\ref{sec:forward-translation}.
    The Substitution Lemma (Lemma~\ref{lem:substitution})
    relies on the base case $M = y$ ($y \neq x$):
    the de~Bruijn index $m$ of~$y$ walks up from the
    definition-site path $\init(p_{\mathrm{def}})$ to a scope
    level different from~$x$'s binder, so the substitution
    $\mathrm{argument} \mapsto \mathcal{T}(V)$ at~$x$'s scope
    has no effect.
    With dynamic scope this separation fails.%
    \footnote{
      Counterexample:
      $(\lambda x.\, \lambda y.\, x)\; V_1\; V_2$.
      The reference to~$x$ inside $\lambda y.\, x$ translates to
      $\dbi{1}.\mathrm{argument}$ (de~Bruijn index~$1$).
      With late binding, $\this$ starts from the definition-site
      scope (the inner $\lambda$) and walks up one level to the
      outer $\lambda$, reaching the slot
      $\mathrm{argument} \mapsto \mathcal{T}(V_1)$.
      With dynamic scope, the reference resolves from the
      use-site context, where the innermost enclosing scope is
      the application $\{C_1.\mathrm{result},\;
      \mathrm{argument} \mapsto \mathcal{T}(V_2)\}$;
      walking up one level reaches this scope's
      $\mathrm{argument}$ slot, returning $\mathcal{T}(V_2)$
      instead of $\mathcal{T}(V_1)$.
      This is the classical variable-capture failure of dynamic
      scoping.
    }
    As with early binding, the mixin cannot serve as a complete
    replacement for functions.

  \item[Single-target $\this$]
    Single-target $\this$ is a promising candidate:
    the Single-Path Lemma~(\ref{lem:single-path}) shows that
    for translated $\lambda$-terms, the frontier set~$S$ in
    $\this$~(\ref{eq:this}) always contains exactly one path,
    so the $\lambda$-calculus embedding does not require
    multi-target resolution.
    However, single-target $\this$ is vulnerable to
    nonextensibility for a different reason.
    When multiple inheritance routes reach the same scope,
    single-target resolution must reject the situation or select
    one route.
    The NixOS module system~\cite{nixosmodules} raises a
    duplicate-declaration error
    (Appendix~\ref{app:scala-multi-target}).
    This rejects nested extensions where two independently authored
    modules define the same option, a composition that arises
    when solving the Expression
    Problem~\cite{wadler1998-expression-problem},
    as Section~\ref{sec:expression-problem} demonstrates.

    Our first three implementations of MIXINv2 used
    standard techniques for $\this$ resolution: closure-based
    fixed points following Cook~\cite{cook1989-denotational-inheritance}, and
    stack-based environment lookup using de~Bruijn
    indices~\cite{debruijn1972-nameless-dummies}.
    All three produced subtle bugs when inheritance introduced
    multiple routes to the same enclosing scope.
    The difficulty was structural.
    A closure captures a single environment, so a fixed-point
    combinator solves for a single $\this$ target.
    A stack-based environment is a linear chain where each scope
    level has one parent.
    Both data structures embed a single-valued assumption
    incompatible with the multi-target situation that multiple
    inheritance naturally produces
    (Appendix~\ref{app:scala-multi-target}).
    Equation~(\ref{eq:this}) emerged from abandoning this
    assumption: it tracks a \emph{frontier set}~$S$ rather than a
    single current scope.
\end{description}

\noindent
The following table summarizes the vulnerabilities of all
alternatives at both levels.

\begin{center}
  \begin{tabular}{llll}
    & \bilingual{\textbf{Candidate}}{\textbf{候选方案}} & \bilingual{\textbf{Vulnerability}}{\textbf{脆弱性}}
    & \bilingual{\textbf{Representative}}{\textbf{代表}} \\
    \hline
    \multirow{5}{*}{\rotatebox[origin=c]{90}{\scriptsize \bilingual{composition}{组合}}}
    & \bilingual{Shallow merge}{浅层合并}
    & \bilingual{nested extensions not composable}{嵌套扩展不可组合}
    & JS spread \\
    & \bilingual{Conflict rejection}{冲突拒绝}
    & \bilingual{same-label extensions rejected}{同名扩展被拒绝}
    & Harper--Pierce~\cite{harper1991-symmetric-record-concatenation} \\
    & \bilingual{Priority merge}{优先级合并}
    & \bilingual{ordering dependency}{顺序依赖}
    & Jsonnet~\texttt{+}~\cite{jsonnet},
    Nix~\texttt{//}~\cite{dolstra2010-nixos} \\
    & $\beta$-\bilingual{reduction}{归约}
    & \bilingual{opt-in extensibility}{按需选择的可扩展性}
    & Java, Haskell, Scala, \ldots \\
    & \bilingual{Conditional}{条件并入}
    & \bilingual{causes divergence}{导致发散}
    & --- \\
    \hline
    \multirow{3}{*}{\rotatebox[origin=c]{90}{\scriptsize \bilingual{resolution}{解析}}}
    & \bilingual{Early binding}{早期绑定}
    & \bilingual{mixin cannot replace function}{mixin 无法替代函数}
    & CUE~\cite{cue2019} \\
    & \bilingual{Dynamic scope}{动态作用域}
    & \bilingual{mixin cannot replace function}{mixin 无法替代函数}
    & Jsonnet~\cite{jsonnet}, union FS \\
    & \bilingual{Single-target $\this$}{单目标 $\this$}
    & \bilingual{nested extensions rejected}{嵌套扩展被拒绝}
    & NixOS modules~\cite{nixosmodules} \\
  \end{tabular}
\end{center}

\noindent
No alternative we examined at either level is immune to
nonextensibility.

We arrived at this semantics by trying each alternative in this section
and finding that it failed:
shallow merge could not compose nested extensions;
priority merge introduced ordering dependencies;
early binding froze references before composition could
supply them;
dynamic scope broke the substitution that the
$\lambda$-calculus embedding relies on;
single-target $\this$ rejected the multi-target compositions
that the Expression Problem demands.
Deep merge with multi-target late-binding $\this$ is what
remained.
\fi

\section{\bilingual{Related Work}{相关工作}}
\label{sec:related-work}

\subsection{\bilingual{Computational Models and \texorpdfstring{$\lambda$}{λ}-Calculus Semantics}{计算模型与 \texorpdfstring{$\lambda$}{λ}-演算语义}}

\paragraph{\bilingual{Minimal Turing-complete models}{最小图灵完备模型}}
\ifInheritanceChinese
Sch\"onfinkel~\cite{schonfinkel1924-combinatory-logic} 与 Curry~\cite{curry1958-combinatory-logic} 证明了 $\lambda$-演算可以归约为三个组合子 $S$、$K$、$I$%
\footnote{
  组合子 $I$ 是冗余的,因为 $I = SKK$。
}。将 $\lambda$-项编译为 SKI 的括号抽象是机械的,但会导致项大小的指数级膨胀;在 $\lambda$-演算中已有的计算模式之外,不会涌现出任何新的计算模式。在语义上,两者共享相同的模型:Meyer~\cite{meyer1982-model-of-lambda-calculus} 证明了满足五个附加条件的组合代数恰好就是 $\lambda$-代数。从未有人展示过一个独立的 SKI 语义域,并从中推导出 $\lambda$-演算的语义。方向始终是反过来的:Scott 的 $D_\infty$ 是为 $\lambda$-演算而构建的,而 SKI 随后被解释于其中。Dolan~\cite{dolan2013-mov-turing-complete} 证明了 x86 的 \texttt{mov} 指令单独就已图灵完备,其方法是利用寻址模式作为可变内存上的隐式算术;由此得到的程序是正确的,但不提供任何超越 RAM 机的抽象机制。图灵机本身虽然是通用的,却缺乏随机访问:读取距离为 $d$ 处的格需要 $d$ 次顺序的磁带移动。这些模型中的每一个,都在某个已有模型的语义框架内实现了图灵完备性,且对其中任何一个都从未展示过独立的语义域。
\else
Sch\"onfinkel~\cite{schonfinkel1924-combinatory-logic} and
Curry~\cite{curry1958-combinatory-logic} showed that the
$\lambda$-calculus reduces to three combinators $S$, $K$, $I$
\footnote{
  The $I$ combinator is redundant since $I = SKK$.
}.
The bracket abstraction that compiles $\lambda$-terms to SKI
is mechanical but incurs exponential blow-up in term size;
no new computational patterns emerge beyond those already
present in the $\lambda$-calculus.
Semantically, the two share the same models:
Meyer~\cite{meyer1982-model-of-lambda-calculus} showed that combinatory algebras
satisfying five additional conditions are exactly lambda
algebras.
No independent semantic domain for SKI has ever been
exhibited from which $\lambda$-calculus semantics could be
derived. The direction has always been the reverse:
Scott's $D_\infty$ was built for the $\lambda$-calculus,
and SKI was interpreted inside it afterwards.
Dolan~\cite{dolan2013-mov-turing-complete} showed that the x86 \texttt{mov}
instruction alone is Turing complete by exploiting
addressing modes as implicit arithmetic on mutable memory;
the resulting programs are correct but provide no
abstraction mechanisms beyond those of the RAM Machine.
The Turing Machine itself, while universal, lacks random
access: reading a cell at distance~$d$ requires $d$
sequential tape movements.
Each of these models achieves Turing completeness
within the semantic framework of an existing model,
and no independent semantic domain has been exhibited
for any of them.
\fi

%% === 审稿人备忘：SF-calculus 与 inheritance-calculus 的关系 ===
%%
%% Barry Jay 的 SF-calculus（Intensional Computation with Higher-Order
%% Functions, TCS 2019, doi:10.1016/j.tcs.2019.02.016；前身为
%% Confusion in the Church-Turing Thesis, arXiv:1410.7103, 2014）
%% 和 inheritance-calculus 都试图解决 λ-calculus 的 extensionality 限制，
%% 但方向完全不同。
%%
%% 共同出发点：λ-calculus 是 extensional 的——只能观察函数的输入输出行为，
%% 无法检查程序的内部结构。两个内部结构不同但行为相同的 λ 项不可区分。
%%   - Jay 的诊断：λ-calculus 缺乏 intensional 操作（检查项的结构），
%%     无法定义结构等式判定。
%%   - 本文的诊断：λ-calculus 的 function body 是 closed 的，无法从外部
%%     添加新的可观察投影（open extension）。
%%
%% 核心机制对比：
%%   SF-calculus：保留 S、K（或 λ），添加 F（factorisation）算子，
%%     可以将闭合范式拆解为组成部分。程序仍然是函数，只是多了检查函数
%%     结构的能力。属于 λ-calculus 的 *扩展*。
%%   Inheritance-calculus：用三个不同的原语（record、definition、
%%     inheritance）替代 λ-calculus 的三个原语。β-reduction 是 deep merge
%%     的特例（Section 3 的翻译 T 将函数应用映射为继承）。
%%     属于 λ-calculus 的 *包含*（subsumption）。
%%
%% 关于 F 算子的编码：
%%   F 做两件事：(1) 判断结构形态（atom vs. compound），
%%   (2) 提取子部分。
%%   在 inheritance-calculus 中：
%%     - 提取子部分是平凡的：所有值都是透明 record，通过 qualified this
%%       投影任意 slot 即可。Theorem 6 的 separating context 就在读取
%%       argument slot（↑².argument）。
%%     - 判断结构形态并分支：需要 Acceptance pattern（Section 4）。
%%       NatVisitor 对 Nat 做的事情与 F 对 SF-calculus 项做的事情
%%       完全对应：VisitZero/SuccessorVisitor 对应 F 的 atom/compound
%%       分支。Visitor 可以通过 open extension 添加到已有 term 上，
%%       不修改原始定义。
%%   Acceptance 在精神上就是 F：两者都是对 closed normal form 做
%%   constructor dispatch。区别在于 F 是不可扩展的 built-in primitive
%%   （硬编码了 S 和 F 两个 atom），而 Acceptance 是可扩展的 encoding
%%   （新增 constructor 只需继承新的 Accepted* 分支，即 Expression Problem
%%   immunity）。
%%
%% 精确差异：
%%   SF-calculus 把 intensional observation 做成了 object-level 的
%%   primitive reduction rule（F 是一个 term）；
%%   inheritance-calculus 把它留在了 meta-level（观察者总能通过
%%   properties(p) 查询任何路径），而 term-level 的分支需要通过
%%   Acceptance pattern 显式编码。
%%   换言之：inheritance-calculus 中 F 的 *读取* 能力是免费的
%%   （所有值都是透明 record），但 F 的 *分支* 能力需要 Acceptance。
%%   SF-calculus 的 F 则把读取和分支捆绑在一个不可扩展的 primitive 里。
%%   Acceptance 是可扩展的（新增 constructor 只需继承新分支，
%%   即 Expression Problem immunity），而 F 硬编码了固定的 atom 集合
%%   （S 和 F）。
%%
%% 关于 Jay 所谓的 "trivial solution of the Halting Problem"：
%%   Jay 用非标准编码将每个程序（包括不停机的）映射到闭合范式（静态数据），
%%   然后结构等式判定变得平凡。这并不推翻图灵定理——它改变了问题的
%%   形式化方式（把程序当静态数据而非可执行代码），而非解决了原来的问题。
%%   Jay 的真正论点是：λ-definability 和 Turing computability 不是同一回事，
%%   Church-Turing Thesis 的等价性证明隐含了编码假设。
%%
%% 正文篇幅不够论述 embeds SF-calculus，暂不展开。
%% 如果审稿人提到 SF-calculus，可以基于上述分析撰写回复或补充段落。
%% ================================================================

\paragraph{\bilingual{Semantic descriptions of the $\lambda$-calculus}{$\lambda$-演算的语义描述}}
\ifInheritanceChinese
$\lambda$-演算的语义可以用几种根本上不同的方式来定义。\emph{操作语义}将 $\beta$-归约 $(\lambda x.\,M)\,N \to M[N/x]$ 作为计算的基本概念~\cite{church1936-lambda-calculus,church1941-calculi-of-lambda-conversion,barendregt1984-lambda-calculus}。这要求捕获-回避替换,其微妙之处促使了 de~Bruijn 索引~\cite{debruijn1972-nameless-dummies} 与显式替换演算~\cite{abadi1991-explicit-substitutions} 的出现。它还要求选择求值策略:按名调用与按值调用~\cite{plotkin1975-call-by-name-call-by-value},或 Abramsky 的惰性策略~\cite{abramsky1990-lazy-lambda-calculus}。\emph{指称语义}~\cite{scott1970-mathematical-theory-of-computation,scott1976-data-types-as-lattices} 将项解释于一个数学域中,但赋予 $\beta$-归约意义需要一个自反域 $D \cong [D \to D]$,由 Cantor 定理知该域无集合论解。Scott 的 $D_\infty$ 通过将 $D$ 构造为连续格的逆极限~\cite{scott1972-continuous-lattices} 来解决这一问题。由此得到的语义是适当的,但不是完全抽象的~\cite{plotkin1977-lcf-as-programming-language,milner1977-fully-abstract-models}:它将观察等价所区分的项等同了起来。\emph{博弈语义}~\cite{abramsky2000-full-abstraction-pcf,hyland2000-full-abstraction-pcf,nickau1994-hereditarily-sequential-functionals} 对 PCF 与惰性 $\lambda$-演算~\cite{abramsky1993-full-abstraction-lazy-lambda,abramsky1995-games-full-abstraction-lazy-lambda} 实现了完全抽象,需要竞技场、策略与无害性条件。在这些方法的每一种中,函数空间 $[D \to D]$ 都是核心障碍:操作语义通过替换隐式地需要它,而指称语义则显式地构造它,博弈语义则通过策略加以精化。这些方法中的每一种都涉及高阶结构%
\footnote{
  指称语义中的函数空间,博弈语义中的策略,以及操作语义中的替换。
}
并在其语义域中保留了函数空间。在每种情形下,指称帐和操作帐都是相互独立的形式体系。三者都继承了 $\beta$-归约在递归程序上的发散:例如,有向图上一个朴素的传递闭包在每个框架下都会发散。
\else
The semantics of the $\lambda$-calculus can be defined in
several fundamentally different ways.
\emph{Operational semantics} takes $\beta$-reduction
$(\lambda x.\,M)\,N \to M[N/x]$ as the primitive notion of
computation~\cite{church1936-lambda-calculus,church1941-calculi-of-lambda-conversion,barendregt1984-lambda-calculus}.
This requires capture-avoiding substitution, a mechanism whose
subtleties motivated de~Bruijn
indices~\cite{debruijn1972-nameless-dummies} and explicit substitution
calculi~\cite{abadi1991-explicit-substitutions}.
It also requires a choice of evaluation
strategy: call-by-name versus
call-by-value~\cite{plotkin1975-call-by-name-call-by-value} or the lazy
strategy of Abramsky~\cite{abramsky1990-lazy-lambda-calculus}.
\emph{Denotational semantics}~\cite{scott1970-mathematical-theory-of-computation,scott1976-data-types-as-lattices}
interprets terms in a mathematical domain, but giving meaning
to $\beta$-reduction requires a reflexive domain
$D \cong [D \to D]$, which has no set-theoretic solution by
Cantor's theorem. Scott's $D_\infty$ resolves this by
constructing $D$ as an inverse limit of continuous
lattices~\cite{scott1972-continuous-lattices}.
The resulting semantics is adequate but not fully
abstract~\cite{plotkin1977-lcf-as-programming-language,milner1977-fully-abstract-models}:
it identifies terms that observational equivalence
distinguishes.
\emph{Game semantics}~\cite{abramsky2000-full-abstraction-pcf,hyland2000-full-abstraction-pcf,nickau1994-hereditarily-sequential-functionals}
achieves full abstraction for PCF and for the lazy
$\lambda$-calculus~\cite{abramsky1993-full-abstraction-lazy-lambda,abramsky1995-games-full-abstraction-lazy-lambda},
requiring arenas, strategies, and innocence conditions.
In each of these approaches, the function space $[D \to D]$
is the central obstacle: operational semantics needs it
implicitly via substitution. In contrast, denotational semantics
constructs it explicitly, and game semantics refines it via
strategies.
Each of these approaches involves higher-order structure%
\footnote{
  Function spaces
  in denotational semantics, strategies
  in game semantics, and substitution
  in operational semantics.
}
and retains the function space in its semantic domain.
The denotational and operational accounts remain separate
formalisms in each case.
All three inherit $\beta$-reduction's divergence on
recursive programs: a na\"ive transitive closure on a cyclic
graph, for instance, diverges in every framework.
\fi

\paragraph{\bilingual{Set-theoretic models of the $\lambda$-calculus}{$\lambda$-演算的集合论模型}}
\ifInheritanceChinese
$\lambda$-演算的几个模型族通过在集合论而非域论框架下工作,避免了 Scott 的 $D_\infty$ 的逆极限构造。这些模型消除了 Scott 拓扑,但并未消除函数空间本身:它们仍然是高阶的,通过注入、交叉类型或多重集关系将 $[D \to D]$ 编码于 $D$ 之中。\emph{图模型}~\cite{engeler1981-algebras-and-combinators,plotkin1993-set-theoretic-lambda-models} 通过幂集格上的 web 注入来编码函数;Bucciarelli 与 Salibra~\cite{bucciarelli2008-graph-lambda-theories} 证明了最大的合理图理论等于 Böhm 树理论 $\mathcal{B}$。\emph{过滤器模型}~\cite{barendregt1983-filter-lambda-model,coppo1978-intersection-type-assignment} 将一个项解释为可从中推导出的交叉类型集合;Dezani-Ciancaglini 等人~\cite{dezani1998-intersection-types-bohm-trees} 证明了 BCD 过滤器模型的 $\lambda$-理论恰好与 Böhm 树相等性重合。\emph{关系模型}~\cite{ehrhard2005-finiteness-spaces,bucciarelli2001-phase-semantics-exponentials} 通过有限多重集从线性逻辑的关系语义中涌现;Ehrhard~\cite{ehrhard2005-finiteness-spaces} 证明了在这一范畴中,纯 $\lambda$-演算的标准自反对象不能存在。这三个模型族都是语义模型:它们将 $\lambda$-项解释于一个外部数学结构中,并在该结构中将 $\beta$-归约验证为一个等式,而不是定义一个具有自身原语的独立演算。每一种都编码函数空间而非消除它,且每一种都需要高阶机制才能为惰性 $\lambda$-演算获得完全抽象。
\else
Several families of $\lambda$-calculus models avoid the
inverse-limit construction of Scott's $D_\infty$ by working
in set-theoretic rather than domain-theoretic settings.
These models eliminate Scott topology but not the function
space itself: they remain higher-order, encoding $[D \to D]$
inside $D$ via injections, intersection types, or multiset
relations.
\emph{Graph models}~\cite{engeler1981-algebras-and-combinators,plotkin1993-set-theoretic-lambda-models}
encode functions via a web injection on powerset lattices;
Bucciarelli and Salibra~\cite{bucciarelli2008-graph-lambda-theories}
proved that the greatest sensible graph theory equals the
B\"ohm tree theory~$\mathcal{B}$.
\emph{Filter models}~\cite{barendregt1983-filter-lambda-model,coppo1978-intersection-type-assignment}
interpret a term as the set of intersection types derivable
for it;
Dezani-Ciancaglini et al.~\cite{dezani1998-intersection-types-bohm-trees}
showed that the BCD filter model's $\lambda$-theory
coincides exactly with B\"ohm tree equality.
\emph{Relational models}~\cite{ehrhard2005-finiteness-spaces,bucciarelli2001-phase-semantics-exponentials}
arise from the relational semantics of linear logic via
finite multisets;
Ehrhard~\cite{ehrhard2005-finiteness-spaces} showed that a standard
reflexive object for the pure $\lambda$-calculus cannot
exist in this category.
All three families are semantic models: they interpret
$\lambda$-terms in an external mathematical structure and
validate $\beta$-reduction as an equation in that structure,
rather than defining an independent calculus with its own
primitives.
Each encodes the function space rather than eliminating it,
and each requires higher-order machinery to obtain full
abstraction for the lazy $\lambda$-calculus.
\fi

\paragraph{\bilingual{Recursive equations and fixpoint semantics}{递归方程与不动点语义}}
\ifInheritanceChinese
Van~Emden 与 Kowalski~\cite{vanemden1976-predicate-logic-semantics} 证明了谓词逻辑作为编程语言的语义允许三种等价的刻画:操作性的\footnote{通常是 SLD 归结。}、模型论的\footnote{基于最小 Herbrand 模型。}以及不动点的\footnote{即基原子幂集上即时后承算子 $T_P$ 的最小不动点。}。Aczel~\cite{aczel1977-inductive-definitions} 通过幂集格上的单调算子对归纳定义给出了系统性处理,为 Datalog 风格的语义所依赖的逻辑基础奠定了基础。Leroy 与 Grall~\cite{leroy2006-coinductive-big-step} 通过余归纳定义处理了大步语义中的发散,需要在归纳求值判断之外引入一个单独的余归纳判断来说明非终止。Van~Emden 与 Kowalski 的 $T_P$ 在一阶域上为递归 Datalog 程序计算 $\mathrm{lfp}$。传统的 $\lambda$-演算语义在语义域中需要函数空间,且递归的 $\lambda$-项在 $\beta$-归约下通常会发散;利用一阶域上的 $T_P$ 为 $\lambda$-演算赋予一种改变哪些项收敛的不动点语义,尚未被探索。在近期工作中,Yang~\cite{yang2026-co-lambda} 将 tabling 求值~\cite{tamaki1986-tabled-resolution} 直接引入纯 $\lambda$-演算:首部规范化被读作一个余代数方程,其状态为被驻留的 $\lambda$-项,即具有可判定同一性的一阶对象,而 tabling 求解该方程,在有限时间内将 $\Omega$ 这样的非生产性循环判定为 $\bot$。该处理保留了标准的惰性 (Lévy--Longo) 树语义,仅改变了求值器;它占据了图~\ref{fig:calculi-matrix} 中不动点 $\lambda$ 节点。它在机制上也有所不同。表以整个被驻留的项为键,而非以进入作用域树的路径为键,因此在路径前缀上重叠的查询之间不共享部分结果。偏序不是集合包含下的幂集格,而是树近似序,其中 $\bot$ 叶被精化为子树,且不动点是唯一的(最小与最大因有护性而重合):一个重入调用被以单一的最终 $\bot$ 回答,而后续任何迭代都不会精化该 $\bot$,因此迭代无法增大一个累加器。其比惰性 $\lambda$-演算更确定的部分因此恰好是图~\ref{fig:calculi-matrix} 中非生产性的 $\bot$,即发散被判定为确定 $\bot$ 答案的那些项,别无其他;而幂集序才使得迭代可以积累并产生新的非 $\bot$ 值。
\else
Van~Emden and Kowalski~\cite{vanemden1976-predicate-logic-semantics} showed that
the semantics of predicate logic as a programming language admits
three equivalent characterizations:
operational\footnote{
  Typically SLD resolution.
}, model-theoretic\footnote{
  Based on least Herbrand models.
}, and fixed-point\footnote{
  Namely, the least fixed point of the immediate consequence operator~$T_P$ on the powerset of ground atoms.
}.
Aczel~\cite{aczel1977-inductive-definitions} gave a systematic treatment of
inductive definitions via monotone operators on powerset lattices,
providing the logical foundations that Datalog-style semantics
rests on.
Leroy and Grall~\cite{leroy2006-coinductive-big-step} treated divergence
in big-step semantics via coinductive definitions,
requiring a separate coinductive judgment to account for
non-termination alongside the inductive evaluation judgment.
Van~Emden and Kowalski's $T_P$ computes
$\mathrm{lfp}$ for recursive Datalog programs over a
first-order domain.
Conventional $\lambda$-calculus semantics requires function
spaces in the semantic domain, and recursive $\lambda$-terms
ordinarily diverge under $\beta$-reduction;
using $T_P$ over a first-order domain to give the
$\lambda$-calculus a fixpoint semantics that changes which
terms converge has not been explored.
In recent work, Yang~\cite{yang2026-co-lambda} brings tabled
evaluation~\cite{tamaki1986-tabled-resolution} to the pure
$\lambda$-calculus directly: head normalisation is read as a
coalgebraic equation whose states are interned
$\lambda$-terms, first-order objects with decidable identity,
and tabling solves it, deciding unproductive loops such
as~$\Omega$ as~$\bot$ in finite time.
That treatment preserves the standard lazy
(L\'evy--Longo) tree semantics and changes only the evaluator;
it occupies the fixpoint~$\lambda$ node of
Figure~\ref{fig:calculi-matrix}.
It also differs in mechanism.
The table is keyed by whole interned terms, not by paths
into a scope tree, so partial results are not shared between
queries that overlap on a path prefix.
The order is not a powerset lattice under set inclusion but
the tree approximation order, in which $\bot$ leaves are
refined into subtrees, and the fixpoint is unique (least and
greatest coincide by guardedness): a re-entrant call is
answered with a single final~$\bot$ that no later iteration
refines, so the iteration cannot grow an accumulator.
Its more-definedness over the lazy $\lambda$-calculus is
therefore exactly the unproductive~$\bot$ of
Figure~\ref{fig:calculi-matrix}, terms whose divergence is
decided as a definite~$\bot$ answer, and nothing more;
the powerset order is what lets iteration accumulate and
produce new non-$\bot$ values.
\fi

\subsection{\bilingual{Inheritance, Objects, and Records}{继承、对象与记录}}

\paragraph{\bilingual{Denotational semantics of inheritance}{继承的指称语义}}
\ifInheritanceChinese
Cook~\cite{cook1989-denotational-inheritance} 给出了继承的第一个指称语义,将对象建模为递归记录。继承是\emph{生成器}%
\footnote{
  这些是从 self 到完整对象的函数。
}与\emph{包装器}%
\footnote{
  包装器是修改生成器的函数。
}的复合,由一个不动点构造来求解。Cook 的关键洞见是,继承是一种可适用于任何形式的递归定义的通用机制,而不仅仅是面向对象的方法。这一洞见是本工作的出发点之一。该模型有四个原语(记录、函数、生成器和包装器),并需要一个不动点组合子。复合是不对称的:包装器修改生成器,反之则不然。不动点构造依赖函数域中的 Scott 连续性。由此得到的指称不能直接执行,因此需要一个单独的操作语义来在运行时定义方法分派与对象构造。
\else
Cook~\cite{cook1989-denotational-inheritance} gave the first denotational
semantics of inheritance, modeling objects as recursive records.
Inheritance is composition of \emph{generators}
\footnote{
  These are functions from self to complete object.
} and
\emph{wrappers}\footnote{
  Wrappers are functions that modify generators.
},
solved by a fixed-point construction.
Cook's key insight is that inheritance is a general mechanism
applicable to any form of recursive definition, not only
object-oriented methods. This insight is one of the starting points of the
present work.
The model has four primitives (records, functions, generators, and
wrappers) and requires a fixed-point combinator. Composition is
asymmetric: a wrapper modifies a generator but not vice versa.
The fixed-point construction relies on Scott continuity in a
domain of functions. The resulting denotations are not directly
executable, so a separate operational semantics is needed to
define method dispatch and object construction at runtime.
\fi

\paragraph{\bilingual{Mixins and traits}{Mixin 与 trait}}
\ifInheritanceChinese
Bracha 与 Cook~\cite{bracha1990-mixin-inheritance} 将 mixin 形式化为抽象子类,即从超类参数到子类的函数。由于 mixin 应用是函数复合,它\emph{既不可交换也不幂等};C3 线性化算法~\cite{barrett1996-c3-linearization} 后来被开发出来以施加一个确定性顺序。Schärli 等人~\cite{scharli2003-traits} 引入了 trait,其复合是可交换的,但拒绝同名冲突~\cite{ducasse2006-traits-fine-grained-reuse}。两种机制都在\emph{方法层级}操作:对同名的两个定义进行复合会导致覆盖或冲突,而非嵌套结构的递归合并。两者都预设了 $\lambda$-演算:方法体是函数,而复合机制本身不是图灵完备的。Mixin 与 trait 均不支持深合并。
\else
Bracha and Cook~\cite{bracha1990-mixin-inheritance} formalized mixins as
abstract subclasses, that is, functions from a superclass parameter to a
subclass.
Because mixin application is function composition, it is
\emph{neither commutative nor idempotent}; the C3 linearization
algorithm~\cite{barrett1996-c3-linearization} was later developed to impose a
deterministic order.
Sch\"arli et al.~\cite{scharli2003-traits} introduced traits,
whose composition is commutative but rejects same-name
conflicts~\cite{ducasse2006-traits-fine-grained-reuse}.
Both mechanisms operate at the \emph{method level}: composing two
definitions of the same name results in an override or a conflict,
not a recursive merge of nested structure.
Both presuppose the $\lambda$-calculus: method bodies are functions,
and the composition mechanism itself is not Turing complete.
Neither mixins nor traits are deep-mergeable.
\fi

\paragraph{\bilingual{Object and record calculi}{对象演算与记录演算}}
\ifInheritanceChinese
Abadi 与 Cardelli~\cite{abadi1996-theory-of-objects} 将对象作为基本概念,但方法是以 self 为参数的函数,且对象扩展是不对称的。Boudol~\cite{boudol2004-recursive-record-semantics} 证明了自引用记录需要一个不安全的不动点算子,其适定性依赖于求值策略。Harper 与 Pierce~\cite{harper1991-symmetric-record-concatenation} 给出了一个具有对称连接的记录演算,该演算拒绝同标签的复合;Cardelli~\cite{cardelli1992-extensible-records} 研究了带子类型的可扩展记录;Rémy~\cite{remy1989-row-polymorphism} 为 ML 中的记录和变体引入了行多态。这些演算全都在\emph{平坦}记录上操作:没有嵌套结构的递归合并,没有惰性观察,$\this$ 也需要一个显式的不动点组合子。
\else
Abadi and Cardelli~\cite{abadi1996-theory-of-objects} treat objects as the
primitive notion, but methods are functions that take self as a
parameter and object extension is asymmetric.
Boudol~\cite{boudol2004-recursive-record-semantics} showed that self-referential
records require an unsafe fixed-point operator whose
well-definedness depends on the evaluation strategy.
Harper and Pierce~\cite{harper1991-symmetric-record-concatenation} gave a record
calculus with symmetric concatenation that rejects same-label
composition;
Cardelli~\cite{cardelli1992-extensible-records} studied extensible records
with subtyping;
R\'emy~\cite{remy1989-row-polymorphism} introduced row polymorphism for
records and variants in ML.
All of these calculi operate on \emph{flat} records: there is
no recursive merging of nested structure, no lazy observation,
and $\this$ requires an explicit fixed-point combinator.
\fi

\paragraph{\bilingual{Family polymorphism and DOT}{族多态与 DOT}}
\ifInheritanceChinese
Ernst~\cite{ernst2001-family-polymorphism} 引入了族多态;Ernst、Ostermann 与 Cook~\cite{ernst2006-virtual-class-calculus} 在虚类演算中将其形式化,其中类通过路径表达式 \texttt{this.out.C} 来访问。Amin 等人~\cite{amin2016-dependent-object-types} 在 DOT 演算中形式化了路径依赖类型,这是 Scala 的理论基础。在这两个框架中,每条路径都解析为\emph{单一}封闭对象中的\emph{单一}类。当多条继承路径通向同一作用域时,单值解析必须选择其中一条路径并丢弃其余;DOT 的 self 变量 $\nu(x : T)\, d$ 将 $x$ 绑定到单一对象,而 Scala 在编译时拒绝由此产生的交叉类型(附录~\ref{app:scala-multi-target})。两个框架均未提供集合值解析的机制。
\else
Ernst~\cite{ernst2001-family-polymorphism} introduced family polymorphism;
Ernst, Ostermann, and Cook~\cite{ernst2006-virtual-class-calculus}
formalized it in the virtual class calculus, where classes are
accessed via path expressions \texttt{this.out.C}.
Amin et al.~\cite{amin2016-dependent-object-types} formalized path-dependent types
in the DOT calculus, the theoretical foundation of Scala.
In both frameworks, each path resolves to a
\emph{single} class in a \emph{single} enclosing object.
When multiple inheritance routes lead to the same scope,
single-valued resolution must choose one route and discard the
others; DOT's self variable $\nu(x : T)\, d$ binds $x$ to a
single object, and Scala rejects the resulting intersection
type at compile time
(Appendix~\ref{app:scala-multi-target}).
Neither framework provides a mechanism for set-valued
resolution.
\fi

\subsection{\bilingual{Configuration Languages and Module Systems}{配置语言与模块系统}}

\paragraph{\bilingual{Configuration languages}{配置语言}}
\ifInheritanceChinese
CUE~\cite{cue2019} 在单一格中统一了类型与值,其中以 \texttt{\&} 表示的合一是可交换、可结合且幂等的。CUE 的设计受到 Google 的 Borg 配置语言 (GCL) 的启发,后者使用了图合一。CUE 的格包含标量类型,其冲突语义规定合一不相容标量会得到 $\bot$;标量对于配置语言是否必要,还是仅凭树结构就已足够,是一个有趣的设计问题。Jsonnet~\cite{jsonnet} 通过 \texttt{+} 运算符提供了具有 mixin 语义的对象继承:复合\emph{不是}可交换的,因为在标量冲突上右侧胜出,而深合并需要在每个字段上通过 \texttt{+:} 语法显式选择加入,因此深合并虽可用但并非默认行为。Dhall~\cite{dhall2017} 采取函数式方法:它是一个带记录的类型化 $\lambda$-演算,其中复合是记录合并算子,而非继承机制。
\else
CUE~\cite{cue2019} unifies types and values in a single lattice
where unification, denoted by \texttt{\&}, is commutative, associative, and
idempotent.
CUE's design was motivated by Google's Borg Configuration
Language (GCL), which used graph unification.
CUE's lattice includes scalar types with a conflict
semantics where unifying incompatible scalars yields $\bot$; whether
scalars are necessary for a configuration language, or whether
tree structure alone suffices, is an interesting design question.
Jsonnet~\cite{jsonnet} provides object inheritance via the
\texttt{+} operator with mixin semantics: composition is
\emph{not} commutative since the right-hand side wins on scalar
conflicts, and deep merging requires explicit opt-in via the \texttt{+:}
syntax on a per-field basis, so deep merge is
available but not the default.
Dhall~\cite{dhall2017} takes a functional approach: it is a typed
$\lambda$-calculus with records, where composition is a record merge
operator, not an inheritance mechanism.
\fi

\paragraph{\bilingual{Module systems and deep merge}{模块系统与深合并}}
\ifInheritanceChinese
NixOS 模块系统~\cite{dolstra2010-nixos,nixosmodules} 通过\emph{递归合并嵌套属性集}来组合模块,而非通过线性化或方法层级的覆盖,这与先前的 mixin 和 trait 演算中区分树层级继承与方法层级复合的深合并语义相同。当两个模块定义了相同的嵌套路径时,其子树被合并而非覆盖。NixOS 模块系统集成了深合并、递归的 $\this$、模块的延迟求值,以及其上的丰富类型检查层。同一机制为 NixOS 系统配置、Home Manager~\cite{homemanager}、nix-darwin~\cite{nixdarwin}、flake-parts~\cite{flakeparts} 以及数十个其他生态系统提供支持,共同管理任意复杂度的配置。一个似乎尚未被完全解决的方面是 $\this$ 在多重继承下的语义。当两个各自继承自同一延迟模块的独立模块被组合时,每次导入都会引入对共享选项自己的声明;模块系统的 \texttt{merge\-Option\-Decls} 以静态错误拒绝重复声明(附录~\ref{app:scala-multi-target})。模块系统无法表达两条路径对同一选项的贡献是均等的,即递归合并自然产生的多目标情形,正是这一情形促成了第~\ref{sec:mixin-trees} 节的观测语义。
\else
The NixOS module system~\cite{dolstra2010-nixos,nixosmodules}
composes modules by \emph{recursively merging nested attribute sets},
not by linearization or method-level override, which is the same deep-merge
semantics that distinguishes tree-level inheritance from method-level
composition in prior mixin and trait calculi.
When two modules define the same nested path, their subtrees are
merged rather than overridden.
The NixOS module system incorporates deep merge,
recursive $\this$, deferred evaluation of modules, and a
rich type-checking layer on top.
The same mechanism powers
NixOS system configuration, Home
Manager~\cite{homemanager}, nix-darwin~\cite{nixdarwin},
flake-parts~\cite{flakeparts}, and dozens of other ecosystems,
collectively managing configurations of arbitrary complexity.
The one aspect that appears not to have been fully addressed is the
semantics of $\this$ under multiple inheritance.
When two independent modules that inherit from the same
deferred module are composed, each import introduces its own
declaration of the shared options; the module system's
\texttt{merge\-Option\-Decls} rejects duplicate declarations
with a static error (Appendix~\ref{app:scala-multi-target}).
The module system cannot express that both routes contribute
equally to the same option, that is, the multi-target situation that
recursive merging naturally gives rise to, which motivated the
observational semantics of Section~\ref{sec:mixin-trees}.
\fi

\subsection{\bilingual{Our Companion Work}{我们的配套工作}}

\paragraph{\bilingual{Modular software architectures}{模块化软件架构}}
\label{sec:practical-modular}
\ifInheritanceChinese
插件系统、组件框架与依赖注入框架通过声明式继承实现了软件的可扩展性。第~\ref{sec:mixin-trees} 节所确立的继承的对称性与可结合性解释了这类系统为何有效:组件可以以任意顺序继承而不改变行为。MIXINv2\anon[~(补充材料)]{~\cite{mixin2025}} 本身就是一个依赖注入框架:每个作用域将其依赖声明为具名槽,复合跨越作用域边界按名解析它们,从而使继承演算的模式显式化而非随意拼凑。
\else
Plugin systems, component frameworks, and dependency injection frameworks
enable software extensibility through declarative inheritance. The
symmetric, associative nature of inheritance established in Section~\ref{sec:mixin-trees} explains why such
systems work: components can be inherited in any order without changing
behavior.
MIXINv2\anon[~(supplementary material)]{~\cite{mixin2025}}
is itself a dependency injection framework:
each scope declares its dependencies as named slots, and
composition resolves them by name across scope boundaries,
making the inheritance-calculus patterns explicit rather than ad~hoc.
\fi

\paragraph{\bilingual{Union file systems}{联合文件系统}}
\label{sec:practical-union-fs}
\ifInheritanceChinese
联合文件系统,如 UnionFS、OverlayFS 与 AUFS,将多个目录层次叠加以呈现统一视图。其语义与继承演算相似,表现在后面的层覆盖前面的层且所有层的文件均可访问,但其设计在具体之处有所不同:继承是不对称的,标量文件的冲突解决是不可交换的,且层与层之间不存在晚绑定或动态分派。我们将继承演算直接实现为联合文件系统\anon[~(作为补充材料附上)]{~\cite{ratarmount2025}},仅凭三个原语就获得了文件查询的晚绑定语义与动态分派,无需任何额外机制。在软件包管理器、构建系统与操作系统发行版中,外部工具链将配置转换为文件树;而在这里,文件系统作为自身的配置,转换即是继承:这是一种\emph{文件系统即编译器},其中目录树既是源程序又是编译目标。
\else
Union file systems, such as UnionFS, OverlayFS, and AUFS, layer multiple directory
hierarchies to present a unified view. Their semantics resembles
inheritance-calculus in that later layers override earlier layers and files from all
layers remain accessible, but their designs differ in specific ways:
inheritance is asymmetric, conflict resolution for scalar files
is noncommutative, and there is no late-binding or dynamic
dispatch across layers.
We implemented inheritance-calculus directly as a union file
system\anon[~(included as supplementary material)]{~\cite{ratarmount2025}},
obtaining late-binding semantics and dynamic dispatch for file lookups
with no additional mechanism beyond the three primitives.
In package managers, build systems, and OS distributions,
external toolchains transform configuration into file trees;
here, the file system serves as configuration to itself,
and the transformation is inheritance: a
\emph{file-system-as-compiler} in which the directory tree is
both the source program and the compilation target.
\fi

\section{\bilingual{Future Work}{未来工作}}
\label{sec:future-work}

\ifInheritanceChinese
\anon[辅助实现]{MIXINv2~\cite{mixin2025}}{} 是继承演算的一个可执行实现,它已经超越了本文所呈现的无类型演算:它包含编译期检查,以验证所有引用都解析到 mixin 树中的有效路径,以及一个从宿主语言引入标量值的外部函数接口 (FFI)。以下所描述的未来工作,关注的是这一实现与一门完全实用的语言之间的差距,横跨理论基础与库层面的编码两个方面。
\else
\anon[The supplementary implementation]{MIXINv2~\cite{mixin2025}}{} is an executable implementation of inheritance-calculus
that already goes beyond the untyped calculus presented here:
it includes compile-time checking that all references resolve to
valid paths in the mixin tree, and a foreign-function interface (FFI)
that introduces scalar values from the host language.
The future work described below concerns the distance between this
implementation and a fully practical language, spanning both
theoretical foundations and library-level encodings.
\fi

\subsection{\bilingual{Type System}{类型系统}}
\ifInheritanceChinese
MIXINv2 现有的编译期检查验证每条引用路径都解析到继承结果中存在的属性。将这些检查的可靠性形式化,即证明类型良好的程序在运行时不产生悬空引用,是未来工作的一个方向。这样的形式化与 DOT 演算~\cite{amin2016-dependent-object-types} 的关系,类似于继承演算与 $\lambda$-演算的关系:有类型的继承演算可以看作不含 $\lambda$ 的 DOT,它保留了路径依赖类型,同时以 mixin 继承取代函数。正式的类型系统将适合在 Coq 和 Agda 等证明助手中机械化,以验证类型安全、进展和保持等性质。

在现有检查的可靠性之外,额外的类型系统特性将强化这门语言。一个关键的例子是\emph{完全性检查}:验证一个继承为所有必需槽位提供了实现,而不仅仅是引用能够解析。在无类型的继承演算中,写下 $\mathrm{magnitude} \mapsto \{\}$ 纯属说明性质:它定义了一个结构槽位,但不对填充它的内容施加任何约束。完全性检查将在静态上强制这类声明,拒绝使必需槽位保持未填充状态的继承。
\else
MIXINv2's existing compile-time checks verify that every
reference path resolves to a property that exists in the inherited
result. Formalizing the soundness of these checks, that is, proving that
well-typed programs do not produce dangling references at
runtime, is one direction of future work.
Such a formalization would relate to the DOT
calculus~\cite{amin2016-dependent-object-types} in a manner analogous to how
inheritance-calculus relates to the $\lambda$-calculus: a Typed
inheritance-calculus can be viewed as DOT without $\lambda$, retaining
path-dependent types while replacing functions with
mixin inheritance.
A formal type system would be amenable to mechanization in proof
assistants, such as Coq and Agda, for verifying type safety, progress, and
preservation properties.

Beyond soundness of the existing checks, additional type system
features would strengthen the language.
A key example is \emph{totality checking}: verifying that an
inheritance provides implementations for all required slots, not
merely that references resolve.
In untyped inheritance-calculus, writing $\mathrm{magnitude} \mapsto \{\}$
is purely documentary: it defines a structural slot but imposes no
constraint on what must fill it.
Totality checking would enforce such declarations statically,
rejecting inheritances that leave required slots unfilled.
\fi

\subsection{\bilingual{Standard Library}{标准库}}
\ifInheritanceChinese
其余几项涉及标准库:它们不需要对演算或类型系统进行扩展,但需要在现有框架内进行非平凡的编码。

继承演算\emph{默认是开放的}:继承可以自由合并任意两个 mixin,一个值可以同时具有多个数据构造子(例如,$\{\mathrm{Zero},\; \mathrm{Odd},\; \mathrm{half} \mapsto \{\mathrm{Zero}\}\}$)。这是该演算自然的 trie 语义,也是第~\ref{sec:expression-problem}~节求解表达式问题的基础。然而,许多操作(例如等值测试)假设值是\emph{线性的},即恰好具有一个数据构造子。第~\ref{sec:case-study}~节的观察者模式并不强制这一点:将其应用于一个具有多个数据构造子的值时,所有回调都会触发,其结果被继承在一起。

\emph{模式验证}可以在库层面使用现有原语来实现。观察者可以测试某个特定的数据构造子是否存在,返回一个布尔值;第二次布尔派发随后基于该结果进行分支。通过链接这样的测试,工厂可以验证一个值恰好匹配一个数据构造子并携带所需的属性。这并不限制继承本身(继承仍然不受约束),但提供了一种在不变式违反传播之前检测到它的方式。

同样的技术,即观察者返回布尔值后接布尔派发,可以产生\emph{封闭模式匹配}。第~\ref{sec:case-study}~节的观察者模式本质上是\emph{开放的}:新的数据构造子情况可以通过继承添加。封闭匹配(其中恰好执行一个分支)可以编码为一个\emph{访问者链}:一个由 if-then-else 节点组成的链表,每个节点测试一个数据构造子,在不匹配时向后传递。这类似于 GHC.Generics 的 $\mathrm{(:+:)}$ 和类型表示~\cite{magalhaes2010-generic-deriving},其中每个链节对应一个求和项。一个重要的推论是,这样的链\emph{不}可交换,因为它们具有固定的顺序,因此无法通过开放继承来扩展,而这正是封闭派发的正确语义。
\else
The remaining items concern the standard library: they require
no extensions to the calculus or type system, but involve nontrivial
encodings within the existing framework.

Inheritance-calculus is \emph{open by default}: inheritance
can freely merge any two mixins, and a value may simultaneously
inhabit multiple data constructors (e.g.,
$\{\mathrm{Zero},\; \mathrm{Odd},\; \mathrm{half} \mapsto \{\mathrm{Zero}\}\}$).
This is the natural trie semantics of the
calculus, and the basis for solving the Expression Problem
in Section~\ref{sec:expression-problem}.
However, many operations, such as equality testing, assume that values
are \emph{linear}, inhabiting exactly one data constructor.
The observer pattern of Section~\ref{sec:case-study} does not enforce
this: applied to a multi-data-constructor value, all callbacks fire and
their results are inherited.

\emph{Schema validation} can be implemented at the library level
using existing primitives.
An observer can test whether a particular data constructor is present,
returning a Boolean; a second Boolean dispatch then branches on the
result.
By chaining such tests, a factory can validate that a value matches
exactly one data constructor and carries the required properties.
This does not restrict inheritance itself, which remains
unconstrained, but provides a way to detect invariant violations
before they propagate.

The same technique, an observer returning Boolean followed by Boolean
dispatch, yields \emph{closed pattern matching}.
The observer pattern of Section~\ref{sec:case-study} is inherently
\emph{open}: new data-constructor cases can be added through inheritance.
Closed matching, where exactly one branch is taken, can be encoded as
a \emph{visitor chain}: a linked list of if-then-else nodes, each
testing one data constructor and falling through on mismatch.
This is analogous to GHC.Generics' $\mathrm{(:+:)}$ sum
representation~\cite{magalhaes2010-generic-deriving}, where each link
corresponds to one summand.
An important consequence is that such chains do \emph{not}
commute, as they have a fixed order, and therefore cannot be extended
through open inheritance, which is the correct semantics for closed
dispatch.
\fi

\section{\bilingual{Conclusion}{结论}}

\ifInheritanceChinese
对不可扩展性的免疫指导了继承演算的设计。第~\ref{sec:semantic-variants}~节排除了每一种替代的组合机制和引用解析,使得带有多目标晚绑定限定 $\this$ 的深合并成为唯一的幸存者。mixin 充当了统一模块、类、对象、方法、字段和局部绑定的单一抽象;mixin 树的深合并使继承具有可交换性、幂等性和可结合性,从而多重继承的线性化问题不再出现。第~\ref{sec:mixin-trees}~节的六个一阶方程,具有指称语义且可由 tabling~\cite{tamaki1986-tabled-resolution} 计算,足以定义语义;观察者通过查询一棵惰性构造的树中的路径来驱动计算。这些方程是对标准面向对象语义的最小偏离:唯一的改变是 $\overrides$ 方程~(\ref{eq:overrides}),它以子树的深合并取代了方法层面的遮蔽。面向对象编程传统上被归类为命令式的;去掉可变状态和函数体,便得到纯面向对象编程,%
\footnote{
  ``纯''具有双重含义:不纯粹意味着不纯洁,即值是不可变的且组合没有副作用;而不混合意味着组合完全依赖继承,没有函数式抽象,例如 $\lambda$ 或 $\beta$-归约。
}
这是声明式编程的一个最小计算模型。

由于 $\overrides$ 合并的是子树而非方法体,函数不再是原语。$\lambda$-演算在第~\ref{sec:translation}~节通过六条翻译规则宏表达到继承演算中。第~\ref{sec:bohm-tree}~节的 Böhm 树对应将这一一阶语义扩展到了 $\lambda$-演算:先前的高阶语义框架被证明忠实于继承演算一阶方程的一个子语言的首次迭代近似。

继承演算进一步获得了 $\lambda$-演算无法表达的三种现象;事实上,正如第~\ref{sec:asymmetry}~节所示,在 Felleisen 意义下~\cite{felleisen1991-expressive-power},继承演算严格地比惰性 $\lambda$-演算更具表达力。第~\ref{sec:expression-problem}~节在表达式问题~\cite{wadler1998-expression-problem}意义下建立了对不可扩展性的免疫。第~\ref{sec:case-study}~节的案例研究表明,普通算术产生了逻辑编程的关系语义:循环继承层次结构通过与递归 Datalog 所支配的相同的升链条件收敛,附录~\ref{app:datalog-encoding}给出了一个系统性的编码。第~\ref{sec:practical-function-color}~节展示了函数颜色盲性~\cite{nystrom2015-function-color}。

单个 mixin 同时是数据、程序、配置、函数、关系和对象。它宏表达了 $\lambda$-演算,并恢复了逻辑编程的关系语义,同时对两者均无法逃脱的不可扩展性保持免疫。

引言提出了配置本身是否能成为实用的通用编程的问题。答案是肯定的。我们提出继承演算作为实用通用编程的基础,我们相信它优于任何已知的计算模型。
\else
Immunity to nonextensibility guided the design of
inheritance-calculus.
Section~\ref{sec:semantic-variants} eliminated every
alternative composition mechanism and reference resolution,
leaving deep merge with multi-target late-binding qualified
$\this$ as the unique survivor.
The mixin served as the single abstraction unifying module,
class, object, method, field, and local binding;
deep merge of mixin trees made inheritance commutative, idempotent, and associative,
so the linearization problem of multiple inheritance did not arise.
The six first-order equations of Section~\ref{sec:mixin-trees},
denotational and computable by
tabling~\cite{tamaki1986-tabled-resolution}, suffice to
define the semantics; the observer drives computation by
querying paths in a lazily constructed tree.
These equations are a minimal departure from standard
object-oriented semantics: the only change is the $\overrides$ equation~(\ref{eq:overrides}),
which replaces method-level shadowing with deep merge of subtrees.
Object-oriented programming is traditionally classified as
imperative; removing mutable state and function bodies yields
pure object-oriented programming,%
\footnote{
  ``Pure'' carries a double meaning:
  not impure, since values are immutable and composition has no
  side effects; and not mixed, since composition relies entirely on
  inheritance, without functional abstractions such
  as~$\lambda$ or~$\beta$-reduction.
}
a minimal computational model for declarative programming.

Since $\overrides$ merges subtrees rather than method bodies,
functions are no longer primitive.
The $\lambda$-calculus macro-expressed into
inheritance-calculus in six translation rules in
Section~\ref{sec:translation}.
The B\"ohm tree correspondence of
Section~\ref{sec:bohm-tree} extended this first-order
semantics to the $\lambda$-calculus:
the prior higher-order semantic frameworks proved faithful
to the first-iteration approximation of a sublanguage of
inheritance-calculus's first-order equations.

Inheritance-calculus further gains three phenomena
that the $\lambda$-calculus cannot express;
indeed, inheritance-calculus is strictly more expressive
than the lazy $\lambda$-calculus in Felleisen's
sense~\cite{felleisen1991-expressive-power},
as shown in Section~\ref{sec:asymmetry}.
Section~\ref{sec:expression-problem} established immunity
to nonextensibility in the sense of the
Expression Problem~\cite{wadler1998-expression-problem}.
The case study of Section~\ref{sec:case-study} showed that
ordinary arithmetic yields the relational semantics of logic
programming: cyclic inheritance hierarchies converge by
the same ascending chain condition that governs recursive
Datalog, and Appendix~\ref{app:datalog-encoding} gave a
systematic encoding.
Section~\ref{sec:practical-function-color} exhibited
function color
blindness~\cite{nystrom2015-function-color}.

A single mixin is at once data, program, configuration,
function, relation, and object.
It macro-expresses the $\lambda$-calculus and recovers the
relational semantics of logic programming, while remaining
immune to the nonextensibility that neither escapes.

The introduction asked whether configuration alone can be
practical general-purpose programming.
It can.
We propose inheritance-calculus as a foundation for practical
general-purpose programming, one we believe better than any
known computational model.
\fi

\section{\bilingual{Data Availability Statement}{数据可用性声明}}

\ifInheritanceChinese
可执行实现 (MIXINv2) 以及本文中的所有示例以 MIT 许可证作为开源软件公开发布\anon[{} (匿名审查期间省略 URL)]{~\cite{mixin2025}}。每个以继承演算记法呈现的示例都有对应的可执行 MIXINv2 源码,且一套测试套件复现了论文中的每个示例。案例研究 (第~\ref{sec:case-study}~节) 的布尔逻辑和一元自然数算术、附录~\ref{app:datalog-encoding}~的 Datalog 编码 (包括双跳、常量连接和多变量连接关系程序以及重复访问者自连接回归示例),以及函数颜色盲性段落 (第~\ref{sec:practical-function-color}~节) 中讨论的同步和异步运行时实例,均包含在这些源码中。第~\ref{sec:practical-union-fs}~节所讨论的 ratarmount 补丁\anon[也已公开发布 (匿名审查期间省略 URL)]{作为一个开放的拉取请求~\cite{ratarmount2025}可以获取}。\anon[以上所有内容也作为辅助材料一并附上。]{}
\else
The executable implementation (MIXINv2) and all examples from
this paper are openly available as open-source software under
the MIT license\anon[{} (URL omitted for anonymous review)]{~\cite{mixin2025}}.
Every example rendered in inheritance-calculus notation has a
corresponding executable MIXINv2 source, and a test suite
reproduces every example in the paper.
The boolean logic and unary natural number arithmetic of the case
study (Section~\ref{sec:case-study}), the Datalog encodings of
Appendix~\ref{app:datalog-encoding} (including the two-hop,
constant-join, and multi-variable-join relational programs and the
repeated-visitor self-join regression example), and the synchronous
and asynchronous runtime instances discussed in the
function-color-blindness paragraph
(Section~\ref{sec:practical-function-color}) are all included among
these sources.
The ratarmount patch discussed in
Section~\ref{sec:practical-union-fs} is
\anon[also openly available (URL omitted for anonymous review)]{available as an open pull
request~\cite{ratarmount2025}}.
\anon[All of the above are also included as supplementary material.]{}
\fi

\begin{acks}
  Claude Code, GitHub Copilot, Gemini Code Assist, and Codex were utilized to generate sections of this Work,
  including text, tables, graphs, code, formulas, citations, etc.
\end{acks}
\fi % end \ifInheritanceBody (the main matter)

% References follow the main matter and precede the appendix. This copy fires
% whenever the main matter is present (preprint, submission); the appendix-only
% build emits its copy after the appendix instead.
\ifInheritanceBody
\bibliographystyle{ACM-Reference-Format}
\bibliography{references}
\fi

\ifInheritanceAppendix
\appendix

\section{\bilingual{Well-Definedness of the Semantic Functions}{语义函数的良定义性}}
\label{app:well-definedness}

\ifInheritanceChinese
六个语义方程~(\ref{eq:properties})--(\ref{eq:this}) 是一阶的:其输入与输出均为路径及路径的集合,不涉及函数空间、闭包或高阶构造。
本节通过证明这六个方程定义了一个单调算子来确立良定义性;该算子的最小不动点由 tabling 求值~\cite{tamaki1986-tabled-resolution} 可靠地计算。
\else
The six semantic
equations~(\ref{eq:properties})--(\ref{eq:this}) are
first-order: their inputs and outputs are paths and sets of
paths, with no function spaces, closures, or higher-order
constructs.
This section establishes well-definedness by showing that
the six equations define a monotone operator whose least
fixed point is computed soundly by tabled
evaluation~\cite{tamaki1986-tabled-resolution}.
\fi

\begin{definition}[\bilingual{Dependency graph}{依赖图}]\label{def:dependency-graph}
\ifInheritanceChinese
  固定一棵 AST 及其原始函数 $\defines$ 与 $\inherits$。
  \emph{依赖图}~$G$ 以由该 AST 产生的所有语义函数应用
  ($\properties(p)$、\allowbreak$\supers(p)$、\allowbreak$\overrides(p)$、\allowbreak$\bases(p)$、\allowbreak
  $\resolve(\ldots)$、\allowbreak$\this(\ldots)$) 为顶点。
  当通过方程~(\ref{eq:properties})--(\ref{eq:this}) 计算~$u$
  需要~$v$ 的值时,从顶点~$u$ 到顶点~$v$ 有一条有向边。
\else
  Fix an AST with its primitive functions $\defines$ and
  $\inherits$.
  The \emph{dependency graph}~$G$ has as vertices all
  semantic-function applications
  ($\properties(p)$, $\supers(p)$, $\overrides(p)$,
    $\bases(p)$, $\resolve(\ldots)$,
  $\this(\ldots)$) that arise from the AST\@.
  There is a directed edge from vertex~$u$ to vertex~$v$
  whenever computing~$u$ via
  equations~(\ref{eq:properties})--(\ref{eq:this})
  requires the value of~$v$.
\fi
\end{definition}

\begin{definition}[\bilingual{Immediate consequence operator $T_P$}{直接后果算子 $T_P$}]%
  \label{def:tp-operator}
\ifInheritanceChinese
  设 $\mathcal{Q}$ 为所有语义函数查询(即~$G$ 的顶点)的集合。
  对每个查询 $q \in \mathcal{Q}$,设 $\mathcal{V}_q$ 为形状与~$q$ 相匹配的答案值的集合
  (依定义~$q$ 的方程,为路径、路径的集合或路径对的集合)。
  \emph{查询--答案映射}是逐点乘积
  $\prod_{q \in \mathcal{Q}} \mathcal{P}(\mathcal{V}_q)$ 中的一个元素;
  等价地,它为每个查询~$q$ 指派一个候选答案集合 $M(q) \subseteq \mathcal{V}_q$。
  \emph{直接后果算子}~\cite{vanemden1976-predicate-logic-semantics}
  $T_P$ 将查询--答案映射~$M$ 映射为新的查询--答案映射~$T_P(M)$,方式是:对每个查询~$q$,
  求值定义~$q$ 的那条方程的右侧,并以当前近似~$M(q')$ 解释每个子查询~$q'$。
\else
  Let $\mathcal{Q}$ be the set of all
  semantic-function queries
  (vertices of~$G$).
  For each query $q \in \mathcal{Q}$, let
  $\mathcal{V}_q$ be the set of answer values of the
  appropriate shape for~$q$
  (paths, sets of paths, or sets of pairs of paths,
  depending on the equation defining~$q$).
  A \emph{query--answer map} is an element of the
  pointwise product
  $\prod_{q \in \mathcal{Q}} \mathcal{P}(\mathcal{V}_q)$;
  equivalently, it assigns to each query~$q$ a set
  $M(q) \subseteq \mathcal{V}_q$ of candidate answers.
  The \emph{immediate consequence
  operator}~\cite{vanemden1976-predicate-logic-semantics}
  $T_P$ maps a query--answer map~$M$ to a new
  query--answer map~$T_P(M)$ by evaluating, for each
  query~$q$, the right-hand side of the equation that
  defines~$q$, interpreting every sub-query~$q'$ by
  the current approximation~$M(q')$.
\fi
\end{definition}

\begin{lemma}[\bilingual{Monotonicity}{单调性}]\label{lem:monotonicity}
  \ifInheritanceChinese
  算子 $T_P$ 在完备格
  $\Bigl(\prod_{q \in \mathcal{Q}} \mathcal{P}(\mathcal{V}_q),\; \subseteq\Bigr)$
  上是单调的。
  \else
  The operator $T_P$ is monotone on the complete lattice
  $\Bigl(\prod_{q \in \mathcal{Q}} \mathcal{P}(\mathcal{V}_q),\; \subseteq\Bigr)$.
  \fi
\end{lemma}

\begin{proof}
\ifInheritanceChinese
  每条方程的右侧均由 $\cup$、$\in$、等式测试以及正集合
  概括 $\{\, x \mid x \in S, \ldots \,\}$ 构成。
  六条方程均不使用否定、集合差或补集。
  因此,若 $M \subseteq M'$(逐点),则 $T_P(M) \subseteq T_P(M')$。
\else
  Each equation's right-hand side is built from
  $\cup$, $\in$, equality tests, and positive set
  comprehension $\{\, x \mid x \in S, \ldots \,\}$.
  None of the six equations uses negation, set
  difference, or complement.
  Therefore, if $M \subseteq M'$ (pointwise), then
  $T_P(M) \subseteq T_P(M')$.
\fi
\end{proof}

\begin{theorem}[\bilingual{Existence and uniqueness of
  $\mathrm{lfp}(T_P)$}{$\mathrm{lfp}(T_P)$ 的存在性与唯一性}]\label{thm:lfp-exists}
  \ifInheritanceChinese
  算子 $T_P$ 有唯一的最小不动点 $\mathrm{lfp}(T_P)$。
  \else
  The operator $T_P$ has a unique least fixed point
  $\mathrm{lfp}(T_P)$.
  \fi
\end{theorem}

\begin{proof}
\ifInheritanceChinese
  每个幂集格
  $(\mathcal{P}(\mathcal{V}_q), \subseteq)$
  都是完备的,
  因此逐点乘积格
  $\Bigl(\prod_{q \in \mathcal{Q}} \mathcal{P}(\mathcal{V}_q), \subseteq\Bigr)$
  也是完备的。
  由 Knaster--Tarski 定理~\cite{tarski1955-lattice-fixpoint-theorem},
  完备格上的每个单调函数都有唯一的最小不动点。
\else
  Each powerset lattice
  $(\mathcal{P}(\mathcal{V}_q), \subseteq)$
  is complete,
  so the pointwise product lattice
  $\Bigl(\prod_{q \in \mathcal{Q}} \mathcal{P}(\mathcal{V}_q), \subseteq\Bigr)$
  is also complete.
  By the Knaster--Tarski
  theorem~\cite{tarski1955-lattice-fixpoint-theorem},
  every monotone function on a complete lattice has a
  unique least fixed point.
\fi
\end{proof}

\begin{theorem}[\bilingual{Well-definedness}{良定义性}]\label{thm:well-defined}
  \ifInheritanceChinese
  语义函数 $\properties$、\allowbreak$\supers$、\allowbreak$\overrides$、\allowbreak$\bases$、\allowbreak
  $\resolve$、\allowbreak$\this$ 是良定义的:
  其指称语义为 $\mathrm{lfp}(T_P)$,即直接后果算子的唯一最小不动点。
  \else
  The semantic functions
  $\properties$, $\supers$, $\overrides$, $\bases$,
  $\resolve$, $\this$ are well-defined:
  their denotational semantics is
  $\mathrm{lfp}(T_P)$, the unique least fixed point of
  the immediate consequence operator.
  \fi
\end{theorem}

\begin{theorem}[\bilingual{Tabling soundness}{Tabling 可靠性}]%
  \label{thm:tabling-soundness}
  \ifInheritanceChinese
  即使依赖图~$G$ 是无穷的,方程~(\ref{eq:properties})--(\ref{eq:this}) 的
  tabling 求值也绝不会返回错误答案。具体而言:
  \begin{enumerate}
    \item \textbf{可靠性(无假正例)。}
      tabling 返回的每个事实均由某条方程的右侧在其输入上求值产生。
      对求值轨迹进行归纳,每个返回的事实都属于 $\mathrm{lfp}(T_P)$。

    \item \textbf{单调积累。}
      tabling 累加器在各轮迭代之间只通过集合并增长,任何事实都不会被删除。

    \item \textbf{不动点检测的正确性。}
      若求值过程中没有子查询被重入,则所有子查询均在无近似的情况下被完整计算,结果是精确的。
      若发生了重入但累加器是稳定的(即没有新事实被添加),则累加器等于
      $\mathrm{lfp}(T_P)$ 在可达查询上的限制,因为它是一个从~$\bot$ 上升到达的前不动点。

    \item \textbf{可靠的失败。}
      当 tabling 无法确定答案时(无论是因为迭代预算耗尽还是因为存在无穷非循环链),
      它会抛出错误而非返回错误答案。
  \end{enumerate}
  \else
  Tabled evaluation of
  equations~(\ref{eq:properties})--(\ref{eq:this})
  never returns a wrong answer, even when the
  dependency graph~$G$ is infinite.
  Specifically:
  \begin{enumerate}
    \item \textbf{Soundness (no false positives).}
      Every fact returned by tabling was produced by
      evaluating some equation's right-hand side on its
      inputs.
      By induction on the evaluation trace, every
      returned fact belongs to $\mathrm{lfp}(T_P)$.

    \item \textbf{Monotone accumulation.}
      The tabling accumulator only grows between
      iterations via set union; no fact is ever
      removed.

    \item \textbf{Correct fixpoint detection.}
      If no sub-query was re-entered during evaluation,
      all sub-queries were fully computed with no
      approximation, and the result is exact.
      If re-entry occurred but the accumulator is
      stable, meaning no new facts were added, the
      accumulator equals
      $\mathrm{lfp}(T_P)$ restricted to the reachable
      queries, because it is a pre-fixed point reached
      by ascending from~$\bot$.

    \item \textbf{Sound failure.}
      When tabling cannot determine the answer,
      whether because the iteration budget is exhausted
      or because of an infinite non-cyclic chain,
      it raises an error rather than
      returning a wrong answer.
  \end{enumerate}
  \fi
\end{theorem}

\begin{proof}
\ifInheritanceChinese
  对于~(1),每个返回的事实都是某条方程的右侧作用于先前返回的事实上的输出;
  归纳地,所有输入均属于 $\mathrm{lfp}(T_P)$,因此输出也是(由单调性和不动点性质)。

  对于~(2),实现通过 $A \leftarrow A \cup T_P(A)$ 进行积累,这对~$A$ 是单调的。

  对于~(3),当没有重入时,求值是有限无环可达子图上的普通记忆化递归,
  因此每个子查询都被精确计算,不涉及任何近似。
  当发生重入时,设 $R \subseteq \mathcal{Q}$ 为从初始查询可达的查询集合,
  设 $T_P^R$ 表示 $T_P$ 到~$R$ 上映射的限制。
  若累加器稳定,即 $A = A \cup T_P^R(A)$,则
  $T_P^R(A) \subseteq A$,从而~$A$ 是~$T_P^R$ 的一个前不动点。

  由~(1),曾被添加到累加器中的每个事实均由某次右侧求值产生,
  因此属于 $\mathrm{lfp}(T_P^R)$。
  故 $A \subseteq \mathrm{lfp}(T_P^R)$。

  反之,由 Knaster--Tarski 将最小不动点刻画为最小前不动点,
  $\mathrm{lfp}(T_P^R)$ 包含于~$T_P^R$ 的每个前不动点中,特别地包含于~$A$ 中。
  从而 $\mathrm{lfp}(T_P^R) \subseteq A$。

  因此 $A = \mathrm{lfp}(T_P^R)$,即稳定的累加器等于可达查询上的精确指称。

  对于~(4),预算耗尽与栈溢出均产生异常,而非静默的错误答案。
\else
  For~(1), each returned fact is the output of some
  equation's right-hand side applied to previously
  returned facts; by induction, all inputs are in
  $\mathrm{lfp}(T_P)$, so the output is too (by
  monotonicity and the fixed-point property).

  For~(2), the implementation accumulates via
  $A \leftarrow A \cup T_P(A)$, which is monotone in~$A$.

  For~(3), when no reentry occurs, the evaluation is
  ordinary memoised recursion on a finite acyclic
  reachable subgraph, so every sub-query is computed
  exactly and no approximation is involved.
  When reentry occurs, let $R \subseteq \mathcal{Q}$ be
  the set of queries reachable from the initial query,
  and let $T_P^R$ denote the restriction of $T_P$ to
  maps on~$R$.
  If the accumulator is stable, that is,
  $A = A \cup T_P^R(A)$, then
  $T_P^R(A) \subseteq A$, so $A$ is a pre-fixed point
  of~$T_P^R$.

  By~(1), every fact ever added to the accumulator was
  produced by some right-hand side evaluation and
  therefore belongs to $\mathrm{lfp}(T_P^R)$.
  Hence $A \subseteq \mathrm{lfp}(T_P^R)$.

  Conversely, by the Knaster--Tarski characterization
  of the least fixed point as the least pre-fixed
  point, $\mathrm{lfp}(T_P^R)$ is contained in every
  pre-fixed point of~$T_P^R$, in particular in~$A$.
  Thus $\mathrm{lfp}(T_P^R) \subseteq A$.

  Therefore $A = \mathrm{lfp}(T_P^R)$, that is, the stable
  accumulator equals the exact denotation on the
  reachable queries.

  For~(4), budget exhaustion and stack overflow both
  produce exceptions, never silent wrong answers.
\fi
\end{proof}

% 终止性依赖程序结构:无环有限子图→一趟即止;有环有限逐查询值域→有限格上升链条件收敛;
% 无穷可达子图→固有发散,属图灵完备性,非求值策略的缺陷。
\paragraph{\bilingual{Termination}{终止性}}
\ifInheritanceChinese
求值是否终止取决于程序:
\begin{itemize}
  \item \textbf{无环、有限可达子图。}
    一趟即足;求值终止。
  \item \textbf{有环、逐查询值域有限。}
    多趟迭代在有限格的上升链条件下收敛~\cite{bancilhon1986-recursive-query-strategies}。
    这是 Datalog 的情形(第~\ref{sec:recursive-rules} 节)。
  \item \textbf{无穷可达子图。}
    求值可能发散;这是图灵完备性的固有属性,而非求值策略的缺陷。
\end{itemize}
\else
Whether evaluation terminates depends on the
program:
\begin{itemize}
  \item \textbf{Acyclic, finite reachable subgraph.}
    One pass suffices; evaluation terminates.
  \item \textbf{Cyclic, finite per-query domain.}
    Multiple passes converge by the ascending chain
    condition on a finite
    lattice~\cite{bancilhon1986-recursive-query-strategies}.
    This is the Datalog case
    (Section~\ref{sec:recursive-rules}).
  \item \textbf{Infinite reachable subgraph.}
    Evaluation may diverge; this is inherent to
    Turing completeness and not a defect of the
    evaluation strategy.
\end{itemize}
\fi

\section{B\"ohm Tree Correspondence: Proofs}
\label{app:bohm-tree-proofs}

This appendix proves two results about the sublanguage of
inheritance-calculus corresponding to the $\lambda$-calculus
(Section~\ref{sec:bohm-tree}).
First, the first-iteration approximation
$T_P\!\uparrow\!1$ is fully abstract with respect to
B\"ohm tree equivalence
(Theorems~\ref{thm:adequacy} and~\ref{thm:first-order-semantics}),
establishing that the first iteration coincides with
the lazy $\lambda$-calculus.
Second, the full $\mathrm{lfp}(T_P)$ semantics is strictly
more defined in the powerset sense of
Figure~\ref{fig:calculi-matrix}: translated $\lambda$-terms
that diverge under $\beta$-reduction may converge to definite
sets, as shown in Section~\ref{sec:bohm-tree}.

\subsection{B\"ohm Trees}

We briefly recall the definition of B\"ohm
trees~\cite{barendregt1984-lambda-calculus}.
A \emph{head reduction} $M \to_h M'$ contracts
the outermost $\beta$-redex only: if $M$ has the form
$(\lambda x.\, e)\; v\; M_1 \cdots M_k$, then
$M \to_h e[v/x]\; M_1 \cdots M_k$.
A term $M$ is in \emph{head normal form} (HNF) if it has no
head redex, that is, $M = \lambda x_1 \ldots x_m.\;
y\; M_1 \cdots M_k$ where $y$ is a variable.

The \emph{B\"ohm tree} $\mathrm{BT}(M)$ of a $\lambda$-term $M$
is an infinite labeled tree defined by:
\[
  \mathrm{BT}(M) =
  \begin{cases}
    \bot & \text{if } M \text{ has no HNF} \\[6pt]
    \lambda x_1 \ldots x_m.\;
    y\bigl(\mathrm{BT}(M_1),\; \ldots,\; \mathrm{BT}(M_k)\bigr)
    &
    \begin{aligned}[t]
      &\text{if } M \to^*_h \\[-2pt]
      &\lambda x_1 \ldots x_m.\; y\; M_1 \cdots M_k
    \end{aligned}
  \end{cases}
\]
B\"ohm tree equivalence $\mathrm{BT}(M) = \mathrm{BT}(N)$ is the
standard observational equivalence of the lazy
$\lambda$-calculus~\cite{abramsky1993-full-abstraction-lazy-lambda,abramsky1995-games-full-abstraction-lazy-lambda}.

\subsection{Path Encoding}

We define a correspondence between positions in the B\"ohm tree
and paths in the mixin tree.
When the translation $\mathcal{T}$ is applied to an ANF
$\lambda$-term, the resulting mixin tree has a specific shape:
\begin{itemize}
  \item An abstraction $\lambda x.\, M$ translates to a record
    with own properties $\{\mathrm{argument}, \mathrm{result}\}$,
    which is referred to as the \emph{abstraction shape}.
  \item A $\mathbf{let}$-binding
    $\mathbf{let}\; x = V_1\; V_2 \;\mathbf{in}\; M$ translates
    to a record with own properties $\{x, \mathrm{result}\}$,
    where $x$ holds the encapsulated application
    $\{\mathcal{T}(V_1),\; \mathrm{argument} \mapsto
    \mathcal{T}(V_2)\}$.
  \item A tail call $V_1\; V_2$ translates to a record with own
    properties $\{\mathrm{tailCall}, \mathrm{result}\}$, where
    $\mathrm{tailCall}$ holds the encapsulated application and
    $\mathrm{result}$ projects
    $\mathrm{tailCall}.\mathrm{result}$.
  \item A variable $x$ with de~Bruijn index $n$ translates to
    an indexed reference $\dbi{n}.\mathrm{argument}$.
\end{itemize}

\noindent
Each B\"ohm tree position is a sequence of navigation steps from
the root.
In the B\"ohm tree, the possible steps at a node
$\lambda x_1 \ldots x_m.\; y\; M_1 \cdots M_k$ are:
entering the body by peeling off one abstraction layer, and
entering the $i$-th argument $M_i$ of the head variable.
In the mixin tree, these correspond to following the
$\mathrm{result}$ label (for entering the body) and the
$\mathrm{argument}$ label after $\this$ resolution (for
entering an argument position).

Rather than comparing absolute paths across different mixin
trees, which would be sensitive to internal structure, we
observe only the \emph{convergence} behavior accessible via
the $\mathrm{result}$ projection
(Definition~\ref{def:convergence}).
Each $\mathrm{result}$-following step corresponds to one head
reduction step in the $\lambda$-calculus: a tail call's
$\mathrm{result}$ projects through the encapsulated
application, and a $\mathbf{let}$-binding's $\mathrm{result}$
enters the continuation.
The abstraction shape at the end signals that a head
normal form has been reached.

\subsection{Single-Path Lemma}

The $\this$ function (equation~\ref{eq:this}) is designed for the
general case of multi-target $\this$ resolution.
For translated $\lambda$-terms, we show it degenerates to
single-target resolution.

\begin{lemma}[Single-Path]\label{lem:single-path}
  For any closed ANF $\lambda$-term $M$, at every invocation of
  $\this(S,\; p_{\mathrm{def}},\; n)$ during evaluation of
  $\mathcal{T}(M)$, the frontier set $S$ contains exactly one
  path.
\end{lemma}

\begin{proof}
  By case analysis on the translation rules of~$\mathcal{T}$.
  Since $\mathcal{T}$ is compositional, the structural invariant
  established for each rule propagates to all inherited subtrees:
  every subtree appearing in the mixin tree is itself the image of
  some sub-term under~$\mathcal{T}$, so the invariant holds
  throughout the runtime composition.

  \paragraph{Case $\mathcal{T}(\lambda x.\, M)
    = \{\mathrm{argument} \mapsto \{\},\;
  \mathrm{result} \mapsto \mathcal{T}(M)\}$}
  This is a record literal with
  $\defines = \{\mathrm{argument}, \mathrm{result}\}$ and
  no inheritance sources, so $\inherits = \varnothing$.
  At the root path $p$, we have $\overrides(p) = \{p\}$ (a single
  path) since there are no inheritance sources to introduce
  additional branches.
  Inside $\mathcal{T}(M)$, any reference to $x$ is
  $\dbi{n}.\mathrm{argument}$ where $n$ is the de~Bruijn index
  pointing to this scope level.
  Since $\mathcal{T}(\lambda x.\, M)$ introduces exactly one
  scope level, there is exactly one
  mixin at that level, hence $\this$ finds exactly one matching
  pair $(p_{\mathrm{site}},\; p_{\mathrm{override}})$ in $\supers$.

  \paragraph{Case
    $\mathcal{T}(\mathbf{let}\; x = V_1\; V_2
    \;\mathbf{in}\; M)
    = \{x \mapsto \{\mathcal{T}(V_1),\;
      \mathrm{argument} \mapsto \mathcal{T}(V_2)\},\;
  \mathrm{result} \mapsto \mathcal{T}(M)\}$}
  The outer record has own properties $\{x, \mathrm{result}\}$
  and no inheritance sources at the outer level
  ($\inherits = \varnothing$), so $\overrides(\text{root})
  = \{\text{root}\}$.
  Inside the $x$ subtree, the inheritance
  $\{\mathcal{T}(V_1),\; \mathrm{argument} \mapsto
  \mathcal{T}(V_2)\}$ has exactly one inheritance source
  ($\mathcal{T}(V_1)$), creating exactly one $\bases$ entry.
  By induction on $V_1$, the single-path property holds inside
  $\mathcal{T}(V_1)$.
  Inside $\mathcal{T}(M)$, the induction hypothesis applies
  directly.

  \paragraph{Case
    $\mathcal{T}(\mathbf{let}\; x = V
    \;\mathbf{in}\; M)
    = \{x \mapsto \{\mathrm{result} \mapsto \mathcal{T}(V)\},\;
  \mathrm{result} \mapsto \mathcal{T}(M)\}$}
  The outer record has own properties $\{x, \mathrm{result}\}$
  and no inheritance sources at the outer level, so
  $\overrides(\text{root}) = \{\text{root}\}$.
  The $x$ subtree is itself a record literal with own
  property $\{\mathrm{result}\}$ and no inheritance
  sources.
  Thus neither entering the bound value nor entering the
  body introduces branching in $\supers$.
  By induction, the single-path property holds inside both
  $\mathcal{T}(V)$ and $\mathcal{T}(M)$.

  \paragraph{Case $\mathcal{T}(V_1\; V_2) =
    \{\mathrm{tailCall} \mapsto
      \{\mathcal{T}(V_1),\; \mathrm{argument} \mapsto
      \mathcal{T}(V_2)\},\;
      \mathrm{result} \mapsto
  \{\mathrm{tailCall}.\mathrm{result}\}\}$}
  Identical to the $\mathbf{let}$-binding case: the outer record
  has own properties $\{\mathrm{tailCall}, \mathrm{result}\}$
  with no inheritance sources at the outer level, and the
  $\mathrm{tailCall}$ subtree contains exactly one inheritance
  source.

  \paragraph{Case $\mathcal{T}(x)
  = \dbi{n}.\mathrm{argument}$}
  This is a reference, not a record.
  The $\this$ function walks up $n$ steps from the enclosing scope.
  By the compositionality argument above, every scope level
  encountered during the walk, including inherited subtrees
  reached through reference resolution, is the image of some
  sub-term under~$\mathcal{T}$, so each step of $\this$
  encounters exactly one matching pair in $\supers$.

  \medskip\noindent
  The key structural invariant is the following.
  For every subtree generated by~$\mathcal{T}$,
  its root has either no inheritance source or exactly one
  inheritance source.
  Moreover, the only roots with an inheritance source are the
  application wrappers introduced by the application
  $\mathbf{let}$-binding and tail-call rules, and in both cases
  that source is unique and encapsulated under a distinguished
  child ($x$ or $\mathrm{tailCall}$), never at the enclosing
  scope level.
  Therefore every scope level traversed by $\this$ contributes
  at most one matching pair in $\supers$, so the frontier set
  remains a singleton throughout evaluation.
\end{proof}

\subsection{Substitution Lemma}

The core mechanism of the translation is that inheritance with
$\mathrm{argument} \mapsto \mathcal{T}(V)$ plays the role of
substitution.  We make this precise.

\begin{lemma}[Substitution]\label{lem:substitution}
  Let $M$ be an ANF term in which $x$ may occur free,
  and let $V$ be a closed value.%
  \footnote{
    The proof uses the Single-Path Lemma
    (Lemma~\ref{lem:single-path}), which requires the inherited tree
    to occur within a closed term. This holds whenever
    $\lambda x.\, M$ applied to $V$ is a subterm of a closed ANF
    program.
  }
  Define the \emph{inherited tree}
  $C = \{\mathcal{T}(\lambda x.\, M),\;
  \mathrm{argument} \mapsto \mathcal{T}(V)\}$.
  Then for every path $p$ under $C \snoc \mathrm{result}$ and
  every label $\ell$:
  \[
    \ell \in \properties(p)
    \text{ in } C \snoc \mathrm{result}
    \quad\Longleftrightarrow\quad
    \ell \in \properties(p)
    \text{ in } \mathcal{T}(M[V/x])
  \]
  where $M[V/x]$ is the usual capture-avoiding substitution.
\end{lemma}

\begin{proof}
  By structural induction on $M$.

  \paragraph{Base case: $M = x$}
  Then $\mathcal{T}(x) =
  \dbi{n}.\mathrm{argument}$ (where $n$ is the de~Bruijn index
  of $x$), and
  $C \snoc \mathrm{result}$ contains this reference.
  In the inherited tree $C$, $\resolve$ and $\this$ resolve the reference by
  navigating from the reference's enclosing scope up $n$ steps
  to the binding $\lambda$, where the inheritance
  $\{\mathcal{T}(\lambda x.\, x),\;
  \mathrm{argument} \mapsto \mathcal{T}(V)\}$ inherits into the
  $\mathrm{argument}$ slot, providing $\mathcal{T}(V)$.
  Since $C$ is the $\mathrm{tailCall}$ subtree of
  $\mathcal{T}\bigl((\lambda x.\, M)\; V\bigr)$, which is a
  closed ANF term, Lemma~\ref{lem:single-path} applies and
  $\this$ finds exactly one path,
  so $\resolve$ returns $\mathcal{T}(V)$'s subtree.
  Since $M[V/x] = V$, we have
  $\mathcal{T}(M[V/x]) = \mathcal{T}(V)$, and the properties
  coincide.

  \paragraph{Base case: $M = y$ where $y \neq x$ and $y$ is
  $\lambda$-bound}
  Then $\mathcal{T}(y) = \dbi{m}.\mathrm{argument}$ where $m$
  is the de~Bruijn index of~$y$.
  This walks up to the scope that binds~$y$, which is different
  from the scope that binds~$x$, so the inheritance with
  $\mathrm{argument} \mapsto \mathcal{T}(V)$ at $x$'s binding
  scope has no effect.
  Since $M[V/x] = y$, the properties are identical.

  \paragraph{Base case: $M = y$ where $y \neq x$ and $y$ is
  let-bound}
  Then $\mathcal{T}(y) = y.\mathrm{result}$.
  This is a lexical reference to the sibling property~$y$ in the
  enclosing record, so the inherited tree~$C$ does not alter its
  resolution.
  Since $M[V/x] = y$, the properties are identical.

  \paragraph{Inductive case: $M = \lambda y.\, M'$}
  Then $\mathcal{T}(M) =
  \{\mathrm{argument} \mapsto \{\},\;
  \mathrm{result} \mapsto \mathcal{T}(M')\}$.
  The subtree at $C \snoc \mathrm{result}$ is another
  abstraction-shaped record.
  References to $x$ inside $\mathcal{T}(M')$ resolve through
  $\this$ in exactly the same way (with one additional scope
  level from the inner $\lambda$), and by induction on $M'$, the
  properties under
  $C \snoc \mathrm{result} \snoc \mathrm{result}$
  coincide with those of
  $\mathcal{T}(M'[V/x])$.
  Since
  $(\lambda y.\, M')[V/x] = \lambda y.\, (M'[V/x])$
  (assuming $y$ is fresh), the result follows.

  \paragraph{Inductive case:
    $M = \mathbf{let}\; z = V_1\; V_2$
  $\mathbf{in}\; M'$}
  Then
  \[
    \mathcal{T}(M) =
    \{z \mapsto \{\mathcal{T}(V_1),\;
      \mathrm{argument} \mapsto \mathcal{T}(V_2)\},\;
    \mathrm{result} \mapsto \mathcal{T}(M')\}.
  \]
  The substitution distributes:
  $M[V/x] =
  \mathbf{let}\; z = V_1[V/x]\; V_2[V/x]
  \;\mathbf{in}\; M'[V/x]$.
  By induction on $V_1$, $V_2$, and $M'$, the properties under
  each subtree ($z$ and $\mathrm{result}$) coincide,
  since references to $x$ inside each subtree resolve to
  $\mathcal{T}(V)$ via the same $\this$ mechanism.

  \paragraph{Inductive case:
    $M = \mathbf{let}\; z = V'
  \;\mathbf{in}\; M'$}
  Then
  \[
    \mathcal{T}(M) =
    \{z \mapsto \{\mathrm{result} \mapsto \mathcal{T}(V')\},\;
    \mathrm{result} \mapsto \mathcal{T}(M')\}.
  \]
  The substitution distributes:
  $M[V/x] =
  \mathbf{let}\; z = V'[V/x]
  \;\mathbf{in}\; M'[V/x]$.
  By induction on $V'$ and $M'$, the properties under both
  subtrees, specifically $z$ and $\mathrm{result}$, coincide.
  In particular, occurrences of~$z$ in $M'$ are translated as the
  lexical reference $z.\mathrm{result}$ on both sides, so the
  wrapper introduced by the value $\mathbf{let}$-binding is
  preserved by the substitution.

  \paragraph{Inductive case: $M = V_1\; V_2$ (tail call)}
  Analogous to the $\mathbf{let}$-binding case with
  $\mathrm{tailCall}$ in place of $z$.
\end{proof}

\subsection{Convergence Preservation}

The convergence criterion (Definition~\ref{def:convergence})
follows the $\mathrm{result}$ chain from the root.
Each step in this chain corresponds to one head reduction step
in the $\lambda$-calculus.
We make this precise.

\begin{lemma}[Result-step]\label{lem:result-step}
  Let $(\lambda x.\, M)\; V$ be a tail call, which is a $\beta$-redex.
  Then for every path $p$ and label $\ell$:
  \[
    \ell \in \properties\!\bigl(\,
    \text{root} \snoc \mathrm{result} \snoc p\,\bigr)
    \text{ in } \mathcal{T}\bigl((\lambda x.\, M)\; V\bigr)
    \;\Longleftrightarrow\;
    \ell \in \properties\!\bigl(\,
    \text{root} \snoc p\,\bigr)
    \text{ in } \mathcal{T}(M[V/x])
  \]
  That is, following one $\mathrm{result}$ projection in the
  tail call's mixin tree yields the same properties as the
  root of the reduct's mixin tree.
\end{lemma}

\begin{proof}
  The translation gives:
  \[
    \mathcal{T}\bigl((\lambda x.\, M)\; V\bigr) =
    \{\mathrm{tailCall} \mapsto
      \{\mathcal{T}(\lambda x.\, M),\;
      \mathrm{argument} \mapsto \mathcal{T}(V)\},\;
      \mathrm{result} \mapsto
    \{\mathrm{tailCall}.\mathrm{result}\}\}
  \]
  The $\mathrm{result}$ label at the root is defined as
  $\mathrm{tailCall}.\mathrm{result}$, so
  $\text{root} \snoc \mathrm{result}$ resolves to the
  $\mathrm{result}$ path inside the inheritance
  $\mathrm{tailCall} = \{\mathcal{T}(\lambda x.\, M),\;
  \mathrm{argument} \mapsto \mathcal{T}(V)\}$.
  This inheritance is exactly the inherited tree $C$ of
  Lemma~\ref{lem:substitution}, and
  $\mathrm{tailCall} \snoc \mathrm{result}$ is
  $C \snoc \mathrm{result}$.
  By Lemma~\ref{lem:substitution}, the properties under
  $C \snoc \mathrm{result}$ coincide with those of
  $\mathcal{T}(M[V/x])$.

  \medskip\noindent
  \emph{Let-binding variant.}
  An analogous result holds for
  $\mathbf{let}\; x = V_1\; V_2 \;\mathbf{in}\; M'$
  where $V_1 = \lambda y.\, M''$.
  The translation gives:
  \[
    \mathcal{T}\bigl(\mathbf{let}\; x = V_1\; V_2
    \;\mathbf{in}\; M'\bigr) =
    \{x \mapsto \{\mathcal{T}(V_1),\;
      \mathrm{argument} \mapsto \mathcal{T}(V_2)\},\;
    \mathrm{result} \mapsto \mathcal{T}(M')\}
  \]
  Since $\mathrm{result} \mapsto \mathcal{T}(M')$ is a property
  definition, the subtree at
  $\text{root} \snoc \mathrm{result}$ is $\mathcal{T}(M')$.
  Inside $\mathcal{T}(M')$, each reference to $x$ is
  $x.\mathrm{result}$, which resolves to the $\mathrm{result}$
  path inside the $x$ subtree
  $\{\mathcal{T}(\lambda y.\, M''),\;
  \mathrm{argument} \mapsto \mathcal{T}(V_2)\}$.
  This is the inherited tree $C$ of
  Lemma~\ref{lem:substitution}, so
  $x.\mathrm{result}$ has the same properties as
  $\mathcal{T}(M''[V_2/y])$.
  In the $\lambda$-calculus,
  $M = (\lambda x.\, M')(V_1\; V_2)$, and the head reduct is
  $N = M'[(V_1\; V_2)/x]$.
  In $\mathcal{T}(N)$, each occurrence of $x$ is replaced by
  $(V_1\; V_2)$, whose result, by the tail-call variant above,
  also yields $\mathcal{T}(M''[V_2/y])$.
  Therefore the properties at
  $\text{root} \snoc \mathrm{result} \snoc p$ in
  $\mathcal{T}(M)$ coincide with those at
  $\text{root} \snoc p$ in $\mathcal{T}(N)$ for all $p$.

  \medskip\noindent
  \emph{Value let-binding variant.}
  An analogous result holds for
  $\mathbf{let}\; x = V \;\mathbf{in}\; M'$.
  The translation gives:
  \[
    \mathcal{T}\bigl(\mathbf{let}\; x = V \;\mathbf{in}\; M'\bigr) =
    \{x \mapsto \{\mathrm{result} \mapsto \mathcal{T}(V)\},\;
    \mathrm{result} \mapsto \mathcal{T}(M')\}.
  \]
  Since $\mathrm{result} \mapsto \mathcal{T}(M')$ is a property
  definition, the subtree at
  $\text{root} \snoc \mathrm{result}$ is $\mathcal{T}(M')$.
  Inside $\mathcal{T}(M')$, each occurrence of $x$ is translated
  as the lexical reference $x.\mathrm{result}$, which resolves to
  the subtree $\mathcal{T}(V)$.
  This is exactly the translation of the reduct
  $M'[V/x]$.
  Therefore the properties at
  $\text{root} \snoc \mathrm{result} \snoc p$ in
  $\mathcal{T}(\mathbf{let}\; x = V \;\mathbf{in}\; M')$
  coincide with those at
  $\text{root} \snoc p$ in $\mathcal{T}(M'[V/x])$ for all $p$.
\end{proof}

\begin{theorem}[Convergence preservation]%
  \label{thm:convergence-preservation}
  If $M \to_h N$ (head reduction), then
  $\mathcal{T}(M){\Downarrow}$ if and only if
  $\mathcal{T}(N){\Downarrow}$.
\end{theorem}

\begin{proof}
  A head reduction step in ANF takes one of three forms:
  a tail call $(\lambda x.\, M')\; V \to_h M'[V/x]$,
  an application $\mathbf{let}$-binding
  $\mathbf{let}\; x = V_1\; V_2 \;\mathbf{in}\; M'
  \to_h M'[(V_1\; V_2)/x]$,
  or a value $\mathbf{let}$-binding
  $\mathbf{let}\; x = V \;\mathbf{in}\; M'
  \to_h M'[V/x]$.
  The second and third forms are simply the ANF
  presentation of one $\beta$-step in the underlying
  $\lambda$-term.
  In all three cases, Lemma~\ref{lem:result-step}
  (tail-call, application let-binding, and value
  let-binding variants) shows that
  the properties at $\text{root} \snoc \mathrm{result}^n$ in
  $\mathcal{T}(M)$ coincide with the properties at
  $\text{root} \snoc \mathrm{result}^{n-1}$ in
  $\mathcal{T}(N)$ for all $n \ge 1$, where $N$ is the
  head reduct.
  Hence the abstraction shape appears at depth $n$ in the
  pre-reduct if and only if it appears at depth $n - 1$ in
  the post-reduct.
\end{proof}

\begin{remark}[Iteration levels]
  \label{rem:convergence-lfp}
  Section~\ref{sec:bohm-tree} showed that the B\"ohm tree
  correspondence holds at the first iteration
  $T_P\!\uparrow\!1$ and that $\mathrm{lfp}(T_P)$ is
  strictly more defined in the powerset sense of
  Figure~\ref{fig:calculi-matrix}.
  Here we make the connection precise.
  Tabled evaluation~\cite{tamaki1986-tabled-resolution}
  computes $\mathrm{lfp}(T_P)$ by iterating:
  $T_P\!\uparrow\!0 = \bot$,
  $T_P\!\uparrow\!1 = T_P(\bot)$,
  $T_P\!\uparrow\!2 = T_P(T_P(\bot))$,
  \ldots\ until stabilisation
  (Theorem~\ref{thm:tabling-soundness}).

  \begin{description}
    \item[$T_P\!\uparrow\!1$ corresponds to
      $\beta$-reduction.]
      $T_P\!\uparrow\!1$ evaluates each equation once,
      with re-entrant calls returning~$\bot$
      (the value from $T_P\!\uparrow\!0$).
      The Result-Step Lemma
      (Lemma~\ref{lem:result-step}) establishes a
      structural isomorphism between the equation
      systems of $\mathcal{T}(M)$ and
      $\mathcal{T}(N)$, shifted by one
      $\mathrm{result}$ step.
      This isomorphism is purely structural and holds
      for any number of iterations.
      Theorems~\ref{thm:adequacy}
      and~\ref{thm:first-order-semantics}
      are proved at this level.
  \end{description}

  \paragraph{Observation granularity}%
  \label{rem:set-vs-domain}
  The set-valued semantics at $T_P\!\uparrow\!1$ is
  strictly finer than the domain-theoretic semantics of
  the lazy $\lambda$-calculus.
  In the latter, a non-terminating computation yields a
  single undifferentiated~$\bot$;
  in the former, one can in principle distinguish
  $\varnothing$ (a definite empty set, arrived at in
  finite time) from non-stabilisation of the Kleene
  iteration (no finite approximation suffices).
  The convergence criterion
  (Definition~\ref{def:convergence}) deliberately
  projects away this distinction: it observes only
  whether the $\mathrm{result}$ chain reaches an abstraction
  shape, reducing the observation to a single bit.
  Theorems~\ref{thm:adequacy}
  and~\ref{thm:first-order-semantics}
  are stated at this coarsened observation level.

  \begin{description}
    \item[$\mathrm{lfp}(T_P)$ restores full precision.]
      At $\mathrm{lfp}(T_P)$, the full precision of the
      set-valued semantics is restored: re-entrant calls
      that returned~$\varnothing$ at $T_P\!\uparrow\!1$
      may now converge to non-empty sets, so the semantics
      distinguishes ``definitionally empty'' from
      ``convergent after further iteration,''
      and $M{\Downarrow}_\lambda$ implies
      $\mathcal{T}(M){\Downarrow}_{\mathrm{lfp}}$,
      but not conversely.
  \end{description}
\end{remark}

\subsection{Adequacy}

\begin{definition}[Inheritance-convergence]\label{def:convergence}
  A mixin tree $T$ \emph{converges}, written $T{\Downarrow}$,
  if there exists $n \ge 0$ such that:
  \[
    \{\mathrm{argument},\, \mathrm{result}\}
    \;\subseteq\;
    \properties\!\bigl(\,
      \underbrace{\text{root} \snoc \mathrm{result}
      \snoc \cdots \snoc \mathrm{result}}_{n}
    \,\bigr)
  \]
  in the semantics of
  Section~\ref{sec:mixin-trees}.
\end{definition}

\begin{theorem}[Adequacy]\label{thm:adequacy}
  A closed ANF $\lambda$-term $M$ has a head normal form if and
  only if $\mathcal{T}(M){\Downarrow}$.
\end{theorem}

A closed ANF term at the top level is either an abstraction
(a value) or a computation (an application
  $\mathbf{let}$-binding, a value $\mathbf{let}$-binding,
or a tail call).
An abstraction translates to a record whose root immediately
has the abstraction shape; a computation's root has the form
$\{x, \mathrm{result}\}$ or
$\{\mathrm{tailCall}, \mathrm{result}\}$, and one must follow
the $\mathrm{result}$ chain to find the eventual value.

\begin{proof}[Proof of Theorem~\ref{thm:adequacy}]
  ($\Rightarrow$)
  Suppose $M \to^*_h \lambda x.\, M'$ in $k$ head reduction
  steps.
  We show $\mathcal{T}(M)$ converges at depth $\le k$ by
  induction on $k$.

  \paragraph{Base case ($k = 0$)}
  $M$ is already an abstraction
  $\lambda x.\, M'$.
  Then $\mathcal{T}(M) =
  \{\mathrm{argument} \mapsto \{\},\;
  \mathrm{result} \mapsto \mathcal{T}(M')\}$,
  whose root has
  $\defines = \{\mathrm{argument}, \mathrm{result}\}$.
  These labels are directly in $\defines$ of the root, so
  the recursive evaluation returns them without further
  recursion, and $\mathcal{T}(M)$ converges at depth $n = 0$.

  \paragraph{Inductive step ($k \ge 1$)}
  $M$ is not an abstraction,
  so $M$ is either a tail call, an application
  $\mathbf{let}$-binding, or a value
  $\mathbf{let}$-binding.
  Each form corresponds to one head $\beta$-step in the
  underlying $\lambda$-term.
  Let $N$ be the head reduct, so
  $M \to_h N \to^{k-1}_h \lambda x.\, M'$.
  Note that $N$ is closed: head reduction substitutes a
  closed value into a term whose only free variable is the
  substitution target, preserving closedness.
  By Lemma~\ref{lem:result-step} (tail-call,
    application let-binding, or value let-binding
  variant), the properties at
  $\text{root} \snoc \mathrm{result} \snoc p$
  in $\mathcal{T}(M)$ coincide with those at
  $\text{root} \snoc p$ in $\mathcal{T}(N)$.
  By the induction hypothesis, $\mathcal{T}(N)$
  converges at depth $\le k - 1$, so $\mathcal{T}(M)$
  converges at depth $\le k$.

  \medskip\noindent
  ($\Leftarrow$)
  Suppose $\mathcal{T}(M){\Downarrow}$ at depth $n$.
  We show $M$ has a head normal form by induction on $n$.

  \paragraph{Base case ($n = 0$)}
  The root of $\mathcal{T}(M)$ has
  abstraction shape
  $\{\mathrm{argument}, \mathrm{result}\}$.
  Only the abstraction rule $\mathcal{T}(\lambda x.\, M')$
  produces a root with own property $\mathrm{argument}$.
  (For tail calls and $\mathbf{let}$-bindings, the root has
    $\defines = \{\mathrm{tailCall}, \mathrm{result}\}$
    or $\{x, \mathrm{result}\}$, and $\mathrm{argument}$ cannot
    appear via $\supers$ because
    $\inherits = \varnothing$ at the root, giving
  $\bases(\text{root}) = \varnothing$.)
  Therefore $M$ is an abstraction, hence already in head normal
  form.

  \paragraph{Inductive step ($n \ge 1$)}
  The abstraction shape
  appears at depth $n$ but not at depth $0$.
  By the base case argument, $M$ is not an abstraction, so at the
  head position of this closed top-level ANF term, $M$ is either
  a tail call, an application $\mathbf{let}$-binding, or a value
  $\mathbf{let}$-binding.
  In each case, the ANF grammar requires the function
  position to be a value~$V_1$
  ($V ::= x \mid \lambda x.\, M$).
  Since $M$ is closed, $V_1$ cannot be a free variable,
  so $V_1$ must be a $\lambda$-abstraction;
  the head form is therefore a $\beta$-redex.
  By Lemma~\ref{lem:result-step} (tail-call,
    application let-binding, or value let-binding
  variant), the properties at depth
  $n - 1$ in $\mathcal{T}(N)$ match those at depth $n$
  in $\mathcal{T}(M)$, where $N$ is the head reduct.
  So $\mathcal{T}(N)$ converges at depth $n - 1$.
  Since $M$ is closed and head reduction preserves
  closedness (it substitutes a closed value for the only
  free variable), $N$ is a closed ANF term.
  By the induction hypothesis, $N$ has a head normal
  form, and since $M \to_h N$, so does $M$.
\end{proof}

\subsection{First-Order Semantics of the Lazy \texorpdfstring{$\lambda$}{λ}-Calculus}

Adequacy relates a single term's head normal form to its
inheritance-convergence.
The following theorem lifts this to an equivalence between
terms, establishing that the translation $\mathcal{T}$
gives the lazy $\lambda$-calculus a first-order semantics
that preserves and reflects B\"ohm tree equivalence.

\begin{definition}[$\lambda$-contextual equivalence under $\mathcal{T}$]%
  \label{def:ctx-equiv}
  For closed $\lambda$-terms $M$ and $N$, write
  $\mathcal{T}(M) \approx_\lambda \mathcal{T}(N)$ if for every
  closing $\lambda$-calculus context $C[\cdot]$:
  \[
    \mathcal{T}(C[M]){\Downarrow}
    \;\Longleftrightarrow\;
    \mathcal{T}(C[N]){\Downarrow}
  \]
\end{definition}

\noindent
The subscript $\lambda$ emphasizes that the quantification ranges
over $\lambda$-calculus contexts, not over all inheritance-calculus contexts.
This definition observes only convergence\footnote{
  Convergence provides exactly a single bit of information.
}, but
quantifying over all $\lambda$-contexts makes it a fine-grained
equivalence: different contexts can supply different arguments
and project different results, probing every aspect of the
term's behavior within the $\lambda$-definable fragment.

\begin{theorem}[First-order semantics of the lazy $\lambda$-calculus]%
  \label{thm:first-order-semantics}
  For closed $\lambda$-terms $M$ and $N$:
  \[
    \mathcal{T}(M) \approx_\lambda \mathcal{T}(N)
    \quad\Longleftrightarrow\quad
    \mathrm{BT}(M) = \mathrm{BT}(N)
  \]
\end{theorem}

\begin{proof}[Proof of Theorem~\ref{thm:first-order-semantics}]
  The translation $\mathcal{T}$ is compositional: for every
  $\lambda$-calculus context $C[\cdot]$, the mixin tree
  $\mathcal{T}(C[M])$ is determined by the mixin tree
  $\mathcal{T}(M)$ and the translation of the context.
  This means $\mathcal{T}$ preserves the context structure
  needed for the argument.

  \medskip\noindent
  ($\Leftarrow$)
  Suppose $\mathrm{BT}(M) = \mathrm{BT}(N)$.
  B\"ohm tree equivalence is a congruence~\cite{barendregt1984-lambda-calculus},
  so for every context $C[\cdot]$,
  $\mathrm{BT}(C[M]) = \mathrm{BT}(C[N])$.
  In particular, $C[M]$ has a head normal form iff $C[N]$ does.
  By Theorem~\ref{thm:adequacy},
  $\mathcal{T}(C[M]){\Downarrow}$ iff
  $\mathcal{T}(C[N]){\Downarrow}$.
  Hence $\mathcal{T}(M) \approx_\lambda \mathcal{T}(N)$.

  \medskip\noindent
  ($\Rightarrow$)
  Suppose $\mathrm{BT}(M) \neq \mathrm{BT}(N)$.
  B\"ohm tree equivalence coincides with observational
  equivalence for the lazy
  $\lambda$-calculus~\cite{abramsky1993-full-abstraction-lazy-lambda,abramsky1995-games-full-abstraction-lazy-lambda}:
  there exists a context $C[\cdot]$ such that $C[M]$ converges
  and $C[N]$ diverges, or vice versa.
  By Theorem~\ref{thm:adequacy},
  $\mathcal{T}(C[M]){\Downarrow}$ and
  $\mathcal{T}(C[N]){\Uparrow}$.
  Hence $\mathcal{T}(M) \not\approx_\lambda \mathcal{T}(N)$.
\end{proof}

\noindent
The proof rests on two pillars: Adequacy
(Theorem~\ref{thm:adequacy}), which is our contribution, and
the classical result that B\"ohm tree equivalence equals
observational equivalence for the lazy
$\lambda$-calculus~\cite{abramsky1993-full-abstraction-lazy-lambda,abramsky1995-games-full-abstraction-lazy-lambda}.
Neither pillar requires defining or comparing
absolute paths inside mixin trees.

\section{\bilingual{Expressive Asymmetry: Proofs}{表达力的不对称:证明}}
\label{app:expressiveness}

\begin{definition}[\bilingual{Macro-expressibility, after Felleisen~\cite{felleisen1991-expressive-power}}{宏可表达性,据 Felleisen~\cite{felleisen1991-expressive-power}}]
  \label{def:macro-expressibility}
  \ifInheritanceChinese
  语言 $\mathscr{L}_1$ 的一个构造 $C$ 在语言 $\mathscr{L}_0$ 中是\emph{宏可表达的},若存在一个翻译 $\mathcal{E}$,使得对每一个含有 $C$ 的出现的程序 $P$,翻译 $\mathcal{E}(P)$ 是通过将 $C$ 的每次出现替换为一个仅依赖于 $C$ 的子表达式的 $\mathscr{L}_0$ 表达式而得到的,并保持 $P$ 的所有其他构造不变。
  \else
  A construct $C$ of language $\mathscr{L}_1$ is \emph{macro-expressible}
  in language $\mathscr{L}_0$ if there exists a translation $\mathcal{E}$
  such that for every program $P$ containing occurrences of $C$,
  the translation $\mathcal{E}(P)$ is obtained by replacing each
  occurrence of $C$ with an $\mathscr{L}_0$ expression that depends only
  on $C$'s subexpressions, leaving all other constructs of $P$
  unchanged.
  \fi
\end{definition}

\begin{lemma}[\bilingual{Composition of macro-expressible translations}{宏可表达翻译的复合}]
  \label{lem:macro-composition}
  \ifInheritanceChinese
  若 $\mathcal{E}_1$ 在 $\mathscr{L}_0$ 中宏可表达 $\mathscr{L}_1$ 的构造,且 $\mathcal{E}_2$ 在 $\mathscr{L}_1$ 中宏可表达 $\mathscr{L}_2$ 的构造,则复合 $\mathcal{E}_1 \circ \mathcal{E}_2$ 在 $\mathscr{L}_0$ 中宏可表达 $\mathscr{L}_2$ 的构造。
  \else
  If $\mathcal{E}_1$ macro-expresses the constructs of
  $\mathscr{L}_1$ in $\mathscr{L}_0$ and
  $\mathcal{E}_2$ macro-expresses the constructs of
  $\mathscr{L}_2$ in $\mathscr{L}_1$, then the composition
  $\mathcal{E}_1 \circ \mathcal{E}_2$ macro-expresses the
  constructs of $\mathscr{L}_2$ in $\mathscr{L}_0$.
  \fi
\end{lemma}

\begin{proof}
  \ifInheritanceChinese
  $\mathscr{L}_2$ 的每个构造 $F$ 在 $\mathscr{L}_1$ 上有一个固定的语法抽象 $A_F^{(2)}$。
  出现在 $A_F^{(2)}$ 中的每个 $\mathscr{L}_1$-构造 $G$ 在 $\mathscr{L}_0$ 上有一个固定的语法抽象 $A_G^{(1)}$。
  将 $A_F^{(2)}$ 中的每个 $G$ 替换为 $A_G^{(1)}$,得到 $F$ 在 $\mathscr{L}_0$ 上的一个固定语法抽象,从而满足 E4。
  \else
  Each construct~$F$ of~$\mathscr{L}_2$ has a fixed
  syntactic abstraction~$A_F^{(2)}$
  over~$\mathscr{L}_1$.
  Each $\mathscr{L}_1$-construct~$G$ appearing
  in~$A_F^{(2)}$ has a fixed
  syntactic abstraction~$A_G^{(1)}$
  over~$\mathscr{L}_0$.
  Substituting~$A_G^{(1)}$ for every~$G$
  in~$A_F^{(2)}$ yields a fixed syntactic
  abstraction over~$\mathscr{L}_0$ for~$F$,
  satisfying~E4.
  \fi
\end{proof}

\begin{definition}[\bilingual{Language universe}{语言全集}]
  \label{def:language-universe}
  \ifInheritanceChinese
  设 $\mathscr{U}$ 为一个语言,其构造子是继承演算的构造子、惰性 $\lambda$-演算的构造子(包括 $\lambda$-抽象、应用与变量)以及 $\mathbf{let}$-绑定的并集。
  $\mathscr{U}$ 的语义通过将每个 $\mathscr{U}$-程序翻译为继承演算来给出,分三步:
  \begin{enumerate}
    \item 将每个极大 $\lambda$-演算子表达式(即根节点与所有内部节点仅使用 $\lambda$-演算构造子的每棵子树)转化为 A-范式~\cite{flanagan1993-essence-compiling-continuations}。
    \item 将翻译 $\mathcal{T}$(第~\ref{sec:forward-translation} 节)应用于每个 ANF $\lambda$-演算子表达式。
    \item 按第~\ref{sec:mixin-trees} 节的规则对所得继承演算程序求值。
  \end{enumerate}
  当一个 $\mathscr{U}$-表达式交错使用 $\lambda$-演算与继承演算的构造子(例如,一个 $\mathbf{let}$-绑定,其被绑定的值是一个继承演算表达式)时,翻译从叶到根组合地应用:每个 $\lambda$-演算构造子通过 $\mathcal{T}$ 替换为其继承演算等价物,同时将继承演算子表达式视为不透明的值。
  由于 $\mathcal{T}$ 的每条规则(第~\ref{sec:forward-translation} 节)都是一个固定的语法抽象,仅依赖于被翻译的子表达式,因此这种自底向上的应用是良定义的,并为步骤~(3) 产生一个纯继承演算程序。

  若一个 $\mathscr{U}$-程序 $P$ 的翻译收敛,则 $P$ 收敛,记为 $\mathrm{eval}_{\mathscr{U}}(P)$。
  \else
  Let $\mathscr{U}$ be the language whose constructors are
  the union of the constructors of inheritance-calculus,
  the constructors of the lazy $\lambda$-calculus
  , including $\lambda$-abstraction, application, and variable, and
  $\mathbf{let}$-binding.
  The semantics of $\mathscr{U}$ is given by translating
  every $\mathscr{U}$-program to inheritance-calculus in
  three steps:
  \begin{enumerate}
    \item Convert every maximal $\lambda$-calculus
      subexpression (that is, every subtree whose root
        and all internal nodes use only $\lambda$-calculus
      constructors) to A-normal
      form~\cite{flanagan1993-essence-compiling-continuations}.
    \item Apply the translation $\mathcal{T}$
      (Section~\ref{sec:forward-translation}) to every ANF
      $\lambda$-calculus subexpression.
    \item Evaluate the resulting inheritance-calculus
      program by the rules of
      Section~\ref{sec:mixin-trees}.
  \end{enumerate}
  When a $\mathscr{U}$-expression interleaves
  $\lambda$-calculus and inheritance-calculus constructors
  (e.g.\ a $\mathbf{let}$-binding whose bound value is an
  inheritance-calculus expression), the translation is
  applied compositionally from leaves to root: each
  $\lambda$-calculus construct is replaced by its
  inheritance-calculus equivalent via~$\mathcal{T}$,
  treating inheritance-calculus subexpressions as opaque
  values.
  Since each rule of~$\mathcal{T}$
  (Section~\ref{sec:forward-translation}) is a fixed
  syntactic abstraction depending only on the translated
  subexpressions, this bottom-up application is
  well-defined and yields a pure inheritance-calculus
  program for step~(3).

  A $\mathscr{U}$-program $P$ converges,
  $\mathrm{eval}_{\mathscr{U}}(P)$, iff its translation
  converges.
  \fi
\end{definition}

\begin{remark}\label{rem:universe-properties}
  \ifInheritanceChinese
  惰性 $\lambda$-演算与继承演算是 $\mathscr{U}$ 的保守性限制:各自通过去掉另一方的构造子来得到。
  对于纯继承演算程序,步骤~(1) 与~(2) 是平凡的,因此 $\mathrm{eval}_{\mathscr{U}}$ 与继承演算的语义一致。
  对于纯 $\lambda$-演算程序,步骤~(1)--(3) 的复合由充分性(定理~\ref{thm:adequacy})与惰性求值一致。
  \else
  The lazy $\lambda$-calculus and inheritance-calculus are
  conservative restrictions of $\mathscr{U}$: each is
  obtained by removing the other's constructors.
  For a pure inheritance-calculus program, steps~(1)
  and~(2) are vacuous, so
  $\mathrm{eval}_{\mathscr{U}}$ agrees with the
  inheritance-calculus semantics.
  For a pure $\lambda$-calculus program, the composition
  of steps~(1)--(3) agrees with lazy evaluation by
  Adequacy (Theorem~\ref{thm:adequacy}).
  \fi
\end{remark}

\ifInheritanceChinese
三步翻译是有效的(ANF 转换与 $\mathcal{T}$ 是语法变换),而继承演算的递归求值(第~\ref{sec:mixin-trees} 节)是一个半判定过程,因此 $\mathrm{eval}_{\mathscr{U}}$ 是递归可枚举的,满足 Felleisen 框架~\cite{felleisen1991-expressive-power} 的要求。
\else
The three-step translation is effective (ANF conversion
and~$\mathcal{T}$ are syntactic transformations) and the
recursive evaluation of inheritance-calculus
(Section~\ref{sec:mixin-trees}) is a semi-decision
procedure, so $\mathrm{eval}_{\mathscr{U}}$ is
recursively enumerable, as required by
Felleisen's framework~\cite{felleisen1991-expressive-power}.
\fi

\begin{lemma}[\bilingual{$\mathbf{let}$-binding is macro-expressible
  in the lazy $\lambda$-calculus}{$\mathbf{let}$-绑定在惰性 $\lambda$-演算中是宏可表达的}]
  \label{lem:let-macro}
  \ifInheritanceChinese
  翻译 $\mathbf{let}\; x = e_1 \;\mathbf{in}\; e_2 \;\mapsto\; (\lambda x.\, e_2)\; e_1$ 是 $\mathbf{let}$ 在惰性 $\lambda$-演算中的一个宏可表达的消去。
  \else
  The translation
  $\mathbf{let}\; x = e_1 \;\mathbf{in}\; e_2
  \;\mapsto\; (\lambda x.\, e_2)\; e_1$
  is a macro-expressible elimination of $\mathbf{let}$
  in the lazy $\lambda$-calculus.
  \fi
\end{lemma}

\begin{proof}
  \ifInheritanceChinese
  该翻译将每个 $\mathbf{let}$-绑定替换为一个仅依赖于两个子表达式 $e_1$ 与 $e_2$ 的 $\lambda$-表达式。
  这是 Felleisen~\cite{felleisen1991-expressive-power} 的命题 3.7。
  \else
  The translation replaces each $\mathbf{let}$-binding
  with a $\lambda$-expression that depends only on the
  two subexpressions $e_1$ and~$e_2$.
  This is Proposition~3.7 of
  Felleisen~\cite{felleisen1991-expressive-power}.
  \fi
\end{proof}

\begin{lemma}[\bilingual{Equal expressiveness of the lazy
  $\lambda$-calculus and its ANF fragment}{惰性 $\lambda$-演算与其 ANF 片段的等同表达力}]
  \label{lem:anf-equiv}
  \ifInheritanceChinese
  在 $\mathscr{U}$ 中,惰性 $\lambda$-演算与 ANF 惰性 $\lambda$-演算(第~\ref{sec:forward-translation} 节)具有相同的宏可表达力。
  \else
  In $\mathscr{U}$, the lazy $\lambda$-calculus and the
  ANF lazy $\lambda$-calculus
  (Section~\ref{sec:forward-translation}) are equally
  macro-expressive.
  \fi
\end{lemma}

\begin{proof}
  \ifInheritanceChinese
  每个 ANF 项都是一个 $\lambda$-项,因此 ANF 片段嵌入完整 $\lambda$-演算的映射是恒等变换,平凡地宏可表达。
  对于反方向,ANF 变换~\cite{flanagan1993-essence-compiling-continuations} 由两步构成,每步都是宏可表达的:
  \begin{enumerate}
    \item \emph{\bilingual{Application flattening}{应用展开}。}
      将每个嵌套应用 $(e_1\; e_2)\; e_3$ 替换为 $\mathbf{let}\; x = e_1\; e_2 \;\mathbf{in}\; x\; e_3$。
      其语法抽象为 $A(\cdot_1, \cdot_2, \cdot_3) = \mathbf{let}\; x = \cdot_1\; \cdot_2 \;\mathbf{in}\; x\; \cdot_3$,仅依赖三个子表达式。
    \item \emph{\bilingual{Value lifting}{值提升}。}
      将应用位置上每个自由值 $V$ 替换为 $\mathbf{let}\; x = V \;\mathbf{in}\; x$。
      其语法抽象为 $A(\cdot_1) = \mathbf{let}\; x = \cdot_1 \;\mathbf{in}\; x$。
  \end{enumerate}
  由引理~\ref{lem:let-macro},每个所得 $\mathbf{let}$-绑定本身在 $\lambda$-演算中是宏可表达的。
  由引理~\ref{lem:macro-composition},复合这些宏可表达翻译,得到一个从完整 $\lambda$-演算到其 ANF 片段的宏可表达翻译。
  \else
  Every ANF term is a $\lambda$-term, so the embedding
  of the ANF fragment into the full $\lambda$-calculus is
  the identity, which is trivially macro-expressible.
  For the reverse direction, the ANF
  transformation~\cite{flanagan1993-essence-compiling-continuations} consists of
  two steps, each macro-expressible:
  \begin{enumerate}
    \item \emph{Application flattening.}
      Replace each nested application
      $(e_1\; e_2)\; e_3$ with
      $\mathbf{let}\; x = e_1\; e_2 \;\mathbf{in}\;
      x\; e_3$.
      The syntactic abstraction is
      $A(\cdot_1, \cdot_2, \cdot_3)
      = \mathbf{let}\; x = \cdot_1\; \cdot_2
      \;\mathbf{in}\; x\; \cdot_3$,
      depending only on the three subexpressions.
    \item \emph{Value lifting.}
      Replace each free value~$V$ used in an application
      position with
      $\mathbf{let}\; x = V \;\mathbf{in}\; x$.
      The syntactic abstraction is
      $A(\cdot_1) = \mathbf{let}\; x = \cdot_1
      \;\mathbf{in}\; x$.
  \end{enumerate}
  By Lemma~\ref{lem:let-macro}, each resulting
  $\mathbf{let}$-binding is itself
  macro-expressible in the $\lambda$-calculus.
  By Lemma~\ref{lem:macro-composition}, composing
  these macro-expressible translations
  yields a macro-expressible translation from the full
  $\lambda$-calculus to its ANF fragment.
  \fi
\end{proof}

\begin{theorem}[\bilingual{Forward macro-expressibility}{正向宏可表达性}]
  \label{thm:forward-macro}
  \ifInheritanceChinese
  翻译 $\mathcal{T}$(第~\ref{sec:forward-translation} 节)是 $\lambda$-演算嵌入继承演算的一个宏可表达的嵌入。
  \else
  The translation $\mathcal{T}$
  (Section~\ref{sec:forward-translation}) is a macro-expressible
  embedding of the $\lambda$-calculus into inheritance-calculus.
  \fi
\end{theorem}

\begin{proof}
  \ifInheritanceChinese
  我们验证 Felleisen~\cite[Definition~3.11]{felleisen1991-expressive-power} 的条件 E4:对于每个元数为 $a$ 的 $\lambda$-演算构造子 $F$,必须存在一个固定的 $a$ 元语法抽象 $A_F(\cdot_1,\ldots,\cdot_a)$ 使得 $\mathcal{T}(F(e_1,\ldots,e_a)) = A_F(\mathcal{T}(e_1),\ldots,\mathcal{T}(e_a))$。
  我们逐条检验 $\mathcal{T}$ 的规则(第~\ref{sec:forward-translation} 节):
  \begin{enumerate}
    \item \emph{\bilingual{$\lambda$-bound variable}{$\lambda$-绑定变量}}($x$,de~Bruijn 索引为 $n'$)。
      这是一个零元构造子,因此没有子表达式。
      语法抽象是常量 $A = \dbi{n'}.\mathrm{argument}$。
    \item \emph{\bilingual{Let-bound variable}{Let-绑定变量}}($x$)。
      零元。$A = x.\mathrm{result}$。
    \item \emph{\bilingual{$\lambda$-abstraction}{$\lambda$-抽象}},记为 $\lambda x.\, M$。
      关于体 $M$ 是一元的。
      $A(\cdot_1) = \{\mathrm{argument} \mapsto \{\},\; \mathrm{result} \mapsto \cdot_1\}$。
    \item \emph{\bilingual{Application let-binding}{应用 Let-绑定}}($\mathbf{let}\; x = V_1\; V_2 \;\mathbf{in}\; M$)。
      关于 $V_1$、$V_2$ 与 $M$ 是三元的。
      $A(\cdot_1, \cdot_2, \cdot_3) = \{x \mapsto \{\cdot_1,\; \mathrm{argument} \mapsto \cdot_2\},\; \mathrm{result} \mapsto \cdot_3\}$。
    \item \emph{\bilingual{Value let-binding}{值 Let-绑定}}($\mathbf{let}\; x = V \;\mathbf{in}\; M$)。
      关于 $V$ 与 $M$ 是二元的。
      $A(\cdot_1, \cdot_2) = \{x \mapsto \{\mathrm{result} \mapsto \cdot_1\},\; \mathrm{result} \mapsto \cdot_2\}$。
    \item \emph{\bilingual{Tail call}{尾调用}}($V_1\; V_2$)。
      二元。\newline
      $A(\cdot_1, \cdot_2) = \{\mathrm{tailCall} \mapsto \{\cdot_1,\; \mathrm{argument} \mapsto \cdot_2\},\; \mathrm{result} \mapsto \{\mathrm{tailCall}.\mathrm{result}\}\}$。
  \end{enumerate}
  在每种情况下,语法抽象 $A_F$ 都是一个固定的继承演算上下文,其空位由子表达式的递归翻译填充,从而满足 E4。
  \else
  We verify condition~E4 of
  Felleisen~\cite[Definition~3.11]{felleisen1991-expressive-power}:
  for each $\lambda$-calculus construct~$F$ of arity~$a$,
  there must exist a fixed $a$-ary syntactic abstraction
  $A_F(\cdot_1,\ldots,\cdot_a)$ over
  inheritance-calculus such that
  $\mathcal{T}(F(e_1,\ldots,e_a)) =
  A_F(\mathcal{T}(e_1),\ldots,\mathcal{T}(e_a))$.
  We check each rule of $\mathcal{T}$
  (Section~\ref{sec:forward-translation}):
  \begin{enumerate}
    \item \emph{$\lambda$-bound variable}
      ($x$ with de~Bruijn index $n'$).
      This is a nullary construct and therefore has no subexpressions.
      The syntactic abstraction is the constant
      $A = \dbi{n'}.\mathrm{argument}$.
    \item \emph{Let-bound variable} ($x$).
      Nullary.
      $A = x.\mathrm{result}$.
    \item \emph{$\lambda$-abstraction}, denoted $\lambda x.\, M$.
      Unary in the body~$M$.
      $A(\cdot_1) = \{\mathrm{argument} \mapsto \{\},\;
      \mathrm{result} \mapsto \cdot_1\}$.
    \item \emph{Application let-binding}
      ($\mathbf{let}\; x = V_1\; V_2
      \;\mathbf{in}\; M$).
      Ternary in~$V_1$, $V_2$, and~$M$.
      $A(\cdot_1, \cdot_2, \cdot_3) =
      \{x \mapsto \{\cdot_1,\;
        \mathrm{argument} \mapsto \cdot_2\},\;
      \mathrm{result} \mapsto \cdot_3\}$.
    \item \emph{Value let-binding}
      ($\mathbf{let}\; x = V \;\mathbf{in}\; M$).
      Binary in~$V$ and~$M$.
      $A(\cdot_1, \cdot_2) =
      \{x \mapsto \{\mathrm{result} \mapsto \cdot_1\},\;
      \mathrm{result} \mapsto \cdot_2\}$.
    \item \emph{Tail call} ($V_1\; V_2$).
      Binary.\newline
      $A(\cdot_1, \cdot_2) =
      \{\mathrm{tailCall} \mapsto
        \{\cdot_1,\;
        \mathrm{argument} \mapsto \cdot_2\},\;
        \mathrm{result} \mapsto
      \{\mathrm{tailCall}.\mathrm{result}\}\}$.
  \end{enumerate}
  In each case the syntactic abstraction~$A_F$ is a fixed
  inheritance-calculus context whose holes are filled
  by the recursive translations of the subexpressions,
  satisfying~E4.
  \fi
\end{proof}

\begin{theorem}[\bilingual{Nonexpressibility of inheritance}{继承的不可宏表达性}]
  \label{thm:cartesian-nonexpressibility}
  \ifInheritanceChinese
  惰性 $\lambda$-演算不能宏可表达 $\mathscr{U}$(定义~\ref{def:language-universe})的继承构造子。
  \else
  The lazy $\lambda$-calculus cannot macro-express the
  inheritance constructors of $\mathscr{U}$
  (Definition~\ref{def:language-universe}).
  \fi
\end{theorem}

\begin{proof}
  \ifInheritanceChinese
  我们对 $L_0 = \text{惰性 $\lambda$-演算}$,$L_1 = \mathscr{U}$ 应用 Felleisen~\cite{felleisen1991-expressive-power} 的定理 3.14(i)。
  由定义~\ref{def:language-universe},$\mathscr{U} = L_0 + \{$mixin-树构造子$\}$ 是 $L_0$ 的一个保守性扩展(Felleisen~\cite{felleisen1991-expressive-power} 的定义 3.2)。
  我们验证四个条件:
  \begin{enumerate}
    \item[(i)] \emph{$L_0 \subseteq L_1$:} $\mathscr{U}$ 的语法包含所有 $L_0$-短语。
    \item[(ii)] \emph{不引入新的 $L_0$-短语:} mixin-树构造子不在 $L_0$ 中引入新短语。
    \item[(iii)] \emph{构造子不相交:} mixin-树构造子与 $L_0$-构造子不相交。
    \item[(iv)] \emph{语义保持:} 对纯 $\lambda$-项,定义~\ref{def:language-universe} 的步骤~(1) 与~(2) 是平凡的,步骤~(3) 由充分性(定理~\ref{thm:adequacy})给出惰性求值,因此 $\mathrm{eval}_{\mathscr{U}}$ 在 $L_0$-短语上与 $\mathrm{eval}_{L_0}$ 一致。
  \end{enumerate}
  该定理要求我们给出两个 $L_0$-项 $M$ 与 $N$,满足 $M \cong_0 N$ 但 $M \not\cong_1 N$,其中 $\cong_0$ 是 $L_0$ 中的操作等价,$\cong_1$ 是 $\mathscr{U}$ 中的操作等价。

  \paragraph{\bilingual{Operational equivalence in $L_0$}{$L_0$ 中的操作等价}}
  由 Böhm 树等价~\cite{abramsky1993-full-abstraction-lazy-lambda,abramsky1995-games-full-abstraction-lazy-lambda}(定理~\ref{thm:first-order-semantics}),对于封闭 $\lambda$-项 $M$ 与 $N$:$M \cong_0 N$ 当且仅当 $\mathrm{BT}(M) = \mathrm{BT}(N)$。

  \paragraph{\bilingual{Separating pair}{分离对}}
  定义 Church 布尔值与 Church 布尔相等:$\mathrm{true} = \lambda t.\,\lambda f.\, t$,$\mathrm{false} = \lambda t.\,\lambda f.\, f$,以及 $\mathrm{eq} = \lambda a.\,\lambda b.\, (a\;b)\;(b\;\mathrm{false}\;\mathrm{true})$。
  令
  \[
    M = \mathrm{eq}\;\mathrm{false}\;\mathrm{false}
    \qquad
    N = \mathrm{eq}\;\mathrm{true}\;\mathrm{true}.
  \]
  两者在惰性 $\lambda$-演算中均归约为 Church~$\mathrm{true}$,因此 $\mathrm{BT}(M) = \mathrm{BT}(N)$,从而 $M \cong_0 N$。

  \paragraph{\bilingual{Distinguishing $\mathscr{U}$-context}{区分 $\mathscr{U}$-上下文}}
  我们现在构造一个区分 $M$ 与 $N$ 的 $\mathscr{U}$-上下文。
  $M$ 与 $N$ 是语法封闭的 $\lambda$-项;由定义~\ref{def:language-universe},它们在 $\mathscr{U}$ 中的语义通过 ANF 转换后接 $\mathcal{T}$ 将其翻译为继承演算来给出。
  在扩展的 ANF 约定(第~\ref{sec:forward-translation} 节)下,应用位置上的每个自由值首先被绑定到一个名字:
  \else
  We apply Theorem~3.14(i) of
  Felleisen~\cite{felleisen1991-expressive-power} to the pair
  $L_0 = \text{lazy $\lambda$-calculus}$,
  $L_1 = \mathscr{U}$.
  By Definition~\ref{def:language-universe},
  $\mathscr{U} = L_0 + \{$mixin-tree constructors$\}$
  is a conservative extension of~$L_0$
  (Definition~3.2 of
  Felleisen~\cite{felleisen1991-expressive-power}).
  We verify the four conditions:
  \begin{enumerate}
    \item[(i)] \emph{$L_0 \subseteq L_1$:}
      The syntax of~$\mathscr{U}$ includes all
      $L_0$-phrases.
    \item[(ii)] \emph{No new $L_0$-phrases:}
      The mixin-tree constructors do not introduce new
      phrases in~$L_0$.
    \item[(iii)] \emph{Disjoint constructors:}
      The mixin-tree constructors are disjoint from the
      $L_0$-constructors.
    \item[(iv)] \emph{Semantics preserved:}
      For pure $\lambda$-terms, steps~(1) and~(2) of
      Definition~\ref{def:language-universe} are vacuous
      and step~(3) yields lazy evaluation by Adequacy
      (Theorem~\ref{thm:adequacy}), so
      $\mathrm{eval}_{\mathscr{U}}$ agrees with
      $\mathrm{eval}_{L_0}$ on $L_0$-phrases.
  \end{enumerate}
  The theorem requires us to exhibit two $L_0$-terms
  $M$ and $N$ with $M \cong_0 N$ but
  $M \not\cong_1 N$, where $\cong_0$ is operational
  equivalence in $L_0$ and $\cong_1$ is operational
  equivalence in $\mathscr{U}$.

  \paragraph{Operational equivalence in $L_0$}
  By B\"ohm tree
  equivalence~\cite{abramsky1993-full-abstraction-lazy-lambda,abramsky1995-games-full-abstraction-lazy-lambda}
  (Theorem~\ref{thm:first-order-semantics}), for closed
  $\lambda$-terms $M$ and $N$:
  $M \cong_0 N$ iff $\mathrm{BT}(M) = \mathrm{BT}(N)$.

  \paragraph{Separating pair}
  Define Church booleans and Church boolean equality:
  $\mathrm{true} = \lambda t.\,\lambda f.\, t$,
  $\mathrm{false} = \lambda t.\,\lambda f.\, f$, and
  $\mathrm{eq} = \lambda a.\,\lambda b.\,
  (a\;b)\;(b\;\mathrm{false}\;\mathrm{true})$.
  Let
  \[
    M = \mathrm{eq}\;\mathrm{false}\;\mathrm{false}
    \qquad
    N = \mathrm{eq}\;\mathrm{true}\;\mathrm{true}.
  \]
  Both reduce to Church~$\mathrm{true}$ in the
  lazy $\lambda$-calculus, so
  $\mathrm{BT}(M) = \mathrm{BT}(N)$ and hence
  $M \cong_0 N$.

  \paragraph{Distinguishing $\mathscr{U}$-context}
  We now construct a $\mathscr{U}$-context that
  separates $M$ and $N$.
  $M$ and $N$ are syntactically closed $\lambda$-terms;
  by Definition~\ref{def:language-universe}, their
  semantics in $\mathscr{U}$ is given by translating
  them to inheritance-calculus via ANF conversion followed
  by~$\mathcal{T}$.
  Under the extended ANF convention
  (Section~\ref{sec:forward-translation}), every free
  value used in an application is first bound to a name:
  \fi
  \begin{align*}
    \mathrm{ANF}(M) ={} &
    \mathbf{let}\;\mathrm{eq} =
    \lambda a.\,\lambda b.\,
    \mathbf{let}\; x_1 = a\;b \;\mathbf{in}\; \\
    &\qquad
    \mathbf{let}\; f_1 = \mathrm{false}
    \;\mathbf{in}\;
    \mathbf{let}\; x_2 = b\;f_1 \;\mathbf{in}\; \\
    &\qquad
    \mathbf{let}\; t_1 = \mathrm{true}
    \;\mathbf{in}\;
    \mathbf{let}\; x_3 = x_2\;t_1
    \;\mathbf{in}\;
    x_1\;x_3 \\
    & \mathbf{in}\;\;
    \mathbf{let}\; f_1 = \mathrm{false}
    \;\mathbf{in}\;
    \mathbf{let}\;\mathrm{eqFalse} = \mathrm{eq}\;f_1 \\
    &\quad
    \;\mathbf{in}\;
    \mathbf{let}\; f_2 = \mathrm{false}
    \;\mathbf{in}\;
    \mathrm{eqFalse}\;f_2
  \end{align*}
  \ifInheritanceChinese
  应用 $\mathcal{T}$:
  \else
  Applying~$\mathcal{T}$:
  \fi
  \begin{align*}
    \mathcal{T}(\mathrm{ANF}(M)) \;=\; \{&
      \mathrm{eq} \mapsto \{\mathrm{result} \mapsto
      \mathcal{T}(\lambda a.\,\lambda b.\,\cdots)\}, \\
      & f_1 \mapsto \{\mathrm{result} \mapsto
      \mathcal{T}(\mathrm{false})\}, \\
      & \mathrm{eqFalse} \mapsto
      \{\mathrm{eq}.\mathrm{result},\;
        \mathrm{argument} \mapsto
      \{f_1.\mathrm{result}\}\}, \\
      & f_2 \mapsto \{\mathrm{result} \mapsto
      \mathcal{T}(\mathrm{false})\}, \\
      & \mathrm{tailCall} \mapsto
      \{\mathrm{eqFalse}.\mathrm{result},\;
        \mathrm{argument} \mapsto
      \{f_2.\mathrm{result}\}\}, \\
      & \mathrm{result} \mapsto
      \{\mathrm{tailCall}.\mathrm{result}\}
    \;\}
  \end{align*}
  \ifInheritanceChinese
  对 $N$ 类似,以 $\mathrm{true}$ 替换 $\mathrm{false}$。
  注意 $\mathrm{eq}$ 是 $\mathcal{T}(\mathrm{ANF}(M))$ 的一个具名子项,通过值 let-绑定绑定。
  在 $\mathrm{eq}$ 的翻译内部,$\lambda b$ 作用域包含 $x_1$、$f_1$、$x_2$、$t_1$、$x_3$ 的 let-绑定以及一个尾调用。
  特别地,$\dbi{1}.\mathrm{argument}$ 到达外层 $\lambda a$ 作用域,而 $\dbi{0}.\mathrm{argument}$ 到达 $\lambda b$ 自身。\footnote{
    这些索引使用压缩约定,仅计数 $\lambda$-作用域。
    如翻译脚注(第~\ref{sec:forward-translation} 节)所述,每个 let-绑定与尾调用在继承演算中引入一个额外的作用域层级;形式索引相应地更大,但目标作用域相同。
  }

  $\mathscr{U}$-上下文使用继承在 $\mathrm{eq}$ 的 $\lambda b$ 作用域内添加一个 $\mathrm{repr}$ 记录,该记录投影第一个操作数。
  定义上下文 $C[\cdot]$ 为:
  \else
  and symmetrically for $N$ with $\mathrm{true}$ in place
  of $\mathrm{false}$.
  Note that $\mathrm{eq}$ is a named child
  of $\mathcal{T}(\mathrm{ANF}(M))$, bound via a value
  let-binding.
  Inside the translation of $\mathrm{eq}$, the inner
  $\lambda b$ scope contains the let-bindings
  for $x_1$, $f_1$, $x_2$, $t_1$, $x_3$, and a tail call.
  In particular, $\dbi{1}.\mathrm{argument}$ reaches
  the outer $\lambda a$ scope and
  $\dbi{0}.\mathrm{argument}$ reaches $\lambda b$
  itself.\footnote{
    These indices use the condensed
    convention, counting only $\lambda$-scopes.
    As noted in the translation footnote
    (Section~\ref{sec:forward-translation}),
    each let-binding and tail call introduces an
    additional scope level in inheritance-calculus;
    the formal index is correspondingly larger,
    but the target scope is the same.
  }

  The $\mathscr{U}$-context uses inheritance to add
  a $\mathrm{repr}$ record inside the $\lambda b$ scope
  of $\mathrm{eq}$ that projects the first operand.
  Define the context $C[\cdot]$ as:
  \fi
  \begin{align*}
    C[\cdot] ={} & \mathbf{let}\;
    \mathrm{extended} = \{[\cdot],\;
      \mathrm{eq} \mapsto \{
        \mathrm{result} \mapsto \{ \\
          &\quad
          \mathrm{result} \mapsto \{
            \mathrm{repr} \mapsto \{
              \mathrm{firstOperand}
              \mapsto \{\dbi{2}.\mathrm{argument}\}
    \}\}\}\}\} \\
    & \mathbf{in}\;
    \mathrm{extended}.\mathrm{tailCall}.\mathrm{repr}.\mathrm{firstOperand}
  \end{align*}
  \ifInheritanceChinese
  这里 $[\cdot]$ 是空位。
  注意 mixin 中 $\mathrm{result}$ 的两层嵌套:外层 $\mathrm{result}$ 进入 $\mathrm{eq}$ 的值 let-绑定包装器,而内层 $\mathrm{result}$ 进入 $\lambda a$ 作用域。
  第二个 $\mathrm{result}$ 则包含 $\lambda b$ 作用域,因此 $\mathrm{repr}$ 被放置在 $\lambda b$ 作用域层级。
  mixin $\{\mathrm{eq} \mapsto \{\mathrm{result} \mapsto \{\mathrm{result} \mapsto \{\mathrm{repr} \mapsto \cdots\}\}\}\}$ 是一个继承演算表达式:它仅使用 $\mathscr{U}$ 的 mixin-树构造子。
  外层 $\mathbf{let}$-绑定与投影是 $\lambda$-演算构造(引理~\ref{lem:let-macro})。
  因此 $C[\cdot]$ 是一个合法的 $\mathscr{U}$-上下文。

  $\mathrm{repr}$ 中的 de~Bruijn 索引 $\dbi{2}$ 计数作用域层级(在压缩约定下):$\dbi{0}$ 是 $\mathrm{repr}$ 自身,$\dbi{1}$ 是 $\lambda b$ 作用域,而 $\dbi{2}$ 是 $\lambda a$ 作用域。
  因此 $\mathrm{firstOperand}$ 投影 $\lambda a$ 的参数。

  由于 $\mathrm{eqFalse}$(相应地,$\mathrm{eqTrue}$)继承 $\mathrm{eq}.\mathrm{result}$,而 $\mathrm{tailCall}$ 继承 $\mathrm{eqFalse}.\mathrm{result}$($\lambda b$ 作用域),mixin 在 $\lambda b$ 作用域层级放置的 $\mathrm{repr}$ 记录通过继承链传播到 $\mathrm{tailCall}$。
  在 $C[M]$ 中:$\mathrm{eqFalse}$ 绑定向 $\lambda a$ 作用域提供 $\mathrm{argument} \mapsto \{f_1.\mathrm{result}\}$(即 $\mathrm{false}$),因此 $\mathrm{firstOperand}$ 求值为 $\mathrm{false} = \lambda t.\,\lambda f.\, f$。
  在 $C[N]$ 中:$\mathrm{eqTrue}$ 绑定提供 $\mathrm{argument} \mapsto \mathrm{true}$,因此 $\mathrm{firstOperand}$ 求值为 $\mathrm{true} = \lambda t.\,\lambda f.\, t$。

  定义 $D[\cdot] = (\lambda x.\, x\;\Omega\;I)\; C[\cdot]$。
  则 $\mathrm{eval}_{\mathscr{U}}$ 对 $D[M]$ 成立(因为 $\mathrm{false}\;\Omega\;I \to I$,收敛),但对 $D[N]$ 不成立(因为 $\mathrm{true}\;\Omega\;I \to \Omega$,发散)。
  因此 $M \not\cong_1 N$。

  \paragraph{\bilingual{Conclusion}{结论}}
  我们有 $M \cong_0 N$ 但 $M \not\cong_1 N$,故 ${\cong_0} \neq {(\cong_1|_{L_0})}$。
  由 Felleisen~\cite{felleisen1991-expressive-power} 的定理 3.14(i),惰性 $\lambda$-演算不能宏可表达 $\mathscr{U}$ 的继承构造子。
  \else
  Here $[\cdot]$ is the hole.
  Note the two levels of $\mathrm{result}$ nesting in the
  mixin: the outer $\mathrm{result}$ enters the value
  let-binding wrapper for $\mathrm{eq}$, and the inner
  $\mathrm{result}$ enters the $\lambda a$ scope.
  The second $\mathrm{result}$ then contains the
  $\lambda b$ scope, so $\mathrm{repr}$ is placed at
  the $\lambda b$ scope level.
  The mixin $\{\mathrm{eq} \mapsto \{\mathrm{result}
      \mapsto \{\mathrm{result} \mapsto
  \{\mathrm{repr} \mapsto \cdots\}\}\}\}$ is an
  inheritance-calculus expression: it uses only the
  mixin-tree constructors of $\mathscr{U}$.
  The outer $\mathbf{let}$-binding and projection are
  $\lambda$-calculus constructs (Lemma~\ref{lem:let-macro}).
  Thus $C[\cdot]$ is a valid $\mathscr{U}$-context.

  The de~Bruijn index $\dbi{2}$ in
  $\mathrm{repr}$ counts scope levels
  (in the condensed convention):
  $\dbi{0}$ is $\mathrm{repr}$ itself,
  $\dbi{1}$ is the $\lambda b$ scope, and
  $\dbi{2}$ is the $\lambda a$ scope.
  So $\mathrm{firstOperand}$ projects
  $\lambda a$'s argument.

  Since $\mathrm{eqFalse}$ (respectively $\mathrm{eqTrue}$)
  inherits $\mathrm{eq}.\mathrm{result}$,
  and $\mathrm{tailCall}$ inherits
  $\mathrm{eqFalse}.\mathrm{result}$ (the $\lambda b$
  scope), the $\mathrm{repr}$ record placed by the mixin
  at the $\lambda b$ scope level propagates through
  the inheritance chain to $\mathrm{tailCall}$.
  In $C[M]$: the $\mathrm{eqFalse}$ binding supplies
  $\mathrm{argument} \mapsto \{f_1.\mathrm{result}\}$
  (that is, $\mathrm{false}$) to the $\lambda a$ scope, so
  $\mathrm{firstOperand}$ evaluates to
  $\mathrm{false} = \lambda t.\,\lambda f.\, f$.
  In $C[N]$: the $\mathrm{eqTrue}$ binding supplies
  $\mathrm{argument} \mapsto \mathrm{true}$, so
  $\mathrm{firstOperand}$ evaluates to
  $\mathrm{true} = \lambda t.\,\lambda f.\, t$.

  Define $D[\cdot] = (\lambda x.\, x\;\Omega\;I)\;
  C[\cdot]$.
  Then $\mathrm{eval}_{\mathscr{U}}$ holds for $D[M]$
  (since $\mathrm{false}\;\Omega\;I \to I$, which
  converges)
  but not for $D[N]$
  (since $\mathrm{true}\;\Omega\;I \to \Omega$, which
  diverges).
  Therefore $M \not\cong_1 N$.

  \paragraph{Conclusion}
  We have $M \cong_0 N$ but $M \not\cong_1 N$, so
  ${\cong_0} \neq {(\cong_1|_{L_0})}$.
  By Theorem~3.14(i) of
  Felleisen~\cite{felleisen1991-expressive-power},
  the lazy $\lambda$-calculus cannot macro-express the
  inheritance constructors of~$\mathscr{U}$.
  \fi
\end{proof}

\begin{theorem}[\bilingual{Expressive asymmetry}{表达力的不对称}]
  \label{thm:expressive-asymmetry}
  \ifInheritanceChinese
  在语言全集 $\mathscr{U}$(定义~\ref{def:language-universe})中,继承演算比惰性 $\lambda$-演算具有严格更强的表达力:它能宏可表达 $\lambda$-演算能宏可表达的每个 $\mathscr{U}$-构造,但反之不然(Felleisen~\cite{felleisen1991-expressive-power} 的定义 3.17)。
  \else
  In the language universe $\mathscr{U}$
  (Definition~\ref{def:language-universe}),
  inheritance-calculus is strictly more expressive than
  the lazy $\lambda$-calculus: it can macro-express every
  $\mathscr{U}$-construct that the $\lambda$-calculus
  can, but not conversely
  (Definition~3.17 of
  Felleisen~\cite{felleisen1991-expressive-power}).
  \fi
\end{theorem}

\begin{proof}
  \ifInheritanceChinese
  我们须证明:(a)~继承演算能宏可表达惰性 $\lambda$-演算能宏可表达的每个 $\mathscr{U}$-构造,以及 (b)~反之不成立。

  \paragraph{\bilingual{Part (b): the lazy $\lambda$-calculus $\not\geq$
  inheritance-calculus}{第 (b) 部分:惰性 $\lambda$-演算 $\not\geq$ 继承演算}}
  $\lambda$-演算包含其自身的构造子($\lambda$、应用、变量),并能宏可表达 $\mathbf{let}$-绑定(引理~\ref{lem:let-macro}),但不能宏可表达 $\mathscr{U}$ 的继承构造子(定理~\ref{thm:cartesian-nonexpressibility})。
  继承演算包含那些构造子。
  因此,惰性 $\lambda$-演算对于 $\mathscr{U}$ 而言表达力不及继承演算。

  \paragraph{\bilingual{Part (a): inheritance-calculus $\geq$ the lazy
  $\lambda$-calculus}{第 (a) 部分:继承演算 $\geq$ 惰性 $\lambda$-演算}}
  我们证明继承演算能宏可表达 $\lambda$-演算所包含的或能宏可表达的每个 $\mathscr{U}$-构造。
  $\lambda$-演算包含 $\lambda$-抽象、应用与变量;它能宏可表达 $\mathbf{let}$-绑定(引理~\ref{lem:let-macro})。
  我们须证明继承演算能宏可表达这四个构造。
  翻译分两个阶段进行:
  \begin{enumerate}
    \item \emph{\bilingual{ANF conversion}{ANF 转换}。}
      A-范式变换~\cite{flanagan1993-essence-compiling-continuations} 将嵌套应用替换为 $\mathbf{let}$-绑定;由引理~\ref{lem:anf-equiv},它在 $\lambda$-演算中是宏可表达的。
      在反方向,$\mathbf{let}$-绑定由 Felleisen~\cite{felleisen1991-expressive-power} 的命题 3.7(重述为引理~\ref{lem:let-macro})在 $\lambda$-演算中是宏可表达的。
      因此,ANF 转换是 $\lambda$-演算与其 ANF 片段之间的一个宏可表达的等价。
    \item \emph{\bilingual{The translation $\mathcal{T}$}{翻译 $\mathcal{T}$}。}
      每个 ANF $\lambda$-演算构造子映射到继承演算上的一个固定语法抽象(定理~\ref{thm:forward-macro}),满足 Felleisen~\cite[Definition~3.11]{felleisen1991-expressive-power} 的条件 E4。
  \end{enumerate}
  由引理~\ref{lem:macro-composition},复合两个宏可表达翻译,得到一个从 $\lambda$-演算(含 $\mathbf{let}$)到继承演算的宏可表达翻译。
  由于继承演算平凡地包含其自身的构造子,它能宏可表达 $\lambda$-演算所包含的或能宏可表达的每个 $\mathscr{U}$-构造。
  \else
  We must show:
  (a)~inheritance-calculus can macro-express every
  $\mathscr{U}$-construct that the lazy $\lambda$-calculus
  can macro-express, and
  (b)~the converse fails.

  \paragraph{Part (b): the lazy $\lambda$-calculus $\not\geq$
  inheritance-calculus}
  The $\lambda$-calculus contains its own constructors
  ($\lambda$, application, variable) and can macro-express
  $\mathbf{let}$-binding (Lemma~\ref{lem:let-macro}), but
  cannot macro-express the inheritance constructors of
  $\mathscr{U}$
  (Theorem~\ref{thm:cartesian-nonexpressibility}).
  Inheritance-calculus contains those constructors.
  Therefore the lazy $\lambda$-calculus is not at least
  as expressive as inheritance-calculus with respect
  to~$\mathscr{U}$.

  \paragraph{Part (a): inheritance-calculus $\geq$ the lazy
  $\lambda$-calculus}
  We show that inheritance-calculus can macro-express
  every $\mathscr{U}$-construct that the
  $\lambda$-calculus contains or can macro-express.
  The $\lambda$-calculus contains $\lambda$-abstraction,
  application, and variable; it can macro-express
  $\mathbf{let}$-binding
  (Lemma~\ref{lem:let-macro}).
  We must show that inheritance-calculus can
  macro-express all four of these constructs.
  The translation proceeds in two stages:
  \begin{enumerate}
    \item \emph{ANF conversion.}
      The A-normal form
      transformation~\cite{flanagan1993-essence-compiling-continuations}
      replaces nested applications with
      $\mathbf{let}$-bindings; it is macro-expressible
      in the $\lambda$-calculus by
      Lemma~\ref{lem:anf-equiv}.
      In the reverse direction,
      $\mathbf{let}$-bindings are macro-expressible in
      the $\lambda$-calculus by
      Proposition~3.7 of
      Felleisen~\cite{felleisen1991-expressive-power}
      (restated as Lemma~\ref{lem:let-macro}).
      Thus ANF conversion is a macro-expressible
      equivalence between the $\lambda$-calculus and its
      ANF fragment.
    \item \emph{The translation $\mathcal{T}$.}
      Each ANF $\lambda$-calculus construct maps to a
      fixed syntactic abstraction over
      inheritance-calculus
      (Theorem~\ref{thm:forward-macro}), satisfying
      condition~E4 of
      Felleisen~\cite[Definition~3.11]{felleisen1991-expressive-power}.
  \end{enumerate}
  By Lemma~\ref{lem:macro-composition}, composing the
  two macro-expressible translations gives a
  macro-expressible translation from
  the $\lambda$-calculus (with $\mathbf{let}$) into
  inheritance-calculus.
  Since inheritance-calculus trivially contains its own
  constructors, it can macro-express every
  $\mathscr{U}$-construct that the $\lambda$-calculus
  contains or can macro-express.
  \fi
\end{proof}

\section{\bilingual{Multi-Target \texorpdfstring{$\this$}{this} Resolution in Other Systems}{其他系统中的多目标 \texorpdfstring{$\this$}{this} 解析}}
\label{app:scala-multi-target}

\ifInheritanceChinese
本附录说明两个具有代表性的系统——Scala~3 和 NixOS 模块系统——拒绝了继承演算能够自然处理的多目标 $\this$ 解析模式。
\else
This appendix demonstrates that two representative systems, Scala~3
and the NixOS module system, reject the multi-target $\this$ resolution
pattern that inheritance-calculus handles naturally.
\fi

\subsection*{\bilingual{Scala~3}{Scala~3}}

\ifInheritanceChinese
考虑两个独立扩展同一外部类的对象，每个对象各自提供内部 trait 的一份副本：
\else
Consider two objects that independently extend an outer class,
each providing its own copy of an inner trait:
\fi

\begin{lstlisting}[style=scala]
class MyOuter:
  trait MyInner:
    def outer = MyOuter.this

object Object1 extends MyOuter
object Object2 extends MyOuter
object HasMultipleOuters extends Object1.MyInner
                     with Object2.MyInner
\end{lstlisting}

\noindent
\ifInheritanceChinese
Scala~3 以如下错误拒绝了 \texttt{HasMultipleOuters}：
\else
Scala~3 rejects \texttt{HasMultipleOuters} with the error:
\fi

\begin{lstlisting}
trait MyInner is extended twice
object HasMultipleOuters cannot be instantiated since
  it has conflicting base types
  Object1.MyInner and Object2.MyInner
\end{lstlisting}

\noindent
\ifInheritanceChinese
该拒绝并非表面层面的限制，而是 DOT 类型论基础的必然后果~\cite{amin2016-dependent-object-types}。在 DOT 中，一个对象被构造为 $\nu(x : T)\, d$，其中 self 变量~$x$ 绑定到\emph{单一}对象。\texttt{Object1.MyInner} 和 \texttt{Object2.MyInner} 是不同的路径依赖类型，每个都将外部 $\this$ 约束到不同的对象。它们的交集要求 \texttt{outer} 同时返回 \texttt{Object1} 和 \texttt{Object2}，但 DOT 的 self 变量是单值的，使得这一交集不可实现。
\else
The rejection is not a surface-level restriction but a consequence
of DOT's type-theoretic foundations~\cite{amin2016-dependent-object-types}.
In DOT, an object is constructed as $\nu(x : T)\, d$, where the
self variable~$x$ binds to a \emph{single} object.
\texttt{Object1.MyInner} and \texttt{Object2.MyInner} are distinct
path-dependent types, each constraining the outer $\this$
to a different object.
Their intersection requires \texttt{outer} to simultaneously
return \texttt{Object1} and \texttt{Object2}, but DOT's self
variable is single-valued, making this intersection unrealizable.
\fi

\ifInheritanceChinese
等价的继承演算定义如下：
\else
The equivalent inheritance-calculus definition is:
\fi
\begin{align*}
  \{ \quad \mathrm{MyOuter} &\mapsto \{
      \mathrm{MyInner} \mapsto \{
    \mathrm{outer} \mapsto \{\qualifiedthis{\mathrm{MyOuter}}\}\}\}, \\
    \mathrm{Object1} &\mapsto \{\mathrm{MyOuter}\}, \\
    \mathrm{Object2} &\mapsto \{\mathrm{MyOuter}\}, \\
    \mathrm{HasMultipleOuters} &\mapsto \{
      \mathrm{Object1}.\mathrm{MyInner},\;
    \mathrm{Object2}.\mathrm{MyInner}\}
  \quad \}
\end{align*}
\ifInheritanceChinese
此定义在继承演算中是良定义的。$\mathrm{outer}$ 内部的引用 $\qualifiedthis{\mathrm{MyOuter}}$ 具有 de~Bruijn 索引 $n = 1$：从 $\mathrm{outer}$ 的封闭作用域 $\mathrm{MyInner}$ 出发，经过一步 $\this$ 便到达 $\mathrm{MyOuter}$。当 $\mathrm{HasMultipleOuters}$ 同时继承 $\mathrm{Object1}.\mathrm{MyInner}$ 和 $\mathrm{Object2}.\mathrm{MyInner}$ 时，$\this$ 函数（方程~\ref{eq:this}）通过在 $\supers(\mathrm{HasMultipleOuters})$ 中搜索与 $\mathrm{MyInner}$ 定义位置匹配的覆盖路径来解析 $\qualifiedthis{\mathrm{MyOuter}}$。它找到两条继承位置路径，一条经由 $\mathrm{Object1}$，另一条经由 $\mathrm{Object2}$，并返回两者。两条路径都通向继承自同一 $\mathrm{MyOuter}$ 的记录，因此查询 $\mathrm{HasMultipleOuters}.\mathrm{outer}$ 无论走哪条路由都能得到 $\mathrm{MyOuter}$ 的 $\properties$。由于属性解析通过集合并来聚合结果（这是一种本质上幂等的操作），来自两条路由的相同属性定义无冲突地合并。因此，该演算消除了对任何人为路由优先级或线性化规则的需求。
\else
This is well-defined in inheritance-calculus.
The reference
$\qualifiedthis{\mathrm{MyOuter}}$ inside $\mathrm{outer}$
has de~Bruijn index $n = 1$:
starting from $\mathrm{outer}$'s enclosing scope
$\mathrm{MyInner}$, one $\this$ step reaches
$\mathrm{MyOuter}$.
When $\mathrm{HasMultipleOuters}$ inherits from both
$\mathrm{Object1}.\mathrm{MyInner}$ and
$\mathrm{Object2}.\mathrm{MyInner}$,
the $\this$ function (equation~\ref{eq:this}) resolves
$\qualifiedthis{\mathrm{MyOuter}}$ by searching through
$\supers(\mathrm{HasMultipleOuters})$ for override paths matching
$\mathrm{MyInner}$'s definition site.
It finds two inheritance-site paths, one through
$\mathrm{Object1}$ and one through $\mathrm{Object2}$, and
returns both.
Both paths lead to records that inherit from the same
$\mathrm{MyOuter}$, so querying
$\mathrm{HasMultipleOuters}.\mathrm{outer}$ yields the
$\properties$ of $\mathrm{MyOuter}$ regardless of which
route is taken.
Because attribute resolution aggregates results via set union (which is an intrinsically idempotent operation), identical property definitions from the two routes merge without conflict. Thus, the calculus eliminates the need for any artificial route-priority or linearization rules.
\fi

\subsection*{\bilingual{NixOS Module System}{NixOS 模块系统}}

\ifInheritanceChinese
NixOS 模块系统以静态错误拒绝了相同的模式。该翻译使用了模块系统自身的抽象：\texttt{deferredModule} 对应 trait（即未求值的模块值，被导入到其他不动点中），\texttt{submoduleWith} 对应对象（在其自身的不动点中求值）。
\else
The NixOS module system rejects the same pattern with a static error.
The translation uses the module system's own abstractions:
\texttt{deferredModule} for traits, which are unevaluated module values imported
into other fixpoints, and \texttt{submoduleWith} for objects, which are evaluated
in their own fixpoints.
\fi

\begin{lstlisting}[style=nix]
let
  lib = (import <nixpkgs> { }).lib;
  result = lib.evalModules {
    modules = [
      (toplevel@{ config, ... }: {
        # class MyOuter { trait MyInner {
        #   def outer = MyOuter.this } }
        options.MyOuter = lib.mkOption {
          default = { };
          type = lib.types.deferredModuleWith {
            staticModules = [
              (MyOuter: {
                options.MyInner = lib.mkOption {
                  default = { };
                  type = lib.types.deferredModuleWith {
                    staticModules = [
                      (MyInner: {
                        options.outer = lib.mkOption {
                          default = { };
                          type =
                            lib.types.deferredModuleWith {
                              staticModules =
                                [ toplevel.config.MyOuter ];
                            };
                        };
                      })
                    ];
                  };
                };
              })
            ];
          };
        };
        # object Object1 extends MyOuter
        options.Object1 = lib.mkOption {
          default = { };
          type = lib.types.submoduleWith {
            modules = [ toplevel.config.MyOuter ];
          };
        };
        # object Object2 extends MyOuter
        options.Object2 = lib.mkOption {
          default = { };
          type = lib.types.submoduleWith {
            modules = [ toplevel.config.MyOuter ];
          };
        };
        # object HasMultipleOuters extends
        #   Object1.MyInner with Object2.MyInner
        options.HasMultipleOuters = lib.mkOption {
          default = { };
          type = lib.types.submoduleWith {
            modules = [
            toplevel.config.Object1.MyInner
            toplevel.config.Object2.MyInner
          ];
          };
        };
      })
    ];
  };
in
  builtins.attrNames result.config.HasMultipleOuters.outer
\end{lstlisting}

\noindent
\ifInheritanceChinese
NixOS 模块系统以如下错误拒绝了此代码：
\else
The NixOS module system rejects this with:
\fi

\begin{lstlisting}
error: The option `HasMultipleOuters.outer'
  in `<unknown-file>'
  is already declared
  in `<unknown-file>'.
\end{lstlisting}

\noindent
\ifInheritanceChinese
该拒绝发生的原因是：\texttt{Object1.\allowbreak{}My\-Inner} 和 \texttt{Object2.\allowbreak{}My\-Inner} 各自携带了 \texttt{My\-Inner} 延迟模块的完整 \texttt{static\-Modules}，其中包括声明 \texttt{options.\allowbreak{}outer}。当 \texttt{HasMultipleOuters} 同时导入两者时，模块系统的 \texttt{merge\-Option\-Decls} 遇到了同一选项的两个声明，并产生静态错误；它无法表达两个声明源自同一 trait 且应被统一这一事实。
\else
The rejection occurs because
\texttt{Object1.\allowbreak{}My\-Inner} and
\texttt{Object2.\allowbreak{}My\-Inner} each carry the full
\texttt{static\-Modules} of the \texttt{My\-Inner} deferred
module, including the declaration
\texttt{options.\allowbreak{}outer}.
When \texttt{HasMultipleOuters} imports both, the module system's
\texttt{merge\-Option\-Decls} encounters two declarations of the
same option and raises a static error; it cannot express that
both declarations originate from the same trait and should be
unified.
\fi

\ifInheritanceChinese
该拒绝与上文 Scala 的``冲突基类型''错误类似：两个系统都假设每个选项或类型成员只有单一的声明位置，并拒绝两条继承路由独立引入同一声明的多目标情形。在继承演算中，相同的模式是良定义的，因为 $\overrides$（方程~\ref{eq:overrides}）能够识别两条路由通向同一定义，而 $\supers$（方程~\ref{eq:supers}）则无重复地收集两个继承位置上下文。
\else
The rejection is analogous to Scala's ``conflicting base types''
error above: both systems assume that each option or type member
has a single declaration site, and reject the multi-target situation
where two inheritance routes introduce the same declaration
independently.
In inheritance-calculus, the same pattern is well-defined because
$\overrides$ (equation~\ref{eq:overrides}) recognizes that both
routes lead to the same definition, and $\supers$
(equation~\ref{eq:supers}) collects both inheritance-site
contexts without duplication.
\fi

\section{Encoding Datalog in Inheritance-Calculus}
\label{app:datalog-encoding}

Section~\ref{sec:emergent-phenomena} observed that
inheritance-calculus natively produces the relational semantics
of logic programming.
This appendix makes that observation precise by giving a
systematic encoding of Datalog into inheritance-calculus.
The encoding uses no extensions to the calculus; it is
expressible in the three primitives of Section~\ref{sec:syntax}.
All examples in this appendix have executable MIXINv2
implementations\anon[~(supplementary material)]{~\cite{mixin2025}}.
There are three ingredients to keep separate.
First, relations are stored as tries, so tuple lookup is indexed
by prefixes rather than by scanning a flat set of facts.
Second, the control skeleton is written in the visitor-pattern
style already used in Section~\ref{sec:case-study}; from a
functional-programming viewpoint, this is the same role played by
Scott-encoded case analysis.
Third, inheritance union and tabled evaluation give the resulting
program its relational, least-fixed-point semantics.
The appendix proceeds in that order: data representation first,
then control flow, then recursion.

\subsection{Relations as Tries}

A binary relation over a finite domain $D = \{a, b, c, \ldots\}$
is represented as a \emph{trie}: a mixin tree whose paths
encode tuples.
The trie is built from three base scopes:
$\mathrm{Relation}$ (abstract base),
$\mathrm{Cons}$ (non-empty node), and $\mathrm{Nil}$ (terminator).
Each constant $d \in D$ inherits from $\mathrm{Cons}$ and carries
a $\mathrm{TailMap}$ that provides per-constant sub-tries:
\begin{align*}
  &\mathrm{Relation} \mapsto \{\},\quad
  \mathrm{Cons} \mapsto \{\mathrm{Relation},\;
  \mathrm{Tail} \mapsto \{\mathrm{Relation}\}\},\quad
  \mathrm{Nil} \mapsto \{\mathrm{Relation}\}\\[2pt]
  &d \mapsto \{\mathrm{Cons},\;
    \mathrm{Tail} \mapsto \{\},\;
  \mathrm{TailMap} \mapsto \{d \mapsto \{\mathrm{Tail}\}\}\}
\end{align*}
A fact $R(d_1, d_2)$ is encoded as a trie path of depth~2,
with a $\mathrm{Nil}$ terminator.
The per-constant sub-tries are accessed via
$\mathrm{TailMap}.d$:
\begin{align*}
  R(a, b) \;\leadsto\quad
  R \mapsto \{a,\; \mathrm{TailMap} \mapsto \{a \mapsto
  \{b,\; \mathrm{TailMap} \mapsto \{b \mapsto \{\mathrm{Nil}\}\}\}\}\}
\end{align*}
Multiple facts for the same relation are accumulated by
inheritance (trie union).
Adding $R(b, c)$:
\[
  R \mapsto \left\{
    \begin{aligned}
      &a,\; b,\\
      &\mathrm{TailMap} \mapsto \left\{
        \begin{aligned}
          &a \mapsto \{b,\; \mathrm{TailMap} \mapsto \{b \mapsto \{\mathrm{Nil}\}\}\},\\
          &b \mapsto \{c,\; \mathrm{TailMap} \mapsto \{c \mapsto \{\mathrm{Nil}\}\}\}
      \end{aligned}\right\}
  \end{aligned}\right\}
\]
The point of the trie representation is not merely that it is a
tree-shaped encoding.
It also provides the indexing structure used by the subsequent
joins: once a constant $d$ has been selected, the encoding moves
directly to $\mathrm{TailMap}.d$ and continues from the matching
sub-trie.
This is why the appendix speaks about tries rather than about an
arbitrary encoding of finite sets of tuples.

\subsection{Infrastructure}

Each constant $d$ carries four groups of boilerplate
definitions used by the encoding.
These play a role analogous to derived type-class instances
in Haskell: each constant must provide a fixed set of
properties that the encoding's generic patterns rely on.
Inheritance-calculus has no built-in derivation mechanism,
so the boilerplate is written out manually; however, the
amount of boilerplate per constant is $O(1)$
and does not grow with the domain size $|D|$,
so adding a new constant
does not require modifying any existing scope, preserving
the expression-problem property.

% 我们可以通过引入类似ECMAScript的Symbol机制来消除这些重复定义，但这让语言更复杂了

\begin{description}
  \item[Acceptance]
    The visitor-pattern infrastructure from
    Section~\ref{sec:case-study}.
    Operationally, this is the object-oriented spelling of a
    case split on the current constant.
    Each constant $d$ defines a $\mathrm{VisitorMap}$
    containing a branch keyed by $d$, a $\mathrm{Visitor}$
    with a $\mathrm{Visited}$ slot for per-constant results,
    and an $\mathrm{Accepted}$ property that delegates to the
    matching branch's $\mathrm{Visited}$ value.
    This is the dispatch mechanism: given a scope that may
    contain multiple constants, $\mathrm{Acceptance}$ selects the
    branch corresponding to each constant present and
    collects the per-constant $\mathrm{Visited}$ results.

  \item[WithVisitorTail]
    Maps each constant $d$ to its sub-trie entry
    $\mathrm{TailMap}.d$, labeled as $\mathrm{VisitorTail}$.
    This provides the right-hand side of a join\footnote{
      In the context of the visitor pattern, this acts as the visitor.
    }: when two relations' $\mathrm{VisitorMap}$ entries are
    composed, $\mathrm{VisitorTail}$ carries the joined
    relation's row data for the matched constant.

  \item[WithVisiteeTail]
    Maps each constant $d$ to its sub-trie entry
    $\mathrm{TailMap}.d$, labeled as $\mathrm{VisiteeTail}$.
    This provides the left-hand side\footnote{
      This acts as the visitee.
    }: the
    relation being iterated. Together,
    $\mathrm{WithVisitorTail}$ and $\mathrm{WithVisiteeTail}$
    enable value-level joins where only constants present
    in \emph{both} relations produce output.

  \item[Replacement]
    Given a $\mathrm{NewTail}$ value, produces a
    $\mathrm{Replaced}$ scope that inherits the label $d$
    but substitutes $\mathrm{TailMap}.d$ with
    $\mathrm{NewTail}$.
    This is the output construction mechanism: it
    builds output tuples by substituting new sub-tries into
    constant wrappers.
    Crucially, $\mathrm{Replacement}$ is defined on the
    per-constant $\mathrm{TailMap}$ entry (not on the merged
    scope), ensuring that output construction operates on
    individual constants rather than their union.
\end{description}

As a running example, we use the executable two-hop encoding
$\mathrm{twoHop}(X, Z) \mathrel{{:}\!{-}} \mathrm{edge}(X, Y),\; \mathrm{edge}(Y, Z)$
from
\anon[\nolinkurl{RelationalTwoHopRepeatedVisitor.mixin.yaml}~(supplementary material)]{\nolinkurl{RelationalTwoHopRepeatedVisitor.mixin.yaml}~\cite{mixin2025}}.
Its top-level scopes are named after the stages of the
encoding: \texttt{\_XAcceptance} iterates the first column,
\texttt{\_YAcceptance} performs the shared-variable join,
\texttt{\_NilAcceptance0} confirms the end of the first fact,
\texttt{\_ZAcceptance} iterates the output column, and
\texttt{\_NilAcceptance1} triggers the final
\texttt{Replacement} chain.
The schematic encodings below should be read as abstractions
of that executable MIXINv2 file.
If one temporarily ignores relational accumulation and reads
each dispatch as an ordinary case split, the control flow is:
choose $X$ from the first relation, choose a matching $Y$,
confirm that the first tuple is complete, choose $Z$ from the
second relation's matching row, confirm that the second tuple
is complete, and emit $(X, Z)$.
The rest of the appendix explains how that familiar
single-valued control skeleton is realised by visitor dispatch
over tries and then lifted to relational semantics by
inheritance union.

\subsection{Encoding Rules}

A Datalog rule $H(\bar{X}) \mathrel{{:}\!{-}} B_1(\bar{Y}_1), \ldots, B_n(\bar{Y}_n)$
is encoded by nesting three mechanisms inside that control
skeleton:

\begin{enumerate}
  \item \textbf{Column iteration.}
    For each column position, an $\mathrm{Acceptance}$ dispatch
    iterates over the constants present.
    The $\mathrm{VisitorMap}$ is populated from
    $\mathrm{WithVisiteeTail}$, so each matched constant's
    $\mathrm{Visitor}$ receives $\mathrm{VisiteeTail}$: the
    per-constant sub-trie at that position.
    The result of processing each constant is placed in
    $\mathrm{Visited}$, and $\mathrm{Accepted}$ collects all
    per-constant results.
    (Details in Sections~\ref{sec:encoding-variables},
    \ref{sec:encoding-constants}, and~\ref{sec:encoding-multi}.)

  \item \textbf{Join.}
    For each shared variable between body atoms, the inner
    $\mathrm{Acceptance}$'s $\mathrm{VisitorMap}$ inherits from
    \emph{both} $\mathrm{WithVisitorTail}$ (from the joined
    relation) and $\mathrm{WithVisiteeTail}$ (from the iterated
    relation).
    A constant $d$ produces a $\mathrm{Visitor}$ entry only if
    $d$ appears in both relations, achieving a value-level
    join.
    The $\mathrm{Visitor}$ then receives both
    $\mathrm{VisiteeTail}$ and $\mathrm{VisitorTail}$,
    providing access to both relations' data for the
    matched constant.

  \item \textbf{Output construction.}
    Inside the innermost $\mathrm{Nil}$ confirmation (verifying
    that each body atom's tuple is complete), chain
    $\mathrm{Replacement}$ operations from the innermost
    output column outward, each setting $\mathrm{NewTail}$ to
    the previous $\mathrm{Replaced}$ value, terminated by
    $\mathrm{Nil}$.
    The final $\mathrm{Replaced}$ value is assigned to
    $\mathrm{Visited}$.
\end{enumerate}

\subsection{Variable Join}
\label{sec:encoding-variables}

Consider the simplest join:
$\mathrm{composed}(X, Z) \mathrel{{:}\!{-}} \alpha(X, Y),\; \beta(Y, Z)$.
The executable self-join instance of this pattern is
\anon[\nolinkurl{RelationalTwoHopRepeatedVisitor.mixin.yaml}~(supplementary material)]{\nolinkurl{RelationalTwoHopRepeatedVisitor.mixin.yaml}~\cite{mixin2025}},
whose scope names are shown in parentheses below.
This is the right place to see the division of labour from the
section introduction: visitor dispatch provides the control flow,
the trie provides direct access to the matching tails, and
$\supers$ accumulates the outputs produced by each successful
match.

The shared variable $Y$ requires matching constants
between $\alpha$'s second column and $\beta$'s first column.
The encoding nests five $\mathrm{Acceptance}$ dispatches:
iterate $X$ from $\alpha$, join on $Y$ between $\alpha$
and $\beta$, confirm $\mathrm{Nil}$ after $Y$ in $\alpha$,
iterate $Z$ from $\beta$'s matched row, and confirm
$\mathrm{Nil}$ after $Z$.
We present each level as a named sub-scope.

\paragraph{Level~1: iterate $X$ from $\alpha$}
(Executable scope: \texttt{\_XAcceptance}.)
{\small
  \[
    \mathrm{xAcc} \mapsto \left\{
      \begin{aligned}
        &\alpha.\mathrm{Acceptance},\; \mathrm{VisitorMap} \mapsto \{\alpha.\mathrm{WithVisiteeTail}\},\\
        &\mathrm{Visitor} \mapsto \left\{
          \begin{aligned}
            &\_X \mapsto \{\mathrm{VisiteeTail}\},\; \ldots,\\
            &\mathrm{Visited} \mapsto \{\mathrm{yAcc}.\mathrm{Accepted}\}
        \end{aligned}\right\}
    \end{aligned}\right\}
\]}

\paragraph{Level~2: join $\alpha$'s $Y$-column with $\beta$'s first column}

(Executable scope: \texttt{\_YAcceptance}.)

The $\mathrm{VisitorMap}$ inherits from \emph{both}
$\beta.\mathrm{WithVisitorTail}$ and
$\_X.\mathrm{WithVisiteeTail}$, producing an entry for
constant $d$ only if $d$ appears in both $\alpha$'s
second column (via $\_X$) and $\beta$'s first column.
{\small
  \[
    \mathrm{yAcc} \mapsto \left\{
      \begin{aligned}
        &\_X.\mathrm{Acceptance},\\
        &\mathrm{VisitorMap} \mapsto \{\beta.\mathrm{WithVisitorTail},\; \_X.\mathrm{WithVisiteeTail}\},\\
        &\mathrm{Visitor} \mapsto \left\{
          \begin{aligned}
            &\_Y_0 \mapsto \{\mathrm{VisiteeTail}\},\; \mathrm{VisitorTail} \mapsto \{\beta.\mathrm{Tail}\},\\
            &\ldots,\\
            &\mathrm{Visited} \mapsto \{\mathrm{nilAcc}_0.\mathrm{Accepted}\}
        \end{aligned}\right\}
    \end{aligned}\right\}
\]}

\paragraph{Level~3: confirm $\mathrm{Nil}$ after $Y$ in $\alpha$}

(Executable scope: \texttt{\_NilAcceptance0}.)

This checks that $\alpha(X, Y)$ is a complete two-column
fact (its trie path terminates at $\mathrm{Nil}$ after $Y$).
Inside the $\mathrm{Nil}$ branch, the encoding proceeds to
iterate $\beta$'s second column via $\mathrm{VisitorTail}$:
{\small
  \[
    \mathrm{nilAcc}_0 \mapsto \left\{
      \begin{aligned}
        &\_Y_0.\mathrm{Acceptance},\\
        &\mathrm{VisitorMap} \mapsto \left\{\mathrm{Nil} \mapsto \left\{
            \begin{aligned}
              &\ldots,\\
              &\mathrm{Visited} \mapsto \{\mathrm{zAcc}.\mathrm{Accepted}\}
        \end{aligned}\right\}\right\}
    \end{aligned}\right\}
\]}

\paragraph{Level~4: iterate $Z$ from $\beta$'s matched row}
(Executable scope: \texttt{\_ZAcceptance}.)
{\small
  \[
    \mathrm{zAcc} \mapsto \left\{
      \begin{aligned}
        &\mathrm{VisitorTail}.\mathrm{Acceptance},\; \mathrm{VisitorMap} \mapsto \{\mathrm{VisitorTail}.\mathrm{WithVisiteeTail}\},\\
        &\mathrm{Visitor} \mapsto \left\{
          \begin{aligned}
            &\_Z \mapsto \{\mathrm{VisiteeTail}\},\; \ldots,\\
            &\mathrm{Visited} \mapsto \{\mathrm{nilAcc}_1.\mathrm{Accepted}\}
        \end{aligned}\right\}
    \end{aligned}\right\}
\]}

\paragraph{Level~5: confirm $\mathrm{Nil}$ after $Z$ and construct output}
(Executable scope: \texttt{\_NilAcceptance1}, with
\texttt{\_ZReplacement} and \texttt{\_XReplacement}.)
{\small
  \[
    \mathrm{nilAcc}_1 \mapsto \left\{
      \begin{aligned}
        &\mathrm{VisiteeTail}.\mathrm{Acceptance},\\
        &\mathrm{VisitorMap} \mapsto \left\{\mathrm{Nil} \mapsto \left\{
            \begin{aligned}
              &\mathrm{zR} \mapsto \{\_Z.\mathrm{Replacement},\; \mathrm{NewTail} \mapsto \{\mathrm{Nil}\}\},\\
              &\mathrm{xR} \mapsto \{\_X.\mathrm{Replacement},\; \mathrm{NewTail} \mapsto \{\mathrm{zR}.\mathrm{Replaced}\}\},\\
              &\mathrm{Visited} \mapsto \{\mathrm{xR}.\mathrm{Replaced}\}
        \end{aligned}\right\}\right\}
    \end{aligned}\right\}
\]}

\paragraph{Final result}
\[
  \mathrm{composed} \mapsto \{\mathrm{xAcc}.\mathrm{Accepted}\}
\]

\paragraph{Semantics}
At Level~2, the combined $\mathrm{VisitorMap}$ ensures that a
constant $d$ produces a $\mathrm{Visitor}$ only when $d$
appears in both $\alpha$'s second column \emph{and}
$\beta$'s first column.
The $\mathrm{Visitor}$ for each matched constant $d$
independently receives $\mathrm{VisiteeTail}$ ($\alpha$'s
data for~$d$) and $\mathrm{VisitorTail}$ ($\beta$'s row
for~$d$).
Output construction (Level~5) operates inside this
per-constant $\mathrm{Visitor}$, so $\mathrm{Replacement}$
acts on the \emph{individual} constant's data, not on the
merged scope.
Each matched constant $d$ produces its own output tuples
in $\mathrm{Visited}$, and these are collected by
$\mathrm{Accepted}$ via trie union ($\supers$).
This yields standard Datalog tuple-level semantics: only
constants that actually match between both relations
contribute to the output.
For the concrete instance
$\mathrm{edge} = \{(a,b),\; (b,c),\; (c,a)\}$, as encoded in
\anon[\nolinkurl{RelationalTwoHopRepeatedVisitor.mixin.yaml}~(supplementary material)]{\nolinkurl{RelationalTwoHopRepeatedVisitor.mixin.yaml}~\cite{mixin2025}},
Level~1 selects $a$, Level~2 matches $b$, Level~3 confirms
that $\mathrm{edge}(a,b)$ ends in $\mathrm{Nil}$, Level~4
iterates $\mathrm{edge}.\mathrm{TailMap}.b$ and selects $c$,
and Level~5 confirms $\mathrm{Nil}$ again before the two
$\mathrm{Replacement}$ operations reconstruct $(a,c)$.

\subsection{Constant Match and Navigation}
\label{sec:encoding-constants}

Constants in Datalog atoms use two mechanisms, neither of
which is a join.
They are simpler than variable joins because no value from a
second relation needs to be matched against the currently chosen
constant:

\paragraph{Type existence match}
An atom $R(X, b)$ with a constant in the second position
requires checking that constant $b$ exists in $R$'s
second column at the matched row.
This is encoded as an $\mathrm{Acceptance}$ with a
hand-crafted $\mathrm{VisitorMap}$ containing \emph{only}
the $b$ branch (instead of composing from
  $\mathrm{WithVisitorTail}$, which would create branches for
all constants):
{\small
  \begin{align*}
    \mathrm{VisitorMap} \mapsto \{b \mapsto \{\ldots\}\}
\end{align*}}
If $b$ exists in the dispatching scope, the
$b$ branch fires; otherwise, $\mathrm{Accepted}$
receives no matching branch.

\paragraph{Direct navigation}
An atom $R(b, Z)$ with a constant in the first position
navigates directly to $b$'s sub-trie:
\[
  \mathrm{edgeB}_Z \mapsto \{R.\mathrm{TailMap}.b\}
\]
This bypasses the visitor mechanism entirely, accessing
the per-constant $\mathrm{TailMap}$ entry to restrict the
scope to only $b$'s data.

\paragraph{Example}
$\mathrm{viaB}(X, Z) \mathrel{{:}\!{-}} \mathrm{edge}(X, b),\; \mathrm{edge}(b, Z)$
combines both mechanisms.
The executable MIXINv2 encoding is
\anon[\nolinkurl{RelationalConstant.mixin.yaml}~(supplementary material)]{\nolinkurl{RelationalConstant.mixin.yaml}~\cite{mixin2025}}.
The encoding iterates $\mathrm{edge}$ for $X$ (Level~1),
then checks for constant $b$ in $X$'s row with a
hand-crafted $\mathrm{VisitorMap}$ (Level~2), navigates
$\mathrm{edge}.\mathrm{TailMap}.b$ for $Z$ (Level~3),
confirms $\mathrm{Nil}$ (Level~4), and constructs $(X, Z)$
via $\mathrm{Replacement}$.

\subsection{Multi-Body Rules}
\label{sec:encoding-multi}

Rules with $n > 2$ body atoms chain additional join levels.
Each shared variable introduces one combined
$\mathrm{WithVisitorTail}$/$\mathrm{WithVisiteeTail}$ join
inside the previous level's $\mathrm{Visitor}$,
and the output construction is placed inside the innermost
$\mathrm{Nil}$ confirmation.
So the binary-rule encoding is the whole story structurally:
larger bodies only add more nested copies of the same pattern.

\paragraph{Three-body rule}
$\mathrm{chain}(X, W) \mathrel{{:}\!{-}} \alpha(X, Y),\; \beta(Y, Z),\; \gamma(Z, W)$
requires two joins: on $Y$ (between $\alpha$ and $\beta$)
and on $Z$ (between $\beta$ and $\gamma$).
The structure is: iterate $X$ from $\alpha$, join on $Y$
(getting both $\alpha$'s and $\beta$'s data), confirm
$\mathrm{Nil}$ after $Y$ in $\alpha$, then inside the
$\mathrm{Nil}$ branch join $\beta$'s $Z$-column with
$\gamma$'s first column, iterate $W$ from $\gamma$'s
matched row, confirm $\mathrm{Nil}$, and construct
$(X, W)$.
The $Z$-join is nested inside the $Y$-join's
$\mathrm{Nil}$ confirmation, ensuring that $\beta$'s data
for the matched $Y$ constant (available via
$\mathrm{VisitorTail}$) is used to iterate $\beta$'s
second column.

\paragraph{Mixed constants and variables}
$\mathrm{constJoin}(X, W) \mathrel{{:}\!{-}} \mathrm{edge}(X, Y),\; \mathrm{edge}(Y, b),\; \mathrm{edge}(b, W)$
combines a variable join on $Y$, a constant match on $b$
(hand-crafted $\mathrm{VisitorMap}$ with only the $b$ entry),
and direct navigation via $\mathrm{edge}.\mathrm{TailMap}.b$
for $W$.
The corresponding executable MIXINv2 file is
\anon[\nolinkurl{RelationalConstantJoin.mixin.yaml}~(supplementary material)]{\nolinkurl{RelationalConstantJoin.mixin.yaml}~\cite{mixin2025}}.

\subsection{Multi-Variable Join}

When a rule shares \emph{multiple} variables between two
body atoms, each shared variable requires its own
$\mathrm{Acceptance}$ check.
The checks are chained: the outer $\mathrm{Acceptance}$
joins on one column pair using combined
$\mathrm{WithVisitorTail}$/$\mathrm{WithVisiteeTail}$, and
its $\mathrm{Visitor}$ contains an inner $\mathrm{Acceptance}$
joining on another column pair.

\paragraph{Example}
$\mathrm{multiJoin}(X, Y) \mathrel{{:}\!{-}} r(X, Y),\; s(X, Y)$
shares both $X$ and $Y$.
The executable MIXINv2 encoding is
\anon[\nolinkurl{RelationalMultiVariableJoin.mixin.yaml}~(supplementary material)]{\nolinkurl{RelationalMultiVariableJoin.mixin.yaml}~\cite{mixin2025}}.
The outer $\mathrm{Acceptance}$ joins $r$'s and $s$'s first
columns (variable $X$) using combined
$\mathrm{WithVisitorTail}$/$\mathrm{WithVisiteeTail}$.
Inside the outer $\mathrm{Visitor}$, an inner
$\mathrm{Acceptance}$ joins $r$'s and $s$'s second columns
(variable $Y$) similarly.
Both joins must match for the output to be produced, and
the per-constant $\mathrm{Visited}$ pattern ensures that
only constants present in both relations contribute.

\subsection{Recursive Rules}
\label{sec:recursive-rules}

Recursive Datalog rules, such as transitive closure
$\mathrm{path}(X, Y) \mathrel{{:}\!{-}} \mathrm{edge}(X, Z),\; \mathrm{path}(Z, Y)$,
are encoded identically to non-recursive rules: the
head relation $\mathrm{path}$ appears both as a defined scope
and as a reference in the body (providing
$\mathrm{WithVisitorTail}$ for the $Z$-join).
This should not be read as ordinary $\this$ of a
single-valued lazy program.
Operationally, tabled evaluation
(Section~\ref{sec:mixin-trees},
Appendix~\ref{app:well-definedness})
computes $\mathrm{lfp}(T_P)$: each query into
$\mathrm{path}$ triggers further evaluation through
the recursive body, and the trie union ($\supers$)
accumulates all reachable tuples.
This corresponds exactly to \emph{Naive
evaluation}~\cite{bancilhon1986-recursive-query-strategies} of the
$T_P$
operator~\cite{vanemden1976-predicate-logic-semantics}:
each iteration applies every rule to the current
fact set and adds newly derived facts.
Because Datalog-encoded relations have finitely many
constants per query, the per-query answer domain is
finite, and convergence is guaranteed by the ascending
chain condition on a finite lattice.
\emph{Semi-Naive} evaluation~\cite{bancilhon1986-recursive-query-strategies},
which restricts each iteration to rules that use at
least one newly derived fact, is a natural
optimisation left for future work.

\paragraph{The recursive rule as monadic pseudocode}
The recursive rule above roughly corresponds to the
following pseudocode (Haskell-like syntax with
  \texttt{RebindableSyntax} and \texttt{OverloadedLists}
  semantics, where do-notation desugars to a
\texttt{Set}-based bind and list literals denote sets):
\begin{lstlisting}[style=haskell]
{-# LANGUAGE RebindableSyntax, OverloadedLists #-}
s >>= f = unions [f x | x <- s]
return x = [x]
guard True  = [()]
guard False = []
edge = [("a","b"), ("b","c"), ("c","a")]
path = edge <> do
  (x, y0) <- edge
  (y1, z) <- path
  guard (y0 == y1)
  return (x, z)
\end{lstlisting}
\noindent
Each line of the pseudocode has a direct counterpart in
the inheritance-calculus encoding
(Section~\ref{sec:encoding-variables}):
\begin{description}
  \item[\texttt{edge}] corresponds to the base-fact trie;
  \item[\texttt{path = edge <> do \ldots}]
    corresponds to $\mathrm{path}$ inheriting $\mathrm{edge}$
    (base facts) and the recursive body (derived facts via
    $\supers$);
  \item[\texttt{(x, y0) <- edge}]
    corresponds to Level~1 ($\mathrm{xAcc}$), iterating $X$
    from $\mathrm{edge}$;
  \item[\texttt{(y1, z) <- path}]
    corresponds to Level~2 ($\mathrm{yAcc}$), joining on $Y$
    with $\mathrm{path}$, and Level~4 ($\mathrm{zAcc}$),
    iterating $Z$;
  \item[\texttt{guard (y0 == y1)}]
    corresponds to the $\mathrm{WithVisitorTail} \cap
    \mathrm{WithVisiteeTail}$ intersection in Level~2: a
    constant produces a $\mathrm{Visitor}$ entry only if it
    appears in both relations;
  \item[\texttt{return (x, z)}]
    corresponds to the $\mathrm{Replacement}$ chain in
    Level~5, constructing the output tuple $(X, Z)$.
\end{description}
The only difference is the evaluation strategy:
Haskell evaluates \texttt{path} lazily (diverges),
whereas inheritance-calculus evaluates it via tabled
$T_P$ iteration (converges to $\mathrm{lfp}$).

However, this monadic program does not converge because
\texttt{path} is defined recursively.%
\footnote{
  Enabling the \texttt{RecursiveDo} extension does not help:
  \texttt{mfix} ties a knot at the \emph{element} level via laziness,
  whereas the transitive closure equation requires a
  fixed point at the \emph{container} level, that is, finding a
  set~$S$ such that $S = F(S)$.
}
Inheritance-calculus solves this through tabled evaluation
(Section~\ref{sec:mixin-trees}), which computes the
least fixed point of the $T_P$ operator by iterating
over the lattice of mixin trees, precisely the
container-level fixed point that the pseudocode above cannot express.

\subsection{Semantic Correspondence with Standard Datalog}
\label{sec:datalog-semantics}

The encoding produces standard Datalog tuple-level semantics.
Putting the previous subsections together: tries provide indexed
tuple storage, visitor dispatch supplies the control skeleton,
and per-constant $\mathrm{Replacement}$ keeps output
construction local to each witness.
The key mechanism is the per-constant $\mathrm{Visited}$
pattern: output construction ($\mathrm{Replacement}$)
occurs inside each matched constant's $\mathrm{Visitor}$
scope, operating on the \emph{individual} constant's data
rather than on the merged scope.
The $\mathrm{Accepted}$ property then collects all per-constant
$\mathrm{Visited}$ results via trie union ($\supers$).

For example, given $\mathrm{edge} = \{(a,b),\; (b,c),\; (c,a)\}$
(a~3-cycle), the two-hop self-join
$\mathrm{twoHop}(X, Z) \mathrel{{:}\!{-}} \mathrm{edge}(X, Y),\; \mathrm{edge}(Y, Z)$
produces exactly $3$ pairs:
$\{(a,c),\; (b,a),\; (c,b)\}$, matching standard Datalog.
This is the behavior exhibited by
\anon[\nolinkurl{RelationalTwoHopRepeatedVisitor.mixin.yaml}~(supplementary material)]{\nolinkurl{RelationalTwoHopRepeatedVisitor.mixin.yaml}~\cite{mixin2025}}.
When the $Y$-join matches constant $b$ (present in both
  $\mathrm{edge}$'s second column from $(a,b)$ and
$\mathrm{edge}$'s first column), the $\mathrm{Visitor}$ for
$b$ receives $\mathrm{VisitorTail} = \mathrm{edge}.\mathrm{TailMap}.b$
(containing only~$c$), and $\mathrm{Replacement}$ constructs
a single output tuple $(a, c)$, rather than all of $\{a,b,c\}$.

This value-level precision arises because each constant's
$\mathrm{Replacement}$ operates on $\mathrm{TailMap}.d$
(the per-constant sub-trie) rather than on the merged
trie.
The same point is visible in
the repeated-visitor self-join regression example\anon[~(supplementary material)]{~\cite{mixin2025}},
where repeated visitor dispatch prevents witnesses from
different constants from being merged spuriously.
Although $\bases$ (equation~\ref{eq:bases}) still computes
cross-joins at the primitive level, the encoding routes
computation through per-constant $\mathrm{TailMap}$ entries
before any merging occurs, achieving the same effect as
standard Datalog's tuple-level evaluation.

\section{Evaluation Trace of Multi-Target \texorpdfstring{$\this$}{this} Resolution}
\label{app:evaluation-trace}

This appendix gives a step-by-step evaluation trace of the program of
Appendix~\ref{app:scala-multi-target}, computing
$\properties((\text{``HasMultipleOuters''},\, \text{``outer''}))$ by the six
functions of Section~\ref{sec:definitions}
(equations~\ref{eq:properties} to~\ref{eq:this}).
Every call to $\overrides$, $\bases$, $\resolve$, $\this$, and $\supers$
is shown with its arguments and returned set.
The trace lives in the metalanguage: every argument and result is a
\emph{path}, that is, a sequence of labels
$(\ell_1, \ldots, \ell_n)$ identifying a position in the tree
(Section~\ref{sec:definitions}), not an inheritance-calculus expression.
To keep labels distinct from the metalanguage variables of
Section~\ref{sec:definitions} (the italic $p$, $\ell$, $n$, $S$), we
quote each concrete label, so a path is written as a parenthesized,
comma-separated sequence of quoted labels.
Thus $(\text{``HasMultipleOuters''}, \text{``outer''})$ is the path that
the inheritance-calculus projection
$\mathrm{HasMultipleOuters}.\mathrm{outer}$ identifies, and the root path
is the empty sequence $()$.

\paragraph{Program}
\begin{align*}
  \{ \quad \mathrm{MyOuter} &\mapsto \{
      \mathrm{MyInner} \mapsto \{
    \mathrm{outer} \mapsto \{\qualifiedthis{\mathrm{MyOuter}}\}\}\}, \\
    \mathrm{Object1} &\mapsto \{\mathrm{MyOuter}\}, \\
    \mathrm{Object2} &\mapsto \{\mathrm{MyOuter}\}, \\
    \mathrm{HasMultipleOuters} &\mapsto \{
      \mathrm{Object1}.\mathrm{MyInner},\;
    \mathrm{Object2}.\mathrm{MyInner}\}
  \quad \}
\end{align*}

\paragraph{AST primitives}
Parsing populates $\defines$ and $\inherits$
(Section~\ref{sec:definitions}):
\begin{align*}
  \defines(()) &= \{
    \text{``MyOuter''},\;
    \text{``Object1''},\;
    \text{``Object2''},\;
  \text{``HasMultipleOuters''}\}, \\
  \defines((\text{``MyOuter''})) &= \{\text{``MyInner''}\}, \\
  \defines((\text{``MyOuter''}, \text{``MyInner''})) &= \{\text{``outer''}\}, \\
  \inherits((\text{``MyOuter''}, \text{``MyInner''}, \text{``outer''})) &=
  \{(1,\, ())\}, \\
  \inherits((\text{``Object1''})) = \inherits((\text{``Object2''})) &=
  \{(0,\, (\text{``MyOuter''}))\}, \\
  \inherits((\text{``HasMultipleOuters''})) &=
  \{(0,\, (\text{``Object1''}, \text{``MyInner''})),\;
  (0,\, (\text{``Object2''}, \text{``MyInner''}))\}.
\end{align*}
All other $\defines$ and $\inherits$ are empty.
Each $\inherits$ entry is a reference pair $(n,\, w)$,
where $n$ is the number of upward steps and $w$ the
downward projection list (Section~\ref{sec:definitions}).

\paragraph{Trace}
We present the trace as a numbered list of semantic-function calls in
dependency order: a call appears after the calls it depends on, and each
item states, in English, which earlier items it uses, how it computes, and
what it yields. Each distinct call is listed once, so a single item may be
used by several later items, the reuse that an implementation could realize
by tabling~\cite{tamaki1986-tabled-resolution}.
The arguments and results are metalanguage values: a label is a quoted
string such as $\text{``outer''}$, and a path is a tuple of labels such as
$(\text{``HasMultipleOuters''}, \text{``outer''})$, with the root path
written $()$.
The query is
$\properties((\text{``HasMultipleOuters''}, \text{``outer''}))$.

\input{generated-evaluation-trace.tex}

\begin{remark}
  The crux is the $\this$ step (item~\ref{trace:this-Has-MyOuter-MyInner-1}).
  The inheritance reference at
  $(\text{``MyOuter''}, \text{``MyInner''}, \text{``outer''})$ takes one step
  upward, so $\this$ moves one step up from
  $(\text{``MyOuter''}, \text{``MyInner''})$ and yields the frontier
  $\{(\text{``Object1''}), (\text{``Object2''})\}$, the two inheritance
  sites that incorporated $\text{``MyInner''}$. This is precisely the
  multi-target situation that Scala~3 and the NixOS module system reject
  (Appendix~\ref{app:scala-multi-target}): there $\this$ would have to
  select a single site. Reaching the properties of $(\text{``MyOuter''})$
  then happens because both $(\text{``Object1''})$ and
  $(\text{``Object2''})$ have $(\text{``MyOuter''})$ as a base, so the final
  $\supers$ contains $(\text{``MyOuter''})$ through both routes; since
  attribute resolution aggregates by set union, the two routes collapse to
  the single answer
  $\properties((\text{``HasMultipleOuters''}, \text{``outer''}))
  = \{\text{``MyInner''}\} = \properties((\text{``MyOuter''}))$.
\end{remark}

\fi % end \ifInheritanceAppendix

% Appendix-only build (supplement.tex): no main matter ran above, so emit the
% bibliography here, after the appendix.
\ifInheritanceBody\else
\bibliographystyle{ACM-Reference-Format}
\bibliography{references}
\fi

% \clearpage before closing CJK* ships every remaining page and float while the
% CJK active characters still mean CJK (the running head and any deferred float
% are typeset in the output routine, after \end{CJK*} would otherwise restore the
% default UTF-8 meaning).
\ifInheritanceChinese\clearpage\end{CJK*}\fi
\end{document}
.  The
% defaults include everything, so the full build (preprint.tex) sets
% nothing; submission.tex clears the appendices (POPL 2027 requires them
% as separate supplemental material), supplement.tex clears the body.
% The single bibliography is emitted exactly once, after the main matter and
% before the appendix, since both cite it. In the appendix-only build (no main
% matter) it instead follows the appendix; see the two guarded copies below.
\makeatletter
\@ifundefined{ifInheritanceBody}{\newif\ifInheritanceBody\InheritanceBodytrue}{}
\@ifundefined{ifInheritanceAppendix}{\newif\ifInheritanceAppendix\InheritanceAppendixtrue}{}
% Language switch for the reader-facing prose. Defaults to false (English), so the
% existing pdfLaTeX entries are unaffected; preprint-zh.tex sets it true and adds
% CJK support, still under pdfLaTeX. See the paper-writing skill (Bilingual
% Translation) for the convention.
\@ifundefined{ifInheritanceChinese}{\newif\ifInheritanceChinese\InheritanceChinesefalse}{}
% acmart writes \newlabel{tocindent<n>}{<dimen>} into the .aux to align the
% table of contents. Those bare-dimen labels break hyperref's \newlabel when
% another build imports this .aux via xr-hyper; the importing entries
% (submission.tex, supplement.tex) drop them at import time through
% xr-tocindent-filter.tex, so nothing is suppressed here at the source.
\makeatother

% --- bilingual reader-facing prose ----------------------------------
% The reader-facing body exists in two languages in this one file, selected by
% \ifInheritanceChinese: English (default) and a Chinese translation. Both build
% with pdfLaTeX, so the math renders identically. Block prose is guarded by
% \ifInheritanceChinese <chinese> \else <english> \fi; inline switches (title,
% captions, theorem notes interleaved with math) use \bilingual{<en>}{<zh>}.
% Author-side notes stay Chinese % comments in both builds.
\newcommand{\bilingual}[2]{#1}
\ifInheritanceChinese
  % CJKutf8 + the Arphic gbsn font render Chinese under pdfLaTeX (see
  % modules/texlive.nix). The document is wrapped in a CJK* environment opened
  % after \begin{document}; ASCII passes through, so English and math render as in
  % the English build.
  \usepackage{CJKutf8}
  \setlength{\emergencystretch}{2em}
  % Display the Chinese, but feed the English to hyperref's PDF strings, which
  % cannot encode the CJK active characters.
  \renewcommand{\bilingual}[2]{\texorpdfstring{#2}{#1}}
\fi

\begin{document}
% Open CJK* before \title so the Chinese in the title, abstract, and body is read
% inside it; \proofname and the theorem-environment names are localized here
% (their bytes must sit inside a CJK environment, and translating \thmname's
% argument leaves the shared counter and numbering identical to English). Closed
% before \end{document}.
\ifInheritanceChinese
  \begin{CJK*}{UTF8}{gbsn}
  \renewcommand{\proofname}{证明}
  \makeatletter
  \newcommand{\inheritanceThmName}[1]{\@ifundefined{inheritance@thm@#1}{#1}{\csname inheritance@thm@#1\endcsname}}
  \renewcommand{\thmname}[1]{\inheritanceThmName{#1}}
  \@namedef{inheritance@thm@Theorem}{定理}
  \@namedef{inheritance@thm@Lemma}{引理}
  \@namedef{inheritance@thm@Corollary}{推论}
  \@namedef{inheritance@thm@Proposition}{命题}
  \@namedef{inheritance@thm@Conjecture}{猜想}
  \@namedef{inheritance@thm@Definition}{定义}
  \@namedef{inheritance@thm@Example}{例}
  \@namedef{inheritance@thm@Remark}{注}
  \@namedef{inheritance@thm@Notation}{记法}
  \makeatother
\fi

% The optional short title feeds the running head and the PDF metadata; keep it
% ASCII English in both builds, since those are typeset outside CJK*.
\title[A Calculus of Inheritance]{\bilingual{A Calculus of Inheritance}{一种继承的演算}}
% The appendix-only build (supplement.tex) is distributed as a standalone
% artifact; mark it so it is not mistaken for the main paper, which carries
% no subtitle.
\ifInheritanceBody\else
  \subtitle{Supplementary Material}
\fi

\author{Bo Yang}
\affiliation{%
  \institution{Figure AI Inc.}
  \city{San Jose}
  \state{California}
  \country{USA}
}
\email{yang-bo@yang-bo.com}
\thanks{This work was conducted independently prior to the author's employment at Figure AI.}

% The abstract, CCS concepts, and keywords are main-matter front matter: they
% belong to the paper, not to the standalone appendix PDF (supplement.tex),
% which sets \ifInheritanceBody false. Without this guard \maketitle typesets
% them in the appendix-only build.
\ifInheritanceBody
\begin{abstract}
  Just as the $\lambda$-calculus uses three primitives (abstraction, application, variable) as the foundation of functional programming, inheritance-calculus uses three primitives (record, definition, inheritance) as the foundation of declarative programming. A record is a set of definitions; by unifying modules, classes, objects, methods, fields, and locals under this single record abstraction, the calculus models inheritance simply as set union of definitions, two definitions of the same name merging recursively rather than one overriding the other. A merge is always meaningful because every value is a record and observation is purely structural: the merged record simply answers queries with the union of what its parts provide, so there is no scalar leaf at which two values could conflict. Consequently, composition is inherently commutative, idempotent, and associative, structurally eliminating the multiple-inheritance linearization problem. Its semantics is first-order, denotational, and computable by tabling (memoized fixpoint iteration as in tabled logic programming), even for cyclic inheritance hierarchies. This first-order, denotational, tabled semantics extends to the $\lambda$-calculus: $\lambda$-terms translate into a sublanguage of a lazy approximation of inheritance-calculus, the approximation that treats cyclic self-reference as divergent instead of iterating it to a fixpoint, and the translation is fully abstract, equating exactly the $\lambda$-terms with equal B\"ohm trees.

  Inheritance-calculus is distilled from MIXINv2, a practical implementation in which the same code can be invoked under different calling conventions (function colors such as synchronous and asynchronous); ordinary arithmetic, encoded within the calculus itself, behaves as multi-directional relations in the sense of logic programming; $\mathtt{this}$ resolves to multiple enclosing targets at once rather than a single receiver; and programs are extensible in both dimensions of the Expression Problem, gaining new variants and new operations by adding code. This makes inheritance-calculus strictly more expressive than the $\lambda$-calculus in Felleisen's sense of macro-expressibility.
\end{abstract}

\begin{CCSXML}
  <ccs2012>
  <concept>
  <concept_id>10003752.10010124.10010131</concept_id>
  <concept_desc>Theory of computation~Program semantics</concept_desc>
  <concept_significance>500</concept_significance>
  </concept>
  <concept>
  <concept_id>10003752.10010124.10010125.10010128</concept_id>
  <concept_desc>Theory of computation~Object oriented constructs</concept_desc>
  <concept_significance>500</concept_significance>
  </concept>
  <concept>
  <concept_id>10003752.10010124.10010125.10010127</concept_id>
  <concept_desc>Theory of computation~Functional constructs</concept_desc>
  <concept_significance>300</concept_significance>
  </concept>
  <concept>
  <concept_id>10003752.10010124.10010125.10010129</concept_id>
  <concept_desc>Theory of computation~Program schemes</concept_desc>
  <concept_significance>100</concept_significance>
  </concept>
  <concept>
  <concept_id>10011007.10011006.10011039</concept_id>
  <concept_desc>Software and its engineering~Formal language definitions</concept_desc>
  <concept_significance>500</concept_significance>
  </concept>
  <concept>
  <concept_id>10011007.10011006.10011008.10011009.10011019</concept_id>
  <concept_desc>Software and its engineering~Extensible languages</concept_desc>
  <concept_significance>500</concept_significance>
  </concept>
  <concept>
  <concept_id>10011007.10011006.10011008.10011009.10011011</concept_id>
  <concept_desc>Software and its engineering~Object oriented languages</concept_desc>
  <concept_significance>100</concept_significance>
  </concept>
  </ccs2012>
\end{CCSXML}

\ccsdesc[500]{Theory of computation~Program semantics}
\ccsdesc[500]{Theory of computation~Object oriented constructs}
\ccsdesc[300]{Theory of computation~Functional constructs}
\ccsdesc[100]{Theory of computation~Program schemes}
\ccsdesc[500]{Software and its engineering~Formal language definitions}
\ccsdesc[500]{Software and its engineering~Extensible languages}
\ccsdesc[100]{Software and its engineering~Object oriented languages}

\keywords{declarative programming, inheritance, fixpoint semantics,
  self-referential records, expression problem,
  first-order semantics,
  \texorpdfstring{$\lambda$-calculus}{lambda-calculus},
\texorpdfstring{B\"ohm trees}{Bohm trees}, configuration languages}
\fi % end \ifInheritanceBody (abstract, CCS concepts, keywords are main-matter only)

\maketitle

\ifInheritanceBody
\section{\bilingual{Introduction}{引言}}
\label{sec:introduction}

\ifInheritanceChinese
诸如 Jsonnet~\cite{jsonnet}、Hydra~\cite{hydra2023}、CUE~\cite{cue2019}、Dhall~\cite{dhall2017}、Kustomize~\cite{kustomize} 与 JSON Patch~\cite{rfc6902-json-patch} 这样的配置语言,都提供了通过继承或合并来组合结构化数据的机制。

在这些之中,Nix 生态系统~\cite{dolstra2010-nixos} 格外突出,无论是就其规模还是其设计的走向而言。它的软件包集合 \texttt{nixpkgs} 收录了超过 113{,}000 个软件包,%
\footnote{
  据 Repology~\cite{repology2025},这接近 Debian 或 Ubuntu 软件包数量的三倍,而贡献者数量相近。并非每个软件包都用到这里所描述的可扩展性机制,但同样的告诫也适用于其他发行版。
}
是现存最大的软件包集合之一。软件包的组合最初依赖一条线性化的 mixin 链,其中每一层都能看见它的前一层,%
\footnote{
  Nix overlay 通过一条线性化的 \texttt{self}/\texttt{super} 链来组合软件包集合,遵循 Bracha--Cook 的 mixin 模式~\cite{bracha1990-mixin-inheritance}:\texttt{self} 晚绑定到不动点结果,而 \texttt{super} 指向前一层。该机制是不对称且依赖顺序的。
}
但该生态系统已日益汇聚到一种对称的替代方案上:NixOS 模块系统~\cite{dolstra2010-nixos,nixosmodules},它通过对嵌套的属性集进行递归合并、配以自引用求值与延迟模块~\cite{nixosmodules} 来组合,而这种组合是可交换、幂等且可结合的。模块系统最初为系统配置而设计,如今已被用于用户环境~\cite{homemanager}、macOS 配置~\cite{nixdarwin}、磁盘分区~\cite{disko}、flake 结构~\cite{flakeparts}、开发者环境~\cite{devenv}、Kubernetes 集群~\cite{kubenix,nixidy},以及多语言软件包构建~\cite{dream2nix},而最后这一项正是 \texttt{nixpkgs} 那套不对称机制原本的领域。Nix 是一门带头等函数的函数式语言;然而一份典型的模块配置几乎完全由嵌套的属性集构成,函数只是偶尔作为模块系统的 DSL 语法出现,例如 \texttt{mkIf} 与 \texttt{mkMerge}。没有任何计算理论刻画了这一机制所展现的可扩展性。
\else
Configuration languages such as Jsonnet~\cite{jsonnet}, Hydra~\cite{hydra2023},
CUE~\cite{cue2019}, Dhall~\cite{dhall2017},
Kustomize~\cite{kustomize}, and JSON
Patch~\cite{rfc6902-json-patch} all provide mechanisms for
composing structured data through inheritance or merging.

Among these, the Nix ecosystem~\cite{dolstra2010-nixos} stands out,
both for scale and for the trajectory of its design.
Its package collection, \texttt{nixpkgs}, contains over
113{,}000 packages,%
\footnote{
  According to Repology~\cite{repology2025},
  nearly three times the package count of Debian or Ubuntu,
  achieved with a similar number of contributors.
  Not every package exercises the extensibility mechanisms
  described here, but the same caveat applies to other distributions.
}
one of the largest in existence.
Package composition originally relied on a linearized mixin
chain in which each layer sees its predecessor,%
\footnote{
  Nix overlays compose package sets via a linearized
  \texttt{self}/\texttt{super} chain following the
  Bracha--Cook mixin pattern~\cite{bracha1990-mixin-inheritance}:
  \texttt{self} is late-bound to the fixed-point result
  and \texttt{super} refers to the previous layer.
  The mechanism is asymmetric and order-dependent.
}
but the ecosystem has increasingly converged on a symmetric
alternative: the NixOS module
system~\cite{dolstra2010-nixos,nixosmodules}, which composes by
recursively merging nested attribute sets with self-referential
evaluation and deferred modules~\cite{nixosmodules}---commutative,
idempotent,
and associative.
Originally designed for system configuration, the module system
has been adopted for user environments~\cite{homemanager},
macOS configuration~\cite{nixdarwin},
disk partitioning~\cite{disko},
flake structure~\cite{flakeparts},
developer environments~\cite{devenv},
Kubernetes clusters~\cite{kubenix,nixidy},
and multi-language package building~\cite{dream2nix}---the
original domain of \texttt{nixpkgs}'s asymmetric mechanism.
Nix is a functional language with first-class functions;
yet a typical module configuration consists almost entirely of
nested attribute sets, with functions appearing only
occasionally as DSL syntax of the module system
such as \texttt{mkIf} and \texttt{mkMerge}.
No computational theory captures the extensibility that this
mechanism exhibits.
\fi

The $\lambda$-calculus has no analogue for declarative
programming. This gap matters because the two paradigms differ in
fundamental ways, which raises a question:

\begin{quote}
  Can configuration alone be practical general-purpose programming?
\end{quote}

\noindent
To make the question precise, we operationalize it with four
\emph{opinionated} proxies, grouped under the question's two halves:
\emph{configuration} and \emph{practical general-purpose} programming.

\begin{itemize}
  \item Immutable and first-order%
    \footnote{
      Here \emph{first-order} holds in three senses at once:
      the semantic domain contains only data, no function spaces,
      closures, or
      strategies~\cite{vanemden1976-predicate-logic-semantics};
      the defining equations map data to
      data~\cite{reynolds1972-definitional-interpreters};
      and the semantics is a first-order inductive
      definition~\cite{aczel1977-inductive-definitions} whose
      quantifiers range only over data.
    }
    is our proxy for
    \emph{configuration}: a configuration is read, not a function that
    reads an argument.
  \item Turing-completeness together with extensibility is our proxy for
    \emph{practical general-purpose} programming: a language must be able
    to compute anything and to grow by adding code rather than rewriting
    it, the Expression
    Problem~\cite{wadler1998-expression-problem}.
\end{itemize}

\noindent
Each classical model satisfies some of these proxies but fails at least
one:

\begin{table}[h]
  \centering
  \small
  \begin{tabular}{lcccl}
    \textbf{Model} & \textbf{Immutable} & \textbf{First-order} & \textbf{Turing complete} & \textbf{Extensibility} \\
    \hline
    Turing Machine        & $\times$     & $\checkmark$ & $\checkmark$ & $\times$ \\
    $\lambda$-calculus    & $\checkmark$ & $\times$     & $\checkmark$ & opt-in \\
    RAM Machine           & $\times$     & $\checkmark$ & $\checkmark$ & $\times$ \\
    Prolog / Horn clauses & $\checkmark$ & $\checkmark$ & $\checkmark$ & disjunct-only \\
  \end{tabular}
  \caption{The question operationalized as four proxies: immutable and
    first-order (configuration) versus Turing-complete and extensible
  (practical general-purpose). No classical model satisfies all four.}
  \label{tab:computational-models}
\end{table}

\noindent
The Turing Machine and RAM Machine require mutable state and have no native
extension mechanism;
the $\lambda$-calculus is higher-order rather than first-order, and its
extensibility is opt-in: a function is extensible only where its author
anticipated extension, for example by threading open-recursion
parameters, while a closed function body cannot be extended without
rewriting it;
Prolog meets the other three proxies but is extensible only by adding
disjuncts (alternative clauses) to existing predicates, not by attaching
new fields or methods to existing values.
No known computational model satisfies all four proxies simultaneously.

The Nix ecosystem already suggests that such a model may exist:
its inheritance-based composition over recursive records, without
explicit functions, is expressive enough to configure entire
operating systems and manage one of the largest package
collections in existence.
Most configuration languages operate on finite, well-founded
structures (initial algebras in the sense of universal algebra).
Nix is an exception: its values are lazily observable
and possibly infinite, allowing any single package or
configuration option to be evaluated without materializing the
entire set of hundreds of thousands of packages and options.
Guided by this observation, we set out to reduce the
NixOS module system to a minimal set of primitives.
The reduction produced MIXIN, an early implementation in Nix,
and subsequently
MIXINv2\anon[~(included as supplementary material)]{~\cite{mixin2025}}, now implemented in Python: a declarative programming
language with only three constructs (record, definition,
and inheritance) and no functions or let-bindings.
MIXINv2 obtains scalars such as floating-point numbers through a
foreign-function interface; inheritance-calculus is what remains after
removing that interface, so names like $\mathrm{PI}$ and
$\mathrm{Float}$ in Figure~\ref{fig:cheatsheet} denote imported
records whose internals the examples never inspect: scalar
computation such as the figure's additions happens behind the
interface, while the calculus itself only ever observes structure;
the case study of Section~\ref{sec:case-study} instead encodes
naturals and their arithmetic purely in the calculus, with no FFI.
When two merged records define the same name with different contents,
the result is simply a record carrying both structures; the calculus
never compares or rejects them
(Section~\ref{sec:mixin-trees} makes the merge precise).
No such collision occurs in Figure~\ref{fig:cheatsheet}: each
instance defines each scalar field exactly once, and the type record
it merges with contributes structure, not a competing value.

\begin{figure*}
  \centering
  \begin{minipage}[t]{0.34\textwidth}
    \textbf{Python}\par
\begin{lstlisting}[style=python]
# complex_module.py
from builtins import *
from math import *
from operator import *
from dataclasses import *
@dataclass(kw_only=True)
class Complex:
    re: float
    im: float
    def plus(self, *,
             other: "Complex"):
        sumRe = add(
            self.re,
            other.re)
        sumIm = add(
            self.im,
            other.im)
        return Complex(
            re=sumRe,
            im=sumIm)
a = Complex(
    re=pi,
    im=e)
b = Complex(
    re=tau,
    im=pi)
total = a.plus(
    other=b)
\end{lstlisting}
  \end{minipage}
  \hfill
  \begin{minipage}[t]{0.65\textwidth}
    \textbf{Inheritance-calculus}\par\smallskip
    {\small\(
        \figureAnnotationAnchor
        \begin{aligned}
          &\mathrm{ComplexModule} \mapsto {} \figureAnnotation{module} \\
          &\left\{
            \begin{aligned}
              &\mathrm{BuiltIns}, \figureAnnotation{import} \\
              &\mathrm{Math}, \figureAnnotation{import} \\
              &\mathrm{Operator}, \figureAnnotation{import} \\
              &\mathrm{Complex} \mapsto {} \figureAnnotation{class} \\
              &\left\{
                \begin{aligned}
                  &\mathrm{re} \mapsto {} \figureAnnotation{field} \\
                  &\left\{ \qualifiedthis{\mathrm{ComplexModule}}.\mathrm{Float} \right\}, \figureAnnotation{imported type} \\
                  &\mathrm{im} \mapsto {} \figureAnnotation{field} \\
                  &\left\{ \qualifiedthis{\mathrm{ComplexModule}}.\mathrm{Float} \right\}, \figureAnnotation{imported type} \\
                  &\mathrm{Plus} \mapsto {} \figureAnnotation{method} \\
                  &\left\{
                    \begin{aligned}
                      &\mathrm{other} \mapsto \left\{ \mathrm{Complex} \right\}, \figureAnnotation{parameter} \\
                      &\mathrm{\_sumReCall} \mapsto {} \figureAnnotation{function application} \\
                      &\left\{
                        \begin{aligned}
                          &\qualifiedthis{\mathrm{ComplexModule}}.\mathrm{Add}, \figureAnnotation{imported function} \\
                          &\mathrm{operand0} \mapsto \left\{ \qualifiedthis{\mathrm{Complex}}.\mathrm{re} \right\}, \figureAnnotation{keyword argument} \\
                          &\mathrm{operand1} \mapsto \left\{ \qualifiedthis{\mathrm{Plus}}.\mathrm{other}.\mathrm{re} \right\} \figureAnnotation{keyword argument} \\
                      \end{aligned}\right\}, \\
                      &\mathrm{\_sumRe} \mapsto \left\{ \mathrm{\_sumReCall}.\mathrm{result} \right\}, \figureAnnotation{local} \\
                      &\mathrm{\_sumImCall} \mapsto {} \figureAnnotation{function application} \\
                      &\left\{
                        \begin{aligned}
                          &\qualifiedthis{\mathrm{ComplexModule}}.\mathrm{Add}, \figureAnnotation{imported function} \\
                          &\mathrm{operand0} \mapsto \left\{ \qualifiedthis{\mathrm{Complex}}.\mathrm{im} \right\}, \figureAnnotation{keyword argument} \\
                          &\mathrm{operand1} \mapsto \left\{ \qualifiedthis{\mathrm{Plus}}.\mathrm{other}.\mathrm{im} \right\} \figureAnnotation{keyword argument} \\
                      \end{aligned}\right\}, \\
                      &\mathrm{\_sumIm} \mapsto \left\{ \mathrm{\_sumImCall}.\mathrm{result} \right\}, \figureAnnotation{local} \\
                      &\mathrm{result} \mapsto {} \figureAnnotation{return} \\
                      &\left\{
                        \begin{aligned}
                          &\mathrm{Complex}, \figureAnnotation{new instance} \\
                          &\mathrm{re} \mapsto \left\{ \mathrm{\_sumRe} \right\}, \figureAnnotation{keyword argument} \\
                          &\mathrm{im} \mapsto \left\{ \mathrm{\_sumIm} \right\} \figureAnnotation{keyword argument} \\
                      \end{aligned}\right\} \\
                  \end{aligned}\right\} \\
              \end{aligned}\right\}, \\
              &\mathrm{a} \mapsto {} \figureAnnotation{module-level variable} \\
              &\left\{
                \begin{aligned}
                  &\mathrm{Complex}, \figureAnnotation{new instance} \\
                  &\mathrm{re} \mapsto \left\{ \qualifiedthis{\mathrm{ComplexModule}}.\mathrm{PI} \right\}, \figureAnnotation{imported constant} \\
                  &\mathrm{im} \mapsto \left\{ \qualifiedthis{\mathrm{ComplexModule}}.\mathrm{E} \right\} \figureAnnotation{imported constant} \\
              \end{aligned}\right\}, \\
              &\mathrm{b} \mapsto {} \figureAnnotation{module-level variable} \\
              &\left\{
                \begin{aligned}
                  &\mathrm{Complex}, \figureAnnotation{new instance} \\
                  &\mathrm{re} \mapsto \left\{ \qualifiedthis{\mathrm{ComplexModule}}.\mathrm{Tau} \right\}, \figureAnnotation{imported constant} \\
                  &\mathrm{im} \mapsto \left\{ \qualifiedthis{\mathrm{ComplexModule}}.\mathrm{PI} \right\} \figureAnnotation{imported constant} \\
              \end{aligned}\right\}, \\
              &\mathrm{\_totalCall} \mapsto {} \figureAnnotation{method application} \\
              &\left\{
                \begin{aligned}
                  &\mathrm{a}.\mathrm{Plus}, \figureAnnotation{method reference} \\
                  &\mathrm{other} \mapsto \left\{ \mathrm{b} \right\} \figureAnnotation{keyword argument} \\
              \end{aligned}\right\}, \\
              &\mathrm{total} \mapsto \left\{ \mathrm{\_totalCall}.\mathrm{result} \right\} \figureAnnotation{module-level variable} \\
          \end{aligned}\right\}
        \end{aligned}
    \)}
  \end{minipage}
  \caption{A first look at inheritance-calculus (right), beside Python
    (left): the same complex-number class with an addition method.
    An entry $name \mapsto \{\dots\}$ is a definition; a bare name
    inside braces (such as $\mathrm{BuiltIns}$ or $\mathrm{Complex}$) is
    an inheritance clause that copies that record's definitions into the
    enclosing record.
    $\mathrm{BuiltIns}$, $\mathrm{Math}$, and $\mathrm{Operator}$ are
    defined in an enclosing scope, elided here.
    A bare name resolves lexically to the innermost enclosing record
    that defines it directly; qualified $\this$ reaches members that
    arrive by inheritance, such as the imported $\mathrm{Float}$ and
    $\mathrm{Add}$.
    Types and values share the same record syntax: a field such as
    $\mathrm{re}$ records its type by inheriting the type's record,
    with no checking force in the untyped calculus.
    Both sides, typing their components with \texttt{builtins.float},
    drawing their constants from the \texttt{math} module, and adding with
    \texttt{operator.add}, compute the complex
    number $(\pi+\tau) + (e+\pi)i$.
    $\mathrm{\_sumReCall}$ and $\mathrm{\_sumImCall}$ are function
    applications, and $\mathrm{\_totalCall}$ a method application, each
    encoded as a command pattern: a record that inherits the callee and
    defines its arguments, the call's result read back from a
    $\mathrm{result}$ field that the callee defines.
    Idempotence only means that listing the same inheritance source
    twice in one record is a no-op; distinct records inheriting the
    same source, like the two call sites here or the two instances
    $\mathrm{a}$ and $\mathrm{b}$, remain distinct.}
  \label{fig:cheatsheet}
\end{figure*}



In developing real programs in MIXINv2, including a database-backed web server whose Python foreign-function interface only faithfully wraps individual host-library calls, we found that module, class, object, method, field, and local binding can all be expressed as the same construct: a record of definitions composed by inheritance. This paper distills this minimal computational model into \textbf{inheritance-calculus}, shown beside Python in Figure~\ref{fig:cheatsheet}.
The key to reading the figure is that a reference
$\qualifiedthis{X}.w$ behaves like a late-bound Java-style
\texttt{Outer.this}: it denotes the enclosing record labeled $X$ as
seen from wherever the reference itself ends up in the fully merged
tree.
Inheriting a record does not duplicate it; it makes the record's
definitions, references included, visible at the inheriting position
as well, and a reference resolves relative to the position from which
it is viewed.
The set-union picture and this late-binding picture are the same
mechanism seen globally and locally: the union says which definitions
a position carries, and late binding says how a definition behaves at
each position that carries it.
Qualified references therefore see members that arrive by merging,
such as the imported $\mathrm{Float}$, while bare names such as
$\mathrm{Complex}$ resolve lexically against directly written
definitions only.
This late binding is what makes inheritance act as function
application: the call record $\mathrm{\_totalCall}$ inherits
$\mathrm{Plus}$, so the merged tree presents
$\mathrm{Plus}$'s body inside $\mathrm{\_totalCall}$ as well, and
viewed from there the reference
$\qualifiedthis{\mathrm{Plus}}.\mathrm{other}$
denotes that copy's enclosing record, $\mathrm{\_totalCall}$ itself,
where the argument $\mathrm{other}$ is defined.
Receivers bind the same way: the dotted bare name
$\mathrm{a}.\mathrm{Plus}$ resolves its first label lexically and
projects the rest, denoting the copy of $\mathrm{Plus}$ inside the
instance $\mathrm{a}$, in which
$\qualifiedthis{\mathrm{Complex}}.\mathrm{re}$ denotes
$\mathrm{a}$'s own field.
Section~\ref{sec:syntax} gives its formal syntax, and
Section~\ref{sec:mixin-trees} expands this first look into the full
semantics.
The linearization problem of mixin-based
systems~\cite{bracha1990-mixin-inheritance,barrett1996-c3-linearization}
did not arise: inheritance is inherently commutative, idempotent,
and associative, and $\this$ naturally resolves to
\emph{multiple targets} rather than one: when inheritance places the
same definition inside several enclosing records, a reference through
$\this$ inherits from all of them at once, their contents merged,%
\footnote{
  Resolution is anchored at the position from which the reference is
  observed. When both $\mathrm{A}$ and $\mathrm{B}$ inherit
  $\mathrm{Base}$, a reference in $\mathrm{Base}$ observed inside
  $\mathrm{A}$ sees only $\mathrm{A}$; the two call sites of
  Figure~\ref{fig:cheatsheet} therefore stay independent. Multiple
  targets arise when a single observation position inherits the same
  definition through two routes, say $\mathrm{C}$ inheriting
  $\mathrm{Base}$ via both $\mathrm{A}$ and $\mathrm{B}$: observed at
  $\mathrm{C}$, a $\this$ reference in $\mathrm{Base}$ resolves to
  $\mathrm{A}$ and $\mathrm{B}$ at once, and projection through it
  yields the union of what they provide.
}
where
prior systems
such as Scala and the NixOS module system reject the multi-target
situation.
Section~\ref{sec:case-study} traces how boolean logic and
arithmetic, encoded purely in the calculus as separate files, extend in both
dimensions of the Expression
Problem~\cite{wadler1998-expression-problem} without any dedicated
mechanism, and how ordinary arithmetic yields the relational
semantics of logic
programming~\cite{vanemden1976-predicate-logic-semantics}.
Section~\ref{sec:emergent-phenomena} discusses these and further
emergent phenomena, including function color
blindness~\cite{nystrom2015-function-color}.
Together these sections answer the four proxies: inheritance-calculus
is immutable and first-order by construction
(Section~\ref{sec:definitions}), Turing-complete because the
$\lambda$-calculus embeds into it
(Section~\ref{sec:forward-translation}), and extensible in both
dimensions of the Expression Problem (Section~\ref{sec:case-study}).

\begin{figure*}
  \centering
  \begin{tikzpicture}[
      calc/.style={align=center, inner sep=2pt, font=\small},
      more/.style={-{Stealth[length=2mm]}, semithick},
      iso/.style={{Stealth[length=2mm]}-{Stealth[length=2mm]}, semithick},
      elbl/.style={font=\scriptsize\itshape, midway, fill=white, inner sep=1pt},
      x=1mm, y=1mm,
    ]
    % rows: bottom = lambda calculi, middle = lambda sublanguages, top = inheritance-calculi
    % columns: eager (x=0), lazy (x=42), fixpoint (x=84)
    \node[calc] (lam)  at (0,0)   {$\lambda$};
    \node[calc] (llam) at (42,0)  {lazy $\lambda$};
    \node[calc] (flam) at (84,0)  {fixpoint $\lambda$~\cite{yang2026-co-lambda}};

    \node[calc] (lsub) at (42,19)  {$\lambda$-lazy-IC (ours)};
    \node[calc] (sub)  at (84,19)  {$\lambda$-IC};

    \node[calc] (lic)  at (42,38)  {lazy IC (ours)};
    \node[calc] (ic)   at (84,38)  {IC (ours)};

    % horizontal edges (convergence): "more defined"
    \draw[more] (lam)  -- node[elbl,below]{more defined (non-strict)} (llam);
    \draw[more] (llam) -- node[elbl,below]{more defined (unproductive $\bot$)} (flam);
    \draw[more] (lsub) -- node[elbl,above]{more defined (powerset)} (sub);
    \draw[more] (lic)  -- node[elbl,above]{more defined (powerset)} (ic);

    % vertical edges (expressiveness): "strictly more expressive"
    \draw[more] (lsub) -- node[elbl,left]{strictly more expressive} (lic);
    \draw[more] (sub)  -- node[elbl,right]{strictly more expressive} (ic);

    % full abstraction: vertical, in the lazy column
    \draw[iso] (llam) -- node[elbl,left]{fully abstract} (lsub);
  \end{tikzpicture}
  \caption{The calculi of this paper; nodes marked
    (ours) are contributions of this paper.
    Abbreviations:
    $\lambda$ is the $\lambda$-calculus;
    the lazy prefix denotes the approximation that stops after the first
    iteration of the fixpoint of this paper's semantic equations
    (written $T_P\!\uparrow\!1$ in
    logic-programming notation): lookups that re-enter themselves
    through cyclic self-reference return divergence instead of being
    iterated further.
    Continuing:
    fixpoint $\lambda$ is the tabled evaluation of the
    $\lambda$-calculus of Yang~\cite{yang2026-co-lambda};
    IC the inheritance-calculus;
    and $\lambda$-IC and $\lambda$-lazy-IC the $\lambda$ sublanguages of
    IC and of lazy IC, respectively.
    Edge labels:
    \emph{more defined} means every term that converges in the source
    calculus also converges in the target calculus, but not conversely.
    The parenthesized qualifiers distinguish the mechanism:
    \emph{non-strict} more-definedness gains convergence by deeming
    abstractions values without evaluating their bodies;
    \emph{unproductive $\bot$} more-definedness gains convergence by
    deciding unproductive loops such as $\Omega$ as the definite answer
    $\bot$ in finite time, preserving the lazy tree semantics;
    and \emph{powerset} more-definedness gains convergence by
    resolving cyclic self-references to definite sets, as least fixed
    points over the powerset lattice.
    Continuing the edge labels:
    \emph{strictly more expressive} is meant in Felleisen's sense of
    macro-expressibility;
    and \emph{fully abstract} marks full abstraction.}
  \label{fig:calculi-matrix}
\end{figure*}

Figure~\ref{fig:calculi-matrix} situates inheritance-calculus among
several related calculi developed below.
Section~\ref{sec:definitions} defines the semantics via six
mutually recursive first-order functions on paths and sets of paths,
with no function spaces, closures, or higher-order constructs;
the functions are computable by
tabling~\cite{tamaki1986-tabled-resolution}.
Section~\ref{sec:forward-translation} shows that
these functions extend to the $\lambda$-calculus, embedding it.
Section~\ref{sec:asymmetry} establishes that the lazy
$\lambda$-calculus cannot macro-express the inheritance constructors,
so inheritance-calculus is strictly more expressive in
Felleisen's sense~\cite{felleisen1991-expressive-power}.

\section{\bilingual{Syntax}{语法}}
\label{sec:syntax}

\ifInheritanceChinese
设 $e$ 表示一个表达式,$s$ 表示一个元素列表,$c$ 表示一个元素,
$q$ 表示一个限定符,$w$ 表示一个投影列表,$\ell$ 表示一个充当
属性名的标签,$n$ 表示一个 de~Bruijn 索引。
空串 $\varepsilon$ 是空元素列表。
\else
Let $e$ denote an expression, $s$ an element list, $c$ an element,
$q$ a qualifier, $w$ a projection list, $\ell$ a label serving as a
property name, and $n$ a de~Bruijn index.
The empty string $\varepsilon$ is the empty element list.
\fi

\begin{align*}
  e \quad ::= \quad & \{\, s \,\}
  && \text{\bilingual{(record)}{(记录)}} \\[6pt]
  s \quad ::= \quad & \varepsilon
  && \text{\bilingual{(empty)}{(空)}} \\*
  \mid\quad & c
  && \text{\bilingual{(one element)}{(一个元素)}} \\*
  \mid\quad & c,\; s
  && \text{\bilingual{(element, then more)}{(一个元素,然后更多)}} \\[6pt]
  c \quad ::= \quad & \ell \mapsto e
  && \text{\bilingual{(definition)}{(定义)}} \\*
  \mid\quad & q.\, w
  && \text{\bilingual{(inheritance: qualifier and projections)}{(继承:限定符与投影)}} \\*
  \mid\quad & q
  && \text{\bilingual{(inheritance: qualifier only)}{(继承:仅限定符)}} \\*
  \mid\quad & w
  && \text{\bilingual{(inheritance: projections only)}{(继承:仅投影)}} \\[6pt]
  q \quad ::= \quad & \qualifiedthis{\ell_{\mathrm{qualifier}}}
  && \text{\bilingual{(named qualifier)}{(具名限定符)}} \\*
  \mid\quad & \dbi{n}
  && \text{\bilingual{(indexed qualifier)}{(索引限定符)}} \\[6pt]
  w \quad ::= \quad & \ell_{\mathrm{projection}}
  && \text{\bilingual{(one projection)}{(一个投影)}} \\*
  \mid\quad & \ell_{\mathrm{projection}}.\, w
  && \text{\bilingual{(projection, then more)}{(一个投影,然后更多)}}
\end{align*}

\ifInheritanceChinese
一个表达式将一个元素列表 $s$ 包裹在花括号中,而每个元素
$c$ 要么是一个定义 $\ell \mapsto e$,要么是一个继承源。
语法把 $s$ 写成一个递归列表,但该列表被读作一个
\emph{集合}:元素的顺序无关紧要,重复元素也没有任何作用,
因此两个仅在重排或重复上有所不同的元素列表表示
同一个表达式。由于表达式是集合,继承本身就是
可交换、幂等且可结合的。
图~\ref{fig:cheatsheet} 中的记录 $\mathrm{Complex} \mapsto \{ \mathrm{re} \mapsto \ldots,\;
\mathrm{im} \mapsto \ldots,\; \mathrm{Plus} \mapsto \ldots \}$
就是一个这样的多元素列表。

$\ell \mapsto e$ 中的 $\mapsto$ 定义了一个名为 $\ell$、其值为
$e$ 的属性;例如,图~\ref{fig:cheatsheet} 中的 $\mathrm{re} \mapsto \{\ldots\}$ 与
$\mathrm{Complex} \mapsto \{\ldots\}$。
MIXINv2 还通过一个
外部函数接口(FFI)提供标量类型;继承演算就是
去掉 FFI 之后所剩下的部分。
本文中没有任何例子用到 FFI 或任何标量值。

\paragraph{\bilingual{Inheritance}{继承}}
一个继承通过一个限定符 $q$ 引用一个作用域;该限定符
沿作用域层级向上导航,然后通过一个投影列表 $w$
向下投影属性。限定符 $q$ 容许两种等价的
记法:
\begin{description}
  \item[\bilingual{Named form}{具名形式}]
    $\qualifiedthis{\ell_{\mathrm{qualifier}}}.\, w$,
    类比于 Java 的 \texttt{Outer.this.field},其中 $w$ 是一个
    投影列表。
    每个 $\{\ldots\}$ 记录在它的外围作用域中都有一个标签;
    $\ell_{\mathrm{qualifier}}$ 命名目标外围作用域。
    例如,图~\ref{fig:cheatsheet} 中的
    $\qualifiedthis{\mathrm{ComplexModule}}.\mathrm{Float}$
    投影单个标签,而
    $\qualifiedthis{\mathrm{Plus}}.\mathrm{other}.\mathrm{re}$ 展示了一个
    多标签的投影列表 $w$。
  \item[\bilingual{Indexed form}{索引形式}]
    $\dbi{n}.\, w$,
    其中 $n \ge 0$ 是一个 de~Bruijn 索引~\cite{debruijn1972-nameless-dummies}:
    $n = 0$ 指向该引用的外围作用域,
    而 $n$ 每加一,就向外移动一个作用域层级。
\end{description}
两种形式都先取出目标作用域,即所有继承都应用完毕之后
那个完全继承后的记录,然后从中投影由投影列表 $w$
所命名的属性。
具名形式在解析期间脱糖为索引形式,
如第~\ref{sec:mixin-trees} 节所述。

每个 $\{\ldots\}$ 记录字面量都会创建一个新的作用域层级。%
\footnote{
  我们用\emph{记录}指那个语法构造,用\emph{作用域}指
  它所创建的名字解析语境,用\emph{mixin} 指它
  作为可组合单元的角色;在继承演算中这些是同一个实体的三种
  视角。
}
该作用域包含那个 mixin 的所有属性,包括那些
通过继承源继承而来的属性。
例如,$\{a \mapsto \{\},\; b \mapsto \{\}\}$ 是单个 mixin,其作用域
同时包含 $a$ 与 $b$;$\{r,\; a \mapsto \{\}\}$ 也是如此,其中 $r$
是对一个包含 $b$ 的记录的引用。
继承源不会创建额外的作用域层级。
在图~\ref{fig:cheatsheet} 中,绑定到 $\mathrm{Plus}$ 的记录
就是一个这样的作用域,它包含 $\mathrm{other}$、$\mathrm{\_sumReCall}$、
$\mathrm{\_sumRe}$ 以及在其内部直接定义的其他属性。

\begin{itemize}
  \item $\qualifiedthis{\mathrm{Foo}}$(等价地,$\dbi{n}$,其中
    $\mathrm{Foo}$ 在外围作用域之上 $n$ 步)
    取出名为 $\mathrm{Foo}$ 的记录。
  \item $\qualifiedthis{\mathrm{Foo}}.\ell$,或等价地,$\dbi{n}.\ell$,
    从 $\mathrm{Foo}$ 投影属性 $\ell$。
  \item 一个仅限定符的源 $q$(没有投影列表)继承
    整个外围作用域,这可能产生一棵无限深的
    树;图~\ref{fig:cheatsheet} 中裸的 $\mathrm{BuiltIns}$、$\mathrm{Math}$ 与
    $\mathrm{Operator}$ 源就是此类
    仅限定符的导入。
\end{itemize}

当限定符无歧义时,它可以省略:这种简写是
一个裸的投影列表 $w$。它的首个投影针对
那个在其记录字面量中直接定义首个投影标签的
最内层外围作用域来解析,而该源脱糖为
对应的索引形式。
例如,图~\ref{fig:cheatsheet} 中的 $\mathrm{\_sumReCall}.\mathrm{result}$、
$\mathrm{a}.\mathrm{Plus}$ 与 $\mathrm{\_totalCall}.\mathrm{result}$
就是此类省略限定符的简写。
由于这一搜索只检查记录字面量的定义,该
简写只能访问在记录字面量中定义的属性,而不能访问
那些通过继承而继承来的属性;对于继承来的属性或要绕过变量遮蔽,
都需要两种限定形式之一。
\else
An expression encloses an element list $s$ in braces, and each element
$c$ is either a definition $\ell \mapsto e$ or an inheritance source.
The grammar writes $s$ as a recursive list, but this list is read as a
\emph{set}: element order is irrelevant and duplicates have no effect,
so two element lists differing only by reordering or repetition denote
the same expression. Because expressions are sets, inheritance is
inherently commutative, idempotent, and associative.
The record $\mathrm{Complex} \mapsto \{ \mathrm{re} \mapsto \ldots,\;
\mathrm{im} \mapsto \ldots,\; \mathrm{Plus} \mapsto \ldots \}$ in
Figure~\ref{fig:cheatsheet} is one such multi-element list.

The $\mapsto$ in $\ell \mapsto e$ defines a property named $\ell$ with
value $e$; for example, $\mathrm{re} \mapsto \{\ldots\}$ and
$\mathrm{Complex} \mapsto \{\ldots\}$ in Figure~\ref{fig:cheatsheet}.
MIXINv2 also provides scalar types through a
foreign-function interface (FFI); inheritance-calculus is what remains
after removing the FFI.
None of the examples in this paper uses the FFI or any scalar values.

\paragraph{Inheritance}
An inheritance references a scope through a qualifier $q$ that
navigates the scope hierarchy upward, then projects properties downward
through a projection list $w$. The qualifier $q$ admits two equivalent
notations:
\begin{description}
  \item[Named form]
    $\qualifiedthis{\ell_{\mathrm{qualifier}}}.\, w$,
    analogous to Java's \texttt{Outer.this.field}, where $w$ is a
    projection list.
    Each $\{\ldots\}$ record has a label in its enclosing scope;
    $\ell_{\mathrm{qualifier}}$ names the target enclosing scope.
    For example, $\qualifiedthis{\mathrm{ComplexModule}}.\mathrm{Float}$
    in Figure~\ref{fig:cheatsheet} projects a single label, while
    $\qualifiedthis{\mathrm{Plus}}.\mathrm{other}.\mathrm{re}$ shows a
    multi-label projection list $w$.
  \item[Indexed form]
    $\dbi{n}.\, w$,
    where $n \ge 0$ is a de~Bruijn index~\cite{debruijn1972-nameless-dummies}:
    $n = 0$ refers to the reference's enclosing scope,
    and each increment of $n$ moves one scope level further out.
\end{description}
Both forms retrieve the target scope, that is, the fully inherited record
after all inheritance has been applied, and then project the properties
named by the projection list $w$ from it.
The named form desugars to the indexed form during parsing,
as described in Section~\ref{sec:mixin-trees}.

A new scope level is created at each $\{\ldots\}$ record literal.%
\footnote{
  We use \emph{record} for the syntactic construct, \emph{scope}
  for the name-resolution context it creates, and \emph{mixin} for its
  role as a composable unit; in inheritance-calculus these are three views
  of the same entity.
}
The scope contains all properties of that mixin, including those
inherited via inheritance sources.
For example, $\{a \mapsto \{\},\; b \mapsto \{\}\}$ is a single mixin whose scope
contains both $a$ and $b$; so is $\{r,\; a \mapsto \{\}\}$ where $r$
is a reference to a record containing $b$.
Inheritance sources do not create additional scope levels.
In Figure~\ref{fig:cheatsheet}, the record bound to $\mathrm{Plus}$ is
one such scope, containing $\mathrm{other}$, $\mathrm{\_sumReCall}$,
$\mathrm{\_sumRe}$, and the other properties defined directly within it.

\begin{itemize}
  \item $\qualifiedthis{\mathrm{Foo}}$ (equivalently, $\dbi{n}$ where
    $\mathrm{Foo}$ is $n$ steps above the enclosing scope)
    retrieves the record named $\mathrm{Foo}$.
  \item $\qualifiedthis{\mathrm{Foo}}.\ell$, or equivalently, $\dbi{n}.\ell$,
    projects property $\ell$ from $\mathrm{Foo}$.
  \item A qualifier-only source $q$ (with no projection list) inherits
    the entire enclosing scope, which may produce an infinitely deep
    tree; the bare $\mathrm{BuiltIns}$, $\mathrm{Math}$, and
    $\mathrm{Operator}$ sources in Figure~\ref{fig:cheatsheet} are
    qualifier-only imports of this kind.
\end{itemize}

When the qualifier is unambiguous, it may be omitted: the shorthand is
a bare projection list $w$. Its leading projection is resolved against
the innermost enclosing scope that defines the leading projection's
label directly in its record literal, and the source desugars to the
corresponding indexed form.
For example, $\mathrm{\_sumReCall}.\mathrm{result}$,
$\mathrm{a}.\mathrm{Plus}$, and $\mathrm{\_totalCall}.\mathrm{result}$
in Figure~\ref{fig:cheatsheet} are such qualifier-omitted shorthands.
Because this search inspects only record-literal definitions, the
shorthand can only access properties defined in record literals, not
those inherited through inheritance; one of the two qualified forms is
required for inherited properties or to bypass variable shadowing.
\fi

\section{\bilingual{Mixin Trees}{Mixin 树}}
\label{sec:mixin-trees}

\ifInheritanceChinese
在面向对象编程中,一个
\emph{mixin}~\cite{bracha1990-mixin-inheritance} 是一段行为片段;
它可以通过继承组合到一个类上。
继承演算只有一个抽象,即
\emph{深可合并 mixin}(deep-mergeable mixin),其成员是
\emph{深可合并属性}(deep-mergeable properties)。%
\footnote{
  在不引起歧义之处,我们简写为``mixin''与``property''。
  正如图~\ref{fig:cheatsheet} 所示,二者各自推广了它们在先前
  工作中的同名物:Bracha--Cook 的 mixin~\cite{bracha1990-mixin-inheritance} 与
  面向对象语言的属性。
}
它可以同时充当模块、类、对象、方法、字段以及
局部绑定,正如图~\ref{fig:cheatsheet} 所展示的那样,因为
没有不透明的方法:每个值都是
一个透明的记录,而继承总是对子树执行一次深度
合并。
我们把所得到的结构称为一棵\emph{mixin 树}。
它的语义是纯\emph{观察式}(observational)的:给定一个
继承演算程序,我们查询某个
标签 $\ell$ 是否是某个给定位置 $p$ 处的一个属性。

本节通过六个相互递归的
路径上的一阶函数来定义语义。
这六个函数把面向对象的概念映射为
集合论的运算:$\properties$ 与 $\supers$
推广成员查找与继承链,
$\overrides$ 与 $\bases$ 捕捉方法覆盖与
基类,而 $\resolve$ 与 $\this$ 处理名字
解析与限定 $\this$ 解析。
\else
In object-oriented programming, a
\emph{mixin}~\cite{bracha1990-mixin-inheritance} is a fragment of behavior
that can be composed onto a class via inheritance.
Inheritance-calculus has a single abstraction, the
\emph{deep-mergeable mixin}, whose members are
\emph{deep-mergeable properties}.%
\footnote{
  Where no ambiguity arises, we write simply ``mixin'' and ``property.''
  As Figure~\ref{fig:cheatsheet} shows, each generalizes its namesake in prior
  work: the Bracha--Cook mixin~\cite{bracha1990-mixin-inheritance} and the
  property of object-oriented languages.
}
which can simultaneously serve as module, class, object, method, field, and
local binding, as Figure~\ref{fig:cheatsheet} illustrated, because there
are no opaque methods: every value is
a transparent record, and inheritance always performs a deep
merge of subtrees.
We call the resulting structure a \emph{mixin tree}.
Its semantics is purely \emph{observational}: given an
inheritance-calculus program, one queries whether a
label $\ell$ is a property at a given position $p$.

This section defines the semantics via six mutually recursive
first-order functions on paths.
The six functions map object-oriented concepts to
set-theoretic operations: $\properties$ and $\supers$
generalize member lookup and inheritance chains,
$\overrides$ and $\bases$ capture method override and
base classes, and $\resolve$ and $\this$ handle name
resolution and qualified $\this$ resolution.
\fi

\subsection{\bilingual{Definitions}{定义}}
\label{sec:definitions}

\paragraph{\bilingual{Path}{路径}}
\ifInheritanceChinese
一个 $\Path$ 是一个标签序列
$(\ell_1, \ell_2, \ldots, \ell_n)$;它标识树中的一个
位置。
我们用 $p$ 表示一条路径。
\emph{根路径}是空序列 $()$。
给定一条路径 $p$ 与一个标签 $\ell$,
$p \snoc \ell$ 是序列 $p$ 用 $\ell$ 扩展后的结果。
对于非根路径,父路径 $\init(p)$ 是 $p$ 去掉
其最后一个元素后的结果,
而末尾标签 $\last(p)$ 是 $p$ 的最后一个元素。
为方便起见,当路径在上下文中清楚时,我们有时
用 $\ell$ 表示 $\last(p)$。
有了路径,我们就能陈述 AST 在
每条路径处提供什么。
\else
A $\Path$ is a sequence of labels
$(\ell_1, \ell_2, \ldots, \ell_n)$ that identifies a position in
the tree.
We write $p$ for a path.
The \emph{root path} is the empty sequence $()$.
Given a path $p$ and a label $\ell$,
$p \snoc \ell$ is the sequence $p$ extended with $\ell$.
For a nonroot path, the parent path $\init(p)$ is $p$ with
its last element removed,
and the final label $\last(p)$ is the last element of $p$.
For convenience, we sometimes write $\ell$ for $\last(p)$
when the path is clear from context.
With paths in hand, we can state what the AST provides
at each path.
\fi

\paragraph{AST}
\ifInheritanceChinese
来自第~\ref{sec:syntax} 节的一个表达式 $e$ 被解析成一棵 AST。
我们以两种等价的风格给出 AST:首先是一个固定其形状的
代数数据类型,然后是两个派生的投影;本节其余部分的
语义函数消费它们。

\emph{代数数据类型风格}把规范化后的 AST 给出为一个由节点项
构成的纯语法。一个节点 $N$ 把一个定义列表 $d$ 与一个
引用列表 $b$ 配成对;每个定义把一个标签绑定到一个子节点,而
每个引用是一个 de~Bruijn 索引 $n$ 连同一个投影列表
$w$:
\else
An expression $e$ from Section~\ref{sec:syntax} is parsed into an AST.
We give the AST in two equivalent styles: first an algebraic
datatype that fixes its shape, then two derived projections that
the semantic functions in the rest of this section consume.

The \emph{algebraic-datatype style} gives the normalized AST as a pure
grammar of node terms. A node $N$ pairs a definition list $d$ with a
reference list $b$; each definition binds a label to a child node, and
each reference is a de~Bruijn index $n$ together with a projection list
$w$:
\fi
\begin{align*}
  N \quad ::= \quad & \langle\, d \,;\; b \,\rangle
  && \text{\bilingual{(node: definitions and references)}{(节点:定义与引用)}} \\[6pt]
  d \quad ::= \quad & \varepsilon
  && \text{\bilingual{(no definitions)}{(没有定义)}} \\*
  \mid\quad & (\ell \mapsto N),\; d
  && \text{\bilingual{(definition, then more)}{(一个定义,然后更多)}} \\[6pt]
  b \quad ::= \quad & \varepsilon
  && \text{\bilingual{(no references)}{(没有引用)}} \\*
  \mid\quad & (n,\; w),\; b
  && \text{\bilingual{(reference, then more)}{(一个引用,然后更多)}} \\[6pt]
  w \quad ::= \quad & \varepsilon
  && \text{\bilingual{(no projections)}{(没有投影)}} \\*
  \mid\quad & \ell.\, w
  && \text{\bilingual{(projection, then more)}{(一个投影,然后更多)}}
\end{align*}
\ifInheritanceChinese
与表达式的元素列表一样,定义列表 $d$ 是
递归写成的,但读取时不计重排与重复:它
表示从其各绑定收集而来的有限偏映射 $\defs : \Label \rightharpoonup \Node$;该映射
对每个被定义的标签恰有一条目,而
$N.\defs(\ell)$ 是绑定到 $\ell$ 的子节点(当 $\ell$
未绑定时无定义)。同样地,引用列表 $b$ 表示由继承源
$(n,\; w)$ 构成的集合
$\refs \subseteq \mathbb{N} \times \Label^{*}$,读取时不计重排与重复。由于继承
是幂等的且记录是集合,这些商不丢失任何
信息。
该语法只固定形状;一个节点项是良构的,当它
满足:
\else
Like the element list of an expression, the definition list $d$ is
written recursively but read up to reordering and duplication: it
denotes the finite partial map $\defs : \Label \rightharpoonup \Node$
collected from its bindings, with one entry per defined label, and
$N.\defs(\ell)$ is the child bound to $\ell$ (undefined when $\ell$ is
unbound). Likewise the reference list $b$ denotes the set
$\refs \subseteq \mathbb{N} \times \Label^{*}$ of inheritance sources
$(n,\; w)$, read up to reordering and duplication. Because inheritance
is idempotent and a record is a set, these quotients lose no
information.
The grammar fixes only the shape; a node term is well-formed when it
satisfies:
\fi
\begin{enumerate}
  \item \ifInheritanceChinese 路径 $p$ 处的每个源 $(n,\; w)$ 满足 $n < |p|$。\else Each source $(n,\; w)$ at path $p$ satisfies $n < |p|$.\fi
  \item \ifInheritanceChinese 一个 $w$ 为空的源 $(n,\; \varepsilon)$ 是
    一个仅限定符源的规范化;第~\ref{sec:syntax} 节的表达式语法
    没有 $w$ 为空的仅投影源。\else A source $(n,\; \varepsilon)$ with empty $w$ is the
    normalization of a qualifier-only source; the expression grammar of
    Section~\ref{sec:syntax} has no projection-only source with empty
    $w$.\fi
\end{enumerate}
\ifInheritanceChinese
该节点语法被\emph{余归纳地}(coinductively)读取:通过
继承,一个子节点可能展开成一棵无界乃至
无限深的树,所以这些产生式描述的是最大不动点,
而非有限生成的最小不动点。整个程序是单个
全局节点 $\ASTroot$,即顶层记录的规范化 AST。

\emph{余代数风格}(coalgebraic style)用两个由路径 $p$ 索引的
原始函数替换全局树,每次观察一个节点。它们通过沿 $p$
导航从 $\ASTroot$ 派生而来:
从 $\ASTroot$ 出发,在 $p$ 的每个标签处跟随 $\defs$ 映射
到达 $p$ 处的节点,而这两个投影读出其
字段。
\else
The node grammar is read \emph{coinductively}: through
inheritance a child node may expand to a tree of unbounded or even
infinite depth, so the productions describe the greatest fixed point,
not a finitely generated least one. The whole program is a single
global node $\ASTroot$, the normalized AST of the top-level record.

The \emph{coalgebraic style} replaces the global tree by two
primitive functions indexed by a path $p$, observing one node at a
time. They are derived from $\ASTroot$ by navigating along $p$:
starting at $\ASTroot$ and following the $\defs$ map at each label
of $p$ reaches the node at $p$, and the two projections read off its
fields.
\fi
\begin{align}
  \defines(p) &= \operatorname{dom}\bigl(
    \ASTroot.\defs(\ell_1).\defs(\ell_2)\cdots\defs(\ell_n).\defs
  \bigr)
  \label{eq:defines}\\
  \inherits(p) &= \ASTroot.\defs(\ell_1).\defs(\ell_2)\cdots\defs(\ell_n).\refs
  \label{eq:inherits}
\end{align}
\ifInheritanceChinese
其中 $p = (\ell_1, \ldots, \ell_n)$,从而 $\defines(p)$ 是
那些使得 $p$ 包含一个定义 $\ell \mapsto e$ 的标签 $\ell$ 的集合,
而 $\inherits(p)$ 是 $p$ 处的引用对
$(n,\; w)$ 的集合。
在根路径 $()$ 处导航为空,给出
$\defines(()) = \operatorname{dom}(\ASTroot.\defs)$ 以及
$\inherits(()) = \ASTroot.\refs$。
下文的语义函数只使用 $\defines$ 与 $\inherits$,
所以任一风格皆可取作原始。

在解析期间,三种语法形式全都被解析为
de~Bruijn 索引对 $(n,\; w)$。
\emph{索引形式} $\dbi{n}.\, w$
已经直接携带 de~Bruijn 索引 $n$;
不需要解析。
路径 $p$ 处的一个\emph{具名引用} $\qualifiedthis{\ell_{\mathrm{qualifier}}}.\, w$
通过在 $p$ 的标签中找到
$\ell_{\mathrm{qualifier}}$ 的最后一次出现来解析。
设 $p_{\mathrm{target}}$ 为 $p$ 的前缀,直到并包括
那次出现。
de~Bruijn 索引为 $n = |p| - |p_{\mathrm{target}}| - 1$,
而投影列表 $w$ 原样保留。
路径 $p$ 处一个首投影为 $\ell$ 的\emph{词法引用} $w$
通过找到 $p$ 的最近前缀 $p'$
使得 $\ell \in \defines(p')$ 来解析。
de~Bruijn 索引为 $n = |p| - |p'| - 1$,而投影列表
$w$ 原样保留。
三种形式都产生同一种表示;下文的语义函数
只在 de~Bruijn 索引对上操作。

AST 是纯数据,由程序文本固定,而不是
运行期状态的函数。%
\footnote{
  一个实现可以非严格地解析 AST:$\ASTroot$ 不必
  预先被完全物化。每个 $\Node$ 把它的源
  文本保持为未解析,直到某个查询观察它,从而读取
  $\defines(p)$ 或 $\inherits(p)$ 时按需强制
  沿 $p$ 的子树。这只影响解析何时发生,而不影响其结果,
  因为 AST 是引用透明的。
}
有了两个原始函数 $\defines$ 与 $\inherits$,我们现在可以回答
本节开头提出的那个问题:
某个标签 $\ell$ 是否是某个给定位置 $p$ 处的一个属性。
\else
for $p = (\ell_1, \ldots, \ell_n)$, so that $\defines(p)$ is the set
of labels $\ell$ for which $p$ contains a definition $\ell \mapsto e$,
and $\inherits(p)$ is the set of reference pairs
$(n,\; w)$ at $p$.
At the root path $()$ the navigation is empty, giving
$\defines(()) = \operatorname{dom}(\ASTroot.\defs)$ and
$\inherits(()) = \ASTroot.\refs$.
The semantic functions below use only $\defines$ and $\inherits$,
so either style may be taken as primitive.

During parsing, all three syntactic forms are resolved to
de~Bruijn index pairs $(n,\; w)$.
The \emph{indexed form} $\dbi{n}.\, w$
already carries the de~Bruijn index $n$ directly;
no resolution is needed.
A \emph{named reference} $\qualifiedthis{\ell_{\mathrm{qualifier}}}.\, w$
at path $p$ is resolved by finding the last occurrence of
$\ell_{\mathrm{qualifier}}$ among the labels of $p$.
Let $p_{\mathrm{target}}$ be the prefix of $p$ up to and including
that occurrence.
The de~Bruijn index is $n = |p| - |p_{\mathrm{target}}| - 1$
and the projection list $w$ is carried over unchanged.
A \emph{lexical reference} $w$ at path $p$, with leading
projection $\ell$, is resolved by finding the nearest prefix $p'$ of $p$
such that $\ell \in \defines(p')$.
The de~Bruijn index is $n = |p| - |p'| - 1$ and the projection list
$w$ is carried over unchanged.
All three forms produce the same representation; the semantic functions
below operate only on de~Bruijn index pairs.

The AST is pure data, fixed by the program text and not a function
of runtime state.%
\footnote{
  An implementation may parse the AST non-strictly: $\ASTroot$ need
  not be materialized in full up front. Each $\Node$ keeps its source
  text unparsed until a query observes it, so that reading
  $\defines(p)$ or $\inherits(p)$ forces the subtree along $p$ on
  demand. This affects only when parsing happens, not its result,
  since the AST is referentially transparent.
}
Given the two primitives $\defines$ and $\inherits$, we can now answer
the question posed at the beginning of this section:
whether a label $\ell$ is a property at a given position $p$.
\fi

\paragraph{\bilingual{Properties}{属性}}
\ifInheritanceChinese
一条路径的属性是它自身的属性,连同
从所有 supers 继承而来的属性:
\else
The properties of a path are its own properties together with
those inherited from all supers:
\fi
\begin{equation}\label{eq:properties}
  \properties(p) =
  \bigl\{\; \ell \;\big|\;
    (\_,\; p_{\mathrm{override}}) \in \supers(p),\;
    \ell \in \defines(p_{\mathrm{override}})
  \;\bigr\}
\end{equation}
\ifInheritanceChinese
本节其余部分定义 $\supers$,
即传递的继承闭包,
及其依赖项。
\else
The remainder of this section defines $\supers$,
the transitive inheritance closure,
and its dependencies.
\fi

\paragraph{Supers}
\ifInheritanceChinese
直观上,$\supers(p)$ 收集 $p$ 所继承的每一条路径:
$p$ 自身的 $\overrides$,
加上 $p$ 的每个
直接基 $\bases$ 的 $\overrides$,如此传递下去。
每个结果与到达它所经由的继承点语境
$\init(p_{\mathrm{base}})$ 配成对。
我们下文定义的限定 $\this$ 解析需要这一来源,
以便把一个定义点 mixin 映射回
那些纳入它的继承点路径:
\else
Intuitively, $\supers(p)$ collects every path that $p$ inherits from:
the $\overrides$ of $p$ itself,
plus the $\overrides$ of each
direct base $\bases$ of $p$, and so on transitively.
Each result is paired with the inheritance-site context
$\init(p_{\mathrm{base}})$ through which it is reached.
This provenance is needed by qualified $\this$ resolution,
which we define below, to map a definition-site mixin back to
the inheritance-site paths that incorporate it:
\fi
\begin{equation}\label{eq:supers}
  \supers(p) =
  \bigl\{\; (\init(p_{\mathrm{base}}),\; p_{\mathrm{override}}) \;\big|\;
    p_{\mathrm{base}} \in \bases^*(p),\;
    p_{\mathrm{override}} \in \overrides(p_{\mathrm{base}})
  \;\bigr\}
\end{equation}
\ifInheritanceChinese
这里 $\bases^*$ 表示 $\bases$ 的自反传递闭包。
$\supers$ 公式依赖两个函数:
$\overrides$ 与单跳引用目标 $\bases$。
我们先定义 $\overrides$。
\else
Here $\bases^*$ denotes the reflexive-transitive closure of $\bases$.
The $\supers$ formula depends on two functions:
$\overrides$ and the one-hop reference targets $\bases$.
We define $\overrides$ first.
\fi

\paragraph{Overrides}
\ifInheritanceChinese
由于一个 mixin 可以扮演一个方法的角色,
一个子作用域应当能够
通过父作用域的标签找到它,即便该子作用域
并不重新定义那个标签;这就是方法覆盖。
具体地说,继承可以在同一个作用域层级引入同一标签的
多个定义;
$\overrides(p)$ 收集所有与 $p$ 共享
\emph{相同身份}(same identity)的这类路径,从而使它们的定义
相互合并而非相互遮蔽,以实现深度合并。
$p$ 的 overrides 包含 $p$ 自身
以及对于 $\init(p)$ 的每个
也定义 $\last(p)$ 的分支 $p_{\mathrm{branch}}$,$p_{\mathrm{branch}} \snoc \last(p)$:
\else
Since a mixin can play the role of a method,
a child scope should be
able to find it through the parent's label even if the child
does not redefine that label; this is method override.
Concretely, inheritance can introduce multiple
definitions of the same label at the same scope level;
$\overrides(p)$ collects all such paths that share the
\emph{same identity} as $p$, so that their definitions
merge rather than shadow each other, enabling deep merge.
The overrides of $p$ include $p$ itself
and $p_{\mathrm{branch}} \snoc \last(p)$ for every
branch $p_{\mathrm{branch}}$ of $\init(p)$ that also defines $\last(p)$:
\fi
\begin{equation}\label{eq:overrides}
  \overrides(p) =
  \begin{cases}
    \{p\} & \text{if } p = ()  \\[6pt]
    \{p\}
    \;\cup\;
    \left\{\; p_{\mathrm{branch}} \snoc \last(p) \;\middle|\;
      \begin{aligned}
        &(\_,\; p_{\mathrm{branch}}) \in \supers(\init(p)), \\
        &\text{s.t.}\; \last(p) \in \defines(p_{\mathrm{branch}})
      \end{aligned}
    \right\}
    & \text{if } p \neq ()
  \end{cases}
\end{equation}
\ifInheritanceChinese
还剩下要定义 $\bases$,即 $\supers$ 的另一个依赖项。
\else
It remains to define $\bases$, the other dependency of $\supers$.
\fi

\paragraph{Bases}
\ifInheritanceChinese
直观上,$\bases(p)$ 是 $p$ 通过引用直接
继承的那些路径,类比于面向对象语言中的直接
基类。
具体地说,$\bases$ 把 $p$ 的 $\overrides$ 中的每个引用
解析一步:
\else
Intuitively, $\bases(p)$ are the paths that $p$ directly
inherits from via references, analogous to the direct base
classes in object-oriented languages.
Concretely, $\bases$ resolves every reference in
$p$'s $\overrides$ one step:
\fi
\begin{equation}\label{eq:bases}
  \bases(p) =
  \left\{\; p_{\mathrm{target}} \;\middle|\;
    \begin{aligned}
      &p_{\mathrm{override}} \in \overrides(p), \\
      &(n,\; w) \in \inherits(p_{\mathrm{override}}), \\
      &p_{\mathrm{target}} \in \resolve(\init(p),\; p_{\mathrm{override}},\; n,\; w)
    \end{aligned}
  \right\}
\end{equation}
\ifInheritanceChinese
$\bases$ 公式调用引用解析函数 $\resolve$,我们接下来定义它。
\else
The $\bases$ formula calls the reference resolution function $\resolve$, which we define next.
\fi

\paragraph{\bilingual{Reference resolution}{引用解析}}
\ifInheritanceChinese
直观上,$\resolve$ 把一个语法引用转化为它在
完全继承后的树中所指向的路径集合。
它取一条继承点路径 $p_{\mathrm{site}}$、
一条定义点路径 $p_{\mathrm{def}}$、
一个 de~Bruijn 索引 $n$,
以及投影列表 $w$。
解析分两个阶段进行:
$\this$ 从外围作用域 $\init(p_{\mathrm{def}})$ 出发
执行 $n$ 次向上步进,
把定义点路径映射为继承点路径;
然后追加
$w = \ell_{\mathrm{projection},1}\ldots\ell_{\mathrm{projection},k}$
的向下投影:
\else
Intuitively, $\resolve$ turns a syntactic reference into the
set of paths it points to in the fully inherited tree.
It takes an inheritance-site path $p_{\mathrm{site}}$,
a definition-site path $p_{\mathrm{def}}$,
a de~Bruijn index $n$,
and the projection list $w$.
Resolution proceeds in two phases:
$\this$ performs $n$ upward steps
starting from the enclosing scope $\init(p_{\mathrm{def}})$,
mapping the definition-site path to inheritance-site paths;
then the downward projections of
$w = \ell_{\mathrm{projection},1}\ldots\ell_{\mathrm{projection},k}$
are appended:
\fi
\begin{multline}\label{eq:resolve}
  \resolve(p_{\mathrm{site}},\; p_{\mathrm{def}},\; n,\; w)
  = \\
  \bigl\{\;
    p_{\mathrm{current}} \snoc \ell_{\mathrm{projection},1} \snoc \cdots \snoc \ell_{\mathrm{projection},k}
    \;\big|
    p_{\mathrm{current}} \in
    \this(\{p_{\mathrm{site}}\},\; \init(p_{\mathrm{def}}),\; n)
  \;\bigr\}
\end{multline}
\ifInheritanceChinese
$\resolve$ 返回一个集合而非单条路径,因为
$\this$ 可能解析到多个目标。
我们现在定义 $\this$。
\else
$\resolve$ returns a set rather than a single path because
$\this$ may resolve to multiple targets.
We now define $\this$.
\fi

\paragraph{\bilingual{Qualified $\this$ resolution}{限定 $\this$ 解析}}
\ifInheritanceChinese
$\this$ 回答这个问题:
``在完全继承后的树中,定义点作用域
$p_{\mathrm{def}}$ 究竟住在哪里?''
每一步在前沿 $S$ 中每条路径的 supers 当中,
找出那些其 override 分量匹配
$p_{\mathrm{def}}$ 的,并把对应的
继承点路径收集为新的前沿。
经过 $n$ 步后,前沿就包含答案:
\else
$\this$ answers the question:
``in the fully inherited tree, where does
the definition-site scope $p_{\mathrm{def}}$ actually live?''
Each step finds, among the supers of every path in the
frontier $S$, those whose override component matches
$p_{\mathrm{def}}$, and collects the corresponding
inheritance-site paths as the new frontier.
After $n$ steps the frontier contains the answer:
\fi
{\small
  \begin{equation}\label{eq:this}
    \this(S,\; p_{\mathrm{def}},\; n) =
    \begin{cases}
      S & \text{if } n = 0 \\[6pt]
      \this\!\left(
        \left\{\; p_{\mathrm{site}} \;\middle|\;
          \begin{aligned}
            &p_{\mathrm{current}} \in S, \\
            &(p_{\mathrm{site}},\; p_{\mathrm{override}})
            \in \supers(p_{\mathrm{current}}), \\
            &\text{s.t.}\; p_{\mathrm{override}} = p_{\mathrm{def}}
          \end{aligned}
        \right\},\;
        \init(p_{\mathrm{def}}),\;
        n - 1
      \right)
      & \text{if } n > 0
    \end{cases}
\end{equation}}%
\ifInheritanceChinese
每一步 $\init$ 把 $p_{\mathrm{def}}$ 缩短一个标签,
而 $n$ 减少一。
由于 $n$ 是一个非负整数,该递归终止。%
\footnote{
  若在某一步不存在匹配对
  $(p_{\mathrm{site}},\; p_{\mathrm{override}})$,
  则前沿变为空集,而
  $\resolve$ 返回 $\varnothing$:该引用没有
  目标;不继承任何属性。
  这是无类型演算最为自洽的
  语义。
  MIXINv2\anon[~(补充材料)]{~\cite{mixin2025}} 更为严格:它的解析器拒绝
  前沿变为空的引用,在编译期报告一个
  解析错误。
}

这就完成了计算 $\properties(p)$ 所需的整条定义链。

附录~\ref{app:evaluation-trace} 把这六个方程逐步
应用到一个突出 $\resolve$ 集合值本质的例子上;
在那里前沿 $S$ 携带多个目标,而
$\this$ 同时解析到它们,而非如 Scala~3 与 NixOS 模块
这类面向对象语言中所见的单个目标
(附录~\ref{app:scala-multi-target})。
\else
At each step $\init$ shortens $p_{\mathrm{def}}$ by one label
and $n$ decreases by one.
Since $n$ is a nonnegative integer, the recursion terminates.%
\footnote{
  If no matching pair
  $(p_{\mathrm{site}},\; p_{\mathrm{override}})$ exists
  at some step, the frontier becomes the empty set, and
  $\resolve$ returns $\varnothing$: the reference has no
  target; no properties are inherited.
  This is the most self-consistent semantics for the
  untyped calculus.
  MIXINv2\anon[~(supplementary material)]{~\cite{mixin2025}} is stricter: its resolver rejects
  references whose frontier becomes empty, reporting a
  resolution error at compile time.
}

This completes the chain of definitions needed to compute $\properties(p)$.

Appendix~\ref{app:evaluation-trace} applies the six equations step by
step to an example highlighting the set-valued nature of $\resolve$,
where the frontier $S$ carries multiple targets that
$\this$ resolves to simultaneously rather than the single target
seen in object-oriented languages such as Scala~3 and NixOS modules
(Appendix~\ref{app:scala-multi-target}).
\fi

\subsection{\bilingual{Recursive Evaluation}{递归求值}}
\ifInheritanceChinese
六个方程~(\ref{eq:properties})--(\ref{eq:this})
通过相互递归定义 $\properties(p)$:
$\ell \in \properties(p)$ 成立,当且仅当它能
由这些函数的有限次应用链
确立。
这些函数\emph{就是}解释器:观察者
通过选择检查哪条路径来驱动计算,
而每个方程按需展开。
由于这些函数是纯的且与顺序无关,
(\ref{eq:supers}) 中的自反传递闭包 $\bases^*$
可以按广度优先或
深度优先来探索而不影响结果;
路径可以被 intern,从而结构相等
退化为指针相等;
而这六个函数构成一个以路径为键的可记忆化
动态规划递推式。

在操作上,求值使用\emph{tabled
求值}~\cite{tamaki1986-tabled-resolution}:
每个查询在首次遇到时被记忆化,而
重入调用,即对应于依赖图中的环的调用,
返回当前的累加器而非发散。
在指称上,六个方程定义了一个单调的
直接后承
算子~$T_P$~\cite{vanemden1976-predicate-logic-semantics};
它的右端只使用 $\cup$、$\in$、相等以及
正向集合概括,没有否定、集合
差或补。
由 Knaster--Tarski
定理~\cite{tarski1955-lattice-fixpoint-theorem},$T_P$ 有一个唯一的
最小不动点 $\mathrm{lfp}(T_P)$;它充当
指称语义。
这一指称刻画把上文引入的观察式
语义形式化:观察者驱动的查询
由同一组不动点方程计算。
求值是否终止取决于程序。
回忆 $\supers$(方程~\ref{eq:supers})取
$\bases$(方程~\ref{eq:bases})的自反传递
闭包;所得的图就是
\emph{继承层级}(inheritance hierarchy)。
由于每个继承源在 mixin 树中占据一个不同的节点,
这一层级可从
树结构中直接读出。
继承层级无环的程序在
一趟内终止。
当层级包含环时,
当每个查询有有限的答案域时求值
终止:有限格上的升链条件保证
收敛~\cite{bancilhon1986-recursive-query-strategies},
如同在 Datalog~\cite{vanemden1976-predicate-logic-semantics} 中那样。
当层级无限时,求值可能
发散;这是图灵完备性所固有的。
附录~\ref{app:well-definedness} 证明 tabled
求值绝不返回错误答案,并把
这些终止条件形式化。
\else
The six equations~(\ref{eq:properties})--(\ref{eq:this})
define $\properties(p)$ by mutual recursion:
$\ell \in \properties(p)$ holds if and only if it can
be established by a finite chain of applications of
these functions.
The functions \emph{are} the interpreter: the observer
drives the computation by choosing which path to
inspect, and each equation unfolds on demand.
Because the functions are pure and order-independent,
the reflexive-transitive closure $\bases^*$
in~(\ref{eq:supers}) may be explored breadth-first or
depth-first without affecting the result;
paths may be interned so that structural equality
reduces to pointer equality;
and the six functions form a memoizable
dynamic-programming recurrence keyed by paths.

Operationally, evaluation uses \emph{tabled
evaluation}~\cite{tamaki1986-tabled-resolution}:
each query is memoised on first encounter, and
re-entrant calls, which correspond to cycles in the dependency graph,
return the current accumulator rather than diverging.
Denotationally, the six equations define a monotone
immediate consequence
operator~$T_P$~\cite{vanemden1976-predicate-logic-semantics}
whose right-hand sides use only $\cup$, $\in$, equality, and
positive set comprehension, with no negation, set
difference, or complement.
By the Knaster--Tarski
theorem~\cite{tarski1955-lattice-fixpoint-theorem}, $T_P$ has a unique
least fixed point $\mathrm{lfp}(T_P)$, which serves as
the denotational semantics.
This denotational characterization formalizes the observational
semantics introduced above: the observer-driven queries are
computed by the same fixed-point equations.
Whether evaluation terminates depends on the program.
Recall that $\supers$ (equation~\ref{eq:supers}) takes the
reflexive-transitive closure of $\bases$
(equation~\ref{eq:bases}); the resulting graph is the
\emph{inheritance hierarchy}.
Since each inheritance source occupies a distinct node in
the mixin tree, this hierarchy is directly readable from
the tree structure.
Programs whose inheritance hierarchy is acyclic terminate in
one pass.
When the hierarchy contains cycles,
evaluation terminates when each query has a finite
answer domain: the ascending chain condition on the finite
lattice guarantees
convergence~\cite{bancilhon1986-recursive-query-strategies},
as in Datalog~\cite{vanemden1976-predicate-logic-semantics}.
When the hierarchy is infinite, evaluation may
diverge; this is inherent to Turing completeness.
Appendix~\ref{app:well-definedness} proves that tabled
evaluation never returns a wrong answer and formalizes
these termination conditions.
\fi

%% === 审稿人备忘：六个方程的渐近复杂度分析 ===
%%
%% 论文声称语义是一阶的，审稿人可能追问：一阶性是否带来额外的
%% 计算开销？下面按场景分析。
%%
%% 关键观察：set comprehension 自动去重（集合语义），加上路径
%% intern（结构相等退化为指针相等）和 memoization，相同的
%% (path, query) 对只计算一次。这不是可选的优化，而是集合语义
%% 的固有性质——集合不会包含重复元素。
%%
%% --- 场景 1：λ-calculus 翻译 ---
%% Single-Path Lemma (Lemma 3) 保证 this 的 frontier 集合 S
%% 始终只有一个元素。每个路径的 overrides 也只有一个元素。
%% bases 每个节点最多一个继承源。supers 的传递闭包是线性链。
%% 复杂度：O(n)，其中 n 是 ANF term 的大小。
%% 与传统环境机（environment machine）一致。
%%
%% --- 场景 2：Object algebra（k 个 constructor，t 个 operation 文件，m 个字段） ---
%% 多个 operation 文件各自继承同一个 data 文件，然后在最终组合时
%% merge。每个 constructor 在每个 operation 文件里各有一个定义，
%% 通过 overrides 合并，所以 overrides 大小为 O(t)。
%%
%% 关键点：this 的 frontier 是 1，不是"去重后 O(1)"，就是 1。
%% 原因：qualified this 解析的是外层 scope（如 NatFactory），
%% 而这个外层 scope 在继承链上只有一条路径到达——因为是从一个
%% 具体的 operation 定义点出发往上走的。multi-target this 的前提
%% 是同一个定义点被继承到多个不同的 site，这只在 trie union 时
%% 发生。Object algebra 的继承是把多个 operation merge 到同一个
%% factory 上，不是把多个 factory merge 成一个值。方向不同。
%%
%% 这也解释了与 Scala 的关系：Scala 拒绝 multi-target self-reference resolution
%%（Appendix B 的 "conflicting base types"），但在 object algebra
%% 场景下 this 本来就是 single-target 的，所以 Scala 的拒绝是
%% 过度保守——它拒绝了一类去重后（实际上根本不需要去重，因为
%% frontier 就是 1）完全合法的组合。
%%
%% 复杂度：O(k · m · t)，与 Scala trait linearization 一致。
%%
%% --- 场景 3：Nat/BinNat/Boolean 算术 ---
%% 数值的递归深度为 n（Nat 值的大小），每次递归触发一次 this
%% 解析（single-target），遍历 O(t) 个 override。
%% Nat 算术复杂度：O(n · t)，即 Peano 算术 + 常数因子。
%% BinNat 算术复杂度：O(log n · t)，即二进制算术 + 常数因子。
%%
%% --- 场景 4：Cartesian product / 逻辑编程 ---
%% 这是唯一出现真正 multi-target this 的场景。
%% trie union 的 bases 包含多条路径，每条路径的 addend 通过
%% this 解析到另一个 trie union，路径数相乘。
%% 对于 |R₁| × |R₂| 的 join，复杂度为 O(|R₁| · |R₂|)
%% （memoization 下）。
%%
%% 这与 B-Tree 引擎的 Datalog nested-loop join 复杂度一致。
%% trie 天然有索引——路径本身是前缀树，共享前缀的值自动共享
%% 计算（memoization 命中）。这与 Datalog 引擎对 B-Tree 做
%% prefix scan 在精神上一致。
%%
%% --- 总结 ---
%%
%% | 场景               | 复杂度              | 对比               |
%% |--------------------|---------------------|--------------------|
%% | λ-calculus         | O(n)                | = 环境机           |
%% | Object algebra     | O(k·m·t)            | = Scala线性化      |
%% | Nat 算术           | O(n·t)              | = Peano 算术       |
%% | BinNat 算术        | O(log n · t)        | = 二进制算术       |
%% | Cartesian product  | O(|R₁|·|R₂|)       | = nested-loop join |
%%
%% 每个场景都与对应的专用引擎复杂度一致。一阶性没有带来额外开销。
%%
%% Single-Path Lemma 的适用范围比论文明示的更大：不仅 λ-calculus
%% 翻译是 single-path，所有不涉及 trie union 的正常 OOP/object
%% algebra 代码都是 single-path。multi-target this 是 trie union
%% 的专属现象，而 trie union 对应的恰好是关系代数的 join，其
%% 复杂度是固有的，与引擎无关。
%% ================================================================

\section{\bilingual{Embedding the \texorpdfstring{$\lambda$}{λ}-Calculus}{嵌入 \texorpdfstring{$\lambda$}{λ}-演算}}
\label{sec:translation}
\label{sec:forward-translation}

\ifInheritanceChinese
以下翻译将 A-正规形式（ANF）~\cite{flanagan1993-essence-compiling-continuations}的 $\lambda$-演算映射到继承演算。
在 ANF 中,每个中间结果都绑定到一个名字,且每个表达式至多包含一次应用。
我们用一个\emph{值绑定} $\mathbf{let}\; x = V \;\mathbf{in}\; M$ 来扩展标准 ANF 语法,该值绑定将一个并非应用结果的值命名:%
\footnote{
  标准 ANF 允许将值内联用作操作数;值绑定产生式为其命名,使得翻译后程序的每个子表达式都可通过词法引用访问(附录~\ref{app:expressiveness})。
}
\[
  \begin{array}{r@{\;::=\;}l}
    M & \mathbf{let}\; x = V_1\; V_2 \;\mathbf{in}\; M
    \mid \mathbf{let}\; x = V \;\mathbf{in}\; M
    \mid V_1\; V_2
    \mid V \\
    V & x \mid \lambda x.\, M
  \end{array}
\]
每个 $\lambda$-项都可以通过对中间结果命名的方式机械地转换为 ANF。\footnote{%
  这可以选择性地包括对值的命名;由于每个 ANF 项也是一个 $\lambda$-项,ANF 既不是 $\lambda$-演算的限制,也不是其扩展。
}
ANF 结构与继承演算天然吻合:每个 $\mathbf{let}$-绑定成为一个命名的记录属性,每次应用则通过 $.\mathrm{result}$ 包装并投影。

设 $\mathcal{T}$ 为翻译函数。每个 $\lambda$-抽象在翻译后的 mixin 树中引入一个作用域层级。
\else
The following translation maps the $\lambda$-calculus in
A-normal form (ANF)~\cite{flanagan1993-essence-compiling-continuations} to inheritance-calculus.
In ANF, every intermediate result is bound to a name and each
expression contains at most one application.
We extend the standard ANF grammar with a \emph{value binding}
$\mathbf{let}\; x = V \;\mathbf{in}\; M$, which gives a name to a
value that is not an application result:%
\footnote{
  Standard ANF allows values inline as operands;
  the value-binding production names them so that every
  subexpression of the translated program is accessible
  via a lexical reference (Appendix~\ref{app:expressiveness}).
}
\[
  \begin{array}{r@{\;::=\;}l}
    M & \mathbf{let}\; x = V_1\; V_2 \;\mathbf{in}\; M
    \mid \mathbf{let}\; x = V \;\mathbf{in}\; M
    \mid V_1\; V_2
    \mid V \\
    V & x \mid \lambda x.\, M
  \end{array}
\]
Every $\lambda$-term can be mechanically converted to ANF by
naming intermediate results.\footnote{%
  This can optionally include naming values as well;
  since every ANF term is also a $\lambda$-term,
  ANF is neither a restriction nor an extension of the $\lambda$-calculus.
}
The ANF structure matches inheritance-calculus naturally: each
$\mathbf{let}$-binding becomes a named record property, and
each application is wrapped and projected via
$.\mathrm{result}$.

Let $\mathcal{T}$ denote the translation function.
Each $\lambda$-abstraction introduces one scope level
in the translated mixin tree.
\fi

\medskip
\begin{center}
  \small
  \begin{tabular}{l@{\quad$\longrightarrow$\quad}l}
    $x$ \text{\scriptsize(\bilingual{$\lambda$-bound, index $n'$}{$\lambda$-绑定,索引 $n'$})} &
    $\dbi{n'}.\mathrm{argument}$ \\[4pt]
    $x$ \text{\scriptsize(\bilingual{let-bound}{let-绑定})} &
    $x.\mathrm{result}$ \\[4pt]
    $\lambda x.\, M$ &
    $\{\mathrm{argument} \mapsto \{\},\;
    \mathrm{result} \mapsto \mathcal{T}(M)\}$ \\[4pt]
    $\mathbf{let}\; x = V_1\; V_2
    \;\mathbf{in}\; M$ &
    $\{x \mapsto \{\mathcal{T}(V_1),\;
      \mathrm{argument} \mapsto \mathcal{T}(V_2)\},$ \\
    & $\;\mathrm{result} \mapsto \mathcal{T}(M)\}$ \\[4pt]
    $\mathbf{let}\; x = V
    \;\mathbf{in}\; M$ &
    $\{x \mapsto \{\mathrm{result} \mapsto \mathcal{T}(V)\},\;
    \mathrm{result} \mapsto \mathcal{T}(M)\}$ \\[4pt]
    $V_1\; V_2$ \text{\scriptsize(\bilingual{tail call}{尾调用})} &
    $\{\mathrm{tailCall} \mapsto \{\mathcal{T}(V_1),\;
      \mathrm{argument} \mapsto \mathcal{T}(V_2)\},$ \\
    & $\;\mathrm{result} \mapsto \{\mathrm{tailCall}.\mathrm{result}\}\}$
  \end{tabular}
\end{center}
\medskip

\ifInheritanceChinese
带 de~Bruijn 索引 $n$ 的 $\lambda$-绑定变量 $x$(即从引用处到其绑定者之间包围的 $\lambda$ 数目)翻译为 $\dbi{n'}.\mathrm{argument}$,该引用到达绑定 $x$ 的抽象并访问其 $\mathrm{argument}$ 槽。%
\footnote{\label{fn:debruijn-adjust}在 $\lambda$-演算中,de~Bruijn 索引只计数 $\lambda$-抽象。在继承演算中,de~Bruijn 索引计数所有包围的作用域层级,包括由 $\mathbf{let}$-绑定和尾调用引入的层级。因此翻译会调整该索引:若一个 $\lambda$-绑定变量跨越 $k$ 个中间的 $\mathbf{let}$-绑定或尾调用作用域(即 $k$ 计数从变量出现处到其绑定的 $\lambda$ 在 ANF 语法树中路径上的 $\mathbf{let}$- 和尾调用构造数),则其继承演算索引为 $n' = n + k$,而非 $n$。let-绑定变量不受影响,因为它们使用词法引用(如 $x.\mathrm{result}$)而非 de~Bruijn 索引。我们假设翻译引入的合成标签($\mathrm{argument}$、$\mathrm{result}$、$\mathrm{tailCall}$)与用户定义的标签(如 let-绑定名 $x$)来自不相交的字母表,因此不会发生冲突。}
let-绑定变量 $x$ 翻译为词法引用 $x.\mathrm{result}$,从封闭记录中的兄弟属性 $x$ 投影应用结果。

抽象 $\lambda x.\, M$ 翻译为具有自身属性 $\{\mathrm{argument}, \mathrm{result}\}$ 的记录;我们称之为\emph{抽象形状}。在应用 $\mathbf{let}$-绑定规则中,名字 $x$ 绑定到继承记录 $\{\mathcal{T}(V_1),\;\mathrm{argument} \mapsto \mathcal{T}(V_2)\}$,而 $\mathcal{T}(M)$ 可引用 $x.\mathrm{result}$ 以获取应用的值。在值 $\mathbf{let}$-绑定规则中,$x$ 将 $\mathcal{T}(V)$ 包装在一个 $\mathrm{result}$ 子节点之下,使得统一访问器 $x.\mathrm{result}$ 产生 $\mathcal{T}(V)$ 本身,即翻译后的值。这个额外的包装在语义上是惰性的,但与证明相关:它使得用于应用的同一 $\mathrm{result}$-追踪论证在充分性证明中能均匀地扩展到值 $\mathbf{let}$-绑定(附录~\ref{app:bohm-tree-proofs})。在尾调用规则中,新标签 $\mathrm{tailCall}$ 起到相同作用,而 $\mathrm{result}$ 投影出答案。在应用和尾调用两种情形中,执行应用的继承被封装在一个命名属性之后,使其内部结构(被调用者的 $\mathrm{argument}$ 和 $\mathrm{result}$ 标签)不会泄漏到封闭作用域。

以上六条规则构成了从 $\lambda$-演算(经由 ANF)到继承演算的完整翻译:每个 ANF 项都有一个像,且构造是复合的。由于 $\lambda$-演算是图灵完备的,继承演算也是。%
\footnote{
  MIXINv2 将符号解析(解析词法引用和限定 $\this$ 引用)与运行时 mixin 求值分离,但符号解析是逐符号即时执行的,而非提前执行,因此它不提供传统意义上的静态类型检查。尽管如此,MIXINv2 的解析阶段比无类型的继承演算更严格。本文中的所有示例均从 MIXINv2 移植而来,并同时符合两种语义。为满足解析器而添加的元素对此处呈现的算法无害,且可能有助于可读性。
}
然而,图灵完备性并未说明翻译保留了什么;下一小节将精确刻画语义对应关系。
\else
A $\lambda$-bound variable $x$ with de~Bruijn index $n$
(the number of enclosing $\lambda$s between
the reference and its binder) becomes
$\dbi{n'}.\mathrm{argument}$,
which reaches the abstraction that binds $x$ and
accesses its $\mathrm{argument}$ slot.%
\footnote{\label{fn:debruijn-adjust}In the $\lambda$-calculus, de~Bruijn
  indices count only $\lambda$-abstractions.
  In inheritance-calculus, de~Bruijn indices count all
  enclosing scope levels, including those introduced by
  $\mathbf{let}$-bindings and tail calls.
  The translation therefore adjusts the index:
  if a $\lambda$-bound variable crosses $k$ intervening
  $\mathbf{let}$-binding or tail-call scopes
  (that is, $k$ counts the $\mathbf{let}$- and tail-call
    constructs on the path from the variable occurrence
  to its binding $\lambda$ in the ANF syntax tree),
  its inheritance-calculus
  index is $n' = n + k$,
  not~$n$.
  Let-bound variables are unaffected, as they use
  lexical references (such as $x.\mathrm{result}$) rather than
  de~Bruijn indices.
  We assume that the synthetic labels introduced by the
  translation ($\mathrm{argument}$, $\mathrm{result}$,
  $\mathrm{tailCall}$) and the user-defined labels
  (let-bound names such as~$x$) are drawn from disjoint
  alphabets, so no collision can arise.
}
A let-bound variable $x$ becomes the lexical reference
$x.\mathrm{result}$, projecting the application result
from the sibling property $x$ in the enclosing record.

An abstraction $\lambda x.\, M$ translates to a record with own
properties $\{\mathrm{argument}, \mathrm{result}\}$; we call this
the \emph{abstraction shape}.
In the application $\mathbf{let}$-binding rule, the name $x$
binds to the inherited record
$\{\mathcal{T}(V_1),\;\mathrm{argument} \mapsto \mathcal{T}(V_2)\}$,
and $\mathcal{T}(M)$ may reference $x.\mathrm{result}$ to obtain
the value of the application.
In the value $\mathbf{let}$-binding rule, $x$ wraps $\mathcal{T}(V)$
under a $\mathrm{result}$ child so that the uniform accessor
$x.\mathrm{result}$ yields $\mathcal{T}(V)$, the translated value
itself.
This extra wrapper is semantically inert but proof-relevant:
it lets the same $\mathrm{result}$-following argument used for
applications extend uniformly to value
$\mathbf{let}$-bindings in the adequacy proof
(Appendix~\ref{app:bohm-tree-proofs}).
In the tail-call rule, the fresh label $\mathrm{tailCall}$ serves
the same purpose, and $\mathrm{result}$ projects the answer.
In the application and tail-call cases, the inheritance that
performs the application is
encapsulated behind a named property, so its internal
structure (the $\mathrm{argument}$ and $\mathrm{result}$ labels
of the callee) does not leak to the enclosing scope.

The six rules above are a complete translation from the
$\lambda$-calculus (via ANF) to inheritance-calculus: every ANF
term has an image, and the construction is compositional.
Since the $\lambda$-calculus is Turing complete, so is
inheritance-calculus.%
\footnote{
  MIXINv2 separates symbol resolution (which resolves
  lexical references and qualified $\this$ references) from runtime
  mixin evaluation, but symbol resolution is performed
  just-in-time per symbol rather than ahead of time, so it does
  not provide static type checking in the traditional sense.
  Nevertheless, MIXINv2's resolution phase is stricter than the
  untyped inheritance-calculus.
  All examples in this paper are ported from MIXINv2
  and conform to both semantics.
  Elements added to satisfy the resolver are harmless
  for the algorithms presented here and may aid readability.
}
Turing completeness, however, says nothing about what the
translation preserves; the next subsection pins down
the semantic correspondence.
\fi

\subsection{\bilingual{B\"ohm Tree Correspondence}{Böhm 树对应}}
\label{sec:bohm-tree}

\ifInheritanceChinese
翻译 $\mathcal{T}$ 将 $\lambda$-演算嵌入继承演算的一个子语言中。直觉上,mixin 树以附加的语义结构增广了 AST;对于翻译后的 $\lambda$-项,这一附加结构的投影与 Böhm 树同构。若从根出发追踪有限次 $\mathrm{result}$ 投影后到达一个同时拥有 $\mathrm{argument}$ 和 $\mathrm{result}$ 的节点(即 $\mathcal{T}(\lambda x.\, M)$ 产生的\emph{抽象形状}),则称 mixin 树 $T$ \emph{收敛},记为 $T{\Downarrow}$。翻译保持并反映收敛性(充分性)。记 $\mathcal{T}(M) \approx_\lambda \mathcal{T}(N)$ 当对所有封闭的 $\lambda$-演算上下文 $C[\cdot]$ 均有 $\mathcal{T}(C[M]){\Downarrow} \Leftrightarrow \mathcal{T}(C[N]){\Downarrow}$;则 $\approx_\lambda$ 恰好与 Böhm 树等价~\cite{barendregt1984-lambda-calculus} 重合。形式化陈述和证明见附录~\ref{app:bohm-tree-proofs}。注意 $\approx_\lambda$ 仅在 $\lambda$-演算上下文上量化;继承演算的上下文区分能力严格更强,如第~\ref{sec:asymmetry} 节所证。以上对应关系在首次迭代近似 $T_P\!\uparrow\!1$ 处成立;我们现在考察在完整最小不动点处会发生什么。

在 $T_P\!\uparrow\!1$ 处,每个方程被求值一次,可重入调用返回 $\bot$;这对应于 $\beta$-规约。在完整的 $\mathrm{lfp}(T_P)$ 处,可重入引用改而解析为幂集格上的最小不动点:迭代在\emph{容器}层级计算不动点,求出满足 $S = F(S)$ 的集合 $S$。这正是递归 Datalog 在循环图上收敛的机制(附录~\ref{app:datalog-encoding},第~\ref{sec:recursive-rules} 节);既非不动点组合子(如 $Y$),也非 Haskell 的 \texttt{mfix}(\texttt{RecursiveDo})能够表达它:两者都在项层级打结,而非在语义所累积的集合上打结。破坏 Böhm 树对应关系的这种\emph{幂集}更多确定性(图~\ref{fig:calculi-matrix})的最小见证例是 $\mathbf{let}\; x = x \;\mathbf{in}\; x$:它在 $\beta$-规约下发散,但其翻译 $\{x \mapsto \{\mathrm{result} \mapsto \{x.\mathrm{result}\}\},\; \mathrm{result} \mapsto \{x.\mathrm{result}\}\}$ 在 $\mathrm{lfp}(T_P)$ 下稳定,tabling 求值检测到循环并返回确定结果 $\properties(\mathrm{root}) = \{x, \mathrm{result}\}$。因此本小节的对应关系是 $T_P\!\uparrow\!1$ 独有的性质。下一小节考察继承演算是否也严格更具表达力。
\else
The translation $\mathcal{T}$ embeds the $\lambda$-calculus
into a sublanguage of inheritance-calculus.
Intuitively, a mixin tree augments the AST with
additional semantic structure; for translated
$\lambda$-terms, the projection of this additional
structure is isomorphic to the B\"ohm tree.
A mixin tree $T$ \emph{converges}, written $T{\Downarrow}$,
if following finitely many $\mathrm{result}$ projections from
the root reaches a node with both $\mathrm{argument}$ and
$\mathrm{result}$ among its properties, that is, the \emph{abstraction shape}
produced by $\mathcal{T}(\lambda x.\, M)$.
The translation preserves and reflects convergence
(adequacy).
Write $\mathcal{T}(M) \approx_\lambda \mathcal{T}(N)$
when $\mathcal{T}(C[M]){\Downarrow} \Leftrightarrow
\mathcal{T}(C[N]){\Downarrow}$ for every closing
$\lambda$-calculus context $C[\cdot]$;
then $\approx_\lambda$ coincides with B\"ohm tree
equivalence~\cite{barendregt1984-lambda-calculus}.
Formal statements and proofs are in
Appendix~\ref{app:bohm-tree-proofs}.
Note that $\approx_\lambda$ quantifies over $\lambda$-calculus
contexts only; inheritance-calculus contexts are strictly more
discriminating, as Section~\ref{sec:asymmetry} establishes.
The correspondence above holds at the first-iteration
approximation $T_P\!\uparrow\!1$; we now ask what happens
at the full least fixed point.

At $T_P\!\uparrow\!1$, each equation is evaluated once and
re-entrant calls return~$\bot$; this corresponds to
$\beta$-reduction.
At the full $\mathrm{lfp}(T_P)$, re-entrant references instead
resolve to least fixed points over the powerset lattice: the
iteration computes fixed points at the \emph{container} level,
finding a set~$S$ such that $S = F(S)$.
This is the mechanism by which recursive Datalog converges on
cyclic graphs (Appendix~\ref{app:datalog-encoding},
Section~\ref{sec:recursive-rules}), and neither a fixed-point
combinator such as~$Y$ nor Haskell's \texttt{mfix}
(\texttt{RecursiveDo}) expresses it: both tie the knot at the
term level, not over the sets the semantics accumulates.
The smallest witness that this \emph{powerset} more-definedness
(Figure~\ref{fig:calculi-matrix}) breaks the B\"ohm-tree
correspondence is
$\mathbf{let}\; x = x \;\mathbf{in}\; x$:
it diverges under $\beta$-reduction, but its translation
$\{x \mapsto \{\mathrm{result} \mapsto \{x.\mathrm{result}\}\},\;
\mathrm{result} \mapsto \{x.\mathrm{result}\}\}$
stabilizes under $\mathrm{lfp}(T_P)$, where tabled evaluation
detects the cycle and returns the definite result
$\properties(\mathrm{root}) = \{x, \mathrm{result}\}$.
The correspondence of this subsection is therefore a property
of $T_P\!\uparrow\!1$ alone.
The next subsection examines whether inheritance-calculus is
also strictly more expressive.
\fi

\subsection{\bilingual{Expressive Asymmetry}{表达力不对称}}
\label{sec:asymmetry}
\label{sec:formal-separation}
\ifInheritanceChinese
两个 Böhm 树等价的 $\lambda$-项,在 $\lambda$-演算中因此无法区分,可以被一个混合了 $\lambda$-演算与继承演算构造子的上下文分离;该上下文通过开放扩展在已有定义上继承新的可观测投影。根据 Felleisen 的表达力判据~\cite{felleisen1991-expressive-power},继承演算严格比惰性 $\lambda$-演算更具表达力(定理~\ref{thm:expressive-asymmetry}):$\lambda$-演算可以通过 ANF 转换和 $\mathcal{T}$ 宏表达于继承演算中(引理~\ref{lem:anf-equiv},定理~\ref{thm:forward-macro}),但惰性 $\lambda$-演算无法宏表达继承(定理~\ref{thm:cartesian-nonexpressibility})。形式化定义、构造和证明见附录~\ref{app:expressiveness};该分离构造本身是表达式问题的一个微型实例,它在不修改已有定义的情况下添加了一个新的可观测投影。
\else
Two $\lambda$-terms that are B\"ohm-tree equivalent,
and therefore indistinguishable in the $\lambda$-calculus,
can be
separated by a context that mixes $\lambda$-calculus and
inheritance-calculus constructors, inheriting new observable
projections onto an existing definition via open extension.
By Felleisen's expressiveness
criterion~\cite{felleisen1991-expressive-power}, inheritance-calculus is
strictly more expressive than the lazy $\lambda$-calculus
(Theorem~\ref{thm:expressive-asymmetry}):
the $\lambda$-calculus can be macro-expressed in
inheritance-calculus via ANF conversion
and~$\mathcal{T}$
(Lemma~\ref{lem:anf-equiv}, Theorem~\ref{thm:forward-macro}),
but the lazy $\lambda$-calculus cannot macro-express inheritance
(Theorem~\ref{thm:cartesian-nonexpressibility}).
The formal definitions, constructions, and proofs are given in
Appendix~\ref{app:expressiveness};
the separating construction is itself a miniature instance of
the Expression Problem, adding a new observable projection to
an existing definition without modifying it.
\fi

\section{\bilingual{Case Study: Nat Arithmetic}{案例研究：Nat(自然数)算术}}
\label{sec:case-study}
\label{sec:expression-problem}

\ifInheritanceChinese
我们在继承演算中实现了布尔逻辑与自然数算术,包括一元与二进制两种形式。
本节详细追踪一元 Nat 的情形:其数据表示、加法与相等判断。
该实现遵循普通的 Gang-of-Four 设计模式~\cite{gamma1994-design-patterns},采用声明式面向对象风格,将每项关注点沿两个轴分解为 mixin:操作轴与情形轴。由于演算本身没有方法,所有名称均为名词或形容词:UpperCamelCase 名称扮演模块、类型、数据构造子与方法的角色,lowerCamelCase 名称扮演字段与参数的角色。
\else
We implemented boolean logic and natural number arithmetic, both unary
and binary, in inheritance-calculus.
This section traces the unary Nat case in detail: its data representation,
addition, and equality.
The implementation follows ordinary Gang-of-Four design
patterns~\cite{gamma1994-design-patterns} in a declarative object-oriented
style, with each concern split into mixins along two axes: operations and
cases. Because the calculus has no methods of its own, every name is a noun or
adjective: UpperCamelCase names play the roles of modules, types, data
constructors, and methods, and lowerCamelCase names the roles of fields and
parameters.
\fi

\paragraph{NatData}
\ifInheritanceChinese
mixin $\mathrm{NatData}$ 是自然数接口及其工厂。$\mathrm{NatFactory}$ 是一个工厂,其抽象产品类型被别名为 $\mathrm{Nat}$:
\else
The mixin $\mathrm{NatData}$ is the natural-number interface and its factory.
$\mathrm{NatFactory}$ is a factory whose abstract product type is aliased as $\mathrm{Nat}$:
\fi
\[
  \mathrm{NatData} \mapsto \left\{
    \begin{aligned}
      &\mathrm{NatFactory} \mapsto \{\mathrm{Product} \mapsto \{\}\},\\
      &\mathrm{Nat} \mapsto \{\mathrm{NatFactory}.\mathrm{Product}\}
  \end{aligned}\right\}
\]

\paragraph{NatDataZero}
\ifInheritanceChinese
每个数据构造子都是一个独立的实现 mixin,继承 $\mathrm{NatData}$。一个 $\mathrm{Zero}$ 值继承 $\mathrm{Product}$:
\else
Each data constructor is a separate implementation mixin that inherits
$\mathrm{NatData}$.
A $\mathrm{Zero}$ value inherits $\mathrm{Product}$:
\fi
\[
  \mathrm{NatDataZero} \mapsto \left\{
    \begin{aligned}
      &\mathrm{NatData},\\
      &\mathrm{NatFactory} \mapsto \{\mathrm{Zero} \mapsto \{\qualifiedthis{\mathrm{NatFactory}}.\mathrm{Product}\}\}
  \end{aligned}\right\}
\]

\paragraph{NatDataSuccessor}
\ifInheritanceChinese
一个 $\mathrm{Successor}$ 值继承 $\mathrm{Product}$,并暴露一个类型为 $\mathrm{Product}$ 的 $\mathrm{predecessor}$ 属性:
\else
A $\mathrm{Successor}$ value inherits $\mathrm{Product}$ and
exposes a $\mathrm{predecessor}$ property whose type is $\mathrm{Product}$:
\fi
\[
  \begin{aligned}
    &\mathrm{NatDataSuccessor} \mapsto \\
    &\left\{
    \begin{aligned}
      &\mathrm{NatData},\\
      &\mathrm{NatFactory} \mapsto \\
      &\left\{\mathrm{Successor} \mapsto \left\{
        \begin{aligned}
          &\qualifiedthis{\mathrm{NatDataSuccessor}}.\mathrm{NatFactory}.\mathrm{Product},\\
          &\mathrm{predecessor} \mapsto \{\qualifiedthis{\mathrm{NatDataSuccessor}}.\mathrm{NatFactory}.\mathrm{Product}\}
      \end{aligned}\right\}\right\}
  \end{aligned}\right\}
  \end{aligned}
\]
\ifInheritanceChinese
与 $\mathrm{NatDataZero}$ 合在一起,这以开放的方式完成了自然数的 Peano 编码,即函子 $F(X) = 1 + X$ 的初始代数:两个构造子作为独立的 mixin 加入,而非固定在单个封闭的数据类型定义中。
\else
Together with $\mathrm{NatDataZero}$, this completes the Peano encoding of
the natural numbers, the initial algebra of the functor $F(X) = 1 + X$,
in an open way: the two constructors are added as independent mixins
rather than fixed in a single closed datatype definition.
\fi

\paragraph{NatPlus}
\ifInheritanceChinese
加法操作沿情形轴分解为一个\emph{接口} mixin $\mathrm{NatPlus}$,该接口 mixin 保留关注点的名称并在抽象 $\mathrm{Product}$ 上声明命令,以及每个情形对应一个\emph{实现} mixin,实现 mixin 在名称后附加构造子名称(即下文的 $\mathrm{NatPlusZero}$ 与 $\mathrm{NatPlusSuccessor}$)。$\mathrm{NatPlus}$ 继承 $\mathrm{NatData}$。面向对象方法会用动词命名的地方,命令模式~\cite{gamma1994-design-patterns} 将其具化为以名词命名的命令,因此加法方法被命名为 $\mathrm{Plus}$ 而非 $\mathrm{Add}$,并通过投影属性 $\mathrm{sum}$ 暴露其结果;这里已指定类型,但尚未计算:
\else
The addition operation is split along the case axis into an \emph{interface}
mixin $\mathrm{NatPlus}$, which keeps the concern's name and declares the command
on the abstract $\mathrm{Product}$, and one \emph{implementation} mixin per case,
which appends the constructor name ($\mathrm{NatPlusZero}$ and
$\mathrm{NatPlusSuccessor}$ below). $\mathrm{NatPlus}$ inherits
$\mathrm{NatData}$. Where an
object-oriented method would be named with a verb, the command
pattern~\cite{gamma1994-design-patterns} reifies it as a noun-named command, so
the addition method is named $\mathrm{Plus}$ rather than $\mathrm{Add}$, and it
exposes its result through the projection property $\mathrm{sum}$, here typed but
not yet computed:
\fi
\[
  \mathrm{NatPlus} \mapsto \left\{
    \begin{aligned}
      &\mathrm{NatData},\\
      &\mathrm{NatFactory} \mapsto \left\{\mathrm{Product} \mapsto \left\{\mathrm{Plus} \mapsto \left\{\mathrm{sum} \mapsto \{\mathrm{Product}\}\right\}\right\}\right\}
  \end{aligned}\right\}
\]

\paragraph{NatPlusZero}
\ifInheritanceChinese
基础情形为 $0 + m = m$。这里 $\mathrm{NatPlusZero}$ 充当一个导入其他模块的模块:它继承 $\mathrm{Zero}$ 数据与 $\mathrm{Plus}$ 接口,并直接返回被加数:
\else
The base case is $0 + m = m$. Here $\mathrm{NatPlusZero}$ acts as a module that
imports other modules: it inherits the $\mathrm{Zero}$ data and the $\mathrm{Plus}$
interface, and returns the addend directly:
\fi
\[
  \begin{aligned}
    &\mathrm{NatPlusZero} \mapsto \\
    &\left\{
    \begin{aligned}
      &\mathrm{NatDataZero},\; \mathrm{NatPlus},\\
      &\mathrm{NatFactory} \mapsto \left\{\mathrm{Zero} \mapsto \left\{\mathrm{Plus} \mapsto \left\{
          \begin{aligned}
            &\mathrm{addend} \mapsto \{\qualifiedthis{\mathrm{NatPlusZero}}.\mathrm{NatFactory}.\mathrm{Product}\},\\
            &\mathrm{sum} \mapsto \{\mathrm{addend}\}
      \end{aligned}\right\}\right\}\right\}
  \end{aligned}\right\}
  \end{aligned}
\]

\paragraph{NatPlusSuccessor}
\ifInheritanceChinese
递归情形将 $\mathrm{S}(n_0) + m$ 规约为 $n_0 + \mathrm{S}(m)$,方法是以递增后的被加数委托给 $n_0$ 的 $\mathrm{Plus}$:
\else
The recursive case reduces $\mathrm{S}(n_0) + m$ to $n_0 + \mathrm{S}(m)$
by delegating to $n_0$'s $\mathrm{Plus}$ with an incremented addend:
\fi
\[
  \begin{aligned}
    &\mathrm{NatPlusSuccessor} \mapsto \\
    &\left\{
    \begin{aligned}
      &\mathrm{NatDataSuccessor},\; \mathrm{NatPlus},\\
      &\mathrm{NatFactory} \mapsto \\
      &\left\{\mathrm{Successor} \mapsto \left\{\mathrm{Plus} \mapsto \left\{
              \begin{aligned}
                &\mathrm{addend} \mapsto \{\qualifiedthis{\mathrm{NatFactory}}.\mathrm{Product}\},\\
                &\mathrm{\_increasedAddend} \mapsto  \\
                &\left\{
                  \begin{aligned}
                    &\mathrm{Successor},\\
                    &\mathrm{predecessor} \mapsto \{\mathrm{addend}\}
                \end{aligned}\right\},\\
                &\mathrm{\_recursiveAddition} \mapsto \\
                &\left\{
                  \begin{aligned}
                    &\qualifiedthis{\mathrm{Successor}}.\mathrm{predecessor}.\mathrm{Plus},\\
                    &\mathrm{addend} \mapsto \{\mathrm{\_increasedAddend}\}
                \end{aligned}\right\},\\
                &\mathrm{sum} \mapsto \{\mathrm{\_recursiveAddition}.\mathrm{sum}\}
          \end{aligned}\right\}\right\}\right\}
  \end{aligned}\right\}
  \end{aligned}
\]

\paragraph{NatVisitor}
\ifInheritanceChinese
相等判断需要对数据构造子进行情形分析;我们使用访问者模式~\cite{gamma1994-design-patterns}。$\mathrm{NatVisitor}$ 向 $\mathrm{Product}$ 添加一个 $\mathrm{Acceptance}$ 命令,即访问者模式的 $\mathrm{accept}$ 方法的名词形式。mixin $\mathrm{NatVisitor}$ 将其声明为空,留下 $\mathrm{VisitorMap}$ 与 $\mathrm{Accepted}$ 属性由各构造子的实现来填充:
\else
Equality requires case analysis on the data constructor; we use the
visitor pattern~\cite{gamma1994-design-patterns}.
$\mathrm{NatVisitor}$ adds an $\mathrm{Acceptance}$ command to $\mathrm{Product}$,
the noun form of the visitor pattern's $\mathrm{accept}$ method.
The mixin $\mathrm{NatVisitor}$ declares it empty, leaving a $\mathrm{VisitorMap}$ and an
$\mathrm{Accepted}$ property for the per-constructor implementations to fill in:
\fi
\[
  \begin{aligned}
    &\mathrm{NatVisitor} \mapsto \\
    &\left\{
    \begin{aligned}
      &\mathrm{NatData},\\
      &\mathrm{NatFactory} \mapsto \left\{\mathrm{Product} \mapsto \{\mathrm{Acceptance} \mapsto \{\mathrm{VisitorMap} \mapsto \{\},\; \mathrm{Accepted} \mapsto \{\}\}\}\right\}
  \end{aligned}\right\}
  \end{aligned}
\]

\paragraph{NatVisitorZero}
\ifInheritanceChinese
$\mathrm{Zero}$ 的实现选择 $\mathrm{ZeroVisitor}$ 分支。这里的访问者模式与~\cite{gamma1994-design-patterns} 中的表述不同,那里数据构造子由 $\mathrm{accept}$ 接口固定;此处通过组合允许添加新的构造子,因为每个构造子作为独立的 mixin 贡献其自身的分支:
\else
The $\mathrm{Zero}$ implementation selects the $\mathrm{ZeroVisitor}$ branch.
The visitor pattern here, unlike its formulation in
\cite{gamma1994-design-patterns} where the data constructors are fixed by the
$\mathrm{accept}$ interface, admits new constructors by composition, since each
constructor contributes its own branch as a separate mixin:
\fi
\[
  \begin{aligned}
    &\mathrm{NatVisitorZero} \mapsto \\
    &\left\{
    \begin{aligned}
      &\mathrm{NatDataZero},\; \mathrm{NatVisitor},\\
      &\mathrm{NatFactory} \mapsto \left\{\mathrm{Zero} \mapsto \left\{\mathrm{Acceptance} \mapsto \left\{
                \begin{aligned}
                  &\mathrm{VisitorMap} \mapsto \{\mathrm{ZeroVisitor} \mapsto \{\}\},\\
                  &\mathrm{Accepted} \mapsto \{\mathrm{VisitorMap}.\mathrm{ZeroVisitor}\}
              \end{aligned}\right\}\right\}\right\}
  \end{aligned}\right\}
  \end{aligned}
\]

\paragraph{NatVisitorSuccessor}
\ifInheritanceChinese
$\mathrm{Successor}$ 的实现选择 $\mathrm{SuccessorVisitor}$ 分支:
\else
The $\mathrm{Successor}$ implementation selects the $\mathrm{SuccessorVisitor}$ branch:
\fi
\[
  \begin{aligned}
    &\mathrm{NatVisitorSuccessor} \mapsto \\
    &\left\{
    \begin{aligned}
      &\mathrm{NatDataSuccessor},\; \mathrm{NatVisitor},\\
      &\mathrm{NatFactory} \mapsto \\
      &\left\{\mathrm{Successor} \mapsto \left\{
            \begin{aligned}
              &\mathrm{predecessor} \mapsto \{\qualifiedthis{\mathrm{NatFactory}}.\mathrm{Product}\},\\
              &\mathrm{Acceptance} \mapsto \left\{
                \begin{aligned}
                  &\mathrm{VisitorMap} \mapsto \{\mathrm{SuccessorVisitor} \mapsto \{\}\},\\
                  &\mathrm{Accepted} \mapsto \{\mathrm{VisitorMap}.\mathrm{SuccessorVisitor}\}
              \end{aligned}\right\}
          \end{aligned}\right\}\right\}
  \end{aligned}\right\}
  \end{aligned}
\]

\paragraph{BooleanData}
\ifInheritanceChinese
$\mathrm{BooleanData}$ 定义输出的域:
\else
$\mathrm{BooleanData}$ defines the output sort:
\fi
\[
  \mathrm{BooleanData} \mapsto \left\{
    \begin{aligned}
      &\mathrm{BooleanFactory} \mapsto \left\{
        \begin{aligned}
          &\mathrm{Product} \mapsto \{\},\\
          &\mathrm{True} \mapsto \{\mathrm{Product}\},\\
          &\mathrm{False} \mapsto \{\mathrm{Product}\}
      \end{aligned}\right\},\\
      &\mathrm{Boolean} \mapsto \{\mathrm{BooleanFactory}.\mathrm{Product}\}
  \end{aligned}\right\}
\]
\paragraph{NatEquality}
\label{sec:nat-equality}
\ifInheritanceChinese
为了验证 $2 + 3 = 5$,我们需要对 Nat 值进行相等判断。mixin $\mathrm{NatEquality}$ 导入 $\mathrm{NatVisitor}$ 与 $\mathrm{BooleanData}$,并在 $\mathrm{Product}$ 上声明一个 $\mathrm{Equal}$ 命令,由各构造子的实现提供。由于 $\mathrm{NatEquality}$ 继承 $\mathrm{BooleanData}$,限定 $\this$ $\qualifiedthis{\mathrm{NatEquality}}$ 可以到达 $\mathrm{Boolean}$,因此 $\mathrm{Equal}$ 以类型 $\qualifiedthis{\mathrm{NatEquality}}.\mathrm{Boolean}$ 声明其结果 $\mathrm{equal}$:
\else
To test that $2 + 3 = 5$, we needed equality on Nat values.
The mixin $\mathrm{NatEquality}$ imports $\mathrm{NatVisitor}$ and $\mathrm{BooleanData}$,
and declares an $\mathrm{Equal}$ command on $\mathrm{Product}$ that the per-constructor implementations provide.
Because $\mathrm{NatEquality}$ inherits $\mathrm{BooleanData}$,
the qualified this $\qualifiedthis{\mathrm{NatEquality}}$ can reach $\mathrm{Boolean}$,
so $\mathrm{Equal}$ declares its result
$\mathrm{equal}$ with the type $\qualifiedthis{\mathrm{NatEquality}}.\mathrm{Boolean}$:
\fi
\[
  \begin{aligned}
    &\mathrm{NatEquality} \mapsto \\
    &\left\{
    \begin{aligned}
      &\mathrm{NatVisitor},\; \mathrm{BooleanData},\\
      &\mathrm{NatFactory} \mapsto \left\{\mathrm{Product} \mapsto \left\{\mathrm{Equal} \mapsto \left\{
          \begin{aligned}
            &\mathrm{other} \mapsto \{\mathrm{Product}\},\\
            &\mathrm{equal} \mapsto \{\qualifiedthis{\mathrm{NatEquality}}.\mathrm{Boolean}\}
      \end{aligned}\right\}\right\}\right\}
  \end{aligned}\right\}
  \end{aligned}
\]

\paragraph{NatEqualityZero}
\ifInheritanceChinese
$\mathrm{Zero}$ 只与另一个 $\mathrm{Zero}$ 相等。这里的继承执行情形分析:继承 $\mathrm{other}.\mathrm{Acceptance}$ 对 $\mathrm{other}$ 是由 $\mathrm{Zero}$ 还是 $\mathrm{Successor}$ 构建进行分派,在 $\mathrm{ZeroVisitor}$ 分支返回 $\mathrm{True}$,否则返回 $\mathrm{False}$,即 $\mathrm{equal} = \{ \mathrm{True} \text{ 若 } \mathrm{other} = \mathrm{Zero},\; \mathrm{False} \text{ 若 } \mathrm{other} = \mathrm{Successor} \}$。$\mathrm{True}$ 与 $\mathrm{False}$ 均为 $\mathrm{BooleanFactory}$ 的数据构造子;该工厂来自 $\mathrm{BooleanData}$,通过继承的 $\mathrm{NatEquality}$ 间接导入,并以 $\qualifiedthis{\mathrm{NatEqualityZero}}.\mathrm{BooleanFactory}$ 到达。同一继承机制也提供了结果类型:在 $\mathrm{Accepted}$ 中,$\mathrm{equal}$ 继承 $\qualifiedthis{\mathrm{NatEqualityZero}}.\mathrm{Boolean}$,从而声明其类型为 $\mathrm{Boolean}$:
\else
$\mathrm{Zero}$ is equal only to another $\mathrm{Zero}$.
Inheriting here performs case analysis:
inheriting $\mathrm{other}.\mathrm{Acceptance}$ dispatches on whether $\mathrm{other}$
was built by $\mathrm{Zero}$ or $\mathrm{Successor}$, returning $\mathrm{True}$
on the $\mathrm{ZeroVisitor}$ branch and $\mathrm{False}$ otherwise, that is,
$\mathrm{equal} = \{ \mathrm{True} \text{ if } \mathrm{other} = \mathrm{Zero},\; \mathrm{False} \text{ if } \mathrm{other} = \mathrm{Successor} \}$.
Both $\mathrm{True}$ and $\mathrm{False}$ are data constructors of
$\mathrm{BooleanFactory}$, which comes from $\mathrm{BooleanData}$, imported
indirectly through the inherited $\mathrm{NatEquality}$, and is reached as
$\qualifiedthis{\mathrm{NatEqualityZero}}.\mathrm{BooleanFactory}$.
The same inheritance mechanism also supplies the result type: in
$\mathrm{Accepted}$, $\mathrm{equal}$ inherits
$\qualifiedthis{\mathrm{NatEqualityZero}}.\mathrm{Boolean}$, declaring its type to
be $\mathrm{Boolean}$:
\fi
{\small
  \[
    \begin{aligned}
      &\mathrm{NatEqualityZero} \mapsto \\
      &\left\{
      \begin{aligned}
        &\mathrm{NatVisitorZero},\; \mathrm{NatEquality},\\
        &\mathrm{NatFactory} \mapsto \\
        &\left\{
        \begin{aligned}
        &\mathrm{Zero} \mapsto \\
        &\left\{
        \begin{aligned}
        &\mathrm{Equal} \mapsto \\
        &\left\{
                  \begin{aligned}
                    &\mathrm{other} \mapsto \{\qualifiedthis{\mathrm{NatEqualityZero}}.\mathrm{NatFactory}.\mathrm{Product}\},\\
                    &\mathrm{\_OtherAcceptance} \mapsto \\
                    &\left\{
                      \begin{aligned}
                        &\mathrm{other}.\mathrm{Acceptance},\\
                        &\mathrm{VisitorMap} \mapsto \\
                        &\left\{
                          \begin{aligned}
                            &\mathrm{ZeroVisitor} \mapsto \{\mathrm{equal} \mapsto \{\qualifiedthis{\mathrm{NatEqualityZero}}.\mathrm{BooleanFactory}.\mathrm{True}\}\},\\
                            &\mathrm{SuccessorVisitor} \mapsto \{\mathrm{equal} \mapsto \{\qualifiedthis{\mathrm{NatEqualityZero}}.\mathrm{BooleanFactory}.\mathrm{False}\}\}
                        \end{aligned}\right\},\\
                        &\mathrm{Accepted} \mapsto \{\mathrm{equal} \mapsto \{\qualifiedthis{\mathrm{NatEqualityZero}}.\mathrm{Boolean}\}\}
                    \end{aligned}\right\},\\
                    &\mathrm{equal} \mapsto \{\mathrm{\_OtherAcceptance}.\mathrm{Accepted}.\mathrm{equal}\}
                \end{aligned}\right\}
        \end{aligned}\right\}
        \end{aligned}\right\}
    \end{aligned}\right\}
    \end{aligned}
\]}

\paragraph{NatEqualitySuccessor}
\ifInheritanceChinese
$\mathrm{Successor}$ 只与另一个具有相等前驱的 $\mathrm{Successor}$ 相等,因此其实现通过 $\qualifiedthis{\mathrm{Successor}}.\mathrm{predecessor}.\mathrm{Equal}$ 进行递归:
\else
$\mathrm{Successor}$ is equal only to another $\mathrm{Successor}$ with an equal
predecessor, so its implementation recurses through
$\qualifiedthis{\mathrm{Successor}}.\mathrm{predecessor}.\mathrm{Equal}$:
\fi
{\small
  \[
    \begin{aligned}
      &\mathrm{NatEqualitySuccessor} \mapsto \\
      &\left\{
      \begin{aligned}
        &\mathrm{NatVisitorSuccessor},\; \mathrm{NatEquality},\\
        &\mathrm{NatFactory} \mapsto \\
        &\left\{
        \begin{aligned}
        &\mathrm{Successor} \mapsto \\
        &\left\{
        \begin{aligned}
        &\mathrm{Equal} \mapsto \\
        &\left\{
                  \begin{aligned}
                    &\mathrm{other} \mapsto \left\{
                      \begin{aligned}
                        &\qualifiedthis{\mathrm{NatEqualitySuccessor}}.\mathrm{NatFactory}.\mathrm{Product},\\
                        &\mathrm{predecessor} \mapsto \{\qualifiedthis{\mathrm{NatEqualitySuccessor}}.\mathrm{NatFactory}.\mathrm{Product}\}
                    \end{aligned}\right\},\\
                    &\mathrm{\_recursiveEquality} \mapsto \left\{
                      \begin{aligned}
                        &\qualifiedthis{\mathrm{Successor}}.\mathrm{predecessor}.\mathrm{Equal},\\
                        &\mathrm{other} \mapsto \{\qualifiedthis{\mathrm{Equal}}.\mathrm{other}.\mathrm{predecessor}\}
                    \end{aligned}\right\},\\
                    &\mathrm{\_OtherAcceptance} \mapsto \\
                    &\left\{
                      \begin{aligned}
                        &\mathrm{other}.\mathrm{Acceptance},\\
                        &\mathrm{VisitorMap} \mapsto \\
                        &\left\{
                          \begin{aligned}
                            &\mathrm{ZeroVisitor} \mapsto \{\mathrm{equal} \mapsto \{\qualifiedthis{\mathrm{NatEqualitySuccessor}}.\mathrm{BooleanFactory}.\mathrm{False}\}\},\\
                            &\mathrm{SuccessorVisitor} \mapsto \{\mathrm{equal} \mapsto \{\mathrm{\_recursiveEquality}.\mathrm{equal}\}\}
                        \end{aligned}\right\},\\
                        &\mathrm{Accepted} \mapsto \{\mathrm{equal} \mapsto \{\qualifiedthis{\mathrm{NatEqualitySuccessor}}.\mathrm{Boolean}\}\}
                    \end{aligned}\right\},\\
                    &\mathrm{equal} \mapsto \{\mathrm{\_OtherAcceptance}.\mathrm{Accepted}.\mathrm{equal}\}
                \end{aligned}\right\}
        \end{aligned}\right\}
        \end{aligned}\right\}
    \end{aligned}\right\}
    \end{aligned}
\]}
%% === 审稿人备忘：关于代码"冗长"的回应 ===
%%
%% NatEquality 的定义看起来比等价的 Haskell 或 ML 模式匹配更长，
%% 但这种冗长与 point-free / 组合子风格的"简洁"形成的对比是表面的。
%% 需要区分两种"长"：
%%
%% 1. 自描述式的长：每个中间步骤都有名字（如 _recursiveEquality、
%%    OtherAcceptance），步骤写得详细。名字本身是零成本的语义锚点。
%% 2. 组合子膨胀式的长：如 SKI bracket abstraction 的指数级膨胀，
%%    在玩具例子里看起来精炼，但在真实程序中复杂度是 λ-calculus 的
%%    多项式甚至指数倍，而且没有涌现出新的计算模式（论文 Related Work
%%    已指出这一点）。
%%
%% Inheritance-calculus 的"冗长"属于第一种，而且是结构上的必要代价：
%% 每个有名字的中间步骤既是代码，也是 mixin 树上可路径寻址、可扩展的
%% 节点。这正是 open extension 的前提——无法扩展匿名的子表达式。
%% 如果加语法糖把名字隐藏（如同 M-表达式之于 S-表达式），反而会
%% 破坏这一性质。Lisp 社区的历史经验表明，S-表达式的"冗长"最终
%% 被证明是优势（代码即数据，宏系统直接操作统一结构），M-表达式
%% 的语法糖从未被真正需要。
%%
%% 在 AI 辅助编程的语境下，这种显式命名的风格尤其有利：
%% _recursiveEquality 这样的名字对 LLM 而言是自带的 chain-of-thought，
%% 比起让 AI 猜测 point-free 管道中第 n 个组合子的意图，
%% 每个自描述的步骤都降低了生成和理解的认知负荷。
%% 考虑到未来越来越多的代码将由 AI 生成和维护，
%% 多写几个 token 的显式名称不是成本，而是投资。
%%
%% 当然，如果确实需要更紧凑的表面语法，为 inheritance-calculus 加一个
%% ANF 预处理器（自动为匿名子表达式生成名字）并不困难，
%% 正如 ANF 变换本身是机械化的。但这是表面语法的选择，
%% 不影响演算本身的语义。
%% ================================================================
\ifInheritanceChinese
通过继承来组合 $\mathrm{Plus}$ 与 $\mathrm{Equal}$ 的实现,无需修改任何 mixin,即可使所得的 Nat 值自动同时具备 $\mathrm{Plus}$ 与 $\mathrm{Equal}$。具体的数字常量存放在一个只继承两个数据构造子的共享 mixin 中:
\else
Composing the $\mathrm{Plus}$ and $\mathrm{Equal}$ implementations by
inheriting them, with no modifications to any mixin, gives the
resulting Nat values both $\mathrm{Plus}$ and $\mathrm{Equal}$
automatically.
The concrete numerals live in a shared mixin that inherits only the
two data constructors:
\fi
\[
  \mathrm{NatConstants} \mapsto \left\{
    \begin{aligned}
      &\mathrm{NatDataZero},\; \mathrm{NatDataSuccessor},\\
      &\mathrm{One} \mapsto \left\{
        \begin{aligned}
          &\qualifiedthis{\mathrm{NatConstants}}.\mathrm{NatFactory}.\mathrm{Successor},\\
          &\mathrm{predecessor} \mapsto \{\qualifiedthis{\mathrm{NatConstants}}.\mathrm{NatFactory}.\mathrm{Zero}\}
      \end{aligned}\right\},\\
      &\mathrm{Two} \mapsto \left\{
        \begin{aligned}
          &\qualifiedthis{\mathrm{NatConstants}}.\mathrm{NatFactory}.\mathrm{Successor},\\
          &\mathrm{predecessor} \mapsto \{\mathrm{One}\}
      \end{aligned}\right\},\\
      &\mathrm{Three} \mapsto \left\{
        \begin{aligned}
          &\qualifiedthis{\mathrm{NatConstants}}.\mathrm{NatFactory}.\mathrm{Successor},\\
          &\mathrm{predecessor} \mapsto \{\mathrm{Two}\}
      \end{aligned}\right\},\\
      &\mathrm{Four} \mapsto \left\{
        \begin{aligned}
          &\qualifiedthis{\mathrm{NatConstants}}.\mathrm{NatFactory}.\mathrm{Successor},\\
          &\mathrm{predecessor} \mapsto \{\mathrm{Three}\}
      \end{aligned}\right\},\\
      &\mathrm{Five} \mapsto \left\{
        \begin{aligned}
          &\qualifiedthis{\mathrm{NatConstants}}.\mathrm{NatFactory}.\mathrm{Successor},\\
          &\mathrm{predecessor} \mapsto \{\mathrm{Four}\}
      \end{aligned}\right\}
  \end{aligned}\right\}
\]
\ifInheritanceChinese
算术测试继承 $\mathrm{NatConstants}$ 以及它所需的操作:
\else
The arithmetic test inherits $\mathrm{NatConstants}$ together with the operations it needs:
\fi
\[
  \mathrm{Test} \mapsto \left\{
    \begin{aligned}
      &\mathrm{NatConstants},\; \mathrm{NatPlusZero},\; \mathrm{NatPlusSuccessor},\\
      &\mathrm{NatEqualityZero},\; \mathrm{NatEqualitySuccessor},\\
      &\mathrm{TwoPlusThree} \mapsto \left\{
        \begin{aligned}
          &\qualifiedthis{\mathrm{Test}}.\mathrm{Two}.\mathrm{Plus},\\
          &\mathrm{addend} \mapsto \{\qualifiedthis{\mathrm{Test}}.\mathrm{Three}\}
      \end{aligned}\right\},\\
      &\mathrm{FiveEqualTwoPlusThree} \mapsto \left\{
        \begin{aligned}
          &\qualifiedthis{\mathrm{Test}}.\mathrm{Five}.\mathrm{Equal},\\
          &\mathrm{other} \mapsto \{\mathrm{TwoPlusThree}.\mathrm{sum}\}
      \end{aligned}\right\}
  \end{aligned}\right\}
\]
\ifInheritanceChinese
路径 $\mathrm{Test}.\mathrm{FiveEqualTwoPlusThree}.\mathrm{equal}$ 继承 $\mathrm{True}$。%
\footnote{
  读取一个结果,需要借助 FFI 来打印,或者借助元语言观察以在该路径上调用 $\supers$。
}
没有任何现有 mixin 被修改。$\mathrm{NatEquality}$(一项新操作)与 $\mathrm{Boolean}$(一种新数据类型)各自只需要一个新 mixin:沿操作轴与情形轴均可扩展。
\else
The path $\mathrm{Test}.\mathrm{FiveEqualTwoPlusThree}.\mathrm{equal}$
inherits $\mathrm{True}$.%
\footnote{
  Reading a result requires either the FFI to print it or a
  metalanguage observation that invokes $\supers$ on the path.
}
No existing mixin was modified.
$\mathrm{NatEquality}$, a new operation, and $\mathrm{Boolean}$,
a new data type, each required only a new mixin: extensibility along both
the operation axis and the case axis.
\fi

\paragraph{\bilingual{Tree-level inheritance and return types}{树层次的继承与返回类型}}
\label{sec:discussion-expression-problem}
\ifInheritanceChinese
每个操作都是一个独立的 mixin,扩展工厂的子树。继承递归地合并共享同一标签的子树,因此工厂中的每个标签都从所有被继承的 mixin 获得全部操作。关键在于,这同样适用于返回类型:$\mathrm{NatPlus}$ 从 $\mathrm{NatFactory}$ 返回值,而 $\mathrm{NatEquality}$ 独立地向同一工厂添加 $\mathrm{Equal}$。同时继承两者,使得 $\mathrm{Plus}$ 返回的值自动携带 $\mathrm{Equal}$,无需修改任何一个 mixin。在对象代数~\cite{oliveira2012-object-algebras} 与 finally tagless 解释器~\cite{carette2009-finally-tagless} 中,独立定义的操作不会自动丰富彼此的返回类型;需要额外的样板代码或类型类机制。

多目标限定 $\this$(方程~\ref{eq:this})至关重要。在上述测试 mixin 中,$\mathrm{NatPlusSuccessor}$ 与 $\mathrm{NatEqualitySuccessor}$ 均继承 $\mathrm{NatDataSuccessor}$。当两个操作被组合时,所得的 $\mathrm{Successor}$ 通过两条独立路由继承 $\mathrm{NatDataSuccessor}$ 的模式。每个操作内部的引用 $\qualifiedthis{\mathrm{Successor}}.\mathrm{predecessor}$ 通过两条路由解析,由于继承是幂等的,所有路由贡献相同的属性。在 Scala 中,同样的模式,即两个独立定义的内部类扩展同一外部类的特质,会触发``冲突基类型''拒绝(附录~\ref{app:scala-multi-target}),因为 DOT 的 self 变量是单值的,无法同时通过多条继承路由解析。表达式问题在 Scala 中仍可使用对象代数或类似模式求解,但组合独立定义的操作本质上会产生上述对共享模式的多目标 $\this$ 解析。
\else
Each operation is an independent mixin that extends a factory's
subtree.
Inheritance recursively merges subtrees that share a label,
so every label in the factory acquires all operations from all
inherited mixins.
Crucially, this applies to return types: $\mathrm{NatPlus}$ returns
values from $\mathrm{NatFactory}$, and $\mathrm{NatEquality}$
independently adds $\mathrm{Equal}$ to the same factory.
Inheriting both causes the values returned by $\mathrm{Plus}$ to
carry $\mathrm{Equal}$ automatically, without modifying either mixin.
In object algebras~\cite{oliveira2012-object-algebras} and finally
tagless interpreters~\cite{carette2009-finally-tagless}, independently defined
operations do not automatically enrich each other's return types;
additional boilerplate or type-class machinery is required.

Multi-target qualified $\this$ (equation~\ref{eq:this}) is essential.
In the test mixin above, $\mathrm{NatPlusSuccessor}$ and
$\mathrm{NatEqualitySuccessor}$ both inherit
$\mathrm{NatDataSuccessor}$.
When the two operations are composed, the resulting
$\mathrm{Successor}$ inherits the $\mathrm{NatDataSuccessor}$ schema
through two independent routes.
The reference
$\qualifiedthis{\mathrm{Successor}}.\mathrm{predecessor}$
inside each operation resolves through both routes, and since
inheritance is idempotent, all routes contribute the same
properties.
In Scala, the same pattern, where two independently defined inner
classes extend the same outer class's trait, triggers a
``conflicting base types'' rejection
(Appendix~\ref{app:scala-multi-target}), because DOT's self variable
is single-valued and cannot resolve through multiple inheritance
routes simultaneously.
The Expression Problem can still be solved in Scala using
object algebras or similar patterns, but composing independently
defined operations inherently produces the multi-target $\this$
resolution to shared schemas just described.
\fi

\paragraph{\bilingual{Church-encoded Nats are tries}{Church 编码的 Nat 是字典树}}
\label{sec:case-study-cartesian}
\ifInheritanceChinese
Church 编码的 Nat 是一个记录,其结构映射了构建它所用的数据构造子路径:$\mathrm{Zero}$ 是一个平坦记录;$\mathrm{S}(\mathrm{Zero})$ 是一个带有 $\mathrm{predecessor}$ 属性、指向一个 $\mathrm{Zero}$ 记录的记录;依此类推。这恰好是字典树的形状。%
\footnote{
  一元 Nat 的字典树可以被视为一个稀疏链表,其深度等于它所代表的数值:在字典树并中,链表中的某些 $\mathrm{Successor}$ 节点指向 $\mathrm{Zero}$(叶节点),因此链表中并非每个位置都被占用。\anon[{} 补充材料]{{} MIXINv2 源码~\cite{mixin2025}}{} 中包含一个字典树深度为对数级的二进制自然数(BinNat)算术。
}
一个同时继承两个 Nat 值的作用域是一个字典树并。下面的测试将 $\mathrm{Plus}$ 应用于 $\{1,2\}$ 的字典树并,被加数为 $\{3,4\}$ 的字典树并,然后用 $\mathrm{Equal}$ 将结果与自身比较:
\else
A Church-encoded Nat is a record whose structure mirrors the
data-constructor path used to build it: $\mathrm{Zero}$ is a flat record;
$\mathrm{S}(\mathrm{Zero})$ is a record with a $\mathrm{predecessor}$
property pointing to a $\mathrm{Zero}$ record; and so on.
This is exactly the shape of a trie.%
\footnote{
  The unary Nat trie can be viewed as a sparse linked list
  whose depth equals the number it represents:
  in a trie union, some $\mathrm{Successor}$ nodes along
  the list point to $\mathrm{Zero}$ (leaf nodes),
  so not every position in the list is occupied.
  The\anon[{} supplementary material]{{} MIXINv2 source~\cite{mixin2025}}{} includes a binary natural number
  (BinNat) arithmetic whose trie depth is logarithmic.
}
A scope that simultaneously inherits two Nat values is a trie union.
The following test applies $\mathrm{Plus}$ to a trie union of $\{1,2\}$
with addend a trie union of $\{3,4\}$, then compares the result to itself
with $\mathrm{Equal}$:
\fi
\[
  \mathrm{CartesianTest} \mapsto \left\{
    \begin{aligned}
      &\mathrm{NatConstants},\; \mathrm{NatPlusZero},\; \mathrm{NatPlusSuccessor},\\
      &\mathrm{NatEqualityZero},\; \mathrm{NatEqualitySuccessor},\\
      &\mathrm{OneOrTwo} \mapsto \left\{\qualifiedthis{\mathrm{CartesianTest}}.\mathrm{One},\; \qualifiedthis{\mathrm{CartesianTest}}.\mathrm{Two}\right\},\\
      &\mathrm{ThreeOrFour} \mapsto \left\{\qualifiedthis{\mathrm{CartesianTest}}.\mathrm{Three},\; \qualifiedthis{\mathrm{CartesianTest}}.\mathrm{Four}\right\},\\
      &\mathrm{Result} \mapsto \left\{\mathrm{OneOrTwo}.\mathrm{Plus},\; \mathrm{addend} \mapsto \{\mathrm{ThreeOrFour}\}\right\},\\
      &\mathrm{Check} \mapsto \left\{\mathrm{Result}.\mathrm{sum}.\mathrm{Equal},\; \mathrm{other} \mapsto \{\mathrm{Result}.\mathrm{sum}\}\right\}
  \end{aligned}\right\}
\]
\ifInheritanceChinese
$\mathrm{Result}.\mathrm{sum}$ 求值为字典树 $\{4,5,6\}$,即所有两两之和的集合。$\mathrm{Equal}$ 随后对从 $\{4,5,6\} \times \{4,5,6\}$ 中取出的每一对进行应用;路径 $\mathrm{CartesianTest}.\mathrm{Check}.\mathrm{equal}$ 同时继承 $\mathrm{True}$ 与 $\mathrm{False}$。笛卡尔积直接来自语义方程。$\mathrm{OneOrTwo}$ 同时继承 $\mathrm{One}$ 与 $\mathrm{Two}$,因此 $\bases(\mathrm{OneOrTwo})$ 包含两条路径。当 $\mathrm{Result}$ 继承 $\mathrm{OneOrTwo}.\mathrm{Plus}$ 时,$\mathrm{Result}$ 的 $\bases$ 包含 $\mathrm{One}$ 与 $\mathrm{Two}$ 各自的 $\mathrm{Plus}$。每个 $\mathrm{Plus}$ 通过 $\this$ 解析其 $\mathrm{addend}$,到达 $\mathrm{ThreeOrFour}$;该值本身是 $\mathrm{Three}$ 与 $\mathrm{Four}$ 的字典树并,因此 $\mathrm{addend}$ 路径的 $\supers$ 包含两个值。由于 $\mathrm{Successor}.\mathrm{Plus}$ 以递增后的被加数递归委托给 $\mathrm{predecessor}.\mathrm{Plus}$,而 $\mathrm{predecessor}$ 也通过字典树并解析,递归在每个字典树节点处分支。因此结果 $\mathrm{sum}$ 收集了通过两个操作数的任意组合所能到达的全部路径,即笛卡尔积。

同样操作单个值的 $\mathrm{Plus}$ 与 $\mathrm{Equal}$ 文件,无需修改,即产生逻辑编程的关系语义:$\overrides$ 与 $\supers$ 收集所有继承路径,每个操作都在其上分布。
\else
$\mathrm{Result}.\mathrm{sum}$ evaluates to the trie $\{4,5,6\}$, the set of all pairwise sums.
$\mathrm{Equal}$ then applies to every pair drawn from $\{4,5,6\} \times \{4,5,6\}$;
the path $\mathrm{CartesianTest}.\mathrm{Check}.\mathrm{equal}$
inherits both $\mathrm{True}$ and $\mathrm{False}$.
The Cartesian product arises directly from the semantic equations.
$\mathrm{OneOrTwo}$ inherits both $\mathrm{One}$ and $\mathrm{Two}$,
so $\bases(\mathrm{OneOrTwo})$ contains two paths.
When $\mathrm{Result}$ inherits
$\mathrm{OneOrTwo}.\mathrm{Plus}$, the $\bases$ of
$\mathrm{Result}$ include the $\mathrm{Plus}$ of both
$\mathrm{One}$ and $\mathrm{Two}$.
Each $\mathrm{Plus}$ resolves its $\mathrm{addend}$ via
$\this$, reaching $\mathrm{ThreeOrFour}$, which is itself a
trie union of $\mathrm{Three}$ and $\mathrm{Four}$:
the $\supers$ of the $\mathrm{addend}$ path contain both values.
Since $\mathrm{Successor}.\mathrm{Plus}$ recursively delegates to
$\mathrm{predecessor}.\mathrm{Plus}$ with an incremented addend,
and $\mathrm{predecessor}$ also resolves through the trie union,
the recursion branches at every trie node.
The result $\mathrm{sum}$ therefore collects all paths reachable
through any combination of the two operands, which is the
Cartesian product.

The same $\mathrm{Plus}$ and $\mathrm{Equal}$ files that operate on
single values thus produce, without modification, the relational
semantics of logic programming:
$\overrides$ and $\supers$ collect all inheritance paths, and
every operation distributes over them.
\fi

\section{\bilingual{Discussion}{讨论}}

\subsection{\bilingual{Emergent Phenomena}{涌现现象}}
\label{sec:emergent-phenomena}

\ifInheritanceChinese
在 MIXINv2 中开发真实程序的过程中，我们发现了三种现象，每一种都是对某种已有设计模式的非预期用法，超出了该模式的原始能力。

\begin{description}
  \item[关系语义]
    第~\ref{sec:case-study}~节案例研究中的笛卡尔积之所以出现，是因为字典树是一种集合结构，因此对字典树的深度合并即为集合并，而对合并后字典树的操作则遍历其操作数的笛卡尔积。
    这并不局限于 $\mathrm{Nat}$：任何 Datalog 元组都可编码为链表，而深度合并在这些元组上计算集合操作，我们认为这正是 Datalog 语义。
    附录~\ref{app:datalog-encoding}~将 Datalog 编码进继承演算，仅用第~\ref{sec:syntax}~节的三个基本构造即可表达变量连接、常量匹配、多体规则和递归规则。
    这种逻辑编程语义是对编码案例研究算术的工厂模式的非预期用法中涌现出来的。

  \item[对不可扩展性的免疫]\label{item:immunity}
    在第~\ref{sec:case-study}~节的案例研究中，$\mathrm{Nat}$ 编码被证明是无类型表达问题~\cite{wadler1998-expression-problem}的一个实例：%
    \footnote{
      表达问题有两个要求：沿操作轴和数据构造子轴均可扩展，以及静态类型安全。无类型演算无法满足类型安全要求，因此我们只考虑双轴可扩展性要求。
    }
    一个新操作和一个新数据类型各自作为独立的 mixin 添加，沿操作轴和情形轴均进行了扩展，且无需修改任何已有 mixin。
    在其他场景中，访问者模式允许添加新操作但不允许添加新情形；而这里两者都允许，因为第~\ref{sec:syntax}~节的每个构造都可通过深度合并（公式~\ref{eq:overrides}）进行扩展，所以不可扩展的那一半根本无法写出；第~\ref{sec:asymmetry}~节（定理~\ref{thm:cartesian-nonexpressibility}）则形式化了使这一结论无条件成立的不对称性。
    在其他地方恢复第二个扩展维度需要专门的机制~\cite{liang1995-monad-transformers, lammel2003-scrap-your-boilerplate, loh2006-open-data-types,
      swierstra2008-data-types-a-la-carte, carette2009-finally-tagless, oliveira2012-object-algebras,
      kiselyov2013-extensible-effects, wu2014-effect-handlers-in-scope, kiselyov2015-freer-monads,
    wang2016-expression-problem-trivially, leijen2017-row-typed-algebraic-effects, poulsen2023-hefty-algebras}。
    第~\ref{sec:semantic-variants}~节分析了为何深度合并是唯一能实现这一点的组合机制。
    这第二个可扩展性维度是对访问者模式的非预期用法中涌现出来的。

\end{description}

\label{sec:practical-function-color}%
更多现象依赖于 FFI，而本文中的示例排除了 FFI。其中一种是函数颜色盲~\cite{nystrom2015-function-color}：同一个 MIXINv2 文件可以在同步和异步运行时下运行，这是通过继承进行依赖注入的非预期用法。补充材料\anon[]{~\cite{mixin2025}}给出了相关示例。
\else
Developing real programs in MIXINv2 surfaced three phenomena, each an
off-label use of an adopted design pattern beyond its original abilities.

\begin{description}
  \item[Relational semantics]
    The Cartesian product in the case study of
    Section~\ref{sec:case-study} arose because a trie is a set
    structure, so deep merge over tries is set union and an operation
    over the merged trie ranges over the product of its operands.
    This is not specific to $\mathrm{Nat}$: any Datalog tuple encodes
    as a linked list, and deep merge computes set operations over such
    tuples, which we believe is Datalog semantics.
    Appendix~\ref{app:datalog-encoding} encodes Datalog into
    inheritance-calculus, expressing variable joins,
    constant matches, multi-body rules, and recursive rules with
    the three primitives of Section~\ref{sec:syntax} alone.
    This logic-programming semantics emerged from an off-label use of
    the factory pattern that encodes the case study's arithmetic.

  \item[Immunity to nonextensibility]\label{item:immunity}
    In the case study of Section~\ref{sec:case-study}, the
    $\mathrm{Nat}$ encoding turned out to be an instance of the
    untyped Expression
    Problem~\cite{wadler1998-expression-problem}:%
    \footnote{
      The Expression Problem has two requirements: extensibility along
      both the operation and the data-constructor axis, and static
      type safety. An untyped calculus cannot meet the type-safety
      requirement, so we consider only the two-axis extensibility
      requirement.
    }
    a new operation and a new data type were each added as a separate
    mixin, extending along both the operation axis and the case axis,
    with no existing mixin modified.
    Elsewhere the visitor pattern admits new operations but not new
    cases; here it admits both, because every construct of
    Section~\ref{sec:syntax} is extensible via deep merge
    (equation~\ref{eq:overrides}), so the nonextensible half cannot be
    written, and Section~\ref{sec:asymmetry}
    (Theorem~\ref{thm:cartesian-nonexpressibility}) formalizes the
    asymmetry that makes this unconditional.
    Recovering that second dimension elsewhere takes dedicated
    mechanisms~\cite{liang1995-monad-transformers, lammel2003-scrap-your-boilerplate, loh2006-open-data-types,
      swierstra2008-data-types-a-la-carte, carette2009-finally-tagless, oliveira2012-object-algebras,
      kiselyov2013-extensible-effects, wu2014-effect-handlers-in-scope, kiselyov2015-freer-monads,
    wang2016-expression-problem-trivially, leijen2017-row-typed-algebraic-effects, poulsen2023-hefty-algebras}.
    Section~\ref{sec:semantic-variants} analyzes why deep merge
    is the only composition mechanism that achieves this.
    This second extensibility dimension emerged from an off-label use
    of the visitor pattern.

\end{description}

\label{sec:practical-function-color}%
Further phenomena rely on the FFI, which the examples in this paper
exclude. One is function color
blindness~\cite{nystrom2015-function-color}: the same MIXINv2 file runs
under both synchronous and asynchronous runtimes, an off-label use of
dependency injection via inheritance. The
supplementary material\anon[]{~\cite{mixin2025}} gives the example.
\fi

\subsection{\bilingual{Semantic Alternatives}{语义替代方案}}
\label{sec:semantic-variants}

\ifInheritanceChinese
上述\hyperref[item:immunity]{对不可扩展性的免疫}条目确立了继承演算对不可扩展性的免疫性。我们将此作为设计目标，并通过调查我们已知的所有候选方案，询问是否有替代设计能共享这一性质。

\paragraph{\bilingual{Degrees of freedom}{自由度}}
以下分析假设 AST 是一棵\emph{命名树}：从父节点到子节点的每条边都携带一个标签，因此每个节点都由从根节点出发的路径唯一标识。对于配置语言，这自然成立：JSON、YAML 或 Nix 属性集中的每个键都是一个标签，而嵌套产生路径。%
\footnote{
  Nix 是一门具有一等函数的函数式语言；仅凭其属性集并不构成命名树 DSL。然而，NixOS 模块系统~\cite{nixosmodules}将用户界面限制为通过递归合并组合的嵌套属性集，而正是这种 DSL 层面的结构自然地构成了命名树。
}
对于主流编程语言，AST 并不直接是命名树，因为匿名表达式位置和隐式作用域缺乏标签。然而，经过 ANF 变换~\cite{flanagan1993-essence-compiling-continuations}后，每个中间结果都绑定到一个名字，而剩余的匿名位置%
\footnote{
  例如，$\mathbf{let}$ 的主体。
}
可以被分配合成的新鲜标签，例如第~\ref{sec:forward-translation}~节中的 $\mathrm{result}$ 和 $\mathrm{tailCall}$，从而得到一棵命名树。

在此前提下，每种程序是带引用的命名树的语言，都可以由两个有限集来刻画：一个由\emph{节点类型}构成的集合 $\mathcal{N}$，每种节点类型决定子树如何组织其子节点；以及一个由\emph{引用类型}构成的集合 $\mathcal{R}$，每种引用类型决定树中一个位置如何引用另一个位置。语法由 $\mathcal{N}$ 和 $\mathcal{R}$ 的选择固定；语义则由每种节点类型如何构造组合以及每种引用类型如何解析其目标来决定。

{\small
  \begin{center}
    \begin{tabular}{lp{0.38\columnwidth}p{0.33\columnwidth}}
      \bilingual{\textbf{Language}}{\textbf{语言}}
      & \bilingual{\textbf{Node types $\mathcal{N}$}}{\textbf{节点类型 $\mathcal{N}$}}
      & \bilingual{\textbf{Reference types $\mathcal{R}$}}{\textbf{引用类型 $\mathcal{R}$}} \\
      \hline
      Java (OOP)
      & \bilingual{class, interface, method,
      block, conditional, loop,
      \texttt{throw}, lambda, \ldots}{类、接口、方法、块、条件、循环、\texttt{throw}、lambda、\ldots}
      & \bilingual{variable, field access, method call,
      \texttt{extends}/\texttt{implements},
      \texttt{import}, \ldots}{变量、字段访问、方法调用、\texttt{extends}/\texttt{implements}、\texttt{import}、\ldots} \\[3pt]
      Haskell (FP)
      & \bilingual{$\lambda$, application, \texttt{let},
      \texttt{case}, \texttt{data},
      type class, instance, \ldots}{$\lambda$、应用、\texttt{let}、\texttt{case}、\texttt{data}、类型类、实例、\ldots}
      & \bilingual{variable, constructor,
      type class dispatch,
      \texttt{import}, \ldots}{变量、构造子、类型类分派、\texttt{import}、\ldots} \\[3pt]
      Scala\footnote{
        \bilingual{Scala is a multi-paradigm language.}{Scala 是一门多范式语言。}
      }
      & \bilingual{class, trait, object, method,
      block, conditional, \texttt{match},
      \texttt{while}, \texttt{throw}, \ldots}{类、trait、对象、方法、块、条件、\texttt{match}、\texttt{while}、\texttt{throw}、\ldots}
      & \bilingual{variable, member access,
      \texttt{extends}, \texttt{with},
      implicits, \texttt{import}, \ldots}{变量、成员访问、\texttt{extends}、\texttt{with}、隐式参数、\texttt{import}、\ldots} \\
      \hline
      $\lambda$-\bilingual{calculus}{演算}
      & \bilingual{abstraction, application}{抽象、应用}
      & \bilingual{variable}{变量} \\
      \bilingual{Inheritance-calculus}{继承演算}
      & \bilingual{record}{记录}
      & \bilingual{inheritance}{继承} \\
    \end{tabular}
  \end{center}
}

\paragraph{\bilingual{Immunity to nonextensibility}{对不可扩展性的免疫}}
程序 $P$ 中的路径 $p$ 是\emph{可扩展的}，若存在程序 $Q$，使得通过引用将 $P$ 与 $Q$ 组合后，在不修改 $P$ 的前提下，$p$ 处能够观察到新的属性。一种语言\emph{对不可扩展性免疫}，若其每个良构程序的每条路径都是可扩展的。

我们现在对设计空间进行分类，以解释为何我们选择第~\ref{sec:mixin-trees}~节的语义。两个维度保持自由：\emph{组合机制}，决定单一引用类型如何合并定义；以及\emph{引用解析}，决定单一引用类型如何找到其目标。

\subsubsection*{\bilingual{The composition mechanism}{组合机制}}

每种组合机制决定了一个抽象的定义如何被合并进另一个抽象。下表按各候选方案是否满足可交换性~(C)、幂等性~(I) 和可结合性~(A)，以及其可扩展性模型对已知候选方案进行分类。

\begin{center}
  \begin{tabular}{lcccll}
    \bilingual{\textbf{Candidate}}{\textbf{候选方案}} & \textbf{C} & \textbf{I} & \textbf{A}
    & \bilingual{\textbf{Extensibility}}{\textbf{可扩展性}} & \bilingual{\textbf{Representative}}{\textbf{代表}} \\
    \hline
    \bilingual{Deep merge}{深度合并}
    & $\checkmark$ & $\checkmark$ & $\checkmark$
    & \bilingual{universal}{通用} & \bilingual{this paper}{本文} \\
    \bilingual{Shallow merge}{浅层合并}
    & $\times$ & $\checkmark$ & $\checkmark$
    & \bilingual{top-level only}{仅顶层} & JS spread \\
    \bilingual{Conflict rejection}{冲突拒绝}
    & $\checkmark$ & $\checkmark$ & $\checkmark$
    & \bilingual{none on overlap}{重叠时无扩展}
    & Harper--Pierce~\cite{harper1991-symmetric-record-concatenation} \\
    \bilingual{Priority merge}{优先级合并}
    & $\times$ & $\times$ & $\checkmark$
    & \bilingual{universal}{通用} & Jsonnet~\texttt{+}, Nix~\texttt{//} \\
    $\beta$-\bilingual{reduction}{归约}
    & $\times$ & $\times$ & $\times$
    & \bilingual{opt-in}{按需选择} & Java, Haskell, Scala, \ldots \\
    \bilingual{Conditional}{条件并入}
    & \multicolumn{3}{c}{\bilingual{causes divergence}{导致发散}}
    & --- & --- \\
  \end{tabular}
\end{center}

\begin{description}
  \item[\bilingual{Deep merge}{深度合并}]
    对同名定义的合并递归进入其子树（$\overrides$，公式~\ref{eq:overrides}）。每个标签都是一个扩展点，无论原作者是否预期了扩展。

  \item[\bilingual{Shallow merge}{浅层合并}]
    嵌套标签无法被独立扩展。

  \item[\bilingual{Conflict rejection}{冲突拒绝}]
    同名定义被作为错误拒绝。三个等式性质均成立，但为空洞成立：这些等式得到满足，是因为用于验证它们的组合被拒绝了。冲突拒绝阻止了对已有标签的扩展~\cite{harper1991-symmetric-record-concatenation}，而这恰恰是表达问题~\cite{wadler1998-expression-problem}所要求的可组合性。

  \item[\bilingual{Priority merge}{优先级合并}]
    一个来源具有优先权~\cite{jsonnet,dolstra2010-nixos}，因此 $A + B \neq B + A$，引入了顺序依赖。

  \item[$\beta$-\bilingual{reduction}{归约}]
    只有作者放置的参数才是扩展点；函数体是封闭的。要使一个新的方面可扩展，原函数必须被重写以接受一个额外参数。任何基于函数应用的组合机制都继承了这一限制，因为 $\beta$-归约使函数体封闭。表达问题~\cite{wadler1998-expression-problem}在此类语言中之所以困难，原因相同：可扩展性是按需选择的。模拟对不可扩展性免疫的框架（访问者模式、对象代数~\cite{oliveira2012-object-algebras}、finally tagless 解释器~\cite{carette2009-finally-tagless}）正是为绕过这一限制而生的机制。

  \item[\bilingual{Conditional incorporation}{条件并入}]
    根据其他定义的存在与否来移除或有条件地包含定义，会产生依赖循环：某定义是否存在，取决于另一定义是否存在，而后者又取决于前者。当添加和移除交替出现时，递归求值（附录~\ref{app:well-definedness}）无法到达不动点，从而导致发散。这与图灵完备性固有的发散不同：图灵完备程序发散是因为计算无界，而条件并入发散是因为直接后果算子 $T_P$ 振荡而非单调递增。
\end{description}

\noindent
在上述候选方案中，深度合并是我们找到的唯一一个可交换、幂等、可结合且对不可扩展性免疫的方案。其他候选方案各自引入了对不可扩展性的脆弱性，或导致发散。

\subsubsection*{\bilingual{The reference resolution}{引用解析}}

一旦将记录上的深度合并确定为组合机制，六个语义方程中的四个便由语法决定：$\properties$~(\ref{eq:properties}) 读取记录的标签，$\bases$~(\ref{eq:bases}) 读取其继承来源，$\supers$~(\ref{eq:supers}) 取传递闭包，$\overrides$~(\ref{eq:overrides}) 实现合并。唯一剩余的自由度是引用的解析方式：$\resolve$~(\ref{eq:resolve}) 和 $\this$~(\ref{eq:this}) 方程。我们考察三个已知候选方案。

\begin{description}
  \item[\bilingual{Early binding}{早期绑定}]
    使用早期绑定时，引用在定义时即解析到一个固定路径，与记录后来如何被继承无关。CUE~\cite{cue2019}是这方面的典型：字段引用被静态解析，嵌套字段无法引用其封闭上下文中由后续组合贡献的属性。这使得 mixin 无法充当 $\lambda$-演算的调用栈：$\beta$-归约要求被调用方能观察到调用点提供的参数，而那是一次后续组合。使用早期绑定时，引用在此组合发生前就已被冻结。在方程中，$\resolve$~(\ref{eq:resolve}) 调用 $\this$~(\ref{eq:this})，后者遍历继承的树结构；迟绑定正是使 $\this$ 能够观察到后续组合的 mixin 所贡献属性的原因。若没有图灵完备的 mixin，函数必须被重新引入为组合机制，从而再次暴露上述按需选择的可扩展性问题。

  \item[\bilingual{Dynamic scope}{动态作用域}]
    使用动态作用域时，引用在使用点而非定义点处解析。Jsonnet~\cite{jsonnet}采用此方式：\texttt{self} 始终指向最终合并后的对象，而非定义点处的对象。联合文件系统在文件系统层面表现出相同的性质：符号链接的相对路径在解引用时相对于挂载上下文解析，而非相对于定义该链接的层。动态作用域破坏了第~\ref{sec:forward-translation}~节的 $\lambda$-演算嵌入。替换引理（引理~\ref{lem:substitution}）依赖于基本情况 $M = y$（$y \neq x$）：$y$ 的 de~Bruijn 指标 $m$ 从定义点路径 $\init(p_{\mathrm{def}})$ 向上走到一个与 $x$ 的绑定器不同的作用域层级，因此 $x$ 作用域处的替换 $\mathrm{argument} \mapsto \mathcal{T}(V)$ 不产生任何效果。动态作用域下，这种分离不再成立。%
    \footnote{
      反例：$(\lambda x.\, \lambda y.\, x)\; V_1\; V_2$。$\lambda y.\, x$ 内部对 $x$ 的引用翻译为 $\dbi{1}.\mathrm{argument}$（de~Bruijn 指标 $1$）。使用迟绑定时，$\this$ 从定义点作用域（内层 $\lambda$）出发，向上走一层到外层 $\lambda$，到达槽 $\mathrm{argument} \mapsto \mathcal{T}(V_1)$。使用动态作用域时，引用从使用点上下文解析，其中最内层封闭作用域是应用 $\{C_1.\mathrm{result},\; \mathrm{argument} \mapsto \mathcal{T}(V_2)\}$；向上走一层到达该作用域的 $\mathrm{argument}$ 槽，返回 $\mathcal{T}(V_2)$ 而非 $\mathcal{T}(V_1)$。这是动态作用域的经典变量捕获失败。
    }
    与早期绑定一样，mixin 无法完全替代函数。

  \item[\bilingual{Single-target $\this$}{单目标 $\this$}]
    单目标 $\this$ 是一个有前途的候选方案：单路径引理~(\ref{lem:single-path})表明，对于翻译后的 $\lambda$-项，$\this$~(\ref{eq:this}) 中的前沿集 $S$ 始终恰好包含一条路径，因此 $\lambda$-演算嵌入不需要多目标解析。然而，单目标 $\this$ 因另一原因而对不可扩展性脆弱。当多条继承路径到达同一作用域时，单目标解析必须拒绝该情况或选择其中一条路径。NixOS 模块系统~\cite{nixosmodules}会引发重复声明错误（附录~\ref{app:scala-multi-target}）。这拒绝了两个独立编写的模块定义同一选项的嵌套扩展，而这种组合在求解表达问题~\cite{wadler1998-expression-problem}时会出现，如第~\ref{sec:expression-problem}~节所示。

    我们对 MIXINv2 的前三个实现使用了 $\this$ 解析的标准技术：遵循 Cook~\cite{cook1989-denotational-inheritance}的基于闭包的不动点，以及使用 de~Bruijn 指标~\cite{debruijn1972-nameless-dummies}的基于栈的环境查找。当继承引入到同一封闭作用域的多条路径时，三种实现都产生了隐晦的错误。困难在于结构上。闭包捕获单一环境，因此不动点组合子只能为单一 $\this$ 目标求解。基于栈的环境是一条线性链，每个作用域层级只有一个父节点。这两种数据结构都内嵌了一个与多重继承自然产生的多目标情况不相容的单值假设（附录~\ref{app:scala-multi-target}）。公式~(\ref{eq:this})是在放弃这一假设后涌现出来的：它跟踪一个\emph{前沿集} $S$，而非单一当前作用域。
\end{description}

\noindent
下表总结了两个层面上所有替代方案的脆弱性。

\begin{center}
  \begin{tabular}{llll}
    & \bilingual{\textbf{Candidate}}{\textbf{候选方案}} & \bilingual{\textbf{Vulnerability}}{\textbf{脆弱性}}
    & \bilingual{\textbf{Representative}}{\textbf{代表}} \\
    \hline
    \multirow{5}{*}{\rotatebox[origin=c]{90}{\scriptsize \bilingual{composition}{组合}}}
    & \bilingual{Shallow merge}{浅层合并}
    & \bilingual{nested extensions not composable}{嵌套扩展不可组合}
    & JS spread \\
    & \bilingual{Conflict rejection}{冲突拒绝}
    & \bilingual{same-label extensions rejected}{同名扩展被拒绝}
    & Harper--Pierce~\cite{harper1991-symmetric-record-concatenation} \\
    & \bilingual{Priority merge}{优先级合并}
    & \bilingual{ordering dependency}{顺序依赖}
    & Jsonnet~\texttt{+}~\cite{jsonnet},
    Nix~\texttt{//}~\cite{dolstra2010-nixos} \\
    & $\beta$-\bilingual{reduction}{归约}
    & \bilingual{opt-in extensibility}{按需选择的可扩展性}
    & Java, Haskell, Scala, \ldots \\
    & \bilingual{Conditional}{条件并入}
    & \bilingual{causes divergence}{导致发散}
    & --- \\
    \hline
    \multirow{3}{*}{\rotatebox[origin=c]{90}{\scriptsize \bilingual{resolution}{解析}}}
    & \bilingual{Early binding}{早期绑定}
    & \bilingual{mixin cannot replace function}{mixin 无法替代函数}
    & CUE~\cite{cue2019} \\
    & \bilingual{Dynamic scope}{动态作用域}
    & \bilingual{mixin cannot replace function}{mixin 无法替代函数}
    & Jsonnet~\cite{jsonnet}, union FS \\
    & \bilingual{Single-target $\this$}{单目标 $\this$}
    & \bilingual{nested extensions rejected}{嵌套扩展被拒绝}
    & NixOS modules~\cite{nixosmodules} \\
  \end{tabular}
\end{center}

\noindent
我们在两个层面考察的所有替代方案，均对不可扩展性不免疫。

我们通过尝试本节中的每个替代方案并发现其失败，从而得到了这一语义：浅层合并无法组合嵌套扩展；优先级合并引入了顺序依赖；早期绑定在组合能够提供引用目标之前就将引用冻结；动态作用域破坏了 $\lambda$-演算嵌入所依赖的替换；单目标 $\this$ 拒绝了表达问题所要求的多目标组合。深度合并加上多目标迟绑定 $\this$，正是最终剩下的选择。
\else
The \hyperref[item:immunity]{immunity-to-nonextensibility item} above established that inheritance-calculus is immune
to nonextensibility.
We adopt this as a design goal and ask whether an alternative
design could share this property
by surveying every candidate known to us.

\paragraph{\bilingual{Degrees of freedom}{自由度}}
The analysis below assumes that the AST is a \emph{named tree}:%
every edge from parent to child carries a label, so that every
node is uniquely identified by the path from the root.
For configuration languages this is naturally the case: every
key in a JSON, YAML, or Nix attribute set is a label, and nesting
produces paths.%
\footnote{
  Nix is a functional language with first-class functions;
  its attribute sets alone do not form a named-tree DSL.
  The NixOS module system~\cite{nixosmodules}, however,
  restricts the user-facing interface to nested attribute sets
  composed by recursive merging, and it is this DSL-level
  structure that naturally forms a named tree.
}
For mainstream programming languages the AST is not directly a
named tree, since anonymous expression positions and implicit
scopes lack labels.
After ANF
transformation~\cite{flanagan1993-essence-compiling-continuations}, however, every intermediate
result is bound to a name, and the remaining anonymous positions%
\footnote{
  For example, the body of a $\mathbf{let}$.
}
can be assigned synthetic
fresh labels such as $\mathrm{result}$ and $\mathrm{tailCall}$
in Section~\ref{sec:forward-translation}, yielding a named
tree.

Under this premise, every language whose programs are such
named trees with references can be characterized by two finite sets:
a set~$\mathcal{N}$ of \emph{node types}, each determining how a
subtree organizes its children, and a set~$\mathcal{R}$ of
\emph{reference types}, each determining how one position in the
tree refers to another.
The syntax is fixed by the choice of $\mathcal{N}$ and
$\mathcal{R}$; the semantics is determined by how each node type
structures composition and how each reference type resolves its
target.

{\small
  \begin{center}
    \begin{tabular}{lp{0.38\columnwidth}p{0.33\columnwidth}}
      \textbf{Language}
      & \textbf{Node types $\mathcal{N}$}
      & \textbf{Reference types $\mathcal{R}$} \\
      \hline
      Java (OOP)
      & class, interface, method,
      block, conditional, loop,
      \texttt{throw}, lambda, \ldots
      & variable, field access, method call,
      \texttt{extends}/\texttt{implements},
      \texttt{import}, \ldots \\[3pt]
      Haskell (FP)
      & $\lambda$, application, \texttt{let},
      \texttt{case}, \texttt{data},
      type class, instance, \ldots
      & variable, constructor,
      type class dispatch,
      \texttt{import}, \ldots \\[3pt]
      Scala\footnote{
        Scala is a multi-paradigm language.
      }
      & class, trait, object, method,
      block, conditional, \texttt{match},
      \texttt{while}, \texttt{throw}, \ldots
      & variable, member access,
      \texttt{extends}, \texttt{with},
      implicits, \texttt{import}, \ldots \\
      \hline
      $\lambda$-calculus
      & abstraction, application
      & variable \\
      Inheritance-calculus
      & record
      & inheritance \\
    \end{tabular}
  \end{center}
}

\paragraph{\bilingual{Immunity to nonextensibility}{对不可扩展性的免疫}}
A path~$p$ in a program~$P$ is \emph{extensible} if there exists
a program~$Q$ such that composing $P$ with~$Q$ through a reference
makes new properties observable at~$p$ without modifying~$P$.
A language is \emph{immune to nonextensibility} if every path in
every well-formed program is extensible.

We now classify the design space to explain why we chose the
semantics of Section~\ref{sec:mixin-trees}.
Two axes remain free:
the \emph{composition mechanism}, which determines how the single reference
type incorporates definitions;
and the \emph{reference resolution}, which determines how the single reference
type finds its target.

\subsubsection*{\bilingual{The composition mechanism}{组合机制}}

Every composition mechanism determines how the definitions of one
abstraction are incorporated into another.
The following table classifies the known candidates by whether
they satisfy commutativity~(C), idempotence~(I), and
associativity~(A), and by their extensibility model.

\begin{center}
  \begin{tabular}{lcccll}
    \bilingual{\textbf{Candidate}}{\textbf{候选方案}} & \textbf{C} & \textbf{I} & \textbf{A}
    & \bilingual{\textbf{Extensibility}}{\textbf{可扩展性}} & \bilingual{\textbf{Representative}}{\textbf{代表}} \\
    \hline
    \bilingual{Deep merge}{深度合并}
    & $\checkmark$ & $\checkmark$ & $\checkmark$
    & \bilingual{universal}{通用} & \bilingual{this paper}{本文} \\
    \bilingual{Shallow merge}{浅层合并}
    & $\times$ & $\checkmark$ & $\checkmark$
    & \bilingual{top-level only}{仅顶层} & JS spread \\
    \bilingual{Conflict rejection}{冲突拒绝}
    & $\checkmark$ & $\checkmark$ & $\checkmark$
    & \bilingual{none on overlap}{重叠时无扩展}
    & Harper--Pierce~\cite{harper1991-symmetric-record-concatenation} \\
    \bilingual{Priority merge}{优先级合并}
    & $\times$ & $\times$ & $\checkmark$
    & \bilingual{universal}{通用} & Jsonnet~\texttt{+}, Nix~\texttt{//} \\
    $\beta$-\bilingual{reduction}{归约}
    & $\times$ & $\times$ & $\times$
    & \bilingual{opt-in}{按需选择} & Java, Haskell, Scala, \ldots \\
    \bilingual{Conditional}{条件并入}
    & \multicolumn{3}{c}{\bilingual{causes divergence}{导致发散}}
    & --- & --- \\
  \end{tabular}
\end{center}

\begin{description}
  \item[Deep merge]
    Merging same-label definitions recurses into their subtrees
    ($\overrides$, equation~\ref{eq:overrides}).
    Every label is an extension point, whether or not the
    original author anticipated extension.

  \item[Shallow merge]
    Nested labels cannot be independently extended.

  \item[Conflict rejection]
    Same-label definitions are rejected as errors.
    All three identities hold, but vacuously:
    the identities are satisfied because the compositions that
    would test them are refused.
    Conflict rejection prevents extending existing
    labels~\cite{harper1991-symmetric-record-concatenation}, precisely the composability that the
    Expression Problem~\cite{wadler1998-expression-problem} requires.

  \item[Priority merge]
    One source takes
    precedence~\cite{jsonnet,dolstra2010-nixos},
    so $A + B \neq B + A$, introducing
    an ordering dependency.

  \item[$\beta$-reduction]
    Only the parameters placed by the author are extension
    points; the function body is closed.
    To make a new aspect extensible, the original function must
    be rewritten to accept an additional parameter.
    Any composition mechanism based on function application
    inherits this limitation, because $\beta$-reduction
    closes the function body.
    The Expression Problem~\cite{wadler1998-expression-problem} is hard
    in such languages for the same reason:
    extensibility is opt-in.
    The frameworks that simulate immunity to nonextensibility
    (visitor pattern,
      object algebras~\cite{oliveira2012-object-algebras},
    finally tagless interpreters~\cite{carette2009-finally-tagless})
    are the machinery needed to work around this limitation.

  \item[Conditional incorporation]
    Removing or conditionally including definitions
    based on the presence of other definitions
    creates a dependency cycle:
    whether a definition exists depends on
    whether another definition exists, which in turn
    depends on the first.
    The recursive evaluation
    (Appendix~\ref{app:well-definedness}) cannot
    reach a fixed point when addition and removal
    alternate, causing divergence.
    This differs from the divergence inherent to
    Turing completeness: Turing-complete programs
    diverge because computation is unbounded, whereas
    conditional incorporation diverges because the
    immediate consequence operator $T_P$ oscillates
    rather than ascending monotonically.
\end{description}

\noindent
Among the candidates above, deep merge is the only one we found
that is commutative, idempotent, associative, and
immune to nonextensibility.
The other candidates each introduce a vulnerability to
nonextensibility or cause divergence.

\subsubsection*{\bilingual{The reference resolution}{引用解析}}

Once deep merge over records is fixed as the composition
mechanism, four of the six semantic equations are determined
by the syntax:
$\properties$~(\ref{eq:properties}) reads the labels of a
record, $\bases$~(\ref{eq:bases}) reads its inheritance
sources, $\supers$~(\ref{eq:supers}) takes the transitive
closure, and $\overrides$~(\ref{eq:overrides}) implements
the merge.
The only remaining degree of freedom is how references
resolve: the $\resolve$~(\ref{eq:resolve}) and
$\this$~(\ref{eq:this}) equations.
We examine the three known candidates.

\begin{description}
  \item[Early binding]
    With early binding, references resolve at definition time to a
    fixed path, independently of how the record is later inherited.
    CUE~\cite{cue2019} exemplifies this: field references are
    statically resolved, and a nested field cannot refer to
    properties of its enclosing context contributed by later
    composition.
    This prevents the mixin from serving as a call stack for the
    $\lambda$-calculus: $\beta$-reduction requires the callee to
    observe the argument supplied at the call site, which is a
    later composition.
    With early binding, references are frozen before this
    composition occurs.
    In the equations,
    $\resolve$~(\ref{eq:resolve}) calls
    $\this$~(\ref{eq:this}), which walks the inherited
    tree structure; late binding is what allows
    $\this$ to see properties contributed by later-composed
    mixins.
    Without a Turing-complete mixin, functions must be
    reintroduced as the composition mechanism, re-exposing
    the opt-in extensibility problem above.

  \item[Dynamic scope]
    With dynamic scope, references resolve at the point of use
    rather than the definition site.
    Jsonnet~\cite{jsonnet} follows this approach: \texttt{self}
    always refers to the final merged object, not the object at
    the definition site.
    Union file systems exhibit the same property at the
    file-system level: a symbolic link's relative path resolves
    against the mount context at dereference time, not against the
    layer that defined the link.
    Dynamic scope breaks the $\lambda$-calculus embedding
    of Section~\ref{sec:forward-translation}.
    The Substitution Lemma (Lemma~\ref{lem:substitution})
    relies on the base case $M = y$ ($y \neq x$):
    the de~Bruijn index $m$ of~$y$ walks up from the
    definition-site path $\init(p_{\mathrm{def}})$ to a scope
    level different from~$x$'s binder, so the substitution
    $\mathrm{argument} \mapsto \mathcal{T}(V)$ at~$x$'s scope
    has no effect.
    With dynamic scope this separation fails.%
    \footnote{
      Counterexample:
      $(\lambda x.\, \lambda y.\, x)\; V_1\; V_2$.
      The reference to~$x$ inside $\lambda y.\, x$ translates to
      $\dbi{1}.\mathrm{argument}$ (de~Bruijn index~$1$).
      With late binding, $\this$ starts from the definition-site
      scope (the inner $\lambda$) and walks up one level to the
      outer $\lambda$, reaching the slot
      $\mathrm{argument} \mapsto \mathcal{T}(V_1)$.
      With dynamic scope, the reference resolves from the
      use-site context, where the innermost enclosing scope is
      the application $\{C_1.\mathrm{result},\;
      \mathrm{argument} \mapsto \mathcal{T}(V_2)\}$;
      walking up one level reaches this scope's
      $\mathrm{argument}$ slot, returning $\mathcal{T}(V_2)$
      instead of $\mathcal{T}(V_1)$.
      This is the classical variable-capture failure of dynamic
      scoping.
    }
    As with early binding, the mixin cannot serve as a complete
    replacement for functions.

  \item[Single-target $\this$]
    Single-target $\this$ is a promising candidate:
    the Single-Path Lemma~(\ref{lem:single-path}) shows that
    for translated $\lambda$-terms, the frontier set~$S$ in
    $\this$~(\ref{eq:this}) always contains exactly one path,
    so the $\lambda$-calculus embedding does not require
    multi-target resolution.
    However, single-target $\this$ is vulnerable to
    nonextensibility for a different reason.
    When multiple inheritance routes reach the same scope,
    single-target resolution must reject the situation or select
    one route.
    The NixOS module system~\cite{nixosmodules} raises a
    duplicate-declaration error
    (Appendix~\ref{app:scala-multi-target}).
    This rejects nested extensions where two independently authored
    modules define the same option, a composition that arises
    when solving the Expression
    Problem~\cite{wadler1998-expression-problem},
    as Section~\ref{sec:expression-problem} demonstrates.

    Our first three implementations of MIXINv2 used
    standard techniques for $\this$ resolution: closure-based
    fixed points following Cook~\cite{cook1989-denotational-inheritance}, and
    stack-based environment lookup using de~Bruijn
    indices~\cite{debruijn1972-nameless-dummies}.
    All three produced subtle bugs when inheritance introduced
    multiple routes to the same enclosing scope.
    The difficulty was structural.
    A closure captures a single environment, so a fixed-point
    combinator solves for a single $\this$ target.
    A stack-based environment is a linear chain where each scope
    level has one parent.
    Both data structures embed a single-valued assumption
    incompatible with the multi-target situation that multiple
    inheritance naturally produces
    (Appendix~\ref{app:scala-multi-target}).
    Equation~(\ref{eq:this}) emerged from abandoning this
    assumption: it tracks a \emph{frontier set}~$S$ rather than a
    single current scope.
\end{description}

\noindent
The following table summarizes the vulnerabilities of all
alternatives at both levels.

\begin{center}
  \begin{tabular}{llll}
    & \bilingual{\textbf{Candidate}}{\textbf{候选方案}} & \bilingual{\textbf{Vulnerability}}{\textbf{脆弱性}}
    & \bilingual{\textbf{Representative}}{\textbf{代表}} \\
    \hline
    \multirow{5}{*}{\rotatebox[origin=c]{90}{\scriptsize \bilingual{composition}{组合}}}
    & \bilingual{Shallow merge}{浅层合并}
    & \bilingual{nested extensions not composable}{嵌套扩展不可组合}
    & JS spread \\
    & \bilingual{Conflict rejection}{冲突拒绝}
    & \bilingual{same-label extensions rejected}{同名扩展被拒绝}
    & Harper--Pierce~\cite{harper1991-symmetric-record-concatenation} \\
    & \bilingual{Priority merge}{优先级合并}
    & \bilingual{ordering dependency}{顺序依赖}
    & Jsonnet~\texttt{+}~\cite{jsonnet},
    Nix~\texttt{//}~\cite{dolstra2010-nixos} \\
    & $\beta$-\bilingual{reduction}{归约}
    & \bilingual{opt-in extensibility}{按需选择的可扩展性}
    & Java, Haskell, Scala, \ldots \\
    & \bilingual{Conditional}{条件并入}
    & \bilingual{causes divergence}{导致发散}
    & --- \\
    \hline
    \multirow{3}{*}{\rotatebox[origin=c]{90}{\scriptsize \bilingual{resolution}{解析}}}
    & \bilingual{Early binding}{早期绑定}
    & \bilingual{mixin cannot replace function}{mixin 无法替代函数}
    & CUE~\cite{cue2019} \\
    & \bilingual{Dynamic scope}{动态作用域}
    & \bilingual{mixin cannot replace function}{mixin 无法替代函数}
    & Jsonnet~\cite{jsonnet}, union FS \\
    & \bilingual{Single-target $\this$}{单目标 $\this$}
    & \bilingual{nested extensions rejected}{嵌套扩展被拒绝}
    & NixOS modules~\cite{nixosmodules} \\
  \end{tabular}
\end{center}

\noindent
No alternative we examined at either level is immune to
nonextensibility.

We arrived at this semantics by trying each alternative in this section
and finding that it failed:
shallow merge could not compose nested extensions;
priority merge introduced ordering dependencies;
early binding froze references before composition could
supply them;
dynamic scope broke the substitution that the
$\lambda$-calculus embedding relies on;
single-target $\this$ rejected the multi-target compositions
that the Expression Problem demands.
Deep merge with multi-target late-binding $\this$ is what
remained.
\fi

\section{\bilingual{Related Work}{相关工作}}
\label{sec:related-work}

\subsection{\bilingual{Computational Models and \texorpdfstring{$\lambda$}{λ}-Calculus Semantics}{计算模型与 \texorpdfstring{$\lambda$}{λ}-演算语义}}

\paragraph{\bilingual{Minimal Turing-complete models}{最小图灵完备模型}}
\ifInheritanceChinese
Sch\"onfinkel~\cite{schonfinkel1924-combinatory-logic} 与 Curry~\cite{curry1958-combinatory-logic} 证明了 $\lambda$-演算可以归约为三个组合子 $S$、$K$、$I$%
\footnote{
  组合子 $I$ 是冗余的,因为 $I = SKK$。
}。将 $\lambda$-项编译为 SKI 的括号抽象是机械的,但会导致项大小的指数级膨胀;在 $\lambda$-演算中已有的计算模式之外,不会涌现出任何新的计算模式。在语义上,两者共享相同的模型:Meyer~\cite{meyer1982-model-of-lambda-calculus} 证明了满足五个附加条件的组合代数恰好就是 $\lambda$-代数。从未有人展示过一个独立的 SKI 语义域,并从中推导出 $\lambda$-演算的语义。方向始终是反过来的:Scott 的 $D_\infty$ 是为 $\lambda$-演算而构建的,而 SKI 随后被解释于其中。Dolan~\cite{dolan2013-mov-turing-complete} 证明了 x86 的 \texttt{mov} 指令单独就已图灵完备,其方法是利用寻址模式作为可变内存上的隐式算术;由此得到的程序是正确的,但不提供任何超越 RAM 机的抽象机制。图灵机本身虽然是通用的,却缺乏随机访问:读取距离为 $d$ 处的格需要 $d$ 次顺序的磁带移动。这些模型中的每一个,都在某个已有模型的语义框架内实现了图灵完备性,且对其中任何一个都从未展示过独立的语义域。
\else
Sch\"onfinkel~\cite{schonfinkel1924-combinatory-logic} and
Curry~\cite{curry1958-combinatory-logic} showed that the
$\lambda$-calculus reduces to three combinators $S$, $K$, $I$
\footnote{
  The $I$ combinator is redundant since $I = SKK$.
}.
The bracket abstraction that compiles $\lambda$-terms to SKI
is mechanical but incurs exponential blow-up in term size;
no new computational patterns emerge beyond those already
present in the $\lambda$-calculus.
Semantically, the two share the same models:
Meyer~\cite{meyer1982-model-of-lambda-calculus} showed that combinatory algebras
satisfying five additional conditions are exactly lambda
algebras.
No independent semantic domain for SKI has ever been
exhibited from which $\lambda$-calculus semantics could be
derived. The direction has always been the reverse:
Scott's $D_\infty$ was built for the $\lambda$-calculus,
and SKI was interpreted inside it afterwards.
Dolan~\cite{dolan2013-mov-turing-complete} showed that the x86 \texttt{mov}
instruction alone is Turing complete by exploiting
addressing modes as implicit arithmetic on mutable memory;
the resulting programs are correct but provide no
abstraction mechanisms beyond those of the RAM Machine.
The Turing Machine itself, while universal, lacks random
access: reading a cell at distance~$d$ requires $d$
sequential tape movements.
Each of these models achieves Turing completeness
within the semantic framework of an existing model,
and no independent semantic domain has been exhibited
for any of them.
\fi

%% === 审稿人备忘：SF-calculus 与 inheritance-calculus 的关系 ===
%%
%% Barry Jay 的 SF-calculus（Intensional Computation with Higher-Order
%% Functions, TCS 2019, doi:10.1016/j.tcs.2019.02.016；前身为
%% Confusion in the Church-Turing Thesis, arXiv:1410.7103, 2014）
%% 和 inheritance-calculus 都试图解决 λ-calculus 的 extensionality 限制，
%% 但方向完全不同。
%%
%% 共同出发点：λ-calculus 是 extensional 的——只能观察函数的输入输出行为，
%% 无法检查程序的内部结构。两个内部结构不同但行为相同的 λ 项不可区分。
%%   - Jay 的诊断：λ-calculus 缺乏 intensional 操作（检查项的结构），
%%     无法定义结构等式判定。
%%   - 本文的诊断：λ-calculus 的 function body 是 closed 的，无法从外部
%%     添加新的可观察投影（open extension）。
%%
%% 核心机制对比：
%%   SF-calculus：保留 S、K（或 λ），添加 F（factorisation）算子，
%%     可以将闭合范式拆解为组成部分。程序仍然是函数，只是多了检查函数
%%     结构的能力。属于 λ-calculus 的 *扩展*。
%%   Inheritance-calculus：用三个不同的原语（record、definition、
%%     inheritance）替代 λ-calculus 的三个原语。β-reduction 是 deep merge
%%     的特例（Section 3 的翻译 T 将函数应用映射为继承）。
%%     属于 λ-calculus 的 *包含*（subsumption）。
%%
%% 关于 F 算子的编码：
%%   F 做两件事：(1) 判断结构形态（atom vs. compound），
%%   (2) 提取子部分。
%%   在 inheritance-calculus 中：
%%     - 提取子部分是平凡的：所有值都是透明 record，通过 qualified this
%%       投影任意 slot 即可。Theorem 6 的 separating context 就在读取
%%       argument slot（↑².argument）。
%%     - 判断结构形态并分支：需要 Acceptance pattern（Section 4）。
%%       NatVisitor 对 Nat 做的事情与 F 对 SF-calculus 项做的事情
%%       完全对应：VisitZero/SuccessorVisitor 对应 F 的 atom/compound
%%       分支。Visitor 可以通过 open extension 添加到已有 term 上，
%%       不修改原始定义。
%%   Acceptance 在精神上就是 F：两者都是对 closed normal form 做
%%   constructor dispatch。区别在于 F 是不可扩展的 built-in primitive
%%   （硬编码了 S 和 F 两个 atom），而 Acceptance 是可扩展的 encoding
%%   （新增 constructor 只需继承新的 Accepted* 分支，即 Expression Problem
%%   immunity）。
%%
%% 精确差异：
%%   SF-calculus 把 intensional observation 做成了 object-level 的
%%   primitive reduction rule（F 是一个 term）；
%%   inheritance-calculus 把它留在了 meta-level（观察者总能通过
%%   properties(p) 查询任何路径），而 term-level 的分支需要通过
%%   Acceptance pattern 显式编码。
%%   换言之：inheritance-calculus 中 F 的 *读取* 能力是免费的
%%   （所有值都是透明 record），但 F 的 *分支* 能力需要 Acceptance。
%%   SF-calculus 的 F 则把读取和分支捆绑在一个不可扩展的 primitive 里。
%%   Acceptance 是可扩展的（新增 constructor 只需继承新分支，
%%   即 Expression Problem immunity），而 F 硬编码了固定的 atom 集合
%%   （S 和 F）。
%%
%% 关于 Jay 所谓的 "trivial solution of the Halting Problem"：
%%   Jay 用非标准编码将每个程序（包括不停机的）映射到闭合范式（静态数据），
%%   然后结构等式判定变得平凡。这并不推翻图灵定理——它改变了问题的
%%   形式化方式（把程序当静态数据而非可执行代码），而非解决了原来的问题。
%%   Jay 的真正论点是：λ-definability 和 Turing computability 不是同一回事，
%%   Church-Turing Thesis 的等价性证明隐含了编码假设。
%%
%% 正文篇幅不够论述 embeds SF-calculus，暂不展开。
%% 如果审稿人提到 SF-calculus，可以基于上述分析撰写回复或补充段落。
%% ================================================================

\paragraph{\bilingual{Semantic descriptions of the $\lambda$-calculus}{$\lambda$-演算的语义描述}}
\ifInheritanceChinese
$\lambda$-演算的语义可以用几种根本上不同的方式来定义。\emph{操作语义}将 $\beta$-归约 $(\lambda x.\,M)\,N \to M[N/x]$ 作为计算的基本概念~\cite{church1936-lambda-calculus,church1941-calculi-of-lambda-conversion,barendregt1984-lambda-calculus}。这要求捕获-回避替换,其微妙之处促使了 de~Bruijn 索引~\cite{debruijn1972-nameless-dummies} 与显式替换演算~\cite{abadi1991-explicit-substitutions} 的出现。它还要求选择求值策略:按名调用与按值调用~\cite{plotkin1975-call-by-name-call-by-value},或 Abramsky 的惰性策略~\cite{abramsky1990-lazy-lambda-calculus}。\emph{指称语义}~\cite{scott1970-mathematical-theory-of-computation,scott1976-data-types-as-lattices} 将项解释于一个数学域中,但赋予 $\beta$-归约意义需要一个自反域 $D \cong [D \to D]$,由 Cantor 定理知该域无集合论解。Scott 的 $D_\infty$ 通过将 $D$ 构造为连续格的逆极限~\cite{scott1972-continuous-lattices} 来解决这一问题。由此得到的语义是适当的,但不是完全抽象的~\cite{plotkin1977-lcf-as-programming-language,milner1977-fully-abstract-models}:它将观察等价所区分的项等同了起来。\emph{博弈语义}~\cite{abramsky2000-full-abstraction-pcf,hyland2000-full-abstraction-pcf,nickau1994-hereditarily-sequential-functionals} 对 PCF 与惰性 $\lambda$-演算~\cite{abramsky1993-full-abstraction-lazy-lambda,abramsky1995-games-full-abstraction-lazy-lambda} 实现了完全抽象,需要竞技场、策略与无害性条件。在这些方法的每一种中,函数空间 $[D \to D]$ 都是核心障碍:操作语义通过替换隐式地需要它,而指称语义则显式地构造它,博弈语义则通过策略加以精化。这些方法中的每一种都涉及高阶结构%
\footnote{
  指称语义中的函数空间,博弈语义中的策略,以及操作语义中的替换。
}
并在其语义域中保留了函数空间。在每种情形下,指称帐和操作帐都是相互独立的形式体系。三者都继承了 $\beta$-归约在递归程序上的发散:例如,有向图上一个朴素的传递闭包在每个框架下都会发散。
\else
The semantics of the $\lambda$-calculus can be defined in
several fundamentally different ways.
\emph{Operational semantics} takes $\beta$-reduction
$(\lambda x.\,M)\,N \to M[N/x]$ as the primitive notion of
computation~\cite{church1936-lambda-calculus,church1941-calculi-of-lambda-conversion,barendregt1984-lambda-calculus}.
This requires capture-avoiding substitution, a mechanism whose
subtleties motivated de~Bruijn
indices~\cite{debruijn1972-nameless-dummies} and explicit substitution
calculi~\cite{abadi1991-explicit-substitutions}.
It also requires a choice of evaluation
strategy: call-by-name versus
call-by-value~\cite{plotkin1975-call-by-name-call-by-value} or the lazy
strategy of Abramsky~\cite{abramsky1990-lazy-lambda-calculus}.
\emph{Denotational semantics}~\cite{scott1970-mathematical-theory-of-computation,scott1976-data-types-as-lattices}
interprets terms in a mathematical domain, but giving meaning
to $\beta$-reduction requires a reflexive domain
$D \cong [D \to D]$, which has no set-theoretic solution by
Cantor's theorem. Scott's $D_\infty$ resolves this by
constructing $D$ as an inverse limit of continuous
lattices~\cite{scott1972-continuous-lattices}.
The resulting semantics is adequate but not fully
abstract~\cite{plotkin1977-lcf-as-programming-language,milner1977-fully-abstract-models}:
it identifies terms that observational equivalence
distinguishes.
\emph{Game semantics}~\cite{abramsky2000-full-abstraction-pcf,hyland2000-full-abstraction-pcf,nickau1994-hereditarily-sequential-functionals}
achieves full abstraction for PCF and for the lazy
$\lambda$-calculus~\cite{abramsky1993-full-abstraction-lazy-lambda,abramsky1995-games-full-abstraction-lazy-lambda},
requiring arenas, strategies, and innocence conditions.
In each of these approaches, the function space $[D \to D]$
is the central obstacle: operational semantics needs it
implicitly via substitution. In contrast, denotational semantics
constructs it explicitly, and game semantics refines it via
strategies.
Each of these approaches involves higher-order structure%
\footnote{
  Function spaces
  in denotational semantics, strategies
  in game semantics, and substitution
  in operational semantics.
}
and retains the function space in its semantic domain.
The denotational and operational accounts remain separate
formalisms in each case.
All three inherit $\beta$-reduction's divergence on
recursive programs: a na\"ive transitive closure on a cyclic
graph, for instance, diverges in every framework.
\fi

\paragraph{\bilingual{Set-theoretic models of the $\lambda$-calculus}{$\lambda$-演算的集合论模型}}
\ifInheritanceChinese
$\lambda$-演算的几个模型族通过在集合论而非域论框架下工作,避免了 Scott 的 $D_\infty$ 的逆极限构造。这些模型消除了 Scott 拓扑,但并未消除函数空间本身:它们仍然是高阶的,通过注入、交叉类型或多重集关系将 $[D \to D]$ 编码于 $D$ 之中。\emph{图模型}~\cite{engeler1981-algebras-and-combinators,plotkin1993-set-theoretic-lambda-models} 通过幂集格上的 web 注入来编码函数;Bucciarelli 与 Salibra~\cite{bucciarelli2008-graph-lambda-theories} 证明了最大的合理图理论等于 Böhm 树理论 $\mathcal{B}$。\emph{过滤器模型}~\cite{barendregt1983-filter-lambda-model,coppo1978-intersection-type-assignment} 将一个项解释为可从中推导出的交叉类型集合;Dezani-Ciancaglini 等人~\cite{dezani1998-intersection-types-bohm-trees} 证明了 BCD 过滤器模型的 $\lambda$-理论恰好与 Böhm 树相等性重合。\emph{关系模型}~\cite{ehrhard2005-finiteness-spaces,bucciarelli2001-phase-semantics-exponentials} 通过有限多重集从线性逻辑的关系语义中涌现;Ehrhard~\cite{ehrhard2005-finiteness-spaces} 证明了在这一范畴中,纯 $\lambda$-演算的标准自反对象不能存在。这三个模型族都是语义模型:它们将 $\lambda$-项解释于一个外部数学结构中,并在该结构中将 $\beta$-归约验证为一个等式,而不是定义一个具有自身原语的独立演算。每一种都编码函数空间而非消除它,且每一种都需要高阶机制才能为惰性 $\lambda$-演算获得完全抽象。
\else
Several families of $\lambda$-calculus models avoid the
inverse-limit construction of Scott's $D_\infty$ by working
in set-theoretic rather than domain-theoretic settings.
These models eliminate Scott topology but not the function
space itself: they remain higher-order, encoding $[D \to D]$
inside $D$ via injections, intersection types, or multiset
relations.
\emph{Graph models}~\cite{engeler1981-algebras-and-combinators,plotkin1993-set-theoretic-lambda-models}
encode functions via a web injection on powerset lattices;
Bucciarelli and Salibra~\cite{bucciarelli2008-graph-lambda-theories}
proved that the greatest sensible graph theory equals the
B\"ohm tree theory~$\mathcal{B}$.
\emph{Filter models}~\cite{barendregt1983-filter-lambda-model,coppo1978-intersection-type-assignment}
interpret a term as the set of intersection types derivable
for it;
Dezani-Ciancaglini et al.~\cite{dezani1998-intersection-types-bohm-trees}
showed that the BCD filter model's $\lambda$-theory
coincides exactly with B\"ohm tree equality.
\emph{Relational models}~\cite{ehrhard2005-finiteness-spaces,bucciarelli2001-phase-semantics-exponentials}
arise from the relational semantics of linear logic via
finite multisets;
Ehrhard~\cite{ehrhard2005-finiteness-spaces} showed that a standard
reflexive object for the pure $\lambda$-calculus cannot
exist in this category.
All three families are semantic models: they interpret
$\lambda$-terms in an external mathematical structure and
validate $\beta$-reduction as an equation in that structure,
rather than defining an independent calculus with its own
primitives.
Each encodes the function space rather than eliminating it,
and each requires higher-order machinery to obtain full
abstraction for the lazy $\lambda$-calculus.
\fi

\paragraph{\bilingual{Recursive equations and fixpoint semantics}{递归方程与不动点语义}}
\ifInheritanceChinese
Van~Emden 与 Kowalski~\cite{vanemden1976-predicate-logic-semantics} 证明了谓词逻辑作为编程语言的语义允许三种等价的刻画:操作性的\footnote{通常是 SLD 归结。}、模型论的\footnote{基于最小 Herbrand 模型。}以及不动点的\footnote{即基原子幂集上即时后承算子 $T_P$ 的最小不动点。}。Aczel~\cite{aczel1977-inductive-definitions} 通过幂集格上的单调算子对归纳定义给出了系统性处理,为 Datalog 风格的语义所依赖的逻辑基础奠定了基础。Leroy 与 Grall~\cite{leroy2006-coinductive-big-step} 通过余归纳定义处理了大步语义中的发散,需要在归纳求值判断之外引入一个单独的余归纳判断来说明非终止。Van~Emden 与 Kowalski 的 $T_P$ 在一阶域上为递归 Datalog 程序计算 $\mathrm{lfp}$。传统的 $\lambda$-演算语义在语义域中需要函数空间,且递归的 $\lambda$-项在 $\beta$-归约下通常会发散;利用一阶域上的 $T_P$ 为 $\lambda$-演算赋予一种改变哪些项收敛的不动点语义,尚未被探索。在近期工作中,Yang~\cite{yang2026-co-lambda} 将 tabling 求值~\cite{tamaki1986-tabled-resolution} 直接引入纯 $\lambda$-演算:首部规范化被读作一个余代数方程,其状态为被驻留的 $\lambda$-项,即具有可判定同一性的一阶对象,而 tabling 求解该方程,在有限时间内将 $\Omega$ 这样的非生产性循环判定为 $\bot$。该处理保留了标准的惰性 (Lévy--Longo) 树语义,仅改变了求值器;它占据了图~\ref{fig:calculi-matrix} 中不动点 $\lambda$ 节点。它在机制上也有所不同。表以整个被驻留的项为键,而非以进入作用域树的路径为键,因此在路径前缀上重叠的查询之间不共享部分结果。偏序不是集合包含下的幂集格,而是树近似序,其中 $\bot$ 叶被精化为子树,且不动点是唯一的(最小与最大因有护性而重合):一个重入调用被以单一的最终 $\bot$ 回答,而后续任何迭代都不会精化该 $\bot$,因此迭代无法增大一个累加器。其比惰性 $\lambda$-演算更确定的部分因此恰好是图~\ref{fig:calculi-matrix} 中非生产性的 $\bot$,即发散被判定为确定 $\bot$ 答案的那些项,别无其他;而幂集序才使得迭代可以积累并产生新的非 $\bot$ 值。
\else
Van~Emden and Kowalski~\cite{vanemden1976-predicate-logic-semantics} showed that
the semantics of predicate logic as a programming language admits
three equivalent characterizations:
operational\footnote{
  Typically SLD resolution.
}, model-theoretic\footnote{
  Based on least Herbrand models.
}, and fixed-point\footnote{
  Namely, the least fixed point of the immediate consequence operator~$T_P$ on the powerset of ground atoms.
}.
Aczel~\cite{aczel1977-inductive-definitions} gave a systematic treatment of
inductive definitions via monotone operators on powerset lattices,
providing the logical foundations that Datalog-style semantics
rests on.
Leroy and Grall~\cite{leroy2006-coinductive-big-step} treated divergence
in big-step semantics via coinductive definitions,
requiring a separate coinductive judgment to account for
non-termination alongside the inductive evaluation judgment.
Van~Emden and Kowalski's $T_P$ computes
$\mathrm{lfp}$ for recursive Datalog programs over a
first-order domain.
Conventional $\lambda$-calculus semantics requires function
spaces in the semantic domain, and recursive $\lambda$-terms
ordinarily diverge under $\beta$-reduction;
using $T_P$ over a first-order domain to give the
$\lambda$-calculus a fixpoint semantics that changes which
terms converge has not been explored.
In recent work, Yang~\cite{yang2026-co-lambda} brings tabled
evaluation~\cite{tamaki1986-tabled-resolution} to the pure
$\lambda$-calculus directly: head normalisation is read as a
coalgebraic equation whose states are interned
$\lambda$-terms, first-order objects with decidable identity,
and tabling solves it, deciding unproductive loops such
as~$\Omega$ as~$\bot$ in finite time.
That treatment preserves the standard lazy
(L\'evy--Longo) tree semantics and changes only the evaluator;
it occupies the fixpoint~$\lambda$ node of
Figure~\ref{fig:calculi-matrix}.
It also differs in mechanism.
The table is keyed by whole interned terms, not by paths
into a scope tree, so partial results are not shared between
queries that overlap on a path prefix.
The order is not a powerset lattice under set inclusion but
the tree approximation order, in which $\bot$ leaves are
refined into subtrees, and the fixpoint is unique (least and
greatest coincide by guardedness): a re-entrant call is
answered with a single final~$\bot$ that no later iteration
refines, so the iteration cannot grow an accumulator.
Its more-definedness over the lazy $\lambda$-calculus is
therefore exactly the unproductive~$\bot$ of
Figure~\ref{fig:calculi-matrix}, terms whose divergence is
decided as a definite~$\bot$ answer, and nothing more;
the powerset order is what lets iteration accumulate and
produce new non-$\bot$ values.
\fi

\subsection{\bilingual{Inheritance, Objects, and Records}{继承、对象与记录}}

\paragraph{\bilingual{Denotational semantics of inheritance}{继承的指称语义}}
\ifInheritanceChinese
Cook~\cite{cook1989-denotational-inheritance} 给出了继承的第一个指称语义,将对象建模为递归记录。继承是\emph{生成器}%
\footnote{
  这些是从 self 到完整对象的函数。
}与\emph{包装器}%
\footnote{
  包装器是修改生成器的函数。
}的复合,由一个不动点构造来求解。Cook 的关键洞见是,继承是一种可适用于任何形式的递归定义的通用机制,而不仅仅是面向对象的方法。这一洞见是本工作的出发点之一。该模型有四个原语(记录、函数、生成器和包装器),并需要一个不动点组合子。复合是不对称的:包装器修改生成器,反之则不然。不动点构造依赖函数域中的 Scott 连续性。由此得到的指称不能直接执行,因此需要一个单独的操作语义来在运行时定义方法分派与对象构造。
\else
Cook~\cite{cook1989-denotational-inheritance} gave the first denotational
semantics of inheritance, modeling objects as recursive records.
Inheritance is composition of \emph{generators}
\footnote{
  These are functions from self to complete object.
} and
\emph{wrappers}\footnote{
  Wrappers are functions that modify generators.
},
solved by a fixed-point construction.
Cook's key insight is that inheritance is a general mechanism
applicable to any form of recursive definition, not only
object-oriented methods. This insight is one of the starting points of the
present work.
The model has four primitives (records, functions, generators, and
wrappers) and requires a fixed-point combinator. Composition is
asymmetric: a wrapper modifies a generator but not vice versa.
The fixed-point construction relies on Scott continuity in a
domain of functions. The resulting denotations are not directly
executable, so a separate operational semantics is needed to
define method dispatch and object construction at runtime.
\fi

\paragraph{\bilingual{Mixins and traits}{Mixin 与 trait}}
\ifInheritanceChinese
Bracha 与 Cook~\cite{bracha1990-mixin-inheritance} 将 mixin 形式化为抽象子类,即从超类参数到子类的函数。由于 mixin 应用是函数复合,它\emph{既不可交换也不幂等};C3 线性化算法~\cite{barrett1996-c3-linearization} 后来被开发出来以施加一个确定性顺序。Schärli 等人~\cite{scharli2003-traits} 引入了 trait,其复合是可交换的,但拒绝同名冲突~\cite{ducasse2006-traits-fine-grained-reuse}。两种机制都在\emph{方法层级}操作:对同名的两个定义进行复合会导致覆盖或冲突,而非嵌套结构的递归合并。两者都预设了 $\lambda$-演算:方法体是函数,而复合机制本身不是图灵完备的。Mixin 与 trait 均不支持深合并。
\else
Bracha and Cook~\cite{bracha1990-mixin-inheritance} formalized mixins as
abstract subclasses, that is, functions from a superclass parameter to a
subclass.
Because mixin application is function composition, it is
\emph{neither commutative nor idempotent}; the C3 linearization
algorithm~\cite{barrett1996-c3-linearization} was later developed to impose a
deterministic order.
Sch\"arli et al.~\cite{scharli2003-traits} introduced traits,
whose composition is commutative but rejects same-name
conflicts~\cite{ducasse2006-traits-fine-grained-reuse}.
Both mechanisms operate at the \emph{method level}: composing two
definitions of the same name results in an override or a conflict,
not a recursive merge of nested structure.
Both presuppose the $\lambda$-calculus: method bodies are functions,
and the composition mechanism itself is not Turing complete.
Neither mixins nor traits are deep-mergeable.
\fi

\paragraph{\bilingual{Object and record calculi}{对象演算与记录演算}}
\ifInheritanceChinese
Abadi 与 Cardelli~\cite{abadi1996-theory-of-objects} 将对象作为基本概念,但方法是以 self 为参数的函数,且对象扩展是不对称的。Boudol~\cite{boudol2004-recursive-record-semantics} 证明了自引用记录需要一个不安全的不动点算子,其适定性依赖于求值策略。Harper 与 Pierce~\cite{harper1991-symmetric-record-concatenation} 给出了一个具有对称连接的记录演算,该演算拒绝同标签的复合;Cardelli~\cite{cardelli1992-extensible-records} 研究了带子类型的可扩展记录;Rémy~\cite{remy1989-row-polymorphism} 为 ML 中的记录和变体引入了行多态。这些演算全都在\emph{平坦}记录上操作:没有嵌套结构的递归合并,没有惰性观察,$\this$ 也需要一个显式的不动点组合子。
\else
Abadi and Cardelli~\cite{abadi1996-theory-of-objects} treat objects as the
primitive notion, but methods are functions that take self as a
parameter and object extension is asymmetric.
Boudol~\cite{boudol2004-recursive-record-semantics} showed that self-referential
records require an unsafe fixed-point operator whose
well-definedness depends on the evaluation strategy.
Harper and Pierce~\cite{harper1991-symmetric-record-concatenation} gave a record
calculus with symmetric concatenation that rejects same-label
composition;
Cardelli~\cite{cardelli1992-extensible-records} studied extensible records
with subtyping;
R\'emy~\cite{remy1989-row-polymorphism} introduced row polymorphism for
records and variants in ML.
All of these calculi operate on \emph{flat} records: there is
no recursive merging of nested structure, no lazy observation,
and $\this$ requires an explicit fixed-point combinator.
\fi

\paragraph{\bilingual{Family polymorphism and DOT}{族多态与 DOT}}
\ifInheritanceChinese
Ernst~\cite{ernst2001-family-polymorphism} 引入了族多态;Ernst、Ostermann 与 Cook~\cite{ernst2006-virtual-class-calculus} 在虚类演算中将其形式化,其中类通过路径表达式 \texttt{this.out.C} 来访问。Amin 等人~\cite{amin2016-dependent-object-types} 在 DOT 演算中形式化了路径依赖类型,这是 Scala 的理论基础。在这两个框架中,每条路径都解析为\emph{单一}封闭对象中的\emph{单一}类。当多条继承路径通向同一作用域时,单值解析必须选择其中一条路径并丢弃其余;DOT 的 self 变量 $\nu(x : T)\, d$ 将 $x$ 绑定到单一对象,而 Scala 在编译时拒绝由此产生的交叉类型(附录~\ref{app:scala-multi-target})。两个框架均未提供集合值解析的机制。
\else
Ernst~\cite{ernst2001-family-polymorphism} introduced family polymorphism;
Ernst, Ostermann, and Cook~\cite{ernst2006-virtual-class-calculus}
formalized it in the virtual class calculus, where classes are
accessed via path expressions \texttt{this.out.C}.
Amin et al.~\cite{amin2016-dependent-object-types} formalized path-dependent types
in the DOT calculus, the theoretical foundation of Scala.
In both frameworks, each path resolves to a
\emph{single} class in a \emph{single} enclosing object.
When multiple inheritance routes lead to the same scope,
single-valued resolution must choose one route and discard the
others; DOT's self variable $\nu(x : T)\, d$ binds $x$ to a
single object, and Scala rejects the resulting intersection
type at compile time
(Appendix~\ref{app:scala-multi-target}).
Neither framework provides a mechanism for set-valued
resolution.
\fi

\subsection{\bilingual{Configuration Languages and Module Systems}{配置语言与模块系统}}

\paragraph{\bilingual{Configuration languages}{配置语言}}
\ifInheritanceChinese
CUE~\cite{cue2019} 在单一格中统一了类型与值,其中以 \texttt{\&} 表示的合一是可交换、可结合且幂等的。CUE 的设计受到 Google 的 Borg 配置语言 (GCL) 的启发,后者使用了图合一。CUE 的格包含标量类型,其冲突语义规定合一不相容标量会得到 $\bot$;标量对于配置语言是否必要,还是仅凭树结构就已足够,是一个有趣的设计问题。Jsonnet~\cite{jsonnet} 通过 \texttt{+} 运算符提供了具有 mixin 语义的对象继承:复合\emph{不是}可交换的,因为在标量冲突上右侧胜出,而深合并需要在每个字段上通过 \texttt{+:} 语法显式选择加入,因此深合并虽可用但并非默认行为。Dhall~\cite{dhall2017} 采取函数式方法:它是一个带记录的类型化 $\lambda$-演算,其中复合是记录合并算子,而非继承机制。
\else
CUE~\cite{cue2019} unifies types and values in a single lattice
where unification, denoted by \texttt{\&}, is commutative, associative, and
idempotent.
CUE's design was motivated by Google's Borg Configuration
Language (GCL), which used graph unification.
CUE's lattice includes scalar types with a conflict
semantics where unifying incompatible scalars yields $\bot$; whether
scalars are necessary for a configuration language, or whether
tree structure alone suffices, is an interesting design question.
Jsonnet~\cite{jsonnet} provides object inheritance via the
\texttt{+} operator with mixin semantics: composition is
\emph{not} commutative since the right-hand side wins on scalar
conflicts, and deep merging requires explicit opt-in via the \texttt{+:}
syntax on a per-field basis, so deep merge is
available but not the default.
Dhall~\cite{dhall2017} takes a functional approach: it is a typed
$\lambda$-calculus with records, where composition is a record merge
operator, not an inheritance mechanism.
\fi

\paragraph{\bilingual{Module systems and deep merge}{模块系统与深合并}}
\ifInheritanceChinese
NixOS 模块系统~\cite{dolstra2010-nixos,nixosmodules} 通过\emph{递归合并嵌套属性集}来组合模块,而非通过线性化或方法层级的覆盖,这与先前的 mixin 和 trait 演算中区分树层级继承与方法层级复合的深合并语义相同。当两个模块定义了相同的嵌套路径时,其子树被合并而非覆盖。NixOS 模块系统集成了深合并、递归的 $\this$、模块的延迟求值,以及其上的丰富类型检查层。同一机制为 NixOS 系统配置、Home Manager~\cite{homemanager}、nix-darwin~\cite{nixdarwin}、flake-parts~\cite{flakeparts} 以及数十个其他生态系统提供支持,共同管理任意复杂度的配置。一个似乎尚未被完全解决的方面是 $\this$ 在多重继承下的语义。当两个各自继承自同一延迟模块的独立模块被组合时,每次导入都会引入对共享选项自己的声明;模块系统的 \texttt{merge\-Option\-Decls} 以静态错误拒绝重复声明(附录~\ref{app:scala-multi-target})。模块系统无法表达两条路径对同一选项的贡献是均等的,即递归合并自然产生的多目标情形,正是这一情形促成了第~\ref{sec:mixin-trees} 节的观测语义。
\else
The NixOS module system~\cite{dolstra2010-nixos,nixosmodules}
composes modules by \emph{recursively merging nested attribute sets},
not by linearization or method-level override, which is the same deep-merge
semantics that distinguishes tree-level inheritance from method-level
composition in prior mixin and trait calculi.
When two modules define the same nested path, their subtrees are
merged rather than overridden.
The NixOS module system incorporates deep merge,
recursive $\this$, deferred evaluation of modules, and a
rich type-checking layer on top.
The same mechanism powers
NixOS system configuration, Home
Manager~\cite{homemanager}, nix-darwin~\cite{nixdarwin},
flake-parts~\cite{flakeparts}, and dozens of other ecosystems,
collectively managing configurations of arbitrary complexity.
The one aspect that appears not to have been fully addressed is the
semantics of $\this$ under multiple inheritance.
When two independent modules that inherit from the same
deferred module are composed, each import introduces its own
declaration of the shared options; the module system's
\texttt{merge\-Option\-Decls} rejects duplicate declarations
with a static error (Appendix~\ref{app:scala-multi-target}).
The module system cannot express that both routes contribute
equally to the same option, that is, the multi-target situation that
recursive merging naturally gives rise to, which motivated the
observational semantics of Section~\ref{sec:mixin-trees}.
\fi

\subsection{\bilingual{Our Companion Work}{我们的配套工作}}

\paragraph{\bilingual{Modular software architectures}{模块化软件架构}}
\label{sec:practical-modular}
\ifInheritanceChinese
插件系统、组件框架与依赖注入框架通过声明式继承实现了软件的可扩展性。第~\ref{sec:mixin-trees} 节所确立的继承的对称性与可结合性解释了这类系统为何有效:组件可以以任意顺序继承而不改变行为。MIXINv2\anon[~(补充材料)]{~\cite{mixin2025}} 本身就是一个依赖注入框架:每个作用域将其依赖声明为具名槽,复合跨越作用域边界按名解析它们,从而使继承演算的模式显式化而非随意拼凑。
\else
Plugin systems, component frameworks, and dependency injection frameworks
enable software extensibility through declarative inheritance. The
symmetric, associative nature of inheritance established in Section~\ref{sec:mixin-trees} explains why such
systems work: components can be inherited in any order without changing
behavior.
MIXINv2\anon[~(supplementary material)]{~\cite{mixin2025}}
is itself a dependency injection framework:
each scope declares its dependencies as named slots, and
composition resolves them by name across scope boundaries,
making the inheritance-calculus patterns explicit rather than ad~hoc.
\fi

\paragraph{\bilingual{Union file systems}{联合文件系统}}
\label{sec:practical-union-fs}
\ifInheritanceChinese
联合文件系统,如 UnionFS、OverlayFS 与 AUFS,将多个目录层次叠加以呈现统一视图。其语义与继承演算相似,表现在后面的层覆盖前面的层且所有层的文件均可访问,但其设计在具体之处有所不同:继承是不对称的,标量文件的冲突解决是不可交换的,且层与层之间不存在晚绑定或动态分派。我们将继承演算直接实现为联合文件系统\anon[~(作为补充材料附上)]{~\cite{ratarmount2025}},仅凭三个原语就获得了文件查询的晚绑定语义与动态分派,无需任何额外机制。在软件包管理器、构建系统与操作系统发行版中,外部工具链将配置转换为文件树;而在这里,文件系统作为自身的配置,转换即是继承:这是一种\emph{文件系统即编译器},其中目录树既是源程序又是编译目标。
\else
Union file systems, such as UnionFS, OverlayFS, and AUFS, layer multiple directory
hierarchies to present a unified view. Their semantics resembles
inheritance-calculus in that later layers override earlier layers and files from all
layers remain accessible, but their designs differ in specific ways:
inheritance is asymmetric, conflict resolution for scalar files
is noncommutative, and there is no late-binding or dynamic
dispatch across layers.
We implemented inheritance-calculus directly as a union file
system\anon[~(included as supplementary material)]{~\cite{ratarmount2025}},
obtaining late-binding semantics and dynamic dispatch for file lookups
with no additional mechanism beyond the three primitives.
In package managers, build systems, and OS distributions,
external toolchains transform configuration into file trees;
here, the file system serves as configuration to itself,
and the transformation is inheritance: a
\emph{file-system-as-compiler} in which the directory tree is
both the source program and the compilation target.
\fi

\section{\bilingual{Future Work}{未来工作}}
\label{sec:future-work}

\ifInheritanceChinese
\anon[辅助实现]{MIXINv2~\cite{mixin2025}}{} 是继承演算的一个可执行实现,它已经超越了本文所呈现的无类型演算:它包含编译期检查,以验证所有引用都解析到 mixin 树中的有效路径,以及一个从宿主语言引入标量值的外部函数接口 (FFI)。以下所描述的未来工作,关注的是这一实现与一门完全实用的语言之间的差距,横跨理论基础与库层面的编码两个方面。
\else
\anon[The supplementary implementation]{MIXINv2~\cite{mixin2025}}{} is an executable implementation of inheritance-calculus
that already goes beyond the untyped calculus presented here:
it includes compile-time checking that all references resolve to
valid paths in the mixin tree, and a foreign-function interface (FFI)
that introduces scalar values from the host language.
The future work described below concerns the distance between this
implementation and a fully practical language, spanning both
theoretical foundations and library-level encodings.
\fi

\subsection{\bilingual{Type System}{类型系统}}
\ifInheritanceChinese
MIXINv2 现有的编译期检查验证每条引用路径都解析到继承结果中存在的属性。将这些检查的可靠性形式化,即证明类型良好的程序在运行时不产生悬空引用,是未来工作的一个方向。这样的形式化与 DOT 演算~\cite{amin2016-dependent-object-types} 的关系,类似于继承演算与 $\lambda$-演算的关系:有类型的继承演算可以看作不含 $\lambda$ 的 DOT,它保留了路径依赖类型,同时以 mixin 继承取代函数。正式的类型系统将适合在 Coq 和 Agda 等证明助手中机械化,以验证类型安全、进展和保持等性质。

在现有检查的可靠性之外,额外的类型系统特性将强化这门语言。一个关键的例子是\emph{完全性检查}:验证一个继承为所有必需槽位提供了实现,而不仅仅是引用能够解析。在无类型的继承演算中,写下 $\mathrm{magnitude} \mapsto \{\}$ 纯属说明性质:它定义了一个结构槽位,但不对填充它的内容施加任何约束。完全性检查将在静态上强制这类声明,拒绝使必需槽位保持未填充状态的继承。
\else
MIXINv2's existing compile-time checks verify that every
reference path resolves to a property that exists in the inherited
result. Formalizing the soundness of these checks, that is, proving that
well-typed programs do not produce dangling references at
runtime, is one direction of future work.
Such a formalization would relate to the DOT
calculus~\cite{amin2016-dependent-object-types} in a manner analogous to how
inheritance-calculus relates to the $\lambda$-calculus: a Typed
inheritance-calculus can be viewed as DOT without $\lambda$, retaining
path-dependent types while replacing functions with
mixin inheritance.
A formal type system would be amenable to mechanization in proof
assistants, such as Coq and Agda, for verifying type safety, progress, and
preservation properties.

Beyond soundness of the existing checks, additional type system
features would strengthen the language.
A key example is \emph{totality checking}: verifying that an
inheritance provides implementations for all required slots, not
merely that references resolve.
In untyped inheritance-calculus, writing $\mathrm{magnitude} \mapsto \{\}$
is purely documentary: it defines a structural slot but imposes no
constraint on what must fill it.
Totality checking would enforce such declarations statically,
rejecting inheritances that leave required slots unfilled.
\fi

\subsection{\bilingual{Standard Library}{标准库}}
\ifInheritanceChinese
其余几项涉及标准库:它们不需要对演算或类型系统进行扩展,但需要在现有框架内进行非平凡的编码。

继承演算\emph{默认是开放的}:继承可以自由合并任意两个 mixin,一个值可以同时具有多个数据构造子(例如,$\{\mathrm{Zero},\; \mathrm{Odd},\; \mathrm{half} \mapsto \{\mathrm{Zero}\}\}$)。这是该演算自然的 trie 语义,也是第~\ref{sec:expression-problem}~节求解表达式问题的基础。然而,许多操作(例如等值测试)假设值是\emph{线性的},即恰好具有一个数据构造子。第~\ref{sec:case-study}~节的观察者模式并不强制这一点:将其应用于一个具有多个数据构造子的值时,所有回调都会触发,其结果被继承在一起。

\emph{模式验证}可以在库层面使用现有原语来实现。观察者可以测试某个特定的数据构造子是否存在,返回一个布尔值;第二次布尔派发随后基于该结果进行分支。通过链接这样的测试,工厂可以验证一个值恰好匹配一个数据构造子并携带所需的属性。这并不限制继承本身(继承仍然不受约束),但提供了一种在不变式违反传播之前检测到它的方式。

同样的技术,即观察者返回布尔值后接布尔派发,可以产生\emph{封闭模式匹配}。第~\ref{sec:case-study}~节的观察者模式本质上是\emph{开放的}:新的数据构造子情况可以通过继承添加。封闭匹配(其中恰好执行一个分支)可以编码为一个\emph{访问者链}:一个由 if-then-else 节点组成的链表,每个节点测试一个数据构造子,在不匹配时向后传递。这类似于 GHC.Generics 的 $\mathrm{(:+:)}$ 和类型表示~\cite{magalhaes2010-generic-deriving},其中每个链节对应一个求和项。一个重要的推论是,这样的链\emph{不}可交换,因为它们具有固定的顺序,因此无法通过开放继承来扩展,而这正是封闭派发的正确语义。
\else
The remaining items concern the standard library: they require
no extensions to the calculus or type system, but involve nontrivial
encodings within the existing framework.

Inheritance-calculus is \emph{open by default}: inheritance
can freely merge any two mixins, and a value may simultaneously
inhabit multiple data constructors (e.g.,
$\{\mathrm{Zero},\; \mathrm{Odd},\; \mathrm{half} \mapsto \{\mathrm{Zero}\}\}$).
This is the natural trie semantics of the
calculus, and the basis for solving the Expression Problem
in Section~\ref{sec:expression-problem}.
However, many operations, such as equality testing, assume that values
are \emph{linear}, inhabiting exactly one data constructor.
The observer pattern of Section~\ref{sec:case-study} does not enforce
this: applied to a multi-data-constructor value, all callbacks fire and
their results are inherited.

\emph{Schema validation} can be implemented at the library level
using existing primitives.
An observer can test whether a particular data constructor is present,
returning a Boolean; a second Boolean dispatch then branches on the
result.
By chaining such tests, a factory can validate that a value matches
exactly one data constructor and carries the required properties.
This does not restrict inheritance itself, which remains
unconstrained, but provides a way to detect invariant violations
before they propagate.

The same technique, an observer returning Boolean followed by Boolean
dispatch, yields \emph{closed pattern matching}.
The observer pattern of Section~\ref{sec:case-study} is inherently
\emph{open}: new data-constructor cases can be added through inheritance.
Closed matching, where exactly one branch is taken, can be encoded as
a \emph{visitor chain}: a linked list of if-then-else nodes, each
testing one data constructor and falling through on mismatch.
This is analogous to GHC.Generics' $\mathrm{(:+:)}$ sum
representation~\cite{magalhaes2010-generic-deriving}, where each link
corresponds to one summand.
An important consequence is that such chains do \emph{not}
commute, as they have a fixed order, and therefore cannot be extended
through open inheritance, which is the correct semantics for closed
dispatch.
\fi

\section{\bilingual{Conclusion}{结论}}

\ifInheritanceChinese
对不可扩展性的免疫指导了继承演算的设计。第~\ref{sec:semantic-variants}~节排除了每一种替代的组合机制和引用解析,使得带有多目标晚绑定限定 $\this$ 的深合并成为唯一的幸存者。mixin 充当了统一模块、类、对象、方法、字段和局部绑定的单一抽象;mixin 树的深合并使继承具有可交换性、幂等性和可结合性,从而多重继承的线性化问题不再出现。第~\ref{sec:mixin-trees}~节的六个一阶方程,具有指称语义且可由 tabling~\cite{tamaki1986-tabled-resolution} 计算,足以定义语义;观察者通过查询一棵惰性构造的树中的路径来驱动计算。这些方程是对标准面向对象语义的最小偏离:唯一的改变是 $\overrides$ 方程~(\ref{eq:overrides}),它以子树的深合并取代了方法层面的遮蔽。面向对象编程传统上被归类为命令式的;去掉可变状态和函数体,便得到纯面向对象编程,%
\footnote{
  ``纯''具有双重含义:不纯粹意味着不纯洁,即值是不可变的且组合没有副作用;而不混合意味着组合完全依赖继承,没有函数式抽象,例如 $\lambda$ 或 $\beta$-归约。
}
这是声明式编程的一个最小计算模型。

由于 $\overrides$ 合并的是子树而非方法体,函数不再是原语。$\lambda$-演算在第~\ref{sec:translation}~节通过六条翻译规则宏表达到继承演算中。第~\ref{sec:bohm-tree}~节的 Böhm 树对应将这一一阶语义扩展到了 $\lambda$-演算:先前的高阶语义框架被证明忠实于继承演算一阶方程的一个子语言的首次迭代近似。

继承演算进一步获得了 $\lambda$-演算无法表达的三种现象;事实上,正如第~\ref{sec:asymmetry}~节所示,在 Felleisen 意义下~\cite{felleisen1991-expressive-power},继承演算严格地比惰性 $\lambda$-演算更具表达力。第~\ref{sec:expression-problem}~节在表达式问题~\cite{wadler1998-expression-problem}意义下建立了对不可扩展性的免疫。第~\ref{sec:case-study}~节的案例研究表明,普通算术产生了逻辑编程的关系语义:循环继承层次结构通过与递归 Datalog 所支配的相同的升链条件收敛,附录~\ref{app:datalog-encoding}给出了一个系统性的编码。第~\ref{sec:practical-function-color}~节展示了函数颜色盲性~\cite{nystrom2015-function-color}。

单个 mixin 同时是数据、程序、配置、函数、关系和对象。它宏表达了 $\lambda$-演算,并恢复了逻辑编程的关系语义,同时对两者均无法逃脱的不可扩展性保持免疫。

引言提出了配置本身是否能成为实用的通用编程的问题。答案是肯定的。我们提出继承演算作为实用通用编程的基础,我们相信它优于任何已知的计算模型。
\else
Immunity to nonextensibility guided the design of
inheritance-calculus.
Section~\ref{sec:semantic-variants} eliminated every
alternative composition mechanism and reference resolution,
leaving deep merge with multi-target late-binding qualified
$\this$ as the unique survivor.
The mixin served as the single abstraction unifying module,
class, object, method, field, and local binding;
deep merge of mixin trees made inheritance commutative, idempotent, and associative,
so the linearization problem of multiple inheritance did not arise.
The six first-order equations of Section~\ref{sec:mixin-trees},
denotational and computable by
tabling~\cite{tamaki1986-tabled-resolution}, suffice to
define the semantics; the observer drives computation by
querying paths in a lazily constructed tree.
These equations are a minimal departure from standard
object-oriented semantics: the only change is the $\overrides$ equation~(\ref{eq:overrides}),
which replaces method-level shadowing with deep merge of subtrees.
Object-oriented programming is traditionally classified as
imperative; removing mutable state and function bodies yields
pure object-oriented programming,%
\footnote{
  ``Pure'' carries a double meaning:
  not impure, since values are immutable and composition has no
  side effects; and not mixed, since composition relies entirely on
  inheritance, without functional abstractions such
  as~$\lambda$ or~$\beta$-reduction.
}
a minimal computational model for declarative programming.

Since $\overrides$ merges subtrees rather than method bodies,
functions are no longer primitive.
The $\lambda$-calculus macro-expressed into
inheritance-calculus in six translation rules in
Section~\ref{sec:translation}.
The B\"ohm tree correspondence of
Section~\ref{sec:bohm-tree} extended this first-order
semantics to the $\lambda$-calculus:
the prior higher-order semantic frameworks proved faithful
to the first-iteration approximation of a sublanguage of
inheritance-calculus's first-order equations.

Inheritance-calculus further gains three phenomena
that the $\lambda$-calculus cannot express;
indeed, inheritance-calculus is strictly more expressive
than the lazy $\lambda$-calculus in Felleisen's
sense~\cite{felleisen1991-expressive-power},
as shown in Section~\ref{sec:asymmetry}.
Section~\ref{sec:expression-problem} established immunity
to nonextensibility in the sense of the
Expression Problem~\cite{wadler1998-expression-problem}.
The case study of Section~\ref{sec:case-study} showed that
ordinary arithmetic yields the relational semantics of logic
programming: cyclic inheritance hierarchies converge by
the same ascending chain condition that governs recursive
Datalog, and Appendix~\ref{app:datalog-encoding} gave a
systematic encoding.
Section~\ref{sec:practical-function-color} exhibited
function color
blindness~\cite{nystrom2015-function-color}.

A single mixin is at once data, program, configuration,
function, relation, and object.
It macro-expresses the $\lambda$-calculus and recovers the
relational semantics of logic programming, while remaining
immune to the nonextensibility that neither escapes.

The introduction asked whether configuration alone can be
practical general-purpose programming.
It can.
We propose inheritance-calculus as a foundation for practical
general-purpose programming, one we believe better than any
known computational model.
\fi

\section{\bilingual{Data Availability Statement}{数据可用性声明}}

\ifInheritanceChinese
可执行实现 (MIXINv2) 以及本文中的所有示例以 MIT 许可证作为开源软件公开发布\anon[{} (匿名审查期间省略 URL)]{~\cite{mixin2025}}。每个以继承演算记法呈现的示例都有对应的可执行 MIXINv2 源码,且一套测试套件复现了论文中的每个示例。案例研究 (第~\ref{sec:case-study}~节) 的布尔逻辑和一元自然数算术、附录~\ref{app:datalog-encoding}~的 Datalog 编码 (包括双跳、常量连接和多变量连接关系程序以及重复访问者自连接回归示例),以及函数颜色盲性段落 (第~\ref{sec:practical-function-color}~节) 中讨论的同步和异步运行时实例,均包含在这些源码中。第~\ref{sec:practical-union-fs}~节所讨论的 ratarmount 补丁\anon[也已公开发布 (匿名审查期间省略 URL)]{作为一个开放的拉取请求~\cite{ratarmount2025}可以获取}。\anon[以上所有内容也作为辅助材料一并附上。]{}
\else
The executable implementation (MIXINv2) and all examples from
this paper are openly available as open-source software under
the MIT license\anon[{} (URL omitted for anonymous review)]{~\cite{mixin2025}}.
Every example rendered in inheritance-calculus notation has a
corresponding executable MIXINv2 source, and a test suite
reproduces every example in the paper.
The boolean logic and unary natural number arithmetic of the case
study (Section~\ref{sec:case-study}), the Datalog encodings of
Appendix~\ref{app:datalog-encoding} (including the two-hop,
constant-join, and multi-variable-join relational programs and the
repeated-visitor self-join regression example), and the synchronous
and asynchronous runtime instances discussed in the
function-color-blindness paragraph
(Section~\ref{sec:practical-function-color}) are all included among
these sources.
The ratarmount patch discussed in
Section~\ref{sec:practical-union-fs} is
\anon[also openly available (URL omitted for anonymous review)]{available as an open pull
request~\cite{ratarmount2025}}.
\anon[All of the above are also included as supplementary material.]{}
\fi

\begin{acks}
  Claude Code, GitHub Copilot, Gemini Code Assist, and Codex were utilized to generate sections of this Work,
  including text, tables, graphs, code, formulas, citations, etc.
\end{acks}
\fi % end \ifInheritanceBody (the main matter)

% References follow the main matter and precede the appendix. This copy fires
% whenever the main matter is present (preprint, submission); the appendix-only
% build emits its copy after the appendix instead.
\ifInheritanceBody
\bibliographystyle{ACM-Reference-Format}
\bibliography{references}
\fi

\ifInheritanceAppendix
\appendix

\section{\bilingual{Well-Definedness of the Semantic Functions}{语义函数的良定义性}}
\label{app:well-definedness}

\ifInheritanceChinese
六个语义方程~(\ref{eq:properties})--(\ref{eq:this}) 是一阶的:其输入与输出均为路径及路径的集合,不涉及函数空间、闭包或高阶构造。
本节通过证明这六个方程定义了一个单调算子来确立良定义性;该算子的最小不动点由 tabling 求值~\cite{tamaki1986-tabled-resolution} 可靠地计算。
\else
The six semantic
equations~(\ref{eq:properties})--(\ref{eq:this}) are
first-order: their inputs and outputs are paths and sets of
paths, with no function spaces, closures, or higher-order
constructs.
This section establishes well-definedness by showing that
the six equations define a monotone operator whose least
fixed point is computed soundly by tabled
evaluation~\cite{tamaki1986-tabled-resolution}.
\fi

\begin{definition}[\bilingual{Dependency graph}{依赖图}]\label{def:dependency-graph}
\ifInheritanceChinese
  固定一棵 AST 及其原始函数 $\defines$ 与 $\inherits$。
  \emph{依赖图}~$G$ 以由该 AST 产生的所有语义函数应用
  ($\properties(p)$、\allowbreak$\supers(p)$、\allowbreak$\overrides(p)$、\allowbreak$\bases(p)$、\allowbreak
  $\resolve(\ldots)$、\allowbreak$\this(\ldots)$) 为顶点。
  当通过方程~(\ref{eq:properties})--(\ref{eq:this}) 计算~$u$
  需要~$v$ 的值时,从顶点~$u$ 到顶点~$v$ 有一条有向边。
\else
  Fix an AST with its primitive functions $\defines$ and
  $\inherits$.
  The \emph{dependency graph}~$G$ has as vertices all
  semantic-function applications
  ($\properties(p)$, $\supers(p)$, $\overrides(p)$,
    $\bases(p)$, $\resolve(\ldots)$,
  $\this(\ldots)$) that arise from the AST\@.
  There is a directed edge from vertex~$u$ to vertex~$v$
  whenever computing~$u$ via
  equations~(\ref{eq:properties})--(\ref{eq:this})
  requires the value of~$v$.
\fi
\end{definition}

\begin{definition}[\bilingual{Immediate consequence operator $T_P$}{直接后果算子 $T_P$}]%
  \label{def:tp-operator}
\ifInheritanceChinese
  设 $\mathcal{Q}$ 为所有语义函数查询(即~$G$ 的顶点)的集合。
  对每个查询 $q \in \mathcal{Q}$,设 $\mathcal{V}_q$ 为形状与~$q$ 相匹配的答案值的集合
  (依定义~$q$ 的方程,为路径、路径的集合或路径对的集合)。
  \emph{查询--答案映射}是逐点乘积
  $\prod_{q \in \mathcal{Q}} \mathcal{P}(\mathcal{V}_q)$ 中的一个元素;
  等价地,它为每个查询~$q$ 指派一个候选答案集合 $M(q) \subseteq \mathcal{V}_q$。
  \emph{直接后果算子}~\cite{vanemden1976-predicate-logic-semantics}
  $T_P$ 将查询--答案映射~$M$ 映射为新的查询--答案映射~$T_P(M)$,方式是:对每个查询~$q$,
  求值定义~$q$ 的那条方程的右侧,并以当前近似~$M(q')$ 解释每个子查询~$q'$。
\else
  Let $\mathcal{Q}$ be the set of all
  semantic-function queries
  (vertices of~$G$).
  For each query $q \in \mathcal{Q}$, let
  $\mathcal{V}_q$ be the set of answer values of the
  appropriate shape for~$q$
  (paths, sets of paths, or sets of pairs of paths,
  depending on the equation defining~$q$).
  A \emph{query--answer map} is an element of the
  pointwise product
  $\prod_{q \in \mathcal{Q}} \mathcal{P}(\mathcal{V}_q)$;
  equivalently, it assigns to each query~$q$ a set
  $M(q) \subseteq \mathcal{V}_q$ of candidate answers.
  The \emph{immediate consequence
  operator}~\cite{vanemden1976-predicate-logic-semantics}
  $T_P$ maps a query--answer map~$M$ to a new
  query--answer map~$T_P(M)$ by evaluating, for each
  query~$q$, the right-hand side of the equation that
  defines~$q$, interpreting every sub-query~$q'$ by
  the current approximation~$M(q')$.
\fi
\end{definition}

\begin{lemma}[\bilingual{Monotonicity}{单调性}]\label{lem:monotonicity}
  \ifInheritanceChinese
  算子 $T_P$ 在完备格
  $\Bigl(\prod_{q \in \mathcal{Q}} \mathcal{P}(\mathcal{V}_q),\; \subseteq\Bigr)$
  上是单调的。
  \else
  The operator $T_P$ is monotone on the complete lattice
  $\Bigl(\prod_{q \in \mathcal{Q}} \mathcal{P}(\mathcal{V}_q),\; \subseteq\Bigr)$.
  \fi
\end{lemma}

\begin{proof}
\ifInheritanceChinese
  每条方程的右侧均由 $\cup$、$\in$、等式测试以及正集合
  概括 $\{\, x \mid x \in S, \ldots \,\}$ 构成。
  六条方程均不使用否定、集合差或补集。
  因此,若 $M \subseteq M'$(逐点),则 $T_P(M) \subseteq T_P(M')$。
\else
  Each equation's right-hand side is built from
  $\cup$, $\in$, equality tests, and positive set
  comprehension $\{\, x \mid x \in S, \ldots \,\}$.
  None of the six equations uses negation, set
  difference, or complement.
  Therefore, if $M \subseteq M'$ (pointwise), then
  $T_P(M) \subseteq T_P(M')$.
\fi
\end{proof}

\begin{theorem}[\bilingual{Existence and uniqueness of
  $\mathrm{lfp}(T_P)$}{$\mathrm{lfp}(T_P)$ 的存在性与唯一性}]\label{thm:lfp-exists}
  \ifInheritanceChinese
  算子 $T_P$ 有唯一的最小不动点 $\mathrm{lfp}(T_P)$。
  \else
  The operator $T_P$ has a unique least fixed point
  $\mathrm{lfp}(T_P)$.
  \fi
\end{theorem}

\begin{proof}
\ifInheritanceChinese
  每个幂集格
  $(\mathcal{P}(\mathcal{V}_q), \subseteq)$
  都是完备的,
  因此逐点乘积格
  $\Bigl(\prod_{q \in \mathcal{Q}} \mathcal{P}(\mathcal{V}_q), \subseteq\Bigr)$
  也是完备的。
  由 Knaster--Tarski 定理~\cite{tarski1955-lattice-fixpoint-theorem},
  完备格上的每个单调函数都有唯一的最小不动点。
\else
  Each powerset lattice
  $(\mathcal{P}(\mathcal{V}_q), \subseteq)$
  is complete,
  so the pointwise product lattice
  $\Bigl(\prod_{q \in \mathcal{Q}} \mathcal{P}(\mathcal{V}_q), \subseteq\Bigr)$
  is also complete.
  By the Knaster--Tarski
  theorem~\cite{tarski1955-lattice-fixpoint-theorem},
  every monotone function on a complete lattice has a
  unique least fixed point.
\fi
\end{proof}

\begin{theorem}[\bilingual{Well-definedness}{良定义性}]\label{thm:well-defined}
  \ifInheritanceChinese
  语义函数 $\properties$、\allowbreak$\supers$、\allowbreak$\overrides$、\allowbreak$\bases$、\allowbreak
  $\resolve$、\allowbreak$\this$ 是良定义的:
  其指称语义为 $\mathrm{lfp}(T_P)$,即直接后果算子的唯一最小不动点。
  \else
  The semantic functions
  $\properties$, $\supers$, $\overrides$, $\bases$,
  $\resolve$, $\this$ are well-defined:
  their denotational semantics is
  $\mathrm{lfp}(T_P)$, the unique least fixed point of
  the immediate consequence operator.
  \fi
\end{theorem}

\begin{theorem}[\bilingual{Tabling soundness}{Tabling 可靠性}]%
  \label{thm:tabling-soundness}
  \ifInheritanceChinese
  即使依赖图~$G$ 是无穷的,方程~(\ref{eq:properties})--(\ref{eq:this}) 的
  tabling 求值也绝不会返回错误答案。具体而言:
  \begin{enumerate}
    \item \textbf{可靠性(无假正例)。}
      tabling 返回的每个事实均由某条方程的右侧在其输入上求值产生。
      对求值轨迹进行归纳,每个返回的事实都属于 $\mathrm{lfp}(T_P)$。

    \item \textbf{单调积累。}
      tabling 累加器在各轮迭代之间只通过集合并增长,任何事实都不会被删除。

    \item \textbf{不动点检测的正确性。}
      若求值过程中没有子查询被重入,则所有子查询均在无近似的情况下被完整计算,结果是精确的。
      若发生了重入但累加器是稳定的(即没有新事实被添加),则累加器等于
      $\mathrm{lfp}(T_P)$ 在可达查询上的限制,因为它是一个从~$\bot$ 上升到达的前不动点。

    \item \textbf{可靠的失败。}
      当 tabling 无法确定答案时(无论是因为迭代预算耗尽还是因为存在无穷非循环链),
      它会抛出错误而非返回错误答案。
  \end{enumerate}
  \else
  Tabled evaluation of
  equations~(\ref{eq:properties})--(\ref{eq:this})
  never returns a wrong answer, even when the
  dependency graph~$G$ is infinite.
  Specifically:
  \begin{enumerate}
    \item \textbf{Soundness (no false positives).}
      Every fact returned by tabling was produced by
      evaluating some equation's right-hand side on its
      inputs.
      By induction on the evaluation trace, every
      returned fact belongs to $\mathrm{lfp}(T_P)$.

    \item \textbf{Monotone accumulation.}
      The tabling accumulator only grows between
      iterations via set union; no fact is ever
      removed.

    \item \textbf{Correct fixpoint detection.}
      If no sub-query was re-entered during evaluation,
      all sub-queries were fully computed with no
      approximation, and the result is exact.
      If re-entry occurred but the accumulator is
      stable, meaning no new facts were added, the
      accumulator equals
      $\mathrm{lfp}(T_P)$ restricted to the reachable
      queries, because it is a pre-fixed point reached
      by ascending from~$\bot$.

    \item \textbf{Sound failure.}
      When tabling cannot determine the answer,
      whether because the iteration budget is exhausted
      or because of an infinite non-cyclic chain,
      it raises an error rather than
      returning a wrong answer.
  \end{enumerate}
  \fi
\end{theorem}

\begin{proof}
\ifInheritanceChinese
  对于~(1),每个返回的事实都是某条方程的右侧作用于先前返回的事实上的输出;
  归纳地,所有输入均属于 $\mathrm{lfp}(T_P)$,因此输出也是(由单调性和不动点性质)。

  对于~(2),实现通过 $A \leftarrow A \cup T_P(A)$ 进行积累,这对~$A$ 是单调的。

  对于~(3),当没有重入时,求值是有限无环可达子图上的普通记忆化递归,
  因此每个子查询都被精确计算,不涉及任何近似。
  当发生重入时,设 $R \subseteq \mathcal{Q}$ 为从初始查询可达的查询集合,
  设 $T_P^R$ 表示 $T_P$ 到~$R$ 上映射的限制。
  若累加器稳定,即 $A = A \cup T_P^R(A)$,则
  $T_P^R(A) \subseteq A$,从而~$A$ 是~$T_P^R$ 的一个前不动点。

  由~(1),曾被添加到累加器中的每个事实均由某次右侧求值产生,
  因此属于 $\mathrm{lfp}(T_P^R)$。
  故 $A \subseteq \mathrm{lfp}(T_P^R)$。

  反之,由 Knaster--Tarski 将最小不动点刻画为最小前不动点,
  $\mathrm{lfp}(T_P^R)$ 包含于~$T_P^R$ 的每个前不动点中,特别地包含于~$A$ 中。
  从而 $\mathrm{lfp}(T_P^R) \subseteq A$。

  因此 $A = \mathrm{lfp}(T_P^R)$,即稳定的累加器等于可达查询上的精确指称。

  对于~(4),预算耗尽与栈溢出均产生异常,而非静默的错误答案。
\else
  For~(1), each returned fact is the output of some
  equation's right-hand side applied to previously
  returned facts; by induction, all inputs are in
  $\mathrm{lfp}(T_P)$, so the output is too (by
  monotonicity and the fixed-point property).

  For~(2), the implementation accumulates via
  $A \leftarrow A \cup T_P(A)$, which is monotone in~$A$.

  For~(3), when no reentry occurs, the evaluation is
  ordinary memoised recursion on a finite acyclic
  reachable subgraph, so every sub-query is computed
  exactly and no approximation is involved.
  When reentry occurs, let $R \subseteq \mathcal{Q}$ be
  the set of queries reachable from the initial query,
  and let $T_P^R$ denote the restriction of $T_P$ to
  maps on~$R$.
  If the accumulator is stable, that is,
  $A = A \cup T_P^R(A)$, then
  $T_P^R(A) \subseteq A$, so $A$ is a pre-fixed point
  of~$T_P^R$.

  By~(1), every fact ever added to the accumulator was
  produced by some right-hand side evaluation and
  therefore belongs to $\mathrm{lfp}(T_P^R)$.
  Hence $A \subseteq \mathrm{lfp}(T_P^R)$.

  Conversely, by the Knaster--Tarski characterization
  of the least fixed point as the least pre-fixed
  point, $\mathrm{lfp}(T_P^R)$ is contained in every
  pre-fixed point of~$T_P^R$, in particular in~$A$.
  Thus $\mathrm{lfp}(T_P^R) \subseteq A$.

  Therefore $A = \mathrm{lfp}(T_P^R)$, that is, the stable
  accumulator equals the exact denotation on the
  reachable queries.

  For~(4), budget exhaustion and stack overflow both
  produce exceptions, never silent wrong answers.
\fi
\end{proof}

% 终止性依赖程序结构:无环有限子图→一趟即止;有环有限逐查询值域→有限格上升链条件收敛;
% 无穷可达子图→固有发散,属图灵完备性,非求值策略的缺陷。
\paragraph{\bilingual{Termination}{终止性}}
\ifInheritanceChinese
求值是否终止取决于程序:
\begin{itemize}
  \item \textbf{无环、有限可达子图。}
    一趟即足;求值终止。
  \item \textbf{有环、逐查询值域有限。}
    多趟迭代在有限格的上升链条件下收敛~\cite{bancilhon1986-recursive-query-strategies}。
    这是 Datalog 的情形(第~\ref{sec:recursive-rules} 节)。
  \item \textbf{无穷可达子图。}
    求值可能发散;这是图灵完备性的固有属性,而非求值策略的缺陷。
\end{itemize}
\else
Whether evaluation terminates depends on the
program:
\begin{itemize}
  \item \textbf{Acyclic, finite reachable subgraph.}
    One pass suffices; evaluation terminates.
  \item \textbf{Cyclic, finite per-query domain.}
    Multiple passes converge by the ascending chain
    condition on a finite
    lattice~\cite{bancilhon1986-recursive-query-strategies}.
    This is the Datalog case
    (Section~\ref{sec:recursive-rules}).
  \item \textbf{Infinite reachable subgraph.}
    Evaluation may diverge; this is inherent to
    Turing completeness and not a defect of the
    evaluation strategy.
\end{itemize}
\fi

\section{B\"ohm Tree Correspondence: Proofs}
\label{app:bohm-tree-proofs}

This appendix proves two results about the sublanguage of
inheritance-calculus corresponding to the $\lambda$-calculus
(Section~\ref{sec:bohm-tree}).
First, the first-iteration approximation
$T_P\!\uparrow\!1$ is fully abstract with respect to
B\"ohm tree equivalence
(Theorems~\ref{thm:adequacy} and~\ref{thm:first-order-semantics}),
establishing that the first iteration coincides with
the lazy $\lambda$-calculus.
Second, the full $\mathrm{lfp}(T_P)$ semantics is strictly
more defined in the powerset sense of
Figure~\ref{fig:calculi-matrix}: translated $\lambda$-terms
that diverge under $\beta$-reduction may converge to definite
sets, as shown in Section~\ref{sec:bohm-tree}.

\subsection{B\"ohm Trees}

We briefly recall the definition of B\"ohm
trees~\cite{barendregt1984-lambda-calculus}.
A \emph{head reduction} $M \to_h M'$ contracts
the outermost $\beta$-redex only: if $M$ has the form
$(\lambda x.\, e)\; v\; M_1 \cdots M_k$, then
$M \to_h e[v/x]\; M_1 \cdots M_k$.
A term $M$ is in \emph{head normal form} (HNF) if it has no
head redex, that is, $M = \lambda x_1 \ldots x_m.\;
y\; M_1 \cdots M_k$ where $y$ is a variable.

The \emph{B\"ohm tree} $\mathrm{BT}(M)$ of a $\lambda$-term $M$
is an infinite labeled tree defined by:
\[
  \mathrm{BT}(M) =
  \begin{cases}
    \bot & \text{if } M \text{ has no HNF} \\[6pt]
    \lambda x_1 \ldots x_m.\;
    y\bigl(\mathrm{BT}(M_1),\; \ldots,\; \mathrm{BT}(M_k)\bigr)
    &
    \begin{aligned}[t]
      &\text{if } M \to^*_h \\[-2pt]
      &\lambda x_1 \ldots x_m.\; y\; M_1 \cdots M_k
    \end{aligned}
  \end{cases}
\]
B\"ohm tree equivalence $\mathrm{BT}(M) = \mathrm{BT}(N)$ is the
standard observational equivalence of the lazy
$\lambda$-calculus~\cite{abramsky1993-full-abstraction-lazy-lambda,abramsky1995-games-full-abstraction-lazy-lambda}.

\subsection{Path Encoding}

We define a correspondence between positions in the B\"ohm tree
and paths in the mixin tree.
When the translation $\mathcal{T}$ is applied to an ANF
$\lambda$-term, the resulting mixin tree has a specific shape:
\begin{itemize}
  \item An abstraction $\lambda x.\, M$ translates to a record
    with own properties $\{\mathrm{argument}, \mathrm{result}\}$,
    which is referred to as the \emph{abstraction shape}.
  \item A $\mathbf{let}$-binding
    $\mathbf{let}\; x = V_1\; V_2 \;\mathbf{in}\; M$ translates
    to a record with own properties $\{x, \mathrm{result}\}$,
    where $x$ holds the encapsulated application
    $\{\mathcal{T}(V_1),\; \mathrm{argument} \mapsto
    \mathcal{T}(V_2)\}$.
  \item A tail call $V_1\; V_2$ translates to a record with own
    properties $\{\mathrm{tailCall}, \mathrm{result}\}$, where
    $\mathrm{tailCall}$ holds the encapsulated application and
    $\mathrm{result}$ projects
    $\mathrm{tailCall}.\mathrm{result}$.
  \item A variable $x$ with de~Bruijn index $n$ translates to
    an indexed reference $\dbi{n}.\mathrm{argument}$.
\end{itemize}

\noindent
Each B\"ohm tree position is a sequence of navigation steps from
the root.
In the B\"ohm tree, the possible steps at a node
$\lambda x_1 \ldots x_m.\; y\; M_1 \cdots M_k$ are:
entering the body by peeling off one abstraction layer, and
entering the $i$-th argument $M_i$ of the head variable.
In the mixin tree, these correspond to following the
$\mathrm{result}$ label (for entering the body) and the
$\mathrm{argument}$ label after $\this$ resolution (for
entering an argument position).

Rather than comparing absolute paths across different mixin
trees, which would be sensitive to internal structure, we
observe only the \emph{convergence} behavior accessible via
the $\mathrm{result}$ projection
(Definition~\ref{def:convergence}).
Each $\mathrm{result}$-following step corresponds to one head
reduction step in the $\lambda$-calculus: a tail call's
$\mathrm{result}$ projects through the encapsulated
application, and a $\mathbf{let}$-binding's $\mathrm{result}$
enters the continuation.
The abstraction shape at the end signals that a head
normal form has been reached.

\subsection{Single-Path Lemma}

The $\this$ function (equation~\ref{eq:this}) is designed for the
general case of multi-target $\this$ resolution.
For translated $\lambda$-terms, we show it degenerates to
single-target resolution.

\begin{lemma}[Single-Path]\label{lem:single-path}
  For any closed ANF $\lambda$-term $M$, at every invocation of
  $\this(S,\; p_{\mathrm{def}},\; n)$ during evaluation of
  $\mathcal{T}(M)$, the frontier set $S$ contains exactly one
  path.
\end{lemma}

\begin{proof}
  By case analysis on the translation rules of~$\mathcal{T}$.
  Since $\mathcal{T}$ is compositional, the structural invariant
  established for each rule propagates to all inherited subtrees:
  every subtree appearing in the mixin tree is itself the image of
  some sub-term under~$\mathcal{T}$, so the invariant holds
  throughout the runtime composition.

  \paragraph{Case $\mathcal{T}(\lambda x.\, M)
    = \{\mathrm{argument} \mapsto \{\},\;
  \mathrm{result} \mapsto \mathcal{T}(M)\}$}
  This is a record literal with
  $\defines = \{\mathrm{argument}, \mathrm{result}\}$ and
  no inheritance sources, so $\inherits = \varnothing$.
  At the root path $p$, we have $\overrides(p) = \{p\}$ (a single
  path) since there are no inheritance sources to introduce
  additional branches.
  Inside $\mathcal{T}(M)$, any reference to $x$ is
  $\dbi{n}.\mathrm{argument}$ where $n$ is the de~Bruijn index
  pointing to this scope level.
  Since $\mathcal{T}(\lambda x.\, M)$ introduces exactly one
  scope level, there is exactly one
  mixin at that level, hence $\this$ finds exactly one matching
  pair $(p_{\mathrm{site}},\; p_{\mathrm{override}})$ in $\supers$.

  \paragraph{Case
    $\mathcal{T}(\mathbf{let}\; x = V_1\; V_2
    \;\mathbf{in}\; M)
    = \{x \mapsto \{\mathcal{T}(V_1),\;
      \mathrm{argument} \mapsto \mathcal{T}(V_2)\},\;
  \mathrm{result} \mapsto \mathcal{T}(M)\}$}
  The outer record has own properties $\{x, \mathrm{result}\}$
  and no inheritance sources at the outer level
  ($\inherits = \varnothing$), so $\overrides(\text{root})
  = \{\text{root}\}$.
  Inside the $x$ subtree, the inheritance
  $\{\mathcal{T}(V_1),\; \mathrm{argument} \mapsto
  \mathcal{T}(V_2)\}$ has exactly one inheritance source
  ($\mathcal{T}(V_1)$), creating exactly one $\bases$ entry.
  By induction on $V_1$, the single-path property holds inside
  $\mathcal{T}(V_1)$.
  Inside $\mathcal{T}(M)$, the induction hypothesis applies
  directly.

  \paragraph{Case
    $\mathcal{T}(\mathbf{let}\; x = V
    \;\mathbf{in}\; M)
    = \{x \mapsto \{\mathrm{result} \mapsto \mathcal{T}(V)\},\;
  \mathrm{result} \mapsto \mathcal{T}(M)\}$}
  The outer record has own properties $\{x, \mathrm{result}\}$
  and no inheritance sources at the outer level, so
  $\overrides(\text{root}) = \{\text{root}\}$.
  The $x$ subtree is itself a record literal with own
  property $\{\mathrm{result}\}$ and no inheritance
  sources.
  Thus neither entering the bound value nor entering the
  body introduces branching in $\supers$.
  By induction, the single-path property holds inside both
  $\mathcal{T}(V)$ and $\mathcal{T}(M)$.

  \paragraph{Case $\mathcal{T}(V_1\; V_2) =
    \{\mathrm{tailCall} \mapsto
      \{\mathcal{T}(V_1),\; \mathrm{argument} \mapsto
      \mathcal{T}(V_2)\},\;
      \mathrm{result} \mapsto
  \{\mathrm{tailCall}.\mathrm{result}\}\}$}
  Identical to the $\mathbf{let}$-binding case: the outer record
  has own properties $\{\mathrm{tailCall}, \mathrm{result}\}$
  with no inheritance sources at the outer level, and the
  $\mathrm{tailCall}$ subtree contains exactly one inheritance
  source.

  \paragraph{Case $\mathcal{T}(x)
  = \dbi{n}.\mathrm{argument}$}
  This is a reference, not a record.
  The $\this$ function walks up $n$ steps from the enclosing scope.
  By the compositionality argument above, every scope level
  encountered during the walk, including inherited subtrees
  reached through reference resolution, is the image of some
  sub-term under~$\mathcal{T}$, so each step of $\this$
  encounters exactly one matching pair in $\supers$.

  \medskip\noindent
  The key structural invariant is the following.
  For every subtree generated by~$\mathcal{T}$,
  its root has either no inheritance source or exactly one
  inheritance source.
  Moreover, the only roots with an inheritance source are the
  application wrappers introduced by the application
  $\mathbf{let}$-binding and tail-call rules, and in both cases
  that source is unique and encapsulated under a distinguished
  child ($x$ or $\mathrm{tailCall}$), never at the enclosing
  scope level.
  Therefore every scope level traversed by $\this$ contributes
  at most one matching pair in $\supers$, so the frontier set
  remains a singleton throughout evaluation.
\end{proof}

\subsection{Substitution Lemma}

The core mechanism of the translation is that inheritance with
$\mathrm{argument} \mapsto \mathcal{T}(V)$ plays the role of
substitution.  We make this precise.

\begin{lemma}[Substitution]\label{lem:substitution}
  Let $M$ be an ANF term in which $x$ may occur free,
  and let $V$ be a closed value.%
  \footnote{
    The proof uses the Single-Path Lemma
    (Lemma~\ref{lem:single-path}), which requires the inherited tree
    to occur within a closed term. This holds whenever
    $\lambda x.\, M$ applied to $V$ is a subterm of a closed ANF
    program.
  }
  Define the \emph{inherited tree}
  $C = \{\mathcal{T}(\lambda x.\, M),\;
  \mathrm{argument} \mapsto \mathcal{T}(V)\}$.
  Then for every path $p$ under $C \snoc \mathrm{result}$ and
  every label $\ell$:
  \[
    \ell \in \properties(p)
    \text{ in } C \snoc \mathrm{result}
    \quad\Longleftrightarrow\quad
    \ell \in \properties(p)
    \text{ in } \mathcal{T}(M[V/x])
  \]
  where $M[V/x]$ is the usual capture-avoiding substitution.
\end{lemma}

\begin{proof}
  By structural induction on $M$.

  \paragraph{Base case: $M = x$}
  Then $\mathcal{T}(x) =
  \dbi{n}.\mathrm{argument}$ (where $n$ is the de~Bruijn index
  of $x$), and
  $C \snoc \mathrm{result}$ contains this reference.
  In the inherited tree $C$, $\resolve$ and $\this$ resolve the reference by
  navigating from the reference's enclosing scope up $n$ steps
  to the binding $\lambda$, where the inheritance
  $\{\mathcal{T}(\lambda x.\, x),\;
  \mathrm{argument} \mapsto \mathcal{T}(V)\}$ inherits into the
  $\mathrm{argument}$ slot, providing $\mathcal{T}(V)$.
  Since $C$ is the $\mathrm{tailCall}$ subtree of
  $\mathcal{T}\bigl((\lambda x.\, M)\; V\bigr)$, which is a
  closed ANF term, Lemma~\ref{lem:single-path} applies and
  $\this$ finds exactly one path,
  so $\resolve$ returns $\mathcal{T}(V)$'s subtree.
  Since $M[V/x] = V$, we have
  $\mathcal{T}(M[V/x]) = \mathcal{T}(V)$, and the properties
  coincide.

  \paragraph{Base case: $M = y$ where $y \neq x$ and $y$ is
  $\lambda$-bound}
  Then $\mathcal{T}(y) = \dbi{m}.\mathrm{argument}$ where $m$
  is the de~Bruijn index of~$y$.
  This walks up to the scope that binds~$y$, which is different
  from the scope that binds~$x$, so the inheritance with
  $\mathrm{argument} \mapsto \mathcal{T}(V)$ at $x$'s binding
  scope has no effect.
  Since $M[V/x] = y$, the properties are identical.

  \paragraph{Base case: $M = y$ where $y \neq x$ and $y$ is
  let-bound}
  Then $\mathcal{T}(y) = y.\mathrm{result}$.
  This is a lexical reference to the sibling property~$y$ in the
  enclosing record, so the inherited tree~$C$ does not alter its
  resolution.
  Since $M[V/x] = y$, the properties are identical.

  \paragraph{Inductive case: $M = \lambda y.\, M'$}
  Then $\mathcal{T}(M) =
  \{\mathrm{argument} \mapsto \{\},\;
  \mathrm{result} \mapsto \mathcal{T}(M')\}$.
  The subtree at $C \snoc \mathrm{result}$ is another
  abstraction-shaped record.
  References to $x$ inside $\mathcal{T}(M')$ resolve through
  $\this$ in exactly the same way (with one additional scope
  level from the inner $\lambda$), and by induction on $M'$, the
  properties under
  $C \snoc \mathrm{result} \snoc \mathrm{result}$
  coincide with those of
  $\mathcal{T}(M'[V/x])$.
  Since
  $(\lambda y.\, M')[V/x] = \lambda y.\, (M'[V/x])$
  (assuming $y$ is fresh), the result follows.

  \paragraph{Inductive case:
    $M = \mathbf{let}\; z = V_1\; V_2$
  $\mathbf{in}\; M'$}
  Then
  \[
    \mathcal{T}(M) =
    \{z \mapsto \{\mathcal{T}(V_1),\;
      \mathrm{argument} \mapsto \mathcal{T}(V_2)\},\;
    \mathrm{result} \mapsto \mathcal{T}(M')\}.
  \]
  The substitution distributes:
  $M[V/x] =
  \mathbf{let}\; z = V_1[V/x]\; V_2[V/x]
  \;\mathbf{in}\; M'[V/x]$.
  By induction on $V_1$, $V_2$, and $M'$, the properties under
  each subtree ($z$ and $\mathrm{result}$) coincide,
  since references to $x$ inside each subtree resolve to
  $\mathcal{T}(V)$ via the same $\this$ mechanism.

  \paragraph{Inductive case:
    $M = \mathbf{let}\; z = V'
  \;\mathbf{in}\; M'$}
  Then
  \[
    \mathcal{T}(M) =
    \{z \mapsto \{\mathrm{result} \mapsto \mathcal{T}(V')\},\;
    \mathrm{result} \mapsto \mathcal{T}(M')\}.
  \]
  The substitution distributes:
  $M[V/x] =
  \mathbf{let}\; z = V'[V/x]
  \;\mathbf{in}\; M'[V/x]$.
  By induction on $V'$ and $M'$, the properties under both
  subtrees, specifically $z$ and $\mathrm{result}$, coincide.
  In particular, occurrences of~$z$ in $M'$ are translated as the
  lexical reference $z.\mathrm{result}$ on both sides, so the
  wrapper introduced by the value $\mathbf{let}$-binding is
  preserved by the substitution.

  \paragraph{Inductive case: $M = V_1\; V_2$ (tail call)}
  Analogous to the $\mathbf{let}$-binding case with
  $\mathrm{tailCall}$ in place of $z$.
\end{proof}

\subsection{Convergence Preservation}

The convergence criterion (Definition~\ref{def:convergence})
follows the $\mathrm{result}$ chain from the root.
Each step in this chain corresponds to one head reduction step
in the $\lambda$-calculus.
We make this precise.

\begin{lemma}[Result-step]\label{lem:result-step}
  Let $(\lambda x.\, M)\; V$ be a tail call, which is a $\beta$-redex.
  Then for every path $p$ and label $\ell$:
  \[
    \ell \in \properties\!\bigl(\,
    \text{root} \snoc \mathrm{result} \snoc p\,\bigr)
    \text{ in } \mathcal{T}\bigl((\lambda x.\, M)\; V\bigr)
    \;\Longleftrightarrow\;
    \ell \in \properties\!\bigl(\,
    \text{root} \snoc p\,\bigr)
    \text{ in } \mathcal{T}(M[V/x])
  \]
  That is, following one $\mathrm{result}$ projection in the
  tail call's mixin tree yields the same properties as the
  root of the reduct's mixin tree.
\end{lemma}

\begin{proof}
  The translation gives:
  \[
    \mathcal{T}\bigl((\lambda x.\, M)\; V\bigr) =
    \{\mathrm{tailCall} \mapsto
      \{\mathcal{T}(\lambda x.\, M),\;
      \mathrm{argument} \mapsto \mathcal{T}(V)\},\;
      \mathrm{result} \mapsto
    \{\mathrm{tailCall}.\mathrm{result}\}\}
  \]
  The $\mathrm{result}$ label at the root is defined as
  $\mathrm{tailCall}.\mathrm{result}$, so
  $\text{root} \snoc \mathrm{result}$ resolves to the
  $\mathrm{result}$ path inside the inheritance
  $\mathrm{tailCall} = \{\mathcal{T}(\lambda x.\, M),\;
  \mathrm{argument} \mapsto \mathcal{T}(V)\}$.
  This inheritance is exactly the inherited tree $C$ of
  Lemma~\ref{lem:substitution}, and
  $\mathrm{tailCall} \snoc \mathrm{result}$ is
  $C \snoc \mathrm{result}$.
  By Lemma~\ref{lem:substitution}, the properties under
  $C \snoc \mathrm{result}$ coincide with those of
  $\mathcal{T}(M[V/x])$.

  \medskip\noindent
  \emph{Let-binding variant.}
  An analogous result holds for
  $\mathbf{let}\; x = V_1\; V_2 \;\mathbf{in}\; M'$
  where $V_1 = \lambda y.\, M''$.
  The translation gives:
  \[
    \mathcal{T}\bigl(\mathbf{let}\; x = V_1\; V_2
    \;\mathbf{in}\; M'\bigr) =
    \{x \mapsto \{\mathcal{T}(V_1),\;
      \mathrm{argument} \mapsto \mathcal{T}(V_2)\},\;
    \mathrm{result} \mapsto \mathcal{T}(M')\}
  \]
  Since $\mathrm{result} \mapsto \mathcal{T}(M')$ is a property
  definition, the subtree at
  $\text{root} \snoc \mathrm{result}$ is $\mathcal{T}(M')$.
  Inside $\mathcal{T}(M')$, each reference to $x$ is
  $x.\mathrm{result}$, which resolves to the $\mathrm{result}$
  path inside the $x$ subtree
  $\{\mathcal{T}(\lambda y.\, M''),\;
  \mathrm{argument} \mapsto \mathcal{T}(V_2)\}$.
  This is the inherited tree $C$ of
  Lemma~\ref{lem:substitution}, so
  $x.\mathrm{result}$ has the same properties as
  $\mathcal{T}(M''[V_2/y])$.
  In the $\lambda$-calculus,
  $M = (\lambda x.\, M')(V_1\; V_2)$, and the head reduct is
  $N = M'[(V_1\; V_2)/x]$.
  In $\mathcal{T}(N)$, each occurrence of $x$ is replaced by
  $(V_1\; V_2)$, whose result, by the tail-call variant above,
  also yields $\mathcal{T}(M''[V_2/y])$.
  Therefore the properties at
  $\text{root} \snoc \mathrm{result} \snoc p$ in
  $\mathcal{T}(M)$ coincide with those at
  $\text{root} \snoc p$ in $\mathcal{T}(N)$ for all $p$.

  \medskip\noindent
  \emph{Value let-binding variant.}
  An analogous result holds for
  $\mathbf{let}\; x = V \;\mathbf{in}\; M'$.
  The translation gives:
  \[
    \mathcal{T}\bigl(\mathbf{let}\; x = V \;\mathbf{in}\; M'\bigr) =
    \{x \mapsto \{\mathrm{result} \mapsto \mathcal{T}(V)\},\;
    \mathrm{result} \mapsto \mathcal{T}(M')\}.
  \]
  Since $\mathrm{result} \mapsto \mathcal{T}(M')$ is a property
  definition, the subtree at
  $\text{root} \snoc \mathrm{result}$ is $\mathcal{T}(M')$.
  Inside $\mathcal{T}(M')$, each occurrence of $x$ is translated
  as the lexical reference $x.\mathrm{result}$, which resolves to
  the subtree $\mathcal{T}(V)$.
  This is exactly the translation of the reduct
  $M'[V/x]$.
  Therefore the properties at
  $\text{root} \snoc \mathrm{result} \snoc p$ in
  $\mathcal{T}(\mathbf{let}\; x = V \;\mathbf{in}\; M')$
  coincide with those at
  $\text{root} \snoc p$ in $\mathcal{T}(M'[V/x])$ for all $p$.
\end{proof}

\begin{theorem}[Convergence preservation]%
  \label{thm:convergence-preservation}
  If $M \to_h N$ (head reduction), then
  $\mathcal{T}(M){\Downarrow}$ if and only if
  $\mathcal{T}(N){\Downarrow}$.
\end{theorem}

\begin{proof}
  A head reduction step in ANF takes one of three forms:
  a tail call $(\lambda x.\, M')\; V \to_h M'[V/x]$,
  an application $\mathbf{let}$-binding
  $\mathbf{let}\; x = V_1\; V_2 \;\mathbf{in}\; M'
  \to_h M'[(V_1\; V_2)/x]$,
  or a value $\mathbf{let}$-binding
  $\mathbf{let}\; x = V \;\mathbf{in}\; M'
  \to_h M'[V/x]$.
  The second and third forms are simply the ANF
  presentation of one $\beta$-step in the underlying
  $\lambda$-term.
  In all three cases, Lemma~\ref{lem:result-step}
  (tail-call, application let-binding, and value
  let-binding variants) shows that
  the properties at $\text{root} \snoc \mathrm{result}^n$ in
  $\mathcal{T}(M)$ coincide with the properties at
  $\text{root} \snoc \mathrm{result}^{n-1}$ in
  $\mathcal{T}(N)$ for all $n \ge 1$, where $N$ is the
  head reduct.
  Hence the abstraction shape appears at depth $n$ in the
  pre-reduct if and only if it appears at depth $n - 1$ in
  the post-reduct.
\end{proof}

\begin{remark}[Iteration levels]
  \label{rem:convergence-lfp}
  Section~\ref{sec:bohm-tree} showed that the B\"ohm tree
  correspondence holds at the first iteration
  $T_P\!\uparrow\!1$ and that $\mathrm{lfp}(T_P)$ is
  strictly more defined in the powerset sense of
  Figure~\ref{fig:calculi-matrix}.
  Here we make the connection precise.
  Tabled evaluation~\cite{tamaki1986-tabled-resolution}
  computes $\mathrm{lfp}(T_P)$ by iterating:
  $T_P\!\uparrow\!0 = \bot$,
  $T_P\!\uparrow\!1 = T_P(\bot)$,
  $T_P\!\uparrow\!2 = T_P(T_P(\bot))$,
  \ldots\ until stabilisation
  (Theorem~\ref{thm:tabling-soundness}).

  \begin{description}
    \item[$T_P\!\uparrow\!1$ corresponds to
      $\beta$-reduction.]
      $T_P\!\uparrow\!1$ evaluates each equation once,
      with re-entrant calls returning~$\bot$
      (the value from $T_P\!\uparrow\!0$).
      The Result-Step Lemma
      (Lemma~\ref{lem:result-step}) establishes a
      structural isomorphism between the equation
      systems of $\mathcal{T}(M)$ and
      $\mathcal{T}(N)$, shifted by one
      $\mathrm{result}$ step.
      This isomorphism is purely structural and holds
      for any number of iterations.
      Theorems~\ref{thm:adequacy}
      and~\ref{thm:first-order-semantics}
      are proved at this level.
  \end{description}

  \paragraph{Observation granularity}%
  \label{rem:set-vs-domain}
  The set-valued semantics at $T_P\!\uparrow\!1$ is
  strictly finer than the domain-theoretic semantics of
  the lazy $\lambda$-calculus.
  In the latter, a non-terminating computation yields a
  single undifferentiated~$\bot$;
  in the former, one can in principle distinguish
  $\varnothing$ (a definite empty set, arrived at in
  finite time) from non-stabilisation of the Kleene
  iteration (no finite approximation suffices).
  The convergence criterion
  (Definition~\ref{def:convergence}) deliberately
  projects away this distinction: it observes only
  whether the $\mathrm{result}$ chain reaches an abstraction
  shape, reducing the observation to a single bit.
  Theorems~\ref{thm:adequacy}
  and~\ref{thm:first-order-semantics}
  are stated at this coarsened observation level.

  \begin{description}
    \item[$\mathrm{lfp}(T_P)$ restores full precision.]
      At $\mathrm{lfp}(T_P)$, the full precision of the
      set-valued semantics is restored: re-entrant calls
      that returned~$\varnothing$ at $T_P\!\uparrow\!1$
      may now converge to non-empty sets, so the semantics
      distinguishes ``definitionally empty'' from
      ``convergent after further iteration,''
      and $M{\Downarrow}_\lambda$ implies
      $\mathcal{T}(M){\Downarrow}_{\mathrm{lfp}}$,
      but not conversely.
  \end{description}
\end{remark}

\subsection{Adequacy}

\begin{definition}[Inheritance-convergence]\label{def:convergence}
  A mixin tree $T$ \emph{converges}, written $T{\Downarrow}$,
  if there exists $n \ge 0$ such that:
  \[
    \{\mathrm{argument},\, \mathrm{result}\}
    \;\subseteq\;
    \properties\!\bigl(\,
      \underbrace{\text{root} \snoc \mathrm{result}
      \snoc \cdots \snoc \mathrm{result}}_{n}
    \,\bigr)
  \]
  in the semantics of
  Section~\ref{sec:mixin-trees}.
\end{definition}

\begin{theorem}[Adequacy]\label{thm:adequacy}
  A closed ANF $\lambda$-term $M$ has a head normal form if and
  only if $\mathcal{T}(M){\Downarrow}$.
\end{theorem}

A closed ANF term at the top level is either an abstraction
(a value) or a computation (an application
  $\mathbf{let}$-binding, a value $\mathbf{let}$-binding,
or a tail call).
An abstraction translates to a record whose root immediately
has the abstraction shape; a computation's root has the form
$\{x, \mathrm{result}\}$ or
$\{\mathrm{tailCall}, \mathrm{result}\}$, and one must follow
the $\mathrm{result}$ chain to find the eventual value.

\begin{proof}[Proof of Theorem~\ref{thm:adequacy}]
  ($\Rightarrow$)
  Suppose $M \to^*_h \lambda x.\, M'$ in $k$ head reduction
  steps.
  We show $\mathcal{T}(M)$ converges at depth $\le k$ by
  induction on $k$.

  \paragraph{Base case ($k = 0$)}
  $M$ is already an abstraction
  $\lambda x.\, M'$.
  Then $\mathcal{T}(M) =
  \{\mathrm{argument} \mapsto \{\},\;
  \mathrm{result} \mapsto \mathcal{T}(M')\}$,
  whose root has
  $\defines = \{\mathrm{argument}, \mathrm{result}\}$.
  These labels are directly in $\defines$ of the root, so
  the recursive evaluation returns them without further
  recursion, and $\mathcal{T}(M)$ converges at depth $n = 0$.

  \paragraph{Inductive step ($k \ge 1$)}
  $M$ is not an abstraction,
  so $M$ is either a tail call, an application
  $\mathbf{let}$-binding, or a value
  $\mathbf{let}$-binding.
  Each form corresponds to one head $\beta$-step in the
  underlying $\lambda$-term.
  Let $N$ be the head reduct, so
  $M \to_h N \to^{k-1}_h \lambda x.\, M'$.
  Note that $N$ is closed: head reduction substitutes a
  closed value into a term whose only free variable is the
  substitution target, preserving closedness.
  By Lemma~\ref{lem:result-step} (tail-call,
    application let-binding, or value let-binding
  variant), the properties at
  $\text{root} \snoc \mathrm{result} \snoc p$
  in $\mathcal{T}(M)$ coincide with those at
  $\text{root} \snoc p$ in $\mathcal{T}(N)$.
  By the induction hypothesis, $\mathcal{T}(N)$
  converges at depth $\le k - 1$, so $\mathcal{T}(M)$
  converges at depth $\le k$.

  \medskip\noindent
  ($\Leftarrow$)
  Suppose $\mathcal{T}(M){\Downarrow}$ at depth $n$.
  We show $M$ has a head normal form by induction on $n$.

  \paragraph{Base case ($n = 0$)}
  The root of $\mathcal{T}(M)$ has
  abstraction shape
  $\{\mathrm{argument}, \mathrm{result}\}$.
  Only the abstraction rule $\mathcal{T}(\lambda x.\, M')$
  produces a root with own property $\mathrm{argument}$.
  (For tail calls and $\mathbf{let}$-bindings, the root has
    $\defines = \{\mathrm{tailCall}, \mathrm{result}\}$
    or $\{x, \mathrm{result}\}$, and $\mathrm{argument}$ cannot
    appear via $\supers$ because
    $\inherits = \varnothing$ at the root, giving
  $\bases(\text{root}) = \varnothing$.)
  Therefore $M$ is an abstraction, hence already in head normal
  form.

  \paragraph{Inductive step ($n \ge 1$)}
  The abstraction shape
  appears at depth $n$ but not at depth $0$.
  By the base case argument, $M$ is not an abstraction, so at the
  head position of this closed top-level ANF term, $M$ is either
  a tail call, an application $\mathbf{let}$-binding, or a value
  $\mathbf{let}$-binding.
  In each case, the ANF grammar requires the function
  position to be a value~$V_1$
  ($V ::= x \mid \lambda x.\, M$).
  Since $M$ is closed, $V_1$ cannot be a free variable,
  so $V_1$ must be a $\lambda$-abstraction;
  the head form is therefore a $\beta$-redex.
  By Lemma~\ref{lem:result-step} (tail-call,
    application let-binding, or value let-binding
  variant), the properties at depth
  $n - 1$ in $\mathcal{T}(N)$ match those at depth $n$
  in $\mathcal{T}(M)$, where $N$ is the head reduct.
  So $\mathcal{T}(N)$ converges at depth $n - 1$.
  Since $M$ is closed and head reduction preserves
  closedness (it substitutes a closed value for the only
  free variable), $N$ is a closed ANF term.
  By the induction hypothesis, $N$ has a head normal
  form, and since $M \to_h N$, so does $M$.
\end{proof}

\subsection{First-Order Semantics of the Lazy \texorpdfstring{$\lambda$}{λ}-Calculus}

Adequacy relates a single term's head normal form to its
inheritance-convergence.
The following theorem lifts this to an equivalence between
terms, establishing that the translation $\mathcal{T}$
gives the lazy $\lambda$-calculus a first-order semantics
that preserves and reflects B\"ohm tree equivalence.

\begin{definition}[$\lambda$-contextual equivalence under $\mathcal{T}$]%
  \label{def:ctx-equiv}
  For closed $\lambda$-terms $M$ and $N$, write
  $\mathcal{T}(M) \approx_\lambda \mathcal{T}(N)$ if for every
  closing $\lambda$-calculus context $C[\cdot]$:
  \[
    \mathcal{T}(C[M]){\Downarrow}
    \;\Longleftrightarrow\;
    \mathcal{T}(C[N]){\Downarrow}
  \]
\end{definition}

\noindent
The subscript $\lambda$ emphasizes that the quantification ranges
over $\lambda$-calculus contexts, not over all inheritance-calculus contexts.
This definition observes only convergence\footnote{
  Convergence provides exactly a single bit of information.
}, but
quantifying over all $\lambda$-contexts makes it a fine-grained
equivalence: different contexts can supply different arguments
and project different results, probing every aspect of the
term's behavior within the $\lambda$-definable fragment.

\begin{theorem}[First-order semantics of the lazy $\lambda$-calculus]%
  \label{thm:first-order-semantics}
  For closed $\lambda$-terms $M$ and $N$:
  \[
    \mathcal{T}(M) \approx_\lambda \mathcal{T}(N)
    \quad\Longleftrightarrow\quad
    \mathrm{BT}(M) = \mathrm{BT}(N)
  \]
\end{theorem}

\begin{proof}[Proof of Theorem~\ref{thm:first-order-semantics}]
  The translation $\mathcal{T}$ is compositional: for every
  $\lambda$-calculus context $C[\cdot]$, the mixin tree
  $\mathcal{T}(C[M])$ is determined by the mixin tree
  $\mathcal{T}(M)$ and the translation of the context.
  This means $\mathcal{T}$ preserves the context structure
  needed for the argument.

  \medskip\noindent
  ($\Leftarrow$)
  Suppose $\mathrm{BT}(M) = \mathrm{BT}(N)$.
  B\"ohm tree equivalence is a congruence~\cite{barendregt1984-lambda-calculus},
  so for every context $C[\cdot]$,
  $\mathrm{BT}(C[M]) = \mathrm{BT}(C[N])$.
  In particular, $C[M]$ has a head normal form iff $C[N]$ does.
  By Theorem~\ref{thm:adequacy},
  $\mathcal{T}(C[M]){\Downarrow}$ iff
  $\mathcal{T}(C[N]){\Downarrow}$.
  Hence $\mathcal{T}(M) \approx_\lambda \mathcal{T}(N)$.

  \medskip\noindent
  ($\Rightarrow$)
  Suppose $\mathrm{BT}(M) \neq \mathrm{BT}(N)$.
  B\"ohm tree equivalence coincides with observational
  equivalence for the lazy
  $\lambda$-calculus~\cite{abramsky1993-full-abstraction-lazy-lambda,abramsky1995-games-full-abstraction-lazy-lambda}:
  there exists a context $C[\cdot]$ such that $C[M]$ converges
  and $C[N]$ diverges, or vice versa.
  By Theorem~\ref{thm:adequacy},
  $\mathcal{T}(C[M]){\Downarrow}$ and
  $\mathcal{T}(C[N]){\Uparrow}$.
  Hence $\mathcal{T}(M) \not\approx_\lambda \mathcal{T}(N)$.
\end{proof}

\noindent
The proof rests on two pillars: Adequacy
(Theorem~\ref{thm:adequacy}), which is our contribution, and
the classical result that B\"ohm tree equivalence equals
observational equivalence for the lazy
$\lambda$-calculus~\cite{abramsky1993-full-abstraction-lazy-lambda,abramsky1995-games-full-abstraction-lazy-lambda}.
Neither pillar requires defining or comparing
absolute paths inside mixin trees.

\section{\bilingual{Expressive Asymmetry: Proofs}{表达力的不对称:证明}}
\label{app:expressiveness}

\begin{definition}[\bilingual{Macro-expressibility, after Felleisen~\cite{felleisen1991-expressive-power}}{宏可表达性,据 Felleisen~\cite{felleisen1991-expressive-power}}]
  \label{def:macro-expressibility}
  \ifInheritanceChinese
  语言 $\mathscr{L}_1$ 的一个构造 $C$ 在语言 $\mathscr{L}_0$ 中是\emph{宏可表达的},若存在一个翻译 $\mathcal{E}$,使得对每一个含有 $C$ 的出现的程序 $P$,翻译 $\mathcal{E}(P)$ 是通过将 $C$ 的每次出现替换为一个仅依赖于 $C$ 的子表达式的 $\mathscr{L}_0$ 表达式而得到的,并保持 $P$ 的所有其他构造不变。
  \else
  A construct $C$ of language $\mathscr{L}_1$ is \emph{macro-expressible}
  in language $\mathscr{L}_0$ if there exists a translation $\mathcal{E}$
  such that for every program $P$ containing occurrences of $C$,
  the translation $\mathcal{E}(P)$ is obtained by replacing each
  occurrence of $C$ with an $\mathscr{L}_0$ expression that depends only
  on $C$'s subexpressions, leaving all other constructs of $P$
  unchanged.
  \fi
\end{definition}

\begin{lemma}[\bilingual{Composition of macro-expressible translations}{宏可表达翻译的复合}]
  \label{lem:macro-composition}
  \ifInheritanceChinese
  若 $\mathcal{E}_1$ 在 $\mathscr{L}_0$ 中宏可表达 $\mathscr{L}_1$ 的构造,且 $\mathcal{E}_2$ 在 $\mathscr{L}_1$ 中宏可表达 $\mathscr{L}_2$ 的构造,则复合 $\mathcal{E}_1 \circ \mathcal{E}_2$ 在 $\mathscr{L}_0$ 中宏可表达 $\mathscr{L}_2$ 的构造。
  \else
  If $\mathcal{E}_1$ macro-expresses the constructs of
  $\mathscr{L}_1$ in $\mathscr{L}_0$ and
  $\mathcal{E}_2$ macro-expresses the constructs of
  $\mathscr{L}_2$ in $\mathscr{L}_1$, then the composition
  $\mathcal{E}_1 \circ \mathcal{E}_2$ macro-expresses the
  constructs of $\mathscr{L}_2$ in $\mathscr{L}_0$.
  \fi
\end{lemma}

\begin{proof}
  \ifInheritanceChinese
  $\mathscr{L}_2$ 的每个构造 $F$ 在 $\mathscr{L}_1$ 上有一个固定的语法抽象 $A_F^{(2)}$。
  出现在 $A_F^{(2)}$ 中的每个 $\mathscr{L}_1$-构造 $G$ 在 $\mathscr{L}_0$ 上有一个固定的语法抽象 $A_G^{(1)}$。
  将 $A_F^{(2)}$ 中的每个 $G$ 替换为 $A_G^{(1)}$,得到 $F$ 在 $\mathscr{L}_0$ 上的一个固定语法抽象,从而满足 E4。
  \else
  Each construct~$F$ of~$\mathscr{L}_2$ has a fixed
  syntactic abstraction~$A_F^{(2)}$
  over~$\mathscr{L}_1$.
  Each $\mathscr{L}_1$-construct~$G$ appearing
  in~$A_F^{(2)}$ has a fixed
  syntactic abstraction~$A_G^{(1)}$
  over~$\mathscr{L}_0$.
  Substituting~$A_G^{(1)}$ for every~$G$
  in~$A_F^{(2)}$ yields a fixed syntactic
  abstraction over~$\mathscr{L}_0$ for~$F$,
  satisfying~E4.
  \fi
\end{proof}

\begin{definition}[\bilingual{Language universe}{语言全集}]
  \label{def:language-universe}
  \ifInheritanceChinese
  设 $\mathscr{U}$ 为一个语言,其构造子是继承演算的构造子、惰性 $\lambda$-演算的构造子(包括 $\lambda$-抽象、应用与变量)以及 $\mathbf{let}$-绑定的并集。
  $\mathscr{U}$ 的语义通过将每个 $\mathscr{U}$-程序翻译为继承演算来给出,分三步:
  \begin{enumerate}
    \item 将每个极大 $\lambda$-演算子表达式(即根节点与所有内部节点仅使用 $\lambda$-演算构造子的每棵子树)转化为 A-范式~\cite{flanagan1993-essence-compiling-continuations}。
    \item 将翻译 $\mathcal{T}$(第~\ref{sec:forward-translation} 节)应用于每个 ANF $\lambda$-演算子表达式。
    \item 按第~\ref{sec:mixin-trees} 节的规则对所得继承演算程序求值。
  \end{enumerate}
  当一个 $\mathscr{U}$-表达式交错使用 $\lambda$-演算与继承演算的构造子(例如,一个 $\mathbf{let}$-绑定,其被绑定的值是一个继承演算表达式)时,翻译从叶到根组合地应用:每个 $\lambda$-演算构造子通过 $\mathcal{T}$ 替换为其继承演算等价物,同时将继承演算子表达式视为不透明的值。
  由于 $\mathcal{T}$ 的每条规则(第~\ref{sec:forward-translation} 节)都是一个固定的语法抽象,仅依赖于被翻译的子表达式,因此这种自底向上的应用是良定义的,并为步骤~(3) 产生一个纯继承演算程序。

  若一个 $\mathscr{U}$-程序 $P$ 的翻译收敛,则 $P$ 收敛,记为 $\mathrm{eval}_{\mathscr{U}}(P)$。
  \else
  Let $\mathscr{U}$ be the language whose constructors are
  the union of the constructors of inheritance-calculus,
  the constructors of the lazy $\lambda$-calculus
  , including $\lambda$-abstraction, application, and variable, and
  $\mathbf{let}$-binding.
  The semantics of $\mathscr{U}$ is given by translating
  every $\mathscr{U}$-program to inheritance-calculus in
  three steps:
  \begin{enumerate}
    \item Convert every maximal $\lambda$-calculus
      subexpression (that is, every subtree whose root
        and all internal nodes use only $\lambda$-calculus
      constructors) to A-normal
      form~\cite{flanagan1993-essence-compiling-continuations}.
    \item Apply the translation $\mathcal{T}$
      (Section~\ref{sec:forward-translation}) to every ANF
      $\lambda$-calculus subexpression.
    \item Evaluate the resulting inheritance-calculus
      program by the rules of
      Section~\ref{sec:mixin-trees}.
  \end{enumerate}
  When a $\mathscr{U}$-expression interleaves
  $\lambda$-calculus and inheritance-calculus constructors
  (e.g.\ a $\mathbf{let}$-binding whose bound value is an
  inheritance-calculus expression), the translation is
  applied compositionally from leaves to root: each
  $\lambda$-calculus construct is replaced by its
  inheritance-calculus equivalent via~$\mathcal{T}$,
  treating inheritance-calculus subexpressions as opaque
  values.
  Since each rule of~$\mathcal{T}$
  (Section~\ref{sec:forward-translation}) is a fixed
  syntactic abstraction depending only on the translated
  subexpressions, this bottom-up application is
  well-defined and yields a pure inheritance-calculus
  program for step~(3).

  A $\mathscr{U}$-program $P$ converges,
  $\mathrm{eval}_{\mathscr{U}}(P)$, iff its translation
  converges.
  \fi
\end{definition}

\begin{remark}\label{rem:universe-properties}
  \ifInheritanceChinese
  惰性 $\lambda$-演算与继承演算是 $\mathscr{U}$ 的保守性限制:各自通过去掉另一方的构造子来得到。
  对于纯继承演算程序,步骤~(1) 与~(2) 是平凡的,因此 $\mathrm{eval}_{\mathscr{U}}$ 与继承演算的语义一致。
  对于纯 $\lambda$-演算程序,步骤~(1)--(3) 的复合由充分性(定理~\ref{thm:adequacy})与惰性求值一致。
  \else
  The lazy $\lambda$-calculus and inheritance-calculus are
  conservative restrictions of $\mathscr{U}$: each is
  obtained by removing the other's constructors.
  For a pure inheritance-calculus program, steps~(1)
  and~(2) are vacuous, so
  $\mathrm{eval}_{\mathscr{U}}$ agrees with the
  inheritance-calculus semantics.
  For a pure $\lambda$-calculus program, the composition
  of steps~(1)--(3) agrees with lazy evaluation by
  Adequacy (Theorem~\ref{thm:adequacy}).
  \fi
\end{remark}

\ifInheritanceChinese
三步翻译是有效的(ANF 转换与 $\mathcal{T}$ 是语法变换),而继承演算的递归求值(第~\ref{sec:mixin-trees} 节)是一个半判定过程,因此 $\mathrm{eval}_{\mathscr{U}}$ 是递归可枚举的,满足 Felleisen 框架~\cite{felleisen1991-expressive-power} 的要求。
\else
The three-step translation is effective (ANF conversion
and~$\mathcal{T}$ are syntactic transformations) and the
recursive evaluation of inheritance-calculus
(Section~\ref{sec:mixin-trees}) is a semi-decision
procedure, so $\mathrm{eval}_{\mathscr{U}}$ is
recursively enumerable, as required by
Felleisen's framework~\cite{felleisen1991-expressive-power}.
\fi

\begin{lemma}[\bilingual{$\mathbf{let}$-binding is macro-expressible
  in the lazy $\lambda$-calculus}{$\mathbf{let}$-绑定在惰性 $\lambda$-演算中是宏可表达的}]
  \label{lem:let-macro}
  \ifInheritanceChinese
  翻译 $\mathbf{let}\; x = e_1 \;\mathbf{in}\; e_2 \;\mapsto\; (\lambda x.\, e_2)\; e_1$ 是 $\mathbf{let}$ 在惰性 $\lambda$-演算中的一个宏可表达的消去。
  \else
  The translation
  $\mathbf{let}\; x = e_1 \;\mathbf{in}\; e_2
  \;\mapsto\; (\lambda x.\, e_2)\; e_1$
  is a macro-expressible elimination of $\mathbf{let}$
  in the lazy $\lambda$-calculus.
  \fi
\end{lemma}

\begin{proof}
  \ifInheritanceChinese
  该翻译将每个 $\mathbf{let}$-绑定替换为一个仅依赖于两个子表达式 $e_1$ 与 $e_2$ 的 $\lambda$-表达式。
  这是 Felleisen~\cite{felleisen1991-expressive-power} 的命题 3.7。
  \else
  The translation replaces each $\mathbf{let}$-binding
  with a $\lambda$-expression that depends only on the
  two subexpressions $e_1$ and~$e_2$.
  This is Proposition~3.7 of
  Felleisen~\cite{felleisen1991-expressive-power}.
  \fi
\end{proof}

\begin{lemma}[\bilingual{Equal expressiveness of the lazy
  $\lambda$-calculus and its ANF fragment}{惰性 $\lambda$-演算与其 ANF 片段的等同表达力}]
  \label{lem:anf-equiv}
  \ifInheritanceChinese
  在 $\mathscr{U}$ 中,惰性 $\lambda$-演算与 ANF 惰性 $\lambda$-演算(第~\ref{sec:forward-translation} 节)具有相同的宏可表达力。
  \else
  In $\mathscr{U}$, the lazy $\lambda$-calculus and the
  ANF lazy $\lambda$-calculus
  (Section~\ref{sec:forward-translation}) are equally
  macro-expressive.
  \fi
\end{lemma}

\begin{proof}
  \ifInheritanceChinese
  每个 ANF 项都是一个 $\lambda$-项,因此 ANF 片段嵌入完整 $\lambda$-演算的映射是恒等变换,平凡地宏可表达。
  对于反方向,ANF 变换~\cite{flanagan1993-essence-compiling-continuations} 由两步构成,每步都是宏可表达的:
  \begin{enumerate}
    \item \emph{\bilingual{Application flattening}{应用展开}。}
      将每个嵌套应用 $(e_1\; e_2)\; e_3$ 替换为 $\mathbf{let}\; x = e_1\; e_2 \;\mathbf{in}\; x\; e_3$。
      其语法抽象为 $A(\cdot_1, \cdot_2, \cdot_3) = \mathbf{let}\; x = \cdot_1\; \cdot_2 \;\mathbf{in}\; x\; \cdot_3$,仅依赖三个子表达式。
    \item \emph{\bilingual{Value lifting}{值提升}。}
      将应用位置上每个自由值 $V$ 替换为 $\mathbf{let}\; x = V \;\mathbf{in}\; x$。
      其语法抽象为 $A(\cdot_1) = \mathbf{let}\; x = \cdot_1 \;\mathbf{in}\; x$。
  \end{enumerate}
  由引理~\ref{lem:let-macro},每个所得 $\mathbf{let}$-绑定本身在 $\lambda$-演算中是宏可表达的。
  由引理~\ref{lem:macro-composition},复合这些宏可表达翻译,得到一个从完整 $\lambda$-演算到其 ANF 片段的宏可表达翻译。
  \else
  Every ANF term is a $\lambda$-term, so the embedding
  of the ANF fragment into the full $\lambda$-calculus is
  the identity, which is trivially macro-expressible.
  For the reverse direction, the ANF
  transformation~\cite{flanagan1993-essence-compiling-continuations} consists of
  two steps, each macro-expressible:
  \begin{enumerate}
    \item \emph{Application flattening.}
      Replace each nested application
      $(e_1\; e_2)\; e_3$ with
      $\mathbf{let}\; x = e_1\; e_2 \;\mathbf{in}\;
      x\; e_3$.
      The syntactic abstraction is
      $A(\cdot_1, \cdot_2, \cdot_3)
      = \mathbf{let}\; x = \cdot_1\; \cdot_2
      \;\mathbf{in}\; x\; \cdot_3$,
      depending only on the three subexpressions.
    \item \emph{Value lifting.}
      Replace each free value~$V$ used in an application
      position with
      $\mathbf{let}\; x = V \;\mathbf{in}\; x$.
      The syntactic abstraction is
      $A(\cdot_1) = \mathbf{let}\; x = \cdot_1
      \;\mathbf{in}\; x$.
  \end{enumerate}
  By Lemma~\ref{lem:let-macro}, each resulting
  $\mathbf{let}$-binding is itself
  macro-expressible in the $\lambda$-calculus.
  By Lemma~\ref{lem:macro-composition}, composing
  these macro-expressible translations
  yields a macro-expressible translation from the full
  $\lambda$-calculus to its ANF fragment.
  \fi
\end{proof}

\begin{theorem}[\bilingual{Forward macro-expressibility}{正向宏可表达性}]
  \label{thm:forward-macro}
  \ifInheritanceChinese
  翻译 $\mathcal{T}$(第~\ref{sec:forward-translation} 节)是 $\lambda$-演算嵌入继承演算的一个宏可表达的嵌入。
  \else
  The translation $\mathcal{T}$
  (Section~\ref{sec:forward-translation}) is a macro-expressible
  embedding of the $\lambda$-calculus into inheritance-calculus.
  \fi
\end{theorem}

\begin{proof}
  \ifInheritanceChinese
  我们验证 Felleisen~\cite[Definition~3.11]{felleisen1991-expressive-power} 的条件 E4:对于每个元数为 $a$ 的 $\lambda$-演算构造子 $F$,必须存在一个固定的 $a$ 元语法抽象 $A_F(\cdot_1,\ldots,\cdot_a)$ 使得 $\mathcal{T}(F(e_1,\ldots,e_a)) = A_F(\mathcal{T}(e_1),\ldots,\mathcal{T}(e_a))$。
  我们逐条检验 $\mathcal{T}$ 的规则(第~\ref{sec:forward-translation} 节):
  \begin{enumerate}
    \item \emph{\bilingual{$\lambda$-bound variable}{$\lambda$-绑定变量}}($x$,de~Bruijn 索引为 $n'$)。
      这是一个零元构造子,因此没有子表达式。
      语法抽象是常量 $A = \dbi{n'}.\mathrm{argument}$。
    \item \emph{\bilingual{Let-bound variable}{Let-绑定变量}}($x$)。
      零元。$A = x.\mathrm{result}$。
    \item \emph{\bilingual{$\lambda$-abstraction}{$\lambda$-抽象}},记为 $\lambda x.\, M$。
      关于体 $M$ 是一元的。
      $A(\cdot_1) = \{\mathrm{argument} \mapsto \{\},\; \mathrm{result} \mapsto \cdot_1\}$。
    \item \emph{\bilingual{Application let-binding}{应用 Let-绑定}}($\mathbf{let}\; x = V_1\; V_2 \;\mathbf{in}\; M$)。
      关于 $V_1$、$V_2$ 与 $M$ 是三元的。
      $A(\cdot_1, \cdot_2, \cdot_3) = \{x \mapsto \{\cdot_1,\; \mathrm{argument} \mapsto \cdot_2\},\; \mathrm{result} \mapsto \cdot_3\}$。
    \item \emph{\bilingual{Value let-binding}{值 Let-绑定}}($\mathbf{let}\; x = V \;\mathbf{in}\; M$)。
      关于 $V$ 与 $M$ 是二元的。
      $A(\cdot_1, \cdot_2) = \{x \mapsto \{\mathrm{result} \mapsto \cdot_1\},\; \mathrm{result} \mapsto \cdot_2\}$。
    \item \emph{\bilingual{Tail call}{尾调用}}($V_1\; V_2$)。
      二元。\newline
      $A(\cdot_1, \cdot_2) = \{\mathrm{tailCall} \mapsto \{\cdot_1,\; \mathrm{argument} \mapsto \cdot_2\},\; \mathrm{result} \mapsto \{\mathrm{tailCall}.\mathrm{result}\}\}$。
  \end{enumerate}
  在每种情况下,语法抽象 $A_F$ 都是一个固定的继承演算上下文,其空位由子表达式的递归翻译填充,从而满足 E4。
  \else
  We verify condition~E4 of
  Felleisen~\cite[Definition~3.11]{felleisen1991-expressive-power}:
  for each $\lambda$-calculus construct~$F$ of arity~$a$,
  there must exist a fixed $a$-ary syntactic abstraction
  $A_F(\cdot_1,\ldots,\cdot_a)$ over
  inheritance-calculus such that
  $\mathcal{T}(F(e_1,\ldots,e_a)) =
  A_F(\mathcal{T}(e_1),\ldots,\mathcal{T}(e_a))$.
  We check each rule of $\mathcal{T}$
  (Section~\ref{sec:forward-translation}):
  \begin{enumerate}
    \item \emph{$\lambda$-bound variable}
      ($x$ with de~Bruijn index $n'$).
      This is a nullary construct and therefore has no subexpressions.
      The syntactic abstraction is the constant
      $A = \dbi{n'}.\mathrm{argument}$.
    \item \emph{Let-bound variable} ($x$).
      Nullary.
      $A = x.\mathrm{result}$.
    \item \emph{$\lambda$-abstraction}, denoted $\lambda x.\, M$.
      Unary in the body~$M$.
      $A(\cdot_1) = \{\mathrm{argument} \mapsto \{\},\;
      \mathrm{result} \mapsto \cdot_1\}$.
    \item \emph{Application let-binding}
      ($\mathbf{let}\; x = V_1\; V_2
      \;\mathbf{in}\; M$).
      Ternary in~$V_1$, $V_2$, and~$M$.
      $A(\cdot_1, \cdot_2, \cdot_3) =
      \{x \mapsto \{\cdot_1,\;
        \mathrm{argument} \mapsto \cdot_2\},\;
      \mathrm{result} \mapsto \cdot_3\}$.
    \item \emph{Value let-binding}
      ($\mathbf{let}\; x = V \;\mathbf{in}\; M$).
      Binary in~$V$ and~$M$.
      $A(\cdot_1, \cdot_2) =
      \{x \mapsto \{\mathrm{result} \mapsto \cdot_1\},\;
      \mathrm{result} \mapsto \cdot_2\}$.
    \item \emph{Tail call} ($V_1\; V_2$).
      Binary.\newline
      $A(\cdot_1, \cdot_2) =
      \{\mathrm{tailCall} \mapsto
        \{\cdot_1,\;
        \mathrm{argument} \mapsto \cdot_2\},\;
        \mathrm{result} \mapsto
      \{\mathrm{tailCall}.\mathrm{result}\}\}$.
  \end{enumerate}
  In each case the syntactic abstraction~$A_F$ is a fixed
  inheritance-calculus context whose holes are filled
  by the recursive translations of the subexpressions,
  satisfying~E4.
  \fi
\end{proof}

\begin{theorem}[\bilingual{Nonexpressibility of inheritance}{继承的不可宏表达性}]
  \label{thm:cartesian-nonexpressibility}
  \ifInheritanceChinese
  惰性 $\lambda$-演算不能宏可表达 $\mathscr{U}$(定义~\ref{def:language-universe})的继承构造子。
  \else
  The lazy $\lambda$-calculus cannot macro-express the
  inheritance constructors of $\mathscr{U}$
  (Definition~\ref{def:language-universe}).
  \fi
\end{theorem}

\begin{proof}
  \ifInheritanceChinese
  我们对 $L_0 = \text{惰性 $\lambda$-演算}$,$L_1 = \mathscr{U}$ 应用 Felleisen~\cite{felleisen1991-expressive-power} 的定理 3.14(i)。
  由定义~\ref{def:language-universe},$\mathscr{U} = L_0 + \{$mixin-树构造子$\}$ 是 $L_0$ 的一个保守性扩展(Felleisen~\cite{felleisen1991-expressive-power} 的定义 3.2)。
  我们验证四个条件:
  \begin{enumerate}
    \item[(i)] \emph{$L_0 \subseteq L_1$:} $\mathscr{U}$ 的语法包含所有 $L_0$-短语。
    \item[(ii)] \emph{不引入新的 $L_0$-短语:} mixin-树构造子不在 $L_0$ 中引入新短语。
    \item[(iii)] \emph{构造子不相交:} mixin-树构造子与 $L_0$-构造子不相交。
    \item[(iv)] \emph{语义保持:} 对纯 $\lambda$-项,定义~\ref{def:language-universe} 的步骤~(1) 与~(2) 是平凡的,步骤~(3) 由充分性(定理~\ref{thm:adequacy})给出惰性求值,因此 $\mathrm{eval}_{\mathscr{U}}$ 在 $L_0$-短语上与 $\mathrm{eval}_{L_0}$ 一致。
  \end{enumerate}
  该定理要求我们给出两个 $L_0$-项 $M$ 与 $N$,满足 $M \cong_0 N$ 但 $M \not\cong_1 N$,其中 $\cong_0$ 是 $L_0$ 中的操作等价,$\cong_1$ 是 $\mathscr{U}$ 中的操作等价。

  \paragraph{\bilingual{Operational equivalence in $L_0$}{$L_0$ 中的操作等价}}
  由 Böhm 树等价~\cite{abramsky1993-full-abstraction-lazy-lambda,abramsky1995-games-full-abstraction-lazy-lambda}(定理~\ref{thm:first-order-semantics}),对于封闭 $\lambda$-项 $M$ 与 $N$:$M \cong_0 N$ 当且仅当 $\mathrm{BT}(M) = \mathrm{BT}(N)$。

  \paragraph{\bilingual{Separating pair}{分离对}}
  定义 Church 布尔值与 Church 布尔相等:$\mathrm{true} = \lambda t.\,\lambda f.\, t$,$\mathrm{false} = \lambda t.\,\lambda f.\, f$,以及 $\mathrm{eq} = \lambda a.\,\lambda b.\, (a\;b)\;(b\;\mathrm{false}\;\mathrm{true})$。
  令
  \[
    M = \mathrm{eq}\;\mathrm{false}\;\mathrm{false}
    \qquad
    N = \mathrm{eq}\;\mathrm{true}\;\mathrm{true}.
  \]
  两者在惰性 $\lambda$-演算中均归约为 Church~$\mathrm{true}$,因此 $\mathrm{BT}(M) = \mathrm{BT}(N)$,从而 $M \cong_0 N$。

  \paragraph{\bilingual{Distinguishing $\mathscr{U}$-context}{区分 $\mathscr{U}$-上下文}}
  我们现在构造一个区分 $M$ 与 $N$ 的 $\mathscr{U}$-上下文。
  $M$ 与 $N$ 是语法封闭的 $\lambda$-项;由定义~\ref{def:language-universe},它们在 $\mathscr{U}$ 中的语义通过 ANF 转换后接 $\mathcal{T}$ 将其翻译为继承演算来给出。
  在扩展的 ANF 约定(第~\ref{sec:forward-translation} 节)下,应用位置上的每个自由值首先被绑定到一个名字:
  \else
  We apply Theorem~3.14(i) of
  Felleisen~\cite{felleisen1991-expressive-power} to the pair
  $L_0 = \text{lazy $\lambda$-calculus}$,
  $L_1 = \mathscr{U}$.
  By Definition~\ref{def:language-universe},
  $\mathscr{U} = L_0 + \{$mixin-tree constructors$\}$
  is a conservative extension of~$L_0$
  (Definition~3.2 of
  Felleisen~\cite{felleisen1991-expressive-power}).
  We verify the four conditions:
  \begin{enumerate}
    \item[(i)] \emph{$L_0 \subseteq L_1$:}
      The syntax of~$\mathscr{U}$ includes all
      $L_0$-phrases.
    \item[(ii)] \emph{No new $L_0$-phrases:}
      The mixin-tree constructors do not introduce new
      phrases in~$L_0$.
    \item[(iii)] \emph{Disjoint constructors:}
      The mixin-tree constructors are disjoint from the
      $L_0$-constructors.
    \item[(iv)] \emph{Semantics preserved:}
      For pure $\lambda$-terms, steps~(1) and~(2) of
      Definition~\ref{def:language-universe} are vacuous
      and step~(3) yields lazy evaluation by Adequacy
      (Theorem~\ref{thm:adequacy}), so
      $\mathrm{eval}_{\mathscr{U}}$ agrees with
      $\mathrm{eval}_{L_0}$ on $L_0$-phrases.
  \end{enumerate}
  The theorem requires us to exhibit two $L_0$-terms
  $M$ and $N$ with $M \cong_0 N$ but
  $M \not\cong_1 N$, where $\cong_0$ is operational
  equivalence in $L_0$ and $\cong_1$ is operational
  equivalence in $\mathscr{U}$.

  \paragraph{Operational equivalence in $L_0$}
  By B\"ohm tree
  equivalence~\cite{abramsky1993-full-abstraction-lazy-lambda,abramsky1995-games-full-abstraction-lazy-lambda}
  (Theorem~\ref{thm:first-order-semantics}), for closed
  $\lambda$-terms $M$ and $N$:
  $M \cong_0 N$ iff $\mathrm{BT}(M) = \mathrm{BT}(N)$.

  \paragraph{Separating pair}
  Define Church booleans and Church boolean equality:
  $\mathrm{true} = \lambda t.\,\lambda f.\, t$,
  $\mathrm{false} = \lambda t.\,\lambda f.\, f$, and
  $\mathrm{eq} = \lambda a.\,\lambda b.\,
  (a\;b)\;(b\;\mathrm{false}\;\mathrm{true})$.
  Let
  \[
    M = \mathrm{eq}\;\mathrm{false}\;\mathrm{false}
    \qquad
    N = \mathrm{eq}\;\mathrm{true}\;\mathrm{true}.
  \]
  Both reduce to Church~$\mathrm{true}$ in the
  lazy $\lambda$-calculus, so
  $\mathrm{BT}(M) = \mathrm{BT}(N)$ and hence
  $M \cong_0 N$.

  \paragraph{Distinguishing $\mathscr{U}$-context}
  We now construct a $\mathscr{U}$-context that
  separates $M$ and $N$.
  $M$ and $N$ are syntactically closed $\lambda$-terms;
  by Definition~\ref{def:language-universe}, their
  semantics in $\mathscr{U}$ is given by translating
  them to inheritance-calculus via ANF conversion followed
  by~$\mathcal{T}$.
  Under the extended ANF convention
  (Section~\ref{sec:forward-translation}), every free
  value used in an application is first bound to a name:
  \fi
  \begin{align*}
    \mathrm{ANF}(M) ={} &
    \mathbf{let}\;\mathrm{eq} =
    \lambda a.\,\lambda b.\,
    \mathbf{let}\; x_1 = a\;b \;\mathbf{in}\; \\
    &\qquad
    \mathbf{let}\; f_1 = \mathrm{false}
    \;\mathbf{in}\;
    \mathbf{let}\; x_2 = b\;f_1 \;\mathbf{in}\; \\
    &\qquad
    \mathbf{let}\; t_1 = \mathrm{true}
    \;\mathbf{in}\;
    \mathbf{let}\; x_3 = x_2\;t_1
    \;\mathbf{in}\;
    x_1\;x_3 \\
    & \mathbf{in}\;\;
    \mathbf{let}\; f_1 = \mathrm{false}
    \;\mathbf{in}\;
    \mathbf{let}\;\mathrm{eqFalse} = \mathrm{eq}\;f_1 \\
    &\quad
    \;\mathbf{in}\;
    \mathbf{let}\; f_2 = \mathrm{false}
    \;\mathbf{in}\;
    \mathrm{eqFalse}\;f_2
  \end{align*}
  \ifInheritanceChinese
  应用 $\mathcal{T}$:
  \else
  Applying~$\mathcal{T}$:
  \fi
  \begin{align*}
    \mathcal{T}(\mathrm{ANF}(M)) \;=\; \{&
      \mathrm{eq} \mapsto \{\mathrm{result} \mapsto
      \mathcal{T}(\lambda a.\,\lambda b.\,\cdots)\}, \\
      & f_1 \mapsto \{\mathrm{result} \mapsto
      \mathcal{T}(\mathrm{false})\}, \\
      & \mathrm{eqFalse} \mapsto
      \{\mathrm{eq}.\mathrm{result},\;
        \mathrm{argument} \mapsto
      \{f_1.\mathrm{result}\}\}, \\
      & f_2 \mapsto \{\mathrm{result} \mapsto
      \mathcal{T}(\mathrm{false})\}, \\
      & \mathrm{tailCall} \mapsto
      \{\mathrm{eqFalse}.\mathrm{result},\;
        \mathrm{argument} \mapsto
      \{f_2.\mathrm{result}\}\}, \\
      & \mathrm{result} \mapsto
      \{\mathrm{tailCall}.\mathrm{result}\}
    \;\}
  \end{align*}
  \ifInheritanceChinese
  对 $N$ 类似,以 $\mathrm{true}$ 替换 $\mathrm{false}$。
  注意 $\mathrm{eq}$ 是 $\mathcal{T}(\mathrm{ANF}(M))$ 的一个具名子项,通过值 let-绑定绑定。
  在 $\mathrm{eq}$ 的翻译内部,$\lambda b$ 作用域包含 $x_1$、$f_1$、$x_2$、$t_1$、$x_3$ 的 let-绑定以及一个尾调用。
  特别地,$\dbi{1}.\mathrm{argument}$ 到达外层 $\lambda a$ 作用域,而 $\dbi{0}.\mathrm{argument}$ 到达 $\lambda b$ 自身。\footnote{
    这些索引使用压缩约定,仅计数 $\lambda$-作用域。
    如翻译脚注(第~\ref{sec:forward-translation} 节)所述,每个 let-绑定与尾调用在继承演算中引入一个额外的作用域层级;形式索引相应地更大,但目标作用域相同。
  }

  $\mathscr{U}$-上下文使用继承在 $\mathrm{eq}$ 的 $\lambda b$ 作用域内添加一个 $\mathrm{repr}$ 记录,该记录投影第一个操作数。
  定义上下文 $C[\cdot]$ 为:
  \else
  and symmetrically for $N$ with $\mathrm{true}$ in place
  of $\mathrm{false}$.
  Note that $\mathrm{eq}$ is a named child
  of $\mathcal{T}(\mathrm{ANF}(M))$, bound via a value
  let-binding.
  Inside the translation of $\mathrm{eq}$, the inner
  $\lambda b$ scope contains the let-bindings
  for $x_1$, $f_1$, $x_2$, $t_1$, $x_3$, and a tail call.
  In particular, $\dbi{1}.\mathrm{argument}$ reaches
  the outer $\lambda a$ scope and
  $\dbi{0}.\mathrm{argument}$ reaches $\lambda b$
  itself.\footnote{
    These indices use the condensed
    convention, counting only $\lambda$-scopes.
    As noted in the translation footnote
    (Section~\ref{sec:forward-translation}),
    each let-binding and tail call introduces an
    additional scope level in inheritance-calculus;
    the formal index is correspondingly larger,
    but the target scope is the same.
  }

  The $\mathscr{U}$-context uses inheritance to add
  a $\mathrm{repr}$ record inside the $\lambda b$ scope
  of $\mathrm{eq}$ that projects the first operand.
  Define the context $C[\cdot]$ as:
  \fi
  \begin{align*}
    C[\cdot] ={} & \mathbf{let}\;
    \mathrm{extended} = \{[\cdot],\;
      \mathrm{eq} \mapsto \{
        \mathrm{result} \mapsto \{ \\
          &\quad
          \mathrm{result} \mapsto \{
            \mathrm{repr} \mapsto \{
              \mathrm{firstOperand}
              \mapsto \{\dbi{2}.\mathrm{argument}\}
    \}\}\}\}\} \\
    & \mathbf{in}\;
    \mathrm{extended}.\mathrm{tailCall}.\mathrm{repr}.\mathrm{firstOperand}
  \end{align*}
  \ifInheritanceChinese
  这里 $[\cdot]$ 是空位。
  注意 mixin 中 $\mathrm{result}$ 的两层嵌套:外层 $\mathrm{result}$ 进入 $\mathrm{eq}$ 的值 let-绑定包装器,而内层 $\mathrm{result}$ 进入 $\lambda a$ 作用域。
  第二个 $\mathrm{result}$ 则包含 $\lambda b$ 作用域,因此 $\mathrm{repr}$ 被放置在 $\lambda b$ 作用域层级。
  mixin $\{\mathrm{eq} \mapsto \{\mathrm{result} \mapsto \{\mathrm{result} \mapsto \{\mathrm{repr} \mapsto \cdots\}\}\}\}$ 是一个继承演算表达式:它仅使用 $\mathscr{U}$ 的 mixin-树构造子。
  外层 $\mathbf{let}$-绑定与投影是 $\lambda$-演算构造(引理~\ref{lem:let-macro})。
  因此 $C[\cdot]$ 是一个合法的 $\mathscr{U}$-上下文。

  $\mathrm{repr}$ 中的 de~Bruijn 索引 $\dbi{2}$ 计数作用域层级(在压缩约定下):$\dbi{0}$ 是 $\mathrm{repr}$ 自身,$\dbi{1}$ 是 $\lambda b$ 作用域,而 $\dbi{2}$ 是 $\lambda a$ 作用域。
  因此 $\mathrm{firstOperand}$ 投影 $\lambda a$ 的参数。

  由于 $\mathrm{eqFalse}$(相应地,$\mathrm{eqTrue}$)继承 $\mathrm{eq}.\mathrm{result}$,而 $\mathrm{tailCall}$ 继承 $\mathrm{eqFalse}.\mathrm{result}$($\lambda b$ 作用域),mixin 在 $\lambda b$ 作用域层级放置的 $\mathrm{repr}$ 记录通过继承链传播到 $\mathrm{tailCall}$。
  在 $C[M]$ 中:$\mathrm{eqFalse}$ 绑定向 $\lambda a$ 作用域提供 $\mathrm{argument} \mapsto \{f_1.\mathrm{result}\}$(即 $\mathrm{false}$),因此 $\mathrm{firstOperand}$ 求值为 $\mathrm{false} = \lambda t.\,\lambda f.\, f$。
  在 $C[N]$ 中:$\mathrm{eqTrue}$ 绑定提供 $\mathrm{argument} \mapsto \mathrm{true}$,因此 $\mathrm{firstOperand}$ 求值为 $\mathrm{true} = \lambda t.\,\lambda f.\, t$。

  定义 $D[\cdot] = (\lambda x.\, x\;\Omega\;I)\; C[\cdot]$。
  则 $\mathrm{eval}_{\mathscr{U}}$ 对 $D[M]$ 成立(因为 $\mathrm{false}\;\Omega\;I \to I$,收敛),但对 $D[N]$ 不成立(因为 $\mathrm{true}\;\Omega\;I \to \Omega$,发散)。
  因此 $M \not\cong_1 N$。

  \paragraph{\bilingual{Conclusion}{结论}}
  我们有 $M \cong_0 N$ 但 $M \not\cong_1 N$,故 ${\cong_0} \neq {(\cong_1|_{L_0})}$。
  由 Felleisen~\cite{felleisen1991-expressive-power} 的定理 3.14(i),惰性 $\lambda$-演算不能宏可表达 $\mathscr{U}$ 的继承构造子。
  \else
  Here $[\cdot]$ is the hole.
  Note the two levels of $\mathrm{result}$ nesting in the
  mixin: the outer $\mathrm{result}$ enters the value
  let-binding wrapper for $\mathrm{eq}$, and the inner
  $\mathrm{result}$ enters the $\lambda a$ scope.
  The second $\mathrm{result}$ then contains the
  $\lambda b$ scope, so $\mathrm{repr}$ is placed at
  the $\lambda b$ scope level.
  The mixin $\{\mathrm{eq} \mapsto \{\mathrm{result}
      \mapsto \{\mathrm{result} \mapsto
  \{\mathrm{repr} \mapsto \cdots\}\}\}\}$ is an
  inheritance-calculus expression: it uses only the
  mixin-tree constructors of $\mathscr{U}$.
  The outer $\mathbf{let}$-binding and projection are
  $\lambda$-calculus constructs (Lemma~\ref{lem:let-macro}).
  Thus $C[\cdot]$ is a valid $\mathscr{U}$-context.

  The de~Bruijn index $\dbi{2}$ in
  $\mathrm{repr}$ counts scope levels
  (in the condensed convention):
  $\dbi{0}$ is $\mathrm{repr}$ itself,
  $\dbi{1}$ is the $\lambda b$ scope, and
  $\dbi{2}$ is the $\lambda a$ scope.
  So $\mathrm{firstOperand}$ projects
  $\lambda a$'s argument.

  Since $\mathrm{eqFalse}$ (respectively $\mathrm{eqTrue}$)
  inherits $\mathrm{eq}.\mathrm{result}$,
  and $\mathrm{tailCall}$ inherits
  $\mathrm{eqFalse}.\mathrm{result}$ (the $\lambda b$
  scope), the $\mathrm{repr}$ record placed by the mixin
  at the $\lambda b$ scope level propagates through
  the inheritance chain to $\mathrm{tailCall}$.
  In $C[M]$: the $\mathrm{eqFalse}$ binding supplies
  $\mathrm{argument} \mapsto \{f_1.\mathrm{result}\}$
  (that is, $\mathrm{false}$) to the $\lambda a$ scope, so
  $\mathrm{firstOperand}$ evaluates to
  $\mathrm{false} = \lambda t.\,\lambda f.\, f$.
  In $C[N]$: the $\mathrm{eqTrue}$ binding supplies
  $\mathrm{argument} \mapsto \mathrm{true}$, so
  $\mathrm{firstOperand}$ evaluates to
  $\mathrm{true} = \lambda t.\,\lambda f.\, t$.

  Define $D[\cdot] = (\lambda x.\, x\;\Omega\;I)\;
  C[\cdot]$.
  Then $\mathrm{eval}_{\mathscr{U}}$ holds for $D[M]$
  (since $\mathrm{false}\;\Omega\;I \to I$, which
  converges)
  but not for $D[N]$
  (since $\mathrm{true}\;\Omega\;I \to \Omega$, which
  diverges).
  Therefore $M \not\cong_1 N$.

  \paragraph{Conclusion}
  We have $M \cong_0 N$ but $M \not\cong_1 N$, so
  ${\cong_0} \neq {(\cong_1|_{L_0})}$.
  By Theorem~3.14(i) of
  Felleisen~\cite{felleisen1991-expressive-power},
  the lazy $\lambda$-calculus cannot macro-express the
  inheritance constructors of~$\mathscr{U}$.
  \fi
\end{proof}

\begin{theorem}[\bilingual{Expressive asymmetry}{表达力的不对称}]
  \label{thm:expressive-asymmetry}
  \ifInheritanceChinese
  在语言全集 $\mathscr{U}$(定义~\ref{def:language-universe})中,继承演算比惰性 $\lambda$-演算具有严格更强的表达力:它能宏可表达 $\lambda$-演算能宏可表达的每个 $\mathscr{U}$-构造,但反之不然(Felleisen~\cite{felleisen1991-expressive-power} 的定义 3.17)。
  \else
  In the language universe $\mathscr{U}$
  (Definition~\ref{def:language-universe}),
  inheritance-calculus is strictly more expressive than
  the lazy $\lambda$-calculus: it can macro-express every
  $\mathscr{U}$-construct that the $\lambda$-calculus
  can, but not conversely
  (Definition~3.17 of
  Felleisen~\cite{felleisen1991-expressive-power}).
  \fi
\end{theorem}

\begin{proof}
  \ifInheritanceChinese
  我们须证明:(a)~继承演算能宏可表达惰性 $\lambda$-演算能宏可表达的每个 $\mathscr{U}$-构造,以及 (b)~反之不成立。

  \paragraph{\bilingual{Part (b): the lazy $\lambda$-calculus $\not\geq$
  inheritance-calculus}{第 (b) 部分:惰性 $\lambda$-演算 $\not\geq$ 继承演算}}
  $\lambda$-演算包含其自身的构造子($\lambda$、应用、变量),并能宏可表达 $\mathbf{let}$-绑定(引理~\ref{lem:let-macro}),但不能宏可表达 $\mathscr{U}$ 的继承构造子(定理~\ref{thm:cartesian-nonexpressibility})。
  继承演算包含那些构造子。
  因此,惰性 $\lambda$-演算对于 $\mathscr{U}$ 而言表达力不及继承演算。

  \paragraph{\bilingual{Part (a): inheritance-calculus $\geq$ the lazy
  $\lambda$-calculus}{第 (a) 部分:继承演算 $\geq$ 惰性 $\lambda$-演算}}
  我们证明继承演算能宏可表达 $\lambda$-演算所包含的或能宏可表达的每个 $\mathscr{U}$-构造。
  $\lambda$-演算包含 $\lambda$-抽象、应用与变量;它能宏可表达 $\mathbf{let}$-绑定(引理~\ref{lem:let-macro})。
  我们须证明继承演算能宏可表达这四个构造。
  翻译分两个阶段进行:
  \begin{enumerate}
    \item \emph{\bilingual{ANF conversion}{ANF 转换}。}
      A-范式变换~\cite{flanagan1993-essence-compiling-continuations} 将嵌套应用替换为 $\mathbf{let}$-绑定;由引理~\ref{lem:anf-equiv},它在 $\lambda$-演算中是宏可表达的。
      在反方向,$\mathbf{let}$-绑定由 Felleisen~\cite{felleisen1991-expressive-power} 的命题 3.7(重述为引理~\ref{lem:let-macro})在 $\lambda$-演算中是宏可表达的。
      因此,ANF 转换是 $\lambda$-演算与其 ANF 片段之间的一个宏可表达的等价。
    \item \emph{\bilingual{The translation $\mathcal{T}$}{翻译 $\mathcal{T}$}。}
      每个 ANF $\lambda$-演算构造子映射到继承演算上的一个固定语法抽象(定理~\ref{thm:forward-macro}),满足 Felleisen~\cite[Definition~3.11]{felleisen1991-expressive-power} 的条件 E4。
  \end{enumerate}
  由引理~\ref{lem:macro-composition},复合两个宏可表达翻译,得到一个从 $\lambda$-演算(含 $\mathbf{let}$)到继承演算的宏可表达翻译。
  由于继承演算平凡地包含其自身的构造子,它能宏可表达 $\lambda$-演算所包含的或能宏可表达的每个 $\mathscr{U}$-构造。
  \else
  We must show:
  (a)~inheritance-calculus can macro-express every
  $\mathscr{U}$-construct that the lazy $\lambda$-calculus
  can macro-express, and
  (b)~the converse fails.

  \paragraph{Part (b): the lazy $\lambda$-calculus $\not\geq$
  inheritance-calculus}
  The $\lambda$-calculus contains its own constructors
  ($\lambda$, application, variable) and can macro-express
  $\mathbf{let}$-binding (Lemma~\ref{lem:let-macro}), but
  cannot macro-express the inheritance constructors of
  $\mathscr{U}$
  (Theorem~\ref{thm:cartesian-nonexpressibility}).
  Inheritance-calculus contains those constructors.
  Therefore the lazy $\lambda$-calculus is not at least
  as expressive as inheritance-calculus with respect
  to~$\mathscr{U}$.

  \paragraph{Part (a): inheritance-calculus $\geq$ the lazy
  $\lambda$-calculus}
  We show that inheritance-calculus can macro-express
  every $\mathscr{U}$-construct that the
  $\lambda$-calculus contains or can macro-express.
  The $\lambda$-calculus contains $\lambda$-abstraction,
  application, and variable; it can macro-express
  $\mathbf{let}$-binding
  (Lemma~\ref{lem:let-macro}).
  We must show that inheritance-calculus can
  macro-express all four of these constructs.
  The translation proceeds in two stages:
  \begin{enumerate}
    \item \emph{ANF conversion.}
      The A-normal form
      transformation~\cite{flanagan1993-essence-compiling-continuations}
      replaces nested applications with
      $\mathbf{let}$-bindings; it is macro-expressible
      in the $\lambda$-calculus by
      Lemma~\ref{lem:anf-equiv}.
      In the reverse direction,
      $\mathbf{let}$-bindings are macro-expressible in
      the $\lambda$-calculus by
      Proposition~3.7 of
      Felleisen~\cite{felleisen1991-expressive-power}
      (restated as Lemma~\ref{lem:let-macro}).
      Thus ANF conversion is a macro-expressible
      equivalence between the $\lambda$-calculus and its
      ANF fragment.
    \item \emph{The translation $\mathcal{T}$.}
      Each ANF $\lambda$-calculus construct maps to a
      fixed syntactic abstraction over
      inheritance-calculus
      (Theorem~\ref{thm:forward-macro}), satisfying
      condition~E4 of
      Felleisen~\cite[Definition~3.11]{felleisen1991-expressive-power}.
  \end{enumerate}
  By Lemma~\ref{lem:macro-composition}, composing the
  two macro-expressible translations gives a
  macro-expressible translation from
  the $\lambda$-calculus (with $\mathbf{let}$) into
  inheritance-calculus.
  Since inheritance-calculus trivially contains its own
  constructors, it can macro-express every
  $\mathscr{U}$-construct that the $\lambda$-calculus
  contains or can macro-express.
  \fi
\end{proof}

\section{\bilingual{Multi-Target \texorpdfstring{$\this$}{this} Resolution in Other Systems}{其他系统中的多目标 \texorpdfstring{$\this$}{this} 解析}}
\label{app:scala-multi-target}

\ifInheritanceChinese
本附录说明两个具有代表性的系统——Scala~3 和 NixOS 模块系统——拒绝了继承演算能够自然处理的多目标 $\this$ 解析模式。
\else
This appendix demonstrates that two representative systems, Scala~3
and the NixOS module system, reject the multi-target $\this$ resolution
pattern that inheritance-calculus handles naturally.
\fi

\subsection*{\bilingual{Scala~3}{Scala~3}}

\ifInheritanceChinese
考虑两个独立扩展同一外部类的对象，每个对象各自提供内部 trait 的一份副本：
\else
Consider two objects that independently extend an outer class,
each providing its own copy of an inner trait:
\fi

\begin{lstlisting}[style=scala]
class MyOuter:
  trait MyInner:
    def outer = MyOuter.this

object Object1 extends MyOuter
object Object2 extends MyOuter
object HasMultipleOuters extends Object1.MyInner
                     with Object2.MyInner
\end{lstlisting}

\noindent
\ifInheritanceChinese
Scala~3 以如下错误拒绝了 \texttt{HasMultipleOuters}：
\else
Scala~3 rejects \texttt{HasMultipleOuters} with the error:
\fi

\begin{lstlisting}
trait MyInner is extended twice
object HasMultipleOuters cannot be instantiated since
  it has conflicting base types
  Object1.MyInner and Object2.MyInner
\end{lstlisting}

\noindent
\ifInheritanceChinese
该拒绝并非表面层面的限制，而是 DOT 类型论基础的必然后果~\cite{amin2016-dependent-object-types}。在 DOT 中，一个对象被构造为 $\nu(x : T)\, d$，其中 self 变量~$x$ 绑定到\emph{单一}对象。\texttt{Object1.MyInner} 和 \texttt{Object2.MyInner} 是不同的路径依赖类型，每个都将外部 $\this$ 约束到不同的对象。它们的交集要求 \texttt{outer} 同时返回 \texttt{Object1} 和 \texttt{Object2}，但 DOT 的 self 变量是单值的，使得这一交集不可实现。
\else
The rejection is not a surface-level restriction but a consequence
of DOT's type-theoretic foundations~\cite{amin2016-dependent-object-types}.
In DOT, an object is constructed as $\nu(x : T)\, d$, where the
self variable~$x$ binds to a \emph{single} object.
\texttt{Object1.MyInner} and \texttt{Object2.MyInner} are distinct
path-dependent types, each constraining the outer $\this$
to a different object.
Their intersection requires \texttt{outer} to simultaneously
return \texttt{Object1} and \texttt{Object2}, but DOT's self
variable is single-valued, making this intersection unrealizable.
\fi

\ifInheritanceChinese
等价的继承演算定义如下：
\else
The equivalent inheritance-calculus definition is:
\fi
\begin{align*}
  \{ \quad \mathrm{MyOuter} &\mapsto \{
      \mathrm{MyInner} \mapsto \{
    \mathrm{outer} \mapsto \{\qualifiedthis{\mathrm{MyOuter}}\}\}\}, \\
    \mathrm{Object1} &\mapsto \{\mathrm{MyOuter}\}, \\
    \mathrm{Object2} &\mapsto \{\mathrm{MyOuter}\}, \\
    \mathrm{HasMultipleOuters} &\mapsto \{
      \mathrm{Object1}.\mathrm{MyInner},\;
    \mathrm{Object2}.\mathrm{MyInner}\}
  \quad \}
\end{align*}
\ifInheritanceChinese
此定义在继承演算中是良定义的。$\mathrm{outer}$ 内部的引用 $\qualifiedthis{\mathrm{MyOuter}}$ 具有 de~Bruijn 索引 $n = 1$：从 $\mathrm{outer}$ 的封闭作用域 $\mathrm{MyInner}$ 出发，经过一步 $\this$ 便到达 $\mathrm{MyOuter}$。当 $\mathrm{HasMultipleOuters}$ 同时继承 $\mathrm{Object1}.\mathrm{MyInner}$ 和 $\mathrm{Object2}.\mathrm{MyInner}$ 时，$\this$ 函数（方程~\ref{eq:this}）通过在 $\supers(\mathrm{HasMultipleOuters})$ 中搜索与 $\mathrm{MyInner}$ 定义位置匹配的覆盖路径来解析 $\qualifiedthis{\mathrm{MyOuter}}$。它找到两条继承位置路径，一条经由 $\mathrm{Object1}$，另一条经由 $\mathrm{Object2}$，并返回两者。两条路径都通向继承自同一 $\mathrm{MyOuter}$ 的记录，因此查询 $\mathrm{HasMultipleOuters}.\mathrm{outer}$ 无论走哪条路由都能得到 $\mathrm{MyOuter}$ 的 $\properties$。由于属性解析通过集合并来聚合结果（这是一种本质上幂等的操作），来自两条路由的相同属性定义无冲突地合并。因此，该演算消除了对任何人为路由优先级或线性化规则的需求。
\else
This is well-defined in inheritance-calculus.
The reference
$\qualifiedthis{\mathrm{MyOuter}}$ inside $\mathrm{outer}$
has de~Bruijn index $n = 1$:
starting from $\mathrm{outer}$'s enclosing scope
$\mathrm{MyInner}$, one $\this$ step reaches
$\mathrm{MyOuter}$.
When $\mathrm{HasMultipleOuters}$ inherits from both
$\mathrm{Object1}.\mathrm{MyInner}$ and
$\mathrm{Object2}.\mathrm{MyInner}$,
the $\this$ function (equation~\ref{eq:this}) resolves
$\qualifiedthis{\mathrm{MyOuter}}$ by searching through
$\supers(\mathrm{HasMultipleOuters})$ for override paths matching
$\mathrm{MyInner}$'s definition site.
It finds two inheritance-site paths, one through
$\mathrm{Object1}$ and one through $\mathrm{Object2}$, and
returns both.
Both paths lead to records that inherit from the same
$\mathrm{MyOuter}$, so querying
$\mathrm{HasMultipleOuters}.\mathrm{outer}$ yields the
$\properties$ of $\mathrm{MyOuter}$ regardless of which
route is taken.
Because attribute resolution aggregates results via set union (which is an intrinsically idempotent operation), identical property definitions from the two routes merge without conflict. Thus, the calculus eliminates the need for any artificial route-priority or linearization rules.
\fi

\subsection*{\bilingual{NixOS Module System}{NixOS 模块系统}}

\ifInheritanceChinese
NixOS 模块系统以静态错误拒绝了相同的模式。该翻译使用了模块系统自身的抽象：\texttt{deferredModule} 对应 trait（即未求值的模块值，被导入到其他不动点中），\texttt{submoduleWith} 对应对象（在其自身的不动点中求值）。
\else
The NixOS module system rejects the same pattern with a static error.
The translation uses the module system's own abstractions:
\texttt{deferredModule} for traits, which are unevaluated module values imported
into other fixpoints, and \texttt{submoduleWith} for objects, which are evaluated
in their own fixpoints.
\fi

\begin{lstlisting}[style=nix]
let
  lib = (import <nixpkgs> { }).lib;
  result = lib.evalModules {
    modules = [
      (toplevel@{ config, ... }: {
        # class MyOuter { trait MyInner {
        #   def outer = MyOuter.this } }
        options.MyOuter = lib.mkOption {
          default = { };
          type = lib.types.deferredModuleWith {
            staticModules = [
              (MyOuter: {
                options.MyInner = lib.mkOption {
                  default = { };
                  type = lib.types.deferredModuleWith {
                    staticModules = [
                      (MyInner: {
                        options.outer = lib.mkOption {
                          default = { };
                          type =
                            lib.types.deferredModuleWith {
                              staticModules =
                                [ toplevel.config.MyOuter ];
                            };
                        };
                      })
                    ];
                  };
                };
              })
            ];
          };
        };
        # object Object1 extends MyOuter
        options.Object1 = lib.mkOption {
          default = { };
          type = lib.types.submoduleWith {
            modules = [ toplevel.config.MyOuter ];
          };
        };
        # object Object2 extends MyOuter
        options.Object2 = lib.mkOption {
          default = { };
          type = lib.types.submoduleWith {
            modules = [ toplevel.config.MyOuter ];
          };
        };
        # object HasMultipleOuters extends
        #   Object1.MyInner with Object2.MyInner
        options.HasMultipleOuters = lib.mkOption {
          default = { };
          type = lib.types.submoduleWith {
            modules = [
            toplevel.config.Object1.MyInner
            toplevel.config.Object2.MyInner
          ];
          };
        };
      })
    ];
  };
in
  builtins.attrNames result.config.HasMultipleOuters.outer
\end{lstlisting}

\noindent
\ifInheritanceChinese
NixOS 模块系统以如下错误拒绝了此代码：
\else
The NixOS module system rejects this with:
\fi

\begin{lstlisting}
error: The option `HasMultipleOuters.outer'
  in `<unknown-file>'
  is already declared
  in `<unknown-file>'.
\end{lstlisting}

\noindent
\ifInheritanceChinese
该拒绝发生的原因是：\texttt{Object1.\allowbreak{}My\-Inner} 和 \texttt{Object2.\allowbreak{}My\-Inner} 各自携带了 \texttt{My\-Inner} 延迟模块的完整 \texttt{static\-Modules}，其中包括声明 \texttt{options.\allowbreak{}outer}。当 \texttt{HasMultipleOuters} 同时导入两者时，模块系统的 \texttt{merge\-Option\-Decls} 遇到了同一选项的两个声明，并产生静态错误；它无法表达两个声明源自同一 trait 且应被统一这一事实。
\else
The rejection occurs because
\texttt{Object1.\allowbreak{}My\-Inner} and
\texttt{Object2.\allowbreak{}My\-Inner} each carry the full
\texttt{static\-Modules} of the \texttt{My\-Inner} deferred
module, including the declaration
\texttt{options.\allowbreak{}outer}.
When \texttt{HasMultipleOuters} imports both, the module system's
\texttt{merge\-Option\-Decls} encounters two declarations of the
same option and raises a static error; it cannot express that
both declarations originate from the same trait and should be
unified.
\fi

\ifInheritanceChinese
该拒绝与上文 Scala 的``冲突基类型''错误类似：两个系统都假设每个选项或类型成员只有单一的声明位置，并拒绝两条继承路由独立引入同一声明的多目标情形。在继承演算中，相同的模式是良定义的，因为 $\overrides$（方程~\ref{eq:overrides}）能够识别两条路由通向同一定义，而 $\supers$（方程~\ref{eq:supers}）则无重复地收集两个继承位置上下文。
\else
The rejection is analogous to Scala's ``conflicting base types''
error above: both systems assume that each option or type member
has a single declaration site, and reject the multi-target situation
where two inheritance routes introduce the same declaration
independently.
In inheritance-calculus, the same pattern is well-defined because
$\overrides$ (equation~\ref{eq:overrides}) recognizes that both
routes lead to the same definition, and $\supers$
(equation~\ref{eq:supers}) collects both inheritance-site
contexts without duplication.
\fi

\section{Encoding Datalog in Inheritance-Calculus}
\label{app:datalog-encoding}

Section~\ref{sec:emergent-phenomena} observed that
inheritance-calculus natively produces the relational semantics
of logic programming.
This appendix makes that observation precise by giving a
systematic encoding of Datalog into inheritance-calculus.
The encoding uses no extensions to the calculus; it is
expressible in the three primitives of Section~\ref{sec:syntax}.
All examples in this appendix have executable MIXINv2
implementations\anon[~(supplementary material)]{~\cite{mixin2025}}.
There are three ingredients to keep separate.
First, relations are stored as tries, so tuple lookup is indexed
by prefixes rather than by scanning a flat set of facts.
Second, the control skeleton is written in the visitor-pattern
style already used in Section~\ref{sec:case-study}; from a
functional-programming viewpoint, this is the same role played by
Scott-encoded case analysis.
Third, inheritance union and tabled evaluation give the resulting
program its relational, least-fixed-point semantics.
The appendix proceeds in that order: data representation first,
then control flow, then recursion.

\subsection{Relations as Tries}

A binary relation over a finite domain $D = \{a, b, c, \ldots\}$
is represented as a \emph{trie}: a mixin tree whose paths
encode tuples.
The trie is built from three base scopes:
$\mathrm{Relation}$ (abstract base),
$\mathrm{Cons}$ (non-empty node), and $\mathrm{Nil}$ (terminator).
Each constant $d \in D$ inherits from $\mathrm{Cons}$ and carries
a $\mathrm{TailMap}$ that provides per-constant sub-tries:
\begin{align*}
  &\mathrm{Relation} \mapsto \{\},\quad
  \mathrm{Cons} \mapsto \{\mathrm{Relation},\;
  \mathrm{Tail} \mapsto \{\mathrm{Relation}\}\},\quad
  \mathrm{Nil} \mapsto \{\mathrm{Relation}\}\\[2pt]
  &d \mapsto \{\mathrm{Cons},\;
    \mathrm{Tail} \mapsto \{\},\;
  \mathrm{TailMap} \mapsto \{d \mapsto \{\mathrm{Tail}\}\}\}
\end{align*}
A fact $R(d_1, d_2)$ is encoded as a trie path of depth~2,
with a $\mathrm{Nil}$ terminator.
The per-constant sub-tries are accessed via
$\mathrm{TailMap}.d$:
\begin{align*}
  R(a, b) \;\leadsto\quad
  R \mapsto \{a,\; \mathrm{TailMap} \mapsto \{a \mapsto
  \{b,\; \mathrm{TailMap} \mapsto \{b \mapsto \{\mathrm{Nil}\}\}\}\}\}
\end{align*}
Multiple facts for the same relation are accumulated by
inheritance (trie union).
Adding $R(b, c)$:
\[
  R \mapsto \left\{
    \begin{aligned}
      &a,\; b,\\
      &\mathrm{TailMap} \mapsto \left\{
        \begin{aligned}
          &a \mapsto \{b,\; \mathrm{TailMap} \mapsto \{b \mapsto \{\mathrm{Nil}\}\}\},\\
          &b \mapsto \{c,\; \mathrm{TailMap} \mapsto \{c \mapsto \{\mathrm{Nil}\}\}\}
      \end{aligned}\right\}
  \end{aligned}\right\}
\]
The point of the trie representation is not merely that it is a
tree-shaped encoding.
It also provides the indexing structure used by the subsequent
joins: once a constant $d$ has been selected, the encoding moves
directly to $\mathrm{TailMap}.d$ and continues from the matching
sub-trie.
This is why the appendix speaks about tries rather than about an
arbitrary encoding of finite sets of tuples.

\subsection{Infrastructure}

Each constant $d$ carries four groups of boilerplate
definitions used by the encoding.
These play a role analogous to derived type-class instances
in Haskell: each constant must provide a fixed set of
properties that the encoding's generic patterns rely on.
Inheritance-calculus has no built-in derivation mechanism,
so the boilerplate is written out manually; however, the
amount of boilerplate per constant is $O(1)$
and does not grow with the domain size $|D|$,
so adding a new constant
does not require modifying any existing scope, preserving
the expression-problem property.

% 我们可以通过引入类似ECMAScript的Symbol机制来消除这些重复定义，但这让语言更复杂了

\begin{description}
  \item[Acceptance]
    The visitor-pattern infrastructure from
    Section~\ref{sec:case-study}.
    Operationally, this is the object-oriented spelling of a
    case split on the current constant.
    Each constant $d$ defines a $\mathrm{VisitorMap}$
    containing a branch keyed by $d$, a $\mathrm{Visitor}$
    with a $\mathrm{Visited}$ slot for per-constant results,
    and an $\mathrm{Accepted}$ property that delegates to the
    matching branch's $\mathrm{Visited}$ value.
    This is the dispatch mechanism: given a scope that may
    contain multiple constants, $\mathrm{Acceptance}$ selects the
    branch corresponding to each constant present and
    collects the per-constant $\mathrm{Visited}$ results.

  \item[WithVisitorTail]
    Maps each constant $d$ to its sub-trie entry
    $\mathrm{TailMap}.d$, labeled as $\mathrm{VisitorTail}$.
    This provides the right-hand side of a join\footnote{
      In the context of the visitor pattern, this acts as the visitor.
    }: when two relations' $\mathrm{VisitorMap}$ entries are
    composed, $\mathrm{VisitorTail}$ carries the joined
    relation's row data for the matched constant.

  \item[WithVisiteeTail]
    Maps each constant $d$ to its sub-trie entry
    $\mathrm{TailMap}.d$, labeled as $\mathrm{VisiteeTail}$.
    This provides the left-hand side\footnote{
      This acts as the visitee.
    }: the
    relation being iterated. Together,
    $\mathrm{WithVisitorTail}$ and $\mathrm{WithVisiteeTail}$
    enable value-level joins where only constants present
    in \emph{both} relations produce output.

  \item[Replacement]
    Given a $\mathrm{NewTail}$ value, produces a
    $\mathrm{Replaced}$ scope that inherits the label $d$
    but substitutes $\mathrm{TailMap}.d$ with
    $\mathrm{NewTail}$.
    This is the output construction mechanism: it
    builds output tuples by substituting new sub-tries into
    constant wrappers.
    Crucially, $\mathrm{Replacement}$ is defined on the
    per-constant $\mathrm{TailMap}$ entry (not on the merged
    scope), ensuring that output construction operates on
    individual constants rather than their union.
\end{description}

As a running example, we use the executable two-hop encoding
$\mathrm{twoHop}(X, Z) \mathrel{{:}\!{-}} \mathrm{edge}(X, Y),\; \mathrm{edge}(Y, Z)$
from
\anon[\nolinkurl{RelationalTwoHopRepeatedVisitor.mixin.yaml}~(supplementary material)]{\nolinkurl{RelationalTwoHopRepeatedVisitor.mixin.yaml}~\cite{mixin2025}}.
Its top-level scopes are named after the stages of the
encoding: \texttt{\_XAcceptance} iterates the first column,
\texttt{\_YAcceptance} performs the shared-variable join,
\texttt{\_NilAcceptance0} confirms the end of the first fact,
\texttt{\_ZAcceptance} iterates the output column, and
\texttt{\_NilAcceptance1} triggers the final
\texttt{Replacement} chain.
The schematic encodings below should be read as abstractions
of that executable MIXINv2 file.
If one temporarily ignores relational accumulation and reads
each dispatch as an ordinary case split, the control flow is:
choose $X$ from the first relation, choose a matching $Y$,
confirm that the first tuple is complete, choose $Z$ from the
second relation's matching row, confirm that the second tuple
is complete, and emit $(X, Z)$.
The rest of the appendix explains how that familiar
single-valued control skeleton is realised by visitor dispatch
over tries and then lifted to relational semantics by
inheritance union.

\subsection{Encoding Rules}

A Datalog rule $H(\bar{X}) \mathrel{{:}\!{-}} B_1(\bar{Y}_1), \ldots, B_n(\bar{Y}_n)$
is encoded by nesting three mechanisms inside that control
skeleton:

\begin{enumerate}
  \item \textbf{Column iteration.}
    For each column position, an $\mathrm{Acceptance}$ dispatch
    iterates over the constants present.
    The $\mathrm{VisitorMap}$ is populated from
    $\mathrm{WithVisiteeTail}$, so each matched constant's
    $\mathrm{Visitor}$ receives $\mathrm{VisiteeTail}$: the
    per-constant sub-trie at that position.
    The result of processing each constant is placed in
    $\mathrm{Visited}$, and $\mathrm{Accepted}$ collects all
    per-constant results.
    (Details in Sections~\ref{sec:encoding-variables},
    \ref{sec:encoding-constants}, and~\ref{sec:encoding-multi}.)

  \item \textbf{Join.}
    For each shared variable between body atoms, the inner
    $\mathrm{Acceptance}$'s $\mathrm{VisitorMap}$ inherits from
    \emph{both} $\mathrm{WithVisitorTail}$ (from the joined
    relation) and $\mathrm{WithVisiteeTail}$ (from the iterated
    relation).
    A constant $d$ produces a $\mathrm{Visitor}$ entry only if
    $d$ appears in both relations, achieving a value-level
    join.
    The $\mathrm{Visitor}$ then receives both
    $\mathrm{VisiteeTail}$ and $\mathrm{VisitorTail}$,
    providing access to both relations' data for the
    matched constant.

  \item \textbf{Output construction.}
    Inside the innermost $\mathrm{Nil}$ confirmation (verifying
    that each body atom's tuple is complete), chain
    $\mathrm{Replacement}$ operations from the innermost
    output column outward, each setting $\mathrm{NewTail}$ to
    the previous $\mathrm{Replaced}$ value, terminated by
    $\mathrm{Nil}$.
    The final $\mathrm{Replaced}$ value is assigned to
    $\mathrm{Visited}$.
\end{enumerate}

\subsection{Variable Join}
\label{sec:encoding-variables}

Consider the simplest join:
$\mathrm{composed}(X, Z) \mathrel{{:}\!{-}} \alpha(X, Y),\; \beta(Y, Z)$.
The executable self-join instance of this pattern is
\anon[\nolinkurl{RelationalTwoHopRepeatedVisitor.mixin.yaml}~(supplementary material)]{\nolinkurl{RelationalTwoHopRepeatedVisitor.mixin.yaml}~\cite{mixin2025}},
whose scope names are shown in parentheses below.
This is the right place to see the division of labour from the
section introduction: visitor dispatch provides the control flow,
the trie provides direct access to the matching tails, and
$\supers$ accumulates the outputs produced by each successful
match.

The shared variable $Y$ requires matching constants
between $\alpha$'s second column and $\beta$'s first column.
The encoding nests five $\mathrm{Acceptance}$ dispatches:
iterate $X$ from $\alpha$, join on $Y$ between $\alpha$
and $\beta$, confirm $\mathrm{Nil}$ after $Y$ in $\alpha$,
iterate $Z$ from $\beta$'s matched row, and confirm
$\mathrm{Nil}$ after $Z$.
We present each level as a named sub-scope.

\paragraph{Level~1: iterate $X$ from $\alpha$}
(Executable scope: \texttt{\_XAcceptance}.)
{\small
  \[
    \mathrm{xAcc} \mapsto \left\{
      \begin{aligned}
        &\alpha.\mathrm{Acceptance},\; \mathrm{VisitorMap} \mapsto \{\alpha.\mathrm{WithVisiteeTail}\},\\
        &\mathrm{Visitor} \mapsto \left\{
          \begin{aligned}
            &\_X \mapsto \{\mathrm{VisiteeTail}\},\; \ldots,\\
            &\mathrm{Visited} \mapsto \{\mathrm{yAcc}.\mathrm{Accepted}\}
        \end{aligned}\right\}
    \end{aligned}\right\}
\]}

\paragraph{Level~2: join $\alpha$'s $Y$-column with $\beta$'s first column}

(Executable scope: \texttt{\_YAcceptance}.)

The $\mathrm{VisitorMap}$ inherits from \emph{both}
$\beta.\mathrm{WithVisitorTail}$ and
$\_X.\mathrm{WithVisiteeTail}$, producing an entry for
constant $d$ only if $d$ appears in both $\alpha$'s
second column (via $\_X$) and $\beta$'s first column.
{\small
  \[
    \mathrm{yAcc} \mapsto \left\{
      \begin{aligned}
        &\_X.\mathrm{Acceptance},\\
        &\mathrm{VisitorMap} \mapsto \{\beta.\mathrm{WithVisitorTail},\; \_X.\mathrm{WithVisiteeTail}\},\\
        &\mathrm{Visitor} \mapsto \left\{
          \begin{aligned}
            &\_Y_0 \mapsto \{\mathrm{VisiteeTail}\},\; \mathrm{VisitorTail} \mapsto \{\beta.\mathrm{Tail}\},\\
            &\ldots,\\
            &\mathrm{Visited} \mapsto \{\mathrm{nilAcc}_0.\mathrm{Accepted}\}
        \end{aligned}\right\}
    \end{aligned}\right\}
\]}

\paragraph{Level~3: confirm $\mathrm{Nil}$ after $Y$ in $\alpha$}

(Executable scope: \texttt{\_NilAcceptance0}.)

This checks that $\alpha(X, Y)$ is a complete two-column
fact (its trie path terminates at $\mathrm{Nil}$ after $Y$).
Inside the $\mathrm{Nil}$ branch, the encoding proceeds to
iterate $\beta$'s second column via $\mathrm{VisitorTail}$:
{\small
  \[
    \mathrm{nilAcc}_0 \mapsto \left\{
      \begin{aligned}
        &\_Y_0.\mathrm{Acceptance},\\
        &\mathrm{VisitorMap} \mapsto \left\{\mathrm{Nil} \mapsto \left\{
            \begin{aligned}
              &\ldots,\\
              &\mathrm{Visited} \mapsto \{\mathrm{zAcc}.\mathrm{Accepted}\}
        \end{aligned}\right\}\right\}
    \end{aligned}\right\}
\]}

\paragraph{Level~4: iterate $Z$ from $\beta$'s matched row}
(Executable scope: \texttt{\_ZAcceptance}.)
{\small
  \[
    \mathrm{zAcc} \mapsto \left\{
      \begin{aligned}
        &\mathrm{VisitorTail}.\mathrm{Acceptance},\; \mathrm{VisitorMap} \mapsto \{\mathrm{VisitorTail}.\mathrm{WithVisiteeTail}\},\\
        &\mathrm{Visitor} \mapsto \left\{
          \begin{aligned}
            &\_Z \mapsto \{\mathrm{VisiteeTail}\},\; \ldots,\\
            &\mathrm{Visited} \mapsto \{\mathrm{nilAcc}_1.\mathrm{Accepted}\}
        \end{aligned}\right\}
    \end{aligned}\right\}
\]}

\paragraph{Level~5: confirm $\mathrm{Nil}$ after $Z$ and construct output}
(Executable scope: \texttt{\_NilAcceptance1}, with
\texttt{\_ZReplacement} and \texttt{\_XReplacement}.)
{\small
  \[
    \mathrm{nilAcc}_1 \mapsto \left\{
      \begin{aligned}
        &\mathrm{VisiteeTail}.\mathrm{Acceptance},\\
        &\mathrm{VisitorMap} \mapsto \left\{\mathrm{Nil} \mapsto \left\{
            \begin{aligned}
              &\mathrm{zR} \mapsto \{\_Z.\mathrm{Replacement},\; \mathrm{NewTail} \mapsto \{\mathrm{Nil}\}\},\\
              &\mathrm{xR} \mapsto \{\_X.\mathrm{Replacement},\; \mathrm{NewTail} \mapsto \{\mathrm{zR}.\mathrm{Replaced}\}\},\\
              &\mathrm{Visited} \mapsto \{\mathrm{xR}.\mathrm{Replaced}\}
        \end{aligned}\right\}\right\}
    \end{aligned}\right\}
\]}

\paragraph{Final result}
\[
  \mathrm{composed} \mapsto \{\mathrm{xAcc}.\mathrm{Accepted}\}
\]

\paragraph{Semantics}
At Level~2, the combined $\mathrm{VisitorMap}$ ensures that a
constant $d$ produces a $\mathrm{Visitor}$ only when $d$
appears in both $\alpha$'s second column \emph{and}
$\beta$'s first column.
The $\mathrm{Visitor}$ for each matched constant $d$
independently receives $\mathrm{VisiteeTail}$ ($\alpha$'s
data for~$d$) and $\mathrm{VisitorTail}$ ($\beta$'s row
for~$d$).
Output construction (Level~5) operates inside this
per-constant $\mathrm{Visitor}$, so $\mathrm{Replacement}$
acts on the \emph{individual} constant's data, not on the
merged scope.
Each matched constant $d$ produces its own output tuples
in $\mathrm{Visited}$, and these are collected by
$\mathrm{Accepted}$ via trie union ($\supers$).
This yields standard Datalog tuple-level semantics: only
constants that actually match between both relations
contribute to the output.
For the concrete instance
$\mathrm{edge} = \{(a,b),\; (b,c),\; (c,a)\}$, as encoded in
\anon[\nolinkurl{RelationalTwoHopRepeatedVisitor.mixin.yaml}~(supplementary material)]{\nolinkurl{RelationalTwoHopRepeatedVisitor.mixin.yaml}~\cite{mixin2025}},
Level~1 selects $a$, Level~2 matches $b$, Level~3 confirms
that $\mathrm{edge}(a,b)$ ends in $\mathrm{Nil}$, Level~4
iterates $\mathrm{edge}.\mathrm{TailMap}.b$ and selects $c$,
and Level~5 confirms $\mathrm{Nil}$ again before the two
$\mathrm{Replacement}$ operations reconstruct $(a,c)$.

\subsection{Constant Match and Navigation}
\label{sec:encoding-constants}

Constants in Datalog atoms use two mechanisms, neither of
which is a join.
They are simpler than variable joins because no value from a
second relation needs to be matched against the currently chosen
constant:

\paragraph{Type existence match}
An atom $R(X, b)$ with a constant in the second position
requires checking that constant $b$ exists in $R$'s
second column at the matched row.
This is encoded as an $\mathrm{Acceptance}$ with a
hand-crafted $\mathrm{VisitorMap}$ containing \emph{only}
the $b$ branch (instead of composing from
  $\mathrm{WithVisitorTail}$, which would create branches for
all constants):
{\small
  \begin{align*}
    \mathrm{VisitorMap} \mapsto \{b \mapsto \{\ldots\}\}
\end{align*}}
If $b$ exists in the dispatching scope, the
$b$ branch fires; otherwise, $\mathrm{Accepted}$
receives no matching branch.

\paragraph{Direct navigation}
An atom $R(b, Z)$ with a constant in the first position
navigates directly to $b$'s sub-trie:
\[
  \mathrm{edgeB}_Z \mapsto \{R.\mathrm{TailMap}.b\}
\]
This bypasses the visitor mechanism entirely, accessing
the per-constant $\mathrm{TailMap}$ entry to restrict the
scope to only $b$'s data.

\paragraph{Example}
$\mathrm{viaB}(X, Z) \mathrel{{:}\!{-}} \mathrm{edge}(X, b),\; \mathrm{edge}(b, Z)$
combines both mechanisms.
The executable MIXINv2 encoding is
\anon[\nolinkurl{RelationalConstant.mixin.yaml}~(supplementary material)]{\nolinkurl{RelationalConstant.mixin.yaml}~\cite{mixin2025}}.
The encoding iterates $\mathrm{edge}$ for $X$ (Level~1),
then checks for constant $b$ in $X$'s row with a
hand-crafted $\mathrm{VisitorMap}$ (Level~2), navigates
$\mathrm{edge}.\mathrm{TailMap}.b$ for $Z$ (Level~3),
confirms $\mathrm{Nil}$ (Level~4), and constructs $(X, Z)$
via $\mathrm{Replacement}$.

\subsection{Multi-Body Rules}
\label{sec:encoding-multi}

Rules with $n > 2$ body atoms chain additional join levels.
Each shared variable introduces one combined
$\mathrm{WithVisitorTail}$/$\mathrm{WithVisiteeTail}$ join
inside the previous level's $\mathrm{Visitor}$,
and the output construction is placed inside the innermost
$\mathrm{Nil}$ confirmation.
So the binary-rule encoding is the whole story structurally:
larger bodies only add more nested copies of the same pattern.

\paragraph{Three-body rule}
$\mathrm{chain}(X, W) \mathrel{{:}\!{-}} \alpha(X, Y),\; \beta(Y, Z),\; \gamma(Z, W)$
requires two joins: on $Y$ (between $\alpha$ and $\beta$)
and on $Z$ (between $\beta$ and $\gamma$).
The structure is: iterate $X$ from $\alpha$, join on $Y$
(getting both $\alpha$'s and $\beta$'s data), confirm
$\mathrm{Nil}$ after $Y$ in $\alpha$, then inside the
$\mathrm{Nil}$ branch join $\beta$'s $Z$-column with
$\gamma$'s first column, iterate $W$ from $\gamma$'s
matched row, confirm $\mathrm{Nil}$, and construct
$(X, W)$.
The $Z$-join is nested inside the $Y$-join's
$\mathrm{Nil}$ confirmation, ensuring that $\beta$'s data
for the matched $Y$ constant (available via
$\mathrm{VisitorTail}$) is used to iterate $\beta$'s
second column.

\paragraph{Mixed constants and variables}
$\mathrm{constJoin}(X, W) \mathrel{{:}\!{-}} \mathrm{edge}(X, Y),\; \mathrm{edge}(Y, b),\; \mathrm{edge}(b, W)$
combines a variable join on $Y$, a constant match on $b$
(hand-crafted $\mathrm{VisitorMap}$ with only the $b$ entry),
and direct navigation via $\mathrm{edge}.\mathrm{TailMap}.b$
for $W$.
The corresponding executable MIXINv2 file is
\anon[\nolinkurl{RelationalConstantJoin.mixin.yaml}~(supplementary material)]{\nolinkurl{RelationalConstantJoin.mixin.yaml}~\cite{mixin2025}}.

\subsection{Multi-Variable Join}

When a rule shares \emph{multiple} variables between two
body atoms, each shared variable requires its own
$\mathrm{Acceptance}$ check.
The checks are chained: the outer $\mathrm{Acceptance}$
joins on one column pair using combined
$\mathrm{WithVisitorTail}$/$\mathrm{WithVisiteeTail}$, and
its $\mathrm{Visitor}$ contains an inner $\mathrm{Acceptance}$
joining on another column pair.

\paragraph{Example}
$\mathrm{multiJoin}(X, Y) \mathrel{{:}\!{-}} r(X, Y),\; s(X, Y)$
shares both $X$ and $Y$.
The executable MIXINv2 encoding is
\anon[\nolinkurl{RelationalMultiVariableJoin.mixin.yaml}~(supplementary material)]{\nolinkurl{RelationalMultiVariableJoin.mixin.yaml}~\cite{mixin2025}}.
The outer $\mathrm{Acceptance}$ joins $r$'s and $s$'s first
columns (variable $X$) using combined
$\mathrm{WithVisitorTail}$/$\mathrm{WithVisiteeTail}$.
Inside the outer $\mathrm{Visitor}$, an inner
$\mathrm{Acceptance}$ joins $r$'s and $s$'s second columns
(variable $Y$) similarly.
Both joins must match for the output to be produced, and
the per-constant $\mathrm{Visited}$ pattern ensures that
only constants present in both relations contribute.

\subsection{Recursive Rules}
\label{sec:recursive-rules}

Recursive Datalog rules, such as transitive closure
$\mathrm{path}(X, Y) \mathrel{{:}\!{-}} \mathrm{edge}(X, Z),\; \mathrm{path}(Z, Y)$,
are encoded identically to non-recursive rules: the
head relation $\mathrm{path}$ appears both as a defined scope
and as a reference in the body (providing
$\mathrm{WithVisitorTail}$ for the $Z$-join).
This should not be read as ordinary $\this$ of a
single-valued lazy program.
Operationally, tabled evaluation
(Section~\ref{sec:mixin-trees},
Appendix~\ref{app:well-definedness})
computes $\mathrm{lfp}(T_P)$: each query into
$\mathrm{path}$ triggers further evaluation through
the recursive body, and the trie union ($\supers$)
accumulates all reachable tuples.
This corresponds exactly to \emph{Naive
evaluation}~\cite{bancilhon1986-recursive-query-strategies} of the
$T_P$
operator~\cite{vanemden1976-predicate-logic-semantics}:
each iteration applies every rule to the current
fact set and adds newly derived facts.
Because Datalog-encoded relations have finitely many
constants per query, the per-query answer domain is
finite, and convergence is guaranteed by the ascending
chain condition on a finite lattice.
\emph{Semi-Naive} evaluation~\cite{bancilhon1986-recursive-query-strategies},
which restricts each iteration to rules that use at
least one newly derived fact, is a natural
optimisation left for future work.

\paragraph{The recursive rule as monadic pseudocode}
The recursive rule above roughly corresponds to the
following pseudocode (Haskell-like syntax with
  \texttt{RebindableSyntax} and \texttt{OverloadedLists}
  semantics, where do-notation desugars to a
\texttt{Set}-based bind and list literals denote sets):
\begin{lstlisting}[style=haskell]
{-# LANGUAGE RebindableSyntax, OverloadedLists #-}
s >>= f = unions [f x | x <- s]
return x = [x]
guard True  = [()]
guard False = []
edge = [("a","b"), ("b","c"), ("c","a")]
path = edge <> do
  (x, y0) <- edge
  (y1, z) <- path
  guard (y0 == y1)
  return (x, z)
\end{lstlisting}
\noindent
Each line of the pseudocode has a direct counterpart in
the inheritance-calculus encoding
(Section~\ref{sec:encoding-variables}):
\begin{description}
  \item[\texttt{edge}] corresponds to the base-fact trie;
  \item[\texttt{path = edge <> do \ldots}]
    corresponds to $\mathrm{path}$ inheriting $\mathrm{edge}$
    (base facts) and the recursive body (derived facts via
    $\supers$);
  \item[\texttt{(x, y0) <- edge}]
    corresponds to Level~1 ($\mathrm{xAcc}$), iterating $X$
    from $\mathrm{edge}$;
  \item[\texttt{(y1, z) <- path}]
    corresponds to Level~2 ($\mathrm{yAcc}$), joining on $Y$
    with $\mathrm{path}$, and Level~4 ($\mathrm{zAcc}$),
    iterating $Z$;
  \item[\texttt{guard (y0 == y1)}]
    corresponds to the $\mathrm{WithVisitorTail} \cap
    \mathrm{WithVisiteeTail}$ intersection in Level~2: a
    constant produces a $\mathrm{Visitor}$ entry only if it
    appears in both relations;
  \item[\texttt{return (x, z)}]
    corresponds to the $\mathrm{Replacement}$ chain in
    Level~5, constructing the output tuple $(X, Z)$.
\end{description}
The only difference is the evaluation strategy:
Haskell evaluates \texttt{path} lazily (diverges),
whereas inheritance-calculus evaluates it via tabled
$T_P$ iteration (converges to $\mathrm{lfp}$).

However, this monadic program does not converge because
\texttt{path} is defined recursively.%
\footnote{
  Enabling the \texttt{RecursiveDo} extension does not help:
  \texttt{mfix} ties a knot at the \emph{element} level via laziness,
  whereas the transitive closure equation requires a
  fixed point at the \emph{container} level, that is, finding a
  set~$S$ such that $S = F(S)$.
}
Inheritance-calculus solves this through tabled evaluation
(Section~\ref{sec:mixin-trees}), which computes the
least fixed point of the $T_P$ operator by iterating
over the lattice of mixin trees, precisely the
container-level fixed point that the pseudocode above cannot express.

\subsection{Semantic Correspondence with Standard Datalog}
\label{sec:datalog-semantics}

The encoding produces standard Datalog tuple-level semantics.
Putting the previous subsections together: tries provide indexed
tuple storage, visitor dispatch supplies the control skeleton,
and per-constant $\mathrm{Replacement}$ keeps output
construction local to each witness.
The key mechanism is the per-constant $\mathrm{Visited}$
pattern: output construction ($\mathrm{Replacement}$)
occurs inside each matched constant's $\mathrm{Visitor}$
scope, operating on the \emph{individual} constant's data
rather than on the merged scope.
The $\mathrm{Accepted}$ property then collects all per-constant
$\mathrm{Visited}$ results via trie union ($\supers$).

For example, given $\mathrm{edge} = \{(a,b),\; (b,c),\; (c,a)\}$
(a~3-cycle), the two-hop self-join
$\mathrm{twoHop}(X, Z) \mathrel{{:}\!{-}} \mathrm{edge}(X, Y),\; \mathrm{edge}(Y, Z)$
produces exactly $3$ pairs:
$\{(a,c),\; (b,a),\; (c,b)\}$, matching standard Datalog.
This is the behavior exhibited by
\anon[\nolinkurl{RelationalTwoHopRepeatedVisitor.mixin.yaml}~(supplementary material)]{\nolinkurl{RelationalTwoHopRepeatedVisitor.mixin.yaml}~\cite{mixin2025}}.
When the $Y$-join matches constant $b$ (present in both
  $\mathrm{edge}$'s second column from $(a,b)$ and
$\mathrm{edge}$'s first column), the $\mathrm{Visitor}$ for
$b$ receives $\mathrm{VisitorTail} = \mathrm{edge}.\mathrm{TailMap}.b$
(containing only~$c$), and $\mathrm{Replacement}$ constructs
a single output tuple $(a, c)$, rather than all of $\{a,b,c\}$.

This value-level precision arises because each constant's
$\mathrm{Replacement}$ operates on $\mathrm{TailMap}.d$
(the per-constant sub-trie) rather than on the merged
trie.
The same point is visible in
the repeated-visitor self-join regression example\anon[~(supplementary material)]{~\cite{mixin2025}},
where repeated visitor dispatch prevents witnesses from
different constants from being merged spuriously.
Although $\bases$ (equation~\ref{eq:bases}) still computes
cross-joins at the primitive level, the encoding routes
computation through per-constant $\mathrm{TailMap}$ entries
before any merging occurs, achieving the same effect as
standard Datalog's tuple-level evaluation.

\section{Evaluation Trace of Multi-Target \texorpdfstring{$\this$}{this} Resolution}
\label{app:evaluation-trace}

This appendix gives a step-by-step evaluation trace of the program of
Appendix~\ref{app:scala-multi-target}, computing
$\properties((\text{``HasMultipleOuters''},\, \text{``outer''}))$ by the six
functions of Section~\ref{sec:definitions}
(equations~\ref{eq:properties} to~\ref{eq:this}).
Every call to $\overrides$, $\bases$, $\resolve$, $\this$, and $\supers$
is shown with its arguments and returned set.
The trace lives in the metalanguage: every argument and result is a
\emph{path}, that is, a sequence of labels
$(\ell_1, \ldots, \ell_n)$ identifying a position in the tree
(Section~\ref{sec:definitions}), not an inheritance-calculus expression.
To keep labels distinct from the metalanguage variables of
Section~\ref{sec:definitions} (the italic $p$, $\ell$, $n$, $S$), we
quote each concrete label, so a path is written as a parenthesized,
comma-separated sequence of quoted labels.
Thus $(\text{``HasMultipleOuters''}, \text{``outer''})$ is the path that
the inheritance-calculus projection
$\mathrm{HasMultipleOuters}.\mathrm{outer}$ identifies, and the root path
is the empty sequence $()$.

\paragraph{Program}
\begin{align*}
  \{ \quad \mathrm{MyOuter} &\mapsto \{
      \mathrm{MyInner} \mapsto \{
    \mathrm{outer} \mapsto \{\qualifiedthis{\mathrm{MyOuter}}\}\}\}, \\
    \mathrm{Object1} &\mapsto \{\mathrm{MyOuter}\}, \\
    \mathrm{Object2} &\mapsto \{\mathrm{MyOuter}\}, \\
    \mathrm{HasMultipleOuters} &\mapsto \{
      \mathrm{Object1}.\mathrm{MyInner},\;
    \mathrm{Object2}.\mathrm{MyInner}\}
  \quad \}
\end{align*}

\paragraph{AST primitives}
Parsing populates $\defines$ and $\inherits$
(Section~\ref{sec:definitions}):
\begin{align*}
  \defines(()) &= \{
    \text{``MyOuter''},\;
    \text{``Object1''},\;
    \text{``Object2''},\;
  \text{``HasMultipleOuters''}\}, \\
  \defines((\text{``MyOuter''})) &= \{\text{``MyInner''}\}, \\
  \defines((\text{``MyOuter''}, \text{``MyInner''})) &= \{\text{``outer''}\}, \\
  \inherits((\text{``MyOuter''}, \text{``MyInner''}, \text{``outer''})) &=
  \{(1,\, ())\}, \\
  \inherits((\text{``Object1''})) = \inherits((\text{``Object2''})) &=
  \{(0,\, (\text{``MyOuter''}))\}, \\
  \inherits((\text{``HasMultipleOuters''})) &=
  \{(0,\, (\text{``Object1''}, \text{``MyInner''})),\;
  (0,\, (\text{``Object2''}, \text{``MyInner''}))\}.
\end{align*}
All other $\defines$ and $\inherits$ are empty.
Each $\inherits$ entry is a reference pair $(n,\, w)$,
where $n$ is the number of upward steps and $w$ the
downward projection list (Section~\ref{sec:definitions}).

\paragraph{Trace}
We present the trace as a numbered list of semantic-function calls in
dependency order: a call appears after the calls it depends on, and each
item states, in English, which earlier items it uses, how it computes, and
what it yields. Each distinct call is listed once, so a single item may be
used by several later items, the reuse that an implementation could realize
by tabling~\cite{tamaki1986-tabled-resolution}.
The arguments and results are metalanguage values: a label is a quoted
string such as $\text{``outer''}$, and a path is a tuple of labels such as
$(\text{``HasMultipleOuters''}, \text{``outer''})$, with the root path
written $()$.
The query is
$\properties((\text{``HasMultipleOuters''}, \text{``outer''}))$.

\begin{enumerate}
  \item \label{trace:overrides-MyOuter-MyInner}Compute $\overrides((\text{``MyOuter''}, \text{``MyInner''}))$, the paths sharing the identity of $(\text{``MyOuter''}, \text{``MyInner''})$. It keeps $(\text{``MyOuter''}, \text{``MyInner''})$ and adds any same-label definition reached through the supers of its enclosing scope. Result: $\{(\text{``MyOuter''}, \text{``MyInner''})\}$.
  \item \label{trace:supers-MyOuter-MyInner}Compute $\supers((\text{``MyOuter''}, \text{``MyInner''}))$ using item~\ref{trace:overrides-MyOuter-MyInner}. It takes the reflexive-transitive closure of $\bases$ from $(\text{``MyOuter''}, \text{``MyInner''})$, pairing each reachable override with the inheritance site through which it is reached. Result: $\{((\text{``MyOuter''}), (\text{``MyOuter''}, \text{``MyInner''}))\}$.
  \item \label{trace:overrides-Has-outer}Compute $\overrides((\text{``HasMultipleOuters''}, \text{``outer''}))$, the paths sharing the identity of $(\text{``HasMultipleOuters''}, \text{``outer''})$ using item~\ref{trace:supers-MyOuter-MyInner}. It keeps $(\text{``HasMultipleOuters''}, \text{``outer''})$ and adds any same-label definition reached through the supers of its enclosing scope. Result: $\{(\text{``HasMultipleOuters''}, \text{``outer''}), (\text{``MyOuter''}, \text{``MyInner''}, \text{``outer''})\}$.
  \item \label{trace:bases-Has-outer}Compute $\bases((\text{``HasMultipleOuters''}, \text{``outer''}))$. $(\text{``HasMultipleOuters''}, \text{``outer''})$ carries no inheritance reference ($\inherits = \varnothing$), so it has no direct base. Result: $\varnothing$.
  \item \label{trace:bases-MyOuter-MyInner-outer}Compute $\bases((\text{``MyOuter''}, \text{``MyInner''}, \text{``outer''}))$. $(\text{``MyOuter''}, \text{``MyInner''}, \text{``outer''})$ carries the reference pair $(1, ())$ (take one step upward, then project the empty projection), whose target is obtained by item~\ref{trace:resolve-Has-MyOuter-MyInner-1}.
  \item \label{trace:this-Has-MyOuter-MyInner-1}Compute one $\this$ step: from the frontier $\{(\text{``HasMultipleOuters''})\}$, take the supers of each frontier path and keep those whose definition site is $(\text{``MyOuter''}, \text{``MyInner''})$, collecting their inheritance sites as the new frontier. Result: $\{(\text{``Object1''}), (\text{``Object2''})\}$.
  \item \label{trace:resolve-Has-MyOuter-MyInner-1}Compute $\resolve$ for the reference pair $(1, ())$ defined at $(\text{``MyOuter''}, \text{``MyInner''})$, reached from inheritance site $(\text{``HasMultipleOuters''})$ using item~\ref{trace:this-Has-MyOuter-MyInner-1}. It takes one step upward through $\this$ and then appends the projection. Result: $\{(\text{``Object1''}), (\text{``Object2''})\}$.
  \item \label{trace:resolve-0-MyOuter}Compute $\resolve$ for the reference pair $(0, (\text{``MyOuter''}))$ defined at $()$, reached from inheritance site $()$. It takes no step upward through $\this$ and then appends the projection. Result: $\{(\text{``MyOuter''})\}$.
  \item \label{trace:supers-Has-outer}Compute $\supers((\text{``HasMultipleOuters''}, \text{``outer''}))$ using items~\ref{trace:overrides-Has-outer}, \ref{trace:bases-Has-outer}, \ref{trace:bases-MyOuter-MyInner-outer}, \ref{trace:resolve-Has-MyOuter-MyInner-1} and~\ref{trace:resolve-0-MyOuter}. It takes the reflexive-transitive closure of $\bases$ from $(\text{``HasMultipleOuters''}, \text{``outer''})$, pairing each reachable override with the inheritance site through which it is reached. Result: $\{((), (\text{``MyOuter''})), ((), (\text{``Object1''})), ((), (\text{``Object2''})), ((\text{``HasMultipleOuters''}), (\text{``HasMultipleOuters''}, \text{``outer''})), ((\text{``HasMultipleOuters''}), (\text{``MyOuter''}, \text{``MyInner''}, \text{``outer''}))\}$.
  \item \label{trace:properties}Compute $\properties((\text{``HasMultipleOuters''}, \text{``outer''}))$ using item~\ref{trace:supers-Has-outer}. It unions the $\defines$ of every override path in that $\supers$ set. Only $(\text{``MyOuter''})$ defines a label, namely $\text{``MyInner''}$, reached through both $(\text{``Object1''})$ and $(\text{``Object2''})$; set union collapses the two routes. Result: $\{\text{``MyInner''}\}$.
\end{enumerate}


\begin{remark}
  The crux is the $\this$ step (item~\ref{trace:this-Has-MyOuter-MyInner-1}).
  The inheritance reference at
  $(\text{``MyOuter''}, \text{``MyInner''}, \text{``outer''})$ takes one step
  upward, so $\this$ moves one step up from
  $(\text{``MyOuter''}, \text{``MyInner''})$ and yields the frontier
  $\{(\text{``Object1''}), (\text{``Object2''})\}$, the two inheritance
  sites that incorporated $\text{``MyInner''}$. This is precisely the
  multi-target situation that Scala~3 and the NixOS module system reject
  (Appendix~\ref{app:scala-multi-target}): there $\this$ would have to
  select a single site. Reaching the properties of $(\text{``MyOuter''})$
  then happens because both $(\text{``Object1''})$ and
  $(\text{``Object2''})$ have $(\text{``MyOuter''})$ as a base, so the final
  $\supers$ contains $(\text{``MyOuter''})$ through both routes; since
  attribute resolution aggregates by set union, the two routes collapse to
  the single answer
  $\properties((\text{``HasMultipleOuters''}, \text{``outer''}))
  = \{\text{``MyInner''}\} = \properties((\text{``MyOuter''}))$.
\end{remark}

\fi % end \ifInheritanceAppendix

% Appendix-only build (supplement.tex): no main matter ran above, so emit the
% bibliography here, after the appendix.
\ifInheritanceBody\else
\bibliographystyle{ACM-Reference-Format}
\bibliography{references}
\fi

% \clearpage before closing CJK* ships every remaining page and float while the
% CJK active characters still mean CJK (the running head and any deferred float
% are typeset in the output routine, after \end{CJK*} would otherwise restore the
% default UTF-8 meaning).
\ifInheritanceChinese\clearpage\end{CJK*}\fi
\end{document}
.  The
% defaults include everything, so the full build (preprint.tex) sets
% nothing; submission.tex clears the appendices (POPL 2027 requires them
% as separate supplemental material), supplement.tex clears the body.
% The single bibliography is emitted exactly once, after the main matter and
% before the appendix, since both cite it. In the appendix-only build (no main
% matter) it instead follows the appendix; see the two guarded copies below.
\makeatletter
\@ifundefined{ifInheritanceBody}{\newif\ifInheritanceBody\InheritanceBodytrue}{}
\@ifundefined{ifInheritanceAppendix}{\newif\ifInheritanceAppendix\InheritanceAppendixtrue}{}
% Language switch for the reader-facing prose. Defaults to false (English), so the
% existing pdfLaTeX entries are unaffected; preprint-zh.tex sets it true and adds
% CJK support, still under pdfLaTeX. See the paper-writing skill (Bilingual
% Translation) for the convention.
\@ifundefined{ifInheritanceChinese}{\newif\ifInheritanceChinese\InheritanceChinesefalse}{}
% acmart writes \newlabel{tocindent<n>}{<dimen>} into the .aux to align the
% table of contents. Those bare-dimen labels break hyperref's \newlabel when
% another build imports this .aux via xr-hyper; the importing entries
% (submission.tex, supplement.tex) drop them at import time through
% xr-tocindent-filter.tex, so nothing is suppressed here at the source.
\makeatother

% --- bilingual reader-facing prose ----------------------------------
% The reader-facing body exists in two languages in this one file, selected by
% \ifInheritanceChinese: English (default) and a Chinese translation. Both build
% with pdfLaTeX, so the math renders identically. Block prose is guarded by
% \ifInheritanceChinese <chinese> \else <english> \fi; inline switches (title,
% captions, theorem notes interleaved with math) use \bilingual{<en>}{<zh>}.
% Author-side notes stay Chinese % comments in both builds.
\newcommand{\bilingual}[2]{#1}
\ifInheritanceChinese
  % CJKutf8 + the Arphic gbsn font render Chinese under pdfLaTeX (see
  % modules/texlive.nix). The document is wrapped in a CJK* environment opened
  % after \begin{document}; ASCII passes through, so English and math render as in
  % the English build.
  \usepackage{CJKutf8}
  \setlength{\emergencystretch}{2em}
  % Display the Chinese, but feed the English to hyperref's PDF strings, which
  % cannot encode the CJK active characters.
  \renewcommand{\bilingual}[2]{\texorpdfstring{#2}{#1}}
\fi

\begin{document}
% Open CJK* before \title so the Chinese in the title, abstract, and body is read
% inside it; \proofname and the theorem-environment names are localized here
% (their bytes must sit inside a CJK environment, and translating \thmname's
% argument leaves the shared counter and numbering identical to English). Closed
% before \end{document}.
\ifInheritanceChinese
  \begin{CJK*}{UTF8}{gbsn}
  \renewcommand{\proofname}{证明}
  \makeatletter
  \newcommand{\inheritanceThmName}[1]{\@ifundefined{inheritance@thm@#1}{#1}{\csname inheritance@thm@#1\endcsname}}
  \renewcommand{\thmname}[1]{\inheritanceThmName{#1}}
  \@namedef{inheritance@thm@Theorem}{定理}
  \@namedef{inheritance@thm@Lemma}{引理}
  \@namedef{inheritance@thm@Corollary}{推论}
  \@namedef{inheritance@thm@Proposition}{命题}
  \@namedef{inheritance@thm@Conjecture}{猜想}
  \@namedef{inheritance@thm@Definition}{定义}
  \@namedef{inheritance@thm@Example}{例}
  \@namedef{inheritance@thm@Remark}{注}
  \@namedef{inheritance@thm@Notation}{记法}
  \makeatother
\fi

% The optional short title feeds the running head and the PDF metadata; keep it
% ASCII English in both builds, since those are typeset outside CJK*.
\title[A Calculus of Inheritance]{\bilingual{A Calculus of Inheritance}{一种继承的演算}}
% The appendix-only build (supplement.tex) is distributed as a standalone
% artifact; mark it so it is not mistaken for the main paper, which carries
% no subtitle.
\ifInheritanceBody\else
  \subtitle{Supplementary Material}
\fi

\author{Bo Yang}
\affiliation{%
  \institution{Figure AI Inc.}
  \city{San Jose}
  \state{California}
  \country{USA}
}
\email{yang-bo@yang-bo.com}
\thanks{This work was conducted independently prior to the author's employment at Figure AI.}

% The abstract, CCS concepts, and keywords are main-matter front matter: they
% belong to the paper, not to the standalone appendix PDF (supplement.tex),
% which sets \ifInheritanceBody false. Without this guard \maketitle typesets
% them in the appendix-only build.
\ifInheritanceBody
\begin{abstract}
  Just as the $\lambda$-calculus uses three primitives (abstraction, application, variable) as the foundation of functional programming, inheritance-calculus uses three primitives (record, definition, inheritance) as the foundation of declarative programming. A record is a set of definitions; by unifying modules, classes, objects, methods, fields, and locals under this single record abstraction, the calculus models inheritance simply as set union of definitions, two definitions of the same name merging recursively rather than one overriding the other. A merge is always meaningful because every value is a record and observation is purely structural: the merged record simply answers queries with the union of what its parts provide, so there is no scalar leaf at which two values could conflict. Consequently, composition is inherently commutative, idempotent, and associative, structurally eliminating the multiple-inheritance linearization problem. Its semantics is first-order, denotational, and computable by tabling (memoized fixpoint iteration as in tabled logic programming), even for cyclic inheritance hierarchies. This first-order, denotational, tabled semantics extends to the $\lambda$-calculus: $\lambda$-terms translate into a sublanguage of a lazy approximation of inheritance-calculus, the approximation that treats cyclic self-reference as divergent instead of iterating it to a fixpoint, and the translation is fully abstract, equating exactly the $\lambda$-terms with equal B\"ohm trees.

  Inheritance-calculus is distilled from MIXINv2, a practical implementation in which the same code can be invoked under different calling conventions (function colors such as synchronous and asynchronous); ordinary arithmetic, encoded within the calculus itself, behaves as multi-directional relations in the sense of logic programming; $\mathtt{this}$ resolves to multiple enclosing targets at once rather than a single receiver; and programs are extensible in both dimensions of the Expression Problem, gaining new variants and new operations by adding code. This makes inheritance-calculus strictly more expressive than the $\lambda$-calculus in Felleisen's sense of macro-expressibility.
\end{abstract}

\begin{CCSXML}
  <ccs2012>
  <concept>
  <concept_id>10003752.10010124.10010131</concept_id>
  <concept_desc>Theory of computation~Program semantics</concept_desc>
  <concept_significance>500</concept_significance>
  </concept>
  <concept>
  <concept_id>10003752.10010124.10010125.10010128</concept_id>
  <concept_desc>Theory of computation~Object oriented constructs</concept_desc>
  <concept_significance>500</concept_significance>
  </concept>
  <concept>
  <concept_id>10003752.10010124.10010125.10010127</concept_id>
  <concept_desc>Theory of computation~Functional constructs</concept_desc>
  <concept_significance>300</concept_significance>
  </concept>
  <concept>
  <concept_id>10003752.10010124.10010125.10010129</concept_id>
  <concept_desc>Theory of computation~Program schemes</concept_desc>
  <concept_significance>100</concept_significance>
  </concept>
  <concept>
  <concept_id>10011007.10011006.10011039</concept_id>
  <concept_desc>Software and its engineering~Formal language definitions</concept_desc>
  <concept_significance>500</concept_significance>
  </concept>
  <concept>
  <concept_id>10011007.10011006.10011008.10011009.10011019</concept_id>
  <concept_desc>Software and its engineering~Extensible languages</concept_desc>
  <concept_significance>500</concept_significance>
  </concept>
  <concept>
  <concept_id>10011007.10011006.10011008.10011009.10011011</concept_id>
  <concept_desc>Software and its engineering~Object oriented languages</concept_desc>
  <concept_significance>100</concept_significance>
  </concept>
  </ccs2012>
\end{CCSXML}

\ccsdesc[500]{Theory of computation~Program semantics}
\ccsdesc[500]{Theory of computation~Object oriented constructs}
\ccsdesc[300]{Theory of computation~Functional constructs}
\ccsdesc[100]{Theory of computation~Program schemes}
\ccsdesc[500]{Software and its engineering~Formal language definitions}
\ccsdesc[500]{Software and its engineering~Extensible languages}
\ccsdesc[100]{Software and its engineering~Object oriented languages}

\keywords{declarative programming, inheritance, fixpoint semantics,
  self-referential records, expression problem,
  first-order semantics,
  \texorpdfstring{$\lambda$-calculus}{lambda-calculus},
\texorpdfstring{B\"ohm trees}{Bohm trees}, configuration languages}
\fi % end \ifInheritanceBody (abstract, CCS concepts, keywords are main-matter only)

\maketitle

\ifInheritanceBody
\section{\bilingual{Introduction}{引言}}
\label{sec:introduction}

\ifInheritanceChinese
诸如 Jsonnet~\cite{jsonnet}、Hydra~\cite{hydra2023}、CUE~\cite{cue2019}、Dhall~\cite{dhall2017}、Kustomize~\cite{kustomize} 与 JSON Patch~\cite{rfc6902-json-patch} 这样的配置语言,都提供了通过继承或合并来组合结构化数据的机制。

在这些之中,Nix 生态系统~\cite{dolstra2010-nixos} 格外突出,无论是就其规模还是其设计的走向而言。它的软件包集合 \texttt{nixpkgs} 收录了超过 113{,}000 个软件包,%
\footnote{
  据 Repology~\cite{repology2025},这接近 Debian 或 Ubuntu 软件包数量的三倍,而贡献者数量相近。并非每个软件包都用到这里所描述的可扩展性机制,但同样的告诫也适用于其他发行版。
}
是现存最大的软件包集合之一。软件包的组合最初依赖一条线性化的 mixin 链,其中每一层都能看见它的前一层,%
\footnote{
  Nix overlay 通过一条线性化的 \texttt{self}/\texttt{super} 链来组合软件包集合,遵循 Bracha--Cook 的 mixin 模式~\cite{bracha1990-mixin-inheritance}:\texttt{self} 晚绑定到不动点结果,而 \texttt{super} 指向前一层。该机制是不对称且依赖顺序的。
}
但该生态系统已日益汇聚到一种对称的替代方案上:NixOS 模块系统~\cite{dolstra2010-nixos,nixosmodules},它通过对嵌套的属性集进行递归合并、配以自引用求值与延迟模块~\cite{nixosmodules} 来组合,而这种组合是可交换、幂等且可结合的。模块系统最初为系统配置而设计,如今已被用于用户环境~\cite{homemanager}、macOS 配置~\cite{nixdarwin}、磁盘分区~\cite{disko}、flake 结构~\cite{flakeparts}、开发者环境~\cite{devenv}、Kubernetes 集群~\cite{kubenix,nixidy},以及多语言软件包构建~\cite{dream2nix},而最后这一项正是 \texttt{nixpkgs} 那套不对称机制原本的领域。Nix 是一门带头等函数的函数式语言;然而一份典型的模块配置几乎完全由嵌套的属性集构成,函数只是偶尔作为模块系统的 DSL 语法出现,例如 \texttt{mkIf} 与 \texttt{mkMerge}。没有任何计算理论刻画了这一机制所展现的可扩展性。
\else
Configuration languages such as Jsonnet~\cite{jsonnet}, Hydra~\cite{hydra2023},
CUE~\cite{cue2019}, Dhall~\cite{dhall2017},
Kustomize~\cite{kustomize}, and JSON
Patch~\cite{rfc6902-json-patch} all provide mechanisms for
composing structured data through inheritance or merging.

Among these, the Nix ecosystem~\cite{dolstra2010-nixos} stands out,
both for scale and for the trajectory of its design.
Its package collection, \texttt{nixpkgs}, contains over
113{,}000 packages,%
\footnote{
  According to Repology~\cite{repology2025},
  nearly three times the package count of Debian or Ubuntu,
  achieved with a similar number of contributors.
  Not every package exercises the extensibility mechanisms
  described here, but the same caveat applies to other distributions.
}
one of the largest in existence.
Package composition originally relied on a linearized mixin
chain in which each layer sees its predecessor,%
\footnote{
  Nix overlays compose package sets via a linearized
  \texttt{self}/\texttt{super} chain following the
  Bracha--Cook mixin pattern~\cite{bracha1990-mixin-inheritance}:
  \texttt{self} is late-bound to the fixed-point result
  and \texttt{super} refers to the previous layer.
  The mechanism is asymmetric and order-dependent.
}
but the ecosystem has increasingly converged on a symmetric
alternative: the NixOS module
system~\cite{dolstra2010-nixos,nixosmodules}, which composes by
recursively merging nested attribute sets with self-referential
evaluation and deferred modules~\cite{nixosmodules}---commutative,
idempotent,
and associative.
Originally designed for system configuration, the module system
has been adopted for user environments~\cite{homemanager},
macOS configuration~\cite{nixdarwin},
disk partitioning~\cite{disko},
flake structure~\cite{flakeparts},
developer environments~\cite{devenv},
Kubernetes clusters~\cite{kubenix,nixidy},
and multi-language package building~\cite{dream2nix}---the
original domain of \texttt{nixpkgs}'s asymmetric mechanism.
Nix is a functional language with first-class functions;
yet a typical module configuration consists almost entirely of
nested attribute sets, with functions appearing only
occasionally as DSL syntax of the module system
such as \texttt{mkIf} and \texttt{mkMerge}.
No computational theory captures the extensibility that this
mechanism exhibits.
\fi

The $\lambda$-calculus has no analogue for declarative
programming. This gap matters because the two paradigms differ in
fundamental ways, which raises a question:

\begin{quote}
  Can configuration alone be practical general-purpose programming?
\end{quote}

\noindent
To make the question precise, we operationalize it with four
\emph{opinionated} proxies, grouped under the question's two halves:
\emph{configuration} and \emph{practical general-purpose} programming.

\begin{itemize}
  \item Immutable and first-order%
    \footnote{
      Here \emph{first-order} holds in three senses at once:
      the semantic domain contains only data, no function spaces,
      closures, or
      strategies~\cite{vanemden1976-predicate-logic-semantics};
      the defining equations map data to
      data~\cite{reynolds1972-definitional-interpreters};
      and the semantics is a first-order inductive
      definition~\cite{aczel1977-inductive-definitions} whose
      quantifiers range only over data.
    }
    is our proxy for
    \emph{configuration}: a configuration is read, not a function that
    reads an argument.
  \item Turing-completeness together with extensibility is our proxy for
    \emph{practical general-purpose} programming: a language must be able
    to compute anything and to grow by adding code rather than rewriting
    it, the Expression
    Problem~\cite{wadler1998-expression-problem}.
\end{itemize}

\noindent
Each classical model satisfies some of these proxies but fails at least
one:

\begin{table}[h]
  \centering
  \small
  \begin{tabular}{lcccl}
    \textbf{Model} & \textbf{Immutable} & \textbf{First-order} & \textbf{Turing complete} & \textbf{Extensibility} \\
    \hline
    Turing Machine        & $\times$     & $\checkmark$ & $\checkmark$ & $\times$ \\
    $\lambda$-calculus    & $\checkmark$ & $\times$     & $\checkmark$ & opt-in \\
    RAM Machine           & $\times$     & $\checkmark$ & $\checkmark$ & $\times$ \\
    Prolog / Horn clauses & $\checkmark$ & $\checkmark$ & $\checkmark$ & disjunct-only \\
  \end{tabular}
  \caption{The question operationalized as four proxies: immutable and
    first-order (configuration) versus Turing-complete and extensible
  (practical general-purpose). No classical model satisfies all four.}
  \label{tab:computational-models}
\end{table}

\noindent
The Turing Machine and RAM Machine require mutable state and have no native
extension mechanism;
the $\lambda$-calculus is higher-order rather than first-order, and its
extensibility is opt-in: a function is extensible only where its author
anticipated extension, for example by threading open-recursion
parameters, while a closed function body cannot be extended without
rewriting it;
Prolog meets the other three proxies but is extensible only by adding
disjuncts (alternative clauses) to existing predicates, not by attaching
new fields or methods to existing values.
No known computational model satisfies all four proxies simultaneously.

The Nix ecosystem already suggests that such a model may exist:
its inheritance-based composition over recursive records, without
explicit functions, is expressive enough to configure entire
operating systems and manage one of the largest package
collections in existence.
Most configuration languages operate on finite, well-founded
structures (initial algebras in the sense of universal algebra).
Nix is an exception: its values are lazily observable
and possibly infinite, allowing any single package or
configuration option to be evaluated without materializing the
entire set of hundreds of thousands of packages and options.
Guided by this observation, we set out to reduce the
NixOS module system to a minimal set of primitives.
The reduction produced MIXIN, an early implementation in Nix,
and subsequently
MIXINv2\anon[~(included as supplementary material)]{~\cite{mixin2025}}, now implemented in Python: a declarative programming
language with only three constructs (record, definition,
and inheritance) and no functions or let-bindings.
MIXINv2 obtains scalars such as floating-point numbers through a
foreign-function interface; inheritance-calculus is what remains after
removing that interface, so names like $\mathrm{PI}$ and
$\mathrm{Float}$ in Figure~\ref{fig:cheatsheet} denote imported
records whose internals the examples never inspect: scalar
computation such as the figure's additions happens behind the
interface, while the calculus itself only ever observes structure;
the case study of Section~\ref{sec:case-study} instead encodes
naturals and their arithmetic purely in the calculus, with no FFI.
When two merged records define the same name with different contents,
the result is simply a record carrying both structures; the calculus
never compares or rejects them
(Section~\ref{sec:mixin-trees} makes the merge precise).
No such collision occurs in Figure~\ref{fig:cheatsheet}: each
instance defines each scalar field exactly once, and the type record
it merges with contributes structure, not a competing value.

\begin{figure*}
  \centering
  \begin{minipage}[t]{0.34\textwidth}
    \textbf{Python}\par
\begin{lstlisting}[style=python]
# complex_module.py
from builtins import *
from math import *
from operator import *
from dataclasses import *
@dataclass(kw_only=True)
class Complex:
    re: float
    im: float
    def plus(self, *,
             other: "Complex"):
        sumRe = add(
            self.re,
            other.re)
        sumIm = add(
            self.im,
            other.im)
        return Complex(
            re=sumRe,
            im=sumIm)
a = Complex(
    re=pi,
    im=e)
b = Complex(
    re=tau,
    im=pi)
total = a.plus(
    other=b)
\end{lstlisting}
  \end{minipage}
  \hfill
  \begin{minipage}[t]{0.65\textwidth}
    \textbf{Inheritance-calculus}\par\smallskip
    {\small\(
        \figureAnnotationAnchor
        \begin{aligned}
          &\mathrm{ComplexModule} \mapsto {} \figureAnnotation{module} \\
          &\left\{
            \begin{aligned}
              &\mathrm{BuiltIns}, \figureAnnotation{import} \\
              &\mathrm{Math}, \figureAnnotation{import} \\
              &\mathrm{Operator}, \figureAnnotation{import} \\
              &\mathrm{Complex} \mapsto {} \figureAnnotation{class} \\
              &\left\{
                \begin{aligned}
                  &\mathrm{re} \mapsto {} \figureAnnotation{field} \\
                  &\left\{ \qualifiedthis{\mathrm{ComplexModule}}.\mathrm{Float} \right\}, \figureAnnotation{imported type} \\
                  &\mathrm{im} \mapsto {} \figureAnnotation{field} \\
                  &\left\{ \qualifiedthis{\mathrm{ComplexModule}}.\mathrm{Float} \right\}, \figureAnnotation{imported type} \\
                  &\mathrm{Plus} \mapsto {} \figureAnnotation{method} \\
                  &\left\{
                    \begin{aligned}
                      &\mathrm{other} \mapsto \left\{ \mathrm{Complex} \right\}, \figureAnnotation{parameter} \\
                      &\mathrm{\_sumReCall} \mapsto {} \figureAnnotation{function application} \\
                      &\left\{
                        \begin{aligned}
                          &\qualifiedthis{\mathrm{ComplexModule}}.\mathrm{Add}, \figureAnnotation{imported function} \\
                          &\mathrm{operand0} \mapsto \left\{ \qualifiedthis{\mathrm{Complex}}.\mathrm{re} \right\}, \figureAnnotation{keyword argument} \\
                          &\mathrm{operand1} \mapsto \left\{ \qualifiedthis{\mathrm{Plus}}.\mathrm{other}.\mathrm{re} \right\} \figureAnnotation{keyword argument} \\
                      \end{aligned}\right\}, \\
                      &\mathrm{\_sumRe} \mapsto \left\{ \mathrm{\_sumReCall}.\mathrm{result} \right\}, \figureAnnotation{local} \\
                      &\mathrm{\_sumImCall} \mapsto {} \figureAnnotation{function application} \\
                      &\left\{
                        \begin{aligned}
                          &\qualifiedthis{\mathrm{ComplexModule}}.\mathrm{Add}, \figureAnnotation{imported function} \\
                          &\mathrm{operand0} \mapsto \left\{ \qualifiedthis{\mathrm{Complex}}.\mathrm{im} \right\}, \figureAnnotation{keyword argument} \\
                          &\mathrm{operand1} \mapsto \left\{ \qualifiedthis{\mathrm{Plus}}.\mathrm{other}.\mathrm{im} \right\} \figureAnnotation{keyword argument} \\
                      \end{aligned}\right\}, \\
                      &\mathrm{\_sumIm} \mapsto \left\{ \mathrm{\_sumImCall}.\mathrm{result} \right\}, \figureAnnotation{local} \\
                      &\mathrm{result} \mapsto {} \figureAnnotation{return} \\
                      &\left\{
                        \begin{aligned}
                          &\mathrm{Complex}, \figureAnnotation{new instance} \\
                          &\mathrm{re} \mapsto \left\{ \mathrm{\_sumRe} \right\}, \figureAnnotation{keyword argument} \\
                          &\mathrm{im} \mapsto \left\{ \mathrm{\_sumIm} \right\} \figureAnnotation{keyword argument} \\
                      \end{aligned}\right\} \\
                  \end{aligned}\right\} \\
              \end{aligned}\right\}, \\
              &\mathrm{a} \mapsto {} \figureAnnotation{module-level variable} \\
              &\left\{
                \begin{aligned}
                  &\mathrm{Complex}, \figureAnnotation{new instance} \\
                  &\mathrm{re} \mapsto \left\{ \qualifiedthis{\mathrm{ComplexModule}}.\mathrm{PI} \right\}, \figureAnnotation{imported constant} \\
                  &\mathrm{im} \mapsto \left\{ \qualifiedthis{\mathrm{ComplexModule}}.\mathrm{E} \right\} \figureAnnotation{imported constant} \\
              \end{aligned}\right\}, \\
              &\mathrm{b} \mapsto {} \figureAnnotation{module-level variable} \\
              &\left\{
                \begin{aligned}
                  &\mathrm{Complex}, \figureAnnotation{new instance} \\
                  &\mathrm{re} \mapsto \left\{ \qualifiedthis{\mathrm{ComplexModule}}.\mathrm{Tau} \right\}, \figureAnnotation{imported constant} \\
                  &\mathrm{im} \mapsto \left\{ \qualifiedthis{\mathrm{ComplexModule}}.\mathrm{PI} \right\} \figureAnnotation{imported constant} \\
              \end{aligned}\right\}, \\
              &\mathrm{\_totalCall} \mapsto {} \figureAnnotation{method application} \\
              &\left\{
                \begin{aligned}
                  &\mathrm{a}.\mathrm{Plus}, \figureAnnotation{method reference} \\
                  &\mathrm{other} \mapsto \left\{ \mathrm{b} \right\} \figureAnnotation{keyword argument} \\
              \end{aligned}\right\}, \\
              &\mathrm{total} \mapsto \left\{ \mathrm{\_totalCall}.\mathrm{result} \right\} \figureAnnotation{module-level variable} \\
          \end{aligned}\right\}
        \end{aligned}
    \)}
  \end{minipage}
  \caption{A first look at inheritance-calculus (right), beside Python
    (left): the same complex-number class with an addition method.
    An entry $name \mapsto \{\dots\}$ is a definition; a bare name
    inside braces (such as $\mathrm{BuiltIns}$ or $\mathrm{Complex}$) is
    an inheritance clause that copies that record's definitions into the
    enclosing record.
    $\mathrm{BuiltIns}$, $\mathrm{Math}$, and $\mathrm{Operator}$ are
    defined in an enclosing scope, elided here.
    A bare name resolves lexically to the innermost enclosing record
    that defines it directly; qualified $\this$ reaches members that
    arrive by inheritance, such as the imported $\mathrm{Float}$ and
    $\mathrm{Add}$.
    Types and values share the same record syntax: a field such as
    $\mathrm{re}$ records its type by inheriting the type's record,
    with no checking force in the untyped calculus.
    Both sides, typing their components with \texttt{builtins.float},
    drawing their constants from the \texttt{math} module, and adding with
    \texttt{operator.add}, compute the complex
    number $(\pi+\tau) + (e+\pi)i$.
    $\mathrm{\_sumReCall}$ and $\mathrm{\_sumImCall}$ are function
    applications, and $\mathrm{\_totalCall}$ a method application, each
    encoded as a command pattern: a record that inherits the callee and
    defines its arguments, the call's result read back from a
    $\mathrm{result}$ field that the callee defines.
    Idempotence only means that listing the same inheritance source
    twice in one record is a no-op; distinct records inheriting the
    same source, like the two call sites here or the two instances
    $\mathrm{a}$ and $\mathrm{b}$, remain distinct.}
  \label{fig:cheatsheet}
\end{figure*}



In developing real programs in MIXINv2, including a database-backed web server whose Python foreign-function interface only faithfully wraps individual host-library calls, we found that module, class, object, method, field, and local binding can all be expressed as the same construct: a record of definitions composed by inheritance. This paper distills this minimal computational model into \textbf{inheritance-calculus}, shown beside Python in Figure~\ref{fig:cheatsheet}.
The key to reading the figure is that a reference
$\qualifiedthis{X}.w$ behaves like a late-bound Java-style
\texttt{Outer.this}: it denotes the enclosing record labeled $X$ as
seen from wherever the reference itself ends up in the fully merged
tree.
Inheriting a record does not duplicate it; it makes the record's
definitions, references included, visible at the inheriting position
as well, and a reference resolves relative to the position from which
it is viewed.
The set-union picture and this late-binding picture are the same
mechanism seen globally and locally: the union says which definitions
a position carries, and late binding says how a definition behaves at
each position that carries it.
Qualified references therefore see members that arrive by merging,
such as the imported $\mathrm{Float}$, while bare names such as
$\mathrm{Complex}$ resolve lexically against directly written
definitions only.
This late binding is what makes inheritance act as function
application: the call record $\mathrm{\_totalCall}$ inherits
$\mathrm{Plus}$, so the merged tree presents
$\mathrm{Plus}$'s body inside $\mathrm{\_totalCall}$ as well, and
viewed from there the reference
$\qualifiedthis{\mathrm{Plus}}.\mathrm{other}$
denotes that copy's enclosing record, $\mathrm{\_totalCall}$ itself,
where the argument $\mathrm{other}$ is defined.
Receivers bind the same way: the dotted bare name
$\mathrm{a}.\mathrm{Plus}$ resolves its first label lexically and
projects the rest, denoting the copy of $\mathrm{Plus}$ inside the
instance $\mathrm{a}$, in which
$\qualifiedthis{\mathrm{Complex}}.\mathrm{re}$ denotes
$\mathrm{a}$'s own field.
Section~\ref{sec:syntax} gives its formal syntax, and
Section~\ref{sec:mixin-trees} expands this first look into the full
semantics.
The linearization problem of mixin-based
systems~\cite{bracha1990-mixin-inheritance,barrett1996-c3-linearization}
did not arise: inheritance is inherently commutative, idempotent,
and associative, and $\this$ naturally resolves to
\emph{multiple targets} rather than one: when inheritance places the
same definition inside several enclosing records, a reference through
$\this$ inherits from all of them at once, their contents merged,%
\footnote{
  Resolution is anchored at the position from which the reference is
  observed. When both $\mathrm{A}$ and $\mathrm{B}$ inherit
  $\mathrm{Base}$, a reference in $\mathrm{Base}$ observed inside
  $\mathrm{A}$ sees only $\mathrm{A}$; the two call sites of
  Figure~\ref{fig:cheatsheet} therefore stay independent. Multiple
  targets arise when a single observation position inherits the same
  definition through two routes, say $\mathrm{C}$ inheriting
  $\mathrm{Base}$ via both $\mathrm{A}$ and $\mathrm{B}$: observed at
  $\mathrm{C}$, a $\this$ reference in $\mathrm{Base}$ resolves to
  $\mathrm{A}$ and $\mathrm{B}$ at once, and projection through it
  yields the union of what they provide.
}
where
prior systems
such as Scala and the NixOS module system reject the multi-target
situation.
Section~\ref{sec:case-study} traces how boolean logic and
arithmetic, encoded purely in the calculus as separate files, extend in both
dimensions of the Expression
Problem~\cite{wadler1998-expression-problem} without any dedicated
mechanism, and how ordinary arithmetic yields the relational
semantics of logic
programming~\cite{vanemden1976-predicate-logic-semantics}.
Section~\ref{sec:emergent-phenomena} discusses these and further
emergent phenomena, including function color
blindness~\cite{nystrom2015-function-color}.
Together these sections answer the four proxies: inheritance-calculus
is immutable and first-order by construction
(Section~\ref{sec:definitions}), Turing-complete because the
$\lambda$-calculus embeds into it
(Section~\ref{sec:forward-translation}), and extensible in both
dimensions of the Expression Problem (Section~\ref{sec:case-study}).

\begin{figure*}
  \centering
  \begin{tikzpicture}[
      calc/.style={align=center, inner sep=2pt, font=\small},
      more/.style={-{Stealth[length=2mm]}, semithick},
      iso/.style={{Stealth[length=2mm]}-{Stealth[length=2mm]}, semithick},
      elbl/.style={font=\scriptsize\itshape, midway, fill=white, inner sep=1pt},
      x=1mm, y=1mm,
    ]
    % rows: bottom = lambda calculi, middle = lambda sublanguages, top = inheritance-calculi
    % columns: eager (x=0), lazy (x=42), fixpoint (x=84)
    \node[calc] (lam)  at (0,0)   {$\lambda$};
    \node[calc] (llam) at (42,0)  {lazy $\lambda$};
    \node[calc] (flam) at (84,0)  {fixpoint $\lambda$~\cite{yang2026-co-lambda}};

    \node[calc] (lsub) at (42,19)  {$\lambda$-lazy-IC (ours)};
    \node[calc] (sub)  at (84,19)  {$\lambda$-IC};

    \node[calc] (lic)  at (42,38)  {lazy IC (ours)};
    \node[calc] (ic)   at (84,38)  {IC (ours)};

    % horizontal edges (convergence): "more defined"
    \draw[more] (lam)  -- node[elbl,below]{more defined (non-strict)} (llam);
    \draw[more] (llam) -- node[elbl,below]{more defined (unproductive $\bot$)} (flam);
    \draw[more] (lsub) -- node[elbl,above]{more defined (powerset)} (sub);
    \draw[more] (lic)  -- node[elbl,above]{more defined (powerset)} (ic);

    % vertical edges (expressiveness): "strictly more expressive"
    \draw[more] (lsub) -- node[elbl,left]{strictly more expressive} (lic);
    \draw[more] (sub)  -- node[elbl,right]{strictly more expressive} (ic);

    % full abstraction: vertical, in the lazy column
    \draw[iso] (llam) -- node[elbl,left]{fully abstract} (lsub);
  \end{tikzpicture}
  \caption{The calculi of this paper; nodes marked
    (ours) are contributions of this paper.
    Abbreviations:
    $\lambda$ is the $\lambda$-calculus;
    the lazy prefix denotes the approximation that stops after the first
    iteration of the fixpoint of this paper's semantic equations
    (written $T_P\!\uparrow\!1$ in
    logic-programming notation): lookups that re-enter themselves
    through cyclic self-reference return divergence instead of being
    iterated further.
    Continuing:
    fixpoint $\lambda$ is the tabled evaluation of the
    $\lambda$-calculus of Yang~\cite{yang2026-co-lambda};
    IC the inheritance-calculus;
    and $\lambda$-IC and $\lambda$-lazy-IC the $\lambda$ sublanguages of
    IC and of lazy IC, respectively.
    Edge labels:
    \emph{more defined} means every term that converges in the source
    calculus also converges in the target calculus, but not conversely.
    The parenthesized qualifiers distinguish the mechanism:
    \emph{non-strict} more-definedness gains convergence by deeming
    abstractions values without evaluating their bodies;
    \emph{unproductive $\bot$} more-definedness gains convergence by
    deciding unproductive loops such as $\Omega$ as the definite answer
    $\bot$ in finite time, preserving the lazy tree semantics;
    and \emph{powerset} more-definedness gains convergence by
    resolving cyclic self-references to definite sets, as least fixed
    points over the powerset lattice.
    Continuing the edge labels:
    \emph{strictly more expressive} is meant in Felleisen's sense of
    macro-expressibility;
    and \emph{fully abstract} marks full abstraction.}
  \label{fig:calculi-matrix}
\end{figure*}

Figure~\ref{fig:calculi-matrix} situates inheritance-calculus among
several related calculi developed below.
Section~\ref{sec:definitions} defines the semantics via six
mutually recursive first-order functions on paths and sets of paths,
with no function spaces, closures, or higher-order constructs;
the functions are computable by
tabling~\cite{tamaki1986-tabled-resolution}.
Section~\ref{sec:forward-translation} shows that
these functions extend to the $\lambda$-calculus, embedding it.
Section~\ref{sec:asymmetry} establishes that the lazy
$\lambda$-calculus cannot macro-express the inheritance constructors,
so inheritance-calculus is strictly more expressive in
Felleisen's sense~\cite{felleisen1991-expressive-power}.

\section{\bilingual{Syntax}{语法}}
\label{sec:syntax}

\ifInheritanceChinese
设 $e$ 表示一个表达式,$s$ 表示一个元素列表,$c$ 表示一个元素,
$q$ 表示一个限定符,$w$ 表示一个投影列表,$\ell$ 表示一个充当
属性名的标签,$n$ 表示一个 de~Bruijn 索引。
空串 $\varepsilon$ 是空元素列表。
\else
Let $e$ denote an expression, $s$ an element list, $c$ an element,
$q$ a qualifier, $w$ a projection list, $\ell$ a label serving as a
property name, and $n$ a de~Bruijn index.
The empty string $\varepsilon$ is the empty element list.
\fi

\begin{align*}
  e \quad ::= \quad & \{\, s \,\}
  && \text{\bilingual{(record)}{(记录)}} \\[6pt]
  s \quad ::= \quad & \varepsilon
  && \text{\bilingual{(empty)}{(空)}} \\*
  \mid\quad & c
  && \text{\bilingual{(one element)}{(一个元素)}} \\*
  \mid\quad & c,\; s
  && \text{\bilingual{(element, then more)}{(一个元素,然后更多)}} \\[6pt]
  c \quad ::= \quad & \ell \mapsto e
  && \text{\bilingual{(definition)}{(定义)}} \\*
  \mid\quad & q.\, w
  && \text{\bilingual{(inheritance: qualifier and projections)}{(继承:限定符与投影)}} \\*
  \mid\quad & q
  && \text{\bilingual{(inheritance: qualifier only)}{(继承:仅限定符)}} \\*
  \mid\quad & w
  && \text{\bilingual{(inheritance: projections only)}{(继承:仅投影)}} \\[6pt]
  q \quad ::= \quad & \qualifiedthis{\ell_{\mathrm{qualifier}}}
  && \text{\bilingual{(named qualifier)}{(具名限定符)}} \\*
  \mid\quad & \dbi{n}
  && \text{\bilingual{(indexed qualifier)}{(索引限定符)}} \\[6pt]
  w \quad ::= \quad & \ell_{\mathrm{projection}}
  && \text{\bilingual{(one projection)}{(一个投影)}} \\*
  \mid\quad & \ell_{\mathrm{projection}}.\, w
  && \text{\bilingual{(projection, then more)}{(一个投影,然后更多)}}
\end{align*}

\ifInheritanceChinese
一个表达式将一个元素列表 $s$ 包裹在花括号中,而每个元素
$c$ 要么是一个定义 $\ell \mapsto e$,要么是一个继承源。
语法把 $s$ 写成一个递归列表,但该列表被读作一个
\emph{集合}:元素的顺序无关紧要,重复元素也没有任何作用,
因此两个仅在重排或重复上有所不同的元素列表表示
同一个表达式。由于表达式是集合,继承本身就是
可交换、幂等且可结合的。
图~\ref{fig:cheatsheet} 中的记录 $\mathrm{Complex} \mapsto \{ \mathrm{re} \mapsto \ldots,\;
\mathrm{im} \mapsto \ldots,\; \mathrm{Plus} \mapsto \ldots \}$
就是一个这样的多元素列表。

$\ell \mapsto e$ 中的 $\mapsto$ 定义了一个名为 $\ell$、其值为
$e$ 的属性;例如,图~\ref{fig:cheatsheet} 中的 $\mathrm{re} \mapsto \{\ldots\}$ 与
$\mathrm{Complex} \mapsto \{\ldots\}$。
MIXINv2 还通过一个
外部函数接口(FFI)提供标量类型;继承演算就是
去掉 FFI 之后所剩下的部分。
本文中没有任何例子用到 FFI 或任何标量值。

\paragraph{\bilingual{Inheritance}{继承}}
一个继承通过一个限定符 $q$ 引用一个作用域;该限定符
沿作用域层级向上导航,然后通过一个投影列表 $w$
向下投影属性。限定符 $q$ 容许两种等价的
记法:
\begin{description}
  \item[\bilingual{Named form}{具名形式}]
    $\qualifiedthis{\ell_{\mathrm{qualifier}}}.\, w$,
    类比于 Java 的 \texttt{Outer.this.field},其中 $w$ 是一个
    投影列表。
    每个 $\{\ldots\}$ 记录在它的外围作用域中都有一个标签;
    $\ell_{\mathrm{qualifier}}$ 命名目标外围作用域。
    例如,图~\ref{fig:cheatsheet} 中的
    $\qualifiedthis{\mathrm{ComplexModule}}.\mathrm{Float}$
    投影单个标签,而
    $\qualifiedthis{\mathrm{Plus}}.\mathrm{other}.\mathrm{re}$ 展示了一个
    多标签的投影列表 $w$。
  \item[\bilingual{Indexed form}{索引形式}]
    $\dbi{n}.\, w$,
    其中 $n \ge 0$ 是一个 de~Bruijn 索引~\cite{debruijn1972-nameless-dummies}:
    $n = 0$ 指向该引用的外围作用域,
    而 $n$ 每加一,就向外移动一个作用域层级。
\end{description}
两种形式都先取出目标作用域,即所有继承都应用完毕之后
那个完全继承后的记录,然后从中投影由投影列表 $w$
所命名的属性。
具名形式在解析期间脱糖为索引形式,
如第~\ref{sec:mixin-trees} 节所述。

每个 $\{\ldots\}$ 记录字面量都会创建一个新的作用域层级。%
\footnote{
  我们用\emph{记录}指那个语法构造,用\emph{作用域}指
  它所创建的名字解析语境,用\emph{mixin} 指它
  作为可组合单元的角色;在继承演算中这些是同一个实体的三种
  视角。
}
该作用域包含那个 mixin 的所有属性,包括那些
通过继承源继承而来的属性。
例如,$\{a \mapsto \{\},\; b \mapsto \{\}\}$ 是单个 mixin,其作用域
同时包含 $a$ 与 $b$;$\{r,\; a \mapsto \{\}\}$ 也是如此,其中 $r$
是对一个包含 $b$ 的记录的引用。
继承源不会创建额外的作用域层级。
在图~\ref{fig:cheatsheet} 中,绑定到 $\mathrm{Plus}$ 的记录
就是一个这样的作用域,它包含 $\mathrm{other}$、$\mathrm{\_sumReCall}$、
$\mathrm{\_sumRe}$ 以及在其内部直接定义的其他属性。

\begin{itemize}
  \item $\qualifiedthis{\mathrm{Foo}}$(等价地,$\dbi{n}$,其中
    $\mathrm{Foo}$ 在外围作用域之上 $n$ 步)
    取出名为 $\mathrm{Foo}$ 的记录。
  \item $\qualifiedthis{\mathrm{Foo}}.\ell$,或等价地,$\dbi{n}.\ell$,
    从 $\mathrm{Foo}$ 投影属性 $\ell$。
  \item 一个仅限定符的源 $q$(没有投影列表)继承
    整个外围作用域,这可能产生一棵无限深的
    树;图~\ref{fig:cheatsheet} 中裸的 $\mathrm{BuiltIns}$、$\mathrm{Math}$ 与
    $\mathrm{Operator}$ 源就是此类
    仅限定符的导入。
\end{itemize}

当限定符无歧义时,它可以省略:这种简写是
一个裸的投影列表 $w$。它的首个投影针对
那个在其记录字面量中直接定义首个投影标签的
最内层外围作用域来解析,而该源脱糖为
对应的索引形式。
例如,图~\ref{fig:cheatsheet} 中的 $\mathrm{\_sumReCall}.\mathrm{result}$、
$\mathrm{a}.\mathrm{Plus}$ 与 $\mathrm{\_totalCall}.\mathrm{result}$
就是此类省略限定符的简写。
由于这一搜索只检查记录字面量的定义,该
简写只能访问在记录字面量中定义的属性,而不能访问
那些通过继承而继承来的属性;对于继承来的属性或要绕过变量遮蔽,
都需要两种限定形式之一。
\else
An expression encloses an element list $s$ in braces, and each element
$c$ is either a definition $\ell \mapsto e$ or an inheritance source.
The grammar writes $s$ as a recursive list, but this list is read as a
\emph{set}: element order is irrelevant and duplicates have no effect,
so two element lists differing only by reordering or repetition denote
the same expression. Because expressions are sets, inheritance is
inherently commutative, idempotent, and associative.
The record $\mathrm{Complex} \mapsto \{ \mathrm{re} \mapsto \ldots,\;
\mathrm{im} \mapsto \ldots,\; \mathrm{Plus} \mapsto \ldots \}$ in
Figure~\ref{fig:cheatsheet} is one such multi-element list.

The $\mapsto$ in $\ell \mapsto e$ defines a property named $\ell$ with
value $e$; for example, $\mathrm{re} \mapsto \{\ldots\}$ and
$\mathrm{Complex} \mapsto \{\ldots\}$ in Figure~\ref{fig:cheatsheet}.
MIXINv2 also provides scalar types through a
foreign-function interface (FFI); inheritance-calculus is what remains
after removing the FFI.
None of the examples in this paper uses the FFI or any scalar values.

\paragraph{Inheritance}
An inheritance references a scope through a qualifier $q$ that
navigates the scope hierarchy upward, then projects properties downward
through a projection list $w$. The qualifier $q$ admits two equivalent
notations:
\begin{description}
  \item[Named form]
    $\qualifiedthis{\ell_{\mathrm{qualifier}}}.\, w$,
    analogous to Java's \texttt{Outer.this.field}, where $w$ is a
    projection list.
    Each $\{\ldots\}$ record has a label in its enclosing scope;
    $\ell_{\mathrm{qualifier}}$ names the target enclosing scope.
    For example, $\qualifiedthis{\mathrm{ComplexModule}}.\mathrm{Float}$
    in Figure~\ref{fig:cheatsheet} projects a single label, while
    $\qualifiedthis{\mathrm{Plus}}.\mathrm{other}.\mathrm{re}$ shows a
    multi-label projection list $w$.
  \item[Indexed form]
    $\dbi{n}.\, w$,
    where $n \ge 0$ is a de~Bruijn index~\cite{debruijn1972-nameless-dummies}:
    $n = 0$ refers to the reference's enclosing scope,
    and each increment of $n$ moves one scope level further out.
\end{description}
Both forms retrieve the target scope, that is, the fully inherited record
after all inheritance has been applied, and then project the properties
named by the projection list $w$ from it.
The named form desugars to the indexed form during parsing,
as described in Section~\ref{sec:mixin-trees}.

A new scope level is created at each $\{\ldots\}$ record literal.%
\footnote{
  We use \emph{record} for the syntactic construct, \emph{scope}
  for the name-resolution context it creates, and \emph{mixin} for its
  role as a composable unit; in inheritance-calculus these are three views
  of the same entity.
}
The scope contains all properties of that mixin, including those
inherited via inheritance sources.
For example, $\{a \mapsto \{\},\; b \mapsto \{\}\}$ is a single mixin whose scope
contains both $a$ and $b$; so is $\{r,\; a \mapsto \{\}\}$ where $r$
is a reference to a record containing $b$.
Inheritance sources do not create additional scope levels.
In Figure~\ref{fig:cheatsheet}, the record bound to $\mathrm{Plus}$ is
one such scope, containing $\mathrm{other}$, $\mathrm{\_sumReCall}$,
$\mathrm{\_sumRe}$, and the other properties defined directly within it.

\begin{itemize}
  \item $\qualifiedthis{\mathrm{Foo}}$ (equivalently, $\dbi{n}$ where
    $\mathrm{Foo}$ is $n$ steps above the enclosing scope)
    retrieves the record named $\mathrm{Foo}$.
  \item $\qualifiedthis{\mathrm{Foo}}.\ell$, or equivalently, $\dbi{n}.\ell$,
    projects property $\ell$ from $\mathrm{Foo}$.
  \item A qualifier-only source $q$ (with no projection list) inherits
    the entire enclosing scope, which may produce an infinitely deep
    tree; the bare $\mathrm{BuiltIns}$, $\mathrm{Math}$, and
    $\mathrm{Operator}$ sources in Figure~\ref{fig:cheatsheet} are
    qualifier-only imports of this kind.
\end{itemize}

When the qualifier is unambiguous, it may be omitted: the shorthand is
a bare projection list $w$. Its leading projection is resolved against
the innermost enclosing scope that defines the leading projection's
label directly in its record literal, and the source desugars to the
corresponding indexed form.
For example, $\mathrm{\_sumReCall}.\mathrm{result}$,
$\mathrm{a}.\mathrm{Plus}$, and $\mathrm{\_totalCall}.\mathrm{result}$
in Figure~\ref{fig:cheatsheet} are such qualifier-omitted shorthands.
Because this search inspects only record-literal definitions, the
shorthand can only access properties defined in record literals, not
those inherited through inheritance; one of the two qualified forms is
required for inherited properties or to bypass variable shadowing.
\fi

\section{\bilingual{Mixin Trees}{Mixin 树}}
\label{sec:mixin-trees}

\ifInheritanceChinese
在面向对象编程中,一个
\emph{mixin}~\cite{bracha1990-mixin-inheritance} 是一段行为片段;
它可以通过继承组合到一个类上。
继承演算只有一个抽象,即
\emph{深可合并 mixin}(deep-mergeable mixin),其成员是
\emph{深可合并属性}(deep-mergeable properties)。%
\footnote{
  在不引起歧义之处,我们简写为``mixin''与``property''。
  正如图~\ref{fig:cheatsheet} 所示,二者各自推广了它们在先前
  工作中的同名物:Bracha--Cook 的 mixin~\cite{bracha1990-mixin-inheritance} 与
  面向对象语言的属性。
}
它可以同时充当模块、类、对象、方法、字段以及
局部绑定,正如图~\ref{fig:cheatsheet} 所展示的那样,因为
没有不透明的方法:每个值都是
一个透明的记录,而继承总是对子树执行一次深度
合并。
我们把所得到的结构称为一棵\emph{mixin 树}。
它的语义是纯\emph{观察式}(observational)的:给定一个
继承演算程序,我们查询某个
标签 $\ell$ 是否是某个给定位置 $p$ 处的一个属性。

本节通过六个相互递归的
路径上的一阶函数来定义语义。
这六个函数把面向对象的概念映射为
集合论的运算:$\properties$ 与 $\supers$
推广成员查找与继承链,
$\overrides$ 与 $\bases$ 捕捉方法覆盖与
基类,而 $\resolve$ 与 $\this$ 处理名字
解析与限定 $\this$ 解析。
\else
In object-oriented programming, a
\emph{mixin}~\cite{bracha1990-mixin-inheritance} is a fragment of behavior
that can be composed onto a class via inheritance.
Inheritance-calculus has a single abstraction, the
\emph{deep-mergeable mixin}, whose members are
\emph{deep-mergeable properties}.%
\footnote{
  Where no ambiguity arises, we write simply ``mixin'' and ``property.''
  As Figure~\ref{fig:cheatsheet} shows, each generalizes its namesake in prior
  work: the Bracha--Cook mixin~\cite{bracha1990-mixin-inheritance} and the
  property of object-oriented languages.
}
which can simultaneously serve as module, class, object, method, field, and
local binding, as Figure~\ref{fig:cheatsheet} illustrated, because there
are no opaque methods: every value is
a transparent record, and inheritance always performs a deep
merge of subtrees.
We call the resulting structure a \emph{mixin tree}.
Its semantics is purely \emph{observational}: given an
inheritance-calculus program, one queries whether a
label $\ell$ is a property at a given position $p$.

This section defines the semantics via six mutually recursive
first-order functions on paths.
The six functions map object-oriented concepts to
set-theoretic operations: $\properties$ and $\supers$
generalize member lookup and inheritance chains,
$\overrides$ and $\bases$ capture method override and
base classes, and $\resolve$ and $\this$ handle name
resolution and qualified $\this$ resolution.
\fi

\subsection{\bilingual{Definitions}{定义}}
\label{sec:definitions}

\paragraph{\bilingual{Path}{路径}}
\ifInheritanceChinese
一个 $\Path$ 是一个标签序列
$(\ell_1, \ell_2, \ldots, \ell_n)$;它标识树中的一个
位置。
我们用 $p$ 表示一条路径。
\emph{根路径}是空序列 $()$。
给定一条路径 $p$ 与一个标签 $\ell$,
$p \snoc \ell$ 是序列 $p$ 用 $\ell$ 扩展后的结果。
对于非根路径,父路径 $\init(p)$ 是 $p$ 去掉
其最后一个元素后的结果,
而末尾标签 $\last(p)$ 是 $p$ 的最后一个元素。
为方便起见,当路径在上下文中清楚时,我们有时
用 $\ell$ 表示 $\last(p)$。
有了路径,我们就能陈述 AST 在
每条路径处提供什么。
\else
A $\Path$ is a sequence of labels
$(\ell_1, \ell_2, \ldots, \ell_n)$ that identifies a position in
the tree.
We write $p$ for a path.
The \emph{root path} is the empty sequence $()$.
Given a path $p$ and a label $\ell$,
$p \snoc \ell$ is the sequence $p$ extended with $\ell$.
For a nonroot path, the parent path $\init(p)$ is $p$ with
its last element removed,
and the final label $\last(p)$ is the last element of $p$.
For convenience, we sometimes write $\ell$ for $\last(p)$
when the path is clear from context.
With paths in hand, we can state what the AST provides
at each path.
\fi

\paragraph{AST}
\ifInheritanceChinese
来自第~\ref{sec:syntax} 节的一个表达式 $e$ 被解析成一棵 AST。
我们以两种等价的风格给出 AST:首先是一个固定其形状的
代数数据类型,然后是两个派生的投影;本节其余部分的
语义函数消费它们。

\emph{代数数据类型风格}把规范化后的 AST 给出为一个由节点项
构成的纯语法。一个节点 $N$ 把一个定义列表 $d$ 与一个
引用列表 $b$ 配成对;每个定义把一个标签绑定到一个子节点,而
每个引用是一个 de~Bruijn 索引 $n$ 连同一个投影列表
$w$:
\else
An expression $e$ from Section~\ref{sec:syntax} is parsed into an AST.
We give the AST in two equivalent styles: first an algebraic
datatype that fixes its shape, then two derived projections that
the semantic functions in the rest of this section consume.

The \emph{algebraic-datatype style} gives the normalized AST as a pure
grammar of node terms. A node $N$ pairs a definition list $d$ with a
reference list $b$; each definition binds a label to a child node, and
each reference is a de~Bruijn index $n$ together with a projection list
$w$:
\fi
\begin{align*}
  N \quad ::= \quad & \langle\, d \,;\; b \,\rangle
  && \text{\bilingual{(node: definitions and references)}{(节点:定义与引用)}} \\[6pt]
  d \quad ::= \quad & \varepsilon
  && \text{\bilingual{(no definitions)}{(没有定义)}} \\*
  \mid\quad & (\ell \mapsto N),\; d
  && \text{\bilingual{(definition, then more)}{(一个定义,然后更多)}} \\[6pt]
  b \quad ::= \quad & \varepsilon
  && \text{\bilingual{(no references)}{(没有引用)}} \\*
  \mid\quad & (n,\; w),\; b
  && \text{\bilingual{(reference, then more)}{(一个引用,然后更多)}} \\[6pt]
  w \quad ::= \quad & \varepsilon
  && \text{\bilingual{(no projections)}{(没有投影)}} \\*
  \mid\quad & \ell.\, w
  && \text{\bilingual{(projection, then more)}{(一个投影,然后更多)}}
\end{align*}
\ifInheritanceChinese
与表达式的元素列表一样,定义列表 $d$ 是
递归写成的,但读取时不计重排与重复:它
表示从其各绑定收集而来的有限偏映射 $\defs : \Label \rightharpoonup \Node$;该映射
对每个被定义的标签恰有一条目,而
$N.\defs(\ell)$ 是绑定到 $\ell$ 的子节点(当 $\ell$
未绑定时无定义)。同样地,引用列表 $b$ 表示由继承源
$(n,\; w)$ 构成的集合
$\refs \subseteq \mathbb{N} \times \Label^{*}$,读取时不计重排与重复。由于继承
是幂等的且记录是集合,这些商不丢失任何
信息。
该语法只固定形状;一个节点项是良构的,当它
满足:
\else
Like the element list of an expression, the definition list $d$ is
written recursively but read up to reordering and duplication: it
denotes the finite partial map $\defs : \Label \rightharpoonup \Node$
collected from its bindings, with one entry per defined label, and
$N.\defs(\ell)$ is the child bound to $\ell$ (undefined when $\ell$ is
unbound). Likewise the reference list $b$ denotes the set
$\refs \subseteq \mathbb{N} \times \Label^{*}$ of inheritance sources
$(n,\; w)$, read up to reordering and duplication. Because inheritance
is idempotent and a record is a set, these quotients lose no
information.
The grammar fixes only the shape; a node term is well-formed when it
satisfies:
\fi
\begin{enumerate}
  \item \ifInheritanceChinese 路径 $p$ 处的每个源 $(n,\; w)$ 满足 $n < |p|$。\else Each source $(n,\; w)$ at path $p$ satisfies $n < |p|$.\fi
  \item \ifInheritanceChinese 一个 $w$ 为空的源 $(n,\; \varepsilon)$ 是
    一个仅限定符源的规范化;第~\ref{sec:syntax} 节的表达式语法
    没有 $w$ 为空的仅投影源。\else A source $(n,\; \varepsilon)$ with empty $w$ is the
    normalization of a qualifier-only source; the expression grammar of
    Section~\ref{sec:syntax} has no projection-only source with empty
    $w$.\fi
\end{enumerate}
\ifInheritanceChinese
该节点语法被\emph{余归纳地}(coinductively)读取:通过
继承,一个子节点可能展开成一棵无界乃至
无限深的树,所以这些产生式描述的是最大不动点,
而非有限生成的最小不动点。整个程序是单个
全局节点 $\ASTroot$,即顶层记录的规范化 AST。

\emph{余代数风格}(coalgebraic style)用两个由路径 $p$ 索引的
原始函数替换全局树,每次观察一个节点。它们通过沿 $p$
导航从 $\ASTroot$ 派生而来:
从 $\ASTroot$ 出发,在 $p$ 的每个标签处跟随 $\defs$ 映射
到达 $p$ 处的节点,而这两个投影读出其
字段。
\else
The node grammar is read \emph{coinductively}: through
inheritance a child node may expand to a tree of unbounded or even
infinite depth, so the productions describe the greatest fixed point,
not a finitely generated least one. The whole program is a single
global node $\ASTroot$, the normalized AST of the top-level record.

The \emph{coalgebraic style} replaces the global tree by two
primitive functions indexed by a path $p$, observing one node at a
time. They are derived from $\ASTroot$ by navigating along $p$:
starting at $\ASTroot$ and following the $\defs$ map at each label
of $p$ reaches the node at $p$, and the two projections read off its
fields.
\fi
\begin{align}
  \defines(p) &= \operatorname{dom}\bigl(
    \ASTroot.\defs(\ell_1).\defs(\ell_2)\cdots\defs(\ell_n).\defs
  \bigr)
  \label{eq:defines}\\
  \inherits(p) &= \ASTroot.\defs(\ell_1).\defs(\ell_2)\cdots\defs(\ell_n).\refs
  \label{eq:inherits}
\end{align}
\ifInheritanceChinese
其中 $p = (\ell_1, \ldots, \ell_n)$,从而 $\defines(p)$ 是
那些使得 $p$ 包含一个定义 $\ell \mapsto e$ 的标签 $\ell$ 的集合,
而 $\inherits(p)$ 是 $p$ 处的引用对
$(n,\; w)$ 的集合。
在根路径 $()$ 处导航为空,给出
$\defines(()) = \operatorname{dom}(\ASTroot.\defs)$ 以及
$\inherits(()) = \ASTroot.\refs$。
下文的语义函数只使用 $\defines$ 与 $\inherits$,
所以任一风格皆可取作原始。

在解析期间,三种语法形式全都被解析为
de~Bruijn 索引对 $(n,\; w)$。
\emph{索引形式} $\dbi{n}.\, w$
已经直接携带 de~Bruijn 索引 $n$;
不需要解析。
路径 $p$ 处的一个\emph{具名引用} $\qualifiedthis{\ell_{\mathrm{qualifier}}}.\, w$
通过在 $p$ 的标签中找到
$\ell_{\mathrm{qualifier}}$ 的最后一次出现来解析。
设 $p_{\mathrm{target}}$ 为 $p$ 的前缀,直到并包括
那次出现。
de~Bruijn 索引为 $n = |p| - |p_{\mathrm{target}}| - 1$,
而投影列表 $w$ 原样保留。
路径 $p$ 处一个首投影为 $\ell$ 的\emph{词法引用} $w$
通过找到 $p$ 的最近前缀 $p'$
使得 $\ell \in \defines(p')$ 来解析。
de~Bruijn 索引为 $n = |p| - |p'| - 1$,而投影列表
$w$ 原样保留。
三种形式都产生同一种表示;下文的语义函数
只在 de~Bruijn 索引对上操作。

AST 是纯数据,由程序文本固定,而不是
运行期状态的函数。%
\footnote{
  一个实现可以非严格地解析 AST:$\ASTroot$ 不必
  预先被完全物化。每个 $\Node$ 把它的源
  文本保持为未解析,直到某个查询观察它,从而读取
  $\defines(p)$ 或 $\inherits(p)$ 时按需强制
  沿 $p$ 的子树。这只影响解析何时发生,而不影响其结果,
  因为 AST 是引用透明的。
}
有了两个原始函数 $\defines$ 与 $\inherits$,我们现在可以回答
本节开头提出的那个问题:
某个标签 $\ell$ 是否是某个给定位置 $p$ 处的一个属性。
\else
for $p = (\ell_1, \ldots, \ell_n)$, so that $\defines(p)$ is the set
of labels $\ell$ for which $p$ contains a definition $\ell \mapsto e$,
and $\inherits(p)$ is the set of reference pairs
$(n,\; w)$ at $p$.
At the root path $()$ the navigation is empty, giving
$\defines(()) = \operatorname{dom}(\ASTroot.\defs)$ and
$\inherits(()) = \ASTroot.\refs$.
The semantic functions below use only $\defines$ and $\inherits$,
so either style may be taken as primitive.

During parsing, all three syntactic forms are resolved to
de~Bruijn index pairs $(n,\; w)$.
The \emph{indexed form} $\dbi{n}.\, w$
already carries the de~Bruijn index $n$ directly;
no resolution is needed.
A \emph{named reference} $\qualifiedthis{\ell_{\mathrm{qualifier}}}.\, w$
at path $p$ is resolved by finding the last occurrence of
$\ell_{\mathrm{qualifier}}$ among the labels of $p$.
Let $p_{\mathrm{target}}$ be the prefix of $p$ up to and including
that occurrence.
The de~Bruijn index is $n = |p| - |p_{\mathrm{target}}| - 1$
and the projection list $w$ is carried over unchanged.
A \emph{lexical reference} $w$ at path $p$, with leading
projection $\ell$, is resolved by finding the nearest prefix $p'$ of $p$
such that $\ell \in \defines(p')$.
The de~Bruijn index is $n = |p| - |p'| - 1$ and the projection list
$w$ is carried over unchanged.
All three forms produce the same representation; the semantic functions
below operate only on de~Bruijn index pairs.

The AST is pure data, fixed by the program text and not a function
of runtime state.%
\footnote{
  An implementation may parse the AST non-strictly: $\ASTroot$ need
  not be materialized in full up front. Each $\Node$ keeps its source
  text unparsed until a query observes it, so that reading
  $\defines(p)$ or $\inherits(p)$ forces the subtree along $p$ on
  demand. This affects only when parsing happens, not its result,
  since the AST is referentially transparent.
}
Given the two primitives $\defines$ and $\inherits$, we can now answer
the question posed at the beginning of this section:
whether a label $\ell$ is a property at a given position $p$.
\fi

\paragraph{\bilingual{Properties}{属性}}
\ifInheritanceChinese
一条路径的属性是它自身的属性,连同
从所有 supers 继承而来的属性:
\else
The properties of a path are its own properties together with
those inherited from all supers:
\fi
\begin{equation}\label{eq:properties}
  \properties(p) =
  \bigl\{\; \ell \;\big|\;
    (\_,\; p_{\mathrm{override}}) \in \supers(p),\;
    \ell \in \defines(p_{\mathrm{override}})
  \;\bigr\}
\end{equation}
\ifInheritanceChinese
本节其余部分定义 $\supers$,
即传递的继承闭包,
及其依赖项。
\else
The remainder of this section defines $\supers$,
the transitive inheritance closure,
and its dependencies.
\fi

\paragraph{Supers}
\ifInheritanceChinese
直观上,$\supers(p)$ 收集 $p$ 所继承的每一条路径:
$p$ 自身的 $\overrides$,
加上 $p$ 的每个
直接基 $\bases$ 的 $\overrides$,如此传递下去。
每个结果与到达它所经由的继承点语境
$\init(p_{\mathrm{base}})$ 配成对。
我们下文定义的限定 $\this$ 解析需要这一来源,
以便把一个定义点 mixin 映射回
那些纳入它的继承点路径:
\else
Intuitively, $\supers(p)$ collects every path that $p$ inherits from:
the $\overrides$ of $p$ itself,
plus the $\overrides$ of each
direct base $\bases$ of $p$, and so on transitively.
Each result is paired with the inheritance-site context
$\init(p_{\mathrm{base}})$ through which it is reached.
This provenance is needed by qualified $\this$ resolution,
which we define below, to map a definition-site mixin back to
the inheritance-site paths that incorporate it:
\fi
\begin{equation}\label{eq:supers}
  \supers(p) =
  \bigl\{\; (\init(p_{\mathrm{base}}),\; p_{\mathrm{override}}) \;\big|\;
    p_{\mathrm{base}} \in \bases^*(p),\;
    p_{\mathrm{override}} \in \overrides(p_{\mathrm{base}})
  \;\bigr\}
\end{equation}
\ifInheritanceChinese
这里 $\bases^*$ 表示 $\bases$ 的自反传递闭包。
$\supers$ 公式依赖两个函数:
$\overrides$ 与单跳引用目标 $\bases$。
我们先定义 $\overrides$。
\else
Here $\bases^*$ denotes the reflexive-transitive closure of $\bases$.
The $\supers$ formula depends on two functions:
$\overrides$ and the one-hop reference targets $\bases$.
We define $\overrides$ first.
\fi

\paragraph{Overrides}
\ifInheritanceChinese
由于一个 mixin 可以扮演一个方法的角色,
一个子作用域应当能够
通过父作用域的标签找到它,即便该子作用域
并不重新定义那个标签;这就是方法覆盖。
具体地说,继承可以在同一个作用域层级引入同一标签的
多个定义;
$\overrides(p)$ 收集所有与 $p$ 共享
\emph{相同身份}(same identity)的这类路径,从而使它们的定义
相互合并而非相互遮蔽,以实现深度合并。
$p$ 的 overrides 包含 $p$ 自身
以及对于 $\init(p)$ 的每个
也定义 $\last(p)$ 的分支 $p_{\mathrm{branch}}$,$p_{\mathrm{branch}} \snoc \last(p)$:
\else
Since a mixin can play the role of a method,
a child scope should be
able to find it through the parent's label even if the child
does not redefine that label; this is method override.
Concretely, inheritance can introduce multiple
definitions of the same label at the same scope level;
$\overrides(p)$ collects all such paths that share the
\emph{same identity} as $p$, so that their definitions
merge rather than shadow each other, enabling deep merge.
The overrides of $p$ include $p$ itself
and $p_{\mathrm{branch}} \snoc \last(p)$ for every
branch $p_{\mathrm{branch}}$ of $\init(p)$ that also defines $\last(p)$:
\fi
\begin{equation}\label{eq:overrides}
  \overrides(p) =
  \begin{cases}
    \{p\} & \text{if } p = ()  \\[6pt]
    \{p\}
    \;\cup\;
    \left\{\; p_{\mathrm{branch}} \snoc \last(p) \;\middle|\;
      \begin{aligned}
        &(\_,\; p_{\mathrm{branch}}) \in \supers(\init(p)), \\
        &\text{s.t.}\; \last(p) \in \defines(p_{\mathrm{branch}})
      \end{aligned}
    \right\}
    & \text{if } p \neq ()
  \end{cases}
\end{equation}
\ifInheritanceChinese
还剩下要定义 $\bases$,即 $\supers$ 的另一个依赖项。
\else
It remains to define $\bases$, the other dependency of $\supers$.
\fi

\paragraph{Bases}
\ifInheritanceChinese
直观上,$\bases(p)$ 是 $p$ 通过引用直接
继承的那些路径,类比于面向对象语言中的直接
基类。
具体地说,$\bases$ 把 $p$ 的 $\overrides$ 中的每个引用
解析一步:
\else
Intuitively, $\bases(p)$ are the paths that $p$ directly
inherits from via references, analogous to the direct base
classes in object-oriented languages.
Concretely, $\bases$ resolves every reference in
$p$'s $\overrides$ one step:
\fi
\begin{equation}\label{eq:bases}
  \bases(p) =
  \left\{\; p_{\mathrm{target}} \;\middle|\;
    \begin{aligned}
      &p_{\mathrm{override}} \in \overrides(p), \\
      &(n,\; w) \in \inherits(p_{\mathrm{override}}), \\
      &p_{\mathrm{target}} \in \resolve(\init(p),\; p_{\mathrm{override}},\; n,\; w)
    \end{aligned}
  \right\}
\end{equation}
\ifInheritanceChinese
$\bases$ 公式调用引用解析函数 $\resolve$,我们接下来定义它。
\else
The $\bases$ formula calls the reference resolution function $\resolve$, which we define next.
\fi

\paragraph{\bilingual{Reference resolution}{引用解析}}
\ifInheritanceChinese
直观上,$\resolve$ 把一个语法引用转化为它在
完全继承后的树中所指向的路径集合。
它取一条继承点路径 $p_{\mathrm{site}}$、
一条定义点路径 $p_{\mathrm{def}}$、
一个 de~Bruijn 索引 $n$,
以及投影列表 $w$。
解析分两个阶段进行:
$\this$ 从外围作用域 $\init(p_{\mathrm{def}})$ 出发
执行 $n$ 次向上步进,
把定义点路径映射为继承点路径;
然后追加
$w = \ell_{\mathrm{projection},1}\ldots\ell_{\mathrm{projection},k}$
的向下投影:
\else
Intuitively, $\resolve$ turns a syntactic reference into the
set of paths it points to in the fully inherited tree.
It takes an inheritance-site path $p_{\mathrm{site}}$,
a definition-site path $p_{\mathrm{def}}$,
a de~Bruijn index $n$,
and the projection list $w$.
Resolution proceeds in two phases:
$\this$ performs $n$ upward steps
starting from the enclosing scope $\init(p_{\mathrm{def}})$,
mapping the definition-site path to inheritance-site paths;
then the downward projections of
$w = \ell_{\mathrm{projection},1}\ldots\ell_{\mathrm{projection},k}$
are appended:
\fi
\begin{multline}\label{eq:resolve}
  \resolve(p_{\mathrm{site}},\; p_{\mathrm{def}},\; n,\; w)
  = \\
  \bigl\{\;
    p_{\mathrm{current}} \snoc \ell_{\mathrm{projection},1} \snoc \cdots \snoc \ell_{\mathrm{projection},k}
    \;\big|
    p_{\mathrm{current}} \in
    \this(\{p_{\mathrm{site}}\},\; \init(p_{\mathrm{def}}),\; n)
  \;\bigr\}
\end{multline}
\ifInheritanceChinese
$\resolve$ 返回一个集合而非单条路径,因为
$\this$ 可能解析到多个目标。
我们现在定义 $\this$。
\else
$\resolve$ returns a set rather than a single path because
$\this$ may resolve to multiple targets.
We now define $\this$.
\fi

\paragraph{\bilingual{Qualified $\this$ resolution}{限定 $\this$ 解析}}
\ifInheritanceChinese
$\this$ 回答这个问题:
``在完全继承后的树中,定义点作用域
$p_{\mathrm{def}}$ 究竟住在哪里?''
每一步在前沿 $S$ 中每条路径的 supers 当中,
找出那些其 override 分量匹配
$p_{\mathrm{def}}$ 的,并把对应的
继承点路径收集为新的前沿。
经过 $n$ 步后,前沿就包含答案:
\else
$\this$ answers the question:
``in the fully inherited tree, where does
the definition-site scope $p_{\mathrm{def}}$ actually live?''
Each step finds, among the supers of every path in the
frontier $S$, those whose override component matches
$p_{\mathrm{def}}$, and collects the corresponding
inheritance-site paths as the new frontier.
After $n$ steps the frontier contains the answer:
\fi
{\small
  \begin{equation}\label{eq:this}
    \this(S,\; p_{\mathrm{def}},\; n) =
    \begin{cases}
      S & \text{if } n = 0 \\[6pt]
      \this\!\left(
        \left\{\; p_{\mathrm{site}} \;\middle|\;
          \begin{aligned}
            &p_{\mathrm{current}} \in S, \\
            &(p_{\mathrm{site}},\; p_{\mathrm{override}})
            \in \supers(p_{\mathrm{current}}), \\
            &\text{s.t.}\; p_{\mathrm{override}} = p_{\mathrm{def}}
          \end{aligned}
        \right\},\;
        \init(p_{\mathrm{def}}),\;
        n - 1
      \right)
      & \text{if } n > 0
    \end{cases}
\end{equation}}%
\ifInheritanceChinese
每一步 $\init$ 把 $p_{\mathrm{def}}$ 缩短一个标签,
而 $n$ 减少一。
由于 $n$ 是一个非负整数,该递归终止。%
\footnote{
  若在某一步不存在匹配对
  $(p_{\mathrm{site}},\; p_{\mathrm{override}})$,
  则前沿变为空集,而
  $\resolve$ 返回 $\varnothing$:该引用没有
  目标;不继承任何属性。
  这是无类型演算最为自洽的
  语义。
  MIXINv2\anon[~(补充材料)]{~\cite{mixin2025}} 更为严格:它的解析器拒绝
  前沿变为空的引用,在编译期报告一个
  解析错误。
}

这就完成了计算 $\properties(p)$ 所需的整条定义链。

附录~\ref{app:evaluation-trace} 把这六个方程逐步
应用到一个突出 $\resolve$ 集合值本质的例子上;
在那里前沿 $S$ 携带多个目标,而
$\this$ 同时解析到它们,而非如 Scala~3 与 NixOS 模块
这类面向对象语言中所见的单个目标
(附录~\ref{app:scala-multi-target})。
\else
At each step $\init$ shortens $p_{\mathrm{def}}$ by one label
and $n$ decreases by one.
Since $n$ is a nonnegative integer, the recursion terminates.%
\footnote{
  If no matching pair
  $(p_{\mathrm{site}},\; p_{\mathrm{override}})$ exists
  at some step, the frontier becomes the empty set, and
  $\resolve$ returns $\varnothing$: the reference has no
  target; no properties are inherited.
  This is the most self-consistent semantics for the
  untyped calculus.
  MIXINv2\anon[~(supplementary material)]{~\cite{mixin2025}} is stricter: its resolver rejects
  references whose frontier becomes empty, reporting a
  resolution error at compile time.
}

This completes the chain of definitions needed to compute $\properties(p)$.

Appendix~\ref{app:evaluation-trace} applies the six equations step by
step to an example highlighting the set-valued nature of $\resolve$,
where the frontier $S$ carries multiple targets that
$\this$ resolves to simultaneously rather than the single target
seen in object-oriented languages such as Scala~3 and NixOS modules
(Appendix~\ref{app:scala-multi-target}).
\fi

\subsection{\bilingual{Recursive Evaluation}{递归求值}}
\ifInheritanceChinese
六个方程~(\ref{eq:properties})--(\ref{eq:this})
通过相互递归定义 $\properties(p)$:
$\ell \in \properties(p)$ 成立,当且仅当它能
由这些函数的有限次应用链
确立。
这些函数\emph{就是}解释器:观察者
通过选择检查哪条路径来驱动计算,
而每个方程按需展开。
由于这些函数是纯的且与顺序无关,
(\ref{eq:supers}) 中的自反传递闭包 $\bases^*$
可以按广度优先或
深度优先来探索而不影响结果;
路径可以被 intern,从而结构相等
退化为指针相等;
而这六个函数构成一个以路径为键的可记忆化
动态规划递推式。

在操作上,求值使用\emph{tabled
求值}~\cite{tamaki1986-tabled-resolution}:
每个查询在首次遇到时被记忆化,而
重入调用,即对应于依赖图中的环的调用,
返回当前的累加器而非发散。
在指称上,六个方程定义了一个单调的
直接后承
算子~$T_P$~\cite{vanemden1976-predicate-logic-semantics};
它的右端只使用 $\cup$、$\in$、相等以及
正向集合概括,没有否定、集合
差或补。
由 Knaster--Tarski
定理~\cite{tarski1955-lattice-fixpoint-theorem},$T_P$ 有一个唯一的
最小不动点 $\mathrm{lfp}(T_P)$;它充当
指称语义。
这一指称刻画把上文引入的观察式
语义形式化:观察者驱动的查询
由同一组不动点方程计算。
求值是否终止取决于程序。
回忆 $\supers$(方程~\ref{eq:supers})取
$\bases$(方程~\ref{eq:bases})的自反传递
闭包;所得的图就是
\emph{继承层级}(inheritance hierarchy)。
由于每个继承源在 mixin 树中占据一个不同的节点,
这一层级可从
树结构中直接读出。
继承层级无环的程序在
一趟内终止。
当层级包含环时,
当每个查询有有限的答案域时求值
终止:有限格上的升链条件保证
收敛~\cite{bancilhon1986-recursive-query-strategies},
如同在 Datalog~\cite{vanemden1976-predicate-logic-semantics} 中那样。
当层级无限时,求值可能
发散;这是图灵完备性所固有的。
附录~\ref{app:well-definedness} 证明 tabled
求值绝不返回错误答案,并把
这些终止条件形式化。
\else
The six equations~(\ref{eq:properties})--(\ref{eq:this})
define $\properties(p)$ by mutual recursion:
$\ell \in \properties(p)$ holds if and only if it can
be established by a finite chain of applications of
these functions.
The functions \emph{are} the interpreter: the observer
drives the computation by choosing which path to
inspect, and each equation unfolds on demand.
Because the functions are pure and order-independent,
the reflexive-transitive closure $\bases^*$
in~(\ref{eq:supers}) may be explored breadth-first or
depth-first without affecting the result;
paths may be interned so that structural equality
reduces to pointer equality;
and the six functions form a memoizable
dynamic-programming recurrence keyed by paths.

Operationally, evaluation uses \emph{tabled
evaluation}~\cite{tamaki1986-tabled-resolution}:
each query is memoised on first encounter, and
re-entrant calls, which correspond to cycles in the dependency graph,
return the current accumulator rather than diverging.
Denotationally, the six equations define a monotone
immediate consequence
operator~$T_P$~\cite{vanemden1976-predicate-logic-semantics}
whose right-hand sides use only $\cup$, $\in$, equality, and
positive set comprehension, with no negation, set
difference, or complement.
By the Knaster--Tarski
theorem~\cite{tarski1955-lattice-fixpoint-theorem}, $T_P$ has a unique
least fixed point $\mathrm{lfp}(T_P)$, which serves as
the denotational semantics.
This denotational characterization formalizes the observational
semantics introduced above: the observer-driven queries are
computed by the same fixed-point equations.
Whether evaluation terminates depends on the program.
Recall that $\supers$ (equation~\ref{eq:supers}) takes the
reflexive-transitive closure of $\bases$
(equation~\ref{eq:bases}); the resulting graph is the
\emph{inheritance hierarchy}.
Since each inheritance source occupies a distinct node in
the mixin tree, this hierarchy is directly readable from
the tree structure.
Programs whose inheritance hierarchy is acyclic terminate in
one pass.
When the hierarchy contains cycles,
evaluation terminates when each query has a finite
answer domain: the ascending chain condition on the finite
lattice guarantees
convergence~\cite{bancilhon1986-recursive-query-strategies},
as in Datalog~\cite{vanemden1976-predicate-logic-semantics}.
When the hierarchy is infinite, evaluation may
diverge; this is inherent to Turing completeness.
Appendix~\ref{app:well-definedness} proves that tabled
evaluation never returns a wrong answer and formalizes
these termination conditions.
\fi

%% === 审稿人备忘：六个方程的渐近复杂度分析 ===
%%
%% 论文声称语义是一阶的，审稿人可能追问：一阶性是否带来额外的
%% 计算开销？下面按场景分析。
%%
%% 关键观察：set comprehension 自动去重（集合语义），加上路径
%% intern（结构相等退化为指针相等）和 memoization，相同的
%% (path, query) 对只计算一次。这不是可选的优化，而是集合语义
%% 的固有性质——集合不会包含重复元素。
%%
%% --- 场景 1：λ-calculus 翻译 ---
%% Single-Path Lemma (Lemma 3) 保证 this 的 frontier 集合 S
%% 始终只有一个元素。每个路径的 overrides 也只有一个元素。
%% bases 每个节点最多一个继承源。supers 的传递闭包是线性链。
%% 复杂度：O(n)，其中 n 是 ANF term 的大小。
%% 与传统环境机（environment machine）一致。
%%
%% --- 场景 2：Object algebra（k 个 constructor，t 个 operation 文件，m 个字段） ---
%% 多个 operation 文件各自继承同一个 data 文件，然后在最终组合时
%% merge。每个 constructor 在每个 operation 文件里各有一个定义，
%% 通过 overrides 合并，所以 overrides 大小为 O(t)。
%%
%% 关键点：this 的 frontier 是 1，不是"去重后 O(1)"，就是 1。
%% 原因：qualified this 解析的是外层 scope（如 NatFactory），
%% 而这个外层 scope 在继承链上只有一条路径到达——因为是从一个
%% 具体的 operation 定义点出发往上走的。multi-target this 的前提
%% 是同一个定义点被继承到多个不同的 site，这只在 trie union 时
%% 发生。Object algebra 的继承是把多个 operation merge 到同一个
%% factory 上，不是把多个 factory merge 成一个值。方向不同。
%%
%% 这也解释了与 Scala 的关系：Scala 拒绝 multi-target self-reference resolution
%%（Appendix B 的 "conflicting base types"），但在 object algebra
%% 场景下 this 本来就是 single-target 的，所以 Scala 的拒绝是
%% 过度保守——它拒绝了一类去重后（实际上根本不需要去重，因为
%% frontier 就是 1）完全合法的组合。
%%
%% 复杂度：O(k · m · t)，与 Scala trait linearization 一致。
%%
%% --- 场景 3：Nat/BinNat/Boolean 算术 ---
%% 数值的递归深度为 n（Nat 值的大小），每次递归触发一次 this
%% 解析（single-target），遍历 O(t) 个 override。
%% Nat 算术复杂度：O(n · t)，即 Peano 算术 + 常数因子。
%% BinNat 算术复杂度：O(log n · t)，即二进制算术 + 常数因子。
%%
%% --- 场景 4：Cartesian product / 逻辑编程 ---
%% 这是唯一出现真正 multi-target this 的场景。
%% trie union 的 bases 包含多条路径，每条路径的 addend 通过
%% this 解析到另一个 trie union，路径数相乘。
%% 对于 |R₁| × |R₂| 的 join，复杂度为 O(|R₁| · |R₂|)
%% （memoization 下）。
%%
%% 这与 B-Tree 引擎的 Datalog nested-loop join 复杂度一致。
%% trie 天然有索引——路径本身是前缀树，共享前缀的值自动共享
%% 计算（memoization 命中）。这与 Datalog 引擎对 B-Tree 做
%% prefix scan 在精神上一致。
%%
%% --- 总结 ---
%%
%% | 场景               | 复杂度              | 对比               |
%% |--------------------|---------------------|--------------------|
%% | λ-calculus         | O(n)                | = 环境机           |
%% | Object algebra     | O(k·m·t)            | = Scala线性化      |
%% | Nat 算术           | O(n·t)              | = Peano 算术       |
%% | BinNat 算术        | O(log n · t)        | = 二进制算术       |
%% | Cartesian product  | O(|R₁|·|R₂|)       | = nested-loop join |
%%
%% 每个场景都与对应的专用引擎复杂度一致。一阶性没有带来额外开销。
%%
%% Single-Path Lemma 的适用范围比论文明示的更大：不仅 λ-calculus
%% 翻译是 single-path，所有不涉及 trie union 的正常 OOP/object
%% algebra 代码都是 single-path。multi-target this 是 trie union
%% 的专属现象，而 trie union 对应的恰好是关系代数的 join，其
%% 复杂度是固有的，与引擎无关。
%% ================================================================

\section{\bilingual{Embedding the \texorpdfstring{$\lambda$}{λ}-Calculus}{嵌入 \texorpdfstring{$\lambda$}{λ}-演算}}
\label{sec:translation}
\label{sec:forward-translation}

\ifInheritanceChinese
以下翻译将 A-正规形式（ANF）~\cite{flanagan1993-essence-compiling-continuations}的 $\lambda$-演算映射到继承演算。
在 ANF 中,每个中间结果都绑定到一个名字,且每个表达式至多包含一次应用。
我们用一个\emph{值绑定} $\mathbf{let}\; x = V \;\mathbf{in}\; M$ 来扩展标准 ANF 语法,该值绑定将一个并非应用结果的值命名:%
\footnote{
  标准 ANF 允许将值内联用作操作数;值绑定产生式为其命名,使得翻译后程序的每个子表达式都可通过词法引用访问(附录~\ref{app:expressiveness})。
}
\[
  \begin{array}{r@{\;::=\;}l}
    M & \mathbf{let}\; x = V_1\; V_2 \;\mathbf{in}\; M
    \mid \mathbf{let}\; x = V \;\mathbf{in}\; M
    \mid V_1\; V_2
    \mid V \\
    V & x \mid \lambda x.\, M
  \end{array}
\]
每个 $\lambda$-项都可以通过对中间结果命名的方式机械地转换为 ANF。\footnote{%
  这可以选择性地包括对值的命名;由于每个 ANF 项也是一个 $\lambda$-项,ANF 既不是 $\lambda$-演算的限制,也不是其扩展。
}
ANF 结构与继承演算天然吻合:每个 $\mathbf{let}$-绑定成为一个命名的记录属性,每次应用则通过 $.\mathrm{result}$ 包装并投影。

设 $\mathcal{T}$ 为翻译函数。每个 $\lambda$-抽象在翻译后的 mixin 树中引入一个作用域层级。
\else
The following translation maps the $\lambda$-calculus in
A-normal form (ANF)~\cite{flanagan1993-essence-compiling-continuations} to inheritance-calculus.
In ANF, every intermediate result is bound to a name and each
expression contains at most one application.
We extend the standard ANF grammar with a \emph{value binding}
$\mathbf{let}\; x = V \;\mathbf{in}\; M$, which gives a name to a
value that is not an application result:%
\footnote{
  Standard ANF allows values inline as operands;
  the value-binding production names them so that every
  subexpression of the translated program is accessible
  via a lexical reference (Appendix~\ref{app:expressiveness}).
}
\[
  \begin{array}{r@{\;::=\;}l}
    M & \mathbf{let}\; x = V_1\; V_2 \;\mathbf{in}\; M
    \mid \mathbf{let}\; x = V \;\mathbf{in}\; M
    \mid V_1\; V_2
    \mid V \\
    V & x \mid \lambda x.\, M
  \end{array}
\]
Every $\lambda$-term can be mechanically converted to ANF by
naming intermediate results.\footnote{%
  This can optionally include naming values as well;
  since every ANF term is also a $\lambda$-term,
  ANF is neither a restriction nor an extension of the $\lambda$-calculus.
}
The ANF structure matches inheritance-calculus naturally: each
$\mathbf{let}$-binding becomes a named record property, and
each application is wrapped and projected via
$.\mathrm{result}$.

Let $\mathcal{T}$ denote the translation function.
Each $\lambda$-abstraction introduces one scope level
in the translated mixin tree.
\fi

\medskip
\begin{center}
  \small
  \begin{tabular}{l@{\quad$\longrightarrow$\quad}l}
    $x$ \text{\scriptsize(\bilingual{$\lambda$-bound, index $n'$}{$\lambda$-绑定,索引 $n'$})} &
    $\dbi{n'}.\mathrm{argument}$ \\[4pt]
    $x$ \text{\scriptsize(\bilingual{let-bound}{let-绑定})} &
    $x.\mathrm{result}$ \\[4pt]
    $\lambda x.\, M$ &
    $\{\mathrm{argument} \mapsto \{\},\;
    \mathrm{result} \mapsto \mathcal{T}(M)\}$ \\[4pt]
    $\mathbf{let}\; x = V_1\; V_2
    \;\mathbf{in}\; M$ &
    $\{x \mapsto \{\mathcal{T}(V_1),\;
      \mathrm{argument} \mapsto \mathcal{T}(V_2)\},$ \\
    & $\;\mathrm{result} \mapsto \mathcal{T}(M)\}$ \\[4pt]
    $\mathbf{let}\; x = V
    \;\mathbf{in}\; M$ &
    $\{x \mapsto \{\mathrm{result} \mapsto \mathcal{T}(V)\},\;
    \mathrm{result} \mapsto \mathcal{T}(M)\}$ \\[4pt]
    $V_1\; V_2$ \text{\scriptsize(\bilingual{tail call}{尾调用})} &
    $\{\mathrm{tailCall} \mapsto \{\mathcal{T}(V_1),\;
      \mathrm{argument} \mapsto \mathcal{T}(V_2)\},$ \\
    & $\;\mathrm{result} \mapsto \{\mathrm{tailCall}.\mathrm{result}\}\}$
  \end{tabular}
\end{center}
\medskip

\ifInheritanceChinese
带 de~Bruijn 索引 $n$ 的 $\lambda$-绑定变量 $x$(即从引用处到其绑定者之间包围的 $\lambda$ 数目)翻译为 $\dbi{n'}.\mathrm{argument}$,该引用到达绑定 $x$ 的抽象并访问其 $\mathrm{argument}$ 槽。%
\footnote{\label{fn:debruijn-adjust}在 $\lambda$-演算中,de~Bruijn 索引只计数 $\lambda$-抽象。在继承演算中,de~Bruijn 索引计数所有包围的作用域层级,包括由 $\mathbf{let}$-绑定和尾调用引入的层级。因此翻译会调整该索引:若一个 $\lambda$-绑定变量跨越 $k$ 个中间的 $\mathbf{let}$-绑定或尾调用作用域(即 $k$ 计数从变量出现处到其绑定的 $\lambda$ 在 ANF 语法树中路径上的 $\mathbf{let}$- 和尾调用构造数),则其继承演算索引为 $n' = n + k$,而非 $n$。let-绑定变量不受影响,因为它们使用词法引用(如 $x.\mathrm{result}$)而非 de~Bruijn 索引。我们假设翻译引入的合成标签($\mathrm{argument}$、$\mathrm{result}$、$\mathrm{tailCall}$)与用户定义的标签(如 let-绑定名 $x$)来自不相交的字母表,因此不会发生冲突。}
let-绑定变量 $x$ 翻译为词法引用 $x.\mathrm{result}$,从封闭记录中的兄弟属性 $x$ 投影应用结果。

抽象 $\lambda x.\, M$ 翻译为具有自身属性 $\{\mathrm{argument}, \mathrm{result}\}$ 的记录;我们称之为\emph{抽象形状}。在应用 $\mathbf{let}$-绑定规则中,名字 $x$ 绑定到继承记录 $\{\mathcal{T}(V_1),\;\mathrm{argument} \mapsto \mathcal{T}(V_2)\}$,而 $\mathcal{T}(M)$ 可引用 $x.\mathrm{result}$ 以获取应用的值。在值 $\mathbf{let}$-绑定规则中,$x$ 将 $\mathcal{T}(V)$ 包装在一个 $\mathrm{result}$ 子节点之下,使得统一访问器 $x.\mathrm{result}$ 产生 $\mathcal{T}(V)$ 本身,即翻译后的值。这个额外的包装在语义上是惰性的,但与证明相关:它使得用于应用的同一 $\mathrm{result}$-追踪论证在充分性证明中能均匀地扩展到值 $\mathbf{let}$-绑定(附录~\ref{app:bohm-tree-proofs})。在尾调用规则中,新标签 $\mathrm{tailCall}$ 起到相同作用,而 $\mathrm{result}$ 投影出答案。在应用和尾调用两种情形中,执行应用的继承被封装在一个命名属性之后,使其内部结构(被调用者的 $\mathrm{argument}$ 和 $\mathrm{result}$ 标签)不会泄漏到封闭作用域。

以上六条规则构成了从 $\lambda$-演算(经由 ANF)到继承演算的完整翻译:每个 ANF 项都有一个像,且构造是复合的。由于 $\lambda$-演算是图灵完备的,继承演算也是。%
\footnote{
  MIXINv2 将符号解析(解析词法引用和限定 $\this$ 引用)与运行时 mixin 求值分离,但符号解析是逐符号即时执行的,而非提前执行,因此它不提供传统意义上的静态类型检查。尽管如此,MIXINv2 的解析阶段比无类型的继承演算更严格。本文中的所有示例均从 MIXINv2 移植而来,并同时符合两种语义。为满足解析器而添加的元素对此处呈现的算法无害,且可能有助于可读性。
}
然而,图灵完备性并未说明翻译保留了什么;下一小节将精确刻画语义对应关系。
\else
A $\lambda$-bound variable $x$ with de~Bruijn index $n$
(the number of enclosing $\lambda$s between
the reference and its binder) becomes
$\dbi{n'}.\mathrm{argument}$,
which reaches the abstraction that binds $x$ and
accesses its $\mathrm{argument}$ slot.%
\footnote{\label{fn:debruijn-adjust}In the $\lambda$-calculus, de~Bruijn
  indices count only $\lambda$-abstractions.
  In inheritance-calculus, de~Bruijn indices count all
  enclosing scope levels, including those introduced by
  $\mathbf{let}$-bindings and tail calls.
  The translation therefore adjusts the index:
  if a $\lambda$-bound variable crosses $k$ intervening
  $\mathbf{let}$-binding or tail-call scopes
  (that is, $k$ counts the $\mathbf{let}$- and tail-call
    constructs on the path from the variable occurrence
  to its binding $\lambda$ in the ANF syntax tree),
  its inheritance-calculus
  index is $n' = n + k$,
  not~$n$.
  Let-bound variables are unaffected, as they use
  lexical references (such as $x.\mathrm{result}$) rather than
  de~Bruijn indices.
  We assume that the synthetic labels introduced by the
  translation ($\mathrm{argument}$, $\mathrm{result}$,
  $\mathrm{tailCall}$) and the user-defined labels
  (let-bound names such as~$x$) are drawn from disjoint
  alphabets, so no collision can arise.
}
A let-bound variable $x$ becomes the lexical reference
$x.\mathrm{result}$, projecting the application result
from the sibling property $x$ in the enclosing record.

An abstraction $\lambda x.\, M$ translates to a record with own
properties $\{\mathrm{argument}, \mathrm{result}\}$; we call this
the \emph{abstraction shape}.
In the application $\mathbf{let}$-binding rule, the name $x$
binds to the inherited record
$\{\mathcal{T}(V_1),\;\mathrm{argument} \mapsto \mathcal{T}(V_2)\}$,
and $\mathcal{T}(M)$ may reference $x.\mathrm{result}$ to obtain
the value of the application.
In the value $\mathbf{let}$-binding rule, $x$ wraps $\mathcal{T}(V)$
under a $\mathrm{result}$ child so that the uniform accessor
$x.\mathrm{result}$ yields $\mathcal{T}(V)$, the translated value
itself.
This extra wrapper is semantically inert but proof-relevant:
it lets the same $\mathrm{result}$-following argument used for
applications extend uniformly to value
$\mathbf{let}$-bindings in the adequacy proof
(Appendix~\ref{app:bohm-tree-proofs}).
In the tail-call rule, the fresh label $\mathrm{tailCall}$ serves
the same purpose, and $\mathrm{result}$ projects the answer.
In the application and tail-call cases, the inheritance that
performs the application is
encapsulated behind a named property, so its internal
structure (the $\mathrm{argument}$ and $\mathrm{result}$ labels
of the callee) does not leak to the enclosing scope.

The six rules above are a complete translation from the
$\lambda$-calculus (via ANF) to inheritance-calculus: every ANF
term has an image, and the construction is compositional.
Since the $\lambda$-calculus is Turing complete, so is
inheritance-calculus.%
\footnote{
  MIXINv2 separates symbol resolution (which resolves
  lexical references and qualified $\this$ references) from runtime
  mixin evaluation, but symbol resolution is performed
  just-in-time per symbol rather than ahead of time, so it does
  not provide static type checking in the traditional sense.
  Nevertheless, MIXINv2's resolution phase is stricter than the
  untyped inheritance-calculus.
  All examples in this paper are ported from MIXINv2
  and conform to both semantics.
  Elements added to satisfy the resolver are harmless
  for the algorithms presented here and may aid readability.
}
Turing completeness, however, says nothing about what the
translation preserves; the next subsection pins down
the semantic correspondence.
\fi

\subsection{\bilingual{B\"ohm Tree Correspondence}{Böhm 树对应}}
\label{sec:bohm-tree}

\ifInheritanceChinese
翻译 $\mathcal{T}$ 将 $\lambda$-演算嵌入继承演算的一个子语言中。直觉上,mixin 树以附加的语义结构增广了 AST;对于翻译后的 $\lambda$-项,这一附加结构的投影与 Böhm 树同构。若从根出发追踪有限次 $\mathrm{result}$ 投影后到达一个同时拥有 $\mathrm{argument}$ 和 $\mathrm{result}$ 的节点(即 $\mathcal{T}(\lambda x.\, M)$ 产生的\emph{抽象形状}),则称 mixin 树 $T$ \emph{收敛},记为 $T{\Downarrow}$。翻译保持并反映收敛性(充分性)。记 $\mathcal{T}(M) \approx_\lambda \mathcal{T}(N)$ 当对所有封闭的 $\lambda$-演算上下文 $C[\cdot]$ 均有 $\mathcal{T}(C[M]){\Downarrow} \Leftrightarrow \mathcal{T}(C[N]){\Downarrow}$;则 $\approx_\lambda$ 恰好与 Böhm 树等价~\cite{barendregt1984-lambda-calculus} 重合。形式化陈述和证明见附录~\ref{app:bohm-tree-proofs}。注意 $\approx_\lambda$ 仅在 $\lambda$-演算上下文上量化;继承演算的上下文区分能力严格更强,如第~\ref{sec:asymmetry} 节所证。以上对应关系在首次迭代近似 $T_P\!\uparrow\!1$ 处成立;我们现在考察在完整最小不动点处会发生什么。

在 $T_P\!\uparrow\!1$ 处,每个方程被求值一次,可重入调用返回 $\bot$;这对应于 $\beta$-规约。在完整的 $\mathrm{lfp}(T_P)$ 处,可重入引用改而解析为幂集格上的最小不动点:迭代在\emph{容器}层级计算不动点,求出满足 $S = F(S)$ 的集合 $S$。这正是递归 Datalog 在循环图上收敛的机制(附录~\ref{app:datalog-encoding},第~\ref{sec:recursive-rules} 节);既非不动点组合子(如 $Y$),也非 Haskell 的 \texttt{mfix}(\texttt{RecursiveDo})能够表达它:两者都在项层级打结,而非在语义所累积的集合上打结。破坏 Böhm 树对应关系的这种\emph{幂集}更多确定性(图~\ref{fig:calculi-matrix})的最小见证例是 $\mathbf{let}\; x = x \;\mathbf{in}\; x$:它在 $\beta$-规约下发散,但其翻译 $\{x \mapsto \{\mathrm{result} \mapsto \{x.\mathrm{result}\}\},\; \mathrm{result} \mapsto \{x.\mathrm{result}\}\}$ 在 $\mathrm{lfp}(T_P)$ 下稳定,tabling 求值检测到循环并返回确定结果 $\properties(\mathrm{root}) = \{x, \mathrm{result}\}$。因此本小节的对应关系是 $T_P\!\uparrow\!1$ 独有的性质。下一小节考察继承演算是否也严格更具表达力。
\else
The translation $\mathcal{T}$ embeds the $\lambda$-calculus
into a sublanguage of inheritance-calculus.
Intuitively, a mixin tree augments the AST with
additional semantic structure; for translated
$\lambda$-terms, the projection of this additional
structure is isomorphic to the B\"ohm tree.
A mixin tree $T$ \emph{converges}, written $T{\Downarrow}$,
if following finitely many $\mathrm{result}$ projections from
the root reaches a node with both $\mathrm{argument}$ and
$\mathrm{result}$ among its properties, that is, the \emph{abstraction shape}
produced by $\mathcal{T}(\lambda x.\, M)$.
The translation preserves and reflects convergence
(adequacy).
Write $\mathcal{T}(M) \approx_\lambda \mathcal{T}(N)$
when $\mathcal{T}(C[M]){\Downarrow} \Leftrightarrow
\mathcal{T}(C[N]){\Downarrow}$ for every closing
$\lambda$-calculus context $C[\cdot]$;
then $\approx_\lambda$ coincides with B\"ohm tree
equivalence~\cite{barendregt1984-lambda-calculus}.
Formal statements and proofs are in
Appendix~\ref{app:bohm-tree-proofs}.
Note that $\approx_\lambda$ quantifies over $\lambda$-calculus
contexts only; inheritance-calculus contexts are strictly more
discriminating, as Section~\ref{sec:asymmetry} establishes.
The correspondence above holds at the first-iteration
approximation $T_P\!\uparrow\!1$; we now ask what happens
at the full least fixed point.

At $T_P\!\uparrow\!1$, each equation is evaluated once and
re-entrant calls return~$\bot$; this corresponds to
$\beta$-reduction.
At the full $\mathrm{lfp}(T_P)$, re-entrant references instead
resolve to least fixed points over the powerset lattice: the
iteration computes fixed points at the \emph{container} level,
finding a set~$S$ such that $S = F(S)$.
This is the mechanism by which recursive Datalog converges on
cyclic graphs (Appendix~\ref{app:datalog-encoding},
Section~\ref{sec:recursive-rules}), and neither a fixed-point
combinator such as~$Y$ nor Haskell's \texttt{mfix}
(\texttt{RecursiveDo}) expresses it: both tie the knot at the
term level, not over the sets the semantics accumulates.
The smallest witness that this \emph{powerset} more-definedness
(Figure~\ref{fig:calculi-matrix}) breaks the B\"ohm-tree
correspondence is
$\mathbf{let}\; x = x \;\mathbf{in}\; x$:
it diverges under $\beta$-reduction, but its translation
$\{x \mapsto \{\mathrm{result} \mapsto \{x.\mathrm{result}\}\},\;
\mathrm{result} \mapsto \{x.\mathrm{result}\}\}$
stabilizes under $\mathrm{lfp}(T_P)$, where tabled evaluation
detects the cycle and returns the definite result
$\properties(\mathrm{root}) = \{x, \mathrm{result}\}$.
The correspondence of this subsection is therefore a property
of $T_P\!\uparrow\!1$ alone.
The next subsection examines whether inheritance-calculus is
also strictly more expressive.
\fi

\subsection{\bilingual{Expressive Asymmetry}{表达力不对称}}
\label{sec:asymmetry}
\label{sec:formal-separation}
\ifInheritanceChinese
两个 Böhm 树等价的 $\lambda$-项,在 $\lambda$-演算中因此无法区分,可以被一个混合了 $\lambda$-演算与继承演算构造子的上下文分离;该上下文通过开放扩展在已有定义上继承新的可观测投影。根据 Felleisen 的表达力判据~\cite{felleisen1991-expressive-power},继承演算严格比惰性 $\lambda$-演算更具表达力(定理~\ref{thm:expressive-asymmetry}):$\lambda$-演算可以通过 ANF 转换和 $\mathcal{T}$ 宏表达于继承演算中(引理~\ref{lem:anf-equiv},定理~\ref{thm:forward-macro}),但惰性 $\lambda$-演算无法宏表达继承(定理~\ref{thm:cartesian-nonexpressibility})。形式化定义、构造和证明见附录~\ref{app:expressiveness};该分离构造本身是表达式问题的一个微型实例,它在不修改已有定义的情况下添加了一个新的可观测投影。
\else
Two $\lambda$-terms that are B\"ohm-tree equivalent,
and therefore indistinguishable in the $\lambda$-calculus,
can be
separated by a context that mixes $\lambda$-calculus and
inheritance-calculus constructors, inheriting new observable
projections onto an existing definition via open extension.
By Felleisen's expressiveness
criterion~\cite{felleisen1991-expressive-power}, inheritance-calculus is
strictly more expressive than the lazy $\lambda$-calculus
(Theorem~\ref{thm:expressive-asymmetry}):
the $\lambda$-calculus can be macro-expressed in
inheritance-calculus via ANF conversion
and~$\mathcal{T}$
(Lemma~\ref{lem:anf-equiv}, Theorem~\ref{thm:forward-macro}),
but the lazy $\lambda$-calculus cannot macro-express inheritance
(Theorem~\ref{thm:cartesian-nonexpressibility}).
The formal definitions, constructions, and proofs are given in
Appendix~\ref{app:expressiveness};
the separating construction is itself a miniature instance of
the Expression Problem, adding a new observable projection to
an existing definition without modifying it.
\fi

\section{\bilingual{Case Study: Nat Arithmetic}{案例研究：Nat(自然数)算术}}
\label{sec:case-study}
\label{sec:expression-problem}

\ifInheritanceChinese
我们在继承演算中实现了布尔逻辑与自然数算术,包括一元与二进制两种形式。
本节详细追踪一元 Nat 的情形:其数据表示、加法与相等判断。
该实现遵循普通的 Gang-of-Four 设计模式~\cite{gamma1994-design-patterns},采用声明式面向对象风格,将每项关注点沿两个轴分解为 mixin:操作轴与情形轴。由于演算本身没有方法,所有名称均为名词或形容词:UpperCamelCase 名称扮演模块、类型、数据构造子与方法的角色,lowerCamelCase 名称扮演字段与参数的角色。
\else
We implemented boolean logic and natural number arithmetic, both unary
and binary, in inheritance-calculus.
This section traces the unary Nat case in detail: its data representation,
addition, and equality.
The implementation follows ordinary Gang-of-Four design
patterns~\cite{gamma1994-design-patterns} in a declarative object-oriented
style, with each concern split into mixins along two axes: operations and
cases. Because the calculus has no methods of its own, every name is a noun or
adjective: UpperCamelCase names play the roles of modules, types, data
constructors, and methods, and lowerCamelCase names the roles of fields and
parameters.
\fi

\paragraph{NatData}
\ifInheritanceChinese
mixin $\mathrm{NatData}$ 是自然数接口及其工厂。$\mathrm{NatFactory}$ 是一个工厂,其抽象产品类型被别名为 $\mathrm{Nat}$:
\else
The mixin $\mathrm{NatData}$ is the natural-number interface and its factory.
$\mathrm{NatFactory}$ is a factory whose abstract product type is aliased as $\mathrm{Nat}$:
\fi
\[
  \mathrm{NatData} \mapsto \left\{
    \begin{aligned}
      &\mathrm{NatFactory} \mapsto \{\mathrm{Product} \mapsto \{\}\},\\
      &\mathrm{Nat} \mapsto \{\mathrm{NatFactory}.\mathrm{Product}\}
  \end{aligned}\right\}
\]

\paragraph{NatDataZero}
\ifInheritanceChinese
每个数据构造子都是一个独立的实现 mixin,继承 $\mathrm{NatData}$。一个 $\mathrm{Zero}$ 值继承 $\mathrm{Product}$:
\else
Each data constructor is a separate implementation mixin that inherits
$\mathrm{NatData}$.
A $\mathrm{Zero}$ value inherits $\mathrm{Product}$:
\fi
\[
  \mathrm{NatDataZero} \mapsto \left\{
    \begin{aligned}
      &\mathrm{NatData},\\
      &\mathrm{NatFactory} \mapsto \{\mathrm{Zero} \mapsto \{\qualifiedthis{\mathrm{NatFactory}}.\mathrm{Product}\}\}
  \end{aligned}\right\}
\]

\paragraph{NatDataSuccessor}
\ifInheritanceChinese
一个 $\mathrm{Successor}$ 值继承 $\mathrm{Product}$,并暴露一个类型为 $\mathrm{Product}$ 的 $\mathrm{predecessor}$ 属性:
\else
A $\mathrm{Successor}$ value inherits $\mathrm{Product}$ and
exposes a $\mathrm{predecessor}$ property whose type is $\mathrm{Product}$:
\fi
\[
  \begin{aligned}
    &\mathrm{NatDataSuccessor} \mapsto \\
    &\left\{
    \begin{aligned}
      &\mathrm{NatData},\\
      &\mathrm{NatFactory} \mapsto \\
      &\left\{\mathrm{Successor} \mapsto \left\{
        \begin{aligned}
          &\qualifiedthis{\mathrm{NatDataSuccessor}}.\mathrm{NatFactory}.\mathrm{Product},\\
          &\mathrm{predecessor} \mapsto \{\qualifiedthis{\mathrm{NatDataSuccessor}}.\mathrm{NatFactory}.\mathrm{Product}\}
      \end{aligned}\right\}\right\}
  \end{aligned}\right\}
  \end{aligned}
\]
\ifInheritanceChinese
与 $\mathrm{NatDataZero}$ 合在一起,这以开放的方式完成了自然数的 Peano 编码,即函子 $F(X) = 1 + X$ 的初始代数:两个构造子作为独立的 mixin 加入,而非固定在单个封闭的数据类型定义中。
\else
Together with $\mathrm{NatDataZero}$, this completes the Peano encoding of
the natural numbers, the initial algebra of the functor $F(X) = 1 + X$,
in an open way: the two constructors are added as independent mixins
rather than fixed in a single closed datatype definition.
\fi

\paragraph{NatPlus}
\ifInheritanceChinese
加法操作沿情形轴分解为一个\emph{接口} mixin $\mathrm{NatPlus}$,该接口 mixin 保留关注点的名称并在抽象 $\mathrm{Product}$ 上声明命令,以及每个情形对应一个\emph{实现} mixin,实现 mixin 在名称后附加构造子名称(即下文的 $\mathrm{NatPlusZero}$ 与 $\mathrm{NatPlusSuccessor}$)。$\mathrm{NatPlus}$ 继承 $\mathrm{NatData}$。面向对象方法会用动词命名的地方,命令模式~\cite{gamma1994-design-patterns} 将其具化为以名词命名的命令,因此加法方法被命名为 $\mathrm{Plus}$ 而非 $\mathrm{Add}$,并通过投影属性 $\mathrm{sum}$ 暴露其结果;这里已指定类型,但尚未计算:
\else
The addition operation is split along the case axis into an \emph{interface}
mixin $\mathrm{NatPlus}$, which keeps the concern's name and declares the command
on the abstract $\mathrm{Product}$, and one \emph{implementation} mixin per case,
which appends the constructor name ($\mathrm{NatPlusZero}$ and
$\mathrm{NatPlusSuccessor}$ below). $\mathrm{NatPlus}$ inherits
$\mathrm{NatData}$. Where an
object-oriented method would be named with a verb, the command
pattern~\cite{gamma1994-design-patterns} reifies it as a noun-named command, so
the addition method is named $\mathrm{Plus}$ rather than $\mathrm{Add}$, and it
exposes its result through the projection property $\mathrm{sum}$, here typed but
not yet computed:
\fi
\[
  \mathrm{NatPlus} \mapsto \left\{
    \begin{aligned}
      &\mathrm{NatData},\\
      &\mathrm{NatFactory} \mapsto \left\{\mathrm{Product} \mapsto \left\{\mathrm{Plus} \mapsto \left\{\mathrm{sum} \mapsto \{\mathrm{Product}\}\right\}\right\}\right\}
  \end{aligned}\right\}
\]

\paragraph{NatPlusZero}
\ifInheritanceChinese
基础情形为 $0 + m = m$。这里 $\mathrm{NatPlusZero}$ 充当一个导入其他模块的模块:它继承 $\mathrm{Zero}$ 数据与 $\mathrm{Plus}$ 接口,并直接返回被加数:
\else
The base case is $0 + m = m$. Here $\mathrm{NatPlusZero}$ acts as a module that
imports other modules: it inherits the $\mathrm{Zero}$ data and the $\mathrm{Plus}$
interface, and returns the addend directly:
\fi
\[
  \begin{aligned}
    &\mathrm{NatPlusZero} \mapsto \\
    &\left\{
    \begin{aligned}
      &\mathrm{NatDataZero},\; \mathrm{NatPlus},\\
      &\mathrm{NatFactory} \mapsto \left\{\mathrm{Zero} \mapsto \left\{\mathrm{Plus} \mapsto \left\{
          \begin{aligned}
            &\mathrm{addend} \mapsto \{\qualifiedthis{\mathrm{NatPlusZero}}.\mathrm{NatFactory}.\mathrm{Product}\},\\
            &\mathrm{sum} \mapsto \{\mathrm{addend}\}
      \end{aligned}\right\}\right\}\right\}
  \end{aligned}\right\}
  \end{aligned}
\]

\paragraph{NatPlusSuccessor}
\ifInheritanceChinese
递归情形将 $\mathrm{S}(n_0) + m$ 规约为 $n_0 + \mathrm{S}(m)$,方法是以递增后的被加数委托给 $n_0$ 的 $\mathrm{Plus}$:
\else
The recursive case reduces $\mathrm{S}(n_0) + m$ to $n_0 + \mathrm{S}(m)$
by delegating to $n_0$'s $\mathrm{Plus}$ with an incremented addend:
\fi
\[
  \begin{aligned}
    &\mathrm{NatPlusSuccessor} \mapsto \\
    &\left\{
    \begin{aligned}
      &\mathrm{NatDataSuccessor},\; \mathrm{NatPlus},\\
      &\mathrm{NatFactory} \mapsto \\
      &\left\{\mathrm{Successor} \mapsto \left\{\mathrm{Plus} \mapsto \left\{
              \begin{aligned}
                &\mathrm{addend} \mapsto \{\qualifiedthis{\mathrm{NatFactory}}.\mathrm{Product}\},\\
                &\mathrm{\_increasedAddend} \mapsto  \\
                &\left\{
                  \begin{aligned}
                    &\mathrm{Successor},\\
                    &\mathrm{predecessor} \mapsto \{\mathrm{addend}\}
                \end{aligned}\right\},\\
                &\mathrm{\_recursiveAddition} \mapsto \\
                &\left\{
                  \begin{aligned}
                    &\qualifiedthis{\mathrm{Successor}}.\mathrm{predecessor}.\mathrm{Plus},\\
                    &\mathrm{addend} \mapsto \{\mathrm{\_increasedAddend}\}
                \end{aligned}\right\},\\
                &\mathrm{sum} \mapsto \{\mathrm{\_recursiveAddition}.\mathrm{sum}\}
          \end{aligned}\right\}\right\}\right\}
  \end{aligned}\right\}
  \end{aligned}
\]

\paragraph{NatVisitor}
\ifInheritanceChinese
相等判断需要对数据构造子进行情形分析;我们使用访问者模式~\cite{gamma1994-design-patterns}。$\mathrm{NatVisitor}$ 向 $\mathrm{Product}$ 添加一个 $\mathrm{Acceptance}$ 命令,即访问者模式的 $\mathrm{accept}$ 方法的名词形式。mixin $\mathrm{NatVisitor}$ 将其声明为空,留下 $\mathrm{VisitorMap}$ 与 $\mathrm{Accepted}$ 属性由各构造子的实现来填充:
\else
Equality requires case analysis on the data constructor; we use the
visitor pattern~\cite{gamma1994-design-patterns}.
$\mathrm{NatVisitor}$ adds an $\mathrm{Acceptance}$ command to $\mathrm{Product}$,
the noun form of the visitor pattern's $\mathrm{accept}$ method.
The mixin $\mathrm{NatVisitor}$ declares it empty, leaving a $\mathrm{VisitorMap}$ and an
$\mathrm{Accepted}$ property for the per-constructor implementations to fill in:
\fi
\[
  \begin{aligned}
    &\mathrm{NatVisitor} \mapsto \\
    &\left\{
    \begin{aligned}
      &\mathrm{NatData},\\
      &\mathrm{NatFactory} \mapsto \left\{\mathrm{Product} \mapsto \{\mathrm{Acceptance} \mapsto \{\mathrm{VisitorMap} \mapsto \{\},\; \mathrm{Accepted} \mapsto \{\}\}\}\right\}
  \end{aligned}\right\}
  \end{aligned}
\]

\paragraph{NatVisitorZero}
\ifInheritanceChinese
$\mathrm{Zero}$ 的实现选择 $\mathrm{ZeroVisitor}$ 分支。这里的访问者模式与~\cite{gamma1994-design-patterns} 中的表述不同,那里数据构造子由 $\mathrm{accept}$ 接口固定;此处通过组合允许添加新的构造子,因为每个构造子作为独立的 mixin 贡献其自身的分支:
\else
The $\mathrm{Zero}$ implementation selects the $\mathrm{ZeroVisitor}$ branch.
The visitor pattern here, unlike its formulation in
\cite{gamma1994-design-patterns} where the data constructors are fixed by the
$\mathrm{accept}$ interface, admits new constructors by composition, since each
constructor contributes its own branch as a separate mixin:
\fi
\[
  \begin{aligned}
    &\mathrm{NatVisitorZero} \mapsto \\
    &\left\{
    \begin{aligned}
      &\mathrm{NatDataZero},\; \mathrm{NatVisitor},\\
      &\mathrm{NatFactory} \mapsto \left\{\mathrm{Zero} \mapsto \left\{\mathrm{Acceptance} \mapsto \left\{
                \begin{aligned}
                  &\mathrm{VisitorMap} \mapsto \{\mathrm{ZeroVisitor} \mapsto \{\}\},\\
                  &\mathrm{Accepted} \mapsto \{\mathrm{VisitorMap}.\mathrm{ZeroVisitor}\}
              \end{aligned}\right\}\right\}\right\}
  \end{aligned}\right\}
  \end{aligned}
\]

\paragraph{NatVisitorSuccessor}
\ifInheritanceChinese
$\mathrm{Successor}$ 的实现选择 $\mathrm{SuccessorVisitor}$ 分支:
\else
The $\mathrm{Successor}$ implementation selects the $\mathrm{SuccessorVisitor}$ branch:
\fi
\[
  \begin{aligned}
    &\mathrm{NatVisitorSuccessor} \mapsto \\
    &\left\{
    \begin{aligned}
      &\mathrm{NatDataSuccessor},\; \mathrm{NatVisitor},\\
      &\mathrm{NatFactory} \mapsto \\
      &\left\{\mathrm{Successor} \mapsto \left\{
            \begin{aligned}
              &\mathrm{predecessor} \mapsto \{\qualifiedthis{\mathrm{NatFactory}}.\mathrm{Product}\},\\
              &\mathrm{Acceptance} \mapsto \left\{
                \begin{aligned}
                  &\mathrm{VisitorMap} \mapsto \{\mathrm{SuccessorVisitor} \mapsto \{\}\},\\
                  &\mathrm{Accepted} \mapsto \{\mathrm{VisitorMap}.\mathrm{SuccessorVisitor}\}
              \end{aligned}\right\}
          \end{aligned}\right\}\right\}
  \end{aligned}\right\}
  \end{aligned}
\]

\paragraph{BooleanData}
\ifInheritanceChinese
$\mathrm{BooleanData}$ 定义输出的域:
\else
$\mathrm{BooleanData}$ defines the output sort:
\fi
\[
  \mathrm{BooleanData} \mapsto \left\{
    \begin{aligned}
      &\mathrm{BooleanFactory} \mapsto \left\{
        \begin{aligned}
          &\mathrm{Product} \mapsto \{\},\\
          &\mathrm{True} \mapsto \{\mathrm{Product}\},\\
          &\mathrm{False} \mapsto \{\mathrm{Product}\}
      \end{aligned}\right\},\\
      &\mathrm{Boolean} \mapsto \{\mathrm{BooleanFactory}.\mathrm{Product}\}
  \end{aligned}\right\}
\]
\paragraph{NatEquality}
\label{sec:nat-equality}
\ifInheritanceChinese
为了验证 $2 + 3 = 5$,我们需要对 Nat 值进行相等判断。mixin $\mathrm{NatEquality}$ 导入 $\mathrm{NatVisitor}$ 与 $\mathrm{BooleanData}$,并在 $\mathrm{Product}$ 上声明一个 $\mathrm{Equal}$ 命令,由各构造子的实现提供。由于 $\mathrm{NatEquality}$ 继承 $\mathrm{BooleanData}$,限定 $\this$ $\qualifiedthis{\mathrm{NatEquality}}$ 可以到达 $\mathrm{Boolean}$,因此 $\mathrm{Equal}$ 以类型 $\qualifiedthis{\mathrm{NatEquality}}.\mathrm{Boolean}$ 声明其结果 $\mathrm{equal}$:
\else
To test that $2 + 3 = 5$, we needed equality on Nat values.
The mixin $\mathrm{NatEquality}$ imports $\mathrm{NatVisitor}$ and $\mathrm{BooleanData}$,
and declares an $\mathrm{Equal}$ command on $\mathrm{Product}$ that the per-constructor implementations provide.
Because $\mathrm{NatEquality}$ inherits $\mathrm{BooleanData}$,
the qualified this $\qualifiedthis{\mathrm{NatEquality}}$ can reach $\mathrm{Boolean}$,
so $\mathrm{Equal}$ declares its result
$\mathrm{equal}$ with the type $\qualifiedthis{\mathrm{NatEquality}}.\mathrm{Boolean}$:
\fi
\[
  \begin{aligned}
    &\mathrm{NatEquality} \mapsto \\
    &\left\{
    \begin{aligned}
      &\mathrm{NatVisitor},\; \mathrm{BooleanData},\\
      &\mathrm{NatFactory} \mapsto \left\{\mathrm{Product} \mapsto \left\{\mathrm{Equal} \mapsto \left\{
          \begin{aligned}
            &\mathrm{other} \mapsto \{\mathrm{Product}\},\\
            &\mathrm{equal} \mapsto \{\qualifiedthis{\mathrm{NatEquality}}.\mathrm{Boolean}\}
      \end{aligned}\right\}\right\}\right\}
  \end{aligned}\right\}
  \end{aligned}
\]

\paragraph{NatEqualityZero}
\ifInheritanceChinese
$\mathrm{Zero}$ 只与另一个 $\mathrm{Zero}$ 相等。这里的继承执行情形分析:继承 $\mathrm{other}.\mathrm{Acceptance}$ 对 $\mathrm{other}$ 是由 $\mathrm{Zero}$ 还是 $\mathrm{Successor}$ 构建进行分派,在 $\mathrm{ZeroVisitor}$ 分支返回 $\mathrm{True}$,否则返回 $\mathrm{False}$,即 $\mathrm{equal} = \{ \mathrm{True} \text{ 若 } \mathrm{other} = \mathrm{Zero},\; \mathrm{False} \text{ 若 } \mathrm{other} = \mathrm{Successor} \}$。$\mathrm{True}$ 与 $\mathrm{False}$ 均为 $\mathrm{BooleanFactory}$ 的数据构造子;该工厂来自 $\mathrm{BooleanData}$,通过继承的 $\mathrm{NatEquality}$ 间接导入,并以 $\qualifiedthis{\mathrm{NatEqualityZero}}.\mathrm{BooleanFactory}$ 到达。同一继承机制也提供了结果类型:在 $\mathrm{Accepted}$ 中,$\mathrm{equal}$ 继承 $\qualifiedthis{\mathrm{NatEqualityZero}}.\mathrm{Boolean}$,从而声明其类型为 $\mathrm{Boolean}$:
\else
$\mathrm{Zero}$ is equal only to another $\mathrm{Zero}$.
Inheriting here performs case analysis:
inheriting $\mathrm{other}.\mathrm{Acceptance}$ dispatches on whether $\mathrm{other}$
was built by $\mathrm{Zero}$ or $\mathrm{Successor}$, returning $\mathrm{True}$
on the $\mathrm{ZeroVisitor}$ branch and $\mathrm{False}$ otherwise, that is,
$\mathrm{equal} = \{ \mathrm{True} \text{ if } \mathrm{other} = \mathrm{Zero},\; \mathrm{False} \text{ if } \mathrm{other} = \mathrm{Successor} \}$.
Both $\mathrm{True}$ and $\mathrm{False}$ are data constructors of
$\mathrm{BooleanFactory}$, which comes from $\mathrm{BooleanData}$, imported
indirectly through the inherited $\mathrm{NatEquality}$, and is reached as
$\qualifiedthis{\mathrm{NatEqualityZero}}.\mathrm{BooleanFactory}$.
The same inheritance mechanism also supplies the result type: in
$\mathrm{Accepted}$, $\mathrm{equal}$ inherits
$\qualifiedthis{\mathrm{NatEqualityZero}}.\mathrm{Boolean}$, declaring its type to
be $\mathrm{Boolean}$:
\fi
{\small
  \[
    \begin{aligned}
      &\mathrm{NatEqualityZero} \mapsto \\
      &\left\{
      \begin{aligned}
        &\mathrm{NatVisitorZero},\; \mathrm{NatEquality},\\
        &\mathrm{NatFactory} \mapsto \\
        &\left\{
        \begin{aligned}
        &\mathrm{Zero} \mapsto \\
        &\left\{
        \begin{aligned}
        &\mathrm{Equal} \mapsto \\
        &\left\{
                  \begin{aligned}
                    &\mathrm{other} \mapsto \{\qualifiedthis{\mathrm{NatEqualityZero}}.\mathrm{NatFactory}.\mathrm{Product}\},\\
                    &\mathrm{\_OtherAcceptance} \mapsto \\
                    &\left\{
                      \begin{aligned}
                        &\mathrm{other}.\mathrm{Acceptance},\\
                        &\mathrm{VisitorMap} \mapsto \\
                        &\left\{
                          \begin{aligned}
                            &\mathrm{ZeroVisitor} \mapsto \{\mathrm{equal} \mapsto \{\qualifiedthis{\mathrm{NatEqualityZero}}.\mathrm{BooleanFactory}.\mathrm{True}\}\},\\
                            &\mathrm{SuccessorVisitor} \mapsto \{\mathrm{equal} \mapsto \{\qualifiedthis{\mathrm{NatEqualityZero}}.\mathrm{BooleanFactory}.\mathrm{False}\}\}
                        \end{aligned}\right\},\\
                        &\mathrm{Accepted} \mapsto \{\mathrm{equal} \mapsto \{\qualifiedthis{\mathrm{NatEqualityZero}}.\mathrm{Boolean}\}\}
                    \end{aligned}\right\},\\
                    &\mathrm{equal} \mapsto \{\mathrm{\_OtherAcceptance}.\mathrm{Accepted}.\mathrm{equal}\}
                \end{aligned}\right\}
        \end{aligned}\right\}
        \end{aligned}\right\}
    \end{aligned}\right\}
    \end{aligned}
\]}

\paragraph{NatEqualitySuccessor}
\ifInheritanceChinese
$\mathrm{Successor}$ 只与另一个具有相等前驱的 $\mathrm{Successor}$ 相等,因此其实现通过 $\qualifiedthis{\mathrm{Successor}}.\mathrm{predecessor}.\mathrm{Equal}$ 进行递归:
\else
$\mathrm{Successor}$ is equal only to another $\mathrm{Successor}$ with an equal
predecessor, so its implementation recurses through
$\qualifiedthis{\mathrm{Successor}}.\mathrm{predecessor}.\mathrm{Equal}$:
\fi
{\small
  \[
    \begin{aligned}
      &\mathrm{NatEqualitySuccessor} \mapsto \\
      &\left\{
      \begin{aligned}
        &\mathrm{NatVisitorSuccessor},\; \mathrm{NatEquality},\\
        &\mathrm{NatFactory} \mapsto \\
        &\left\{
        \begin{aligned}
        &\mathrm{Successor} \mapsto \\
        &\left\{
        \begin{aligned}
        &\mathrm{Equal} \mapsto \\
        &\left\{
                  \begin{aligned}
                    &\mathrm{other} \mapsto \left\{
                      \begin{aligned}
                        &\qualifiedthis{\mathrm{NatEqualitySuccessor}}.\mathrm{NatFactory}.\mathrm{Product},\\
                        &\mathrm{predecessor} \mapsto \{\qualifiedthis{\mathrm{NatEqualitySuccessor}}.\mathrm{NatFactory}.\mathrm{Product}\}
                    \end{aligned}\right\},\\
                    &\mathrm{\_recursiveEquality} \mapsto \left\{
                      \begin{aligned}
                        &\qualifiedthis{\mathrm{Successor}}.\mathrm{predecessor}.\mathrm{Equal},\\
                        &\mathrm{other} \mapsto \{\qualifiedthis{\mathrm{Equal}}.\mathrm{other}.\mathrm{predecessor}\}
                    \end{aligned}\right\},\\
                    &\mathrm{\_OtherAcceptance} \mapsto \\
                    &\left\{
                      \begin{aligned}
                        &\mathrm{other}.\mathrm{Acceptance},\\
                        &\mathrm{VisitorMap} \mapsto \\
                        &\left\{
                          \begin{aligned}
                            &\mathrm{ZeroVisitor} \mapsto \{\mathrm{equal} \mapsto \{\qualifiedthis{\mathrm{NatEqualitySuccessor}}.\mathrm{BooleanFactory}.\mathrm{False}\}\},\\
                            &\mathrm{SuccessorVisitor} \mapsto \{\mathrm{equal} \mapsto \{\mathrm{\_recursiveEquality}.\mathrm{equal}\}\}
                        \end{aligned}\right\},\\
                        &\mathrm{Accepted} \mapsto \{\mathrm{equal} \mapsto \{\qualifiedthis{\mathrm{NatEqualitySuccessor}}.\mathrm{Boolean}\}\}
                    \end{aligned}\right\},\\
                    &\mathrm{equal} \mapsto \{\mathrm{\_OtherAcceptance}.\mathrm{Accepted}.\mathrm{equal}\}
                \end{aligned}\right\}
        \end{aligned}\right\}
        \end{aligned}\right\}
    \end{aligned}\right\}
    \end{aligned}
\]}
%% === 审稿人备忘：关于代码"冗长"的回应 ===
%%
%% NatEquality 的定义看起来比等价的 Haskell 或 ML 模式匹配更长，
%% 但这种冗长与 point-free / 组合子风格的"简洁"形成的对比是表面的。
%% 需要区分两种"长"：
%%
%% 1. 自描述式的长：每个中间步骤都有名字（如 _recursiveEquality、
%%    OtherAcceptance），步骤写得详细。名字本身是零成本的语义锚点。
%% 2. 组合子膨胀式的长：如 SKI bracket abstraction 的指数级膨胀，
%%    在玩具例子里看起来精炼，但在真实程序中复杂度是 λ-calculus 的
%%    多项式甚至指数倍，而且没有涌现出新的计算模式（论文 Related Work
%%    已指出这一点）。
%%
%% Inheritance-calculus 的"冗长"属于第一种，而且是结构上的必要代价：
%% 每个有名字的中间步骤既是代码，也是 mixin 树上可路径寻址、可扩展的
%% 节点。这正是 open extension 的前提——无法扩展匿名的子表达式。
%% 如果加语法糖把名字隐藏（如同 M-表达式之于 S-表达式），反而会
%% 破坏这一性质。Lisp 社区的历史经验表明，S-表达式的"冗长"最终
%% 被证明是优势（代码即数据，宏系统直接操作统一结构），M-表达式
%% 的语法糖从未被真正需要。
%%
%% 在 AI 辅助编程的语境下，这种显式命名的风格尤其有利：
%% _recursiveEquality 这样的名字对 LLM 而言是自带的 chain-of-thought，
%% 比起让 AI 猜测 point-free 管道中第 n 个组合子的意图，
%% 每个自描述的步骤都降低了生成和理解的认知负荷。
%% 考虑到未来越来越多的代码将由 AI 生成和维护，
%% 多写几个 token 的显式名称不是成本，而是投资。
%%
%% 当然，如果确实需要更紧凑的表面语法，为 inheritance-calculus 加一个
%% ANF 预处理器（自动为匿名子表达式生成名字）并不困难，
%% 正如 ANF 变换本身是机械化的。但这是表面语法的选择，
%% 不影响演算本身的语义。
%% ================================================================
\ifInheritanceChinese
通过继承来组合 $\mathrm{Plus}$ 与 $\mathrm{Equal}$ 的实现,无需修改任何 mixin,即可使所得的 Nat 值自动同时具备 $\mathrm{Plus}$ 与 $\mathrm{Equal}$。具体的数字常量存放在一个只继承两个数据构造子的共享 mixin 中:
\else
Composing the $\mathrm{Plus}$ and $\mathrm{Equal}$ implementations by
inheriting them, with no modifications to any mixin, gives the
resulting Nat values both $\mathrm{Plus}$ and $\mathrm{Equal}$
automatically.
The concrete numerals live in a shared mixin that inherits only the
two data constructors:
\fi
\[
  \mathrm{NatConstants} \mapsto \left\{
    \begin{aligned}
      &\mathrm{NatDataZero},\; \mathrm{NatDataSuccessor},\\
      &\mathrm{One} \mapsto \left\{
        \begin{aligned}
          &\qualifiedthis{\mathrm{NatConstants}}.\mathrm{NatFactory}.\mathrm{Successor},\\
          &\mathrm{predecessor} \mapsto \{\qualifiedthis{\mathrm{NatConstants}}.\mathrm{NatFactory}.\mathrm{Zero}\}
      \end{aligned}\right\},\\
      &\mathrm{Two} \mapsto \left\{
        \begin{aligned}
          &\qualifiedthis{\mathrm{NatConstants}}.\mathrm{NatFactory}.\mathrm{Successor},\\
          &\mathrm{predecessor} \mapsto \{\mathrm{One}\}
      \end{aligned}\right\},\\
      &\mathrm{Three} \mapsto \left\{
        \begin{aligned}
          &\qualifiedthis{\mathrm{NatConstants}}.\mathrm{NatFactory}.\mathrm{Successor},\\
          &\mathrm{predecessor} \mapsto \{\mathrm{Two}\}
      \end{aligned}\right\},\\
      &\mathrm{Four} \mapsto \left\{
        \begin{aligned}
          &\qualifiedthis{\mathrm{NatConstants}}.\mathrm{NatFactory}.\mathrm{Successor},\\
          &\mathrm{predecessor} \mapsto \{\mathrm{Three}\}
      \end{aligned}\right\},\\
      &\mathrm{Five} \mapsto \left\{
        \begin{aligned}
          &\qualifiedthis{\mathrm{NatConstants}}.\mathrm{NatFactory}.\mathrm{Successor},\\
          &\mathrm{predecessor} \mapsto \{\mathrm{Four}\}
      \end{aligned}\right\}
  \end{aligned}\right\}
\]
\ifInheritanceChinese
算术测试继承 $\mathrm{NatConstants}$ 以及它所需的操作:
\else
The arithmetic test inherits $\mathrm{NatConstants}$ together with the operations it needs:
\fi
\[
  \mathrm{Test} \mapsto \left\{
    \begin{aligned}
      &\mathrm{NatConstants},\; \mathrm{NatPlusZero},\; \mathrm{NatPlusSuccessor},\\
      &\mathrm{NatEqualityZero},\; \mathrm{NatEqualitySuccessor},\\
      &\mathrm{TwoPlusThree} \mapsto \left\{
        \begin{aligned}
          &\qualifiedthis{\mathrm{Test}}.\mathrm{Two}.\mathrm{Plus},\\
          &\mathrm{addend} \mapsto \{\qualifiedthis{\mathrm{Test}}.\mathrm{Three}\}
      \end{aligned}\right\},\\
      &\mathrm{FiveEqualTwoPlusThree} \mapsto \left\{
        \begin{aligned}
          &\qualifiedthis{\mathrm{Test}}.\mathrm{Five}.\mathrm{Equal},\\
          &\mathrm{other} \mapsto \{\mathrm{TwoPlusThree}.\mathrm{sum}\}
      \end{aligned}\right\}
  \end{aligned}\right\}
\]
\ifInheritanceChinese
路径 $\mathrm{Test}.\mathrm{FiveEqualTwoPlusThree}.\mathrm{equal}$ 继承 $\mathrm{True}$。%
\footnote{
  读取一个结果,需要借助 FFI 来打印,或者借助元语言观察以在该路径上调用 $\supers$。
}
没有任何现有 mixin 被修改。$\mathrm{NatEquality}$(一项新操作)与 $\mathrm{Boolean}$(一种新数据类型)各自只需要一个新 mixin:沿操作轴与情形轴均可扩展。
\else
The path $\mathrm{Test}.\mathrm{FiveEqualTwoPlusThree}.\mathrm{equal}$
inherits $\mathrm{True}$.%
\footnote{
  Reading a result requires either the FFI to print it or a
  metalanguage observation that invokes $\supers$ on the path.
}
No existing mixin was modified.
$\mathrm{NatEquality}$, a new operation, and $\mathrm{Boolean}$,
a new data type, each required only a new mixin: extensibility along both
the operation axis and the case axis.
\fi

\paragraph{\bilingual{Tree-level inheritance and return types}{树层次的继承与返回类型}}
\label{sec:discussion-expression-problem}
\ifInheritanceChinese
每个操作都是一个独立的 mixin,扩展工厂的子树。继承递归地合并共享同一标签的子树,因此工厂中的每个标签都从所有被继承的 mixin 获得全部操作。关键在于,这同样适用于返回类型:$\mathrm{NatPlus}$ 从 $\mathrm{NatFactory}$ 返回值,而 $\mathrm{NatEquality}$ 独立地向同一工厂添加 $\mathrm{Equal}$。同时继承两者,使得 $\mathrm{Plus}$ 返回的值自动携带 $\mathrm{Equal}$,无需修改任何一个 mixin。在对象代数~\cite{oliveira2012-object-algebras} 与 finally tagless 解释器~\cite{carette2009-finally-tagless} 中,独立定义的操作不会自动丰富彼此的返回类型;需要额外的样板代码或类型类机制。

多目标限定 $\this$(方程~\ref{eq:this})至关重要。在上述测试 mixin 中,$\mathrm{NatPlusSuccessor}$ 与 $\mathrm{NatEqualitySuccessor}$ 均继承 $\mathrm{NatDataSuccessor}$。当两个操作被组合时,所得的 $\mathrm{Successor}$ 通过两条独立路由继承 $\mathrm{NatDataSuccessor}$ 的模式。每个操作内部的引用 $\qualifiedthis{\mathrm{Successor}}.\mathrm{predecessor}$ 通过两条路由解析,由于继承是幂等的,所有路由贡献相同的属性。在 Scala 中,同样的模式,即两个独立定义的内部类扩展同一外部类的特质,会触发``冲突基类型''拒绝(附录~\ref{app:scala-multi-target}),因为 DOT 的 self 变量是单值的,无法同时通过多条继承路由解析。表达式问题在 Scala 中仍可使用对象代数或类似模式求解,但组合独立定义的操作本质上会产生上述对共享模式的多目标 $\this$ 解析。
\else
Each operation is an independent mixin that extends a factory's
subtree.
Inheritance recursively merges subtrees that share a label,
so every label in the factory acquires all operations from all
inherited mixins.
Crucially, this applies to return types: $\mathrm{NatPlus}$ returns
values from $\mathrm{NatFactory}$, and $\mathrm{NatEquality}$
independently adds $\mathrm{Equal}$ to the same factory.
Inheriting both causes the values returned by $\mathrm{Plus}$ to
carry $\mathrm{Equal}$ automatically, without modifying either mixin.
In object algebras~\cite{oliveira2012-object-algebras} and finally
tagless interpreters~\cite{carette2009-finally-tagless}, independently defined
operations do not automatically enrich each other's return types;
additional boilerplate or type-class machinery is required.

Multi-target qualified $\this$ (equation~\ref{eq:this}) is essential.
In the test mixin above, $\mathrm{NatPlusSuccessor}$ and
$\mathrm{NatEqualitySuccessor}$ both inherit
$\mathrm{NatDataSuccessor}$.
When the two operations are composed, the resulting
$\mathrm{Successor}$ inherits the $\mathrm{NatDataSuccessor}$ schema
through two independent routes.
The reference
$\qualifiedthis{\mathrm{Successor}}.\mathrm{predecessor}$
inside each operation resolves through both routes, and since
inheritance is idempotent, all routes contribute the same
properties.
In Scala, the same pattern, where two independently defined inner
classes extend the same outer class's trait, triggers a
``conflicting base types'' rejection
(Appendix~\ref{app:scala-multi-target}), because DOT's self variable
is single-valued and cannot resolve through multiple inheritance
routes simultaneously.
The Expression Problem can still be solved in Scala using
object algebras or similar patterns, but composing independently
defined operations inherently produces the multi-target $\this$
resolution to shared schemas just described.
\fi

\paragraph{\bilingual{Church-encoded Nats are tries}{Church 编码的 Nat 是字典树}}
\label{sec:case-study-cartesian}
\ifInheritanceChinese
Church 编码的 Nat 是一个记录,其结构映射了构建它所用的数据构造子路径:$\mathrm{Zero}$ 是一个平坦记录;$\mathrm{S}(\mathrm{Zero})$ 是一个带有 $\mathrm{predecessor}$ 属性、指向一个 $\mathrm{Zero}$ 记录的记录;依此类推。这恰好是字典树的形状。%
\footnote{
  一元 Nat 的字典树可以被视为一个稀疏链表,其深度等于它所代表的数值:在字典树并中,链表中的某些 $\mathrm{Successor}$ 节点指向 $\mathrm{Zero}$(叶节点),因此链表中并非每个位置都被占用。\anon[{} 补充材料]{{} MIXINv2 源码~\cite{mixin2025}}{} 中包含一个字典树深度为对数级的二进制自然数(BinNat)算术。
}
一个同时继承两个 Nat 值的作用域是一个字典树并。下面的测试将 $\mathrm{Plus}$ 应用于 $\{1,2\}$ 的字典树并,被加数为 $\{3,4\}$ 的字典树并,然后用 $\mathrm{Equal}$ 将结果与自身比较:
\else
A Church-encoded Nat is a record whose structure mirrors the
data-constructor path used to build it: $\mathrm{Zero}$ is a flat record;
$\mathrm{S}(\mathrm{Zero})$ is a record with a $\mathrm{predecessor}$
property pointing to a $\mathrm{Zero}$ record; and so on.
This is exactly the shape of a trie.%
\footnote{
  The unary Nat trie can be viewed as a sparse linked list
  whose depth equals the number it represents:
  in a trie union, some $\mathrm{Successor}$ nodes along
  the list point to $\mathrm{Zero}$ (leaf nodes),
  so not every position in the list is occupied.
  The\anon[{} supplementary material]{{} MIXINv2 source~\cite{mixin2025}}{} includes a binary natural number
  (BinNat) arithmetic whose trie depth is logarithmic.
}
A scope that simultaneously inherits two Nat values is a trie union.
The following test applies $\mathrm{Plus}$ to a trie union of $\{1,2\}$
with addend a trie union of $\{3,4\}$, then compares the result to itself
with $\mathrm{Equal}$:
\fi
\[
  \mathrm{CartesianTest} \mapsto \left\{
    \begin{aligned}
      &\mathrm{NatConstants},\; \mathrm{NatPlusZero},\; \mathrm{NatPlusSuccessor},\\
      &\mathrm{NatEqualityZero},\; \mathrm{NatEqualitySuccessor},\\
      &\mathrm{OneOrTwo} \mapsto \left\{\qualifiedthis{\mathrm{CartesianTest}}.\mathrm{One},\; \qualifiedthis{\mathrm{CartesianTest}}.\mathrm{Two}\right\},\\
      &\mathrm{ThreeOrFour} \mapsto \left\{\qualifiedthis{\mathrm{CartesianTest}}.\mathrm{Three},\; \qualifiedthis{\mathrm{CartesianTest}}.\mathrm{Four}\right\},\\
      &\mathrm{Result} \mapsto \left\{\mathrm{OneOrTwo}.\mathrm{Plus},\; \mathrm{addend} \mapsto \{\mathrm{ThreeOrFour}\}\right\},\\
      &\mathrm{Check} \mapsto \left\{\mathrm{Result}.\mathrm{sum}.\mathrm{Equal},\; \mathrm{other} \mapsto \{\mathrm{Result}.\mathrm{sum}\}\right\}
  \end{aligned}\right\}
\]
\ifInheritanceChinese
$\mathrm{Result}.\mathrm{sum}$ 求值为字典树 $\{4,5,6\}$,即所有两两之和的集合。$\mathrm{Equal}$ 随后对从 $\{4,5,6\} \times \{4,5,6\}$ 中取出的每一对进行应用;路径 $\mathrm{CartesianTest}.\mathrm{Check}.\mathrm{equal}$ 同时继承 $\mathrm{True}$ 与 $\mathrm{False}$。笛卡尔积直接来自语义方程。$\mathrm{OneOrTwo}$ 同时继承 $\mathrm{One}$ 与 $\mathrm{Two}$,因此 $\bases(\mathrm{OneOrTwo})$ 包含两条路径。当 $\mathrm{Result}$ 继承 $\mathrm{OneOrTwo}.\mathrm{Plus}$ 时,$\mathrm{Result}$ 的 $\bases$ 包含 $\mathrm{One}$ 与 $\mathrm{Two}$ 各自的 $\mathrm{Plus}$。每个 $\mathrm{Plus}$ 通过 $\this$ 解析其 $\mathrm{addend}$,到达 $\mathrm{ThreeOrFour}$;该值本身是 $\mathrm{Three}$ 与 $\mathrm{Four}$ 的字典树并,因此 $\mathrm{addend}$ 路径的 $\supers$ 包含两个值。由于 $\mathrm{Successor}.\mathrm{Plus}$ 以递增后的被加数递归委托给 $\mathrm{predecessor}.\mathrm{Plus}$,而 $\mathrm{predecessor}$ 也通过字典树并解析,递归在每个字典树节点处分支。因此结果 $\mathrm{sum}$ 收集了通过两个操作数的任意组合所能到达的全部路径,即笛卡尔积。

同样操作单个值的 $\mathrm{Plus}$ 与 $\mathrm{Equal}$ 文件,无需修改,即产生逻辑编程的关系语义:$\overrides$ 与 $\supers$ 收集所有继承路径,每个操作都在其上分布。
\else
$\mathrm{Result}.\mathrm{sum}$ evaluates to the trie $\{4,5,6\}$, the set of all pairwise sums.
$\mathrm{Equal}$ then applies to every pair drawn from $\{4,5,6\} \times \{4,5,6\}$;
the path $\mathrm{CartesianTest}.\mathrm{Check}.\mathrm{equal}$
inherits both $\mathrm{True}$ and $\mathrm{False}$.
The Cartesian product arises directly from the semantic equations.
$\mathrm{OneOrTwo}$ inherits both $\mathrm{One}$ and $\mathrm{Two}$,
so $\bases(\mathrm{OneOrTwo})$ contains two paths.
When $\mathrm{Result}$ inherits
$\mathrm{OneOrTwo}.\mathrm{Plus}$, the $\bases$ of
$\mathrm{Result}$ include the $\mathrm{Plus}$ of both
$\mathrm{One}$ and $\mathrm{Two}$.
Each $\mathrm{Plus}$ resolves its $\mathrm{addend}$ via
$\this$, reaching $\mathrm{ThreeOrFour}$, which is itself a
trie union of $\mathrm{Three}$ and $\mathrm{Four}$:
the $\supers$ of the $\mathrm{addend}$ path contain both values.
Since $\mathrm{Successor}.\mathrm{Plus}$ recursively delegates to
$\mathrm{predecessor}.\mathrm{Plus}$ with an incremented addend,
and $\mathrm{predecessor}$ also resolves through the trie union,
the recursion branches at every trie node.
The result $\mathrm{sum}$ therefore collects all paths reachable
through any combination of the two operands, which is the
Cartesian product.

The same $\mathrm{Plus}$ and $\mathrm{Equal}$ files that operate on
single values thus produce, without modification, the relational
semantics of logic programming:
$\overrides$ and $\supers$ collect all inheritance paths, and
every operation distributes over them.
\fi

\section{\bilingual{Discussion}{讨论}}

\subsection{\bilingual{Emergent Phenomena}{涌现现象}}
\label{sec:emergent-phenomena}

\ifInheritanceChinese
在 MIXINv2 中开发真实程序的过程中，我们发现了三种现象，每一种都是对某种已有设计模式的非预期用法，超出了该模式的原始能力。

\begin{description}
  \item[关系语义]
    第~\ref{sec:case-study}~节案例研究中的笛卡尔积之所以出现，是因为字典树是一种集合结构，因此对字典树的深度合并即为集合并，而对合并后字典树的操作则遍历其操作数的笛卡尔积。
    这并不局限于 $\mathrm{Nat}$：任何 Datalog 元组都可编码为链表，而深度合并在这些元组上计算集合操作，我们认为这正是 Datalog 语义。
    附录~\ref{app:datalog-encoding}~将 Datalog 编码进继承演算，仅用第~\ref{sec:syntax}~节的三个基本构造即可表达变量连接、常量匹配、多体规则和递归规则。
    这种逻辑编程语义是对编码案例研究算术的工厂模式的非预期用法中涌现出来的。

  \item[对不可扩展性的免疫]\label{item:immunity}
    在第~\ref{sec:case-study}~节的案例研究中，$\mathrm{Nat}$ 编码被证明是无类型表达问题~\cite{wadler1998-expression-problem}的一个实例：%
    \footnote{
      表达问题有两个要求：沿操作轴和数据构造子轴均可扩展，以及静态类型安全。无类型演算无法满足类型安全要求，因此我们只考虑双轴可扩展性要求。
    }
    一个新操作和一个新数据类型各自作为独立的 mixin 添加，沿操作轴和情形轴均进行了扩展，且无需修改任何已有 mixin。
    在其他场景中，访问者模式允许添加新操作但不允许添加新情形；而这里两者都允许，因为第~\ref{sec:syntax}~节的每个构造都可通过深度合并（公式~\ref{eq:overrides}）进行扩展，所以不可扩展的那一半根本无法写出；第~\ref{sec:asymmetry}~节（定理~\ref{thm:cartesian-nonexpressibility}）则形式化了使这一结论无条件成立的不对称性。
    在其他地方恢复第二个扩展维度需要专门的机制~\cite{liang1995-monad-transformers, lammel2003-scrap-your-boilerplate, loh2006-open-data-types,
      swierstra2008-data-types-a-la-carte, carette2009-finally-tagless, oliveira2012-object-algebras,
      kiselyov2013-extensible-effects, wu2014-effect-handlers-in-scope, kiselyov2015-freer-monads,
    wang2016-expression-problem-trivially, leijen2017-row-typed-algebraic-effects, poulsen2023-hefty-algebras}。
    第~\ref{sec:semantic-variants}~节分析了为何深度合并是唯一能实现这一点的组合机制。
    这第二个可扩展性维度是对访问者模式的非预期用法中涌现出来的。

\end{description}

\label{sec:practical-function-color}%
更多现象依赖于 FFI，而本文中的示例排除了 FFI。其中一种是函数颜色盲~\cite{nystrom2015-function-color}：同一个 MIXINv2 文件可以在同步和异步运行时下运行，这是通过继承进行依赖注入的非预期用法。补充材料\anon[]{~\cite{mixin2025}}给出了相关示例。
\else
Developing real programs in MIXINv2 surfaced three phenomena, each an
off-label use of an adopted design pattern beyond its original abilities.

\begin{description}
  \item[Relational semantics]
    The Cartesian product in the case study of
    Section~\ref{sec:case-study} arose because a trie is a set
    structure, so deep merge over tries is set union and an operation
    over the merged trie ranges over the product of its operands.
    This is not specific to $\mathrm{Nat}$: any Datalog tuple encodes
    as a linked list, and deep merge computes set operations over such
    tuples, which we believe is Datalog semantics.
    Appendix~\ref{app:datalog-encoding} encodes Datalog into
    inheritance-calculus, expressing variable joins,
    constant matches, multi-body rules, and recursive rules with
    the three primitives of Section~\ref{sec:syntax} alone.
    This logic-programming semantics emerged from an off-label use of
    the factory pattern that encodes the case study's arithmetic.

  \item[Immunity to nonextensibility]\label{item:immunity}
    In the case study of Section~\ref{sec:case-study}, the
    $\mathrm{Nat}$ encoding turned out to be an instance of the
    untyped Expression
    Problem~\cite{wadler1998-expression-problem}:%
    \footnote{
      The Expression Problem has two requirements: extensibility along
      both the operation and the data-constructor axis, and static
      type safety. An untyped calculus cannot meet the type-safety
      requirement, so we consider only the two-axis extensibility
      requirement.
    }
    a new operation and a new data type were each added as a separate
    mixin, extending along both the operation axis and the case axis,
    with no existing mixin modified.
    Elsewhere the visitor pattern admits new operations but not new
    cases; here it admits both, because every construct of
    Section~\ref{sec:syntax} is extensible via deep merge
    (equation~\ref{eq:overrides}), so the nonextensible half cannot be
    written, and Section~\ref{sec:asymmetry}
    (Theorem~\ref{thm:cartesian-nonexpressibility}) formalizes the
    asymmetry that makes this unconditional.
    Recovering that second dimension elsewhere takes dedicated
    mechanisms~\cite{liang1995-monad-transformers, lammel2003-scrap-your-boilerplate, loh2006-open-data-types,
      swierstra2008-data-types-a-la-carte, carette2009-finally-tagless, oliveira2012-object-algebras,
      kiselyov2013-extensible-effects, wu2014-effect-handlers-in-scope, kiselyov2015-freer-monads,
    wang2016-expression-problem-trivially, leijen2017-row-typed-algebraic-effects, poulsen2023-hefty-algebras}.
    Section~\ref{sec:semantic-variants} analyzes why deep merge
    is the only composition mechanism that achieves this.
    This second extensibility dimension emerged from an off-label use
    of the visitor pattern.

\end{description}

\label{sec:practical-function-color}%
Further phenomena rely on the FFI, which the examples in this paper
exclude. One is function color
blindness~\cite{nystrom2015-function-color}: the same MIXINv2 file runs
under both synchronous and asynchronous runtimes, an off-label use of
dependency injection via inheritance. The
supplementary material\anon[]{~\cite{mixin2025}} gives the example.
\fi

\subsection{\bilingual{Semantic Alternatives}{语义替代方案}}
\label{sec:semantic-variants}

\ifInheritanceChinese
上述\hyperref[item:immunity]{对不可扩展性的免疫}条目确立了继承演算对不可扩展性的免疫性。我们将此作为设计目标，并通过调查我们已知的所有候选方案，询问是否有替代设计能共享这一性质。

\paragraph{\bilingual{Degrees of freedom}{自由度}}
以下分析假设 AST 是一棵\emph{命名树}：从父节点到子节点的每条边都携带一个标签，因此每个节点都由从根节点出发的路径唯一标识。对于配置语言，这自然成立：JSON、YAML 或 Nix 属性集中的每个键都是一个标签，而嵌套产生路径。%
\footnote{
  Nix 是一门具有一等函数的函数式语言；仅凭其属性集并不构成命名树 DSL。然而，NixOS 模块系统~\cite{nixosmodules}将用户界面限制为通过递归合并组合的嵌套属性集，而正是这种 DSL 层面的结构自然地构成了命名树。
}
对于主流编程语言，AST 并不直接是命名树，因为匿名表达式位置和隐式作用域缺乏标签。然而，经过 ANF 变换~\cite{flanagan1993-essence-compiling-continuations}后，每个中间结果都绑定到一个名字，而剩余的匿名位置%
\footnote{
  例如，$\mathbf{let}$ 的主体。
}
可以被分配合成的新鲜标签，例如第~\ref{sec:forward-translation}~节中的 $\mathrm{result}$ 和 $\mathrm{tailCall}$，从而得到一棵命名树。

在此前提下，每种程序是带引用的命名树的语言，都可以由两个有限集来刻画：一个由\emph{节点类型}构成的集合 $\mathcal{N}$，每种节点类型决定子树如何组织其子节点；以及一个由\emph{引用类型}构成的集合 $\mathcal{R}$，每种引用类型决定树中一个位置如何引用另一个位置。语法由 $\mathcal{N}$ 和 $\mathcal{R}$ 的选择固定；语义则由每种节点类型如何构造组合以及每种引用类型如何解析其目标来决定。

{\small
  \begin{center}
    \begin{tabular}{lp{0.38\columnwidth}p{0.33\columnwidth}}
      \bilingual{\textbf{Language}}{\textbf{语言}}
      & \bilingual{\textbf{Node types $\mathcal{N}$}}{\textbf{节点类型 $\mathcal{N}$}}
      & \bilingual{\textbf{Reference types $\mathcal{R}$}}{\textbf{引用类型 $\mathcal{R}$}} \\
      \hline
      Java (OOP)
      & \bilingual{class, interface, method,
      block, conditional, loop,
      \texttt{throw}, lambda, \ldots}{类、接口、方法、块、条件、循环、\texttt{throw}、lambda、\ldots}
      & \bilingual{variable, field access, method call,
      \texttt{extends}/\texttt{implements},
      \texttt{import}, \ldots}{变量、字段访问、方法调用、\texttt{extends}/\texttt{implements}、\texttt{import}、\ldots} \\[3pt]
      Haskell (FP)
      & \bilingual{$\lambda$, application, \texttt{let},
      \texttt{case}, \texttt{data},
      type class, instance, \ldots}{$\lambda$、应用、\texttt{let}、\texttt{case}、\texttt{data}、类型类、实例、\ldots}
      & \bilingual{variable, constructor,
      type class dispatch,
      \texttt{import}, \ldots}{变量、构造子、类型类分派、\texttt{import}、\ldots} \\[3pt]
      Scala\footnote{
        \bilingual{Scala is a multi-paradigm language.}{Scala 是一门多范式语言。}
      }
      & \bilingual{class, trait, object, method,
      block, conditional, \texttt{match},
      \texttt{while}, \texttt{throw}, \ldots}{类、trait、对象、方法、块、条件、\texttt{match}、\texttt{while}、\texttt{throw}、\ldots}
      & \bilingual{variable, member access,
      \texttt{extends}, \texttt{with},
      implicits, \texttt{import}, \ldots}{变量、成员访问、\texttt{extends}、\texttt{with}、隐式参数、\texttt{import}、\ldots} \\
      \hline
      $\lambda$-\bilingual{calculus}{演算}
      & \bilingual{abstraction, application}{抽象、应用}
      & \bilingual{variable}{变量} \\
      \bilingual{Inheritance-calculus}{继承演算}
      & \bilingual{record}{记录}
      & \bilingual{inheritance}{继承} \\
    \end{tabular}
  \end{center}
}

\paragraph{\bilingual{Immunity to nonextensibility}{对不可扩展性的免疫}}
程序 $P$ 中的路径 $p$ 是\emph{可扩展的}，若存在程序 $Q$，使得通过引用将 $P$ 与 $Q$ 组合后，在不修改 $P$ 的前提下，$p$ 处能够观察到新的属性。一种语言\emph{对不可扩展性免疫}，若其每个良构程序的每条路径都是可扩展的。

我们现在对设计空间进行分类，以解释为何我们选择第~\ref{sec:mixin-trees}~节的语义。两个维度保持自由：\emph{组合机制}，决定单一引用类型如何合并定义；以及\emph{引用解析}，决定单一引用类型如何找到其目标。

\subsubsection*{\bilingual{The composition mechanism}{组合机制}}

每种组合机制决定了一个抽象的定义如何被合并进另一个抽象。下表按各候选方案是否满足可交换性~(C)、幂等性~(I) 和可结合性~(A)，以及其可扩展性模型对已知候选方案进行分类。

\begin{center}
  \begin{tabular}{lcccll}
    \bilingual{\textbf{Candidate}}{\textbf{候选方案}} & \textbf{C} & \textbf{I} & \textbf{A}
    & \bilingual{\textbf{Extensibility}}{\textbf{可扩展性}} & \bilingual{\textbf{Representative}}{\textbf{代表}} \\
    \hline
    \bilingual{Deep merge}{深度合并}
    & $\checkmark$ & $\checkmark$ & $\checkmark$
    & \bilingual{universal}{通用} & \bilingual{this paper}{本文} \\
    \bilingual{Shallow merge}{浅层合并}
    & $\times$ & $\checkmark$ & $\checkmark$
    & \bilingual{top-level only}{仅顶层} & JS spread \\
    \bilingual{Conflict rejection}{冲突拒绝}
    & $\checkmark$ & $\checkmark$ & $\checkmark$
    & \bilingual{none on overlap}{重叠时无扩展}
    & Harper--Pierce~\cite{harper1991-symmetric-record-concatenation} \\
    \bilingual{Priority merge}{优先级合并}
    & $\times$ & $\times$ & $\checkmark$
    & \bilingual{universal}{通用} & Jsonnet~\texttt{+}, Nix~\texttt{//} \\
    $\beta$-\bilingual{reduction}{归约}
    & $\times$ & $\times$ & $\times$
    & \bilingual{opt-in}{按需选择} & Java, Haskell, Scala, \ldots \\
    \bilingual{Conditional}{条件并入}
    & \multicolumn{3}{c}{\bilingual{causes divergence}{导致发散}}
    & --- & --- \\
  \end{tabular}
\end{center}

\begin{description}
  \item[\bilingual{Deep merge}{深度合并}]
    对同名定义的合并递归进入其子树（$\overrides$，公式~\ref{eq:overrides}）。每个标签都是一个扩展点，无论原作者是否预期了扩展。

  \item[\bilingual{Shallow merge}{浅层合并}]
    嵌套标签无法被独立扩展。

  \item[\bilingual{Conflict rejection}{冲突拒绝}]
    同名定义被作为错误拒绝。三个等式性质均成立，但为空洞成立：这些等式得到满足，是因为用于验证它们的组合被拒绝了。冲突拒绝阻止了对已有标签的扩展~\cite{harper1991-symmetric-record-concatenation}，而这恰恰是表达问题~\cite{wadler1998-expression-problem}所要求的可组合性。

  \item[\bilingual{Priority merge}{优先级合并}]
    一个来源具有优先权~\cite{jsonnet,dolstra2010-nixos}，因此 $A + B \neq B + A$，引入了顺序依赖。

  \item[$\beta$-\bilingual{reduction}{归约}]
    只有作者放置的参数才是扩展点；函数体是封闭的。要使一个新的方面可扩展，原函数必须被重写以接受一个额外参数。任何基于函数应用的组合机制都继承了这一限制，因为 $\beta$-归约使函数体封闭。表达问题~\cite{wadler1998-expression-problem}在此类语言中之所以困难，原因相同：可扩展性是按需选择的。模拟对不可扩展性免疫的框架（访问者模式、对象代数~\cite{oliveira2012-object-algebras}、finally tagless 解释器~\cite{carette2009-finally-tagless}）正是为绕过这一限制而生的机制。

  \item[\bilingual{Conditional incorporation}{条件并入}]
    根据其他定义的存在与否来移除或有条件地包含定义，会产生依赖循环：某定义是否存在，取决于另一定义是否存在，而后者又取决于前者。当添加和移除交替出现时，递归求值（附录~\ref{app:well-definedness}）无法到达不动点，从而导致发散。这与图灵完备性固有的发散不同：图灵完备程序发散是因为计算无界，而条件并入发散是因为直接后果算子 $T_P$ 振荡而非单调递增。
\end{description}

\noindent
在上述候选方案中，深度合并是我们找到的唯一一个可交换、幂等、可结合且对不可扩展性免疫的方案。其他候选方案各自引入了对不可扩展性的脆弱性，或导致发散。

\subsubsection*{\bilingual{The reference resolution}{引用解析}}

一旦将记录上的深度合并确定为组合机制，六个语义方程中的四个便由语法决定：$\properties$~(\ref{eq:properties}) 读取记录的标签，$\bases$~(\ref{eq:bases}) 读取其继承来源，$\supers$~(\ref{eq:supers}) 取传递闭包，$\overrides$~(\ref{eq:overrides}) 实现合并。唯一剩余的自由度是引用的解析方式：$\resolve$~(\ref{eq:resolve}) 和 $\this$~(\ref{eq:this}) 方程。我们考察三个已知候选方案。

\begin{description}
  \item[\bilingual{Early binding}{早期绑定}]
    使用早期绑定时，引用在定义时即解析到一个固定路径，与记录后来如何被继承无关。CUE~\cite{cue2019}是这方面的典型：字段引用被静态解析，嵌套字段无法引用其封闭上下文中由后续组合贡献的属性。这使得 mixin 无法充当 $\lambda$-演算的调用栈：$\beta$-归约要求被调用方能观察到调用点提供的参数，而那是一次后续组合。使用早期绑定时，引用在此组合发生前就已被冻结。在方程中，$\resolve$~(\ref{eq:resolve}) 调用 $\this$~(\ref{eq:this})，后者遍历继承的树结构；迟绑定正是使 $\this$ 能够观察到后续组合的 mixin 所贡献属性的原因。若没有图灵完备的 mixin，函数必须被重新引入为组合机制，从而再次暴露上述按需选择的可扩展性问题。

  \item[\bilingual{Dynamic scope}{动态作用域}]
    使用动态作用域时，引用在使用点而非定义点处解析。Jsonnet~\cite{jsonnet}采用此方式：\texttt{self} 始终指向最终合并后的对象，而非定义点处的对象。联合文件系统在文件系统层面表现出相同的性质：符号链接的相对路径在解引用时相对于挂载上下文解析，而非相对于定义该链接的层。动态作用域破坏了第~\ref{sec:forward-translation}~节的 $\lambda$-演算嵌入。替换引理（引理~\ref{lem:substitution}）依赖于基本情况 $M = y$（$y \neq x$）：$y$ 的 de~Bruijn 指标 $m$ 从定义点路径 $\init(p_{\mathrm{def}})$ 向上走到一个与 $x$ 的绑定器不同的作用域层级，因此 $x$ 作用域处的替换 $\mathrm{argument} \mapsto \mathcal{T}(V)$ 不产生任何效果。动态作用域下，这种分离不再成立。%
    \footnote{
      反例：$(\lambda x.\, \lambda y.\, x)\; V_1\; V_2$。$\lambda y.\, x$ 内部对 $x$ 的引用翻译为 $\dbi{1}.\mathrm{argument}$（de~Bruijn 指标 $1$）。使用迟绑定时，$\this$ 从定义点作用域（内层 $\lambda$）出发，向上走一层到外层 $\lambda$，到达槽 $\mathrm{argument} \mapsto \mathcal{T}(V_1)$。使用动态作用域时，引用从使用点上下文解析，其中最内层封闭作用域是应用 $\{C_1.\mathrm{result},\; \mathrm{argument} \mapsto \mathcal{T}(V_2)\}$；向上走一层到达该作用域的 $\mathrm{argument}$ 槽，返回 $\mathcal{T}(V_2)$ 而非 $\mathcal{T}(V_1)$。这是动态作用域的经典变量捕获失败。
    }
    与早期绑定一样，mixin 无法完全替代函数。

  \item[\bilingual{Single-target $\this$}{单目标 $\this$}]
    单目标 $\this$ 是一个有前途的候选方案：单路径引理~(\ref{lem:single-path})表明，对于翻译后的 $\lambda$-项，$\this$~(\ref{eq:this}) 中的前沿集 $S$ 始终恰好包含一条路径，因此 $\lambda$-演算嵌入不需要多目标解析。然而，单目标 $\this$ 因另一原因而对不可扩展性脆弱。当多条继承路径到达同一作用域时，单目标解析必须拒绝该情况或选择其中一条路径。NixOS 模块系统~\cite{nixosmodules}会引发重复声明错误（附录~\ref{app:scala-multi-target}）。这拒绝了两个独立编写的模块定义同一选项的嵌套扩展，而这种组合在求解表达问题~\cite{wadler1998-expression-problem}时会出现，如第~\ref{sec:expression-problem}~节所示。

    我们对 MIXINv2 的前三个实现使用了 $\this$ 解析的标准技术：遵循 Cook~\cite{cook1989-denotational-inheritance}的基于闭包的不动点，以及使用 de~Bruijn 指标~\cite{debruijn1972-nameless-dummies}的基于栈的环境查找。当继承引入到同一封闭作用域的多条路径时，三种实现都产生了隐晦的错误。困难在于结构上。闭包捕获单一环境，因此不动点组合子只能为单一 $\this$ 目标求解。基于栈的环境是一条线性链，每个作用域层级只有一个父节点。这两种数据结构都内嵌了一个与多重继承自然产生的多目标情况不相容的单值假设（附录~\ref{app:scala-multi-target}）。公式~(\ref{eq:this})是在放弃这一假设后涌现出来的：它跟踪一个\emph{前沿集} $S$，而非单一当前作用域。
\end{description}

\noindent
下表总结了两个层面上所有替代方案的脆弱性。

\begin{center}
  \begin{tabular}{llll}
    & \bilingual{\textbf{Candidate}}{\textbf{候选方案}} & \bilingual{\textbf{Vulnerability}}{\textbf{脆弱性}}
    & \bilingual{\textbf{Representative}}{\textbf{代表}} \\
    \hline
    \multirow{5}{*}{\rotatebox[origin=c]{90}{\scriptsize \bilingual{composition}{组合}}}
    & \bilingual{Shallow merge}{浅层合并}
    & \bilingual{nested extensions not composable}{嵌套扩展不可组合}
    & JS spread \\
    & \bilingual{Conflict rejection}{冲突拒绝}
    & \bilingual{same-label extensions rejected}{同名扩展被拒绝}
    & Harper--Pierce~\cite{harper1991-symmetric-record-concatenation} \\
    & \bilingual{Priority merge}{优先级合并}
    & \bilingual{ordering dependency}{顺序依赖}
    & Jsonnet~\texttt{+}~\cite{jsonnet},
    Nix~\texttt{//}~\cite{dolstra2010-nixos} \\
    & $\beta$-\bilingual{reduction}{归约}
    & \bilingual{opt-in extensibility}{按需选择的可扩展性}
    & Java, Haskell, Scala, \ldots \\
    & \bilingual{Conditional}{条件并入}
    & \bilingual{causes divergence}{导致发散}
    & --- \\
    \hline
    \multirow{3}{*}{\rotatebox[origin=c]{90}{\scriptsize \bilingual{resolution}{解析}}}
    & \bilingual{Early binding}{早期绑定}
    & \bilingual{mixin cannot replace function}{mixin 无法替代函数}
    & CUE~\cite{cue2019} \\
    & \bilingual{Dynamic scope}{动态作用域}
    & \bilingual{mixin cannot replace function}{mixin 无法替代函数}
    & Jsonnet~\cite{jsonnet}, union FS \\
    & \bilingual{Single-target $\this$}{单目标 $\this$}
    & \bilingual{nested extensions rejected}{嵌套扩展被拒绝}
    & NixOS modules~\cite{nixosmodules} \\
  \end{tabular}
\end{center}

\noindent
我们在两个层面考察的所有替代方案，均对不可扩展性不免疫。

我们通过尝试本节中的每个替代方案并发现其失败，从而得到了这一语义：浅层合并无法组合嵌套扩展；优先级合并引入了顺序依赖；早期绑定在组合能够提供引用目标之前就将引用冻结；动态作用域破坏了 $\lambda$-演算嵌入所依赖的替换；单目标 $\this$ 拒绝了表达问题所要求的多目标组合。深度合并加上多目标迟绑定 $\this$，正是最终剩下的选择。
\else
The \hyperref[item:immunity]{immunity-to-nonextensibility item} above established that inheritance-calculus is immune
to nonextensibility.
We adopt this as a design goal and ask whether an alternative
design could share this property
by surveying every candidate known to us.

\paragraph{\bilingual{Degrees of freedom}{自由度}}
The analysis below assumes that the AST is a \emph{named tree}:%
every edge from parent to child carries a label, so that every
node is uniquely identified by the path from the root.
For configuration languages this is naturally the case: every
key in a JSON, YAML, or Nix attribute set is a label, and nesting
produces paths.%
\footnote{
  Nix is a functional language with first-class functions;
  its attribute sets alone do not form a named-tree DSL.
  The NixOS module system~\cite{nixosmodules}, however,
  restricts the user-facing interface to nested attribute sets
  composed by recursive merging, and it is this DSL-level
  structure that naturally forms a named tree.
}
For mainstream programming languages the AST is not directly a
named tree, since anonymous expression positions and implicit
scopes lack labels.
After ANF
transformation~\cite{flanagan1993-essence-compiling-continuations}, however, every intermediate
result is bound to a name, and the remaining anonymous positions%
\footnote{
  For example, the body of a $\mathbf{let}$.
}
can be assigned synthetic
fresh labels such as $\mathrm{result}$ and $\mathrm{tailCall}$
in Section~\ref{sec:forward-translation}, yielding a named
tree.

Under this premise, every language whose programs are such
named trees with references can be characterized by two finite sets:
a set~$\mathcal{N}$ of \emph{node types}, each determining how a
subtree organizes its children, and a set~$\mathcal{R}$ of
\emph{reference types}, each determining how one position in the
tree refers to another.
The syntax is fixed by the choice of $\mathcal{N}$ and
$\mathcal{R}$; the semantics is determined by how each node type
structures composition and how each reference type resolves its
target.

{\small
  \begin{center}
    \begin{tabular}{lp{0.38\columnwidth}p{0.33\columnwidth}}
      \textbf{Language}
      & \textbf{Node types $\mathcal{N}$}
      & \textbf{Reference types $\mathcal{R}$} \\
      \hline
      Java (OOP)
      & class, interface, method,
      block, conditional, loop,
      \texttt{throw}, lambda, \ldots
      & variable, field access, method call,
      \texttt{extends}/\texttt{implements},
      \texttt{import}, \ldots \\[3pt]
      Haskell (FP)
      & $\lambda$, application, \texttt{let},
      \texttt{case}, \texttt{data},
      type class, instance, \ldots
      & variable, constructor,
      type class dispatch,
      \texttt{import}, \ldots \\[3pt]
      Scala\footnote{
        Scala is a multi-paradigm language.
      }
      & class, trait, object, method,
      block, conditional, \texttt{match},
      \texttt{while}, \texttt{throw}, \ldots
      & variable, member access,
      \texttt{extends}, \texttt{with},
      implicits, \texttt{import}, \ldots \\
      \hline
      $\lambda$-calculus
      & abstraction, application
      & variable \\
      Inheritance-calculus
      & record
      & inheritance \\
    \end{tabular}
  \end{center}
}

\paragraph{\bilingual{Immunity to nonextensibility}{对不可扩展性的免疫}}
A path~$p$ in a program~$P$ is \emph{extensible} if there exists
a program~$Q$ such that composing $P$ with~$Q$ through a reference
makes new properties observable at~$p$ without modifying~$P$.
A language is \emph{immune to nonextensibility} if every path in
every well-formed program is extensible.

We now classify the design space to explain why we chose the
semantics of Section~\ref{sec:mixin-trees}.
Two axes remain free:
the \emph{composition mechanism}, which determines how the single reference
type incorporates definitions;
and the \emph{reference resolution}, which determines how the single reference
type finds its target.

\subsubsection*{\bilingual{The composition mechanism}{组合机制}}

Every composition mechanism determines how the definitions of one
abstraction are incorporated into another.
The following table classifies the known candidates by whether
they satisfy commutativity~(C), idempotence~(I), and
associativity~(A), and by their extensibility model.

\begin{center}
  \begin{tabular}{lcccll}
    \bilingual{\textbf{Candidate}}{\textbf{候选方案}} & \textbf{C} & \textbf{I} & \textbf{A}
    & \bilingual{\textbf{Extensibility}}{\textbf{可扩展性}} & \bilingual{\textbf{Representative}}{\textbf{代表}} \\
    \hline
    \bilingual{Deep merge}{深度合并}
    & $\checkmark$ & $\checkmark$ & $\checkmark$
    & \bilingual{universal}{通用} & \bilingual{this paper}{本文} \\
    \bilingual{Shallow merge}{浅层合并}
    & $\times$ & $\checkmark$ & $\checkmark$
    & \bilingual{top-level only}{仅顶层} & JS spread \\
    \bilingual{Conflict rejection}{冲突拒绝}
    & $\checkmark$ & $\checkmark$ & $\checkmark$
    & \bilingual{none on overlap}{重叠时无扩展}
    & Harper--Pierce~\cite{harper1991-symmetric-record-concatenation} \\
    \bilingual{Priority merge}{优先级合并}
    & $\times$ & $\times$ & $\checkmark$
    & \bilingual{universal}{通用} & Jsonnet~\texttt{+}, Nix~\texttt{//} \\
    $\beta$-\bilingual{reduction}{归约}
    & $\times$ & $\times$ & $\times$
    & \bilingual{opt-in}{按需选择} & Java, Haskell, Scala, \ldots \\
    \bilingual{Conditional}{条件并入}
    & \multicolumn{3}{c}{\bilingual{causes divergence}{导致发散}}
    & --- & --- \\
  \end{tabular}
\end{center}

\begin{description}
  \item[Deep merge]
    Merging same-label definitions recurses into their subtrees
    ($\overrides$, equation~\ref{eq:overrides}).
    Every label is an extension point, whether or not the
    original author anticipated extension.

  \item[Shallow merge]
    Nested labels cannot be independently extended.

  \item[Conflict rejection]
    Same-label definitions are rejected as errors.
    All three identities hold, but vacuously:
    the identities are satisfied because the compositions that
    would test them are refused.
    Conflict rejection prevents extending existing
    labels~\cite{harper1991-symmetric-record-concatenation}, precisely the composability that the
    Expression Problem~\cite{wadler1998-expression-problem} requires.

  \item[Priority merge]
    One source takes
    precedence~\cite{jsonnet,dolstra2010-nixos},
    so $A + B \neq B + A$, introducing
    an ordering dependency.

  \item[$\beta$-reduction]
    Only the parameters placed by the author are extension
    points; the function body is closed.
    To make a new aspect extensible, the original function must
    be rewritten to accept an additional parameter.
    Any composition mechanism based on function application
    inherits this limitation, because $\beta$-reduction
    closes the function body.
    The Expression Problem~\cite{wadler1998-expression-problem} is hard
    in such languages for the same reason:
    extensibility is opt-in.
    The frameworks that simulate immunity to nonextensibility
    (visitor pattern,
      object algebras~\cite{oliveira2012-object-algebras},
    finally tagless interpreters~\cite{carette2009-finally-tagless})
    are the machinery needed to work around this limitation.

  \item[Conditional incorporation]
    Removing or conditionally including definitions
    based on the presence of other definitions
    creates a dependency cycle:
    whether a definition exists depends on
    whether another definition exists, which in turn
    depends on the first.
    The recursive evaluation
    (Appendix~\ref{app:well-definedness}) cannot
    reach a fixed point when addition and removal
    alternate, causing divergence.
    This differs from the divergence inherent to
    Turing completeness: Turing-complete programs
    diverge because computation is unbounded, whereas
    conditional incorporation diverges because the
    immediate consequence operator $T_P$ oscillates
    rather than ascending monotonically.
\end{description}

\noindent
Among the candidates above, deep merge is the only one we found
that is commutative, idempotent, associative, and
immune to nonextensibility.
The other candidates each introduce a vulnerability to
nonextensibility or cause divergence.

\subsubsection*{\bilingual{The reference resolution}{引用解析}}

Once deep merge over records is fixed as the composition
mechanism, four of the six semantic equations are determined
by the syntax:
$\properties$~(\ref{eq:properties}) reads the labels of a
record, $\bases$~(\ref{eq:bases}) reads its inheritance
sources, $\supers$~(\ref{eq:supers}) takes the transitive
closure, and $\overrides$~(\ref{eq:overrides}) implements
the merge.
The only remaining degree of freedom is how references
resolve: the $\resolve$~(\ref{eq:resolve}) and
$\this$~(\ref{eq:this}) equations.
We examine the three known candidates.

\begin{description}
  \item[Early binding]
    With early binding, references resolve at definition time to a
    fixed path, independently of how the record is later inherited.
    CUE~\cite{cue2019} exemplifies this: field references are
    statically resolved, and a nested field cannot refer to
    properties of its enclosing context contributed by later
    composition.
    This prevents the mixin from serving as a call stack for the
    $\lambda$-calculus: $\beta$-reduction requires the callee to
    observe the argument supplied at the call site, which is a
    later composition.
    With early binding, references are frozen before this
    composition occurs.
    In the equations,
    $\resolve$~(\ref{eq:resolve}) calls
    $\this$~(\ref{eq:this}), which walks the inherited
    tree structure; late binding is what allows
    $\this$ to see properties contributed by later-composed
    mixins.
    Without a Turing-complete mixin, functions must be
    reintroduced as the composition mechanism, re-exposing
    the opt-in extensibility problem above.

  \item[Dynamic scope]
    With dynamic scope, references resolve at the point of use
    rather than the definition site.
    Jsonnet~\cite{jsonnet} follows this approach: \texttt{self}
    always refers to the final merged object, not the object at
    the definition site.
    Union file systems exhibit the same property at the
    file-system level: a symbolic link's relative path resolves
    against the mount context at dereference time, not against the
    layer that defined the link.
    Dynamic scope breaks the $\lambda$-calculus embedding
    of Section~\ref{sec:forward-translation}.
    The Substitution Lemma (Lemma~\ref{lem:substitution})
    relies on the base case $M = y$ ($y \neq x$):
    the de~Bruijn index $m$ of~$y$ walks up from the
    definition-site path $\init(p_{\mathrm{def}})$ to a scope
    level different from~$x$'s binder, so the substitution
    $\mathrm{argument} \mapsto \mathcal{T}(V)$ at~$x$'s scope
    has no effect.
    With dynamic scope this separation fails.%
    \footnote{
      Counterexample:
      $(\lambda x.\, \lambda y.\, x)\; V_1\; V_2$.
      The reference to~$x$ inside $\lambda y.\, x$ translates to
      $\dbi{1}.\mathrm{argument}$ (de~Bruijn index~$1$).
      With late binding, $\this$ starts from the definition-site
      scope (the inner $\lambda$) and walks up one level to the
      outer $\lambda$, reaching the slot
      $\mathrm{argument} \mapsto \mathcal{T}(V_1)$.
      With dynamic scope, the reference resolves from the
      use-site context, where the innermost enclosing scope is
      the application $\{C_1.\mathrm{result},\;
      \mathrm{argument} \mapsto \mathcal{T}(V_2)\}$;
      walking up one level reaches this scope's
      $\mathrm{argument}$ slot, returning $\mathcal{T}(V_2)$
      instead of $\mathcal{T}(V_1)$.
      This is the classical variable-capture failure of dynamic
      scoping.
    }
    As with early binding, the mixin cannot serve as a complete
    replacement for functions.

  \item[Single-target $\this$]
    Single-target $\this$ is a promising candidate:
    the Single-Path Lemma~(\ref{lem:single-path}) shows that
    for translated $\lambda$-terms, the frontier set~$S$ in
    $\this$~(\ref{eq:this}) always contains exactly one path,
    so the $\lambda$-calculus embedding does not require
    multi-target resolution.
    However, single-target $\this$ is vulnerable to
    nonextensibility for a different reason.
    When multiple inheritance routes reach the same scope,
    single-target resolution must reject the situation or select
    one route.
    The NixOS module system~\cite{nixosmodules} raises a
    duplicate-declaration error
    (Appendix~\ref{app:scala-multi-target}).
    This rejects nested extensions where two independently authored
    modules define the same option, a composition that arises
    when solving the Expression
    Problem~\cite{wadler1998-expression-problem},
    as Section~\ref{sec:expression-problem} demonstrates.

    Our first three implementations of MIXINv2 used
    standard techniques for $\this$ resolution: closure-based
    fixed points following Cook~\cite{cook1989-denotational-inheritance}, and
    stack-based environment lookup using de~Bruijn
    indices~\cite{debruijn1972-nameless-dummies}.
    All three produced subtle bugs when inheritance introduced
    multiple routes to the same enclosing scope.
    The difficulty was structural.
    A closure captures a single environment, so a fixed-point
    combinator solves for a single $\this$ target.
    A stack-based environment is a linear chain where each scope
    level has one parent.
    Both data structures embed a single-valued assumption
    incompatible with the multi-target situation that multiple
    inheritance naturally produces
    (Appendix~\ref{app:scala-multi-target}).
    Equation~(\ref{eq:this}) emerged from abandoning this
    assumption: it tracks a \emph{frontier set}~$S$ rather than a
    single current scope.
\end{description}

\noindent
The following table summarizes the vulnerabilities of all
alternatives at both levels.

\begin{center}
  \begin{tabular}{llll}
    & \bilingual{\textbf{Candidate}}{\textbf{候选方案}} & \bilingual{\textbf{Vulnerability}}{\textbf{脆弱性}}
    & \bilingual{\textbf{Representative}}{\textbf{代表}} \\
    \hline
    \multirow{5}{*}{\rotatebox[origin=c]{90}{\scriptsize \bilingual{composition}{组合}}}
    & \bilingual{Shallow merge}{浅层合并}
    & \bilingual{nested extensions not composable}{嵌套扩展不可组合}
    & JS spread \\
    & \bilingual{Conflict rejection}{冲突拒绝}
    & \bilingual{same-label extensions rejected}{同名扩展被拒绝}
    & Harper--Pierce~\cite{harper1991-symmetric-record-concatenation} \\
    & \bilingual{Priority merge}{优先级合并}
    & \bilingual{ordering dependency}{顺序依赖}
    & Jsonnet~\texttt{+}~\cite{jsonnet},
    Nix~\texttt{//}~\cite{dolstra2010-nixos} \\
    & $\beta$-\bilingual{reduction}{归约}
    & \bilingual{opt-in extensibility}{按需选择的可扩展性}
    & Java, Haskell, Scala, \ldots \\
    & \bilingual{Conditional}{条件并入}
    & \bilingual{causes divergence}{导致发散}
    & --- \\
    \hline
    \multirow{3}{*}{\rotatebox[origin=c]{90}{\scriptsize \bilingual{resolution}{解析}}}
    & \bilingual{Early binding}{早期绑定}
    & \bilingual{mixin cannot replace function}{mixin 无法替代函数}
    & CUE~\cite{cue2019} \\
    & \bilingual{Dynamic scope}{动态作用域}
    & \bilingual{mixin cannot replace function}{mixin 无法替代函数}
    & Jsonnet~\cite{jsonnet}, union FS \\
    & \bilingual{Single-target $\this$}{单目标 $\this$}
    & \bilingual{nested extensions rejected}{嵌套扩展被拒绝}
    & NixOS modules~\cite{nixosmodules} \\
  \end{tabular}
\end{center}

\noindent
No alternative we examined at either level is immune to
nonextensibility.

We arrived at this semantics by trying each alternative in this section
and finding that it failed:
shallow merge could not compose nested extensions;
priority merge introduced ordering dependencies;
early binding froze references before composition could
supply them;
dynamic scope broke the substitution that the
$\lambda$-calculus embedding relies on;
single-target $\this$ rejected the multi-target compositions
that the Expression Problem demands.
Deep merge with multi-target late-binding $\this$ is what
remained.
\fi

\section{\bilingual{Related Work}{相关工作}}
\label{sec:related-work}

\subsection{\bilingual{Computational Models and \texorpdfstring{$\lambda$}{λ}-Calculus Semantics}{计算模型与 \texorpdfstring{$\lambda$}{λ}-演算语义}}

\paragraph{\bilingual{Minimal Turing-complete models}{最小图灵完备模型}}
\ifInheritanceChinese
Sch\"onfinkel~\cite{schonfinkel1924-combinatory-logic} 与 Curry~\cite{curry1958-combinatory-logic} 证明了 $\lambda$-演算可以归约为三个组合子 $S$、$K$、$I$%
\footnote{
  组合子 $I$ 是冗余的,因为 $I = SKK$。
}。将 $\lambda$-项编译为 SKI 的括号抽象是机械的,但会导致项大小的指数级膨胀;在 $\lambda$-演算中已有的计算模式之外,不会涌现出任何新的计算模式。在语义上,两者共享相同的模型:Meyer~\cite{meyer1982-model-of-lambda-calculus} 证明了满足五个附加条件的组合代数恰好就是 $\lambda$-代数。从未有人展示过一个独立的 SKI 语义域,并从中推导出 $\lambda$-演算的语义。方向始终是反过来的:Scott 的 $D_\infty$ 是为 $\lambda$-演算而构建的,而 SKI 随后被解释于其中。Dolan~\cite{dolan2013-mov-turing-complete} 证明了 x86 的 \texttt{mov} 指令单独就已图灵完备,其方法是利用寻址模式作为可变内存上的隐式算术;由此得到的程序是正确的,但不提供任何超越 RAM 机的抽象机制。图灵机本身虽然是通用的,却缺乏随机访问:读取距离为 $d$ 处的格需要 $d$ 次顺序的磁带移动。这些模型中的每一个,都在某个已有模型的语义框架内实现了图灵完备性,且对其中任何一个都从未展示过独立的语义域。
\else
Sch\"onfinkel~\cite{schonfinkel1924-combinatory-logic} and
Curry~\cite{curry1958-combinatory-logic} showed that the
$\lambda$-calculus reduces to three combinators $S$, $K$, $I$
\footnote{
  The $I$ combinator is redundant since $I = SKK$.
}.
The bracket abstraction that compiles $\lambda$-terms to SKI
is mechanical but incurs exponential blow-up in term size;
no new computational patterns emerge beyond those already
present in the $\lambda$-calculus.
Semantically, the two share the same models:
Meyer~\cite{meyer1982-model-of-lambda-calculus} showed that combinatory algebras
satisfying five additional conditions are exactly lambda
algebras.
No independent semantic domain for SKI has ever been
exhibited from which $\lambda$-calculus semantics could be
derived. The direction has always been the reverse:
Scott's $D_\infty$ was built for the $\lambda$-calculus,
and SKI was interpreted inside it afterwards.
Dolan~\cite{dolan2013-mov-turing-complete} showed that the x86 \texttt{mov}
instruction alone is Turing complete by exploiting
addressing modes as implicit arithmetic on mutable memory;
the resulting programs are correct but provide no
abstraction mechanisms beyond those of the RAM Machine.
The Turing Machine itself, while universal, lacks random
access: reading a cell at distance~$d$ requires $d$
sequential tape movements.
Each of these models achieves Turing completeness
within the semantic framework of an existing model,
and no independent semantic domain has been exhibited
for any of them.
\fi

%% === 审稿人备忘：SF-calculus 与 inheritance-calculus 的关系 ===
%%
%% Barry Jay 的 SF-calculus（Intensional Computation with Higher-Order
%% Functions, TCS 2019, doi:10.1016/j.tcs.2019.02.016；前身为
%% Confusion in the Church-Turing Thesis, arXiv:1410.7103, 2014）
%% 和 inheritance-calculus 都试图解决 λ-calculus 的 extensionality 限制，
%% 但方向完全不同。
%%
%% 共同出发点：λ-calculus 是 extensional 的——只能观察函数的输入输出行为，
%% 无法检查程序的内部结构。两个内部结构不同但行为相同的 λ 项不可区分。
%%   - Jay 的诊断：λ-calculus 缺乏 intensional 操作（检查项的结构），
%%     无法定义结构等式判定。
%%   - 本文的诊断：λ-calculus 的 function body 是 closed 的，无法从外部
%%     添加新的可观察投影（open extension）。
%%
%% 核心机制对比：
%%   SF-calculus：保留 S、K（或 λ），添加 F（factorisation）算子，
%%     可以将闭合范式拆解为组成部分。程序仍然是函数，只是多了检查函数
%%     结构的能力。属于 λ-calculus 的 *扩展*。
%%   Inheritance-calculus：用三个不同的原语（record、definition、
%%     inheritance）替代 λ-calculus 的三个原语。β-reduction 是 deep merge
%%     的特例（Section 3 的翻译 T 将函数应用映射为继承）。
%%     属于 λ-calculus 的 *包含*（subsumption）。
%%
%% 关于 F 算子的编码：
%%   F 做两件事：(1) 判断结构形态（atom vs. compound），
%%   (2) 提取子部分。
%%   在 inheritance-calculus 中：
%%     - 提取子部分是平凡的：所有值都是透明 record，通过 qualified this
%%       投影任意 slot 即可。Theorem 6 的 separating context 就在读取
%%       argument slot（↑².argument）。
%%     - 判断结构形态并分支：需要 Acceptance pattern（Section 4）。
%%       NatVisitor 对 Nat 做的事情与 F 对 SF-calculus 项做的事情
%%       完全对应：VisitZero/SuccessorVisitor 对应 F 的 atom/compound
%%       分支。Visitor 可以通过 open extension 添加到已有 term 上，
%%       不修改原始定义。
%%   Acceptance 在精神上就是 F：两者都是对 closed normal form 做
%%   constructor dispatch。区别在于 F 是不可扩展的 built-in primitive
%%   （硬编码了 S 和 F 两个 atom），而 Acceptance 是可扩展的 encoding
%%   （新增 constructor 只需继承新的 Accepted* 分支，即 Expression Problem
%%   immunity）。
%%
%% 精确差异：
%%   SF-calculus 把 intensional observation 做成了 object-level 的
%%   primitive reduction rule（F 是一个 term）；
%%   inheritance-calculus 把它留在了 meta-level（观察者总能通过
%%   properties(p) 查询任何路径），而 term-level 的分支需要通过
%%   Acceptance pattern 显式编码。
%%   换言之：inheritance-calculus 中 F 的 *读取* 能力是免费的
%%   （所有值都是透明 record），但 F 的 *分支* 能力需要 Acceptance。
%%   SF-calculus 的 F 则把读取和分支捆绑在一个不可扩展的 primitive 里。
%%   Acceptance 是可扩展的（新增 constructor 只需继承新分支，
%%   即 Expression Problem immunity），而 F 硬编码了固定的 atom 集合
%%   （S 和 F）。
%%
%% 关于 Jay 所谓的 "trivial solution of the Halting Problem"：
%%   Jay 用非标准编码将每个程序（包括不停机的）映射到闭合范式（静态数据），
%%   然后结构等式判定变得平凡。这并不推翻图灵定理——它改变了问题的
%%   形式化方式（把程序当静态数据而非可执行代码），而非解决了原来的问题。
%%   Jay 的真正论点是：λ-definability 和 Turing computability 不是同一回事，
%%   Church-Turing Thesis 的等价性证明隐含了编码假设。
%%
%% 正文篇幅不够论述 embeds SF-calculus，暂不展开。
%% 如果审稿人提到 SF-calculus，可以基于上述分析撰写回复或补充段落。
%% ================================================================

\paragraph{\bilingual{Semantic descriptions of the $\lambda$-calculus}{$\lambda$-演算的语义描述}}
\ifInheritanceChinese
$\lambda$-演算的语义可以用几种根本上不同的方式来定义。\emph{操作语义}将 $\beta$-归约 $(\lambda x.\,M)\,N \to M[N/x]$ 作为计算的基本概念~\cite{church1936-lambda-calculus,church1941-calculi-of-lambda-conversion,barendregt1984-lambda-calculus}。这要求捕获-回避替换,其微妙之处促使了 de~Bruijn 索引~\cite{debruijn1972-nameless-dummies} 与显式替换演算~\cite{abadi1991-explicit-substitutions} 的出现。它还要求选择求值策略:按名调用与按值调用~\cite{plotkin1975-call-by-name-call-by-value},或 Abramsky 的惰性策略~\cite{abramsky1990-lazy-lambda-calculus}。\emph{指称语义}~\cite{scott1970-mathematical-theory-of-computation,scott1976-data-types-as-lattices} 将项解释于一个数学域中,但赋予 $\beta$-归约意义需要一个自反域 $D \cong [D \to D]$,由 Cantor 定理知该域无集合论解。Scott 的 $D_\infty$ 通过将 $D$ 构造为连续格的逆极限~\cite{scott1972-continuous-lattices} 来解决这一问题。由此得到的语义是适当的,但不是完全抽象的~\cite{plotkin1977-lcf-as-programming-language,milner1977-fully-abstract-models}:它将观察等价所区分的项等同了起来。\emph{博弈语义}~\cite{abramsky2000-full-abstraction-pcf,hyland2000-full-abstraction-pcf,nickau1994-hereditarily-sequential-functionals} 对 PCF 与惰性 $\lambda$-演算~\cite{abramsky1993-full-abstraction-lazy-lambda,abramsky1995-games-full-abstraction-lazy-lambda} 实现了完全抽象,需要竞技场、策略与无害性条件。在这些方法的每一种中,函数空间 $[D \to D]$ 都是核心障碍:操作语义通过替换隐式地需要它,而指称语义则显式地构造它,博弈语义则通过策略加以精化。这些方法中的每一种都涉及高阶结构%
\footnote{
  指称语义中的函数空间,博弈语义中的策略,以及操作语义中的替换。
}
并在其语义域中保留了函数空间。在每种情形下,指称帐和操作帐都是相互独立的形式体系。三者都继承了 $\beta$-归约在递归程序上的发散:例如,有向图上一个朴素的传递闭包在每个框架下都会发散。
\else
The semantics of the $\lambda$-calculus can be defined in
several fundamentally different ways.
\emph{Operational semantics} takes $\beta$-reduction
$(\lambda x.\,M)\,N \to M[N/x]$ as the primitive notion of
computation~\cite{church1936-lambda-calculus,church1941-calculi-of-lambda-conversion,barendregt1984-lambda-calculus}.
This requires capture-avoiding substitution, a mechanism whose
subtleties motivated de~Bruijn
indices~\cite{debruijn1972-nameless-dummies} and explicit substitution
calculi~\cite{abadi1991-explicit-substitutions}.
It also requires a choice of evaluation
strategy: call-by-name versus
call-by-value~\cite{plotkin1975-call-by-name-call-by-value} or the lazy
strategy of Abramsky~\cite{abramsky1990-lazy-lambda-calculus}.
\emph{Denotational semantics}~\cite{scott1970-mathematical-theory-of-computation,scott1976-data-types-as-lattices}
interprets terms in a mathematical domain, but giving meaning
to $\beta$-reduction requires a reflexive domain
$D \cong [D \to D]$, which has no set-theoretic solution by
Cantor's theorem. Scott's $D_\infty$ resolves this by
constructing $D$ as an inverse limit of continuous
lattices~\cite{scott1972-continuous-lattices}.
The resulting semantics is adequate but not fully
abstract~\cite{plotkin1977-lcf-as-programming-language,milner1977-fully-abstract-models}:
it identifies terms that observational equivalence
distinguishes.
\emph{Game semantics}~\cite{abramsky2000-full-abstraction-pcf,hyland2000-full-abstraction-pcf,nickau1994-hereditarily-sequential-functionals}
achieves full abstraction for PCF and for the lazy
$\lambda$-calculus~\cite{abramsky1993-full-abstraction-lazy-lambda,abramsky1995-games-full-abstraction-lazy-lambda},
requiring arenas, strategies, and innocence conditions.
In each of these approaches, the function space $[D \to D]$
is the central obstacle: operational semantics needs it
implicitly via substitution. In contrast, denotational semantics
constructs it explicitly, and game semantics refines it via
strategies.
Each of these approaches involves higher-order structure%
\footnote{
  Function spaces
  in denotational semantics, strategies
  in game semantics, and substitution
  in operational semantics.
}
and retains the function space in its semantic domain.
The denotational and operational accounts remain separate
formalisms in each case.
All three inherit $\beta$-reduction's divergence on
recursive programs: a na\"ive transitive closure on a cyclic
graph, for instance, diverges in every framework.
\fi

\paragraph{\bilingual{Set-theoretic models of the $\lambda$-calculus}{$\lambda$-演算的集合论模型}}
\ifInheritanceChinese
$\lambda$-演算的几个模型族通过在集合论而非域论框架下工作,避免了 Scott 的 $D_\infty$ 的逆极限构造。这些模型消除了 Scott 拓扑,但并未消除函数空间本身:它们仍然是高阶的,通过注入、交叉类型或多重集关系将 $[D \to D]$ 编码于 $D$ 之中。\emph{图模型}~\cite{engeler1981-algebras-and-combinators,plotkin1993-set-theoretic-lambda-models} 通过幂集格上的 web 注入来编码函数;Bucciarelli 与 Salibra~\cite{bucciarelli2008-graph-lambda-theories} 证明了最大的合理图理论等于 Böhm 树理论 $\mathcal{B}$。\emph{过滤器模型}~\cite{barendregt1983-filter-lambda-model,coppo1978-intersection-type-assignment} 将一个项解释为可从中推导出的交叉类型集合;Dezani-Ciancaglini 等人~\cite{dezani1998-intersection-types-bohm-trees} 证明了 BCD 过滤器模型的 $\lambda$-理论恰好与 Böhm 树相等性重合。\emph{关系模型}~\cite{ehrhard2005-finiteness-spaces,bucciarelli2001-phase-semantics-exponentials} 通过有限多重集从线性逻辑的关系语义中涌现;Ehrhard~\cite{ehrhard2005-finiteness-spaces} 证明了在这一范畴中,纯 $\lambda$-演算的标准自反对象不能存在。这三个模型族都是语义模型:它们将 $\lambda$-项解释于一个外部数学结构中,并在该结构中将 $\beta$-归约验证为一个等式,而不是定义一个具有自身原语的独立演算。每一种都编码函数空间而非消除它,且每一种都需要高阶机制才能为惰性 $\lambda$-演算获得完全抽象。
\else
Several families of $\lambda$-calculus models avoid the
inverse-limit construction of Scott's $D_\infty$ by working
in set-theoretic rather than domain-theoretic settings.
These models eliminate Scott topology but not the function
space itself: they remain higher-order, encoding $[D \to D]$
inside $D$ via injections, intersection types, or multiset
relations.
\emph{Graph models}~\cite{engeler1981-algebras-and-combinators,plotkin1993-set-theoretic-lambda-models}
encode functions via a web injection on powerset lattices;
Bucciarelli and Salibra~\cite{bucciarelli2008-graph-lambda-theories}
proved that the greatest sensible graph theory equals the
B\"ohm tree theory~$\mathcal{B}$.
\emph{Filter models}~\cite{barendregt1983-filter-lambda-model,coppo1978-intersection-type-assignment}
interpret a term as the set of intersection types derivable
for it;
Dezani-Ciancaglini et al.~\cite{dezani1998-intersection-types-bohm-trees}
showed that the BCD filter model's $\lambda$-theory
coincides exactly with B\"ohm tree equality.
\emph{Relational models}~\cite{ehrhard2005-finiteness-spaces,bucciarelli2001-phase-semantics-exponentials}
arise from the relational semantics of linear logic via
finite multisets;
Ehrhard~\cite{ehrhard2005-finiteness-spaces} showed that a standard
reflexive object for the pure $\lambda$-calculus cannot
exist in this category.
All three families are semantic models: they interpret
$\lambda$-terms in an external mathematical structure and
validate $\beta$-reduction as an equation in that structure,
rather than defining an independent calculus with its own
primitives.
Each encodes the function space rather than eliminating it,
and each requires higher-order machinery to obtain full
abstraction for the lazy $\lambda$-calculus.
\fi

\paragraph{\bilingual{Recursive equations and fixpoint semantics}{递归方程与不动点语义}}
\ifInheritanceChinese
Van~Emden 与 Kowalski~\cite{vanemden1976-predicate-logic-semantics} 证明了谓词逻辑作为编程语言的语义允许三种等价的刻画:操作性的\footnote{通常是 SLD 归结。}、模型论的\footnote{基于最小 Herbrand 模型。}以及不动点的\footnote{即基原子幂集上即时后承算子 $T_P$ 的最小不动点。}。Aczel~\cite{aczel1977-inductive-definitions} 通过幂集格上的单调算子对归纳定义给出了系统性处理,为 Datalog 风格的语义所依赖的逻辑基础奠定了基础。Leroy 与 Grall~\cite{leroy2006-coinductive-big-step} 通过余归纳定义处理了大步语义中的发散,需要在归纳求值判断之外引入一个单独的余归纳判断来说明非终止。Van~Emden 与 Kowalski 的 $T_P$ 在一阶域上为递归 Datalog 程序计算 $\mathrm{lfp}$。传统的 $\lambda$-演算语义在语义域中需要函数空间,且递归的 $\lambda$-项在 $\beta$-归约下通常会发散;利用一阶域上的 $T_P$ 为 $\lambda$-演算赋予一种改变哪些项收敛的不动点语义,尚未被探索。在近期工作中,Yang~\cite{yang2026-co-lambda} 将 tabling 求值~\cite{tamaki1986-tabled-resolution} 直接引入纯 $\lambda$-演算:首部规范化被读作一个余代数方程,其状态为被驻留的 $\lambda$-项,即具有可判定同一性的一阶对象,而 tabling 求解该方程,在有限时间内将 $\Omega$ 这样的非生产性循环判定为 $\bot$。该处理保留了标准的惰性 (Lévy--Longo) 树语义,仅改变了求值器;它占据了图~\ref{fig:calculi-matrix} 中不动点 $\lambda$ 节点。它在机制上也有所不同。表以整个被驻留的项为键,而非以进入作用域树的路径为键,因此在路径前缀上重叠的查询之间不共享部分结果。偏序不是集合包含下的幂集格,而是树近似序,其中 $\bot$ 叶被精化为子树,且不动点是唯一的(最小与最大因有护性而重合):一个重入调用被以单一的最终 $\bot$ 回答,而后续任何迭代都不会精化该 $\bot$,因此迭代无法增大一个累加器。其比惰性 $\lambda$-演算更确定的部分因此恰好是图~\ref{fig:calculi-matrix} 中非生产性的 $\bot$,即发散被判定为确定 $\bot$ 答案的那些项,别无其他;而幂集序才使得迭代可以积累并产生新的非 $\bot$ 值。
\else
Van~Emden and Kowalski~\cite{vanemden1976-predicate-logic-semantics} showed that
the semantics of predicate logic as a programming language admits
three equivalent characterizations:
operational\footnote{
  Typically SLD resolution.
}, model-theoretic\footnote{
  Based on least Herbrand models.
}, and fixed-point\footnote{
  Namely, the least fixed point of the immediate consequence operator~$T_P$ on the powerset of ground atoms.
}.
Aczel~\cite{aczel1977-inductive-definitions} gave a systematic treatment of
inductive definitions via monotone operators on powerset lattices,
providing the logical foundations that Datalog-style semantics
rests on.
Leroy and Grall~\cite{leroy2006-coinductive-big-step} treated divergence
in big-step semantics via coinductive definitions,
requiring a separate coinductive judgment to account for
non-termination alongside the inductive evaluation judgment.
Van~Emden and Kowalski's $T_P$ computes
$\mathrm{lfp}$ for recursive Datalog programs over a
first-order domain.
Conventional $\lambda$-calculus semantics requires function
spaces in the semantic domain, and recursive $\lambda$-terms
ordinarily diverge under $\beta$-reduction;
using $T_P$ over a first-order domain to give the
$\lambda$-calculus a fixpoint semantics that changes which
terms converge has not been explored.
In recent work, Yang~\cite{yang2026-co-lambda} brings tabled
evaluation~\cite{tamaki1986-tabled-resolution} to the pure
$\lambda$-calculus directly: head normalisation is read as a
coalgebraic equation whose states are interned
$\lambda$-terms, first-order objects with decidable identity,
and tabling solves it, deciding unproductive loops such
as~$\Omega$ as~$\bot$ in finite time.
That treatment preserves the standard lazy
(L\'evy--Longo) tree semantics and changes only the evaluator;
it occupies the fixpoint~$\lambda$ node of
Figure~\ref{fig:calculi-matrix}.
It also differs in mechanism.
The table is keyed by whole interned terms, not by paths
into a scope tree, so partial results are not shared between
queries that overlap on a path prefix.
The order is not a powerset lattice under set inclusion but
the tree approximation order, in which $\bot$ leaves are
refined into subtrees, and the fixpoint is unique (least and
greatest coincide by guardedness): a re-entrant call is
answered with a single final~$\bot$ that no later iteration
refines, so the iteration cannot grow an accumulator.
Its more-definedness over the lazy $\lambda$-calculus is
therefore exactly the unproductive~$\bot$ of
Figure~\ref{fig:calculi-matrix}, terms whose divergence is
decided as a definite~$\bot$ answer, and nothing more;
the powerset order is what lets iteration accumulate and
produce new non-$\bot$ values.
\fi

\subsection{\bilingual{Inheritance, Objects, and Records}{继承、对象与记录}}

\paragraph{\bilingual{Denotational semantics of inheritance}{继承的指称语义}}
\ifInheritanceChinese
Cook~\cite{cook1989-denotational-inheritance} 给出了继承的第一个指称语义,将对象建模为递归记录。继承是\emph{生成器}%
\footnote{
  这些是从 self 到完整对象的函数。
}与\emph{包装器}%
\footnote{
  包装器是修改生成器的函数。
}的复合,由一个不动点构造来求解。Cook 的关键洞见是,继承是一种可适用于任何形式的递归定义的通用机制,而不仅仅是面向对象的方法。这一洞见是本工作的出发点之一。该模型有四个原语(记录、函数、生成器和包装器),并需要一个不动点组合子。复合是不对称的:包装器修改生成器,反之则不然。不动点构造依赖函数域中的 Scott 连续性。由此得到的指称不能直接执行,因此需要一个单独的操作语义来在运行时定义方法分派与对象构造。
\else
Cook~\cite{cook1989-denotational-inheritance} gave the first denotational
semantics of inheritance, modeling objects as recursive records.
Inheritance is composition of \emph{generators}
\footnote{
  These are functions from self to complete object.
} and
\emph{wrappers}\footnote{
  Wrappers are functions that modify generators.
},
solved by a fixed-point construction.
Cook's key insight is that inheritance is a general mechanism
applicable to any form of recursive definition, not only
object-oriented methods. This insight is one of the starting points of the
present work.
The model has four primitives (records, functions, generators, and
wrappers) and requires a fixed-point combinator. Composition is
asymmetric: a wrapper modifies a generator but not vice versa.
The fixed-point construction relies on Scott continuity in a
domain of functions. The resulting denotations are not directly
executable, so a separate operational semantics is needed to
define method dispatch and object construction at runtime.
\fi

\paragraph{\bilingual{Mixins and traits}{Mixin 与 trait}}
\ifInheritanceChinese
Bracha 与 Cook~\cite{bracha1990-mixin-inheritance} 将 mixin 形式化为抽象子类,即从超类参数到子类的函数。由于 mixin 应用是函数复合,它\emph{既不可交换也不幂等};C3 线性化算法~\cite{barrett1996-c3-linearization} 后来被开发出来以施加一个确定性顺序。Schärli 等人~\cite{scharli2003-traits} 引入了 trait,其复合是可交换的,但拒绝同名冲突~\cite{ducasse2006-traits-fine-grained-reuse}。两种机制都在\emph{方法层级}操作:对同名的两个定义进行复合会导致覆盖或冲突,而非嵌套结构的递归合并。两者都预设了 $\lambda$-演算:方法体是函数,而复合机制本身不是图灵完备的。Mixin 与 trait 均不支持深合并。
\else
Bracha and Cook~\cite{bracha1990-mixin-inheritance} formalized mixins as
abstract subclasses, that is, functions from a superclass parameter to a
subclass.
Because mixin application is function composition, it is
\emph{neither commutative nor idempotent}; the C3 linearization
algorithm~\cite{barrett1996-c3-linearization} was later developed to impose a
deterministic order.
Sch\"arli et al.~\cite{scharli2003-traits} introduced traits,
whose composition is commutative but rejects same-name
conflicts~\cite{ducasse2006-traits-fine-grained-reuse}.
Both mechanisms operate at the \emph{method level}: composing two
definitions of the same name results in an override or a conflict,
not a recursive merge of nested structure.
Both presuppose the $\lambda$-calculus: method bodies are functions,
and the composition mechanism itself is not Turing complete.
Neither mixins nor traits are deep-mergeable.
\fi

\paragraph{\bilingual{Object and record calculi}{对象演算与记录演算}}
\ifInheritanceChinese
Abadi 与 Cardelli~\cite{abadi1996-theory-of-objects} 将对象作为基本概念,但方法是以 self 为参数的函数,且对象扩展是不对称的。Boudol~\cite{boudol2004-recursive-record-semantics} 证明了自引用记录需要一个不安全的不动点算子,其适定性依赖于求值策略。Harper 与 Pierce~\cite{harper1991-symmetric-record-concatenation} 给出了一个具有对称连接的记录演算,该演算拒绝同标签的复合;Cardelli~\cite{cardelli1992-extensible-records} 研究了带子类型的可扩展记录;Rémy~\cite{remy1989-row-polymorphism} 为 ML 中的记录和变体引入了行多态。这些演算全都在\emph{平坦}记录上操作:没有嵌套结构的递归合并,没有惰性观察,$\this$ 也需要一个显式的不动点组合子。
\else
Abadi and Cardelli~\cite{abadi1996-theory-of-objects} treat objects as the
primitive notion, but methods are functions that take self as a
parameter and object extension is asymmetric.
Boudol~\cite{boudol2004-recursive-record-semantics} showed that self-referential
records require an unsafe fixed-point operator whose
well-definedness depends on the evaluation strategy.
Harper and Pierce~\cite{harper1991-symmetric-record-concatenation} gave a record
calculus with symmetric concatenation that rejects same-label
composition;
Cardelli~\cite{cardelli1992-extensible-records} studied extensible records
with subtyping;
R\'emy~\cite{remy1989-row-polymorphism} introduced row polymorphism for
records and variants in ML.
All of these calculi operate on \emph{flat} records: there is
no recursive merging of nested structure, no lazy observation,
and $\this$ requires an explicit fixed-point combinator.
\fi

\paragraph{\bilingual{Family polymorphism and DOT}{族多态与 DOT}}
\ifInheritanceChinese
Ernst~\cite{ernst2001-family-polymorphism} 引入了族多态;Ernst、Ostermann 与 Cook~\cite{ernst2006-virtual-class-calculus} 在虚类演算中将其形式化,其中类通过路径表达式 \texttt{this.out.C} 来访问。Amin 等人~\cite{amin2016-dependent-object-types} 在 DOT 演算中形式化了路径依赖类型,这是 Scala 的理论基础。在这两个框架中,每条路径都解析为\emph{单一}封闭对象中的\emph{单一}类。当多条继承路径通向同一作用域时,单值解析必须选择其中一条路径并丢弃其余;DOT 的 self 变量 $\nu(x : T)\, d$ 将 $x$ 绑定到单一对象,而 Scala 在编译时拒绝由此产生的交叉类型(附录~\ref{app:scala-multi-target})。两个框架均未提供集合值解析的机制。
\else
Ernst~\cite{ernst2001-family-polymorphism} introduced family polymorphism;
Ernst, Ostermann, and Cook~\cite{ernst2006-virtual-class-calculus}
formalized it in the virtual class calculus, where classes are
accessed via path expressions \texttt{this.out.C}.
Amin et al.~\cite{amin2016-dependent-object-types} formalized path-dependent types
in the DOT calculus, the theoretical foundation of Scala.
In both frameworks, each path resolves to a
\emph{single} class in a \emph{single} enclosing object.
When multiple inheritance routes lead to the same scope,
single-valued resolution must choose one route and discard the
others; DOT's self variable $\nu(x : T)\, d$ binds $x$ to a
single object, and Scala rejects the resulting intersection
type at compile time
(Appendix~\ref{app:scala-multi-target}).
Neither framework provides a mechanism for set-valued
resolution.
\fi

\subsection{\bilingual{Configuration Languages and Module Systems}{配置语言与模块系统}}

\paragraph{\bilingual{Configuration languages}{配置语言}}
\ifInheritanceChinese
CUE~\cite{cue2019} 在单一格中统一了类型与值,其中以 \texttt{\&} 表示的合一是可交换、可结合且幂等的。CUE 的设计受到 Google 的 Borg 配置语言 (GCL) 的启发,后者使用了图合一。CUE 的格包含标量类型,其冲突语义规定合一不相容标量会得到 $\bot$;标量对于配置语言是否必要,还是仅凭树结构就已足够,是一个有趣的设计问题。Jsonnet~\cite{jsonnet} 通过 \texttt{+} 运算符提供了具有 mixin 语义的对象继承:复合\emph{不是}可交换的,因为在标量冲突上右侧胜出,而深合并需要在每个字段上通过 \texttt{+:} 语法显式选择加入,因此深合并虽可用但并非默认行为。Dhall~\cite{dhall2017} 采取函数式方法:它是一个带记录的类型化 $\lambda$-演算,其中复合是记录合并算子,而非继承机制。
\else
CUE~\cite{cue2019} unifies types and values in a single lattice
where unification, denoted by \texttt{\&}, is commutative, associative, and
idempotent.
CUE's design was motivated by Google's Borg Configuration
Language (GCL), which used graph unification.
CUE's lattice includes scalar types with a conflict
semantics where unifying incompatible scalars yields $\bot$; whether
scalars are necessary for a configuration language, or whether
tree structure alone suffices, is an interesting design question.
Jsonnet~\cite{jsonnet} provides object inheritance via the
\texttt{+} operator with mixin semantics: composition is
\emph{not} commutative since the right-hand side wins on scalar
conflicts, and deep merging requires explicit opt-in via the \texttt{+:}
syntax on a per-field basis, so deep merge is
available but not the default.
Dhall~\cite{dhall2017} takes a functional approach: it is a typed
$\lambda$-calculus with records, where composition is a record merge
operator, not an inheritance mechanism.
\fi

\paragraph{\bilingual{Module systems and deep merge}{模块系统与深合并}}
\ifInheritanceChinese
NixOS 模块系统~\cite{dolstra2010-nixos,nixosmodules} 通过\emph{递归合并嵌套属性集}来组合模块,而非通过线性化或方法层级的覆盖,这与先前的 mixin 和 trait 演算中区分树层级继承与方法层级复合的深合并语义相同。当两个模块定义了相同的嵌套路径时,其子树被合并而非覆盖。NixOS 模块系统集成了深合并、递归的 $\this$、模块的延迟求值,以及其上的丰富类型检查层。同一机制为 NixOS 系统配置、Home Manager~\cite{homemanager}、nix-darwin~\cite{nixdarwin}、flake-parts~\cite{flakeparts} 以及数十个其他生态系统提供支持,共同管理任意复杂度的配置。一个似乎尚未被完全解决的方面是 $\this$ 在多重继承下的语义。当两个各自继承自同一延迟模块的独立模块被组合时,每次导入都会引入对共享选项自己的声明;模块系统的 \texttt{merge\-Option\-Decls} 以静态错误拒绝重复声明(附录~\ref{app:scala-multi-target})。模块系统无法表达两条路径对同一选项的贡献是均等的,即递归合并自然产生的多目标情形,正是这一情形促成了第~\ref{sec:mixin-trees} 节的观测语义。
\else
The NixOS module system~\cite{dolstra2010-nixos,nixosmodules}
composes modules by \emph{recursively merging nested attribute sets},
not by linearization or method-level override, which is the same deep-merge
semantics that distinguishes tree-level inheritance from method-level
composition in prior mixin and trait calculi.
When two modules define the same nested path, their subtrees are
merged rather than overridden.
The NixOS module system incorporates deep merge,
recursive $\this$, deferred evaluation of modules, and a
rich type-checking layer on top.
The same mechanism powers
NixOS system configuration, Home
Manager~\cite{homemanager}, nix-darwin~\cite{nixdarwin},
flake-parts~\cite{flakeparts}, and dozens of other ecosystems,
collectively managing configurations of arbitrary complexity.
The one aspect that appears not to have been fully addressed is the
semantics of $\this$ under multiple inheritance.
When two independent modules that inherit from the same
deferred module are composed, each import introduces its own
declaration of the shared options; the module system's
\texttt{merge\-Option\-Decls} rejects duplicate declarations
with a static error (Appendix~\ref{app:scala-multi-target}).
The module system cannot express that both routes contribute
equally to the same option, that is, the multi-target situation that
recursive merging naturally gives rise to, which motivated the
observational semantics of Section~\ref{sec:mixin-trees}.
\fi

\subsection{\bilingual{Our Companion Work}{我们的配套工作}}

\paragraph{\bilingual{Modular software architectures}{模块化软件架构}}
\label{sec:practical-modular}
\ifInheritanceChinese
插件系统、组件框架与依赖注入框架通过声明式继承实现了软件的可扩展性。第~\ref{sec:mixin-trees} 节所确立的继承的对称性与可结合性解释了这类系统为何有效:组件可以以任意顺序继承而不改变行为。MIXINv2\anon[~(补充材料)]{~\cite{mixin2025}} 本身就是一个依赖注入框架:每个作用域将其依赖声明为具名槽,复合跨越作用域边界按名解析它们,从而使继承演算的模式显式化而非随意拼凑。
\else
Plugin systems, component frameworks, and dependency injection frameworks
enable software extensibility through declarative inheritance. The
symmetric, associative nature of inheritance established in Section~\ref{sec:mixin-trees} explains why such
systems work: components can be inherited in any order without changing
behavior.
MIXINv2\anon[~(supplementary material)]{~\cite{mixin2025}}
is itself a dependency injection framework:
each scope declares its dependencies as named slots, and
composition resolves them by name across scope boundaries,
making the inheritance-calculus patterns explicit rather than ad~hoc.
\fi

\paragraph{\bilingual{Union file systems}{联合文件系统}}
\label{sec:practical-union-fs}
\ifInheritanceChinese
联合文件系统,如 UnionFS、OverlayFS 与 AUFS,将多个目录层次叠加以呈现统一视图。其语义与继承演算相似,表现在后面的层覆盖前面的层且所有层的文件均可访问,但其设计在具体之处有所不同:继承是不对称的,标量文件的冲突解决是不可交换的,且层与层之间不存在晚绑定或动态分派。我们将继承演算直接实现为联合文件系统\anon[~(作为补充材料附上)]{~\cite{ratarmount2025}},仅凭三个原语就获得了文件查询的晚绑定语义与动态分派,无需任何额外机制。在软件包管理器、构建系统与操作系统发行版中,外部工具链将配置转换为文件树;而在这里,文件系统作为自身的配置,转换即是继承:这是一种\emph{文件系统即编译器},其中目录树既是源程序又是编译目标。
\else
Union file systems, such as UnionFS, OverlayFS, and AUFS, layer multiple directory
hierarchies to present a unified view. Their semantics resembles
inheritance-calculus in that later layers override earlier layers and files from all
layers remain accessible, but their designs differ in specific ways:
inheritance is asymmetric, conflict resolution for scalar files
is noncommutative, and there is no late-binding or dynamic
dispatch across layers.
We implemented inheritance-calculus directly as a union file
system\anon[~(included as supplementary material)]{~\cite{ratarmount2025}},
obtaining late-binding semantics and dynamic dispatch for file lookups
with no additional mechanism beyond the three primitives.
In package managers, build systems, and OS distributions,
external toolchains transform configuration into file trees;
here, the file system serves as configuration to itself,
and the transformation is inheritance: a
\emph{file-system-as-compiler} in which the directory tree is
both the source program and the compilation target.
\fi

\section{\bilingual{Future Work}{未来工作}}
\label{sec:future-work}

\ifInheritanceChinese
\anon[辅助实现]{MIXINv2~\cite{mixin2025}}{} 是继承演算的一个可执行实现,它已经超越了本文所呈现的无类型演算:它包含编译期检查,以验证所有引用都解析到 mixin 树中的有效路径,以及一个从宿主语言引入标量值的外部函数接口 (FFI)。以下所描述的未来工作,关注的是这一实现与一门完全实用的语言之间的差距,横跨理论基础与库层面的编码两个方面。
\else
\anon[The supplementary implementation]{MIXINv2~\cite{mixin2025}}{} is an executable implementation of inheritance-calculus
that already goes beyond the untyped calculus presented here:
it includes compile-time checking that all references resolve to
valid paths in the mixin tree, and a foreign-function interface (FFI)
that introduces scalar values from the host language.
The future work described below concerns the distance between this
implementation and a fully practical language, spanning both
theoretical foundations and library-level encodings.
\fi

\subsection{\bilingual{Type System}{类型系统}}
\ifInheritanceChinese
MIXINv2 现有的编译期检查验证每条引用路径都解析到继承结果中存在的属性。将这些检查的可靠性形式化,即证明类型良好的程序在运行时不产生悬空引用,是未来工作的一个方向。这样的形式化与 DOT 演算~\cite{amin2016-dependent-object-types} 的关系,类似于继承演算与 $\lambda$-演算的关系:有类型的继承演算可以看作不含 $\lambda$ 的 DOT,它保留了路径依赖类型,同时以 mixin 继承取代函数。正式的类型系统将适合在 Coq 和 Agda 等证明助手中机械化,以验证类型安全、进展和保持等性质。

在现有检查的可靠性之外,额外的类型系统特性将强化这门语言。一个关键的例子是\emph{完全性检查}:验证一个继承为所有必需槽位提供了实现,而不仅仅是引用能够解析。在无类型的继承演算中,写下 $\mathrm{magnitude} \mapsto \{\}$ 纯属说明性质:它定义了一个结构槽位,但不对填充它的内容施加任何约束。完全性检查将在静态上强制这类声明,拒绝使必需槽位保持未填充状态的继承。
\else
MIXINv2's existing compile-time checks verify that every
reference path resolves to a property that exists in the inherited
result. Formalizing the soundness of these checks, that is, proving that
well-typed programs do not produce dangling references at
runtime, is one direction of future work.
Such a formalization would relate to the DOT
calculus~\cite{amin2016-dependent-object-types} in a manner analogous to how
inheritance-calculus relates to the $\lambda$-calculus: a Typed
inheritance-calculus can be viewed as DOT without $\lambda$, retaining
path-dependent types while replacing functions with
mixin inheritance.
A formal type system would be amenable to mechanization in proof
assistants, such as Coq and Agda, for verifying type safety, progress, and
preservation properties.

Beyond soundness of the existing checks, additional type system
features would strengthen the language.
A key example is \emph{totality checking}: verifying that an
inheritance provides implementations for all required slots, not
merely that references resolve.
In untyped inheritance-calculus, writing $\mathrm{magnitude} \mapsto \{\}$
is purely documentary: it defines a structural slot but imposes no
constraint on what must fill it.
Totality checking would enforce such declarations statically,
rejecting inheritances that leave required slots unfilled.
\fi

\subsection{\bilingual{Standard Library}{标准库}}
\ifInheritanceChinese
其余几项涉及标准库:它们不需要对演算或类型系统进行扩展,但需要在现有框架内进行非平凡的编码。

继承演算\emph{默认是开放的}:继承可以自由合并任意两个 mixin,一个值可以同时具有多个数据构造子(例如,$\{\mathrm{Zero},\; \mathrm{Odd},\; \mathrm{half} \mapsto \{\mathrm{Zero}\}\}$)。这是该演算自然的 trie 语义,也是第~\ref{sec:expression-problem}~节求解表达式问题的基础。然而,许多操作(例如等值测试)假设值是\emph{线性的},即恰好具有一个数据构造子。第~\ref{sec:case-study}~节的观察者模式并不强制这一点:将其应用于一个具有多个数据构造子的值时,所有回调都会触发,其结果被继承在一起。

\emph{模式验证}可以在库层面使用现有原语来实现。观察者可以测试某个特定的数据构造子是否存在,返回一个布尔值;第二次布尔派发随后基于该结果进行分支。通过链接这样的测试,工厂可以验证一个值恰好匹配一个数据构造子并携带所需的属性。这并不限制继承本身(继承仍然不受约束),但提供了一种在不变式违反传播之前检测到它的方式。

同样的技术,即观察者返回布尔值后接布尔派发,可以产生\emph{封闭模式匹配}。第~\ref{sec:case-study}~节的观察者模式本质上是\emph{开放的}:新的数据构造子情况可以通过继承添加。封闭匹配(其中恰好执行一个分支)可以编码为一个\emph{访问者链}:一个由 if-then-else 节点组成的链表,每个节点测试一个数据构造子,在不匹配时向后传递。这类似于 GHC.Generics 的 $\mathrm{(:+:)}$ 和类型表示~\cite{magalhaes2010-generic-deriving},其中每个链节对应一个求和项。一个重要的推论是,这样的链\emph{不}可交换,因为它们具有固定的顺序,因此无法通过开放继承来扩展,而这正是封闭派发的正确语义。
\else
The remaining items concern the standard library: they require
no extensions to the calculus or type system, but involve nontrivial
encodings within the existing framework.

Inheritance-calculus is \emph{open by default}: inheritance
can freely merge any two mixins, and a value may simultaneously
inhabit multiple data constructors (e.g.,
$\{\mathrm{Zero},\; \mathrm{Odd},\; \mathrm{half} \mapsto \{\mathrm{Zero}\}\}$).
This is the natural trie semantics of the
calculus, and the basis for solving the Expression Problem
in Section~\ref{sec:expression-problem}.
However, many operations, such as equality testing, assume that values
are \emph{linear}, inhabiting exactly one data constructor.
The observer pattern of Section~\ref{sec:case-study} does not enforce
this: applied to a multi-data-constructor value, all callbacks fire and
their results are inherited.

\emph{Schema validation} can be implemented at the library level
using existing primitives.
An observer can test whether a particular data constructor is present,
returning a Boolean; a second Boolean dispatch then branches on the
result.
By chaining such tests, a factory can validate that a value matches
exactly one data constructor and carries the required properties.
This does not restrict inheritance itself, which remains
unconstrained, but provides a way to detect invariant violations
before they propagate.

The same technique, an observer returning Boolean followed by Boolean
dispatch, yields \emph{closed pattern matching}.
The observer pattern of Section~\ref{sec:case-study} is inherently
\emph{open}: new data-constructor cases can be added through inheritance.
Closed matching, where exactly one branch is taken, can be encoded as
a \emph{visitor chain}: a linked list of if-then-else nodes, each
testing one data constructor and falling through on mismatch.
This is analogous to GHC.Generics' $\mathrm{(:+:)}$ sum
representation~\cite{magalhaes2010-generic-deriving}, where each link
corresponds to one summand.
An important consequence is that such chains do \emph{not}
commute, as they have a fixed order, and therefore cannot be extended
through open inheritance, which is the correct semantics for closed
dispatch.
\fi

\section{\bilingual{Conclusion}{结论}}

\ifInheritanceChinese
对不可扩展性的免疫指导了继承演算的设计。第~\ref{sec:semantic-variants}~节排除了每一种替代的组合机制和引用解析,使得带有多目标晚绑定限定 $\this$ 的深合并成为唯一的幸存者。mixin 充当了统一模块、类、对象、方法、字段和局部绑定的单一抽象;mixin 树的深合并使继承具有可交换性、幂等性和可结合性,从而多重继承的线性化问题不再出现。第~\ref{sec:mixin-trees}~节的六个一阶方程,具有指称语义且可由 tabling~\cite{tamaki1986-tabled-resolution} 计算,足以定义语义;观察者通过查询一棵惰性构造的树中的路径来驱动计算。这些方程是对标准面向对象语义的最小偏离:唯一的改变是 $\overrides$ 方程~(\ref{eq:overrides}),它以子树的深合并取代了方法层面的遮蔽。面向对象编程传统上被归类为命令式的;去掉可变状态和函数体,便得到纯面向对象编程,%
\footnote{
  ``纯''具有双重含义:不纯粹意味着不纯洁,即值是不可变的且组合没有副作用;而不混合意味着组合完全依赖继承,没有函数式抽象,例如 $\lambda$ 或 $\beta$-归约。
}
这是声明式编程的一个最小计算模型。

由于 $\overrides$ 合并的是子树而非方法体,函数不再是原语。$\lambda$-演算在第~\ref{sec:translation}~节通过六条翻译规则宏表达到继承演算中。第~\ref{sec:bohm-tree}~节的 Böhm 树对应将这一一阶语义扩展到了 $\lambda$-演算:先前的高阶语义框架被证明忠实于继承演算一阶方程的一个子语言的首次迭代近似。

继承演算进一步获得了 $\lambda$-演算无法表达的三种现象;事实上,正如第~\ref{sec:asymmetry}~节所示,在 Felleisen 意义下~\cite{felleisen1991-expressive-power},继承演算严格地比惰性 $\lambda$-演算更具表达力。第~\ref{sec:expression-problem}~节在表达式问题~\cite{wadler1998-expression-problem}意义下建立了对不可扩展性的免疫。第~\ref{sec:case-study}~节的案例研究表明,普通算术产生了逻辑编程的关系语义:循环继承层次结构通过与递归 Datalog 所支配的相同的升链条件收敛,附录~\ref{app:datalog-encoding}给出了一个系统性的编码。第~\ref{sec:practical-function-color}~节展示了函数颜色盲性~\cite{nystrom2015-function-color}。

单个 mixin 同时是数据、程序、配置、函数、关系和对象。它宏表达了 $\lambda$-演算,并恢复了逻辑编程的关系语义,同时对两者均无法逃脱的不可扩展性保持免疫。

引言提出了配置本身是否能成为实用的通用编程的问题。答案是肯定的。我们提出继承演算作为实用通用编程的基础,我们相信它优于任何已知的计算模型。
\else
Immunity to nonextensibility guided the design of
inheritance-calculus.
Section~\ref{sec:semantic-variants} eliminated every
alternative composition mechanism and reference resolution,
leaving deep merge with multi-target late-binding qualified
$\this$ as the unique survivor.
The mixin served as the single abstraction unifying module,
class, object, method, field, and local binding;
deep merge of mixin trees made inheritance commutative, idempotent, and associative,
so the linearization problem of multiple inheritance did not arise.
The six first-order equations of Section~\ref{sec:mixin-trees},
denotational and computable by
tabling~\cite{tamaki1986-tabled-resolution}, suffice to
define the semantics; the observer drives computation by
querying paths in a lazily constructed tree.
These equations are a minimal departure from standard
object-oriented semantics: the only change is the $\overrides$ equation~(\ref{eq:overrides}),
which replaces method-level shadowing with deep merge of subtrees.
Object-oriented programming is traditionally classified as
imperative; removing mutable state and function bodies yields
pure object-oriented programming,%
\footnote{
  ``Pure'' carries a double meaning:
  not impure, since values are immutable and composition has no
  side effects; and not mixed, since composition relies entirely on
  inheritance, without functional abstractions such
  as~$\lambda$ or~$\beta$-reduction.
}
a minimal computational model for declarative programming.

Since $\overrides$ merges subtrees rather than method bodies,
functions are no longer primitive.
The $\lambda$-calculus macro-expressed into
inheritance-calculus in six translation rules in
Section~\ref{sec:translation}.
The B\"ohm tree correspondence of
Section~\ref{sec:bohm-tree} extended this first-order
semantics to the $\lambda$-calculus:
the prior higher-order semantic frameworks proved faithful
to the first-iteration approximation of a sublanguage of
inheritance-calculus's first-order equations.

Inheritance-calculus further gains three phenomena
that the $\lambda$-calculus cannot express;
indeed, inheritance-calculus is strictly more expressive
than the lazy $\lambda$-calculus in Felleisen's
sense~\cite{felleisen1991-expressive-power},
as shown in Section~\ref{sec:asymmetry}.
Section~\ref{sec:expression-problem} established immunity
to nonextensibility in the sense of the
Expression Problem~\cite{wadler1998-expression-problem}.
The case study of Section~\ref{sec:case-study} showed that
ordinary arithmetic yields the relational semantics of logic
programming: cyclic inheritance hierarchies converge by
the same ascending chain condition that governs recursive
Datalog, and Appendix~\ref{app:datalog-encoding} gave a
systematic encoding.
Section~\ref{sec:practical-function-color} exhibited
function color
blindness~\cite{nystrom2015-function-color}.

A single mixin is at once data, program, configuration,
function, relation, and object.
It macro-expresses the $\lambda$-calculus and recovers the
relational semantics of logic programming, while remaining
immune to the nonextensibility that neither escapes.

The introduction asked whether configuration alone can be
practical general-purpose programming.
It can.
We propose inheritance-calculus as a foundation for practical
general-purpose programming, one we believe better than any
known computational model.
\fi

\section{\bilingual{Data Availability Statement}{数据可用性声明}}

\ifInheritanceChinese
可执行实现 (MIXINv2) 以及本文中的所有示例以 MIT 许可证作为开源软件公开发布\anon[{} (匿名审查期间省略 URL)]{~\cite{mixin2025}}。每个以继承演算记法呈现的示例都有对应的可执行 MIXINv2 源码,且一套测试套件复现了论文中的每个示例。案例研究 (第~\ref{sec:case-study}~节) 的布尔逻辑和一元自然数算术、附录~\ref{app:datalog-encoding}~的 Datalog 编码 (包括双跳、常量连接和多变量连接关系程序以及重复访问者自连接回归示例),以及函数颜色盲性段落 (第~\ref{sec:practical-function-color}~节) 中讨论的同步和异步运行时实例,均包含在这些源码中。第~\ref{sec:practical-union-fs}~节所讨论的 ratarmount 补丁\anon[也已公开发布 (匿名审查期间省略 URL)]{作为一个开放的拉取请求~\cite{ratarmount2025}可以获取}。\anon[以上所有内容也作为辅助材料一并附上。]{}
\else
The executable implementation (MIXINv2) and all examples from
this paper are openly available as open-source software under
the MIT license\anon[{} (URL omitted for anonymous review)]{~\cite{mixin2025}}.
Every example rendered in inheritance-calculus notation has a
corresponding executable MIXINv2 source, and a test suite
reproduces every example in the paper.
The boolean logic and unary natural number arithmetic of the case
study (Section~\ref{sec:case-study}), the Datalog encodings of
Appendix~\ref{app:datalog-encoding} (including the two-hop,
constant-join, and multi-variable-join relational programs and the
repeated-visitor self-join regression example), and the synchronous
and asynchronous runtime instances discussed in the
function-color-blindness paragraph
(Section~\ref{sec:practical-function-color}) are all included among
these sources.
The ratarmount patch discussed in
Section~\ref{sec:practical-union-fs} is
\anon[also openly available (URL omitted for anonymous review)]{available as an open pull
request~\cite{ratarmount2025}}.
\anon[All of the above are also included as supplementary material.]{}
\fi

\begin{acks}
  Claude Code, GitHub Copilot, Gemini Code Assist, and Codex were utilized to generate sections of this Work,
  including text, tables, graphs, code, formulas, citations, etc.
\end{acks}
\fi % end \ifInheritanceBody (the main matter)

% References follow the main matter and precede the appendix. This copy fires
% whenever the main matter is present (preprint, submission); the appendix-only
% build emits its copy after the appendix instead.
\ifInheritanceBody
\bibliographystyle{ACM-Reference-Format}
\bibliography{references}
\fi

\ifInheritanceAppendix
\appendix

\section{\bilingual{Well-Definedness of the Semantic Functions}{语义函数的良定义性}}
\label{app:well-definedness}

\ifInheritanceChinese
六个语义方程~(\ref{eq:properties})--(\ref{eq:this}) 是一阶的:其输入与输出均为路径及路径的集合,不涉及函数空间、闭包或高阶构造。
本节通过证明这六个方程定义了一个单调算子来确立良定义性;该算子的最小不动点由 tabling 求值~\cite{tamaki1986-tabled-resolution} 可靠地计算。
\else
The six semantic
equations~(\ref{eq:properties})--(\ref{eq:this}) are
first-order: their inputs and outputs are paths and sets of
paths, with no function spaces, closures, or higher-order
constructs.
This section establishes well-definedness by showing that
the six equations define a monotone operator whose least
fixed point is computed soundly by tabled
evaluation~\cite{tamaki1986-tabled-resolution}.
\fi

\begin{definition}[\bilingual{Dependency graph}{依赖图}]\label{def:dependency-graph}
\ifInheritanceChinese
  固定一棵 AST 及其原始函数 $\defines$ 与 $\inherits$。
  \emph{依赖图}~$G$ 以由该 AST 产生的所有语义函数应用
  ($\properties(p)$、\allowbreak$\supers(p)$、\allowbreak$\overrides(p)$、\allowbreak$\bases(p)$、\allowbreak
  $\resolve(\ldots)$、\allowbreak$\this(\ldots)$) 为顶点。
  当通过方程~(\ref{eq:properties})--(\ref{eq:this}) 计算~$u$
  需要~$v$ 的值时,从顶点~$u$ 到顶点~$v$ 有一条有向边。
\else
  Fix an AST with its primitive functions $\defines$ and
  $\inherits$.
  The \emph{dependency graph}~$G$ has as vertices all
  semantic-function applications
  ($\properties(p)$, $\supers(p)$, $\overrides(p)$,
    $\bases(p)$, $\resolve(\ldots)$,
  $\this(\ldots)$) that arise from the AST\@.
  There is a directed edge from vertex~$u$ to vertex~$v$
  whenever computing~$u$ via
  equations~(\ref{eq:properties})--(\ref{eq:this})
  requires the value of~$v$.
\fi
\end{definition}

\begin{definition}[\bilingual{Immediate consequence operator $T_P$}{直接后果算子 $T_P$}]%
  \label{def:tp-operator}
\ifInheritanceChinese
  设 $\mathcal{Q}$ 为所有语义函数查询(即~$G$ 的顶点)的集合。
  对每个查询 $q \in \mathcal{Q}$,设 $\mathcal{V}_q$ 为形状与~$q$ 相匹配的答案值的集合
  (依定义~$q$ 的方程,为路径、路径的集合或路径对的集合)。
  \emph{查询--答案映射}是逐点乘积
  $\prod_{q \in \mathcal{Q}} \mathcal{P}(\mathcal{V}_q)$ 中的一个元素;
  等价地,它为每个查询~$q$ 指派一个候选答案集合 $M(q) \subseteq \mathcal{V}_q$。
  \emph{直接后果算子}~\cite{vanemden1976-predicate-logic-semantics}
  $T_P$ 将查询--答案映射~$M$ 映射为新的查询--答案映射~$T_P(M)$,方式是:对每个查询~$q$,
  求值定义~$q$ 的那条方程的右侧,并以当前近似~$M(q')$ 解释每个子查询~$q'$。
\else
  Let $\mathcal{Q}$ be the set of all
  semantic-function queries
  (vertices of~$G$).
  For each query $q \in \mathcal{Q}$, let
  $\mathcal{V}_q$ be the set of answer values of the
  appropriate shape for~$q$
  (paths, sets of paths, or sets of pairs of paths,
  depending on the equation defining~$q$).
  A \emph{query--answer map} is an element of the
  pointwise product
  $\prod_{q \in \mathcal{Q}} \mathcal{P}(\mathcal{V}_q)$;
  equivalently, it assigns to each query~$q$ a set
  $M(q) \subseteq \mathcal{V}_q$ of candidate answers.
  The \emph{immediate consequence
  operator}~\cite{vanemden1976-predicate-logic-semantics}
  $T_P$ maps a query--answer map~$M$ to a new
  query--answer map~$T_P(M)$ by evaluating, for each
  query~$q$, the right-hand side of the equation that
  defines~$q$, interpreting every sub-query~$q'$ by
  the current approximation~$M(q')$.
\fi
\end{definition}

\begin{lemma}[\bilingual{Monotonicity}{单调性}]\label{lem:monotonicity}
  \ifInheritanceChinese
  算子 $T_P$ 在完备格
  $\Bigl(\prod_{q \in \mathcal{Q}} \mathcal{P}(\mathcal{V}_q),\; \subseteq\Bigr)$
  上是单调的。
  \else
  The operator $T_P$ is monotone on the complete lattice
  $\Bigl(\prod_{q \in \mathcal{Q}} \mathcal{P}(\mathcal{V}_q),\; \subseteq\Bigr)$.
  \fi
\end{lemma}

\begin{proof}
\ifInheritanceChinese
  每条方程的右侧均由 $\cup$、$\in$、等式测试以及正集合
  概括 $\{\, x \mid x \in S, \ldots \,\}$ 构成。
  六条方程均不使用否定、集合差或补集。
  因此,若 $M \subseteq M'$(逐点),则 $T_P(M) \subseteq T_P(M')$。
\else
  Each equation's right-hand side is built from
  $\cup$, $\in$, equality tests, and positive set
  comprehension $\{\, x \mid x \in S, \ldots \,\}$.
  None of the six equations uses negation, set
  difference, or complement.
  Therefore, if $M \subseteq M'$ (pointwise), then
  $T_P(M) \subseteq T_P(M')$.
\fi
\end{proof}

\begin{theorem}[\bilingual{Existence and uniqueness of
  $\mathrm{lfp}(T_P)$}{$\mathrm{lfp}(T_P)$ 的存在性与唯一性}]\label{thm:lfp-exists}
  \ifInheritanceChinese
  算子 $T_P$ 有唯一的最小不动点 $\mathrm{lfp}(T_P)$。
  \else
  The operator $T_P$ has a unique least fixed point
  $\mathrm{lfp}(T_P)$.
  \fi
\end{theorem}

\begin{proof}
\ifInheritanceChinese
  每个幂集格
  $(\mathcal{P}(\mathcal{V}_q), \subseteq)$
  都是完备的,
  因此逐点乘积格
  $\Bigl(\prod_{q \in \mathcal{Q}} \mathcal{P}(\mathcal{V}_q), \subseteq\Bigr)$
  也是完备的。
  由 Knaster--Tarski 定理~\cite{tarski1955-lattice-fixpoint-theorem},
  完备格上的每个单调函数都有唯一的最小不动点。
\else
  Each powerset lattice
  $(\mathcal{P}(\mathcal{V}_q), \subseteq)$
  is complete,
  so the pointwise product lattice
  $\Bigl(\prod_{q \in \mathcal{Q}} \mathcal{P}(\mathcal{V}_q), \subseteq\Bigr)$
  is also complete.
  By the Knaster--Tarski
  theorem~\cite{tarski1955-lattice-fixpoint-theorem},
  every monotone function on a complete lattice has a
  unique least fixed point.
\fi
\end{proof}

\begin{theorem}[\bilingual{Well-definedness}{良定义性}]\label{thm:well-defined}
  \ifInheritanceChinese
  语义函数 $\properties$、\allowbreak$\supers$、\allowbreak$\overrides$、\allowbreak$\bases$、\allowbreak
  $\resolve$、\allowbreak$\this$ 是良定义的:
  其指称语义为 $\mathrm{lfp}(T_P)$,即直接后果算子的唯一最小不动点。
  \else
  The semantic functions
  $\properties$, $\supers$, $\overrides$, $\bases$,
  $\resolve$, $\this$ are well-defined:
  their denotational semantics is
  $\mathrm{lfp}(T_P)$, the unique least fixed point of
  the immediate consequence operator.
  \fi
\end{theorem}

\begin{theorem}[\bilingual{Tabling soundness}{Tabling 可靠性}]%
  \label{thm:tabling-soundness}
  \ifInheritanceChinese
  即使依赖图~$G$ 是无穷的,方程~(\ref{eq:properties})--(\ref{eq:this}) 的
  tabling 求值也绝不会返回错误答案。具体而言:
  \begin{enumerate}
    \item \textbf{可靠性(无假正例)。}
      tabling 返回的每个事实均由某条方程的右侧在其输入上求值产生。
      对求值轨迹进行归纳,每个返回的事实都属于 $\mathrm{lfp}(T_P)$。

    \item \textbf{单调积累。}
      tabling 累加器在各轮迭代之间只通过集合并增长,任何事实都不会被删除。

    \item \textbf{不动点检测的正确性。}
      若求值过程中没有子查询被重入,则所有子查询均在无近似的情况下被完整计算,结果是精确的。
      若发生了重入但累加器是稳定的(即没有新事实被添加),则累加器等于
      $\mathrm{lfp}(T_P)$ 在可达查询上的限制,因为它是一个从~$\bot$ 上升到达的前不动点。

    \item \textbf{可靠的失败。}
      当 tabling 无法确定答案时(无论是因为迭代预算耗尽还是因为存在无穷非循环链),
      它会抛出错误而非返回错误答案。
  \end{enumerate}
  \else
  Tabled evaluation of
  equations~(\ref{eq:properties})--(\ref{eq:this})
  never returns a wrong answer, even when the
  dependency graph~$G$ is infinite.
  Specifically:
  \begin{enumerate}
    \item \textbf{Soundness (no false positives).}
      Every fact returned by tabling was produced by
      evaluating some equation's right-hand side on its
      inputs.
      By induction on the evaluation trace, every
      returned fact belongs to $\mathrm{lfp}(T_P)$.

    \item \textbf{Monotone accumulation.}
      The tabling accumulator only grows between
      iterations via set union; no fact is ever
      removed.

    \item \textbf{Correct fixpoint detection.}
      If no sub-query was re-entered during evaluation,
      all sub-queries were fully computed with no
      approximation, and the result is exact.
      If re-entry occurred but the accumulator is
      stable, meaning no new facts were added, the
      accumulator equals
      $\mathrm{lfp}(T_P)$ restricted to the reachable
      queries, because it is a pre-fixed point reached
      by ascending from~$\bot$.

    \item \textbf{Sound failure.}
      When tabling cannot determine the answer,
      whether because the iteration budget is exhausted
      or because of an infinite non-cyclic chain,
      it raises an error rather than
      returning a wrong answer.
  \end{enumerate}
  \fi
\end{theorem}

\begin{proof}
\ifInheritanceChinese
  对于~(1),每个返回的事实都是某条方程的右侧作用于先前返回的事实上的输出;
  归纳地,所有输入均属于 $\mathrm{lfp}(T_P)$,因此输出也是(由单调性和不动点性质)。

  对于~(2),实现通过 $A \leftarrow A \cup T_P(A)$ 进行积累,这对~$A$ 是单调的。

  对于~(3),当没有重入时,求值是有限无环可达子图上的普通记忆化递归,
  因此每个子查询都被精确计算,不涉及任何近似。
  当发生重入时,设 $R \subseteq \mathcal{Q}$ 为从初始查询可达的查询集合,
  设 $T_P^R$ 表示 $T_P$ 到~$R$ 上映射的限制。
  若累加器稳定,即 $A = A \cup T_P^R(A)$,则
  $T_P^R(A) \subseteq A$,从而~$A$ 是~$T_P^R$ 的一个前不动点。

  由~(1),曾被添加到累加器中的每个事实均由某次右侧求值产生,
  因此属于 $\mathrm{lfp}(T_P^R)$。
  故 $A \subseteq \mathrm{lfp}(T_P^R)$。

  反之,由 Knaster--Tarski 将最小不动点刻画为最小前不动点,
  $\mathrm{lfp}(T_P^R)$ 包含于~$T_P^R$ 的每个前不动点中,特别地包含于~$A$ 中。
  从而 $\mathrm{lfp}(T_P^R) \subseteq A$。

  因此 $A = \mathrm{lfp}(T_P^R)$,即稳定的累加器等于可达查询上的精确指称。

  对于~(4),预算耗尽与栈溢出均产生异常,而非静默的错误答案。
\else
  For~(1), each returned fact is the output of some
  equation's right-hand side applied to previously
  returned facts; by induction, all inputs are in
  $\mathrm{lfp}(T_P)$, so the output is too (by
  monotonicity and the fixed-point property).

  For~(2), the implementation accumulates via
  $A \leftarrow A \cup T_P(A)$, which is monotone in~$A$.

  For~(3), when no reentry occurs, the evaluation is
  ordinary memoised recursion on a finite acyclic
  reachable subgraph, so every sub-query is computed
  exactly and no approximation is involved.
  When reentry occurs, let $R \subseteq \mathcal{Q}$ be
  the set of queries reachable from the initial query,
  and let $T_P^R$ denote the restriction of $T_P$ to
  maps on~$R$.
  If the accumulator is stable, that is,
  $A = A \cup T_P^R(A)$, then
  $T_P^R(A) \subseteq A$, so $A$ is a pre-fixed point
  of~$T_P^R$.

  By~(1), every fact ever added to the accumulator was
  produced by some right-hand side evaluation and
  therefore belongs to $\mathrm{lfp}(T_P^R)$.
  Hence $A \subseteq \mathrm{lfp}(T_P^R)$.

  Conversely, by the Knaster--Tarski characterization
  of the least fixed point as the least pre-fixed
  point, $\mathrm{lfp}(T_P^R)$ is contained in every
  pre-fixed point of~$T_P^R$, in particular in~$A$.
  Thus $\mathrm{lfp}(T_P^R) \subseteq A$.

  Therefore $A = \mathrm{lfp}(T_P^R)$, that is, the stable
  accumulator equals the exact denotation on the
  reachable queries.

  For~(4), budget exhaustion and stack overflow both
  produce exceptions, never silent wrong answers.
\fi
\end{proof}

% 终止性依赖程序结构:无环有限子图→一趟即止;有环有限逐查询值域→有限格上升链条件收敛;
% 无穷可达子图→固有发散,属图灵完备性,非求值策略的缺陷。
\paragraph{\bilingual{Termination}{终止性}}
\ifInheritanceChinese
求值是否终止取决于程序:
\begin{itemize}
  \item \textbf{无环、有限可达子图。}
    一趟即足;求值终止。
  \item \textbf{有环、逐查询值域有限。}
    多趟迭代在有限格的上升链条件下收敛~\cite{bancilhon1986-recursive-query-strategies}。
    这是 Datalog 的情形(第~\ref{sec:recursive-rules} 节)。
  \item \textbf{无穷可达子图。}
    求值可能发散;这是图灵完备性的固有属性,而非求值策略的缺陷。
\end{itemize}
\else
Whether evaluation terminates depends on the
program:
\begin{itemize}
  \item \textbf{Acyclic, finite reachable subgraph.}
    One pass suffices; evaluation terminates.
  \item \textbf{Cyclic, finite per-query domain.}
    Multiple passes converge by the ascending chain
    condition on a finite
    lattice~\cite{bancilhon1986-recursive-query-strategies}.
    This is the Datalog case
    (Section~\ref{sec:recursive-rules}).
  \item \textbf{Infinite reachable subgraph.}
    Evaluation may diverge; this is inherent to
    Turing completeness and not a defect of the
    evaluation strategy.
\end{itemize}
\fi

\section{B\"ohm Tree Correspondence: Proofs}
\label{app:bohm-tree-proofs}

This appendix proves two results about the sublanguage of
inheritance-calculus corresponding to the $\lambda$-calculus
(Section~\ref{sec:bohm-tree}).
First, the first-iteration approximation
$T_P\!\uparrow\!1$ is fully abstract with respect to
B\"ohm tree equivalence
(Theorems~\ref{thm:adequacy} and~\ref{thm:first-order-semantics}),
establishing that the first iteration coincides with
the lazy $\lambda$-calculus.
Second, the full $\mathrm{lfp}(T_P)$ semantics is strictly
more defined in the powerset sense of
Figure~\ref{fig:calculi-matrix}: translated $\lambda$-terms
that diverge under $\beta$-reduction may converge to definite
sets, as shown in Section~\ref{sec:bohm-tree}.

\subsection{B\"ohm Trees}

We briefly recall the definition of B\"ohm
trees~\cite{barendregt1984-lambda-calculus}.
A \emph{head reduction} $M \to_h M'$ contracts
the outermost $\beta$-redex only: if $M$ has the form
$(\lambda x.\, e)\; v\; M_1 \cdots M_k$, then
$M \to_h e[v/x]\; M_1 \cdots M_k$.
A term $M$ is in \emph{head normal form} (HNF) if it has no
head redex, that is, $M = \lambda x_1 \ldots x_m.\;
y\; M_1 \cdots M_k$ where $y$ is a variable.

The \emph{B\"ohm tree} $\mathrm{BT}(M)$ of a $\lambda$-term $M$
is an infinite labeled tree defined by:
\[
  \mathrm{BT}(M) =
  \begin{cases}
    \bot & \text{if } M \text{ has no HNF} \\[6pt]
    \lambda x_1 \ldots x_m.\;
    y\bigl(\mathrm{BT}(M_1),\; \ldots,\; \mathrm{BT}(M_k)\bigr)
    &
    \begin{aligned}[t]
      &\text{if } M \to^*_h \\[-2pt]
      &\lambda x_1 \ldots x_m.\; y\; M_1 \cdots M_k
    \end{aligned}
  \end{cases}
\]
B\"ohm tree equivalence $\mathrm{BT}(M) = \mathrm{BT}(N)$ is the
standard observational equivalence of the lazy
$\lambda$-calculus~\cite{abramsky1993-full-abstraction-lazy-lambda,abramsky1995-games-full-abstraction-lazy-lambda}.

\subsection{Path Encoding}

We define a correspondence between positions in the B\"ohm tree
and paths in the mixin tree.
When the translation $\mathcal{T}$ is applied to an ANF
$\lambda$-term, the resulting mixin tree has a specific shape:
\begin{itemize}
  \item An abstraction $\lambda x.\, M$ translates to a record
    with own properties $\{\mathrm{argument}, \mathrm{result}\}$,
    which is referred to as the \emph{abstraction shape}.
  \item A $\mathbf{let}$-binding
    $\mathbf{let}\; x = V_1\; V_2 \;\mathbf{in}\; M$ translates
    to a record with own properties $\{x, \mathrm{result}\}$,
    where $x$ holds the encapsulated application
    $\{\mathcal{T}(V_1),\; \mathrm{argument} \mapsto
    \mathcal{T}(V_2)\}$.
  \item A tail call $V_1\; V_2$ translates to a record with own
    properties $\{\mathrm{tailCall}, \mathrm{result}\}$, where
    $\mathrm{tailCall}$ holds the encapsulated application and
    $\mathrm{result}$ projects
    $\mathrm{tailCall}.\mathrm{result}$.
  \item A variable $x$ with de~Bruijn index $n$ translates to
    an indexed reference $\dbi{n}.\mathrm{argument}$.
\end{itemize}

\noindent
Each B\"ohm tree position is a sequence of navigation steps from
the root.
In the B\"ohm tree, the possible steps at a node
$\lambda x_1 \ldots x_m.\; y\; M_1 \cdots M_k$ are:
entering the body by peeling off one abstraction layer, and
entering the $i$-th argument $M_i$ of the head variable.
In the mixin tree, these correspond to following the
$\mathrm{result}$ label (for entering the body) and the
$\mathrm{argument}$ label after $\this$ resolution (for
entering an argument position).

Rather than comparing absolute paths across different mixin
trees, which would be sensitive to internal structure, we
observe only the \emph{convergence} behavior accessible via
the $\mathrm{result}$ projection
(Definition~\ref{def:convergence}).
Each $\mathrm{result}$-following step corresponds to one head
reduction step in the $\lambda$-calculus: a tail call's
$\mathrm{result}$ projects through the encapsulated
application, and a $\mathbf{let}$-binding's $\mathrm{result}$
enters the continuation.
The abstraction shape at the end signals that a head
normal form has been reached.

\subsection{Single-Path Lemma}

The $\this$ function (equation~\ref{eq:this}) is designed for the
general case of multi-target $\this$ resolution.
For translated $\lambda$-terms, we show it degenerates to
single-target resolution.

\begin{lemma}[Single-Path]\label{lem:single-path}
  For any closed ANF $\lambda$-term $M$, at every invocation of
  $\this(S,\; p_{\mathrm{def}},\; n)$ during evaluation of
  $\mathcal{T}(M)$, the frontier set $S$ contains exactly one
  path.
\end{lemma}

\begin{proof}
  By case analysis on the translation rules of~$\mathcal{T}$.
  Since $\mathcal{T}$ is compositional, the structural invariant
  established for each rule propagates to all inherited subtrees:
  every subtree appearing in the mixin tree is itself the image of
  some sub-term under~$\mathcal{T}$, so the invariant holds
  throughout the runtime composition.

  \paragraph{Case $\mathcal{T}(\lambda x.\, M)
    = \{\mathrm{argument} \mapsto \{\},\;
  \mathrm{result} \mapsto \mathcal{T}(M)\}$}
  This is a record literal with
  $\defines = \{\mathrm{argument}, \mathrm{result}\}$ and
  no inheritance sources, so $\inherits = \varnothing$.
  At the root path $p$, we have $\overrides(p) = \{p\}$ (a single
  path) since there are no inheritance sources to introduce
  additional branches.
  Inside $\mathcal{T}(M)$, any reference to $x$ is
  $\dbi{n}.\mathrm{argument}$ where $n$ is the de~Bruijn index
  pointing to this scope level.
  Since $\mathcal{T}(\lambda x.\, M)$ introduces exactly one
  scope level, there is exactly one
  mixin at that level, hence $\this$ finds exactly one matching
  pair $(p_{\mathrm{site}},\; p_{\mathrm{override}})$ in $\supers$.

  \paragraph{Case
    $\mathcal{T}(\mathbf{let}\; x = V_1\; V_2
    \;\mathbf{in}\; M)
    = \{x \mapsto \{\mathcal{T}(V_1),\;
      \mathrm{argument} \mapsto \mathcal{T}(V_2)\},\;
  \mathrm{result} \mapsto \mathcal{T}(M)\}$}
  The outer record has own properties $\{x, \mathrm{result}\}$
  and no inheritance sources at the outer level
  ($\inherits = \varnothing$), so $\overrides(\text{root})
  = \{\text{root}\}$.
  Inside the $x$ subtree, the inheritance
  $\{\mathcal{T}(V_1),\; \mathrm{argument} \mapsto
  \mathcal{T}(V_2)\}$ has exactly one inheritance source
  ($\mathcal{T}(V_1)$), creating exactly one $\bases$ entry.
  By induction on $V_1$, the single-path property holds inside
  $\mathcal{T}(V_1)$.
  Inside $\mathcal{T}(M)$, the induction hypothesis applies
  directly.

  \paragraph{Case
    $\mathcal{T}(\mathbf{let}\; x = V
    \;\mathbf{in}\; M)
    = \{x \mapsto \{\mathrm{result} \mapsto \mathcal{T}(V)\},\;
  \mathrm{result} \mapsto \mathcal{T}(M)\}$}
  The outer record has own properties $\{x, \mathrm{result}\}$
  and no inheritance sources at the outer level, so
  $\overrides(\text{root}) = \{\text{root}\}$.
  The $x$ subtree is itself a record literal with own
  property $\{\mathrm{result}\}$ and no inheritance
  sources.
  Thus neither entering the bound value nor entering the
  body introduces branching in $\supers$.
  By induction, the single-path property holds inside both
  $\mathcal{T}(V)$ and $\mathcal{T}(M)$.

  \paragraph{Case $\mathcal{T}(V_1\; V_2) =
    \{\mathrm{tailCall} \mapsto
      \{\mathcal{T}(V_1),\; \mathrm{argument} \mapsto
      \mathcal{T}(V_2)\},\;
      \mathrm{result} \mapsto
  \{\mathrm{tailCall}.\mathrm{result}\}\}$}
  Identical to the $\mathbf{let}$-binding case: the outer record
  has own properties $\{\mathrm{tailCall}, \mathrm{result}\}$
  with no inheritance sources at the outer level, and the
  $\mathrm{tailCall}$ subtree contains exactly one inheritance
  source.

  \paragraph{Case $\mathcal{T}(x)
  = \dbi{n}.\mathrm{argument}$}
  This is a reference, not a record.
  The $\this$ function walks up $n$ steps from the enclosing scope.
  By the compositionality argument above, every scope level
  encountered during the walk, including inherited subtrees
  reached through reference resolution, is the image of some
  sub-term under~$\mathcal{T}$, so each step of $\this$
  encounters exactly one matching pair in $\supers$.

  \medskip\noindent
  The key structural invariant is the following.
  For every subtree generated by~$\mathcal{T}$,
  its root has either no inheritance source or exactly one
  inheritance source.
  Moreover, the only roots with an inheritance source are the
  application wrappers introduced by the application
  $\mathbf{let}$-binding and tail-call rules, and in both cases
  that source is unique and encapsulated under a distinguished
  child ($x$ or $\mathrm{tailCall}$), never at the enclosing
  scope level.
  Therefore every scope level traversed by $\this$ contributes
  at most one matching pair in $\supers$, so the frontier set
  remains a singleton throughout evaluation.
\end{proof}

\subsection{Substitution Lemma}

The core mechanism of the translation is that inheritance with
$\mathrm{argument} \mapsto \mathcal{T}(V)$ plays the role of
substitution.  We make this precise.

\begin{lemma}[Substitution]\label{lem:substitution}
  Let $M$ be an ANF term in which $x$ may occur free,
  and let $V$ be a closed value.%
  \footnote{
    The proof uses the Single-Path Lemma
    (Lemma~\ref{lem:single-path}), which requires the inherited tree
    to occur within a closed term. This holds whenever
    $\lambda x.\, M$ applied to $V$ is a subterm of a closed ANF
    program.
  }
  Define the \emph{inherited tree}
  $C = \{\mathcal{T}(\lambda x.\, M),\;
  \mathrm{argument} \mapsto \mathcal{T}(V)\}$.
  Then for every path $p$ under $C \snoc \mathrm{result}$ and
  every label $\ell$:
  \[
    \ell \in \properties(p)
    \text{ in } C \snoc \mathrm{result}
    \quad\Longleftrightarrow\quad
    \ell \in \properties(p)
    \text{ in } \mathcal{T}(M[V/x])
  \]
  where $M[V/x]$ is the usual capture-avoiding substitution.
\end{lemma}

\begin{proof}
  By structural induction on $M$.

  \paragraph{Base case: $M = x$}
  Then $\mathcal{T}(x) =
  \dbi{n}.\mathrm{argument}$ (where $n$ is the de~Bruijn index
  of $x$), and
  $C \snoc \mathrm{result}$ contains this reference.
  In the inherited tree $C$, $\resolve$ and $\this$ resolve the reference by
  navigating from the reference's enclosing scope up $n$ steps
  to the binding $\lambda$, where the inheritance
  $\{\mathcal{T}(\lambda x.\, x),\;
  \mathrm{argument} \mapsto \mathcal{T}(V)\}$ inherits into the
  $\mathrm{argument}$ slot, providing $\mathcal{T}(V)$.
  Since $C$ is the $\mathrm{tailCall}$ subtree of
  $\mathcal{T}\bigl((\lambda x.\, M)\; V\bigr)$, which is a
  closed ANF term, Lemma~\ref{lem:single-path} applies and
  $\this$ finds exactly one path,
  so $\resolve$ returns $\mathcal{T}(V)$'s subtree.
  Since $M[V/x] = V$, we have
  $\mathcal{T}(M[V/x]) = \mathcal{T}(V)$, and the properties
  coincide.

  \paragraph{Base case: $M = y$ where $y \neq x$ and $y$ is
  $\lambda$-bound}
  Then $\mathcal{T}(y) = \dbi{m}.\mathrm{argument}$ where $m$
  is the de~Bruijn index of~$y$.
  This walks up to the scope that binds~$y$, which is different
  from the scope that binds~$x$, so the inheritance with
  $\mathrm{argument} \mapsto \mathcal{T}(V)$ at $x$'s binding
  scope has no effect.
  Since $M[V/x] = y$, the properties are identical.

  \paragraph{Base case: $M = y$ where $y \neq x$ and $y$ is
  let-bound}
  Then $\mathcal{T}(y) = y.\mathrm{result}$.
  This is a lexical reference to the sibling property~$y$ in the
  enclosing record, so the inherited tree~$C$ does not alter its
  resolution.
  Since $M[V/x] = y$, the properties are identical.

  \paragraph{Inductive case: $M = \lambda y.\, M'$}
  Then $\mathcal{T}(M) =
  \{\mathrm{argument} \mapsto \{\},\;
  \mathrm{result} \mapsto \mathcal{T}(M')\}$.
  The subtree at $C \snoc \mathrm{result}$ is another
  abstraction-shaped record.
  References to $x$ inside $\mathcal{T}(M')$ resolve through
  $\this$ in exactly the same way (with one additional scope
  level from the inner $\lambda$), and by induction on $M'$, the
  properties under
  $C \snoc \mathrm{result} \snoc \mathrm{result}$
  coincide with those of
  $\mathcal{T}(M'[V/x])$.
  Since
  $(\lambda y.\, M')[V/x] = \lambda y.\, (M'[V/x])$
  (assuming $y$ is fresh), the result follows.

  \paragraph{Inductive case:
    $M = \mathbf{let}\; z = V_1\; V_2$
  $\mathbf{in}\; M'$}
  Then
  \[
    \mathcal{T}(M) =
    \{z \mapsto \{\mathcal{T}(V_1),\;
      \mathrm{argument} \mapsto \mathcal{T}(V_2)\},\;
    \mathrm{result} \mapsto \mathcal{T}(M')\}.
  \]
  The substitution distributes:
  $M[V/x] =
  \mathbf{let}\; z = V_1[V/x]\; V_2[V/x]
  \;\mathbf{in}\; M'[V/x]$.
  By induction on $V_1$, $V_2$, and $M'$, the properties under
  each subtree ($z$ and $\mathrm{result}$) coincide,
  since references to $x$ inside each subtree resolve to
  $\mathcal{T}(V)$ via the same $\this$ mechanism.

  \paragraph{Inductive case:
    $M = \mathbf{let}\; z = V'
  \;\mathbf{in}\; M'$}
  Then
  \[
    \mathcal{T}(M) =
    \{z \mapsto \{\mathrm{result} \mapsto \mathcal{T}(V')\},\;
    \mathrm{result} \mapsto \mathcal{T}(M')\}.
  \]
  The substitution distributes:
  $M[V/x] =
  \mathbf{let}\; z = V'[V/x]
  \;\mathbf{in}\; M'[V/x]$.
  By induction on $V'$ and $M'$, the properties under both
  subtrees, specifically $z$ and $\mathrm{result}$, coincide.
  In particular, occurrences of~$z$ in $M'$ are translated as the
  lexical reference $z.\mathrm{result}$ on both sides, so the
  wrapper introduced by the value $\mathbf{let}$-binding is
  preserved by the substitution.

  \paragraph{Inductive case: $M = V_1\; V_2$ (tail call)}
  Analogous to the $\mathbf{let}$-binding case with
  $\mathrm{tailCall}$ in place of $z$.
\end{proof}

\subsection{Convergence Preservation}

The convergence criterion (Definition~\ref{def:convergence})
follows the $\mathrm{result}$ chain from the root.
Each step in this chain corresponds to one head reduction step
in the $\lambda$-calculus.
We make this precise.

\begin{lemma}[Result-step]\label{lem:result-step}
  Let $(\lambda x.\, M)\; V$ be a tail call, which is a $\beta$-redex.
  Then for every path $p$ and label $\ell$:
  \[
    \ell \in \properties\!\bigl(\,
    \text{root} \snoc \mathrm{result} \snoc p\,\bigr)
    \text{ in } \mathcal{T}\bigl((\lambda x.\, M)\; V\bigr)
    \;\Longleftrightarrow\;
    \ell \in \properties\!\bigl(\,
    \text{root} \snoc p\,\bigr)
    \text{ in } \mathcal{T}(M[V/x])
  \]
  That is, following one $\mathrm{result}$ projection in the
  tail call's mixin tree yields the same properties as the
  root of the reduct's mixin tree.
\end{lemma}

\begin{proof}
  The translation gives:
  \[
    \mathcal{T}\bigl((\lambda x.\, M)\; V\bigr) =
    \{\mathrm{tailCall} \mapsto
      \{\mathcal{T}(\lambda x.\, M),\;
      \mathrm{argument} \mapsto \mathcal{T}(V)\},\;
      \mathrm{result} \mapsto
    \{\mathrm{tailCall}.\mathrm{result}\}\}
  \]
  The $\mathrm{result}$ label at the root is defined as
  $\mathrm{tailCall}.\mathrm{result}$, so
  $\text{root} \snoc \mathrm{result}$ resolves to the
  $\mathrm{result}$ path inside the inheritance
  $\mathrm{tailCall} = \{\mathcal{T}(\lambda x.\, M),\;
  \mathrm{argument} \mapsto \mathcal{T}(V)\}$.
  This inheritance is exactly the inherited tree $C$ of
  Lemma~\ref{lem:substitution}, and
  $\mathrm{tailCall} \snoc \mathrm{result}$ is
  $C \snoc \mathrm{result}$.
  By Lemma~\ref{lem:substitution}, the properties under
  $C \snoc \mathrm{result}$ coincide with those of
  $\mathcal{T}(M[V/x])$.

  \medskip\noindent
  \emph{Let-binding variant.}
  An analogous result holds for
  $\mathbf{let}\; x = V_1\; V_2 \;\mathbf{in}\; M'$
  where $V_1 = \lambda y.\, M''$.
  The translation gives:
  \[
    \mathcal{T}\bigl(\mathbf{let}\; x = V_1\; V_2
    \;\mathbf{in}\; M'\bigr) =
    \{x \mapsto \{\mathcal{T}(V_1),\;
      \mathrm{argument} \mapsto \mathcal{T}(V_2)\},\;
    \mathrm{result} \mapsto \mathcal{T}(M')\}
  \]
  Since $\mathrm{result} \mapsto \mathcal{T}(M')$ is a property
  definition, the subtree at
  $\text{root} \snoc \mathrm{result}$ is $\mathcal{T}(M')$.
  Inside $\mathcal{T}(M')$, each reference to $x$ is
  $x.\mathrm{result}$, which resolves to the $\mathrm{result}$
  path inside the $x$ subtree
  $\{\mathcal{T}(\lambda y.\, M''),\;
  \mathrm{argument} \mapsto \mathcal{T}(V_2)\}$.
  This is the inherited tree $C$ of
  Lemma~\ref{lem:substitution}, so
  $x.\mathrm{result}$ has the same properties as
  $\mathcal{T}(M''[V_2/y])$.
  In the $\lambda$-calculus,
  $M = (\lambda x.\, M')(V_1\; V_2)$, and the head reduct is
  $N = M'[(V_1\; V_2)/x]$.
  In $\mathcal{T}(N)$, each occurrence of $x$ is replaced by
  $(V_1\; V_2)$, whose result, by the tail-call variant above,
  also yields $\mathcal{T}(M''[V_2/y])$.
  Therefore the properties at
  $\text{root} \snoc \mathrm{result} \snoc p$ in
  $\mathcal{T}(M)$ coincide with those at
  $\text{root} \snoc p$ in $\mathcal{T}(N)$ for all $p$.

  \medskip\noindent
  \emph{Value let-binding variant.}
  An analogous result holds for
  $\mathbf{let}\; x = V \;\mathbf{in}\; M'$.
  The translation gives:
  \[
    \mathcal{T}\bigl(\mathbf{let}\; x = V \;\mathbf{in}\; M'\bigr) =
    \{x \mapsto \{\mathrm{result} \mapsto \mathcal{T}(V)\},\;
    \mathrm{result} \mapsto \mathcal{T}(M')\}.
  \]
  Since $\mathrm{result} \mapsto \mathcal{T}(M')$ is a property
  definition, the subtree at
  $\text{root} \snoc \mathrm{result}$ is $\mathcal{T}(M')$.
  Inside $\mathcal{T}(M')$, each occurrence of $x$ is translated
  as the lexical reference $x.\mathrm{result}$, which resolves to
  the subtree $\mathcal{T}(V)$.
  This is exactly the translation of the reduct
  $M'[V/x]$.
  Therefore the properties at
  $\text{root} \snoc \mathrm{result} \snoc p$ in
  $\mathcal{T}(\mathbf{let}\; x = V \;\mathbf{in}\; M')$
  coincide with those at
  $\text{root} \snoc p$ in $\mathcal{T}(M'[V/x])$ for all $p$.
\end{proof}

\begin{theorem}[Convergence preservation]%
  \label{thm:convergence-preservation}
  If $M \to_h N$ (head reduction), then
  $\mathcal{T}(M){\Downarrow}$ if and only if
  $\mathcal{T}(N){\Downarrow}$.
\end{theorem}

\begin{proof}
  A head reduction step in ANF takes one of three forms:
  a tail call $(\lambda x.\, M')\; V \to_h M'[V/x]$,
  an application $\mathbf{let}$-binding
  $\mathbf{let}\; x = V_1\; V_2 \;\mathbf{in}\; M'
  \to_h M'[(V_1\; V_2)/x]$,
  or a value $\mathbf{let}$-binding
  $\mathbf{let}\; x = V \;\mathbf{in}\; M'
  \to_h M'[V/x]$.
  The second and third forms are simply the ANF
  presentation of one $\beta$-step in the underlying
  $\lambda$-term.
  In all three cases, Lemma~\ref{lem:result-step}
  (tail-call, application let-binding, and value
  let-binding variants) shows that
  the properties at $\text{root} \snoc \mathrm{result}^n$ in
  $\mathcal{T}(M)$ coincide with the properties at
  $\text{root} \snoc \mathrm{result}^{n-1}$ in
  $\mathcal{T}(N)$ for all $n \ge 1$, where $N$ is the
  head reduct.
  Hence the abstraction shape appears at depth $n$ in the
  pre-reduct if and only if it appears at depth $n - 1$ in
  the post-reduct.
\end{proof}

\begin{remark}[Iteration levels]
  \label{rem:convergence-lfp}
  Section~\ref{sec:bohm-tree} showed that the B\"ohm tree
  correspondence holds at the first iteration
  $T_P\!\uparrow\!1$ and that $\mathrm{lfp}(T_P)$ is
  strictly more defined in the powerset sense of
  Figure~\ref{fig:calculi-matrix}.
  Here we make the connection precise.
  Tabled evaluation~\cite{tamaki1986-tabled-resolution}
  computes $\mathrm{lfp}(T_P)$ by iterating:
  $T_P\!\uparrow\!0 = \bot$,
  $T_P\!\uparrow\!1 = T_P(\bot)$,
  $T_P\!\uparrow\!2 = T_P(T_P(\bot))$,
  \ldots\ until stabilisation
  (Theorem~\ref{thm:tabling-soundness}).

  \begin{description}
    \item[$T_P\!\uparrow\!1$ corresponds to
      $\beta$-reduction.]
      $T_P\!\uparrow\!1$ evaluates each equation once,
      with re-entrant calls returning~$\bot$
      (the value from $T_P\!\uparrow\!0$).
      The Result-Step Lemma
      (Lemma~\ref{lem:result-step}) establishes a
      structural isomorphism between the equation
      systems of $\mathcal{T}(M)$ and
      $\mathcal{T}(N)$, shifted by one
      $\mathrm{result}$ step.
      This isomorphism is purely structural and holds
      for any number of iterations.
      Theorems~\ref{thm:adequacy}
      and~\ref{thm:first-order-semantics}
      are proved at this level.
  \end{description}

  \paragraph{Observation granularity}%
  \label{rem:set-vs-domain}
  The set-valued semantics at $T_P\!\uparrow\!1$ is
  strictly finer than the domain-theoretic semantics of
  the lazy $\lambda$-calculus.
  In the latter, a non-terminating computation yields a
  single undifferentiated~$\bot$;
  in the former, one can in principle distinguish
  $\varnothing$ (a definite empty set, arrived at in
  finite time) from non-stabilisation of the Kleene
  iteration (no finite approximation suffices).
  The convergence criterion
  (Definition~\ref{def:convergence}) deliberately
  projects away this distinction: it observes only
  whether the $\mathrm{result}$ chain reaches an abstraction
  shape, reducing the observation to a single bit.
  Theorems~\ref{thm:adequacy}
  and~\ref{thm:first-order-semantics}
  are stated at this coarsened observation level.

  \begin{description}
    \item[$\mathrm{lfp}(T_P)$ restores full precision.]
      At $\mathrm{lfp}(T_P)$, the full precision of the
      set-valued semantics is restored: re-entrant calls
      that returned~$\varnothing$ at $T_P\!\uparrow\!1$
      may now converge to non-empty sets, so the semantics
      distinguishes ``definitionally empty'' from
      ``convergent after further iteration,''
      and $M{\Downarrow}_\lambda$ implies
      $\mathcal{T}(M){\Downarrow}_{\mathrm{lfp}}$,
      but not conversely.
  \end{description}
\end{remark}

\subsection{Adequacy}

\begin{definition}[Inheritance-convergence]\label{def:convergence}
  A mixin tree $T$ \emph{converges}, written $T{\Downarrow}$,
  if there exists $n \ge 0$ such that:
  \[
    \{\mathrm{argument},\, \mathrm{result}\}
    \;\subseteq\;
    \properties\!\bigl(\,
      \underbrace{\text{root} \snoc \mathrm{result}
      \snoc \cdots \snoc \mathrm{result}}_{n}
    \,\bigr)
  \]
  in the semantics of
  Section~\ref{sec:mixin-trees}.
\end{definition}

\begin{theorem}[Adequacy]\label{thm:adequacy}
  A closed ANF $\lambda$-term $M$ has a head normal form if and
  only if $\mathcal{T}(M){\Downarrow}$.
\end{theorem}

A closed ANF term at the top level is either an abstraction
(a value) or a computation (an application
  $\mathbf{let}$-binding, a value $\mathbf{let}$-binding,
or a tail call).
An abstraction translates to a record whose root immediately
has the abstraction shape; a computation's root has the form
$\{x, \mathrm{result}\}$ or
$\{\mathrm{tailCall}, \mathrm{result}\}$, and one must follow
the $\mathrm{result}$ chain to find the eventual value.

\begin{proof}[Proof of Theorem~\ref{thm:adequacy}]
  ($\Rightarrow$)
  Suppose $M \to^*_h \lambda x.\, M'$ in $k$ head reduction
  steps.
  We show $\mathcal{T}(M)$ converges at depth $\le k$ by
  induction on $k$.

  \paragraph{Base case ($k = 0$)}
  $M$ is already an abstraction
  $\lambda x.\, M'$.
  Then $\mathcal{T}(M) =
  \{\mathrm{argument} \mapsto \{\},\;
  \mathrm{result} \mapsto \mathcal{T}(M')\}$,
  whose root has
  $\defines = \{\mathrm{argument}, \mathrm{result}\}$.
  These labels are directly in $\defines$ of the root, so
  the recursive evaluation returns them without further
  recursion, and $\mathcal{T}(M)$ converges at depth $n = 0$.

  \paragraph{Inductive step ($k \ge 1$)}
  $M$ is not an abstraction,
  so $M$ is either a tail call, an application
  $\mathbf{let}$-binding, or a value
  $\mathbf{let}$-binding.
  Each form corresponds to one head $\beta$-step in the
  underlying $\lambda$-term.
  Let $N$ be the head reduct, so
  $M \to_h N \to^{k-1}_h \lambda x.\, M'$.
  Note that $N$ is closed: head reduction substitutes a
  closed value into a term whose only free variable is the
  substitution target, preserving closedness.
  By Lemma~\ref{lem:result-step} (tail-call,
    application let-binding, or value let-binding
  variant), the properties at
  $\text{root} \snoc \mathrm{result} \snoc p$
  in $\mathcal{T}(M)$ coincide with those at
  $\text{root} \snoc p$ in $\mathcal{T}(N)$.
  By the induction hypothesis, $\mathcal{T}(N)$
  converges at depth $\le k - 1$, so $\mathcal{T}(M)$
  converges at depth $\le k$.

  \medskip\noindent
  ($\Leftarrow$)
  Suppose $\mathcal{T}(M){\Downarrow}$ at depth $n$.
  We show $M$ has a head normal form by induction on $n$.

  \paragraph{Base case ($n = 0$)}
  The root of $\mathcal{T}(M)$ has
  abstraction shape
  $\{\mathrm{argument}, \mathrm{result}\}$.
  Only the abstraction rule $\mathcal{T}(\lambda x.\, M')$
  produces a root with own property $\mathrm{argument}$.
  (For tail calls and $\mathbf{let}$-bindings, the root has
    $\defines = \{\mathrm{tailCall}, \mathrm{result}\}$
    or $\{x, \mathrm{result}\}$, and $\mathrm{argument}$ cannot
    appear via $\supers$ because
    $\inherits = \varnothing$ at the root, giving
  $\bases(\text{root}) = \varnothing$.)
  Therefore $M$ is an abstraction, hence already in head normal
  form.

  \paragraph{Inductive step ($n \ge 1$)}
  The abstraction shape
  appears at depth $n$ but not at depth $0$.
  By the base case argument, $M$ is not an abstraction, so at the
  head position of this closed top-level ANF term, $M$ is either
  a tail call, an application $\mathbf{let}$-binding, or a value
  $\mathbf{let}$-binding.
  In each case, the ANF grammar requires the function
  position to be a value~$V_1$
  ($V ::= x \mid \lambda x.\, M$).
  Since $M$ is closed, $V_1$ cannot be a free variable,
  so $V_1$ must be a $\lambda$-abstraction;
  the head form is therefore a $\beta$-redex.
  By Lemma~\ref{lem:result-step} (tail-call,
    application let-binding, or value let-binding
  variant), the properties at depth
  $n - 1$ in $\mathcal{T}(N)$ match those at depth $n$
  in $\mathcal{T}(M)$, where $N$ is the head reduct.
  So $\mathcal{T}(N)$ converges at depth $n - 1$.
  Since $M$ is closed and head reduction preserves
  closedness (it substitutes a closed value for the only
  free variable), $N$ is a closed ANF term.
  By the induction hypothesis, $N$ has a head normal
  form, and since $M \to_h N$, so does $M$.
\end{proof}

\subsection{First-Order Semantics of the Lazy \texorpdfstring{$\lambda$}{λ}-Calculus}

Adequacy relates a single term's head normal form to its
inheritance-convergence.
The following theorem lifts this to an equivalence between
terms, establishing that the translation $\mathcal{T}$
gives the lazy $\lambda$-calculus a first-order semantics
that preserves and reflects B\"ohm tree equivalence.

\begin{definition}[$\lambda$-contextual equivalence under $\mathcal{T}$]%
  \label{def:ctx-equiv}
  For closed $\lambda$-terms $M$ and $N$, write
  $\mathcal{T}(M) \approx_\lambda \mathcal{T}(N)$ if for every
  closing $\lambda$-calculus context $C[\cdot]$:
  \[
    \mathcal{T}(C[M]){\Downarrow}
    \;\Longleftrightarrow\;
    \mathcal{T}(C[N]){\Downarrow}
  \]
\end{definition}

\noindent
The subscript $\lambda$ emphasizes that the quantification ranges
over $\lambda$-calculus contexts, not over all inheritance-calculus contexts.
This definition observes only convergence\footnote{
  Convergence provides exactly a single bit of information.
}, but
quantifying over all $\lambda$-contexts makes it a fine-grained
equivalence: different contexts can supply different arguments
and project different results, probing every aspect of the
term's behavior within the $\lambda$-definable fragment.

\begin{theorem}[First-order semantics of the lazy $\lambda$-calculus]%
  \label{thm:first-order-semantics}
  For closed $\lambda$-terms $M$ and $N$:
  \[
    \mathcal{T}(M) \approx_\lambda \mathcal{T}(N)
    \quad\Longleftrightarrow\quad
    \mathrm{BT}(M) = \mathrm{BT}(N)
  \]
\end{theorem}

\begin{proof}[Proof of Theorem~\ref{thm:first-order-semantics}]
  The translation $\mathcal{T}$ is compositional: for every
  $\lambda$-calculus context $C[\cdot]$, the mixin tree
  $\mathcal{T}(C[M])$ is determined by the mixin tree
  $\mathcal{T}(M)$ and the translation of the context.
  This means $\mathcal{T}$ preserves the context structure
  needed for the argument.

  \medskip\noindent
  ($\Leftarrow$)
  Suppose $\mathrm{BT}(M) = \mathrm{BT}(N)$.
  B\"ohm tree equivalence is a congruence~\cite{barendregt1984-lambda-calculus},
  so for every context $C[\cdot]$,
  $\mathrm{BT}(C[M]) = \mathrm{BT}(C[N])$.
  In particular, $C[M]$ has a head normal form iff $C[N]$ does.
  By Theorem~\ref{thm:adequacy},
  $\mathcal{T}(C[M]){\Downarrow}$ iff
  $\mathcal{T}(C[N]){\Downarrow}$.
  Hence $\mathcal{T}(M) \approx_\lambda \mathcal{T}(N)$.

  \medskip\noindent
  ($\Rightarrow$)
  Suppose $\mathrm{BT}(M) \neq \mathrm{BT}(N)$.
  B\"ohm tree equivalence coincides with observational
  equivalence for the lazy
  $\lambda$-calculus~\cite{abramsky1993-full-abstraction-lazy-lambda,abramsky1995-games-full-abstraction-lazy-lambda}:
  there exists a context $C[\cdot]$ such that $C[M]$ converges
  and $C[N]$ diverges, or vice versa.
  By Theorem~\ref{thm:adequacy},
  $\mathcal{T}(C[M]){\Downarrow}$ and
  $\mathcal{T}(C[N]){\Uparrow}$.
  Hence $\mathcal{T}(M) \not\approx_\lambda \mathcal{T}(N)$.
\end{proof}

\noindent
The proof rests on two pillars: Adequacy
(Theorem~\ref{thm:adequacy}), which is our contribution, and
the classical result that B\"ohm tree equivalence equals
observational equivalence for the lazy
$\lambda$-calculus~\cite{abramsky1993-full-abstraction-lazy-lambda,abramsky1995-games-full-abstraction-lazy-lambda}.
Neither pillar requires defining or comparing
absolute paths inside mixin trees.

\section{\bilingual{Expressive Asymmetry: Proofs}{表达力的不对称:证明}}
\label{app:expressiveness}

\begin{definition}[\bilingual{Macro-expressibility, after Felleisen~\cite{felleisen1991-expressive-power}}{宏可表达性,据 Felleisen~\cite{felleisen1991-expressive-power}}]
  \label{def:macro-expressibility}
  \ifInheritanceChinese
  语言 $\mathscr{L}_1$ 的一个构造 $C$ 在语言 $\mathscr{L}_0$ 中是\emph{宏可表达的},若存在一个翻译 $\mathcal{E}$,使得对每一个含有 $C$ 的出现的程序 $P$,翻译 $\mathcal{E}(P)$ 是通过将 $C$ 的每次出现替换为一个仅依赖于 $C$ 的子表达式的 $\mathscr{L}_0$ 表达式而得到的,并保持 $P$ 的所有其他构造不变。
  \else
  A construct $C$ of language $\mathscr{L}_1$ is \emph{macro-expressible}
  in language $\mathscr{L}_0$ if there exists a translation $\mathcal{E}$
  such that for every program $P$ containing occurrences of $C$,
  the translation $\mathcal{E}(P)$ is obtained by replacing each
  occurrence of $C$ with an $\mathscr{L}_0$ expression that depends only
  on $C$'s subexpressions, leaving all other constructs of $P$
  unchanged.
  \fi
\end{definition}

\begin{lemma}[\bilingual{Composition of macro-expressible translations}{宏可表达翻译的复合}]
  \label{lem:macro-composition}
  \ifInheritanceChinese
  若 $\mathcal{E}_1$ 在 $\mathscr{L}_0$ 中宏可表达 $\mathscr{L}_1$ 的构造,且 $\mathcal{E}_2$ 在 $\mathscr{L}_1$ 中宏可表达 $\mathscr{L}_2$ 的构造,则复合 $\mathcal{E}_1 \circ \mathcal{E}_2$ 在 $\mathscr{L}_0$ 中宏可表达 $\mathscr{L}_2$ 的构造。
  \else
  If $\mathcal{E}_1$ macro-expresses the constructs of
  $\mathscr{L}_1$ in $\mathscr{L}_0$ and
  $\mathcal{E}_2$ macro-expresses the constructs of
  $\mathscr{L}_2$ in $\mathscr{L}_1$, then the composition
  $\mathcal{E}_1 \circ \mathcal{E}_2$ macro-expresses the
  constructs of $\mathscr{L}_2$ in $\mathscr{L}_0$.
  \fi
\end{lemma}

\begin{proof}
  \ifInheritanceChinese
  $\mathscr{L}_2$ 的每个构造 $F$ 在 $\mathscr{L}_1$ 上有一个固定的语法抽象 $A_F^{(2)}$。
  出现在 $A_F^{(2)}$ 中的每个 $\mathscr{L}_1$-构造 $G$ 在 $\mathscr{L}_0$ 上有一个固定的语法抽象 $A_G^{(1)}$。
  将 $A_F^{(2)}$ 中的每个 $G$ 替换为 $A_G^{(1)}$,得到 $F$ 在 $\mathscr{L}_0$ 上的一个固定语法抽象,从而满足 E4。
  \else
  Each construct~$F$ of~$\mathscr{L}_2$ has a fixed
  syntactic abstraction~$A_F^{(2)}$
  over~$\mathscr{L}_1$.
  Each $\mathscr{L}_1$-construct~$G$ appearing
  in~$A_F^{(2)}$ has a fixed
  syntactic abstraction~$A_G^{(1)}$
  over~$\mathscr{L}_0$.
  Substituting~$A_G^{(1)}$ for every~$G$
  in~$A_F^{(2)}$ yields a fixed syntactic
  abstraction over~$\mathscr{L}_0$ for~$F$,
  satisfying~E4.
  \fi
\end{proof}

\begin{definition}[\bilingual{Language universe}{语言全集}]
  \label{def:language-universe}
  \ifInheritanceChinese
  设 $\mathscr{U}$ 为一个语言,其构造子是继承演算的构造子、惰性 $\lambda$-演算的构造子(包括 $\lambda$-抽象、应用与变量)以及 $\mathbf{let}$-绑定的并集。
  $\mathscr{U}$ 的语义通过将每个 $\mathscr{U}$-程序翻译为继承演算来给出,分三步:
  \begin{enumerate}
    \item 将每个极大 $\lambda$-演算子表达式(即根节点与所有内部节点仅使用 $\lambda$-演算构造子的每棵子树)转化为 A-范式~\cite{flanagan1993-essence-compiling-continuations}。
    \item 将翻译 $\mathcal{T}$(第~\ref{sec:forward-translation} 节)应用于每个 ANF $\lambda$-演算子表达式。
    \item 按第~\ref{sec:mixin-trees} 节的规则对所得继承演算程序求值。
  \end{enumerate}
  当一个 $\mathscr{U}$-表达式交错使用 $\lambda$-演算与继承演算的构造子(例如,一个 $\mathbf{let}$-绑定,其被绑定的值是一个继承演算表达式)时,翻译从叶到根组合地应用:每个 $\lambda$-演算构造子通过 $\mathcal{T}$ 替换为其继承演算等价物,同时将继承演算子表达式视为不透明的值。
  由于 $\mathcal{T}$ 的每条规则(第~\ref{sec:forward-translation} 节)都是一个固定的语法抽象,仅依赖于被翻译的子表达式,因此这种自底向上的应用是良定义的,并为步骤~(3) 产生一个纯继承演算程序。

  若一个 $\mathscr{U}$-程序 $P$ 的翻译收敛,则 $P$ 收敛,记为 $\mathrm{eval}_{\mathscr{U}}(P)$。
  \else
  Let $\mathscr{U}$ be the language whose constructors are
  the union of the constructors of inheritance-calculus,
  the constructors of the lazy $\lambda$-calculus
  , including $\lambda$-abstraction, application, and variable, and
  $\mathbf{let}$-binding.
  The semantics of $\mathscr{U}$ is given by translating
  every $\mathscr{U}$-program to inheritance-calculus in
  three steps:
  \begin{enumerate}
    \item Convert every maximal $\lambda$-calculus
      subexpression (that is, every subtree whose root
        and all internal nodes use only $\lambda$-calculus
      constructors) to A-normal
      form~\cite{flanagan1993-essence-compiling-continuations}.
    \item Apply the translation $\mathcal{T}$
      (Section~\ref{sec:forward-translation}) to every ANF
      $\lambda$-calculus subexpression.
    \item Evaluate the resulting inheritance-calculus
      program by the rules of
      Section~\ref{sec:mixin-trees}.
  \end{enumerate}
  When a $\mathscr{U}$-expression interleaves
  $\lambda$-calculus and inheritance-calculus constructors
  (e.g.\ a $\mathbf{let}$-binding whose bound value is an
  inheritance-calculus expression), the translation is
  applied compositionally from leaves to root: each
  $\lambda$-calculus construct is replaced by its
  inheritance-calculus equivalent via~$\mathcal{T}$,
  treating inheritance-calculus subexpressions as opaque
  values.
  Since each rule of~$\mathcal{T}$
  (Section~\ref{sec:forward-translation}) is a fixed
  syntactic abstraction depending only on the translated
  subexpressions, this bottom-up application is
  well-defined and yields a pure inheritance-calculus
  program for step~(3).

  A $\mathscr{U}$-program $P$ converges,
  $\mathrm{eval}_{\mathscr{U}}(P)$, iff its translation
  converges.
  \fi
\end{definition}

\begin{remark}\label{rem:universe-properties}
  \ifInheritanceChinese
  惰性 $\lambda$-演算与继承演算是 $\mathscr{U}$ 的保守性限制:各自通过去掉另一方的构造子来得到。
  对于纯继承演算程序,步骤~(1) 与~(2) 是平凡的,因此 $\mathrm{eval}_{\mathscr{U}}$ 与继承演算的语义一致。
  对于纯 $\lambda$-演算程序,步骤~(1)--(3) 的复合由充分性(定理~\ref{thm:adequacy})与惰性求值一致。
  \else
  The lazy $\lambda$-calculus and inheritance-calculus are
  conservative restrictions of $\mathscr{U}$: each is
  obtained by removing the other's constructors.
  For a pure inheritance-calculus program, steps~(1)
  and~(2) are vacuous, so
  $\mathrm{eval}_{\mathscr{U}}$ agrees with the
  inheritance-calculus semantics.
  For a pure $\lambda$-calculus program, the composition
  of steps~(1)--(3) agrees with lazy evaluation by
  Adequacy (Theorem~\ref{thm:adequacy}).
  \fi
\end{remark}

\ifInheritanceChinese
三步翻译是有效的(ANF 转换与 $\mathcal{T}$ 是语法变换),而继承演算的递归求值(第~\ref{sec:mixin-trees} 节)是一个半判定过程,因此 $\mathrm{eval}_{\mathscr{U}}$ 是递归可枚举的,满足 Felleisen 框架~\cite{felleisen1991-expressive-power} 的要求。
\else
The three-step translation is effective (ANF conversion
and~$\mathcal{T}$ are syntactic transformations) and the
recursive evaluation of inheritance-calculus
(Section~\ref{sec:mixin-trees}) is a semi-decision
procedure, so $\mathrm{eval}_{\mathscr{U}}$ is
recursively enumerable, as required by
Felleisen's framework~\cite{felleisen1991-expressive-power}.
\fi

\begin{lemma}[\bilingual{$\mathbf{let}$-binding is macro-expressible
  in the lazy $\lambda$-calculus}{$\mathbf{let}$-绑定在惰性 $\lambda$-演算中是宏可表达的}]
  \label{lem:let-macro}
  \ifInheritanceChinese
  翻译 $\mathbf{let}\; x = e_1 \;\mathbf{in}\; e_2 \;\mapsto\; (\lambda x.\, e_2)\; e_1$ 是 $\mathbf{let}$ 在惰性 $\lambda$-演算中的一个宏可表达的消去。
  \else
  The translation
  $\mathbf{let}\; x = e_1 \;\mathbf{in}\; e_2
  \;\mapsto\; (\lambda x.\, e_2)\; e_1$
  is a macro-expressible elimination of $\mathbf{let}$
  in the lazy $\lambda$-calculus.
  \fi
\end{lemma}

\begin{proof}
  \ifInheritanceChinese
  该翻译将每个 $\mathbf{let}$-绑定替换为一个仅依赖于两个子表达式 $e_1$ 与 $e_2$ 的 $\lambda$-表达式。
  这是 Felleisen~\cite{felleisen1991-expressive-power} 的命题 3.7。
  \else
  The translation replaces each $\mathbf{let}$-binding
  with a $\lambda$-expression that depends only on the
  two subexpressions $e_1$ and~$e_2$.
  This is Proposition~3.7 of
  Felleisen~\cite{felleisen1991-expressive-power}.
  \fi
\end{proof}

\begin{lemma}[\bilingual{Equal expressiveness of the lazy
  $\lambda$-calculus and its ANF fragment}{惰性 $\lambda$-演算与其 ANF 片段的等同表达力}]
  \label{lem:anf-equiv}
  \ifInheritanceChinese
  在 $\mathscr{U}$ 中,惰性 $\lambda$-演算与 ANF 惰性 $\lambda$-演算(第~\ref{sec:forward-translation} 节)具有相同的宏可表达力。
  \else
  In $\mathscr{U}$, the lazy $\lambda$-calculus and the
  ANF lazy $\lambda$-calculus
  (Section~\ref{sec:forward-translation}) are equally
  macro-expressive.
  \fi
\end{lemma}

\begin{proof}
  \ifInheritanceChinese
  每个 ANF 项都是一个 $\lambda$-项,因此 ANF 片段嵌入完整 $\lambda$-演算的映射是恒等变换,平凡地宏可表达。
  对于反方向,ANF 变换~\cite{flanagan1993-essence-compiling-continuations} 由两步构成,每步都是宏可表达的:
  \begin{enumerate}
    \item \emph{\bilingual{Application flattening}{应用展开}。}
      将每个嵌套应用 $(e_1\; e_2)\; e_3$ 替换为 $\mathbf{let}\; x = e_1\; e_2 \;\mathbf{in}\; x\; e_3$。
      其语法抽象为 $A(\cdot_1, \cdot_2, \cdot_3) = \mathbf{let}\; x = \cdot_1\; \cdot_2 \;\mathbf{in}\; x\; \cdot_3$,仅依赖三个子表达式。
    \item \emph{\bilingual{Value lifting}{值提升}。}
      将应用位置上每个自由值 $V$ 替换为 $\mathbf{let}\; x = V \;\mathbf{in}\; x$。
      其语法抽象为 $A(\cdot_1) = \mathbf{let}\; x = \cdot_1 \;\mathbf{in}\; x$。
  \end{enumerate}
  由引理~\ref{lem:let-macro},每个所得 $\mathbf{let}$-绑定本身在 $\lambda$-演算中是宏可表达的。
  由引理~\ref{lem:macro-composition},复合这些宏可表达翻译,得到一个从完整 $\lambda$-演算到其 ANF 片段的宏可表达翻译。
  \else
  Every ANF term is a $\lambda$-term, so the embedding
  of the ANF fragment into the full $\lambda$-calculus is
  the identity, which is trivially macro-expressible.
  For the reverse direction, the ANF
  transformation~\cite{flanagan1993-essence-compiling-continuations} consists of
  two steps, each macro-expressible:
  \begin{enumerate}
    \item \emph{Application flattening.}
      Replace each nested application
      $(e_1\; e_2)\; e_3$ with
      $\mathbf{let}\; x = e_1\; e_2 \;\mathbf{in}\;
      x\; e_3$.
      The syntactic abstraction is
      $A(\cdot_1, \cdot_2, \cdot_3)
      = \mathbf{let}\; x = \cdot_1\; \cdot_2
      \;\mathbf{in}\; x\; \cdot_3$,
      depending only on the three subexpressions.
    \item \emph{Value lifting.}
      Replace each free value~$V$ used in an application
      position with
      $\mathbf{let}\; x = V \;\mathbf{in}\; x$.
      The syntactic abstraction is
      $A(\cdot_1) = \mathbf{let}\; x = \cdot_1
      \;\mathbf{in}\; x$.
  \end{enumerate}
  By Lemma~\ref{lem:let-macro}, each resulting
  $\mathbf{let}$-binding is itself
  macro-expressible in the $\lambda$-calculus.
  By Lemma~\ref{lem:macro-composition}, composing
  these macro-expressible translations
  yields a macro-expressible translation from the full
  $\lambda$-calculus to its ANF fragment.
  \fi
\end{proof}

\begin{theorem}[\bilingual{Forward macro-expressibility}{正向宏可表达性}]
  \label{thm:forward-macro}
  \ifInheritanceChinese
  翻译 $\mathcal{T}$(第~\ref{sec:forward-translation} 节)是 $\lambda$-演算嵌入继承演算的一个宏可表达的嵌入。
  \else
  The translation $\mathcal{T}$
  (Section~\ref{sec:forward-translation}) is a macro-expressible
  embedding of the $\lambda$-calculus into inheritance-calculus.
  \fi
\end{theorem}

\begin{proof}
  \ifInheritanceChinese
  我们验证 Felleisen~\cite[Definition~3.11]{felleisen1991-expressive-power} 的条件 E4:对于每个元数为 $a$ 的 $\lambda$-演算构造子 $F$,必须存在一个固定的 $a$ 元语法抽象 $A_F(\cdot_1,\ldots,\cdot_a)$ 使得 $\mathcal{T}(F(e_1,\ldots,e_a)) = A_F(\mathcal{T}(e_1),\ldots,\mathcal{T}(e_a))$。
  我们逐条检验 $\mathcal{T}$ 的规则(第~\ref{sec:forward-translation} 节):
  \begin{enumerate}
    \item \emph{\bilingual{$\lambda$-bound variable}{$\lambda$-绑定变量}}($x$,de~Bruijn 索引为 $n'$)。
      这是一个零元构造子,因此没有子表达式。
      语法抽象是常量 $A = \dbi{n'}.\mathrm{argument}$。
    \item \emph{\bilingual{Let-bound variable}{Let-绑定变量}}($x$)。
      零元。$A = x.\mathrm{result}$。
    \item \emph{\bilingual{$\lambda$-abstraction}{$\lambda$-抽象}},记为 $\lambda x.\, M$。
      关于体 $M$ 是一元的。
      $A(\cdot_1) = \{\mathrm{argument} \mapsto \{\},\; \mathrm{result} \mapsto \cdot_1\}$。
    \item \emph{\bilingual{Application let-binding}{应用 Let-绑定}}($\mathbf{let}\; x = V_1\; V_2 \;\mathbf{in}\; M$)。
      关于 $V_1$、$V_2$ 与 $M$ 是三元的。
      $A(\cdot_1, \cdot_2, \cdot_3) = \{x \mapsto \{\cdot_1,\; \mathrm{argument} \mapsto \cdot_2\},\; \mathrm{result} \mapsto \cdot_3\}$。
    \item \emph{\bilingual{Value let-binding}{值 Let-绑定}}($\mathbf{let}\; x = V \;\mathbf{in}\; M$)。
      关于 $V$ 与 $M$ 是二元的。
      $A(\cdot_1, \cdot_2) = \{x \mapsto \{\mathrm{result} \mapsto \cdot_1\},\; \mathrm{result} \mapsto \cdot_2\}$。
    \item \emph{\bilingual{Tail call}{尾调用}}($V_1\; V_2$)。
      二元。\newline
      $A(\cdot_1, \cdot_2) = \{\mathrm{tailCall} \mapsto \{\cdot_1,\; \mathrm{argument} \mapsto \cdot_2\},\; \mathrm{result} \mapsto \{\mathrm{tailCall}.\mathrm{result}\}\}$。
  \end{enumerate}
  在每种情况下,语法抽象 $A_F$ 都是一个固定的继承演算上下文,其空位由子表达式的递归翻译填充,从而满足 E4。
  \else
  We verify condition~E4 of
  Felleisen~\cite[Definition~3.11]{felleisen1991-expressive-power}:
  for each $\lambda$-calculus construct~$F$ of arity~$a$,
  there must exist a fixed $a$-ary syntactic abstraction
  $A_F(\cdot_1,\ldots,\cdot_a)$ over
  inheritance-calculus such that
  $\mathcal{T}(F(e_1,\ldots,e_a)) =
  A_F(\mathcal{T}(e_1),\ldots,\mathcal{T}(e_a))$.
  We check each rule of $\mathcal{T}$
  (Section~\ref{sec:forward-translation}):
  \begin{enumerate}
    \item \emph{$\lambda$-bound variable}
      ($x$ with de~Bruijn index $n'$).
      This is a nullary construct and therefore has no subexpressions.
      The syntactic abstraction is the constant
      $A = \dbi{n'}.\mathrm{argument}$.
    \item \emph{Let-bound variable} ($x$).
      Nullary.
      $A = x.\mathrm{result}$.
    \item \emph{$\lambda$-abstraction}, denoted $\lambda x.\, M$.
      Unary in the body~$M$.
      $A(\cdot_1) = \{\mathrm{argument} \mapsto \{\},\;
      \mathrm{result} \mapsto \cdot_1\}$.
    \item \emph{Application let-binding}
      ($\mathbf{let}\; x = V_1\; V_2
      \;\mathbf{in}\; M$).
      Ternary in~$V_1$, $V_2$, and~$M$.
      $A(\cdot_1, \cdot_2, \cdot_3) =
      \{x \mapsto \{\cdot_1,\;
        \mathrm{argument} \mapsto \cdot_2\},\;
      \mathrm{result} \mapsto \cdot_3\}$.
    \item \emph{Value let-binding}
      ($\mathbf{let}\; x = V \;\mathbf{in}\; M$).
      Binary in~$V$ and~$M$.
      $A(\cdot_1, \cdot_2) =
      \{x \mapsto \{\mathrm{result} \mapsto \cdot_1\},\;
      \mathrm{result} \mapsto \cdot_2\}$.
    \item \emph{Tail call} ($V_1\; V_2$).
      Binary.\newline
      $A(\cdot_1, \cdot_2) =
      \{\mathrm{tailCall} \mapsto
        \{\cdot_1,\;
        \mathrm{argument} \mapsto \cdot_2\},\;
        \mathrm{result} \mapsto
      \{\mathrm{tailCall}.\mathrm{result}\}\}$.
  \end{enumerate}
  In each case the syntactic abstraction~$A_F$ is a fixed
  inheritance-calculus context whose holes are filled
  by the recursive translations of the subexpressions,
  satisfying~E4.
  \fi
\end{proof}

\begin{theorem}[\bilingual{Nonexpressibility of inheritance}{继承的不可宏表达性}]
  \label{thm:cartesian-nonexpressibility}
  \ifInheritanceChinese
  惰性 $\lambda$-演算不能宏可表达 $\mathscr{U}$(定义~\ref{def:language-universe})的继承构造子。
  \else
  The lazy $\lambda$-calculus cannot macro-express the
  inheritance constructors of $\mathscr{U}$
  (Definition~\ref{def:language-universe}).
  \fi
\end{theorem}

\begin{proof}
  \ifInheritanceChinese
  我们对 $L_0 = \text{惰性 $\lambda$-演算}$,$L_1 = \mathscr{U}$ 应用 Felleisen~\cite{felleisen1991-expressive-power} 的定理 3.14(i)。
  由定义~\ref{def:language-universe},$\mathscr{U} = L_0 + \{$mixin-树构造子$\}$ 是 $L_0$ 的一个保守性扩展(Felleisen~\cite{felleisen1991-expressive-power} 的定义 3.2)。
  我们验证四个条件:
  \begin{enumerate}
    \item[(i)] \emph{$L_0 \subseteq L_1$:} $\mathscr{U}$ 的语法包含所有 $L_0$-短语。
    \item[(ii)] \emph{不引入新的 $L_0$-短语:} mixin-树构造子不在 $L_0$ 中引入新短语。
    \item[(iii)] \emph{构造子不相交:} mixin-树构造子与 $L_0$-构造子不相交。
    \item[(iv)] \emph{语义保持:} 对纯 $\lambda$-项,定义~\ref{def:language-universe} 的步骤~(1) 与~(2) 是平凡的,步骤~(3) 由充分性(定理~\ref{thm:adequacy})给出惰性求值,因此 $\mathrm{eval}_{\mathscr{U}}$ 在 $L_0$-短语上与 $\mathrm{eval}_{L_0}$ 一致。
  \end{enumerate}
  该定理要求我们给出两个 $L_0$-项 $M$ 与 $N$,满足 $M \cong_0 N$ 但 $M \not\cong_1 N$,其中 $\cong_0$ 是 $L_0$ 中的操作等价,$\cong_1$ 是 $\mathscr{U}$ 中的操作等价。

  \paragraph{\bilingual{Operational equivalence in $L_0$}{$L_0$ 中的操作等价}}
  由 Böhm 树等价~\cite{abramsky1993-full-abstraction-lazy-lambda,abramsky1995-games-full-abstraction-lazy-lambda}(定理~\ref{thm:first-order-semantics}),对于封闭 $\lambda$-项 $M$ 与 $N$:$M \cong_0 N$ 当且仅当 $\mathrm{BT}(M) = \mathrm{BT}(N)$。

  \paragraph{\bilingual{Separating pair}{分离对}}
  定义 Church 布尔值与 Church 布尔相等:$\mathrm{true} = \lambda t.\,\lambda f.\, t$,$\mathrm{false} = \lambda t.\,\lambda f.\, f$,以及 $\mathrm{eq} = \lambda a.\,\lambda b.\, (a\;b)\;(b\;\mathrm{false}\;\mathrm{true})$。
  令
  \[
    M = \mathrm{eq}\;\mathrm{false}\;\mathrm{false}
    \qquad
    N = \mathrm{eq}\;\mathrm{true}\;\mathrm{true}.
  \]
  两者在惰性 $\lambda$-演算中均归约为 Church~$\mathrm{true}$,因此 $\mathrm{BT}(M) = \mathrm{BT}(N)$,从而 $M \cong_0 N$。

  \paragraph{\bilingual{Distinguishing $\mathscr{U}$-context}{区分 $\mathscr{U}$-上下文}}
  我们现在构造一个区分 $M$ 与 $N$ 的 $\mathscr{U}$-上下文。
  $M$ 与 $N$ 是语法封闭的 $\lambda$-项;由定义~\ref{def:language-universe},它们在 $\mathscr{U}$ 中的语义通过 ANF 转换后接 $\mathcal{T}$ 将其翻译为继承演算来给出。
  在扩展的 ANF 约定(第~\ref{sec:forward-translation} 节)下,应用位置上的每个自由值首先被绑定到一个名字:
  \else
  We apply Theorem~3.14(i) of
  Felleisen~\cite{felleisen1991-expressive-power} to the pair
  $L_0 = \text{lazy $\lambda$-calculus}$,
  $L_1 = \mathscr{U}$.
  By Definition~\ref{def:language-universe},
  $\mathscr{U} = L_0 + \{$mixin-tree constructors$\}$
  is a conservative extension of~$L_0$
  (Definition~3.2 of
  Felleisen~\cite{felleisen1991-expressive-power}).
  We verify the four conditions:
  \begin{enumerate}
    \item[(i)] \emph{$L_0 \subseteq L_1$:}
      The syntax of~$\mathscr{U}$ includes all
      $L_0$-phrases.
    \item[(ii)] \emph{No new $L_0$-phrases:}
      The mixin-tree constructors do not introduce new
      phrases in~$L_0$.
    \item[(iii)] \emph{Disjoint constructors:}
      The mixin-tree constructors are disjoint from the
      $L_0$-constructors.
    \item[(iv)] \emph{Semantics preserved:}
      For pure $\lambda$-terms, steps~(1) and~(2) of
      Definition~\ref{def:language-universe} are vacuous
      and step~(3) yields lazy evaluation by Adequacy
      (Theorem~\ref{thm:adequacy}), so
      $\mathrm{eval}_{\mathscr{U}}$ agrees with
      $\mathrm{eval}_{L_0}$ on $L_0$-phrases.
  \end{enumerate}
  The theorem requires us to exhibit two $L_0$-terms
  $M$ and $N$ with $M \cong_0 N$ but
  $M \not\cong_1 N$, where $\cong_0$ is operational
  equivalence in $L_0$ and $\cong_1$ is operational
  equivalence in $\mathscr{U}$.

  \paragraph{Operational equivalence in $L_0$}
  By B\"ohm tree
  equivalence~\cite{abramsky1993-full-abstraction-lazy-lambda,abramsky1995-games-full-abstraction-lazy-lambda}
  (Theorem~\ref{thm:first-order-semantics}), for closed
  $\lambda$-terms $M$ and $N$:
  $M \cong_0 N$ iff $\mathrm{BT}(M) = \mathrm{BT}(N)$.

  \paragraph{Separating pair}
  Define Church booleans and Church boolean equality:
  $\mathrm{true} = \lambda t.\,\lambda f.\, t$,
  $\mathrm{false} = \lambda t.\,\lambda f.\, f$, and
  $\mathrm{eq} = \lambda a.\,\lambda b.\,
  (a\;b)\;(b\;\mathrm{false}\;\mathrm{true})$.
  Let
  \[
    M = \mathrm{eq}\;\mathrm{false}\;\mathrm{false}
    \qquad
    N = \mathrm{eq}\;\mathrm{true}\;\mathrm{true}.
  \]
  Both reduce to Church~$\mathrm{true}$ in the
  lazy $\lambda$-calculus, so
  $\mathrm{BT}(M) = \mathrm{BT}(N)$ and hence
  $M \cong_0 N$.

  \paragraph{Distinguishing $\mathscr{U}$-context}
  We now construct a $\mathscr{U}$-context that
  separates $M$ and $N$.
  $M$ and $N$ are syntactically closed $\lambda$-terms;
  by Definition~\ref{def:language-universe}, their
  semantics in $\mathscr{U}$ is given by translating
  them to inheritance-calculus via ANF conversion followed
  by~$\mathcal{T}$.
  Under the extended ANF convention
  (Section~\ref{sec:forward-translation}), every free
  value used in an application is first bound to a name:
  \fi
  \begin{align*}
    \mathrm{ANF}(M) ={} &
    \mathbf{let}\;\mathrm{eq} =
    \lambda a.\,\lambda b.\,
    \mathbf{let}\; x_1 = a\;b \;\mathbf{in}\; \\
    &\qquad
    \mathbf{let}\; f_1 = \mathrm{false}
    \;\mathbf{in}\;
    \mathbf{let}\; x_2 = b\;f_1 \;\mathbf{in}\; \\
    &\qquad
    \mathbf{let}\; t_1 = \mathrm{true}
    \;\mathbf{in}\;
    \mathbf{let}\; x_3 = x_2\;t_1
    \;\mathbf{in}\;
    x_1\;x_3 \\
    & \mathbf{in}\;\;
    \mathbf{let}\; f_1 = \mathrm{false}
    \;\mathbf{in}\;
    \mathbf{let}\;\mathrm{eqFalse} = \mathrm{eq}\;f_1 \\
    &\quad
    \;\mathbf{in}\;
    \mathbf{let}\; f_2 = \mathrm{false}
    \;\mathbf{in}\;
    \mathrm{eqFalse}\;f_2
  \end{align*}
  \ifInheritanceChinese
  应用 $\mathcal{T}$:
  \else
  Applying~$\mathcal{T}$:
  \fi
  \begin{align*}
    \mathcal{T}(\mathrm{ANF}(M)) \;=\; \{&
      \mathrm{eq} \mapsto \{\mathrm{result} \mapsto
      \mathcal{T}(\lambda a.\,\lambda b.\,\cdots)\}, \\
      & f_1 \mapsto \{\mathrm{result} \mapsto
      \mathcal{T}(\mathrm{false})\}, \\
      & \mathrm{eqFalse} \mapsto
      \{\mathrm{eq}.\mathrm{result},\;
        \mathrm{argument} \mapsto
      \{f_1.\mathrm{result}\}\}, \\
      & f_2 \mapsto \{\mathrm{result} \mapsto
      \mathcal{T}(\mathrm{false})\}, \\
      & \mathrm{tailCall} \mapsto
      \{\mathrm{eqFalse}.\mathrm{result},\;
        \mathrm{argument} \mapsto
      \{f_2.\mathrm{result}\}\}, \\
      & \mathrm{result} \mapsto
      \{\mathrm{tailCall}.\mathrm{result}\}
    \;\}
  \end{align*}
  \ifInheritanceChinese
  对 $N$ 类似,以 $\mathrm{true}$ 替换 $\mathrm{false}$。
  注意 $\mathrm{eq}$ 是 $\mathcal{T}(\mathrm{ANF}(M))$ 的一个具名子项,通过值 let-绑定绑定。
  在 $\mathrm{eq}$ 的翻译内部,$\lambda b$ 作用域包含 $x_1$、$f_1$、$x_2$、$t_1$、$x_3$ 的 let-绑定以及一个尾调用。
  特别地,$\dbi{1}.\mathrm{argument}$ 到达外层 $\lambda a$ 作用域,而 $\dbi{0}.\mathrm{argument}$ 到达 $\lambda b$ 自身。\footnote{
    这些索引使用压缩约定,仅计数 $\lambda$-作用域。
    如翻译脚注(第~\ref{sec:forward-translation} 节)所述,每个 let-绑定与尾调用在继承演算中引入一个额外的作用域层级;形式索引相应地更大,但目标作用域相同。
  }

  $\mathscr{U}$-上下文使用继承在 $\mathrm{eq}$ 的 $\lambda b$ 作用域内添加一个 $\mathrm{repr}$ 记录,该记录投影第一个操作数。
  定义上下文 $C[\cdot]$ 为:
  \else
  and symmetrically for $N$ with $\mathrm{true}$ in place
  of $\mathrm{false}$.
  Note that $\mathrm{eq}$ is a named child
  of $\mathcal{T}(\mathrm{ANF}(M))$, bound via a value
  let-binding.
  Inside the translation of $\mathrm{eq}$, the inner
  $\lambda b$ scope contains the let-bindings
  for $x_1$, $f_1$, $x_2$, $t_1$, $x_3$, and a tail call.
  In particular, $\dbi{1}.\mathrm{argument}$ reaches
  the outer $\lambda a$ scope and
  $\dbi{0}.\mathrm{argument}$ reaches $\lambda b$
  itself.\footnote{
    These indices use the condensed
    convention, counting only $\lambda$-scopes.
    As noted in the translation footnote
    (Section~\ref{sec:forward-translation}),
    each let-binding and tail call introduces an
    additional scope level in inheritance-calculus;
    the formal index is correspondingly larger,
    but the target scope is the same.
  }

  The $\mathscr{U}$-context uses inheritance to add
  a $\mathrm{repr}$ record inside the $\lambda b$ scope
  of $\mathrm{eq}$ that projects the first operand.
  Define the context $C[\cdot]$ as:
  \fi
  \begin{align*}
    C[\cdot] ={} & \mathbf{let}\;
    \mathrm{extended} = \{[\cdot],\;
      \mathrm{eq} \mapsto \{
        \mathrm{result} \mapsto \{ \\
          &\quad
          \mathrm{result} \mapsto \{
            \mathrm{repr} \mapsto \{
              \mathrm{firstOperand}
              \mapsto \{\dbi{2}.\mathrm{argument}\}
    \}\}\}\}\} \\
    & \mathbf{in}\;
    \mathrm{extended}.\mathrm{tailCall}.\mathrm{repr}.\mathrm{firstOperand}
  \end{align*}
  \ifInheritanceChinese
  这里 $[\cdot]$ 是空位。
  注意 mixin 中 $\mathrm{result}$ 的两层嵌套:外层 $\mathrm{result}$ 进入 $\mathrm{eq}$ 的值 let-绑定包装器,而内层 $\mathrm{result}$ 进入 $\lambda a$ 作用域。
  第二个 $\mathrm{result}$ 则包含 $\lambda b$ 作用域,因此 $\mathrm{repr}$ 被放置在 $\lambda b$ 作用域层级。
  mixin $\{\mathrm{eq} \mapsto \{\mathrm{result} \mapsto \{\mathrm{result} \mapsto \{\mathrm{repr} \mapsto \cdots\}\}\}\}$ 是一个继承演算表达式:它仅使用 $\mathscr{U}$ 的 mixin-树构造子。
  外层 $\mathbf{let}$-绑定与投影是 $\lambda$-演算构造(引理~\ref{lem:let-macro})。
  因此 $C[\cdot]$ 是一个合法的 $\mathscr{U}$-上下文。

  $\mathrm{repr}$ 中的 de~Bruijn 索引 $\dbi{2}$ 计数作用域层级(在压缩约定下):$\dbi{0}$ 是 $\mathrm{repr}$ 自身,$\dbi{1}$ 是 $\lambda b$ 作用域,而 $\dbi{2}$ 是 $\lambda a$ 作用域。
  因此 $\mathrm{firstOperand}$ 投影 $\lambda a$ 的参数。

  由于 $\mathrm{eqFalse}$(相应地,$\mathrm{eqTrue}$)继承 $\mathrm{eq}.\mathrm{result}$,而 $\mathrm{tailCall}$ 继承 $\mathrm{eqFalse}.\mathrm{result}$($\lambda b$ 作用域),mixin 在 $\lambda b$ 作用域层级放置的 $\mathrm{repr}$ 记录通过继承链传播到 $\mathrm{tailCall}$。
  在 $C[M]$ 中:$\mathrm{eqFalse}$ 绑定向 $\lambda a$ 作用域提供 $\mathrm{argument} \mapsto \{f_1.\mathrm{result}\}$(即 $\mathrm{false}$),因此 $\mathrm{firstOperand}$ 求值为 $\mathrm{false} = \lambda t.\,\lambda f.\, f$。
  在 $C[N]$ 中:$\mathrm{eqTrue}$ 绑定提供 $\mathrm{argument} \mapsto \mathrm{true}$,因此 $\mathrm{firstOperand}$ 求值为 $\mathrm{true} = \lambda t.\,\lambda f.\, t$。

  定义 $D[\cdot] = (\lambda x.\, x\;\Omega\;I)\; C[\cdot]$。
  则 $\mathrm{eval}_{\mathscr{U}}$ 对 $D[M]$ 成立(因为 $\mathrm{false}\;\Omega\;I \to I$,收敛),但对 $D[N]$ 不成立(因为 $\mathrm{true}\;\Omega\;I \to \Omega$,发散)。
  因此 $M \not\cong_1 N$。

  \paragraph{\bilingual{Conclusion}{结论}}
  我们有 $M \cong_0 N$ 但 $M \not\cong_1 N$,故 ${\cong_0} \neq {(\cong_1|_{L_0})}$。
  由 Felleisen~\cite{felleisen1991-expressive-power} 的定理 3.14(i),惰性 $\lambda$-演算不能宏可表达 $\mathscr{U}$ 的继承构造子。
  \else
  Here $[\cdot]$ is the hole.
  Note the two levels of $\mathrm{result}$ nesting in the
  mixin: the outer $\mathrm{result}$ enters the value
  let-binding wrapper for $\mathrm{eq}$, and the inner
  $\mathrm{result}$ enters the $\lambda a$ scope.
  The second $\mathrm{result}$ then contains the
  $\lambda b$ scope, so $\mathrm{repr}$ is placed at
  the $\lambda b$ scope level.
  The mixin $\{\mathrm{eq} \mapsto \{\mathrm{result}
      \mapsto \{\mathrm{result} \mapsto
  \{\mathrm{repr} \mapsto \cdots\}\}\}\}$ is an
  inheritance-calculus expression: it uses only the
  mixin-tree constructors of $\mathscr{U}$.
  The outer $\mathbf{let}$-binding and projection are
  $\lambda$-calculus constructs (Lemma~\ref{lem:let-macro}).
  Thus $C[\cdot]$ is a valid $\mathscr{U}$-context.

  The de~Bruijn index $\dbi{2}$ in
  $\mathrm{repr}$ counts scope levels
  (in the condensed convention):
  $\dbi{0}$ is $\mathrm{repr}$ itself,
  $\dbi{1}$ is the $\lambda b$ scope, and
  $\dbi{2}$ is the $\lambda a$ scope.
  So $\mathrm{firstOperand}$ projects
  $\lambda a$'s argument.

  Since $\mathrm{eqFalse}$ (respectively $\mathrm{eqTrue}$)
  inherits $\mathrm{eq}.\mathrm{result}$,
  and $\mathrm{tailCall}$ inherits
  $\mathrm{eqFalse}.\mathrm{result}$ (the $\lambda b$
  scope), the $\mathrm{repr}$ record placed by the mixin
  at the $\lambda b$ scope level propagates through
  the inheritance chain to $\mathrm{tailCall}$.
  In $C[M]$: the $\mathrm{eqFalse}$ binding supplies
  $\mathrm{argument} \mapsto \{f_1.\mathrm{result}\}$
  (that is, $\mathrm{false}$) to the $\lambda a$ scope, so
  $\mathrm{firstOperand}$ evaluates to
  $\mathrm{false} = \lambda t.\,\lambda f.\, f$.
  In $C[N]$: the $\mathrm{eqTrue}$ binding supplies
  $\mathrm{argument} \mapsto \mathrm{true}$, so
  $\mathrm{firstOperand}$ evaluates to
  $\mathrm{true} = \lambda t.\,\lambda f.\, t$.

  Define $D[\cdot] = (\lambda x.\, x\;\Omega\;I)\;
  C[\cdot]$.
  Then $\mathrm{eval}_{\mathscr{U}}$ holds for $D[M]$
  (since $\mathrm{false}\;\Omega\;I \to I$, which
  converges)
  but not for $D[N]$
  (since $\mathrm{true}\;\Omega\;I \to \Omega$, which
  diverges).
  Therefore $M \not\cong_1 N$.

  \paragraph{Conclusion}
  We have $M \cong_0 N$ but $M \not\cong_1 N$, so
  ${\cong_0} \neq {(\cong_1|_{L_0})}$.
  By Theorem~3.14(i) of
  Felleisen~\cite{felleisen1991-expressive-power},
  the lazy $\lambda$-calculus cannot macro-express the
  inheritance constructors of~$\mathscr{U}$.
  \fi
\end{proof}

\begin{theorem}[\bilingual{Expressive asymmetry}{表达力的不对称}]
  \label{thm:expressive-asymmetry}
  \ifInheritanceChinese
  在语言全集 $\mathscr{U}$(定义~\ref{def:language-universe})中,继承演算比惰性 $\lambda$-演算具有严格更强的表达力:它能宏可表达 $\lambda$-演算能宏可表达的每个 $\mathscr{U}$-构造,但反之不然(Felleisen~\cite{felleisen1991-expressive-power} 的定义 3.17)。
  \else
  In the language universe $\mathscr{U}$
  (Definition~\ref{def:language-universe}),
  inheritance-calculus is strictly more expressive than
  the lazy $\lambda$-calculus: it can macro-express every
  $\mathscr{U}$-construct that the $\lambda$-calculus
  can, but not conversely
  (Definition~3.17 of
  Felleisen~\cite{felleisen1991-expressive-power}).
  \fi
\end{theorem}

\begin{proof}
  \ifInheritanceChinese
  我们须证明:(a)~继承演算能宏可表达惰性 $\lambda$-演算能宏可表达的每个 $\mathscr{U}$-构造,以及 (b)~反之不成立。

  \paragraph{\bilingual{Part (b): the lazy $\lambda$-calculus $\not\geq$
  inheritance-calculus}{第 (b) 部分:惰性 $\lambda$-演算 $\not\geq$ 继承演算}}
  $\lambda$-演算包含其自身的构造子($\lambda$、应用、变量),并能宏可表达 $\mathbf{let}$-绑定(引理~\ref{lem:let-macro}),但不能宏可表达 $\mathscr{U}$ 的继承构造子(定理~\ref{thm:cartesian-nonexpressibility})。
  继承演算包含那些构造子。
  因此,惰性 $\lambda$-演算对于 $\mathscr{U}$ 而言表达力不及继承演算。

  \paragraph{\bilingual{Part (a): inheritance-calculus $\geq$ the lazy
  $\lambda$-calculus}{第 (a) 部分:继承演算 $\geq$ 惰性 $\lambda$-演算}}
  我们证明继承演算能宏可表达 $\lambda$-演算所包含的或能宏可表达的每个 $\mathscr{U}$-构造。
  $\lambda$-演算包含 $\lambda$-抽象、应用与变量;它能宏可表达 $\mathbf{let}$-绑定(引理~\ref{lem:let-macro})。
  我们须证明继承演算能宏可表达这四个构造。
  翻译分两个阶段进行:
  \begin{enumerate}
    \item \emph{\bilingual{ANF conversion}{ANF 转换}。}
      A-范式变换~\cite{flanagan1993-essence-compiling-continuations} 将嵌套应用替换为 $\mathbf{let}$-绑定;由引理~\ref{lem:anf-equiv},它在 $\lambda$-演算中是宏可表达的。
      在反方向,$\mathbf{let}$-绑定由 Felleisen~\cite{felleisen1991-expressive-power} 的命题 3.7(重述为引理~\ref{lem:let-macro})在 $\lambda$-演算中是宏可表达的。
      因此,ANF 转换是 $\lambda$-演算与其 ANF 片段之间的一个宏可表达的等价。
    \item \emph{\bilingual{The translation $\mathcal{T}$}{翻译 $\mathcal{T}$}。}
      每个 ANF $\lambda$-演算构造子映射到继承演算上的一个固定语法抽象(定理~\ref{thm:forward-macro}),满足 Felleisen~\cite[Definition~3.11]{felleisen1991-expressive-power} 的条件 E4。
  \end{enumerate}
  由引理~\ref{lem:macro-composition},复合两个宏可表达翻译,得到一个从 $\lambda$-演算(含 $\mathbf{let}$)到继承演算的宏可表达翻译。
  由于继承演算平凡地包含其自身的构造子,它能宏可表达 $\lambda$-演算所包含的或能宏可表达的每个 $\mathscr{U}$-构造。
  \else
  We must show:
  (a)~inheritance-calculus can macro-express every
  $\mathscr{U}$-construct that the lazy $\lambda$-calculus
  can macro-express, and
  (b)~the converse fails.

  \paragraph{Part (b): the lazy $\lambda$-calculus $\not\geq$
  inheritance-calculus}
  The $\lambda$-calculus contains its own constructors
  ($\lambda$, application, variable) and can macro-express
  $\mathbf{let}$-binding (Lemma~\ref{lem:let-macro}), but
  cannot macro-express the inheritance constructors of
  $\mathscr{U}$
  (Theorem~\ref{thm:cartesian-nonexpressibility}).
  Inheritance-calculus contains those constructors.
  Therefore the lazy $\lambda$-calculus is not at least
  as expressive as inheritance-calculus with respect
  to~$\mathscr{U}$.

  \paragraph{Part (a): inheritance-calculus $\geq$ the lazy
  $\lambda$-calculus}
  We show that inheritance-calculus can macro-express
  every $\mathscr{U}$-construct that the
  $\lambda$-calculus contains or can macro-express.
  The $\lambda$-calculus contains $\lambda$-abstraction,
  application, and variable; it can macro-express
  $\mathbf{let}$-binding
  (Lemma~\ref{lem:let-macro}).
  We must show that inheritance-calculus can
  macro-express all four of these constructs.
  The translation proceeds in two stages:
  \begin{enumerate}
    \item \emph{ANF conversion.}
      The A-normal form
      transformation~\cite{flanagan1993-essence-compiling-continuations}
      replaces nested applications with
      $\mathbf{let}$-bindings; it is macro-expressible
      in the $\lambda$-calculus by
      Lemma~\ref{lem:anf-equiv}.
      In the reverse direction,
      $\mathbf{let}$-bindings are macro-expressible in
      the $\lambda$-calculus by
      Proposition~3.7 of
      Felleisen~\cite{felleisen1991-expressive-power}
      (restated as Lemma~\ref{lem:let-macro}).
      Thus ANF conversion is a macro-expressible
      equivalence between the $\lambda$-calculus and its
      ANF fragment.
    \item \emph{The translation $\mathcal{T}$.}
      Each ANF $\lambda$-calculus construct maps to a
      fixed syntactic abstraction over
      inheritance-calculus
      (Theorem~\ref{thm:forward-macro}), satisfying
      condition~E4 of
      Felleisen~\cite[Definition~3.11]{felleisen1991-expressive-power}.
  \end{enumerate}
  By Lemma~\ref{lem:macro-composition}, composing the
  two macro-expressible translations gives a
  macro-expressible translation from
  the $\lambda$-calculus (with $\mathbf{let}$) into
  inheritance-calculus.
  Since inheritance-calculus trivially contains its own
  constructors, it can macro-express every
  $\mathscr{U}$-construct that the $\lambda$-calculus
  contains or can macro-express.
  \fi
\end{proof}

\section{\bilingual{Multi-Target \texorpdfstring{$\this$}{this} Resolution in Other Systems}{其他系统中的多目标 \texorpdfstring{$\this$}{this} 解析}}
\label{app:scala-multi-target}

\ifInheritanceChinese
本附录说明两个具有代表性的系统——Scala~3 和 NixOS 模块系统——拒绝了继承演算能够自然处理的多目标 $\this$ 解析模式。
\else
This appendix demonstrates that two representative systems, Scala~3
and the NixOS module system, reject the multi-target $\this$ resolution
pattern that inheritance-calculus handles naturally.
\fi

\subsection*{\bilingual{Scala~3}{Scala~3}}

\ifInheritanceChinese
考虑两个独立扩展同一外部类的对象，每个对象各自提供内部 trait 的一份副本：
\else
Consider two objects that independently extend an outer class,
each providing its own copy of an inner trait:
\fi

\begin{lstlisting}[style=scala]
class MyOuter:
  trait MyInner:
    def outer = MyOuter.this

object Object1 extends MyOuter
object Object2 extends MyOuter
object HasMultipleOuters extends Object1.MyInner
                     with Object2.MyInner
\end{lstlisting}

\noindent
\ifInheritanceChinese
Scala~3 以如下错误拒绝了 \texttt{HasMultipleOuters}：
\else
Scala~3 rejects \texttt{HasMultipleOuters} with the error:
\fi

\begin{lstlisting}
trait MyInner is extended twice
object HasMultipleOuters cannot be instantiated since
  it has conflicting base types
  Object1.MyInner and Object2.MyInner
\end{lstlisting}

\noindent
\ifInheritanceChinese
该拒绝并非表面层面的限制，而是 DOT 类型论基础的必然后果~\cite{amin2016-dependent-object-types}。在 DOT 中，一个对象被构造为 $\nu(x : T)\, d$，其中 self 变量~$x$ 绑定到\emph{单一}对象。\texttt{Object1.MyInner} 和 \texttt{Object2.MyInner} 是不同的路径依赖类型，每个都将外部 $\this$ 约束到不同的对象。它们的交集要求 \texttt{outer} 同时返回 \texttt{Object1} 和 \texttt{Object2}，但 DOT 的 self 变量是单值的，使得这一交集不可实现。
\else
The rejection is not a surface-level restriction but a consequence
of DOT's type-theoretic foundations~\cite{amin2016-dependent-object-types}.
In DOT, an object is constructed as $\nu(x : T)\, d$, where the
self variable~$x$ binds to a \emph{single} object.
\texttt{Object1.MyInner} and \texttt{Object2.MyInner} are distinct
path-dependent types, each constraining the outer $\this$
to a different object.
Their intersection requires \texttt{outer} to simultaneously
return \texttt{Object1} and \texttt{Object2}, but DOT's self
variable is single-valued, making this intersection unrealizable.
\fi

\ifInheritanceChinese
等价的继承演算定义如下：
\else
The equivalent inheritance-calculus definition is:
\fi
\begin{align*}
  \{ \quad \mathrm{MyOuter} &\mapsto \{
      \mathrm{MyInner} \mapsto \{
    \mathrm{outer} \mapsto \{\qualifiedthis{\mathrm{MyOuter}}\}\}\}, \\
    \mathrm{Object1} &\mapsto \{\mathrm{MyOuter}\}, \\
    \mathrm{Object2} &\mapsto \{\mathrm{MyOuter}\}, \\
    \mathrm{HasMultipleOuters} &\mapsto \{
      \mathrm{Object1}.\mathrm{MyInner},\;
    \mathrm{Object2}.\mathrm{MyInner}\}
  \quad \}
\end{align*}
\ifInheritanceChinese
此定义在继承演算中是良定义的。$\mathrm{outer}$ 内部的引用 $\qualifiedthis{\mathrm{MyOuter}}$ 具有 de~Bruijn 索引 $n = 1$：从 $\mathrm{outer}$ 的封闭作用域 $\mathrm{MyInner}$ 出发，经过一步 $\this$ 便到达 $\mathrm{MyOuter}$。当 $\mathrm{HasMultipleOuters}$ 同时继承 $\mathrm{Object1}.\mathrm{MyInner}$ 和 $\mathrm{Object2}.\mathrm{MyInner}$ 时，$\this$ 函数（方程~\ref{eq:this}）通过在 $\supers(\mathrm{HasMultipleOuters})$ 中搜索与 $\mathrm{MyInner}$ 定义位置匹配的覆盖路径来解析 $\qualifiedthis{\mathrm{MyOuter}}$。它找到两条继承位置路径，一条经由 $\mathrm{Object1}$，另一条经由 $\mathrm{Object2}$，并返回两者。两条路径都通向继承自同一 $\mathrm{MyOuter}$ 的记录，因此查询 $\mathrm{HasMultipleOuters}.\mathrm{outer}$ 无论走哪条路由都能得到 $\mathrm{MyOuter}$ 的 $\properties$。由于属性解析通过集合并来聚合结果（这是一种本质上幂等的操作），来自两条路由的相同属性定义无冲突地合并。因此，该演算消除了对任何人为路由优先级或线性化规则的需求。
\else
This is well-defined in inheritance-calculus.
The reference
$\qualifiedthis{\mathrm{MyOuter}}$ inside $\mathrm{outer}$
has de~Bruijn index $n = 1$:
starting from $\mathrm{outer}$'s enclosing scope
$\mathrm{MyInner}$, one $\this$ step reaches
$\mathrm{MyOuter}$.
When $\mathrm{HasMultipleOuters}$ inherits from both
$\mathrm{Object1}.\mathrm{MyInner}$ and
$\mathrm{Object2}.\mathrm{MyInner}$,
the $\this$ function (equation~\ref{eq:this}) resolves
$\qualifiedthis{\mathrm{MyOuter}}$ by searching through
$\supers(\mathrm{HasMultipleOuters})$ for override paths matching
$\mathrm{MyInner}$'s definition site.
It finds two inheritance-site paths, one through
$\mathrm{Object1}$ and one through $\mathrm{Object2}$, and
returns both.
Both paths lead to records that inherit from the same
$\mathrm{MyOuter}$, so querying
$\mathrm{HasMultipleOuters}.\mathrm{outer}$ yields the
$\properties$ of $\mathrm{MyOuter}$ regardless of which
route is taken.
Because attribute resolution aggregates results via set union (which is an intrinsically idempotent operation), identical property definitions from the two routes merge without conflict. Thus, the calculus eliminates the need for any artificial route-priority or linearization rules.
\fi

\subsection*{\bilingual{NixOS Module System}{NixOS 模块系统}}

\ifInheritanceChinese
NixOS 模块系统以静态错误拒绝了相同的模式。该翻译使用了模块系统自身的抽象：\texttt{deferredModule} 对应 trait（即未求值的模块值，被导入到其他不动点中），\texttt{submoduleWith} 对应对象（在其自身的不动点中求值）。
\else
The NixOS module system rejects the same pattern with a static error.
The translation uses the module system's own abstractions:
\texttt{deferredModule} for traits, which are unevaluated module values imported
into other fixpoints, and \texttt{submoduleWith} for objects, which are evaluated
in their own fixpoints.
\fi

\begin{lstlisting}[style=nix]
let
  lib = (import <nixpkgs> { }).lib;
  result = lib.evalModules {
    modules = [
      (toplevel@{ config, ... }: {
        # class MyOuter { trait MyInner {
        #   def outer = MyOuter.this } }
        options.MyOuter = lib.mkOption {
          default = { };
          type = lib.types.deferredModuleWith {
            staticModules = [
              (MyOuter: {
                options.MyInner = lib.mkOption {
                  default = { };
                  type = lib.types.deferredModuleWith {
                    staticModules = [
                      (MyInner: {
                        options.outer = lib.mkOption {
                          default = { };
                          type =
                            lib.types.deferredModuleWith {
                              staticModules =
                                [ toplevel.config.MyOuter ];
                            };
                        };
                      })
                    ];
                  };
                };
              })
            ];
          };
        };
        # object Object1 extends MyOuter
        options.Object1 = lib.mkOption {
          default = { };
          type = lib.types.submoduleWith {
            modules = [ toplevel.config.MyOuter ];
          };
        };
        # object Object2 extends MyOuter
        options.Object2 = lib.mkOption {
          default = { };
          type = lib.types.submoduleWith {
            modules = [ toplevel.config.MyOuter ];
          };
        };
        # object HasMultipleOuters extends
        #   Object1.MyInner with Object2.MyInner
        options.HasMultipleOuters = lib.mkOption {
          default = { };
          type = lib.types.submoduleWith {
            modules = [
            toplevel.config.Object1.MyInner
            toplevel.config.Object2.MyInner
          ];
          };
        };
      })
    ];
  };
in
  builtins.attrNames result.config.HasMultipleOuters.outer
\end{lstlisting}

\noindent
\ifInheritanceChinese
NixOS 模块系统以如下错误拒绝了此代码：
\else
The NixOS module system rejects this with:
\fi

\begin{lstlisting}
error: The option `HasMultipleOuters.outer'
  in `<unknown-file>'
  is already declared
  in `<unknown-file>'.
\end{lstlisting}

\noindent
\ifInheritanceChinese
该拒绝发生的原因是：\texttt{Object1.\allowbreak{}My\-Inner} 和 \texttt{Object2.\allowbreak{}My\-Inner} 各自携带了 \texttt{My\-Inner} 延迟模块的完整 \texttt{static\-Modules}，其中包括声明 \texttt{options.\allowbreak{}outer}。当 \texttt{HasMultipleOuters} 同时导入两者时，模块系统的 \texttt{merge\-Option\-Decls} 遇到了同一选项的两个声明，并产生静态错误；它无法表达两个声明源自同一 trait 且应被统一这一事实。
\else
The rejection occurs because
\texttt{Object1.\allowbreak{}My\-Inner} and
\texttt{Object2.\allowbreak{}My\-Inner} each carry the full
\texttt{static\-Modules} of the \texttt{My\-Inner} deferred
module, including the declaration
\texttt{options.\allowbreak{}outer}.
When \texttt{HasMultipleOuters} imports both, the module system's
\texttt{merge\-Option\-Decls} encounters two declarations of the
same option and raises a static error; it cannot express that
both declarations originate from the same trait and should be
unified.
\fi

\ifInheritanceChinese
该拒绝与上文 Scala 的``冲突基类型''错误类似：两个系统都假设每个选项或类型成员只有单一的声明位置，并拒绝两条继承路由独立引入同一声明的多目标情形。在继承演算中，相同的模式是良定义的，因为 $\overrides$（方程~\ref{eq:overrides}）能够识别两条路由通向同一定义，而 $\supers$（方程~\ref{eq:supers}）则无重复地收集两个继承位置上下文。
\else
The rejection is analogous to Scala's ``conflicting base types''
error above: both systems assume that each option or type member
has a single declaration site, and reject the multi-target situation
where two inheritance routes introduce the same declaration
independently.
In inheritance-calculus, the same pattern is well-defined because
$\overrides$ (equation~\ref{eq:overrides}) recognizes that both
routes lead to the same definition, and $\supers$
(equation~\ref{eq:supers}) collects both inheritance-site
contexts without duplication.
\fi

\section{Encoding Datalog in Inheritance-Calculus}
\label{app:datalog-encoding}

Section~\ref{sec:emergent-phenomena} observed that
inheritance-calculus natively produces the relational semantics
of logic programming.
This appendix makes that observation precise by giving a
systematic encoding of Datalog into inheritance-calculus.
The encoding uses no extensions to the calculus; it is
expressible in the three primitives of Section~\ref{sec:syntax}.
All examples in this appendix have executable MIXINv2
implementations\anon[~(supplementary material)]{~\cite{mixin2025}}.
There are three ingredients to keep separate.
First, relations are stored as tries, so tuple lookup is indexed
by prefixes rather than by scanning a flat set of facts.
Second, the control skeleton is written in the visitor-pattern
style already used in Section~\ref{sec:case-study}; from a
functional-programming viewpoint, this is the same role played by
Scott-encoded case analysis.
Third, inheritance union and tabled evaluation give the resulting
program its relational, least-fixed-point semantics.
The appendix proceeds in that order: data representation first,
then control flow, then recursion.

\subsection{Relations as Tries}

A binary relation over a finite domain $D = \{a, b, c, \ldots\}$
is represented as a \emph{trie}: a mixin tree whose paths
encode tuples.
The trie is built from three base scopes:
$\mathrm{Relation}$ (abstract base),
$\mathrm{Cons}$ (non-empty node), and $\mathrm{Nil}$ (terminator).
Each constant $d \in D$ inherits from $\mathrm{Cons}$ and carries
a $\mathrm{TailMap}$ that provides per-constant sub-tries:
\begin{align*}
  &\mathrm{Relation} \mapsto \{\},\quad
  \mathrm{Cons} \mapsto \{\mathrm{Relation},\;
  \mathrm{Tail} \mapsto \{\mathrm{Relation}\}\},\quad
  \mathrm{Nil} \mapsto \{\mathrm{Relation}\}\\[2pt]
  &d \mapsto \{\mathrm{Cons},\;
    \mathrm{Tail} \mapsto \{\},\;
  \mathrm{TailMap} \mapsto \{d \mapsto \{\mathrm{Tail}\}\}\}
\end{align*}
A fact $R(d_1, d_2)$ is encoded as a trie path of depth~2,
with a $\mathrm{Nil}$ terminator.
The per-constant sub-tries are accessed via
$\mathrm{TailMap}.d$:
\begin{align*}
  R(a, b) \;\leadsto\quad
  R \mapsto \{a,\; \mathrm{TailMap} \mapsto \{a \mapsto
  \{b,\; \mathrm{TailMap} \mapsto \{b \mapsto \{\mathrm{Nil}\}\}\}\}\}
\end{align*}
Multiple facts for the same relation are accumulated by
inheritance (trie union).
Adding $R(b, c)$:
\[
  R \mapsto \left\{
    \begin{aligned}
      &a,\; b,\\
      &\mathrm{TailMap} \mapsto \left\{
        \begin{aligned}
          &a \mapsto \{b,\; \mathrm{TailMap} \mapsto \{b \mapsto \{\mathrm{Nil}\}\}\},\\
          &b \mapsto \{c,\; \mathrm{TailMap} \mapsto \{c \mapsto \{\mathrm{Nil}\}\}\}
      \end{aligned}\right\}
  \end{aligned}\right\}
\]
The point of the trie representation is not merely that it is a
tree-shaped encoding.
It also provides the indexing structure used by the subsequent
joins: once a constant $d$ has been selected, the encoding moves
directly to $\mathrm{TailMap}.d$ and continues from the matching
sub-trie.
This is why the appendix speaks about tries rather than about an
arbitrary encoding of finite sets of tuples.

\subsection{Infrastructure}

Each constant $d$ carries four groups of boilerplate
definitions used by the encoding.
These play a role analogous to derived type-class instances
in Haskell: each constant must provide a fixed set of
properties that the encoding's generic patterns rely on.
Inheritance-calculus has no built-in derivation mechanism,
so the boilerplate is written out manually; however, the
amount of boilerplate per constant is $O(1)$
and does not grow with the domain size $|D|$,
so adding a new constant
does not require modifying any existing scope, preserving
the expression-problem property.

% 我们可以通过引入类似ECMAScript的Symbol机制来消除这些重复定义，但这让语言更复杂了

\begin{description}
  \item[Acceptance]
    The visitor-pattern infrastructure from
    Section~\ref{sec:case-study}.
    Operationally, this is the object-oriented spelling of a
    case split on the current constant.
    Each constant $d$ defines a $\mathrm{VisitorMap}$
    containing a branch keyed by $d$, a $\mathrm{Visitor}$
    with a $\mathrm{Visited}$ slot for per-constant results,
    and an $\mathrm{Accepted}$ property that delegates to the
    matching branch's $\mathrm{Visited}$ value.
    This is the dispatch mechanism: given a scope that may
    contain multiple constants, $\mathrm{Acceptance}$ selects the
    branch corresponding to each constant present and
    collects the per-constant $\mathrm{Visited}$ results.

  \item[WithVisitorTail]
    Maps each constant $d$ to its sub-trie entry
    $\mathrm{TailMap}.d$, labeled as $\mathrm{VisitorTail}$.
    This provides the right-hand side of a join\footnote{
      In the context of the visitor pattern, this acts as the visitor.
    }: when two relations' $\mathrm{VisitorMap}$ entries are
    composed, $\mathrm{VisitorTail}$ carries the joined
    relation's row data for the matched constant.

  \item[WithVisiteeTail]
    Maps each constant $d$ to its sub-trie entry
    $\mathrm{TailMap}.d$, labeled as $\mathrm{VisiteeTail}$.
    This provides the left-hand side\footnote{
      This acts as the visitee.
    }: the
    relation being iterated. Together,
    $\mathrm{WithVisitorTail}$ and $\mathrm{WithVisiteeTail}$
    enable value-level joins where only constants present
    in \emph{both} relations produce output.

  \item[Replacement]
    Given a $\mathrm{NewTail}$ value, produces a
    $\mathrm{Replaced}$ scope that inherits the label $d$
    but substitutes $\mathrm{TailMap}.d$ with
    $\mathrm{NewTail}$.
    This is the output construction mechanism: it
    builds output tuples by substituting new sub-tries into
    constant wrappers.
    Crucially, $\mathrm{Replacement}$ is defined on the
    per-constant $\mathrm{TailMap}$ entry (not on the merged
    scope), ensuring that output construction operates on
    individual constants rather than their union.
\end{description}

As a running example, we use the executable two-hop encoding
$\mathrm{twoHop}(X, Z) \mathrel{{:}\!{-}} \mathrm{edge}(X, Y),\; \mathrm{edge}(Y, Z)$
from
\anon[\nolinkurl{RelationalTwoHopRepeatedVisitor.mixin.yaml}~(supplementary material)]{\nolinkurl{RelationalTwoHopRepeatedVisitor.mixin.yaml}~\cite{mixin2025}}.
Its top-level scopes are named after the stages of the
encoding: \texttt{\_XAcceptance} iterates the first column,
\texttt{\_YAcceptance} performs the shared-variable join,
\texttt{\_NilAcceptance0} confirms the end of the first fact,
\texttt{\_ZAcceptance} iterates the output column, and
\texttt{\_NilAcceptance1} triggers the final
\texttt{Replacement} chain.
The schematic encodings below should be read as abstractions
of that executable MIXINv2 file.
If one temporarily ignores relational accumulation and reads
each dispatch as an ordinary case split, the control flow is:
choose $X$ from the first relation, choose a matching $Y$,
confirm that the first tuple is complete, choose $Z$ from the
second relation's matching row, confirm that the second tuple
is complete, and emit $(X, Z)$.
The rest of the appendix explains how that familiar
single-valued control skeleton is realised by visitor dispatch
over tries and then lifted to relational semantics by
inheritance union.

\subsection{Encoding Rules}

A Datalog rule $H(\bar{X}) \mathrel{{:}\!{-}} B_1(\bar{Y}_1), \ldots, B_n(\bar{Y}_n)$
is encoded by nesting three mechanisms inside that control
skeleton:

\begin{enumerate}
  \item \textbf{Column iteration.}
    For each column position, an $\mathrm{Acceptance}$ dispatch
    iterates over the constants present.
    The $\mathrm{VisitorMap}$ is populated from
    $\mathrm{WithVisiteeTail}$, so each matched constant's
    $\mathrm{Visitor}$ receives $\mathrm{VisiteeTail}$: the
    per-constant sub-trie at that position.
    The result of processing each constant is placed in
    $\mathrm{Visited}$, and $\mathrm{Accepted}$ collects all
    per-constant results.
    (Details in Sections~\ref{sec:encoding-variables},
    \ref{sec:encoding-constants}, and~\ref{sec:encoding-multi}.)

  \item \textbf{Join.}
    For each shared variable between body atoms, the inner
    $\mathrm{Acceptance}$'s $\mathrm{VisitorMap}$ inherits from
    \emph{both} $\mathrm{WithVisitorTail}$ (from the joined
    relation) and $\mathrm{WithVisiteeTail}$ (from the iterated
    relation).
    A constant $d$ produces a $\mathrm{Visitor}$ entry only if
    $d$ appears in both relations, achieving a value-level
    join.
    The $\mathrm{Visitor}$ then receives both
    $\mathrm{VisiteeTail}$ and $\mathrm{VisitorTail}$,
    providing access to both relations' data for the
    matched constant.

  \item \textbf{Output construction.}
    Inside the innermost $\mathrm{Nil}$ confirmation (verifying
    that each body atom's tuple is complete), chain
    $\mathrm{Replacement}$ operations from the innermost
    output column outward, each setting $\mathrm{NewTail}$ to
    the previous $\mathrm{Replaced}$ value, terminated by
    $\mathrm{Nil}$.
    The final $\mathrm{Replaced}$ value is assigned to
    $\mathrm{Visited}$.
\end{enumerate}

\subsection{Variable Join}
\label{sec:encoding-variables}

Consider the simplest join:
$\mathrm{composed}(X, Z) \mathrel{{:}\!{-}} \alpha(X, Y),\; \beta(Y, Z)$.
The executable self-join instance of this pattern is
\anon[\nolinkurl{RelationalTwoHopRepeatedVisitor.mixin.yaml}~(supplementary material)]{\nolinkurl{RelationalTwoHopRepeatedVisitor.mixin.yaml}~\cite{mixin2025}},
whose scope names are shown in parentheses below.
This is the right place to see the division of labour from the
section introduction: visitor dispatch provides the control flow,
the trie provides direct access to the matching tails, and
$\supers$ accumulates the outputs produced by each successful
match.

The shared variable $Y$ requires matching constants
between $\alpha$'s second column and $\beta$'s first column.
The encoding nests five $\mathrm{Acceptance}$ dispatches:
iterate $X$ from $\alpha$, join on $Y$ between $\alpha$
and $\beta$, confirm $\mathrm{Nil}$ after $Y$ in $\alpha$,
iterate $Z$ from $\beta$'s matched row, and confirm
$\mathrm{Nil}$ after $Z$.
We present each level as a named sub-scope.

\paragraph{Level~1: iterate $X$ from $\alpha$}
(Executable scope: \texttt{\_XAcceptance}.)
{\small
  \[
    \mathrm{xAcc} \mapsto \left\{
      \begin{aligned}
        &\alpha.\mathrm{Acceptance},\; \mathrm{VisitorMap} \mapsto \{\alpha.\mathrm{WithVisiteeTail}\},\\
        &\mathrm{Visitor} \mapsto \left\{
          \begin{aligned}
            &\_X \mapsto \{\mathrm{VisiteeTail}\},\; \ldots,\\
            &\mathrm{Visited} \mapsto \{\mathrm{yAcc}.\mathrm{Accepted}\}
        \end{aligned}\right\}
    \end{aligned}\right\}
\]}

\paragraph{Level~2: join $\alpha$'s $Y$-column with $\beta$'s first column}

(Executable scope: \texttt{\_YAcceptance}.)

The $\mathrm{VisitorMap}$ inherits from \emph{both}
$\beta.\mathrm{WithVisitorTail}$ and
$\_X.\mathrm{WithVisiteeTail}$, producing an entry for
constant $d$ only if $d$ appears in both $\alpha$'s
second column (via $\_X$) and $\beta$'s first column.
{\small
  \[
    \mathrm{yAcc} \mapsto \left\{
      \begin{aligned}
        &\_X.\mathrm{Acceptance},\\
        &\mathrm{VisitorMap} \mapsto \{\beta.\mathrm{WithVisitorTail},\; \_X.\mathrm{WithVisiteeTail}\},\\
        &\mathrm{Visitor} \mapsto \left\{
          \begin{aligned}
            &\_Y_0 \mapsto \{\mathrm{VisiteeTail}\},\; \mathrm{VisitorTail} \mapsto \{\beta.\mathrm{Tail}\},\\
            &\ldots,\\
            &\mathrm{Visited} \mapsto \{\mathrm{nilAcc}_0.\mathrm{Accepted}\}
        \end{aligned}\right\}
    \end{aligned}\right\}
\]}

\paragraph{Level~3: confirm $\mathrm{Nil}$ after $Y$ in $\alpha$}

(Executable scope: \texttt{\_NilAcceptance0}.)

This checks that $\alpha(X, Y)$ is a complete two-column
fact (its trie path terminates at $\mathrm{Nil}$ after $Y$).
Inside the $\mathrm{Nil}$ branch, the encoding proceeds to
iterate $\beta$'s second column via $\mathrm{VisitorTail}$:
{\small
  \[
    \mathrm{nilAcc}_0 \mapsto \left\{
      \begin{aligned}
        &\_Y_0.\mathrm{Acceptance},\\
        &\mathrm{VisitorMap} \mapsto \left\{\mathrm{Nil} \mapsto \left\{
            \begin{aligned}
              &\ldots,\\
              &\mathrm{Visited} \mapsto \{\mathrm{zAcc}.\mathrm{Accepted}\}
        \end{aligned}\right\}\right\}
    \end{aligned}\right\}
\]}

\paragraph{Level~4: iterate $Z$ from $\beta$'s matched row}
(Executable scope: \texttt{\_ZAcceptance}.)
{\small
  \[
    \mathrm{zAcc} \mapsto \left\{
      \begin{aligned}
        &\mathrm{VisitorTail}.\mathrm{Acceptance},\; \mathrm{VisitorMap} \mapsto \{\mathrm{VisitorTail}.\mathrm{WithVisiteeTail}\},\\
        &\mathrm{Visitor} \mapsto \left\{
          \begin{aligned}
            &\_Z \mapsto \{\mathrm{VisiteeTail}\},\; \ldots,\\
            &\mathrm{Visited} \mapsto \{\mathrm{nilAcc}_1.\mathrm{Accepted}\}
        \end{aligned}\right\}
    \end{aligned}\right\}
\]}

\paragraph{Level~5: confirm $\mathrm{Nil}$ after $Z$ and construct output}
(Executable scope: \texttt{\_NilAcceptance1}, with
\texttt{\_ZReplacement} and \texttt{\_XReplacement}.)
{\small
  \[
    \mathrm{nilAcc}_1 \mapsto \left\{
      \begin{aligned}
        &\mathrm{VisiteeTail}.\mathrm{Acceptance},\\
        &\mathrm{VisitorMap} \mapsto \left\{\mathrm{Nil} \mapsto \left\{
            \begin{aligned}
              &\mathrm{zR} \mapsto \{\_Z.\mathrm{Replacement},\; \mathrm{NewTail} \mapsto \{\mathrm{Nil}\}\},\\
              &\mathrm{xR} \mapsto \{\_X.\mathrm{Replacement},\; \mathrm{NewTail} \mapsto \{\mathrm{zR}.\mathrm{Replaced}\}\},\\
              &\mathrm{Visited} \mapsto \{\mathrm{xR}.\mathrm{Replaced}\}
        \end{aligned}\right\}\right\}
    \end{aligned}\right\}
\]}

\paragraph{Final result}
\[
  \mathrm{composed} \mapsto \{\mathrm{xAcc}.\mathrm{Accepted}\}
\]

\paragraph{Semantics}
At Level~2, the combined $\mathrm{VisitorMap}$ ensures that a
constant $d$ produces a $\mathrm{Visitor}$ only when $d$
appears in both $\alpha$'s second column \emph{and}
$\beta$'s first column.
The $\mathrm{Visitor}$ for each matched constant $d$
independently receives $\mathrm{VisiteeTail}$ ($\alpha$'s
data for~$d$) and $\mathrm{VisitorTail}$ ($\beta$'s row
for~$d$).
Output construction (Level~5) operates inside this
per-constant $\mathrm{Visitor}$, so $\mathrm{Replacement}$
acts on the \emph{individual} constant's data, not on the
merged scope.
Each matched constant $d$ produces its own output tuples
in $\mathrm{Visited}$, and these are collected by
$\mathrm{Accepted}$ via trie union ($\supers$).
This yields standard Datalog tuple-level semantics: only
constants that actually match between both relations
contribute to the output.
For the concrete instance
$\mathrm{edge} = \{(a,b),\; (b,c),\; (c,a)\}$, as encoded in
\anon[\nolinkurl{RelationalTwoHopRepeatedVisitor.mixin.yaml}~(supplementary material)]{\nolinkurl{RelationalTwoHopRepeatedVisitor.mixin.yaml}~\cite{mixin2025}},
Level~1 selects $a$, Level~2 matches $b$, Level~3 confirms
that $\mathrm{edge}(a,b)$ ends in $\mathrm{Nil}$, Level~4
iterates $\mathrm{edge}.\mathrm{TailMap}.b$ and selects $c$,
and Level~5 confirms $\mathrm{Nil}$ again before the two
$\mathrm{Replacement}$ operations reconstruct $(a,c)$.

\subsection{Constant Match and Navigation}
\label{sec:encoding-constants}

Constants in Datalog atoms use two mechanisms, neither of
which is a join.
They are simpler than variable joins because no value from a
second relation needs to be matched against the currently chosen
constant:

\paragraph{Type existence match}
An atom $R(X, b)$ with a constant in the second position
requires checking that constant $b$ exists in $R$'s
second column at the matched row.
This is encoded as an $\mathrm{Acceptance}$ with a
hand-crafted $\mathrm{VisitorMap}$ containing \emph{only}
the $b$ branch (instead of composing from
  $\mathrm{WithVisitorTail}$, which would create branches for
all constants):
{\small
  \begin{align*}
    \mathrm{VisitorMap} \mapsto \{b \mapsto \{\ldots\}\}
\end{align*}}
If $b$ exists in the dispatching scope, the
$b$ branch fires; otherwise, $\mathrm{Accepted}$
receives no matching branch.

\paragraph{Direct navigation}
An atom $R(b, Z)$ with a constant in the first position
navigates directly to $b$'s sub-trie:
\[
  \mathrm{edgeB}_Z \mapsto \{R.\mathrm{TailMap}.b\}
\]
This bypasses the visitor mechanism entirely, accessing
the per-constant $\mathrm{TailMap}$ entry to restrict the
scope to only $b$'s data.

\paragraph{Example}
$\mathrm{viaB}(X, Z) \mathrel{{:}\!{-}} \mathrm{edge}(X, b),\; \mathrm{edge}(b, Z)$
combines both mechanisms.
The executable MIXINv2 encoding is
\anon[\nolinkurl{RelationalConstant.mixin.yaml}~(supplementary material)]{\nolinkurl{RelationalConstant.mixin.yaml}~\cite{mixin2025}}.
The encoding iterates $\mathrm{edge}$ for $X$ (Level~1),
then checks for constant $b$ in $X$'s row with a
hand-crafted $\mathrm{VisitorMap}$ (Level~2), navigates
$\mathrm{edge}.\mathrm{TailMap}.b$ for $Z$ (Level~3),
confirms $\mathrm{Nil}$ (Level~4), and constructs $(X, Z)$
via $\mathrm{Replacement}$.

\subsection{Multi-Body Rules}
\label{sec:encoding-multi}

Rules with $n > 2$ body atoms chain additional join levels.
Each shared variable introduces one combined
$\mathrm{WithVisitorTail}$/$\mathrm{WithVisiteeTail}$ join
inside the previous level's $\mathrm{Visitor}$,
and the output construction is placed inside the innermost
$\mathrm{Nil}$ confirmation.
So the binary-rule encoding is the whole story structurally:
larger bodies only add more nested copies of the same pattern.

\paragraph{Three-body rule}
$\mathrm{chain}(X, W) \mathrel{{:}\!{-}} \alpha(X, Y),\; \beta(Y, Z),\; \gamma(Z, W)$
requires two joins: on $Y$ (between $\alpha$ and $\beta$)
and on $Z$ (between $\beta$ and $\gamma$).
The structure is: iterate $X$ from $\alpha$, join on $Y$
(getting both $\alpha$'s and $\beta$'s data), confirm
$\mathrm{Nil}$ after $Y$ in $\alpha$, then inside the
$\mathrm{Nil}$ branch join $\beta$'s $Z$-column with
$\gamma$'s first column, iterate $W$ from $\gamma$'s
matched row, confirm $\mathrm{Nil}$, and construct
$(X, W)$.
The $Z$-join is nested inside the $Y$-join's
$\mathrm{Nil}$ confirmation, ensuring that $\beta$'s data
for the matched $Y$ constant (available via
$\mathrm{VisitorTail}$) is used to iterate $\beta$'s
second column.

\paragraph{Mixed constants and variables}
$\mathrm{constJoin}(X, W) \mathrel{{:}\!{-}} \mathrm{edge}(X, Y),\; \mathrm{edge}(Y, b),\; \mathrm{edge}(b, W)$
combines a variable join on $Y$, a constant match on $b$
(hand-crafted $\mathrm{VisitorMap}$ with only the $b$ entry),
and direct navigation via $\mathrm{edge}.\mathrm{TailMap}.b$
for $W$.
The corresponding executable MIXINv2 file is
\anon[\nolinkurl{RelationalConstantJoin.mixin.yaml}~(supplementary material)]{\nolinkurl{RelationalConstantJoin.mixin.yaml}~\cite{mixin2025}}.

\subsection{Multi-Variable Join}

When a rule shares \emph{multiple} variables between two
body atoms, each shared variable requires its own
$\mathrm{Acceptance}$ check.
The checks are chained: the outer $\mathrm{Acceptance}$
joins on one column pair using combined
$\mathrm{WithVisitorTail}$/$\mathrm{WithVisiteeTail}$, and
its $\mathrm{Visitor}$ contains an inner $\mathrm{Acceptance}$
joining on another column pair.

\paragraph{Example}
$\mathrm{multiJoin}(X, Y) \mathrel{{:}\!{-}} r(X, Y),\; s(X, Y)$
shares both $X$ and $Y$.
The executable MIXINv2 encoding is
\anon[\nolinkurl{RelationalMultiVariableJoin.mixin.yaml}~(supplementary material)]{\nolinkurl{RelationalMultiVariableJoin.mixin.yaml}~\cite{mixin2025}}.
The outer $\mathrm{Acceptance}$ joins $r$'s and $s$'s first
columns (variable $X$) using combined
$\mathrm{WithVisitorTail}$/$\mathrm{WithVisiteeTail}$.
Inside the outer $\mathrm{Visitor}$, an inner
$\mathrm{Acceptance}$ joins $r$'s and $s$'s second columns
(variable $Y$) similarly.
Both joins must match for the output to be produced, and
the per-constant $\mathrm{Visited}$ pattern ensures that
only constants present in both relations contribute.

\subsection{Recursive Rules}
\label{sec:recursive-rules}

Recursive Datalog rules, such as transitive closure
$\mathrm{path}(X, Y) \mathrel{{:}\!{-}} \mathrm{edge}(X, Z),\; \mathrm{path}(Z, Y)$,
are encoded identically to non-recursive rules: the
head relation $\mathrm{path}$ appears both as a defined scope
and as a reference in the body (providing
$\mathrm{WithVisitorTail}$ for the $Z$-join).
This should not be read as ordinary $\this$ of a
single-valued lazy program.
Operationally, tabled evaluation
(Section~\ref{sec:mixin-trees},
Appendix~\ref{app:well-definedness})
computes $\mathrm{lfp}(T_P)$: each query into
$\mathrm{path}$ triggers further evaluation through
the recursive body, and the trie union ($\supers$)
accumulates all reachable tuples.
This corresponds exactly to \emph{Naive
evaluation}~\cite{bancilhon1986-recursive-query-strategies} of the
$T_P$
operator~\cite{vanemden1976-predicate-logic-semantics}:
each iteration applies every rule to the current
fact set and adds newly derived facts.
Because Datalog-encoded relations have finitely many
constants per query, the per-query answer domain is
finite, and convergence is guaranteed by the ascending
chain condition on a finite lattice.
\emph{Semi-Naive} evaluation~\cite{bancilhon1986-recursive-query-strategies},
which restricts each iteration to rules that use at
least one newly derived fact, is a natural
optimisation left for future work.

\paragraph{The recursive rule as monadic pseudocode}
The recursive rule above roughly corresponds to the
following pseudocode (Haskell-like syntax with
  \texttt{RebindableSyntax} and \texttt{OverloadedLists}
  semantics, where do-notation desugars to a
\texttt{Set}-based bind and list literals denote sets):
\begin{lstlisting}[style=haskell]
{-# LANGUAGE RebindableSyntax, OverloadedLists #-}
s >>= f = unions [f x | x <- s]
return x = [x]
guard True  = [()]
guard False = []
edge = [("a","b"), ("b","c"), ("c","a")]
path = edge <> do
  (x, y0) <- edge
  (y1, z) <- path
  guard (y0 == y1)
  return (x, z)
\end{lstlisting}
\noindent
Each line of the pseudocode has a direct counterpart in
the inheritance-calculus encoding
(Section~\ref{sec:encoding-variables}):
\begin{description}
  \item[\texttt{edge}] corresponds to the base-fact trie;
  \item[\texttt{path = edge <> do \ldots}]
    corresponds to $\mathrm{path}$ inheriting $\mathrm{edge}$
    (base facts) and the recursive body (derived facts via
    $\supers$);
  \item[\texttt{(x, y0) <- edge}]
    corresponds to Level~1 ($\mathrm{xAcc}$), iterating $X$
    from $\mathrm{edge}$;
  \item[\texttt{(y1, z) <- path}]
    corresponds to Level~2 ($\mathrm{yAcc}$), joining on $Y$
    with $\mathrm{path}$, and Level~4 ($\mathrm{zAcc}$),
    iterating $Z$;
  \item[\texttt{guard (y0 == y1)}]
    corresponds to the $\mathrm{WithVisitorTail} \cap
    \mathrm{WithVisiteeTail}$ intersection in Level~2: a
    constant produces a $\mathrm{Visitor}$ entry only if it
    appears in both relations;
  \item[\texttt{return (x, z)}]
    corresponds to the $\mathrm{Replacement}$ chain in
    Level~5, constructing the output tuple $(X, Z)$.
\end{description}
The only difference is the evaluation strategy:
Haskell evaluates \texttt{path} lazily (diverges),
whereas inheritance-calculus evaluates it via tabled
$T_P$ iteration (converges to $\mathrm{lfp}$).

However, this monadic program does not converge because
\texttt{path} is defined recursively.%
\footnote{
  Enabling the \texttt{RecursiveDo} extension does not help:
  \texttt{mfix} ties a knot at the \emph{element} level via laziness,
  whereas the transitive closure equation requires a
  fixed point at the \emph{container} level, that is, finding a
  set~$S$ such that $S = F(S)$.
}
Inheritance-calculus solves this through tabled evaluation
(Section~\ref{sec:mixin-trees}), which computes the
least fixed point of the $T_P$ operator by iterating
over the lattice of mixin trees, precisely the
container-level fixed point that the pseudocode above cannot express.

\subsection{Semantic Correspondence with Standard Datalog}
\label{sec:datalog-semantics}

The encoding produces standard Datalog tuple-level semantics.
Putting the previous subsections together: tries provide indexed
tuple storage, visitor dispatch supplies the control skeleton,
and per-constant $\mathrm{Replacement}$ keeps output
construction local to each witness.
The key mechanism is the per-constant $\mathrm{Visited}$
pattern: output construction ($\mathrm{Replacement}$)
occurs inside each matched constant's $\mathrm{Visitor}$
scope, operating on the \emph{individual} constant's data
rather than on the merged scope.
The $\mathrm{Accepted}$ property then collects all per-constant
$\mathrm{Visited}$ results via trie union ($\supers$).

For example, given $\mathrm{edge} = \{(a,b),\; (b,c),\; (c,a)\}$
(a~3-cycle), the two-hop self-join
$\mathrm{twoHop}(X, Z) \mathrel{{:}\!{-}} \mathrm{edge}(X, Y),\; \mathrm{edge}(Y, Z)$
produces exactly $3$ pairs:
$\{(a,c),\; (b,a),\; (c,b)\}$, matching standard Datalog.
This is the behavior exhibited by
\anon[\nolinkurl{RelationalTwoHopRepeatedVisitor.mixin.yaml}~(supplementary material)]{\nolinkurl{RelationalTwoHopRepeatedVisitor.mixin.yaml}~\cite{mixin2025}}.
When the $Y$-join matches constant $b$ (present in both
  $\mathrm{edge}$'s second column from $(a,b)$ and
$\mathrm{edge}$'s first column), the $\mathrm{Visitor}$ for
$b$ receives $\mathrm{VisitorTail} = \mathrm{edge}.\mathrm{TailMap}.b$
(containing only~$c$), and $\mathrm{Replacement}$ constructs
a single output tuple $(a, c)$, rather than all of $\{a,b,c\}$.

This value-level precision arises because each constant's
$\mathrm{Replacement}$ operates on $\mathrm{TailMap}.d$
(the per-constant sub-trie) rather than on the merged
trie.
The same point is visible in
the repeated-visitor self-join regression example\anon[~(supplementary material)]{~\cite{mixin2025}},
where repeated visitor dispatch prevents witnesses from
different constants from being merged spuriously.
Although $\bases$ (equation~\ref{eq:bases}) still computes
cross-joins at the primitive level, the encoding routes
computation through per-constant $\mathrm{TailMap}$ entries
before any merging occurs, achieving the same effect as
standard Datalog's tuple-level evaluation.

\section{Evaluation Trace of Multi-Target \texorpdfstring{$\this$}{this} Resolution}
\label{app:evaluation-trace}

This appendix gives a step-by-step evaluation trace of the program of
Appendix~\ref{app:scala-multi-target}, computing
$\properties((\text{``HasMultipleOuters''},\, \text{``outer''}))$ by the six
functions of Section~\ref{sec:definitions}
(equations~\ref{eq:properties} to~\ref{eq:this}).
Every call to $\overrides$, $\bases$, $\resolve$, $\this$, and $\supers$
is shown with its arguments and returned set.
The trace lives in the metalanguage: every argument and result is a
\emph{path}, that is, a sequence of labels
$(\ell_1, \ldots, \ell_n)$ identifying a position in the tree
(Section~\ref{sec:definitions}), not an inheritance-calculus expression.
To keep labels distinct from the metalanguage variables of
Section~\ref{sec:definitions} (the italic $p$, $\ell$, $n$, $S$), we
quote each concrete label, so a path is written as a parenthesized,
comma-separated sequence of quoted labels.
Thus $(\text{``HasMultipleOuters''}, \text{``outer''})$ is the path that
the inheritance-calculus projection
$\mathrm{HasMultipleOuters}.\mathrm{outer}$ identifies, and the root path
is the empty sequence $()$.

\paragraph{Program}
\begin{align*}
  \{ \quad \mathrm{MyOuter} &\mapsto \{
      \mathrm{MyInner} \mapsto \{
    \mathrm{outer} \mapsto \{\qualifiedthis{\mathrm{MyOuter}}\}\}\}, \\
    \mathrm{Object1} &\mapsto \{\mathrm{MyOuter}\}, \\
    \mathrm{Object2} &\mapsto \{\mathrm{MyOuter}\}, \\
    \mathrm{HasMultipleOuters} &\mapsto \{
      \mathrm{Object1}.\mathrm{MyInner},\;
    \mathrm{Object2}.\mathrm{MyInner}\}
  \quad \}
\end{align*}

\paragraph{AST primitives}
Parsing populates $\defines$ and $\inherits$
(Section~\ref{sec:definitions}):
\begin{align*}
  \defines(()) &= \{
    \text{``MyOuter''},\;
    \text{``Object1''},\;
    \text{``Object2''},\;
  \text{``HasMultipleOuters''}\}, \\
  \defines((\text{``MyOuter''})) &= \{\text{``MyInner''}\}, \\
  \defines((\text{``MyOuter''}, \text{``MyInner''})) &= \{\text{``outer''}\}, \\
  \inherits((\text{``MyOuter''}, \text{``MyInner''}, \text{``outer''})) &=
  \{(1,\, ())\}, \\
  \inherits((\text{``Object1''})) = \inherits((\text{``Object2''})) &=
  \{(0,\, (\text{``MyOuter''}))\}, \\
  \inherits((\text{``HasMultipleOuters''})) &=
  \{(0,\, (\text{``Object1''}, \text{``MyInner''})),\;
  (0,\, (\text{``Object2''}, \text{``MyInner''}))\}.
\end{align*}
All other $\defines$ and $\inherits$ are empty.
Each $\inherits$ entry is a reference pair $(n,\, w)$,
where $n$ is the number of upward steps and $w$ the
downward projection list (Section~\ref{sec:definitions}).

\paragraph{Trace}
We present the trace as a numbered list of semantic-function calls in
dependency order: a call appears after the calls it depends on, and each
item states, in English, which earlier items it uses, how it computes, and
what it yields. Each distinct call is listed once, so a single item may be
used by several later items, the reuse that an implementation could realize
by tabling~\cite{tamaki1986-tabled-resolution}.
The arguments and results are metalanguage values: a label is a quoted
string such as $\text{``outer''}$, and a path is a tuple of labels such as
$(\text{``HasMultipleOuters''}, \text{``outer''})$, with the root path
written $()$.
The query is
$\properties((\text{``HasMultipleOuters''}, \text{``outer''}))$.

\begin{enumerate}
  \item \label{trace:overrides-MyOuter-MyInner}Compute $\overrides((\text{``MyOuter''}, \text{``MyInner''}))$, the paths sharing the identity of $(\text{``MyOuter''}, \text{``MyInner''})$. It keeps $(\text{``MyOuter''}, \text{``MyInner''})$ and adds any same-label definition reached through the supers of its enclosing scope. Result: $\{(\text{``MyOuter''}, \text{``MyInner''})\}$.
  \item \label{trace:supers-MyOuter-MyInner}Compute $\supers((\text{``MyOuter''}, \text{``MyInner''}))$ using item~\ref{trace:overrides-MyOuter-MyInner}. It takes the reflexive-transitive closure of $\bases$ from $(\text{``MyOuter''}, \text{``MyInner''})$, pairing each reachable override with the inheritance site through which it is reached. Result: $\{((\text{``MyOuter''}), (\text{``MyOuter''}, \text{``MyInner''}))\}$.
  \item \label{trace:overrides-Has-outer}Compute $\overrides((\text{``HasMultipleOuters''}, \text{``outer''}))$, the paths sharing the identity of $(\text{``HasMultipleOuters''}, \text{``outer''})$ using item~\ref{trace:supers-MyOuter-MyInner}. It keeps $(\text{``HasMultipleOuters''}, \text{``outer''})$ and adds any same-label definition reached through the supers of its enclosing scope. Result: $\{(\text{``HasMultipleOuters''}, \text{``outer''}), (\text{``MyOuter''}, \text{``MyInner''}, \text{``outer''})\}$.
  \item \label{trace:bases-Has-outer}Compute $\bases((\text{``HasMultipleOuters''}, \text{``outer''}))$. $(\text{``HasMultipleOuters''}, \text{``outer''})$ carries no inheritance reference ($\inherits = \varnothing$), so it has no direct base. Result: $\varnothing$.
  \item \label{trace:bases-MyOuter-MyInner-outer}Compute $\bases((\text{``MyOuter''}, \text{``MyInner''}, \text{``outer''}))$. $(\text{``MyOuter''}, \text{``MyInner''}, \text{``outer''})$ carries the reference pair $(1, ())$ (take one step upward, then project the empty projection), whose target is obtained by item~\ref{trace:resolve-Has-MyOuter-MyInner-1}.
  \item \label{trace:this-Has-MyOuter-MyInner-1}Compute one $\this$ step: from the frontier $\{(\text{``HasMultipleOuters''})\}$, take the supers of each frontier path and keep those whose definition site is $(\text{``MyOuter''}, \text{``MyInner''})$, collecting their inheritance sites as the new frontier. Result: $\{(\text{``Object1''}), (\text{``Object2''})\}$.
  \item \label{trace:resolve-Has-MyOuter-MyInner-1}Compute $\resolve$ for the reference pair $(1, ())$ defined at $(\text{``MyOuter''}, \text{``MyInner''})$, reached from inheritance site $(\text{``HasMultipleOuters''})$ using item~\ref{trace:this-Has-MyOuter-MyInner-1}. It takes one step upward through $\this$ and then appends the projection. Result: $\{(\text{``Object1''}), (\text{``Object2''})\}$.
  \item \label{trace:resolve-0-MyOuter}Compute $\resolve$ for the reference pair $(0, (\text{``MyOuter''}))$ defined at $()$, reached from inheritance site $()$. It takes no step upward through $\this$ and then appends the projection. Result: $\{(\text{``MyOuter''})\}$.
  \item \label{trace:supers-Has-outer}Compute $\supers((\text{``HasMultipleOuters''}, \text{``outer''}))$ using items~\ref{trace:overrides-Has-outer}, \ref{trace:bases-Has-outer}, \ref{trace:bases-MyOuter-MyInner-outer}, \ref{trace:resolve-Has-MyOuter-MyInner-1} and~\ref{trace:resolve-0-MyOuter}. It takes the reflexive-transitive closure of $\bases$ from $(\text{``HasMultipleOuters''}, \text{``outer''})$, pairing each reachable override with the inheritance site through which it is reached. Result: $\{((), (\text{``MyOuter''})), ((), (\text{``Object1''})), ((), (\text{``Object2''})), ((\text{``HasMultipleOuters''}), (\text{``HasMultipleOuters''}, \text{``outer''})), ((\text{``HasMultipleOuters''}), (\text{``MyOuter''}, \text{``MyInner''}, \text{``outer''}))\}$.
  \item \label{trace:properties}Compute $\properties((\text{``HasMultipleOuters''}, \text{``outer''}))$ using item~\ref{trace:supers-Has-outer}. It unions the $\defines$ of every override path in that $\supers$ set. Only $(\text{``MyOuter''})$ defines a label, namely $\text{``MyInner''}$, reached through both $(\text{``Object1''})$ and $(\text{``Object2''})$; set union collapses the two routes. Result: $\{\text{``MyInner''}\}$.
\end{enumerate}


\begin{remark}
  The crux is the $\this$ step (item~\ref{trace:this-Has-MyOuter-MyInner-1}).
  The inheritance reference at
  $(\text{``MyOuter''}, \text{``MyInner''}, \text{``outer''})$ takes one step
  upward, so $\this$ moves one step up from
  $(\text{``MyOuter''}, \text{``MyInner''})$ and yields the frontier
  $\{(\text{``Object1''}), (\text{``Object2''})\}$, the two inheritance
  sites that incorporated $\text{``MyInner''}$. This is precisely the
  multi-target situation that Scala~3 and the NixOS module system reject
  (Appendix~\ref{app:scala-multi-target}): there $\this$ would have to
  select a single site. Reaching the properties of $(\text{``MyOuter''})$
  then happens because both $(\text{``Object1''})$ and
  $(\text{``Object2''})$ have $(\text{``MyOuter''})$ as a base, so the final
  $\supers$ contains $(\text{``MyOuter''})$ through both routes; since
  attribute resolution aggregates by set union, the two routes collapse to
  the single answer
  $\properties((\text{``HasMultipleOuters''}, \text{``outer''}))
  = \{\text{``MyInner''}\} = \properties((\text{``MyOuter''}))$.
\end{remark}

\fi % end \ifInheritanceAppendix

% Appendix-only build (supplement.tex): no main matter ran above, so emit the
% bibliography here, after the appendix.
\ifInheritanceBody\else
\bibliographystyle{ACM-Reference-Format}
\bibliography{references}
\fi

% \clearpage before closing CJK* ships every remaining page and float while the
% CJK active characters still mean CJK (the running head and any deferred float
% are typeset in the output routine, after \end{CJK*} would otherwise restore the
% default UTF-8 meaning).
\ifInheritanceChinese\clearpage\end{CJK*}\fi
\end{document}
.  The
% defaults include everything, so the full build (preprint.tex) sets
% nothing; submission.tex clears the appendices (POPL 2027 requires them
% as separate supplemental material), supplement.tex clears the body.
% The single bibliography is emitted exactly once, after the main matter and
% before the appendix, since both cite it. In the appendix-only build (no main
% matter) it instead follows the appendix; see the two guarded copies below.
\makeatletter
\@ifundefined{ifInheritanceBody}{\newif\ifInheritanceBody\InheritanceBodytrue}{}
\@ifundefined{ifInheritanceAppendix}{\newif\ifInheritanceAppendix\InheritanceAppendixtrue}{}
% Language switch for the reader-facing prose. Defaults to false (English), so the
% existing pdfLaTeX entries are unaffected; preprint-zh.tex sets it true and adds
% CJK support, still under pdfLaTeX. See the paper-writing skill (Bilingual
% Translation) for the convention.
\@ifundefined{ifInheritanceChinese}{\newif\ifInheritanceChinese\InheritanceChinesefalse}{}
% acmart writes \newlabel{tocindent<n>}{<dimen>} into the .aux to align the
% table of contents. Those bare-dimen labels break hyperref's \newlabel when
% another build imports this .aux via xr-hyper; the importing entries
% (submission.tex, supplement.tex) drop them at import time through
% xr-tocindent-filter.tex, so nothing is suppressed here at the source.
\makeatother

% --- bilingual reader-facing prose ----------------------------------
% The reader-facing body exists in two languages in this one file, selected by
% \ifInheritanceChinese: English (default) and a Chinese translation. Both build
% with pdfLaTeX, so the math renders identically. Block prose is guarded by
% \ifInheritanceChinese <chinese> \else <english> \fi; inline switches (title,
% captions, theorem notes interleaved with math) use \bilingual{<en>}{<zh>}.
% Author-side notes stay Chinese % comments in both builds.
\newcommand{\bilingual}[2]{#1}
\ifInheritanceChinese
  % CJKutf8 + the Arphic gbsn font render Chinese under pdfLaTeX (see
  % modules/texlive.nix). The document is wrapped in a CJK* environment opened
  % after \begin{document}; ASCII passes through, so English and math render as in
  % the English build.
  \usepackage{CJKutf8}
  \setlength{\emergencystretch}{2em}
  % Display the Chinese, but feed the English to hyperref's PDF strings, which
  % cannot encode the CJK active characters.
  \renewcommand{\bilingual}[2]{\texorpdfstring{#2}{#1}}
\fi

\begin{document}
% Open CJK* before \title so the Chinese in the title, abstract, and body is read
% inside it; \proofname and the theorem-environment names are localized here
% (their bytes must sit inside a CJK environment, and translating \thmname's
% argument leaves the shared counter and numbering identical to English). Closed
% before \end{document}.
\ifInheritanceChinese
  \begin{CJK*}{UTF8}{gbsn}
  \renewcommand{\proofname}{证明}
  \makeatletter
  \newcommand{\inheritanceThmName}[1]{\@ifundefined{inheritance@thm@#1}{#1}{\csname inheritance@thm@#1\endcsname}}
  \renewcommand{\thmname}[1]{\inheritanceThmName{#1}}
  \@namedef{inheritance@thm@Theorem}{定理}
  \@namedef{inheritance@thm@Lemma}{引理}
  \@namedef{inheritance@thm@Corollary}{推论}
  \@namedef{inheritance@thm@Proposition}{命题}
  \@namedef{inheritance@thm@Conjecture}{猜想}
  \@namedef{inheritance@thm@Definition}{定义}
  \@namedef{inheritance@thm@Example}{例}
  \@namedef{inheritance@thm@Remark}{注}
  \@namedef{inheritance@thm@Notation}{记法}
  \makeatother
\fi

% The optional short title feeds the running head and the PDF metadata; keep it
% ASCII English in both builds, since those are typeset outside CJK*.
\title[A Calculus of Inheritance]{\bilingual{A Calculus of Inheritance}{一种继承的演算}}
% The appendix-only build (supplement.tex) is distributed as a standalone
% artifact; mark it so it is not mistaken for the main paper, which carries
% no subtitle.
\ifInheritanceBody\else
  \subtitle{Supplementary Material}
\fi

\author{Bo Yang}
\affiliation{%
  \institution{Figure AI Inc.}
  \city{San Jose}
  \state{California}
  \country{USA}
}
\email{yang-bo@yang-bo.com}
\thanks{This work was conducted independently prior to the author's employment at Figure AI.}

% The abstract, CCS concepts, and keywords are main-matter front matter: they
% belong to the paper, not to the standalone appendix PDF (supplement.tex),
% which sets \ifInheritanceBody false. Without this guard \maketitle typesets
% them in the appendix-only build.
\ifInheritanceBody
\begin{abstract}
  Just as the $\lambda$-calculus uses three primitives (abstraction, application, variable) as the foundation of functional programming, inheritance-calculus uses three primitives (record, definition, inheritance) as the foundation of declarative programming. A record is a set of definitions; by unifying modules, classes, objects, methods, fields, and locals under this single record abstraction, the calculus models inheritance simply as set union of definitions, two definitions of the same name merging recursively rather than one overriding the other. A merge is always meaningful because every value is a record and observation is purely structural: the merged record simply answers queries with the union of what its parts provide, so there is no scalar leaf at which two values could conflict. Consequently, composition is inherently commutative, idempotent, and associative, structurally eliminating the multiple-inheritance linearization problem. Its semantics is first-order, denotational, and computable by tabling (memoized fixpoint iteration as in tabled logic programming), even for cyclic inheritance hierarchies. This first-order, denotational, tabled semantics extends to the $\lambda$-calculus: $\lambda$-terms translate into a sublanguage of a lazy approximation of inheritance-calculus, the approximation that treats cyclic self-reference as divergent instead of iterating it to a fixpoint, and the translation is fully abstract, equating exactly the $\lambda$-terms with equal B\"ohm trees.

  Inheritance-calculus is distilled from MIXINv2, a practical implementation in which the same code can be invoked under different calling conventions (function colors such as synchronous and asynchronous); ordinary arithmetic, encoded within the calculus itself, behaves as multi-directional relations in the sense of logic programming; $\mathtt{this}$ resolves to multiple enclosing targets at once rather than a single receiver; and programs are extensible in both dimensions of the Expression Problem, gaining new variants and new operations by adding code. This makes inheritance-calculus strictly more expressive than the $\lambda$-calculus in Felleisen's sense of macro-expressibility.
\end{abstract}

\begin{CCSXML}
  <ccs2012>
  <concept>
  <concept_id>10003752.10010124.10010131</concept_id>
  <concept_desc>Theory of computation~Program semantics</concept_desc>
  <concept_significance>500</concept_significance>
  </concept>
  <concept>
  <concept_id>10003752.10010124.10010125.10010128</concept_id>
  <concept_desc>Theory of computation~Object oriented constructs</concept_desc>
  <concept_significance>500</concept_significance>
  </concept>
  <concept>
  <concept_id>10003752.10010124.10010125.10010127</concept_id>
  <concept_desc>Theory of computation~Functional constructs</concept_desc>
  <concept_significance>300</concept_significance>
  </concept>
  <concept>
  <concept_id>10003752.10010124.10010125.10010129</concept_id>
  <concept_desc>Theory of computation~Program schemes</concept_desc>
  <concept_significance>100</concept_significance>
  </concept>
  <concept>
  <concept_id>10011007.10011006.10011039</concept_id>
  <concept_desc>Software and its engineering~Formal language definitions</concept_desc>
  <concept_significance>500</concept_significance>
  </concept>
  <concept>
  <concept_id>10011007.10011006.10011008.10011009.10011019</concept_id>
  <concept_desc>Software and its engineering~Extensible languages</concept_desc>
  <concept_significance>500</concept_significance>
  </concept>
  <concept>
  <concept_id>10011007.10011006.10011008.10011009.10011011</concept_id>
  <concept_desc>Software and its engineering~Object oriented languages</concept_desc>
  <concept_significance>100</concept_significance>
  </concept>
  </ccs2012>
\end{CCSXML}

\ccsdesc[500]{Theory of computation~Program semantics}
\ccsdesc[500]{Theory of computation~Object oriented constructs}
\ccsdesc[300]{Theory of computation~Functional constructs}
\ccsdesc[100]{Theory of computation~Program schemes}
\ccsdesc[500]{Software and its engineering~Formal language definitions}
\ccsdesc[500]{Software and its engineering~Extensible languages}
\ccsdesc[100]{Software and its engineering~Object oriented languages}

\keywords{declarative programming, inheritance, fixpoint semantics,
  self-referential records, expression problem,
  first-order semantics,
  \texorpdfstring{$\lambda$-calculus}{lambda-calculus},
\texorpdfstring{B\"ohm trees}{Bohm trees}, configuration languages}
\fi % end \ifInheritanceBody (abstract, CCS concepts, keywords are main-matter only)

\maketitle

\ifInheritanceBody
\section{\bilingual{Introduction}{引言}}
\label{sec:introduction}

\ifInheritanceChinese
诸如 Jsonnet~\cite{jsonnet}、Hydra~\cite{hydra2023}、CUE~\cite{cue2019}、Dhall~\cite{dhall2017}、Kustomize~\cite{kustomize} 与 JSON Patch~\cite{rfc6902-json-patch} 这样的配置语言,都提供了通过继承或合并来组合结构化数据的机制。

在这些之中,Nix 生态系统~\cite{dolstra2010-nixos} 格外突出,无论是就其规模还是其设计的走向而言。它的软件包集合 \texttt{nixpkgs} 收录了超过 113{,}000 个软件包,%
\footnote{
  据 Repology~\cite{repology2025},这接近 Debian 或 Ubuntu 软件包数量的三倍,而贡献者数量相近。并非每个软件包都用到这里所描述的可扩展性机制,但同样的告诫也适用于其他发行版。
}
是现存最大的软件包集合之一。软件包的组合最初依赖一条线性化的 mixin 链,其中每一层都能看见它的前一层,%
\footnote{
  Nix overlay 通过一条线性化的 \texttt{self}/\texttt{super} 链来组合软件包集合,遵循 Bracha--Cook 的 mixin 模式~\cite{bracha1990-mixin-inheritance}:\texttt{self} 晚绑定到不动点结果,而 \texttt{super} 指向前一层。该机制是不对称且依赖顺序的。
}
但该生态系统已日益汇聚到一种对称的替代方案上:NixOS 模块系统~\cite{dolstra2010-nixos,nixosmodules},它通过对嵌套的属性集进行递归合并、配以自引用求值与延迟模块~\cite{nixosmodules} 来组合,而这种组合是可交换、幂等且可结合的。模块系统最初为系统配置而设计,如今已被用于用户环境~\cite{homemanager}、macOS 配置~\cite{nixdarwin}、磁盘分区~\cite{disko}、flake 结构~\cite{flakeparts}、开发者环境~\cite{devenv}、Kubernetes 集群~\cite{kubenix,nixidy},以及多语言软件包构建~\cite{dream2nix},而最后这一项正是 \texttt{nixpkgs} 那套不对称机制原本的领域。Nix 是一门带头等函数的函数式语言;然而一份典型的模块配置几乎完全由嵌套的属性集构成,函数只是偶尔作为模块系统的 DSL 语法出现,例如 \texttt{mkIf} 与 \texttt{mkMerge}。没有任何计算理论刻画了这一机制所展现的可扩展性。
\else
Configuration languages such as Jsonnet~\cite{jsonnet}, Hydra~\cite{hydra2023},
CUE~\cite{cue2019}, Dhall~\cite{dhall2017},
Kustomize~\cite{kustomize}, and JSON
Patch~\cite{rfc6902-json-patch} all provide mechanisms for
composing structured data through inheritance or merging.

Among these, the Nix ecosystem~\cite{dolstra2010-nixos} stands out,
both for scale and for the trajectory of its design.
Its package collection, \texttt{nixpkgs}, contains over
113{,}000 packages,%
\footnote{
  According to Repology~\cite{repology2025},
  nearly three times the package count of Debian or Ubuntu,
  achieved with a similar number of contributors.
  Not every package exercises the extensibility mechanisms
  described here, but the same caveat applies to other distributions.
}
one of the largest in existence.
Package composition originally relied on a linearized mixin
chain in which each layer sees its predecessor,%
\footnote{
  Nix overlays compose package sets via a linearized
  \texttt{self}/\texttt{super} chain following the
  Bracha--Cook mixin pattern~\cite{bracha1990-mixin-inheritance}:
  \texttt{self} is late-bound to the fixed-point result
  and \texttt{super} refers to the previous layer.
  The mechanism is asymmetric and order-dependent.
}
but the ecosystem has increasingly converged on a symmetric
alternative: the NixOS module
system~\cite{dolstra2010-nixos,nixosmodules}, which composes by
recursively merging nested attribute sets with self-referential
evaluation and deferred modules~\cite{nixosmodules}---commutative,
idempotent,
and associative.
Originally designed for system configuration, the module system
has been adopted for user environments~\cite{homemanager},
macOS configuration~\cite{nixdarwin},
disk partitioning~\cite{disko},
flake structure~\cite{flakeparts},
developer environments~\cite{devenv},
Kubernetes clusters~\cite{kubenix,nixidy},
and multi-language package building~\cite{dream2nix}---the
original domain of \texttt{nixpkgs}'s asymmetric mechanism.
Nix is a functional language with first-class functions;
yet a typical module configuration consists almost entirely of
nested attribute sets, with functions appearing only
occasionally as DSL syntax of the module system
such as \texttt{mkIf} and \texttt{mkMerge}.
No computational theory captures the extensibility that this
mechanism exhibits.
\fi

The $\lambda$-calculus has no analogue for declarative
programming. This gap matters because the two paradigms differ in
fundamental ways, which raises a question:

\begin{quote}
  Can configuration alone be practical general-purpose programming?
\end{quote}

\noindent
To make the question precise, we operationalize it with four
\emph{opinionated} proxies, grouped under the question's two halves:
\emph{configuration} and \emph{practical general-purpose} programming.

\begin{itemize}
  \item Immutable and first-order%
    \footnote{
      Here \emph{first-order} holds in three senses at once:
      the semantic domain contains only data, no function spaces,
      closures, or
      strategies~\cite{vanemden1976-predicate-logic-semantics};
      the defining equations map data to
      data~\cite{reynolds1972-definitional-interpreters};
      and the semantics is a first-order inductive
      definition~\cite{aczel1977-inductive-definitions} whose
      quantifiers range only over data.
    }
    is our proxy for
    \emph{configuration}: a configuration is read, not a function that
    reads an argument.
  \item Turing-completeness together with extensibility is our proxy for
    \emph{practical general-purpose} programming: a language must be able
    to compute anything and to grow by adding code rather than rewriting
    it, the Expression
    Problem~\cite{wadler1998-expression-problem}.
\end{itemize}

\noindent
Each classical model satisfies some of these proxies but fails at least
one:

\begin{table}[h]
  \centering
  \small
  \begin{tabular}{lcccl}
    \textbf{Model} & \textbf{Immutable} & \textbf{First-order} & \textbf{Turing complete} & \textbf{Extensibility} \\
    \hline
    Turing Machine        & $\times$     & $\checkmark$ & $\checkmark$ & $\times$ \\
    $\lambda$-calculus    & $\checkmark$ & $\times$     & $\checkmark$ & opt-in \\
    RAM Machine           & $\times$     & $\checkmark$ & $\checkmark$ & $\times$ \\
    Prolog / Horn clauses & $\checkmark$ & $\checkmark$ & $\checkmark$ & disjunct-only \\
  \end{tabular}
  \caption{The question operationalized as four proxies: immutable and
    first-order (configuration) versus Turing-complete and extensible
  (practical general-purpose). No classical model satisfies all four.}
  \label{tab:computational-models}
\end{table}

\noindent
The Turing Machine and RAM Machine require mutable state and have no native
extension mechanism;
the $\lambda$-calculus is higher-order rather than first-order, and its
extensibility is opt-in: a function is extensible only where its author
anticipated extension, for example by threading open-recursion
parameters, while a closed function body cannot be extended without
rewriting it;
Prolog meets the other three proxies but is extensible only by adding
disjuncts (alternative clauses) to existing predicates, not by attaching
new fields or methods to existing values.
No known computational model satisfies all four proxies simultaneously.

The Nix ecosystem already suggests that such a model may exist:
its inheritance-based composition over recursive records, without
explicit functions, is expressive enough to configure entire
operating systems and manage one of the largest package
collections in existence.
Most configuration languages operate on finite, well-founded
structures (initial algebras in the sense of universal algebra).
Nix is an exception: its values are lazily observable
and possibly infinite, allowing any single package or
configuration option to be evaluated without materializing the
entire set of hundreds of thousands of packages and options.
Guided by this observation, we set out to reduce the
NixOS module system to a minimal set of primitives.
The reduction produced MIXIN, an early implementation in Nix,
and subsequently
MIXINv2\anon[~(included as supplementary material)]{~\cite{mixin2025}}, now implemented in Python: a declarative programming
language with only three constructs (record, definition,
and inheritance) and no functions or let-bindings.
MIXINv2 obtains scalars such as floating-point numbers through a
foreign-function interface; inheritance-calculus is what remains after
removing that interface, so names like $\mathrm{PI}$ and
$\mathrm{Float}$ in Figure~\ref{fig:cheatsheet} denote imported
records whose internals the examples never inspect: scalar
computation such as the figure's additions happens behind the
interface, while the calculus itself only ever observes structure;
the case study of Section~\ref{sec:case-study} instead encodes
naturals and their arithmetic purely in the calculus, with no FFI.
When two merged records define the same name with different contents,
the result is simply a record carrying both structures; the calculus
never compares or rejects them
(Section~\ref{sec:mixin-trees} makes the merge precise).
No such collision occurs in Figure~\ref{fig:cheatsheet}: each
instance defines each scalar field exactly once, and the type record
it merges with contributes structure, not a competing value.

\begin{figure*}
  \centering
  \begin{minipage}[t]{0.34\textwidth}
    \textbf{Python}\par
\begin{lstlisting}[style=python]
# complex_module.py
from builtins import *
from math import *
from operator import *
from dataclasses import *
@dataclass(kw_only=True)
class Complex:
    re: float
    im: float
    def plus(self, *,
             other: "Complex"):
        sumRe = add(
            self.re,
            other.re)
        sumIm = add(
            self.im,
            other.im)
        return Complex(
            re=sumRe,
            im=sumIm)
a = Complex(
    re=pi,
    im=e)
b = Complex(
    re=tau,
    im=pi)
total = a.plus(
    other=b)
\end{lstlisting}
  \end{minipage}
  \hfill
  \begin{minipage}[t]{0.65\textwidth}
    \textbf{Inheritance-calculus}\par\smallskip
    {\small\(
        \figureAnnotationAnchor
        \begin{aligned}
          &\mathrm{ComplexModule} \mapsto {} \figureAnnotation{module} \\
          &\left\{
            \begin{aligned}
              &\mathrm{BuiltIns}, \figureAnnotation{import} \\
              &\mathrm{Math}, \figureAnnotation{import} \\
              &\mathrm{Operator}, \figureAnnotation{import} \\
              &\mathrm{Complex} \mapsto {} \figureAnnotation{class} \\
              &\left\{
                \begin{aligned}
                  &\mathrm{re} \mapsto {} \figureAnnotation{field} \\
                  &\left\{ \qualifiedthis{\mathrm{ComplexModule}}.\mathrm{Float} \right\}, \figureAnnotation{imported type} \\
                  &\mathrm{im} \mapsto {} \figureAnnotation{field} \\
                  &\left\{ \qualifiedthis{\mathrm{ComplexModule}}.\mathrm{Float} \right\}, \figureAnnotation{imported type} \\
                  &\mathrm{Plus} \mapsto {} \figureAnnotation{method} \\
                  &\left\{
                    \begin{aligned}
                      &\mathrm{other} \mapsto \left\{ \mathrm{Complex} \right\}, \figureAnnotation{parameter} \\
                      &\mathrm{\_sumReCall} \mapsto {} \figureAnnotation{function application} \\
                      &\left\{
                        \begin{aligned}
                          &\qualifiedthis{\mathrm{ComplexModule}}.\mathrm{Add}, \figureAnnotation{imported function} \\
                          &\mathrm{operand0} \mapsto \left\{ \qualifiedthis{\mathrm{Complex}}.\mathrm{re} \right\}, \figureAnnotation{keyword argument} \\
                          &\mathrm{operand1} \mapsto \left\{ \qualifiedthis{\mathrm{Plus}}.\mathrm{other}.\mathrm{re} \right\} \figureAnnotation{keyword argument} \\
                      \end{aligned}\right\}, \\
                      &\mathrm{\_sumRe} \mapsto \left\{ \mathrm{\_sumReCall}.\mathrm{result} \right\}, \figureAnnotation{local} \\
                      &\mathrm{\_sumImCall} \mapsto {} \figureAnnotation{function application} \\
                      &\left\{
                        \begin{aligned}
                          &\qualifiedthis{\mathrm{ComplexModule}}.\mathrm{Add}, \figureAnnotation{imported function} \\
                          &\mathrm{operand0} \mapsto \left\{ \qualifiedthis{\mathrm{Complex}}.\mathrm{im} \right\}, \figureAnnotation{keyword argument} \\
                          &\mathrm{operand1} \mapsto \left\{ \qualifiedthis{\mathrm{Plus}}.\mathrm{other}.\mathrm{im} \right\} \figureAnnotation{keyword argument} \\
                      \end{aligned}\right\}, \\
                      &\mathrm{\_sumIm} \mapsto \left\{ \mathrm{\_sumImCall}.\mathrm{result} \right\}, \figureAnnotation{local} \\
                      &\mathrm{result} \mapsto {} \figureAnnotation{return} \\
                      &\left\{
                        \begin{aligned}
                          &\mathrm{Complex}, \figureAnnotation{new instance} \\
                          &\mathrm{re} \mapsto \left\{ \mathrm{\_sumRe} \right\}, \figureAnnotation{keyword argument} \\
                          &\mathrm{im} \mapsto \left\{ \mathrm{\_sumIm} \right\} \figureAnnotation{keyword argument} \\
                      \end{aligned}\right\} \\
                  \end{aligned}\right\} \\
              \end{aligned}\right\}, \\
              &\mathrm{a} \mapsto {} \figureAnnotation{module-level variable} \\
              &\left\{
                \begin{aligned}
                  &\mathrm{Complex}, \figureAnnotation{new instance} \\
                  &\mathrm{re} \mapsto \left\{ \qualifiedthis{\mathrm{ComplexModule}}.\mathrm{PI} \right\}, \figureAnnotation{imported constant} \\
                  &\mathrm{im} \mapsto \left\{ \qualifiedthis{\mathrm{ComplexModule}}.\mathrm{E} \right\} \figureAnnotation{imported constant} \\
              \end{aligned}\right\}, \\
              &\mathrm{b} \mapsto {} \figureAnnotation{module-level variable} \\
              &\left\{
                \begin{aligned}
                  &\mathrm{Complex}, \figureAnnotation{new instance} \\
                  &\mathrm{re} \mapsto \left\{ \qualifiedthis{\mathrm{ComplexModule}}.\mathrm{Tau} \right\}, \figureAnnotation{imported constant} \\
                  &\mathrm{im} \mapsto \left\{ \qualifiedthis{\mathrm{ComplexModule}}.\mathrm{PI} \right\} \figureAnnotation{imported constant} \\
              \end{aligned}\right\}, \\
              &\mathrm{\_totalCall} \mapsto {} \figureAnnotation{method application} \\
              &\left\{
                \begin{aligned}
                  &\mathrm{a}.\mathrm{Plus}, \figureAnnotation{method reference} \\
                  &\mathrm{other} \mapsto \left\{ \mathrm{b} \right\} \figureAnnotation{keyword argument} \\
              \end{aligned}\right\}, \\
              &\mathrm{total} \mapsto \left\{ \mathrm{\_totalCall}.\mathrm{result} \right\} \figureAnnotation{module-level variable} \\
          \end{aligned}\right\}
        \end{aligned}
    \)}
  \end{minipage}
  \caption{A first look at inheritance-calculus (right), beside Python
    (left): the same complex-number class with an addition method.
    An entry $name \mapsto \{\dots\}$ is a definition; a bare name
    inside braces (such as $\mathrm{BuiltIns}$ or $\mathrm{Complex}$) is
    an inheritance clause that copies that record's definitions into the
    enclosing record.
    $\mathrm{BuiltIns}$, $\mathrm{Math}$, and $\mathrm{Operator}$ are
    defined in an enclosing scope, elided here.
    A bare name resolves lexically to the innermost enclosing record
    that defines it directly; qualified $\this$ reaches members that
    arrive by inheritance, such as the imported $\mathrm{Float}$ and
    $\mathrm{Add}$.
    Types and values share the same record syntax: a field such as
    $\mathrm{re}$ records its type by inheriting the type's record,
    with no checking force in the untyped calculus.
    Both sides, typing their components with \texttt{builtins.float},
    drawing their constants from the \texttt{math} module, and adding with
    \texttt{operator.add}, compute the complex
    number $(\pi+\tau) + (e+\pi)i$.
    $\mathrm{\_sumReCall}$ and $\mathrm{\_sumImCall}$ are function
    applications, and $\mathrm{\_totalCall}$ a method application, each
    encoded as a command pattern: a record that inherits the callee and
    defines its arguments, the call's result read back from a
    $\mathrm{result}$ field that the callee defines.
    Idempotence only means that listing the same inheritance source
    twice in one record is a no-op; distinct records inheriting the
    same source, like the two call sites here or the two instances
    $\mathrm{a}$ and $\mathrm{b}$, remain distinct.}
  \label{fig:cheatsheet}
\end{figure*}



In developing real programs in MIXINv2, including a database-backed web server whose Python foreign-function interface only faithfully wraps individual host-library calls, we found that module, class, object, method, field, and local binding can all be expressed as the same construct: a record of definitions composed by inheritance. This paper distills this minimal computational model into \textbf{inheritance-calculus}, shown beside Python in Figure~\ref{fig:cheatsheet}.
The key to reading the figure is that a reference
$\qualifiedthis{X}.w$ behaves like a late-bound Java-style
\texttt{Outer.this}: it denotes the enclosing record labeled $X$ as
seen from wherever the reference itself ends up in the fully merged
tree.
Inheriting a record does not duplicate it; it makes the record's
definitions, references included, visible at the inheriting position
as well, and a reference resolves relative to the position from which
it is viewed.
The set-union picture and this late-binding picture are the same
mechanism seen globally and locally: the union says which definitions
a position carries, and late binding says how a definition behaves at
each position that carries it.
Qualified references therefore see members that arrive by merging,
such as the imported $\mathrm{Float}$, while bare names such as
$\mathrm{Complex}$ resolve lexically against directly written
definitions only.
This late binding is what makes inheritance act as function
application: the call record $\mathrm{\_totalCall}$ inherits
$\mathrm{Plus}$, so the merged tree presents
$\mathrm{Plus}$'s body inside $\mathrm{\_totalCall}$ as well, and
viewed from there the reference
$\qualifiedthis{\mathrm{Plus}}.\mathrm{other}$
denotes that copy's enclosing record, $\mathrm{\_totalCall}$ itself,
where the argument $\mathrm{other}$ is defined.
Receivers bind the same way: the dotted bare name
$\mathrm{a}.\mathrm{Plus}$ resolves its first label lexically and
projects the rest, denoting the copy of $\mathrm{Plus}$ inside the
instance $\mathrm{a}$, in which
$\qualifiedthis{\mathrm{Complex}}.\mathrm{re}$ denotes
$\mathrm{a}$'s own field.
Section~\ref{sec:syntax} gives its formal syntax, and
Section~\ref{sec:mixin-trees} expands this first look into the full
semantics.
The linearization problem of mixin-based
systems~\cite{bracha1990-mixin-inheritance,barrett1996-c3-linearization}
did not arise: inheritance is inherently commutative, idempotent,
and associative, and $\this$ naturally resolves to
\emph{multiple targets} rather than one: when inheritance places the
same definition inside several enclosing records, a reference through
$\this$ inherits from all of them at once, their contents merged,%
\footnote{
  Resolution is anchored at the position from which the reference is
  observed. When both $\mathrm{A}$ and $\mathrm{B}$ inherit
  $\mathrm{Base}$, a reference in $\mathrm{Base}$ observed inside
  $\mathrm{A}$ sees only $\mathrm{A}$; the two call sites of
  Figure~\ref{fig:cheatsheet} therefore stay independent. Multiple
  targets arise when a single observation position inherits the same
  definition through two routes, say $\mathrm{C}$ inheriting
  $\mathrm{Base}$ via both $\mathrm{A}$ and $\mathrm{B}$: observed at
  $\mathrm{C}$, a $\this$ reference in $\mathrm{Base}$ resolves to
  $\mathrm{A}$ and $\mathrm{B}$ at once, and projection through it
  yields the union of what they provide.
}
where
prior systems
such as Scala and the NixOS module system reject the multi-target
situation.
Section~\ref{sec:case-study} traces how boolean logic and
arithmetic, encoded purely in the calculus as separate files, extend in both
dimensions of the Expression
Problem~\cite{wadler1998-expression-problem} without any dedicated
mechanism, and how ordinary arithmetic yields the relational
semantics of logic
programming~\cite{vanemden1976-predicate-logic-semantics}.
Section~\ref{sec:emergent-phenomena} discusses these and further
emergent phenomena, including function color
blindness~\cite{nystrom2015-function-color}.
Together these sections answer the four proxies: inheritance-calculus
is immutable and first-order by construction
(Section~\ref{sec:definitions}), Turing-complete because the
$\lambda$-calculus embeds into it
(Section~\ref{sec:forward-translation}), and extensible in both
dimensions of the Expression Problem (Section~\ref{sec:case-study}).

\begin{figure*}
  \centering
  \begin{tikzpicture}[
      calc/.style={align=center, inner sep=2pt, font=\small},
      more/.style={-{Stealth[length=2mm]}, semithick},
      iso/.style={{Stealth[length=2mm]}-{Stealth[length=2mm]}, semithick},
      elbl/.style={font=\scriptsize\itshape, midway, fill=white, inner sep=1pt},
      x=1mm, y=1mm,
    ]
    % rows: bottom = lambda calculi, middle = lambda sublanguages, top = inheritance-calculi
    % columns: eager (x=0), lazy (x=42), fixpoint (x=84)
    \node[calc] (lam)  at (0,0)   {$\lambda$};
    \node[calc] (llam) at (42,0)  {lazy $\lambda$};
    \node[calc] (flam) at (84,0)  {fixpoint $\lambda$~\cite{yang2026-co-lambda}};

    \node[calc] (lsub) at (42,19)  {$\lambda$-lazy-IC (ours)};
    \node[calc] (sub)  at (84,19)  {$\lambda$-IC};

    \node[calc] (lic)  at (42,38)  {lazy IC (ours)};
    \node[calc] (ic)   at (84,38)  {IC (ours)};

    % horizontal edges (convergence): "more defined"
    \draw[more] (lam)  -- node[elbl,below]{more defined (non-strict)} (llam);
    \draw[more] (llam) -- node[elbl,below]{more defined (unproductive $\bot$)} (flam);
    \draw[more] (lsub) -- node[elbl,above]{more defined (powerset)} (sub);
    \draw[more] (lic)  -- node[elbl,above]{more defined (powerset)} (ic);

    % vertical edges (expressiveness): "strictly more expressive"
    \draw[more] (lsub) -- node[elbl,left]{strictly more expressive} (lic);
    \draw[more] (sub)  -- node[elbl,right]{strictly more expressive} (ic);

    % full abstraction: vertical, in the lazy column
    \draw[iso] (llam) -- node[elbl,left]{fully abstract} (lsub);
  \end{tikzpicture}
  \caption{The calculi of this paper; nodes marked
    (ours) are contributions of this paper.
    Abbreviations:
    $\lambda$ is the $\lambda$-calculus;
    the lazy prefix denotes the approximation that stops after the first
    iteration of the fixpoint of this paper's semantic equations
    (written $T_P\!\uparrow\!1$ in
    logic-programming notation): lookups that re-enter themselves
    through cyclic self-reference return divergence instead of being
    iterated further.
    Continuing:
    fixpoint $\lambda$ is the tabled evaluation of the
    $\lambda$-calculus of Yang~\cite{yang2026-co-lambda};
    IC the inheritance-calculus;
    and $\lambda$-IC and $\lambda$-lazy-IC the $\lambda$ sublanguages of
    IC and of lazy IC, respectively.
    Edge labels:
    \emph{more defined} means every term that converges in the source
    calculus also converges in the target calculus, but not conversely.
    The parenthesized qualifiers distinguish the mechanism:
    \emph{non-strict} more-definedness gains convergence by deeming
    abstractions values without evaluating their bodies;
    \emph{unproductive $\bot$} more-definedness gains convergence by
    deciding unproductive loops such as $\Omega$ as the definite answer
    $\bot$ in finite time, preserving the lazy tree semantics;
    and \emph{powerset} more-definedness gains convergence by
    resolving cyclic self-references to definite sets, as least fixed
    points over the powerset lattice.
    Continuing the edge labels:
    \emph{strictly more expressive} is meant in Felleisen's sense of
    macro-expressibility;
    and \emph{fully abstract} marks full abstraction.}
  \label{fig:calculi-matrix}
\end{figure*}

Figure~\ref{fig:calculi-matrix} situates inheritance-calculus among
several related calculi developed below.
Section~\ref{sec:definitions} defines the semantics via six
mutually recursive first-order functions on paths and sets of paths,
with no function spaces, closures, or higher-order constructs;
the functions are computable by
tabling~\cite{tamaki1986-tabled-resolution}.
Section~\ref{sec:forward-translation} shows that
these functions extend to the $\lambda$-calculus, embedding it.
Section~\ref{sec:asymmetry} establishes that the lazy
$\lambda$-calculus cannot macro-express the inheritance constructors,
so inheritance-calculus is strictly more expressive in
Felleisen's sense~\cite{felleisen1991-expressive-power}.

\section{\bilingual{Syntax}{语法}}
\label{sec:syntax}

\ifInheritanceChinese
设 $e$ 表示一个表达式,$s$ 表示一个元素列表,$c$ 表示一个元素,
$q$ 表示一个限定符,$w$ 表示一个投影列表,$\ell$ 表示一个充当
属性名的标签,$n$ 表示一个 de~Bruijn 索引。
空串 $\varepsilon$ 是空元素列表。
\else
Let $e$ denote an expression, $s$ an element list, $c$ an element,
$q$ a qualifier, $w$ a projection list, $\ell$ a label serving as a
property name, and $n$ a de~Bruijn index.
The empty string $\varepsilon$ is the empty element list.
\fi

\begin{align*}
  e \quad ::= \quad & \{\, s \,\}
  && \text{\bilingual{(record)}{(记录)}} \\[6pt]
  s \quad ::= \quad & \varepsilon
  && \text{\bilingual{(empty)}{(空)}} \\*
  \mid\quad & c
  && \text{\bilingual{(one element)}{(一个元素)}} \\*
  \mid\quad & c,\; s
  && \text{\bilingual{(element, then more)}{(一个元素,然后更多)}} \\[6pt]
  c \quad ::= \quad & \ell \mapsto e
  && \text{\bilingual{(definition)}{(定义)}} \\*
  \mid\quad & q.\, w
  && \text{\bilingual{(inheritance: qualifier and projections)}{(继承:限定符与投影)}} \\*
  \mid\quad & q
  && \text{\bilingual{(inheritance: qualifier only)}{(继承:仅限定符)}} \\*
  \mid\quad & w
  && \text{\bilingual{(inheritance: projections only)}{(继承:仅投影)}} \\[6pt]
  q \quad ::= \quad & \qualifiedthis{\ell_{\mathrm{qualifier}}}
  && \text{\bilingual{(named qualifier)}{(具名限定符)}} \\*
  \mid\quad & \dbi{n}
  && \text{\bilingual{(indexed qualifier)}{(索引限定符)}} \\[6pt]
  w \quad ::= \quad & \ell_{\mathrm{projection}}
  && \text{\bilingual{(one projection)}{(一个投影)}} \\*
  \mid\quad & \ell_{\mathrm{projection}}.\, w
  && \text{\bilingual{(projection, then more)}{(一个投影,然后更多)}}
\end{align*}

\ifInheritanceChinese
一个表达式将一个元素列表 $s$ 包裹在花括号中,而每个元素
$c$ 要么是一个定义 $\ell \mapsto e$,要么是一个继承源。
语法把 $s$ 写成一个递归列表,但该列表被读作一个
\emph{集合}:元素的顺序无关紧要,重复元素也没有任何作用,
因此两个仅在重排或重复上有所不同的元素列表表示
同一个表达式。由于表达式是集合,继承本身就是
可交换、幂等且可结合的。
图~\ref{fig:cheatsheet} 中的记录 $\mathrm{Complex} \mapsto \{ \mathrm{re} \mapsto \ldots,\;
\mathrm{im} \mapsto \ldots,\; \mathrm{Plus} \mapsto \ldots \}$
就是一个这样的多元素列表。

$\ell \mapsto e$ 中的 $\mapsto$ 定义了一个名为 $\ell$、其值为
$e$ 的属性;例如,图~\ref{fig:cheatsheet} 中的 $\mathrm{re} \mapsto \{\ldots\}$ 与
$\mathrm{Complex} \mapsto \{\ldots\}$。
MIXINv2 还通过一个
外部函数接口(FFI)提供标量类型;继承演算就是
去掉 FFI 之后所剩下的部分。
本文中没有任何例子用到 FFI 或任何标量值。

\paragraph{\bilingual{Inheritance}{继承}}
一个继承通过一个限定符 $q$ 引用一个作用域;该限定符
沿作用域层级向上导航,然后通过一个投影列表 $w$
向下投影属性。限定符 $q$ 容许两种等价的
记法:
\begin{description}
  \item[\bilingual{Named form}{具名形式}]
    $\qualifiedthis{\ell_{\mathrm{qualifier}}}.\, w$,
    类比于 Java 的 \texttt{Outer.this.field},其中 $w$ 是一个
    投影列表。
    每个 $\{\ldots\}$ 记录在它的外围作用域中都有一个标签;
    $\ell_{\mathrm{qualifier}}$ 命名目标外围作用域。
    例如,图~\ref{fig:cheatsheet} 中的
    $\qualifiedthis{\mathrm{ComplexModule}}.\mathrm{Float}$
    投影单个标签,而
    $\qualifiedthis{\mathrm{Plus}}.\mathrm{other}.\mathrm{re}$ 展示了一个
    多标签的投影列表 $w$。
  \item[\bilingual{Indexed form}{索引形式}]
    $\dbi{n}.\, w$,
    其中 $n \ge 0$ 是一个 de~Bruijn 索引~\cite{debruijn1972-nameless-dummies}:
    $n = 0$ 指向该引用的外围作用域,
    而 $n$ 每加一,就向外移动一个作用域层级。
\end{description}
两种形式都先取出目标作用域,即所有继承都应用完毕之后
那个完全继承后的记录,然后从中投影由投影列表 $w$
所命名的属性。
具名形式在解析期间脱糖为索引形式,
如第~\ref{sec:mixin-trees} 节所述。

每个 $\{\ldots\}$ 记录字面量都会创建一个新的作用域层级。%
\footnote{
  我们用\emph{记录}指那个语法构造,用\emph{作用域}指
  它所创建的名字解析语境,用\emph{mixin} 指它
  作为可组合单元的角色;在继承演算中这些是同一个实体的三种
  视角。
}
该作用域包含那个 mixin 的所有属性,包括那些
通过继承源继承而来的属性。
例如,$\{a \mapsto \{\},\; b \mapsto \{\}\}$ 是单个 mixin,其作用域
同时包含 $a$ 与 $b$;$\{r,\; a \mapsto \{\}\}$ 也是如此,其中 $r$
是对一个包含 $b$ 的记录的引用。
继承源不会创建额外的作用域层级。
在图~\ref{fig:cheatsheet} 中,绑定到 $\mathrm{Plus}$ 的记录
就是一个这样的作用域,它包含 $\mathrm{other}$、$\mathrm{\_sumReCall}$、
$\mathrm{\_sumRe}$ 以及在其内部直接定义的其他属性。

\begin{itemize}
  \item $\qualifiedthis{\mathrm{Foo}}$(等价地,$\dbi{n}$,其中
    $\mathrm{Foo}$ 在外围作用域之上 $n$ 步)
    取出名为 $\mathrm{Foo}$ 的记录。
  \item $\qualifiedthis{\mathrm{Foo}}.\ell$,或等价地,$\dbi{n}.\ell$,
    从 $\mathrm{Foo}$ 投影属性 $\ell$。
  \item 一个仅限定符的源 $q$(没有投影列表)继承
    整个外围作用域,这可能产生一棵无限深的
    树;图~\ref{fig:cheatsheet} 中裸的 $\mathrm{BuiltIns}$、$\mathrm{Math}$ 与
    $\mathrm{Operator}$ 源就是此类
    仅限定符的导入。
\end{itemize}

当限定符无歧义时,它可以省略:这种简写是
一个裸的投影列表 $w$。它的首个投影针对
那个在其记录字面量中直接定义首个投影标签的
最内层外围作用域来解析,而该源脱糖为
对应的索引形式。
例如,图~\ref{fig:cheatsheet} 中的 $\mathrm{\_sumReCall}.\mathrm{result}$、
$\mathrm{a}.\mathrm{Plus}$ 与 $\mathrm{\_totalCall}.\mathrm{result}$
就是此类省略限定符的简写。
由于这一搜索只检查记录字面量的定义,该
简写只能访问在记录字面量中定义的属性,而不能访问
那些通过继承而继承来的属性;对于继承来的属性或要绕过变量遮蔽,
都需要两种限定形式之一。
\else
An expression encloses an element list $s$ in braces, and each element
$c$ is either a definition $\ell \mapsto e$ or an inheritance source.
The grammar writes $s$ as a recursive list, but this list is read as a
\emph{set}: element order is irrelevant and duplicates have no effect,
so two element lists differing only by reordering or repetition denote
the same expression. Because expressions are sets, inheritance is
inherently commutative, idempotent, and associative.
The record $\mathrm{Complex} \mapsto \{ \mathrm{re} \mapsto \ldots,\;
\mathrm{im} \mapsto \ldots,\; \mathrm{Plus} \mapsto \ldots \}$ in
Figure~\ref{fig:cheatsheet} is one such multi-element list.

The $\mapsto$ in $\ell \mapsto e$ defines a property named $\ell$ with
value $e$; for example, $\mathrm{re} \mapsto \{\ldots\}$ and
$\mathrm{Complex} \mapsto \{\ldots\}$ in Figure~\ref{fig:cheatsheet}.
MIXINv2 also provides scalar types through a
foreign-function interface (FFI); inheritance-calculus is what remains
after removing the FFI.
None of the examples in this paper uses the FFI or any scalar values.

\paragraph{Inheritance}
An inheritance references a scope through a qualifier $q$ that
navigates the scope hierarchy upward, then projects properties downward
through a projection list $w$. The qualifier $q$ admits two equivalent
notations:
\begin{description}
  \item[Named form]
    $\qualifiedthis{\ell_{\mathrm{qualifier}}}.\, w$,
    analogous to Java's \texttt{Outer.this.field}, where $w$ is a
    projection list.
    Each $\{\ldots\}$ record has a label in its enclosing scope;
    $\ell_{\mathrm{qualifier}}$ names the target enclosing scope.
    For example, $\qualifiedthis{\mathrm{ComplexModule}}.\mathrm{Float}$
    in Figure~\ref{fig:cheatsheet} projects a single label, while
    $\qualifiedthis{\mathrm{Plus}}.\mathrm{other}.\mathrm{re}$ shows a
    multi-label projection list $w$.
  \item[Indexed form]
    $\dbi{n}.\, w$,
    where $n \ge 0$ is a de~Bruijn index~\cite{debruijn1972-nameless-dummies}:
    $n = 0$ refers to the reference's enclosing scope,
    and each increment of $n$ moves one scope level further out.
\end{description}
Both forms retrieve the target scope, that is, the fully inherited record
after all inheritance has been applied, and then project the properties
named by the projection list $w$ from it.
The named form desugars to the indexed form during parsing,
as described in Section~\ref{sec:mixin-trees}.

A new scope level is created at each $\{\ldots\}$ record literal.%
\footnote{
  We use \emph{record} for the syntactic construct, \emph{scope}
  for the name-resolution context it creates, and \emph{mixin} for its
  role as a composable unit; in inheritance-calculus these are three views
  of the same entity.
}
The scope contains all properties of that mixin, including those
inherited via inheritance sources.
For example, $\{a \mapsto \{\},\; b \mapsto \{\}\}$ is a single mixin whose scope
contains both $a$ and $b$; so is $\{r,\; a \mapsto \{\}\}$ where $r$
is a reference to a record containing $b$.
Inheritance sources do not create additional scope levels.
In Figure~\ref{fig:cheatsheet}, the record bound to $\mathrm{Plus}$ is
one such scope, containing $\mathrm{other}$, $\mathrm{\_sumReCall}$,
$\mathrm{\_sumRe}$, and the other properties defined directly within it.

\begin{itemize}
  \item $\qualifiedthis{\mathrm{Foo}}$ (equivalently, $\dbi{n}$ where
    $\mathrm{Foo}$ is $n$ steps above the enclosing scope)
    retrieves the record named $\mathrm{Foo}$.
  \item $\qualifiedthis{\mathrm{Foo}}.\ell$, or equivalently, $\dbi{n}.\ell$,
    projects property $\ell$ from $\mathrm{Foo}$.
  \item A qualifier-only source $q$ (with no projection list) inherits
    the entire enclosing scope, which may produce an infinitely deep
    tree; the bare $\mathrm{BuiltIns}$, $\mathrm{Math}$, and
    $\mathrm{Operator}$ sources in Figure~\ref{fig:cheatsheet} are
    qualifier-only imports of this kind.
\end{itemize}

When the qualifier is unambiguous, it may be omitted: the shorthand is
a bare projection list $w$. Its leading projection is resolved against
the innermost enclosing scope that defines the leading projection's
label directly in its record literal, and the source desugars to the
corresponding indexed form.
For example, $\mathrm{\_sumReCall}.\mathrm{result}$,
$\mathrm{a}.\mathrm{Plus}$, and $\mathrm{\_totalCall}.\mathrm{result}$
in Figure~\ref{fig:cheatsheet} are such qualifier-omitted shorthands.
Because this search inspects only record-literal definitions, the
shorthand can only access properties defined in record literals, not
those inherited through inheritance; one of the two qualified forms is
required for inherited properties or to bypass variable shadowing.
\fi

\section{\bilingual{Mixin Trees}{Mixin 树}}
\label{sec:mixin-trees}

\ifInheritanceChinese
在面向对象编程中,一个
\emph{mixin}~\cite{bracha1990-mixin-inheritance} 是一段行为片段;
它可以通过继承组合到一个类上。
继承演算只有一个抽象,即
\emph{深可合并 mixin}(deep-mergeable mixin),其成员是
\emph{深可合并属性}(deep-mergeable properties)。%
\footnote{
  在不引起歧义之处,我们简写为``mixin''与``property''。
  正如图~\ref{fig:cheatsheet} 所示,二者各自推广了它们在先前
  工作中的同名物:Bracha--Cook 的 mixin~\cite{bracha1990-mixin-inheritance} 与
  面向对象语言的属性。
}
它可以同时充当模块、类、对象、方法、字段以及
局部绑定,正如图~\ref{fig:cheatsheet} 所展示的那样,因为
没有不透明的方法:每个值都是
一个透明的记录,而继承总是对子树执行一次深度
合并。
我们把所得到的结构称为一棵\emph{mixin 树}。
它的语义是纯\emph{观察式}(observational)的:给定一个
继承演算程序,我们查询某个
标签 $\ell$ 是否是某个给定位置 $p$ 处的一个属性。

本节通过六个相互递归的
路径上的一阶函数来定义语义。
这六个函数把面向对象的概念映射为
集合论的运算:$\properties$ 与 $\supers$
推广成员查找与继承链,
$\overrides$ 与 $\bases$ 捕捉方法覆盖与
基类,而 $\resolve$ 与 $\this$ 处理名字
解析与限定 $\this$ 解析。
\else
In object-oriented programming, a
\emph{mixin}~\cite{bracha1990-mixin-inheritance} is a fragment of behavior
that can be composed onto a class via inheritance.
Inheritance-calculus has a single abstraction, the
\emph{deep-mergeable mixin}, whose members are
\emph{deep-mergeable properties}.%
\footnote{
  Where no ambiguity arises, we write simply ``mixin'' and ``property.''
  As Figure~\ref{fig:cheatsheet} shows, each generalizes its namesake in prior
  work: the Bracha--Cook mixin~\cite{bracha1990-mixin-inheritance} and the
  property of object-oriented languages.
}
which can simultaneously serve as module, class, object, method, field, and
local binding, as Figure~\ref{fig:cheatsheet} illustrated, because there
are no opaque methods: every value is
a transparent record, and inheritance always performs a deep
merge of subtrees.
We call the resulting structure a \emph{mixin tree}.
Its semantics is purely \emph{observational}: given an
inheritance-calculus program, one queries whether a
label $\ell$ is a property at a given position $p$.

This section defines the semantics via six mutually recursive
first-order functions on paths.
The six functions map object-oriented concepts to
set-theoretic operations: $\properties$ and $\supers$
generalize member lookup and inheritance chains,
$\overrides$ and $\bases$ capture method override and
base classes, and $\resolve$ and $\this$ handle name
resolution and qualified $\this$ resolution.
\fi

\subsection{\bilingual{Definitions}{定义}}
\label{sec:definitions}

\paragraph{\bilingual{Path}{路径}}
\ifInheritanceChinese
一个 $\Path$ 是一个标签序列
$(\ell_1, \ell_2, \ldots, \ell_n)$;它标识树中的一个
位置。
我们用 $p$ 表示一条路径。
\emph{根路径}是空序列 $()$。
给定一条路径 $p$ 与一个标签 $\ell$,
$p \snoc \ell$ 是序列 $p$ 用 $\ell$ 扩展后的结果。
对于非根路径,父路径 $\init(p)$ 是 $p$ 去掉
其最后一个元素后的结果,
而末尾标签 $\last(p)$ 是 $p$ 的最后一个元素。
为方便起见,当路径在上下文中清楚时,我们有时
用 $\ell$ 表示 $\last(p)$。
有了路径,我们就能陈述 AST 在
每条路径处提供什么。
\else
A $\Path$ is a sequence of labels
$(\ell_1, \ell_2, \ldots, \ell_n)$ that identifies a position in
the tree.
We write $p$ for a path.
The \emph{root path} is the empty sequence $()$.
Given a path $p$ and a label $\ell$,
$p \snoc \ell$ is the sequence $p$ extended with $\ell$.
For a nonroot path, the parent path $\init(p)$ is $p$ with
its last element removed,
and the final label $\last(p)$ is the last element of $p$.
For convenience, we sometimes write $\ell$ for $\last(p)$
when the path is clear from context.
With paths in hand, we can state what the AST provides
at each path.
\fi

\paragraph{AST}
\ifInheritanceChinese
来自第~\ref{sec:syntax} 节的一个表达式 $e$ 被解析成一棵 AST。
我们以两种等价的风格给出 AST:首先是一个固定其形状的
代数数据类型,然后是两个派生的投影;本节其余部分的
语义函数消费它们。

\emph{代数数据类型风格}把规范化后的 AST 给出为一个由节点项
构成的纯语法。一个节点 $N$ 把一个定义列表 $d$ 与一个
引用列表 $b$ 配成对;每个定义把一个标签绑定到一个子节点,而
每个引用是一个 de~Bruijn 索引 $n$ 连同一个投影列表
$w$:
\else
An expression $e$ from Section~\ref{sec:syntax} is parsed into an AST.
We give the AST in two equivalent styles: first an algebraic
datatype that fixes its shape, then two derived projections that
the semantic functions in the rest of this section consume.

The \emph{algebraic-datatype style} gives the normalized AST as a pure
grammar of node terms. A node $N$ pairs a definition list $d$ with a
reference list $b$; each definition binds a label to a child node, and
each reference is a de~Bruijn index $n$ together with a projection list
$w$:
\fi
\begin{align*}
  N \quad ::= \quad & \langle\, d \,;\; b \,\rangle
  && \text{\bilingual{(node: definitions and references)}{(节点:定义与引用)}} \\[6pt]
  d \quad ::= \quad & \varepsilon
  && \text{\bilingual{(no definitions)}{(没有定义)}} \\*
  \mid\quad & (\ell \mapsto N),\; d
  && \text{\bilingual{(definition, then more)}{(一个定义,然后更多)}} \\[6pt]
  b \quad ::= \quad & \varepsilon
  && \text{\bilingual{(no references)}{(没有引用)}} \\*
  \mid\quad & (n,\; w),\; b
  && \text{\bilingual{(reference, then more)}{(一个引用,然后更多)}} \\[6pt]
  w \quad ::= \quad & \varepsilon
  && \text{\bilingual{(no projections)}{(没有投影)}} \\*
  \mid\quad & \ell.\, w
  && \text{\bilingual{(projection, then more)}{(一个投影,然后更多)}}
\end{align*}
\ifInheritanceChinese
与表达式的元素列表一样,定义列表 $d$ 是
递归写成的,但读取时不计重排与重复:它
表示从其各绑定收集而来的有限偏映射 $\defs : \Label \rightharpoonup \Node$;该映射
对每个被定义的标签恰有一条目,而
$N.\defs(\ell)$ 是绑定到 $\ell$ 的子节点(当 $\ell$
未绑定时无定义)。同样地,引用列表 $b$ 表示由继承源
$(n,\; w)$ 构成的集合
$\refs \subseteq \mathbb{N} \times \Label^{*}$,读取时不计重排与重复。由于继承
是幂等的且记录是集合,这些商不丢失任何
信息。
该语法只固定形状;一个节点项是良构的,当它
满足:
\else
Like the element list of an expression, the definition list $d$ is
written recursively but read up to reordering and duplication: it
denotes the finite partial map $\defs : \Label \rightharpoonup \Node$
collected from its bindings, with one entry per defined label, and
$N.\defs(\ell)$ is the child bound to $\ell$ (undefined when $\ell$ is
unbound). Likewise the reference list $b$ denotes the set
$\refs \subseteq \mathbb{N} \times \Label^{*}$ of inheritance sources
$(n,\; w)$, read up to reordering and duplication. Because inheritance
is idempotent and a record is a set, these quotients lose no
information.
The grammar fixes only the shape; a node term is well-formed when it
satisfies:
\fi
\begin{enumerate}
  \item \ifInheritanceChinese 路径 $p$ 处的每个源 $(n,\; w)$ 满足 $n < |p|$。\else Each source $(n,\; w)$ at path $p$ satisfies $n < |p|$.\fi
  \item \ifInheritanceChinese 一个 $w$ 为空的源 $(n,\; \varepsilon)$ 是
    一个仅限定符源的规范化;第~\ref{sec:syntax} 节的表达式语法
    没有 $w$ 为空的仅投影源。\else A source $(n,\; \varepsilon)$ with empty $w$ is the
    normalization of a qualifier-only source; the expression grammar of
    Section~\ref{sec:syntax} has no projection-only source with empty
    $w$.\fi
\end{enumerate}
\ifInheritanceChinese
该节点语法被\emph{余归纳地}(coinductively)读取:通过
继承,一个子节点可能展开成一棵无界乃至
无限深的树,所以这些产生式描述的是最大不动点,
而非有限生成的最小不动点。整个程序是单个
全局节点 $\ASTroot$,即顶层记录的规范化 AST。

\emph{余代数风格}(coalgebraic style)用两个由路径 $p$ 索引的
原始函数替换全局树,每次观察一个节点。它们通过沿 $p$
导航从 $\ASTroot$ 派生而来:
从 $\ASTroot$ 出发,在 $p$ 的每个标签处跟随 $\defs$ 映射
到达 $p$ 处的节点,而这两个投影读出其
字段。
\else
The node grammar is read \emph{coinductively}: through
inheritance a child node may expand to a tree of unbounded or even
infinite depth, so the productions describe the greatest fixed point,
not a finitely generated least one. The whole program is a single
global node $\ASTroot$, the normalized AST of the top-level record.

The \emph{coalgebraic style} replaces the global tree by two
primitive functions indexed by a path $p$, observing one node at a
time. They are derived from $\ASTroot$ by navigating along $p$:
starting at $\ASTroot$ and following the $\defs$ map at each label
of $p$ reaches the node at $p$, and the two projections read off its
fields.
\fi
\begin{align}
  \defines(p) &= \operatorname{dom}\bigl(
    \ASTroot.\defs(\ell_1).\defs(\ell_2)\cdots\defs(\ell_n).\defs
  \bigr)
  \label{eq:defines}\\
  \inherits(p) &= \ASTroot.\defs(\ell_1).\defs(\ell_2)\cdots\defs(\ell_n).\refs
  \label{eq:inherits}
\end{align}
\ifInheritanceChinese
其中 $p = (\ell_1, \ldots, \ell_n)$,从而 $\defines(p)$ 是
那些使得 $p$ 包含一个定义 $\ell \mapsto e$ 的标签 $\ell$ 的集合,
而 $\inherits(p)$ 是 $p$ 处的引用对
$(n,\; w)$ 的集合。
在根路径 $()$ 处导航为空,给出
$\defines(()) = \operatorname{dom}(\ASTroot.\defs)$ 以及
$\inherits(()) = \ASTroot.\refs$。
下文的语义函数只使用 $\defines$ 与 $\inherits$,
所以任一风格皆可取作原始。

在解析期间,三种语法形式全都被解析为
de~Bruijn 索引对 $(n,\; w)$。
\emph{索引形式} $\dbi{n}.\, w$
已经直接携带 de~Bruijn 索引 $n$;
不需要解析。
路径 $p$ 处的一个\emph{具名引用} $\qualifiedthis{\ell_{\mathrm{qualifier}}}.\, w$
通过在 $p$ 的标签中找到
$\ell_{\mathrm{qualifier}}$ 的最后一次出现来解析。
设 $p_{\mathrm{target}}$ 为 $p$ 的前缀,直到并包括
那次出现。
de~Bruijn 索引为 $n = |p| - |p_{\mathrm{target}}| - 1$,
而投影列表 $w$ 原样保留。
路径 $p$ 处一个首投影为 $\ell$ 的\emph{词法引用} $w$
通过找到 $p$ 的最近前缀 $p'$
使得 $\ell \in \defines(p')$ 来解析。
de~Bruijn 索引为 $n = |p| - |p'| - 1$,而投影列表
$w$ 原样保留。
三种形式都产生同一种表示;下文的语义函数
只在 de~Bruijn 索引对上操作。

AST 是纯数据,由程序文本固定,而不是
运行期状态的函数。%
\footnote{
  一个实现可以非严格地解析 AST:$\ASTroot$ 不必
  预先被完全物化。每个 $\Node$ 把它的源
  文本保持为未解析,直到某个查询观察它,从而读取
  $\defines(p)$ 或 $\inherits(p)$ 时按需强制
  沿 $p$ 的子树。这只影响解析何时发生,而不影响其结果,
  因为 AST 是引用透明的。
}
有了两个原始函数 $\defines$ 与 $\inherits$,我们现在可以回答
本节开头提出的那个问题:
某个标签 $\ell$ 是否是某个给定位置 $p$ 处的一个属性。
\else
for $p = (\ell_1, \ldots, \ell_n)$, so that $\defines(p)$ is the set
of labels $\ell$ for which $p$ contains a definition $\ell \mapsto e$,
and $\inherits(p)$ is the set of reference pairs
$(n,\; w)$ at $p$.
At the root path $()$ the navigation is empty, giving
$\defines(()) = \operatorname{dom}(\ASTroot.\defs)$ and
$\inherits(()) = \ASTroot.\refs$.
The semantic functions below use only $\defines$ and $\inherits$,
so either style may be taken as primitive.

During parsing, all three syntactic forms are resolved to
de~Bruijn index pairs $(n,\; w)$.
The \emph{indexed form} $\dbi{n}.\, w$
already carries the de~Bruijn index $n$ directly;
no resolution is needed.
A \emph{named reference} $\qualifiedthis{\ell_{\mathrm{qualifier}}}.\, w$
at path $p$ is resolved by finding the last occurrence of
$\ell_{\mathrm{qualifier}}$ among the labels of $p$.
Let $p_{\mathrm{target}}$ be the prefix of $p$ up to and including
that occurrence.
The de~Bruijn index is $n = |p| - |p_{\mathrm{target}}| - 1$
and the projection list $w$ is carried over unchanged.
A \emph{lexical reference} $w$ at path $p$, with leading
projection $\ell$, is resolved by finding the nearest prefix $p'$ of $p$
such that $\ell \in \defines(p')$.
The de~Bruijn index is $n = |p| - |p'| - 1$ and the projection list
$w$ is carried over unchanged.
All three forms produce the same representation; the semantic functions
below operate only on de~Bruijn index pairs.

The AST is pure data, fixed by the program text and not a function
of runtime state.%
\footnote{
  An implementation may parse the AST non-strictly: $\ASTroot$ need
  not be materialized in full up front. Each $\Node$ keeps its source
  text unparsed until a query observes it, so that reading
  $\defines(p)$ or $\inherits(p)$ forces the subtree along $p$ on
  demand. This affects only when parsing happens, not its result,
  since the AST is referentially transparent.
}
Given the two primitives $\defines$ and $\inherits$, we can now answer
the question posed at the beginning of this section:
whether a label $\ell$ is a property at a given position $p$.
\fi

\paragraph{\bilingual{Properties}{属性}}
\ifInheritanceChinese
一条路径的属性是它自身的属性,连同
从所有 supers 继承而来的属性:
\else
The properties of a path are its own properties together with
those inherited from all supers:
\fi
\begin{equation}\label{eq:properties}
  \properties(p) =
  \bigl\{\; \ell \;\big|\;
    (\_,\; p_{\mathrm{override}}) \in \supers(p),\;
    \ell \in \defines(p_{\mathrm{override}})
  \;\bigr\}
\end{equation}
\ifInheritanceChinese
本节其余部分定义 $\supers$,
即传递的继承闭包,
及其依赖项。
\else
The remainder of this section defines $\supers$,
the transitive inheritance closure,
and its dependencies.
\fi

\paragraph{Supers}
\ifInheritanceChinese
直观上,$\supers(p)$ 收集 $p$ 所继承的每一条路径:
$p$ 自身的 $\overrides$,
加上 $p$ 的每个
直接基 $\bases$ 的 $\overrides$,如此传递下去。
每个结果与到达它所经由的继承点语境
$\init(p_{\mathrm{base}})$ 配成对。
我们下文定义的限定 $\this$ 解析需要这一来源,
以便把一个定义点 mixin 映射回
那些纳入它的继承点路径:
\else
Intuitively, $\supers(p)$ collects every path that $p$ inherits from:
the $\overrides$ of $p$ itself,
plus the $\overrides$ of each
direct base $\bases$ of $p$, and so on transitively.
Each result is paired with the inheritance-site context
$\init(p_{\mathrm{base}})$ through which it is reached.
This provenance is needed by qualified $\this$ resolution,
which we define below, to map a definition-site mixin back to
the inheritance-site paths that incorporate it:
\fi
\begin{equation}\label{eq:supers}
  \supers(p) =
  \bigl\{\; (\init(p_{\mathrm{base}}),\; p_{\mathrm{override}}) \;\big|\;
    p_{\mathrm{base}} \in \bases^*(p),\;
    p_{\mathrm{override}} \in \overrides(p_{\mathrm{base}})
  \;\bigr\}
\end{equation}
\ifInheritanceChinese
这里 $\bases^*$ 表示 $\bases$ 的自反传递闭包。
$\supers$ 公式依赖两个函数:
$\overrides$ 与单跳引用目标 $\bases$。
我们先定义 $\overrides$。
\else
Here $\bases^*$ denotes the reflexive-transitive closure of $\bases$.
The $\supers$ formula depends on two functions:
$\overrides$ and the one-hop reference targets $\bases$.
We define $\overrides$ first.
\fi

\paragraph{Overrides}
\ifInheritanceChinese
由于一个 mixin 可以扮演一个方法的角色,
一个子作用域应当能够
通过父作用域的标签找到它,即便该子作用域
并不重新定义那个标签;这就是方法覆盖。
具体地说,继承可以在同一个作用域层级引入同一标签的
多个定义;
$\overrides(p)$ 收集所有与 $p$ 共享
\emph{相同身份}(same identity)的这类路径,从而使它们的定义
相互合并而非相互遮蔽,以实现深度合并。
$p$ 的 overrides 包含 $p$ 自身
以及对于 $\init(p)$ 的每个
也定义 $\last(p)$ 的分支 $p_{\mathrm{branch}}$,$p_{\mathrm{branch}} \snoc \last(p)$:
\else
Since a mixin can play the role of a method,
a child scope should be
able to find it through the parent's label even if the child
does not redefine that label; this is method override.
Concretely, inheritance can introduce multiple
definitions of the same label at the same scope level;
$\overrides(p)$ collects all such paths that share the
\emph{same identity} as $p$, so that their definitions
merge rather than shadow each other, enabling deep merge.
The overrides of $p$ include $p$ itself
and $p_{\mathrm{branch}} \snoc \last(p)$ for every
branch $p_{\mathrm{branch}}$ of $\init(p)$ that also defines $\last(p)$:
\fi
\begin{equation}\label{eq:overrides}
  \overrides(p) =
  \begin{cases}
    \{p\} & \text{if } p = ()  \\[6pt]
    \{p\}
    \;\cup\;
    \left\{\; p_{\mathrm{branch}} \snoc \last(p) \;\middle|\;
      \begin{aligned}
        &(\_,\; p_{\mathrm{branch}}) \in \supers(\init(p)), \\
        &\text{s.t.}\; \last(p) \in \defines(p_{\mathrm{branch}})
      \end{aligned}
    \right\}
    & \text{if } p \neq ()
  \end{cases}
\end{equation}
\ifInheritanceChinese
还剩下要定义 $\bases$,即 $\supers$ 的另一个依赖项。
\else
It remains to define $\bases$, the other dependency of $\supers$.
\fi

\paragraph{Bases}
\ifInheritanceChinese
直观上,$\bases(p)$ 是 $p$ 通过引用直接
继承的那些路径,类比于面向对象语言中的直接
基类。
具体地说,$\bases$ 把 $p$ 的 $\overrides$ 中的每个引用
解析一步:
\else
Intuitively, $\bases(p)$ are the paths that $p$ directly
inherits from via references, analogous to the direct base
classes in object-oriented languages.
Concretely, $\bases$ resolves every reference in
$p$'s $\overrides$ one step:
\fi
\begin{equation}\label{eq:bases}
  \bases(p) =
  \left\{\; p_{\mathrm{target}} \;\middle|\;
    \begin{aligned}
      &p_{\mathrm{override}} \in \overrides(p), \\
      &(n,\; w) \in \inherits(p_{\mathrm{override}}), \\
      &p_{\mathrm{target}} \in \resolve(\init(p),\; p_{\mathrm{override}},\; n,\; w)
    \end{aligned}
  \right\}
\end{equation}
\ifInheritanceChinese
$\bases$ 公式调用引用解析函数 $\resolve$,我们接下来定义它。
\else
The $\bases$ formula calls the reference resolution function $\resolve$, which we define next.
\fi

\paragraph{\bilingual{Reference resolution}{引用解析}}
\ifInheritanceChinese
直观上,$\resolve$ 把一个语法引用转化为它在
完全继承后的树中所指向的路径集合。
它取一条继承点路径 $p_{\mathrm{site}}$、
一条定义点路径 $p_{\mathrm{def}}$、
一个 de~Bruijn 索引 $n$,
以及投影列表 $w$。
解析分两个阶段进行:
$\this$ 从外围作用域 $\init(p_{\mathrm{def}})$ 出发
执行 $n$ 次向上步进,
把定义点路径映射为继承点路径;
然后追加
$w = \ell_{\mathrm{projection},1}\ldots\ell_{\mathrm{projection},k}$
的向下投影:
\else
Intuitively, $\resolve$ turns a syntactic reference into the
set of paths it points to in the fully inherited tree.
It takes an inheritance-site path $p_{\mathrm{site}}$,
a definition-site path $p_{\mathrm{def}}$,
a de~Bruijn index $n$,
and the projection list $w$.
Resolution proceeds in two phases:
$\this$ performs $n$ upward steps
starting from the enclosing scope $\init(p_{\mathrm{def}})$,
mapping the definition-site path to inheritance-site paths;
then the downward projections of
$w = \ell_{\mathrm{projection},1}\ldots\ell_{\mathrm{projection},k}$
are appended:
\fi
\begin{multline}\label{eq:resolve}
  \resolve(p_{\mathrm{site}},\; p_{\mathrm{def}},\; n,\; w)
  = \\
  \bigl\{\;
    p_{\mathrm{current}} \snoc \ell_{\mathrm{projection},1} \snoc \cdots \snoc \ell_{\mathrm{projection},k}
    \;\big|
    p_{\mathrm{current}} \in
    \this(\{p_{\mathrm{site}}\},\; \init(p_{\mathrm{def}}),\; n)
  \;\bigr\}
\end{multline}
\ifInheritanceChinese
$\resolve$ 返回一个集合而非单条路径,因为
$\this$ 可能解析到多个目标。
我们现在定义 $\this$。
\else
$\resolve$ returns a set rather than a single path because
$\this$ may resolve to multiple targets.
We now define $\this$.
\fi

\paragraph{\bilingual{Qualified $\this$ resolution}{限定 $\this$ 解析}}
\ifInheritanceChinese
$\this$ 回答这个问题:
``在完全继承后的树中,定义点作用域
$p_{\mathrm{def}}$ 究竟住在哪里?''
每一步在前沿 $S$ 中每条路径的 supers 当中,
找出那些其 override 分量匹配
$p_{\mathrm{def}}$ 的,并把对应的
继承点路径收集为新的前沿。
经过 $n$ 步后,前沿就包含答案:
\else
$\this$ answers the question:
``in the fully inherited tree, where does
the definition-site scope $p_{\mathrm{def}}$ actually live?''
Each step finds, among the supers of every path in the
frontier $S$, those whose override component matches
$p_{\mathrm{def}}$, and collects the corresponding
inheritance-site paths as the new frontier.
After $n$ steps the frontier contains the answer:
\fi
{\small
  \begin{equation}\label{eq:this}
    \this(S,\; p_{\mathrm{def}},\; n) =
    \begin{cases}
      S & \text{if } n = 0 \\[6pt]
      \this\!\left(
        \left\{\; p_{\mathrm{site}} \;\middle|\;
          \begin{aligned}
            &p_{\mathrm{current}} \in S, \\
            &(p_{\mathrm{site}},\; p_{\mathrm{override}})
            \in \supers(p_{\mathrm{current}}), \\
            &\text{s.t.}\; p_{\mathrm{override}} = p_{\mathrm{def}}
          \end{aligned}
        \right\},\;
        \init(p_{\mathrm{def}}),\;
        n - 1
      \right)
      & \text{if } n > 0
    \end{cases}
\end{equation}}%
\ifInheritanceChinese
每一步 $\init$ 把 $p_{\mathrm{def}}$ 缩短一个标签,
而 $n$ 减少一。
由于 $n$ 是一个非负整数,该递归终止。%
\footnote{
  若在某一步不存在匹配对
  $(p_{\mathrm{site}},\; p_{\mathrm{override}})$,
  则前沿变为空集,而
  $\resolve$ 返回 $\varnothing$:该引用没有
  目标;不继承任何属性。
  这是无类型演算最为自洽的
  语义。
  MIXINv2\anon[~(补充材料)]{~\cite{mixin2025}} 更为严格:它的解析器拒绝
  前沿变为空的引用,在编译期报告一个
  解析错误。
}

这就完成了计算 $\properties(p)$ 所需的整条定义链。

附录~\ref{app:evaluation-trace} 把这六个方程逐步
应用到一个突出 $\resolve$ 集合值本质的例子上;
在那里前沿 $S$ 携带多个目标,而
$\this$ 同时解析到它们,而非如 Scala~3 与 NixOS 模块
这类面向对象语言中所见的单个目标
(附录~\ref{app:scala-multi-target})。
\else
At each step $\init$ shortens $p_{\mathrm{def}}$ by one label
and $n$ decreases by one.
Since $n$ is a nonnegative integer, the recursion terminates.%
\footnote{
  If no matching pair
  $(p_{\mathrm{site}},\; p_{\mathrm{override}})$ exists
  at some step, the frontier becomes the empty set, and
  $\resolve$ returns $\varnothing$: the reference has no
  target; no properties are inherited.
  This is the most self-consistent semantics for the
  untyped calculus.
  MIXINv2\anon[~(supplementary material)]{~\cite{mixin2025}} is stricter: its resolver rejects
  references whose frontier becomes empty, reporting a
  resolution error at compile time.
}

This completes the chain of definitions needed to compute $\properties(p)$.

Appendix~\ref{app:evaluation-trace} applies the six equations step by
step to an example highlighting the set-valued nature of $\resolve$,
where the frontier $S$ carries multiple targets that
$\this$ resolves to simultaneously rather than the single target
seen in object-oriented languages such as Scala~3 and NixOS modules
(Appendix~\ref{app:scala-multi-target}).
\fi

\subsection{\bilingual{Recursive Evaluation}{递归求值}}
\ifInheritanceChinese
六个方程~(\ref{eq:properties})--(\ref{eq:this})
通过相互递归定义 $\properties(p)$:
$\ell \in \properties(p)$ 成立,当且仅当它能
由这些函数的有限次应用链
确立。
这些函数\emph{就是}解释器:观察者
通过选择检查哪条路径来驱动计算,
而每个方程按需展开。
由于这些函数是纯的且与顺序无关,
(\ref{eq:supers}) 中的自反传递闭包 $\bases^*$
可以按广度优先或
深度优先来探索而不影响结果;
路径可以被 intern,从而结构相等
退化为指针相等;
而这六个函数构成一个以路径为键的可记忆化
动态规划递推式。

在操作上,求值使用\emph{tabled
求值}~\cite{tamaki1986-tabled-resolution}:
每个查询在首次遇到时被记忆化,而
重入调用,即对应于依赖图中的环的调用,
返回当前的累加器而非发散。
在指称上,六个方程定义了一个单调的
直接后承
算子~$T_P$~\cite{vanemden1976-predicate-logic-semantics};
它的右端只使用 $\cup$、$\in$、相等以及
正向集合概括,没有否定、集合
差或补。
由 Knaster--Tarski
定理~\cite{tarski1955-lattice-fixpoint-theorem},$T_P$ 有一个唯一的
最小不动点 $\mathrm{lfp}(T_P)$;它充当
指称语义。
这一指称刻画把上文引入的观察式
语义形式化:观察者驱动的查询
由同一组不动点方程计算。
求值是否终止取决于程序。
回忆 $\supers$(方程~\ref{eq:supers})取
$\bases$(方程~\ref{eq:bases})的自反传递
闭包;所得的图就是
\emph{继承层级}(inheritance hierarchy)。
由于每个继承源在 mixin 树中占据一个不同的节点,
这一层级可从
树结构中直接读出。
继承层级无环的程序在
一趟内终止。
当层级包含环时,
当每个查询有有限的答案域时求值
终止:有限格上的升链条件保证
收敛~\cite{bancilhon1986-recursive-query-strategies},
如同在 Datalog~\cite{vanemden1976-predicate-logic-semantics} 中那样。
当层级无限时,求值可能
发散;这是图灵完备性所固有的。
附录~\ref{app:well-definedness} 证明 tabled
求值绝不返回错误答案,并把
这些终止条件形式化。
\else
The six equations~(\ref{eq:properties})--(\ref{eq:this})
define $\properties(p)$ by mutual recursion:
$\ell \in \properties(p)$ holds if and only if it can
be established by a finite chain of applications of
these functions.
The functions \emph{are} the interpreter: the observer
drives the computation by choosing which path to
inspect, and each equation unfolds on demand.
Because the functions are pure and order-independent,
the reflexive-transitive closure $\bases^*$
in~(\ref{eq:supers}) may be explored breadth-first or
depth-first without affecting the result;
paths may be interned so that structural equality
reduces to pointer equality;
and the six functions form a memoizable
dynamic-programming recurrence keyed by paths.

Operationally, evaluation uses \emph{tabled
evaluation}~\cite{tamaki1986-tabled-resolution}:
each query is memoised on first encounter, and
re-entrant calls, which correspond to cycles in the dependency graph,
return the current accumulator rather than diverging.
Denotationally, the six equations define a monotone
immediate consequence
operator~$T_P$~\cite{vanemden1976-predicate-logic-semantics}
whose right-hand sides use only $\cup$, $\in$, equality, and
positive set comprehension, with no negation, set
difference, or complement.
By the Knaster--Tarski
theorem~\cite{tarski1955-lattice-fixpoint-theorem}, $T_P$ has a unique
least fixed point $\mathrm{lfp}(T_P)$, which serves as
the denotational semantics.
This denotational characterization formalizes the observational
semantics introduced above: the observer-driven queries are
computed by the same fixed-point equations.
Whether evaluation terminates depends on the program.
Recall that $\supers$ (equation~\ref{eq:supers}) takes the
reflexive-transitive closure of $\bases$
(equation~\ref{eq:bases}); the resulting graph is the
\emph{inheritance hierarchy}.
Since each inheritance source occupies a distinct node in
the mixin tree, this hierarchy is directly readable from
the tree structure.
Programs whose inheritance hierarchy is acyclic terminate in
one pass.
When the hierarchy contains cycles,
evaluation terminates when each query has a finite
answer domain: the ascending chain condition on the finite
lattice guarantees
convergence~\cite{bancilhon1986-recursive-query-strategies},
as in Datalog~\cite{vanemden1976-predicate-logic-semantics}.
When the hierarchy is infinite, evaluation may
diverge; this is inherent to Turing completeness.
Appendix~\ref{app:well-definedness} proves that tabled
evaluation never returns a wrong answer and formalizes
these termination conditions.
\fi

%% === 审稿人备忘：六个方程的渐近复杂度分析 ===
%%
%% 论文声称语义是一阶的，审稿人可能追问：一阶性是否带来额外的
%% 计算开销？下面按场景分析。
%%
%% 关键观察：set comprehension 自动去重（集合语义），加上路径
%% intern（结构相等退化为指针相等）和 memoization，相同的
%% (path, query) 对只计算一次。这不是可选的优化，而是集合语义
%% 的固有性质——集合不会包含重复元素。
%%
%% --- 场景 1：λ-calculus 翻译 ---
%% Single-Path Lemma (Lemma 3) 保证 this 的 frontier 集合 S
%% 始终只有一个元素。每个路径的 overrides 也只有一个元素。
%% bases 每个节点最多一个继承源。supers 的传递闭包是线性链。
%% 复杂度：O(n)，其中 n 是 ANF term 的大小。
%% 与传统环境机（environment machine）一致。
%%
%% --- 场景 2：Object algebra（k 个 constructor，t 个 operation 文件，m 个字段） ---
%% 多个 operation 文件各自继承同一个 data 文件，然后在最终组合时
%% merge。每个 constructor 在每个 operation 文件里各有一个定义，
%% 通过 overrides 合并，所以 overrides 大小为 O(t)。
%%
%% 关键点：this 的 frontier 是 1，不是"去重后 O(1)"，就是 1。
%% 原因：qualified this 解析的是外层 scope（如 NatFactory），
%% 而这个外层 scope 在继承链上只有一条路径到达——因为是从一个
%% 具体的 operation 定义点出发往上走的。multi-target this 的前提
%% 是同一个定义点被继承到多个不同的 site，这只在 trie union 时
%% 发生。Object algebra 的继承是把多个 operation merge 到同一个
%% factory 上，不是把多个 factory merge 成一个值。方向不同。
%%
%% 这也解释了与 Scala 的关系：Scala 拒绝 multi-target self-reference resolution
%%（Appendix B 的 "conflicting base types"），但在 object algebra
%% 场景下 this 本来就是 single-target 的，所以 Scala 的拒绝是
%% 过度保守——它拒绝了一类去重后（实际上根本不需要去重，因为
%% frontier 就是 1）完全合法的组合。
%%
%% 复杂度：O(k · m · t)，与 Scala trait linearization 一致。
%%
%% --- 场景 3：Nat/BinNat/Boolean 算术 ---
%% 数值的递归深度为 n（Nat 值的大小），每次递归触发一次 this
%% 解析（single-target），遍历 O(t) 个 override。
%% Nat 算术复杂度：O(n · t)，即 Peano 算术 + 常数因子。
%% BinNat 算术复杂度：O(log n · t)，即二进制算术 + 常数因子。
%%
%% --- 场景 4：Cartesian product / 逻辑编程 ---
%% 这是唯一出现真正 multi-target this 的场景。
%% trie union 的 bases 包含多条路径，每条路径的 addend 通过
%% this 解析到另一个 trie union，路径数相乘。
%% 对于 |R₁| × |R₂| 的 join，复杂度为 O(|R₁| · |R₂|)
%% （memoization 下）。
%%
%% 这与 B-Tree 引擎的 Datalog nested-loop join 复杂度一致。
%% trie 天然有索引——路径本身是前缀树，共享前缀的值自动共享
%% 计算（memoization 命中）。这与 Datalog 引擎对 B-Tree 做
%% prefix scan 在精神上一致。
%%
%% --- 总结 ---
%%
%% | 场景               | 复杂度              | 对比               |
%% |--------------------|---------------------|--------------------|
%% | λ-calculus         | O(n)                | = 环境机           |
%% | Object algebra     | O(k·m·t)            | = Scala线性化      |
%% | Nat 算术           | O(n·t)              | = Peano 算术       |
%% | BinNat 算术        | O(log n · t)        | = 二进制算术       |
%% | Cartesian product  | O(|R₁|·|R₂|)       | = nested-loop join |
%%
%% 每个场景都与对应的专用引擎复杂度一致。一阶性没有带来额外开销。
%%
%% Single-Path Lemma 的适用范围比论文明示的更大：不仅 λ-calculus
%% 翻译是 single-path，所有不涉及 trie union 的正常 OOP/object
%% algebra 代码都是 single-path。multi-target this 是 trie union
%% 的专属现象，而 trie union 对应的恰好是关系代数的 join，其
%% 复杂度是固有的，与引擎无关。
%% ================================================================

\section{\bilingual{Embedding the \texorpdfstring{$\lambda$}{λ}-Calculus}{嵌入 \texorpdfstring{$\lambda$}{λ}-演算}}
\label{sec:translation}
\label{sec:forward-translation}

\ifInheritanceChinese
以下翻译将 A-正规形式（ANF）~\cite{flanagan1993-essence-compiling-continuations}的 $\lambda$-演算映射到继承演算。
在 ANF 中,每个中间结果都绑定到一个名字,且每个表达式至多包含一次应用。
我们用一个\emph{值绑定} $\mathbf{let}\; x = V \;\mathbf{in}\; M$ 来扩展标准 ANF 语法,该值绑定将一个并非应用结果的值命名:%
\footnote{
  标准 ANF 允许将值内联用作操作数;值绑定产生式为其命名,使得翻译后程序的每个子表达式都可通过词法引用访问(附录~\ref{app:expressiveness})。
}
\[
  \begin{array}{r@{\;::=\;}l}
    M & \mathbf{let}\; x = V_1\; V_2 \;\mathbf{in}\; M
    \mid \mathbf{let}\; x = V \;\mathbf{in}\; M
    \mid V_1\; V_2
    \mid V \\
    V & x \mid \lambda x.\, M
  \end{array}
\]
每个 $\lambda$-项都可以通过对中间结果命名的方式机械地转换为 ANF。\footnote{%
  这可以选择性地包括对值的命名;由于每个 ANF 项也是一个 $\lambda$-项,ANF 既不是 $\lambda$-演算的限制,也不是其扩展。
}
ANF 结构与继承演算天然吻合:每个 $\mathbf{let}$-绑定成为一个命名的记录属性,每次应用则通过 $.\mathrm{result}$ 包装并投影。

设 $\mathcal{T}$ 为翻译函数。每个 $\lambda$-抽象在翻译后的 mixin 树中引入一个作用域层级。
\else
The following translation maps the $\lambda$-calculus in
A-normal form (ANF)~\cite{flanagan1993-essence-compiling-continuations} to inheritance-calculus.
In ANF, every intermediate result is bound to a name and each
expression contains at most one application.
We extend the standard ANF grammar with a \emph{value binding}
$\mathbf{let}\; x = V \;\mathbf{in}\; M$, which gives a name to a
value that is not an application result:%
\footnote{
  Standard ANF allows values inline as operands;
  the value-binding production names them so that every
  subexpression of the translated program is accessible
  via a lexical reference (Appendix~\ref{app:expressiveness}).
}
\[
  \begin{array}{r@{\;::=\;}l}
    M & \mathbf{let}\; x = V_1\; V_2 \;\mathbf{in}\; M
    \mid \mathbf{let}\; x = V \;\mathbf{in}\; M
    \mid V_1\; V_2
    \mid V \\
    V & x \mid \lambda x.\, M
  \end{array}
\]
Every $\lambda$-term can be mechanically converted to ANF by
naming intermediate results.\footnote{%
  This can optionally include naming values as well;
  since every ANF term is also a $\lambda$-term,
  ANF is neither a restriction nor an extension of the $\lambda$-calculus.
}
The ANF structure matches inheritance-calculus naturally: each
$\mathbf{let}$-binding becomes a named record property, and
each application is wrapped and projected via
$.\mathrm{result}$.

Let $\mathcal{T}$ denote the translation function.
Each $\lambda$-abstraction introduces one scope level
in the translated mixin tree.
\fi

\medskip
\begin{center}
  \small
  \begin{tabular}{l@{\quad$\longrightarrow$\quad}l}
    $x$ \text{\scriptsize(\bilingual{$\lambda$-bound, index $n'$}{$\lambda$-绑定,索引 $n'$})} &
    $\dbi{n'}.\mathrm{argument}$ \\[4pt]
    $x$ \text{\scriptsize(\bilingual{let-bound}{let-绑定})} &
    $x.\mathrm{result}$ \\[4pt]
    $\lambda x.\, M$ &
    $\{\mathrm{argument} \mapsto \{\},\;
    \mathrm{result} \mapsto \mathcal{T}(M)\}$ \\[4pt]
    $\mathbf{let}\; x = V_1\; V_2
    \;\mathbf{in}\; M$ &
    $\{x \mapsto \{\mathcal{T}(V_1),\;
      \mathrm{argument} \mapsto \mathcal{T}(V_2)\},$ \\
    & $\;\mathrm{result} \mapsto \mathcal{T}(M)\}$ \\[4pt]
    $\mathbf{let}\; x = V
    \;\mathbf{in}\; M$ &
    $\{x \mapsto \{\mathrm{result} \mapsto \mathcal{T}(V)\},\;
    \mathrm{result} \mapsto \mathcal{T}(M)\}$ \\[4pt]
    $V_1\; V_2$ \text{\scriptsize(\bilingual{tail call}{尾调用})} &
    $\{\mathrm{tailCall} \mapsto \{\mathcal{T}(V_1),\;
      \mathrm{argument} \mapsto \mathcal{T}(V_2)\},$ \\
    & $\;\mathrm{result} \mapsto \{\mathrm{tailCall}.\mathrm{result}\}\}$
  \end{tabular}
\end{center}
\medskip

\ifInheritanceChinese
带 de~Bruijn 索引 $n$ 的 $\lambda$-绑定变量 $x$(即从引用处到其绑定者之间包围的 $\lambda$ 数目)翻译为 $\dbi{n'}.\mathrm{argument}$,该引用到达绑定 $x$ 的抽象并访问其 $\mathrm{argument}$ 槽。%
\footnote{\label{fn:debruijn-adjust}在 $\lambda$-演算中,de~Bruijn 索引只计数 $\lambda$-抽象。在继承演算中,de~Bruijn 索引计数所有包围的作用域层级,包括由 $\mathbf{let}$-绑定和尾调用引入的层级。因此翻译会调整该索引:若一个 $\lambda$-绑定变量跨越 $k$ 个中间的 $\mathbf{let}$-绑定或尾调用作用域(即 $k$ 计数从变量出现处到其绑定的 $\lambda$ 在 ANF 语法树中路径上的 $\mathbf{let}$- 和尾调用构造数),则其继承演算索引为 $n' = n + k$,而非 $n$。let-绑定变量不受影响,因为它们使用词法引用(如 $x.\mathrm{result}$)而非 de~Bruijn 索引。我们假设翻译引入的合成标签($\mathrm{argument}$、$\mathrm{result}$、$\mathrm{tailCall}$)与用户定义的标签(如 let-绑定名 $x$)来自不相交的字母表,因此不会发生冲突。}
let-绑定变量 $x$ 翻译为词法引用 $x.\mathrm{result}$,从封闭记录中的兄弟属性 $x$ 投影应用结果。

抽象 $\lambda x.\, M$ 翻译为具有自身属性 $\{\mathrm{argument}, \mathrm{result}\}$ 的记录;我们称之为\emph{抽象形状}。在应用 $\mathbf{let}$-绑定规则中,名字 $x$ 绑定到继承记录 $\{\mathcal{T}(V_1),\;\mathrm{argument} \mapsto \mathcal{T}(V_2)\}$,而 $\mathcal{T}(M)$ 可引用 $x.\mathrm{result}$ 以获取应用的值。在值 $\mathbf{let}$-绑定规则中,$x$ 将 $\mathcal{T}(V)$ 包装在一个 $\mathrm{result}$ 子节点之下,使得统一访问器 $x.\mathrm{result}$ 产生 $\mathcal{T}(V)$ 本身,即翻译后的值。这个额外的包装在语义上是惰性的,但与证明相关:它使得用于应用的同一 $\mathrm{result}$-追踪论证在充分性证明中能均匀地扩展到值 $\mathbf{let}$-绑定(附录~\ref{app:bohm-tree-proofs})。在尾调用规则中,新标签 $\mathrm{tailCall}$ 起到相同作用,而 $\mathrm{result}$ 投影出答案。在应用和尾调用两种情形中,执行应用的继承被封装在一个命名属性之后,使其内部结构(被调用者的 $\mathrm{argument}$ 和 $\mathrm{result}$ 标签)不会泄漏到封闭作用域。

以上六条规则构成了从 $\lambda$-演算(经由 ANF)到继承演算的完整翻译:每个 ANF 项都有一个像,且构造是复合的。由于 $\lambda$-演算是图灵完备的,继承演算也是。%
\footnote{
  MIXINv2 将符号解析(解析词法引用和限定 $\this$ 引用)与运行时 mixin 求值分离,但符号解析是逐符号即时执行的,而非提前执行,因此它不提供传统意义上的静态类型检查。尽管如此,MIXINv2 的解析阶段比无类型的继承演算更严格。本文中的所有示例均从 MIXINv2 移植而来,并同时符合两种语义。为满足解析器而添加的元素对此处呈现的算法无害,且可能有助于可读性。
}
然而,图灵完备性并未说明翻译保留了什么;下一小节将精确刻画语义对应关系。
\else
A $\lambda$-bound variable $x$ with de~Bruijn index $n$
(the number of enclosing $\lambda$s between
the reference and its binder) becomes
$\dbi{n'}.\mathrm{argument}$,
which reaches the abstraction that binds $x$ and
accesses its $\mathrm{argument}$ slot.%
\footnote{\label{fn:debruijn-adjust}In the $\lambda$-calculus, de~Bruijn
  indices count only $\lambda$-abstractions.
  In inheritance-calculus, de~Bruijn indices count all
  enclosing scope levels, including those introduced by
  $\mathbf{let}$-bindings and tail calls.
  The translation therefore adjusts the index:
  if a $\lambda$-bound variable crosses $k$ intervening
  $\mathbf{let}$-binding or tail-call scopes
  (that is, $k$ counts the $\mathbf{let}$- and tail-call
    constructs on the path from the variable occurrence
  to its binding $\lambda$ in the ANF syntax tree),
  its inheritance-calculus
  index is $n' = n + k$,
  not~$n$.
  Let-bound variables are unaffected, as they use
  lexical references (such as $x.\mathrm{result}$) rather than
  de~Bruijn indices.
  We assume that the synthetic labels introduced by the
  translation ($\mathrm{argument}$, $\mathrm{result}$,
  $\mathrm{tailCall}$) and the user-defined labels
  (let-bound names such as~$x$) are drawn from disjoint
  alphabets, so no collision can arise.
}
A let-bound variable $x$ becomes the lexical reference
$x.\mathrm{result}$, projecting the application result
from the sibling property $x$ in the enclosing record.

An abstraction $\lambda x.\, M$ translates to a record with own
properties $\{\mathrm{argument}, \mathrm{result}\}$; we call this
the \emph{abstraction shape}.
In the application $\mathbf{let}$-binding rule, the name $x$
binds to the inherited record
$\{\mathcal{T}(V_1),\;\mathrm{argument} \mapsto \mathcal{T}(V_2)\}$,
and $\mathcal{T}(M)$ may reference $x.\mathrm{result}$ to obtain
the value of the application.
In the value $\mathbf{let}$-binding rule, $x$ wraps $\mathcal{T}(V)$
under a $\mathrm{result}$ child so that the uniform accessor
$x.\mathrm{result}$ yields $\mathcal{T}(V)$, the translated value
itself.
This extra wrapper is semantically inert but proof-relevant:
it lets the same $\mathrm{result}$-following argument used for
applications extend uniformly to value
$\mathbf{let}$-bindings in the adequacy proof
(Appendix~\ref{app:bohm-tree-proofs}).
In the tail-call rule, the fresh label $\mathrm{tailCall}$ serves
the same purpose, and $\mathrm{result}$ projects the answer.
In the application and tail-call cases, the inheritance that
performs the application is
encapsulated behind a named property, so its internal
structure (the $\mathrm{argument}$ and $\mathrm{result}$ labels
of the callee) does not leak to the enclosing scope.

The six rules above are a complete translation from the
$\lambda$-calculus (via ANF) to inheritance-calculus: every ANF
term has an image, and the construction is compositional.
Since the $\lambda$-calculus is Turing complete, so is
inheritance-calculus.%
\footnote{
  MIXINv2 separates symbol resolution (which resolves
  lexical references and qualified $\this$ references) from runtime
  mixin evaluation, but symbol resolution is performed
  just-in-time per symbol rather than ahead of time, so it does
  not provide static type checking in the traditional sense.
  Nevertheless, MIXINv2's resolution phase is stricter than the
  untyped inheritance-calculus.
  All examples in this paper are ported from MIXINv2
  and conform to both semantics.
  Elements added to satisfy the resolver are harmless
  for the algorithms presented here and may aid readability.
}
Turing completeness, however, says nothing about what the
translation preserves; the next subsection pins down
the semantic correspondence.
\fi

\subsection{\bilingual{B\"ohm Tree Correspondence}{Böhm 树对应}}
\label{sec:bohm-tree}

\ifInheritanceChinese
翻译 $\mathcal{T}$ 将 $\lambda$-演算嵌入继承演算的一个子语言中。直觉上,mixin 树以附加的语义结构增广了 AST;对于翻译后的 $\lambda$-项,这一附加结构的投影与 Böhm 树同构。若从根出发追踪有限次 $\mathrm{result}$ 投影后到达一个同时拥有 $\mathrm{argument}$ 和 $\mathrm{result}$ 的节点(即 $\mathcal{T}(\lambda x.\, M)$ 产生的\emph{抽象形状}),则称 mixin 树 $T$ \emph{收敛},记为 $T{\Downarrow}$。翻译保持并反映收敛性(充分性)。记 $\mathcal{T}(M) \approx_\lambda \mathcal{T}(N)$ 当对所有封闭的 $\lambda$-演算上下文 $C[\cdot]$ 均有 $\mathcal{T}(C[M]){\Downarrow} \Leftrightarrow \mathcal{T}(C[N]){\Downarrow}$;则 $\approx_\lambda$ 恰好与 Böhm 树等价~\cite{barendregt1984-lambda-calculus} 重合。形式化陈述和证明见附录~\ref{app:bohm-tree-proofs}。注意 $\approx_\lambda$ 仅在 $\lambda$-演算上下文上量化;继承演算的上下文区分能力严格更强,如第~\ref{sec:asymmetry} 节所证。以上对应关系在首次迭代近似 $T_P\!\uparrow\!1$ 处成立;我们现在考察在完整最小不动点处会发生什么。

在 $T_P\!\uparrow\!1$ 处,每个方程被求值一次,可重入调用返回 $\bot$;这对应于 $\beta$-规约。在完整的 $\mathrm{lfp}(T_P)$ 处,可重入引用改而解析为幂集格上的最小不动点:迭代在\emph{容器}层级计算不动点,求出满足 $S = F(S)$ 的集合 $S$。这正是递归 Datalog 在循环图上收敛的机制(附录~\ref{app:datalog-encoding},第~\ref{sec:recursive-rules} 节);既非不动点组合子(如 $Y$),也非 Haskell 的 \texttt{mfix}(\texttt{RecursiveDo})能够表达它:两者都在项层级打结,而非在语义所累积的集合上打结。破坏 Böhm 树对应关系的这种\emph{幂集}更多确定性(图~\ref{fig:calculi-matrix})的最小见证例是 $\mathbf{let}\; x = x \;\mathbf{in}\; x$:它在 $\beta$-规约下发散,但其翻译 $\{x \mapsto \{\mathrm{result} \mapsto \{x.\mathrm{result}\}\},\; \mathrm{result} \mapsto \{x.\mathrm{result}\}\}$ 在 $\mathrm{lfp}(T_P)$ 下稳定,tabling 求值检测到循环并返回确定结果 $\properties(\mathrm{root}) = \{x, \mathrm{result}\}$。因此本小节的对应关系是 $T_P\!\uparrow\!1$ 独有的性质。下一小节考察继承演算是否也严格更具表达力。
\else
The translation $\mathcal{T}$ embeds the $\lambda$-calculus
into a sublanguage of inheritance-calculus.
Intuitively, a mixin tree augments the AST with
additional semantic structure; for translated
$\lambda$-terms, the projection of this additional
structure is isomorphic to the B\"ohm tree.
A mixin tree $T$ \emph{converges}, written $T{\Downarrow}$,
if following finitely many $\mathrm{result}$ projections from
the root reaches a node with both $\mathrm{argument}$ and
$\mathrm{result}$ among its properties, that is, the \emph{abstraction shape}
produced by $\mathcal{T}(\lambda x.\, M)$.
The translation preserves and reflects convergence
(adequacy).
Write $\mathcal{T}(M) \approx_\lambda \mathcal{T}(N)$
when $\mathcal{T}(C[M]){\Downarrow} \Leftrightarrow
\mathcal{T}(C[N]){\Downarrow}$ for every closing
$\lambda$-calculus context $C[\cdot]$;
then $\approx_\lambda$ coincides with B\"ohm tree
equivalence~\cite{barendregt1984-lambda-calculus}.
Formal statements and proofs are in
Appendix~\ref{app:bohm-tree-proofs}.
Note that $\approx_\lambda$ quantifies over $\lambda$-calculus
contexts only; inheritance-calculus contexts are strictly more
discriminating, as Section~\ref{sec:asymmetry} establishes.
The correspondence above holds at the first-iteration
approximation $T_P\!\uparrow\!1$; we now ask what happens
at the full least fixed point.

At $T_P\!\uparrow\!1$, each equation is evaluated once and
re-entrant calls return~$\bot$; this corresponds to
$\beta$-reduction.
At the full $\mathrm{lfp}(T_P)$, re-entrant references instead
resolve to least fixed points over the powerset lattice: the
iteration computes fixed points at the \emph{container} level,
finding a set~$S$ such that $S = F(S)$.
This is the mechanism by which recursive Datalog converges on
cyclic graphs (Appendix~\ref{app:datalog-encoding},
Section~\ref{sec:recursive-rules}), and neither a fixed-point
combinator such as~$Y$ nor Haskell's \texttt{mfix}
(\texttt{RecursiveDo}) expresses it: both tie the knot at the
term level, not over the sets the semantics accumulates.
The smallest witness that this \emph{powerset} more-definedness
(Figure~\ref{fig:calculi-matrix}) breaks the B\"ohm-tree
correspondence is
$\mathbf{let}\; x = x \;\mathbf{in}\; x$:
it diverges under $\beta$-reduction, but its translation
$\{x \mapsto \{\mathrm{result} \mapsto \{x.\mathrm{result}\}\},\;
\mathrm{result} \mapsto \{x.\mathrm{result}\}\}$
stabilizes under $\mathrm{lfp}(T_P)$, where tabled evaluation
detects the cycle and returns the definite result
$\properties(\mathrm{root}) = \{x, \mathrm{result}\}$.
The correspondence of this subsection is therefore a property
of $T_P\!\uparrow\!1$ alone.
The next subsection examines whether inheritance-calculus is
also strictly more expressive.
\fi

\subsection{\bilingual{Expressive Asymmetry}{表达力不对称}}
\label{sec:asymmetry}
\label{sec:formal-separation}
\ifInheritanceChinese
两个 Böhm 树等价的 $\lambda$-项,在 $\lambda$-演算中因此无法区分,可以被一个混合了 $\lambda$-演算与继承演算构造子的上下文分离;该上下文通过开放扩展在已有定义上继承新的可观测投影。根据 Felleisen 的表达力判据~\cite{felleisen1991-expressive-power},继承演算严格比惰性 $\lambda$-演算更具表达力(定理~\ref{thm:expressive-asymmetry}):$\lambda$-演算可以通过 ANF 转换和 $\mathcal{T}$ 宏表达于继承演算中(引理~\ref{lem:anf-equiv},定理~\ref{thm:forward-macro}),但惰性 $\lambda$-演算无法宏表达继承(定理~\ref{thm:cartesian-nonexpressibility})。形式化定义、构造和证明见附录~\ref{app:expressiveness};该分离构造本身是表达式问题的一个微型实例,它在不修改已有定义的情况下添加了一个新的可观测投影。
\else
Two $\lambda$-terms that are B\"ohm-tree equivalent,
and therefore indistinguishable in the $\lambda$-calculus,
can be
separated by a context that mixes $\lambda$-calculus and
inheritance-calculus constructors, inheriting new observable
projections onto an existing definition via open extension.
By Felleisen's expressiveness
criterion~\cite{felleisen1991-expressive-power}, inheritance-calculus is
strictly more expressive than the lazy $\lambda$-calculus
(Theorem~\ref{thm:expressive-asymmetry}):
the $\lambda$-calculus can be macro-expressed in
inheritance-calculus via ANF conversion
and~$\mathcal{T}$
(Lemma~\ref{lem:anf-equiv}, Theorem~\ref{thm:forward-macro}),
but the lazy $\lambda$-calculus cannot macro-express inheritance
(Theorem~\ref{thm:cartesian-nonexpressibility}).
The formal definitions, constructions, and proofs are given in
Appendix~\ref{app:expressiveness};
the separating construction is itself a miniature instance of
the Expression Problem, adding a new observable projection to
an existing definition without modifying it.
\fi

\section{\bilingual{Case Study: Nat Arithmetic}{案例研究：Nat(自然数)算术}}
\label{sec:case-study}
\label{sec:expression-problem}

\ifInheritanceChinese
我们在继承演算中实现了布尔逻辑与自然数算术,包括一元与二进制两种形式。
本节详细追踪一元 Nat 的情形:其数据表示、加法与相等判断。
该实现遵循普通的 Gang-of-Four 设计模式~\cite{gamma1994-design-patterns},采用声明式面向对象风格,将每项关注点沿两个轴分解为 mixin:操作轴与情形轴。由于演算本身没有方法,所有名称均为名词或形容词:UpperCamelCase 名称扮演模块、类型、数据构造子与方法的角色,lowerCamelCase 名称扮演字段与参数的角色。
\else
We implemented boolean logic and natural number arithmetic, both unary
and binary, in inheritance-calculus.
This section traces the unary Nat case in detail: its data representation,
addition, and equality.
The implementation follows ordinary Gang-of-Four design
patterns~\cite{gamma1994-design-patterns} in a declarative object-oriented
style, with each concern split into mixins along two axes: operations and
cases. Because the calculus has no methods of its own, every name is a noun or
adjective: UpperCamelCase names play the roles of modules, types, data
constructors, and methods, and lowerCamelCase names the roles of fields and
parameters.
\fi

\paragraph{NatData}
\ifInheritanceChinese
mixin $\mathrm{NatData}$ 是自然数接口及其工厂。$\mathrm{NatFactory}$ 是一个工厂,其抽象产品类型被别名为 $\mathrm{Nat}$:
\else
The mixin $\mathrm{NatData}$ is the natural-number interface and its factory.
$\mathrm{NatFactory}$ is a factory whose abstract product type is aliased as $\mathrm{Nat}$:
\fi
\[
  \mathrm{NatData} \mapsto \left\{
    \begin{aligned}
      &\mathrm{NatFactory} \mapsto \{\mathrm{Product} \mapsto \{\}\},\\
      &\mathrm{Nat} \mapsto \{\mathrm{NatFactory}.\mathrm{Product}\}
  \end{aligned}\right\}
\]

\paragraph{NatDataZero}
\ifInheritanceChinese
每个数据构造子都是一个独立的实现 mixin,继承 $\mathrm{NatData}$。一个 $\mathrm{Zero}$ 值继承 $\mathrm{Product}$:
\else
Each data constructor is a separate implementation mixin that inherits
$\mathrm{NatData}$.
A $\mathrm{Zero}$ value inherits $\mathrm{Product}$:
\fi
\[
  \mathrm{NatDataZero} \mapsto \left\{
    \begin{aligned}
      &\mathrm{NatData},\\
      &\mathrm{NatFactory} \mapsto \{\mathrm{Zero} \mapsto \{\qualifiedthis{\mathrm{NatFactory}}.\mathrm{Product}\}\}
  \end{aligned}\right\}
\]

\paragraph{NatDataSuccessor}
\ifInheritanceChinese
一个 $\mathrm{Successor}$ 值继承 $\mathrm{Product}$,并暴露一个类型为 $\mathrm{Product}$ 的 $\mathrm{predecessor}$ 属性:
\else
A $\mathrm{Successor}$ value inherits $\mathrm{Product}$ and
exposes a $\mathrm{predecessor}$ property whose type is $\mathrm{Product}$:
\fi
\[
  \begin{aligned}
    &\mathrm{NatDataSuccessor} \mapsto \\
    &\left\{
    \begin{aligned}
      &\mathrm{NatData},\\
      &\mathrm{NatFactory} \mapsto \\
      &\left\{\mathrm{Successor} \mapsto \left\{
        \begin{aligned}
          &\qualifiedthis{\mathrm{NatDataSuccessor}}.\mathrm{NatFactory}.\mathrm{Product},\\
          &\mathrm{predecessor} \mapsto \{\qualifiedthis{\mathrm{NatDataSuccessor}}.\mathrm{NatFactory}.\mathrm{Product}\}
      \end{aligned}\right\}\right\}
  \end{aligned}\right\}
  \end{aligned}
\]
\ifInheritanceChinese
与 $\mathrm{NatDataZero}$ 合在一起,这以开放的方式完成了自然数的 Peano 编码,即函子 $F(X) = 1 + X$ 的初始代数:两个构造子作为独立的 mixin 加入,而非固定在单个封闭的数据类型定义中。
\else
Together with $\mathrm{NatDataZero}$, this completes the Peano encoding of
the natural numbers, the initial algebra of the functor $F(X) = 1 + X$,
in an open way: the two constructors are added as independent mixins
rather than fixed in a single closed datatype definition.
\fi

\paragraph{NatPlus}
\ifInheritanceChinese
加法操作沿情形轴分解为一个\emph{接口} mixin $\mathrm{NatPlus}$,该接口 mixin 保留关注点的名称并在抽象 $\mathrm{Product}$ 上声明命令,以及每个情形对应一个\emph{实现} mixin,实现 mixin 在名称后附加构造子名称(即下文的 $\mathrm{NatPlusZero}$ 与 $\mathrm{NatPlusSuccessor}$)。$\mathrm{NatPlus}$ 继承 $\mathrm{NatData}$。面向对象方法会用动词命名的地方,命令模式~\cite{gamma1994-design-patterns} 将其具化为以名词命名的命令,因此加法方法被命名为 $\mathrm{Plus}$ 而非 $\mathrm{Add}$,并通过投影属性 $\mathrm{sum}$ 暴露其结果;这里已指定类型,但尚未计算:
\else
The addition operation is split along the case axis into an \emph{interface}
mixin $\mathrm{NatPlus}$, which keeps the concern's name and declares the command
on the abstract $\mathrm{Product}$, and one \emph{implementation} mixin per case,
which appends the constructor name ($\mathrm{NatPlusZero}$ and
$\mathrm{NatPlusSuccessor}$ below). $\mathrm{NatPlus}$ inherits
$\mathrm{NatData}$. Where an
object-oriented method would be named with a verb, the command
pattern~\cite{gamma1994-design-patterns} reifies it as a noun-named command, so
the addition method is named $\mathrm{Plus}$ rather than $\mathrm{Add}$, and it
exposes its result through the projection property $\mathrm{sum}$, here typed but
not yet computed:
\fi
\[
  \mathrm{NatPlus} \mapsto \left\{
    \begin{aligned}
      &\mathrm{NatData},\\
      &\mathrm{NatFactory} \mapsto \left\{\mathrm{Product} \mapsto \left\{\mathrm{Plus} \mapsto \left\{\mathrm{sum} \mapsto \{\mathrm{Product}\}\right\}\right\}\right\}
  \end{aligned}\right\}
\]

\paragraph{NatPlusZero}
\ifInheritanceChinese
基础情形为 $0 + m = m$。这里 $\mathrm{NatPlusZero}$ 充当一个导入其他模块的模块:它继承 $\mathrm{Zero}$ 数据与 $\mathrm{Plus}$ 接口,并直接返回被加数:
\else
The base case is $0 + m = m$. Here $\mathrm{NatPlusZero}$ acts as a module that
imports other modules: it inherits the $\mathrm{Zero}$ data and the $\mathrm{Plus}$
interface, and returns the addend directly:
\fi
\[
  \begin{aligned}
    &\mathrm{NatPlusZero} \mapsto \\
    &\left\{
    \begin{aligned}
      &\mathrm{NatDataZero},\; \mathrm{NatPlus},\\
      &\mathrm{NatFactory} \mapsto \left\{\mathrm{Zero} \mapsto \left\{\mathrm{Plus} \mapsto \left\{
          \begin{aligned}
            &\mathrm{addend} \mapsto \{\qualifiedthis{\mathrm{NatPlusZero}}.\mathrm{NatFactory}.\mathrm{Product}\},\\
            &\mathrm{sum} \mapsto \{\mathrm{addend}\}
      \end{aligned}\right\}\right\}\right\}
  \end{aligned}\right\}
  \end{aligned}
\]

\paragraph{NatPlusSuccessor}
\ifInheritanceChinese
递归情形将 $\mathrm{S}(n_0) + m$ 规约为 $n_0 + \mathrm{S}(m)$,方法是以递增后的被加数委托给 $n_0$ 的 $\mathrm{Plus}$:
\else
The recursive case reduces $\mathrm{S}(n_0) + m$ to $n_0 + \mathrm{S}(m)$
by delegating to $n_0$'s $\mathrm{Plus}$ with an incremented addend:
\fi
\[
  \begin{aligned}
    &\mathrm{NatPlusSuccessor} \mapsto \\
    &\left\{
    \begin{aligned}
      &\mathrm{NatDataSuccessor},\; \mathrm{NatPlus},\\
      &\mathrm{NatFactory} \mapsto \\
      &\left\{\mathrm{Successor} \mapsto \left\{\mathrm{Plus} \mapsto \left\{
              \begin{aligned}
                &\mathrm{addend} \mapsto \{\qualifiedthis{\mathrm{NatFactory}}.\mathrm{Product}\},\\
                &\mathrm{\_increasedAddend} \mapsto  \\
                &\left\{
                  \begin{aligned}
                    &\mathrm{Successor},\\
                    &\mathrm{predecessor} \mapsto \{\mathrm{addend}\}
                \end{aligned}\right\},\\
                &\mathrm{\_recursiveAddition} \mapsto \\
                &\left\{
                  \begin{aligned}
                    &\qualifiedthis{\mathrm{Successor}}.\mathrm{predecessor}.\mathrm{Plus},\\
                    &\mathrm{addend} \mapsto \{\mathrm{\_increasedAddend}\}
                \end{aligned}\right\},\\
                &\mathrm{sum} \mapsto \{\mathrm{\_recursiveAddition}.\mathrm{sum}\}
          \end{aligned}\right\}\right\}\right\}
  \end{aligned}\right\}
  \end{aligned}
\]

\paragraph{NatVisitor}
\ifInheritanceChinese
相等判断需要对数据构造子进行情形分析;我们使用访问者模式~\cite{gamma1994-design-patterns}。$\mathrm{NatVisitor}$ 向 $\mathrm{Product}$ 添加一个 $\mathrm{Acceptance}$ 命令,即访问者模式的 $\mathrm{accept}$ 方法的名词形式。mixin $\mathrm{NatVisitor}$ 将其声明为空,留下 $\mathrm{VisitorMap}$ 与 $\mathrm{Accepted}$ 属性由各构造子的实现来填充:
\else
Equality requires case analysis on the data constructor; we use the
visitor pattern~\cite{gamma1994-design-patterns}.
$\mathrm{NatVisitor}$ adds an $\mathrm{Acceptance}$ command to $\mathrm{Product}$,
the noun form of the visitor pattern's $\mathrm{accept}$ method.
The mixin $\mathrm{NatVisitor}$ declares it empty, leaving a $\mathrm{VisitorMap}$ and an
$\mathrm{Accepted}$ property for the per-constructor implementations to fill in:
\fi
\[
  \begin{aligned}
    &\mathrm{NatVisitor} \mapsto \\
    &\left\{
    \begin{aligned}
      &\mathrm{NatData},\\
      &\mathrm{NatFactory} \mapsto \left\{\mathrm{Product} \mapsto \{\mathrm{Acceptance} \mapsto \{\mathrm{VisitorMap} \mapsto \{\},\; \mathrm{Accepted} \mapsto \{\}\}\}\right\}
  \end{aligned}\right\}
  \end{aligned}
\]

\paragraph{NatVisitorZero}
\ifInheritanceChinese
$\mathrm{Zero}$ 的实现选择 $\mathrm{ZeroVisitor}$ 分支。这里的访问者模式与~\cite{gamma1994-design-patterns} 中的表述不同,那里数据构造子由 $\mathrm{accept}$ 接口固定;此处通过组合允许添加新的构造子,因为每个构造子作为独立的 mixin 贡献其自身的分支:
\else
The $\mathrm{Zero}$ implementation selects the $\mathrm{ZeroVisitor}$ branch.
The visitor pattern here, unlike its formulation in
\cite{gamma1994-design-patterns} where the data constructors are fixed by the
$\mathrm{accept}$ interface, admits new constructors by composition, since each
constructor contributes its own branch as a separate mixin:
\fi
\[
  \begin{aligned}
    &\mathrm{NatVisitorZero} \mapsto \\
    &\left\{
    \begin{aligned}
      &\mathrm{NatDataZero},\; \mathrm{NatVisitor},\\
      &\mathrm{NatFactory} \mapsto \left\{\mathrm{Zero} \mapsto \left\{\mathrm{Acceptance} \mapsto \left\{
                \begin{aligned}
                  &\mathrm{VisitorMap} \mapsto \{\mathrm{ZeroVisitor} \mapsto \{\}\},\\
                  &\mathrm{Accepted} \mapsto \{\mathrm{VisitorMap}.\mathrm{ZeroVisitor}\}
              \end{aligned}\right\}\right\}\right\}
  \end{aligned}\right\}
  \end{aligned}
\]

\paragraph{NatVisitorSuccessor}
\ifInheritanceChinese
$\mathrm{Successor}$ 的实现选择 $\mathrm{SuccessorVisitor}$ 分支:
\else
The $\mathrm{Successor}$ implementation selects the $\mathrm{SuccessorVisitor}$ branch:
\fi
\[
  \begin{aligned}
    &\mathrm{NatVisitorSuccessor} \mapsto \\
    &\left\{
    \begin{aligned}
      &\mathrm{NatDataSuccessor},\; \mathrm{NatVisitor},\\
      &\mathrm{NatFactory} \mapsto \\
      &\left\{\mathrm{Successor} \mapsto \left\{
            \begin{aligned}
              &\mathrm{predecessor} \mapsto \{\qualifiedthis{\mathrm{NatFactory}}.\mathrm{Product}\},\\
              &\mathrm{Acceptance} \mapsto \left\{
                \begin{aligned}
                  &\mathrm{VisitorMap} \mapsto \{\mathrm{SuccessorVisitor} \mapsto \{\}\},\\
                  &\mathrm{Accepted} \mapsto \{\mathrm{VisitorMap}.\mathrm{SuccessorVisitor}\}
              \end{aligned}\right\}
          \end{aligned}\right\}\right\}
  \end{aligned}\right\}
  \end{aligned}
\]

\paragraph{BooleanData}
\ifInheritanceChinese
$\mathrm{BooleanData}$ 定义输出的域:
\else
$\mathrm{BooleanData}$ defines the output sort:
\fi
\[
  \mathrm{BooleanData} \mapsto \left\{
    \begin{aligned}
      &\mathrm{BooleanFactory} \mapsto \left\{
        \begin{aligned}
          &\mathrm{Product} \mapsto \{\},\\
          &\mathrm{True} \mapsto \{\mathrm{Product}\},\\
          &\mathrm{False} \mapsto \{\mathrm{Product}\}
      \end{aligned}\right\},\\
      &\mathrm{Boolean} \mapsto \{\mathrm{BooleanFactory}.\mathrm{Product}\}
  \end{aligned}\right\}
\]
\paragraph{NatEquality}
\label{sec:nat-equality}
\ifInheritanceChinese
为了验证 $2 + 3 = 5$,我们需要对 Nat 值进行相等判断。mixin $\mathrm{NatEquality}$ 导入 $\mathrm{NatVisitor}$ 与 $\mathrm{BooleanData}$,并在 $\mathrm{Product}$ 上声明一个 $\mathrm{Equal}$ 命令,由各构造子的实现提供。由于 $\mathrm{NatEquality}$ 继承 $\mathrm{BooleanData}$,限定 $\this$ $\qualifiedthis{\mathrm{NatEquality}}$ 可以到达 $\mathrm{Boolean}$,因此 $\mathrm{Equal}$ 以类型 $\qualifiedthis{\mathrm{NatEquality}}.\mathrm{Boolean}$ 声明其结果 $\mathrm{equal}$:
\else
To test that $2 + 3 = 5$, we needed equality on Nat values.
The mixin $\mathrm{NatEquality}$ imports $\mathrm{NatVisitor}$ and $\mathrm{BooleanData}$,
and declares an $\mathrm{Equal}$ command on $\mathrm{Product}$ that the per-constructor implementations provide.
Because $\mathrm{NatEquality}$ inherits $\mathrm{BooleanData}$,
the qualified this $\qualifiedthis{\mathrm{NatEquality}}$ can reach $\mathrm{Boolean}$,
so $\mathrm{Equal}$ declares its result
$\mathrm{equal}$ with the type $\qualifiedthis{\mathrm{NatEquality}}.\mathrm{Boolean}$:
\fi
\[
  \begin{aligned}
    &\mathrm{NatEquality} \mapsto \\
    &\left\{
    \begin{aligned}
      &\mathrm{NatVisitor},\; \mathrm{BooleanData},\\
      &\mathrm{NatFactory} \mapsto \left\{\mathrm{Product} \mapsto \left\{\mathrm{Equal} \mapsto \left\{
          \begin{aligned}
            &\mathrm{other} \mapsto \{\mathrm{Product}\},\\
            &\mathrm{equal} \mapsto \{\qualifiedthis{\mathrm{NatEquality}}.\mathrm{Boolean}\}
      \end{aligned}\right\}\right\}\right\}
  \end{aligned}\right\}
  \end{aligned}
\]

\paragraph{NatEqualityZero}
\ifInheritanceChinese
$\mathrm{Zero}$ 只与另一个 $\mathrm{Zero}$ 相等。这里的继承执行情形分析:继承 $\mathrm{other}.\mathrm{Acceptance}$ 对 $\mathrm{other}$ 是由 $\mathrm{Zero}$ 还是 $\mathrm{Successor}$ 构建进行分派,在 $\mathrm{ZeroVisitor}$ 分支返回 $\mathrm{True}$,否则返回 $\mathrm{False}$,即 $\mathrm{equal} = \{ \mathrm{True} \text{ 若 } \mathrm{other} = \mathrm{Zero},\; \mathrm{False} \text{ 若 } \mathrm{other} = \mathrm{Successor} \}$。$\mathrm{True}$ 与 $\mathrm{False}$ 均为 $\mathrm{BooleanFactory}$ 的数据构造子;该工厂来自 $\mathrm{BooleanData}$,通过继承的 $\mathrm{NatEquality}$ 间接导入,并以 $\qualifiedthis{\mathrm{NatEqualityZero}}.\mathrm{BooleanFactory}$ 到达。同一继承机制也提供了结果类型:在 $\mathrm{Accepted}$ 中,$\mathrm{equal}$ 继承 $\qualifiedthis{\mathrm{NatEqualityZero}}.\mathrm{Boolean}$,从而声明其类型为 $\mathrm{Boolean}$:
\else
$\mathrm{Zero}$ is equal only to another $\mathrm{Zero}$.
Inheriting here performs case analysis:
inheriting $\mathrm{other}.\mathrm{Acceptance}$ dispatches on whether $\mathrm{other}$
was built by $\mathrm{Zero}$ or $\mathrm{Successor}$, returning $\mathrm{True}$
on the $\mathrm{ZeroVisitor}$ branch and $\mathrm{False}$ otherwise, that is,
$\mathrm{equal} = \{ \mathrm{True} \text{ if } \mathrm{other} = \mathrm{Zero},\; \mathrm{False} \text{ if } \mathrm{other} = \mathrm{Successor} \}$.
Both $\mathrm{True}$ and $\mathrm{False}$ are data constructors of
$\mathrm{BooleanFactory}$, which comes from $\mathrm{BooleanData}$, imported
indirectly through the inherited $\mathrm{NatEquality}$, and is reached as
$\qualifiedthis{\mathrm{NatEqualityZero}}.\mathrm{BooleanFactory}$.
The same inheritance mechanism also supplies the result type: in
$\mathrm{Accepted}$, $\mathrm{equal}$ inherits
$\qualifiedthis{\mathrm{NatEqualityZero}}.\mathrm{Boolean}$, declaring its type to
be $\mathrm{Boolean}$:
\fi
{\small
  \[
    \begin{aligned}
      &\mathrm{NatEqualityZero} \mapsto \\
      &\left\{
      \begin{aligned}
        &\mathrm{NatVisitorZero},\; \mathrm{NatEquality},\\
        &\mathrm{NatFactory} \mapsto \\
        &\left\{
        \begin{aligned}
        &\mathrm{Zero} \mapsto \\
        &\left\{
        \begin{aligned}
        &\mathrm{Equal} \mapsto \\
        &\left\{
                  \begin{aligned}
                    &\mathrm{other} \mapsto \{\qualifiedthis{\mathrm{NatEqualityZero}}.\mathrm{NatFactory}.\mathrm{Product}\},\\
                    &\mathrm{\_OtherAcceptance} \mapsto \\
                    &\left\{
                      \begin{aligned}
                        &\mathrm{other}.\mathrm{Acceptance},\\
                        &\mathrm{VisitorMap} \mapsto \\
                        &\left\{
                          \begin{aligned}
                            &\mathrm{ZeroVisitor} \mapsto \{\mathrm{equal} \mapsto \{\qualifiedthis{\mathrm{NatEqualityZero}}.\mathrm{BooleanFactory}.\mathrm{True}\}\},\\
                            &\mathrm{SuccessorVisitor} \mapsto \{\mathrm{equal} \mapsto \{\qualifiedthis{\mathrm{NatEqualityZero}}.\mathrm{BooleanFactory}.\mathrm{False}\}\}
                        \end{aligned}\right\},\\
                        &\mathrm{Accepted} \mapsto \{\mathrm{equal} \mapsto \{\qualifiedthis{\mathrm{NatEqualityZero}}.\mathrm{Boolean}\}\}
                    \end{aligned}\right\},\\
                    &\mathrm{equal} \mapsto \{\mathrm{\_OtherAcceptance}.\mathrm{Accepted}.\mathrm{equal}\}
                \end{aligned}\right\}
        \end{aligned}\right\}
        \end{aligned}\right\}
    \end{aligned}\right\}
    \end{aligned}
\]}

\paragraph{NatEqualitySuccessor}
\ifInheritanceChinese
$\mathrm{Successor}$ 只与另一个具有相等前驱的 $\mathrm{Successor}$ 相等,因此其实现通过 $\qualifiedthis{\mathrm{Successor}}.\mathrm{predecessor}.\mathrm{Equal}$ 进行递归:
\else
$\mathrm{Successor}$ is equal only to another $\mathrm{Successor}$ with an equal
predecessor, so its implementation recurses through
$\qualifiedthis{\mathrm{Successor}}.\mathrm{predecessor}.\mathrm{Equal}$:
\fi
{\small
  \[
    \begin{aligned}
      &\mathrm{NatEqualitySuccessor} \mapsto \\
      &\left\{
      \begin{aligned}
        &\mathrm{NatVisitorSuccessor},\; \mathrm{NatEquality},\\
        &\mathrm{NatFactory} \mapsto \\
        &\left\{
        \begin{aligned}
        &\mathrm{Successor} \mapsto \\
        &\left\{
        \begin{aligned}
        &\mathrm{Equal} \mapsto \\
        &\left\{
                  \begin{aligned}
                    &\mathrm{other} \mapsto \left\{
                      \begin{aligned}
                        &\qualifiedthis{\mathrm{NatEqualitySuccessor}}.\mathrm{NatFactory}.\mathrm{Product},\\
                        &\mathrm{predecessor} \mapsto \{\qualifiedthis{\mathrm{NatEqualitySuccessor}}.\mathrm{NatFactory}.\mathrm{Product}\}
                    \end{aligned}\right\},\\
                    &\mathrm{\_recursiveEquality} \mapsto \left\{
                      \begin{aligned}
                        &\qualifiedthis{\mathrm{Successor}}.\mathrm{predecessor}.\mathrm{Equal},\\
                        &\mathrm{other} \mapsto \{\qualifiedthis{\mathrm{Equal}}.\mathrm{other}.\mathrm{predecessor}\}
                    \end{aligned}\right\},\\
                    &\mathrm{\_OtherAcceptance} \mapsto \\
                    &\left\{
                      \begin{aligned}
                        &\mathrm{other}.\mathrm{Acceptance},\\
                        &\mathrm{VisitorMap} \mapsto \\
                        &\left\{
                          \begin{aligned}
                            &\mathrm{ZeroVisitor} \mapsto \{\mathrm{equal} \mapsto \{\qualifiedthis{\mathrm{NatEqualitySuccessor}}.\mathrm{BooleanFactory}.\mathrm{False}\}\},\\
                            &\mathrm{SuccessorVisitor} \mapsto \{\mathrm{equal} \mapsto \{\mathrm{\_recursiveEquality}.\mathrm{equal}\}\}
                        \end{aligned}\right\},\\
                        &\mathrm{Accepted} \mapsto \{\mathrm{equal} \mapsto \{\qualifiedthis{\mathrm{NatEqualitySuccessor}}.\mathrm{Boolean}\}\}
                    \end{aligned}\right\},\\
                    &\mathrm{equal} \mapsto \{\mathrm{\_OtherAcceptance}.\mathrm{Accepted}.\mathrm{equal}\}
                \end{aligned}\right\}
        \end{aligned}\right\}
        \end{aligned}\right\}
    \end{aligned}\right\}
    \end{aligned}
\]}
%% === 审稿人备忘：关于代码"冗长"的回应 ===
%%
%% NatEquality 的定义看起来比等价的 Haskell 或 ML 模式匹配更长，
%% 但这种冗长与 point-free / 组合子风格的"简洁"形成的对比是表面的。
%% 需要区分两种"长"：
%%
%% 1. 自描述式的长：每个中间步骤都有名字（如 _recursiveEquality、
%%    OtherAcceptance），步骤写得详细。名字本身是零成本的语义锚点。
%% 2. 组合子膨胀式的长：如 SKI bracket abstraction 的指数级膨胀，
%%    在玩具例子里看起来精炼，但在真实程序中复杂度是 λ-calculus 的
%%    多项式甚至指数倍，而且没有涌现出新的计算模式（论文 Related Work
%%    已指出这一点）。
%%
%% Inheritance-calculus 的"冗长"属于第一种，而且是结构上的必要代价：
%% 每个有名字的中间步骤既是代码，也是 mixin 树上可路径寻址、可扩展的
%% 节点。这正是 open extension 的前提——无法扩展匿名的子表达式。
%% 如果加语法糖把名字隐藏（如同 M-表达式之于 S-表达式），反而会
%% 破坏这一性质。Lisp 社区的历史经验表明，S-表达式的"冗长"最终
%% 被证明是优势（代码即数据，宏系统直接操作统一结构），M-表达式
%% 的语法糖从未被真正需要。
%%
%% 在 AI 辅助编程的语境下，这种显式命名的风格尤其有利：
%% _recursiveEquality 这样的名字对 LLM 而言是自带的 chain-of-thought，
%% 比起让 AI 猜测 point-free 管道中第 n 个组合子的意图，
%% 每个自描述的步骤都降低了生成和理解的认知负荷。
%% 考虑到未来越来越多的代码将由 AI 生成和维护，
%% 多写几个 token 的显式名称不是成本，而是投资。
%%
%% 当然，如果确实需要更紧凑的表面语法，为 inheritance-calculus 加一个
%% ANF 预处理器（自动为匿名子表达式生成名字）并不困难，
%% 正如 ANF 变换本身是机械化的。但这是表面语法的选择，
%% 不影响演算本身的语义。
%% ================================================================
\ifInheritanceChinese
通过继承来组合 $\mathrm{Plus}$ 与 $\mathrm{Equal}$ 的实现,无需修改任何 mixin,即可使所得的 Nat 值自动同时具备 $\mathrm{Plus}$ 与 $\mathrm{Equal}$。具体的数字常量存放在一个只继承两个数据构造子的共享 mixin 中:
\else
Composing the $\mathrm{Plus}$ and $\mathrm{Equal}$ implementations by
inheriting them, with no modifications to any mixin, gives the
resulting Nat values both $\mathrm{Plus}$ and $\mathrm{Equal}$
automatically.
The concrete numerals live in a shared mixin that inherits only the
two data constructors:
\fi
\[
  \mathrm{NatConstants} \mapsto \left\{
    \begin{aligned}
      &\mathrm{NatDataZero},\; \mathrm{NatDataSuccessor},\\
      &\mathrm{One} \mapsto \left\{
        \begin{aligned}
          &\qualifiedthis{\mathrm{NatConstants}}.\mathrm{NatFactory}.\mathrm{Successor},\\
          &\mathrm{predecessor} \mapsto \{\qualifiedthis{\mathrm{NatConstants}}.\mathrm{NatFactory}.\mathrm{Zero}\}
      \end{aligned}\right\},\\
      &\mathrm{Two} \mapsto \left\{
        \begin{aligned}
          &\qualifiedthis{\mathrm{NatConstants}}.\mathrm{NatFactory}.\mathrm{Successor},\\
          &\mathrm{predecessor} \mapsto \{\mathrm{One}\}
      \end{aligned}\right\},\\
      &\mathrm{Three} \mapsto \left\{
        \begin{aligned}
          &\qualifiedthis{\mathrm{NatConstants}}.\mathrm{NatFactory}.\mathrm{Successor},\\
          &\mathrm{predecessor} \mapsto \{\mathrm{Two}\}
      \end{aligned}\right\},\\
      &\mathrm{Four} \mapsto \left\{
        \begin{aligned}
          &\qualifiedthis{\mathrm{NatConstants}}.\mathrm{NatFactory}.\mathrm{Successor},\\
          &\mathrm{predecessor} \mapsto \{\mathrm{Three}\}
      \end{aligned}\right\},\\
      &\mathrm{Five} \mapsto \left\{
        \begin{aligned}
          &\qualifiedthis{\mathrm{NatConstants}}.\mathrm{NatFactory}.\mathrm{Successor},\\
          &\mathrm{predecessor} \mapsto \{\mathrm{Four}\}
      \end{aligned}\right\}
  \end{aligned}\right\}
\]
\ifInheritanceChinese
算术测试继承 $\mathrm{NatConstants}$ 以及它所需的操作:
\else
The arithmetic test inherits $\mathrm{NatConstants}$ together with the operations it needs:
\fi
\[
  \mathrm{Test} \mapsto \left\{
    \begin{aligned}
      &\mathrm{NatConstants},\; \mathrm{NatPlusZero},\; \mathrm{NatPlusSuccessor},\\
      &\mathrm{NatEqualityZero},\; \mathrm{NatEqualitySuccessor},\\
      &\mathrm{TwoPlusThree} \mapsto \left\{
        \begin{aligned}
          &\qualifiedthis{\mathrm{Test}}.\mathrm{Two}.\mathrm{Plus},\\
          &\mathrm{addend} \mapsto \{\qualifiedthis{\mathrm{Test}}.\mathrm{Three}\}
      \end{aligned}\right\},\\
      &\mathrm{FiveEqualTwoPlusThree} \mapsto \left\{
        \begin{aligned}
          &\qualifiedthis{\mathrm{Test}}.\mathrm{Five}.\mathrm{Equal},\\
          &\mathrm{other} \mapsto \{\mathrm{TwoPlusThree}.\mathrm{sum}\}
      \end{aligned}\right\}
  \end{aligned}\right\}
\]
\ifInheritanceChinese
路径 $\mathrm{Test}.\mathrm{FiveEqualTwoPlusThree}.\mathrm{equal}$ 继承 $\mathrm{True}$。%
\footnote{
  读取一个结果,需要借助 FFI 来打印,或者借助元语言观察以在该路径上调用 $\supers$。
}
没有任何现有 mixin 被修改。$\mathrm{NatEquality}$(一项新操作)与 $\mathrm{Boolean}$(一种新数据类型)各自只需要一个新 mixin:沿操作轴与情形轴均可扩展。
\else
The path $\mathrm{Test}.\mathrm{FiveEqualTwoPlusThree}.\mathrm{equal}$
inherits $\mathrm{True}$.%
\footnote{
  Reading a result requires either the FFI to print it or a
  metalanguage observation that invokes $\supers$ on the path.
}
No existing mixin was modified.
$\mathrm{NatEquality}$, a new operation, and $\mathrm{Boolean}$,
a new data type, each required only a new mixin: extensibility along both
the operation axis and the case axis.
\fi

\paragraph{\bilingual{Tree-level inheritance and return types}{树层次的继承与返回类型}}
\label{sec:discussion-expression-problem}
\ifInheritanceChinese
每个操作都是一个独立的 mixin,扩展工厂的子树。继承递归地合并共享同一标签的子树,因此工厂中的每个标签都从所有被继承的 mixin 获得全部操作。关键在于,这同样适用于返回类型:$\mathrm{NatPlus}$ 从 $\mathrm{NatFactory}$ 返回值,而 $\mathrm{NatEquality}$ 独立地向同一工厂添加 $\mathrm{Equal}$。同时继承两者,使得 $\mathrm{Plus}$ 返回的值自动携带 $\mathrm{Equal}$,无需修改任何一个 mixin。在对象代数~\cite{oliveira2012-object-algebras} 与 finally tagless 解释器~\cite{carette2009-finally-tagless} 中,独立定义的操作不会自动丰富彼此的返回类型;需要额外的样板代码或类型类机制。

多目标限定 $\this$(方程~\ref{eq:this})至关重要。在上述测试 mixin 中,$\mathrm{NatPlusSuccessor}$ 与 $\mathrm{NatEqualitySuccessor}$ 均继承 $\mathrm{NatDataSuccessor}$。当两个操作被组合时,所得的 $\mathrm{Successor}$ 通过两条独立路由继承 $\mathrm{NatDataSuccessor}$ 的模式。每个操作内部的引用 $\qualifiedthis{\mathrm{Successor}}.\mathrm{predecessor}$ 通过两条路由解析,由于继承是幂等的,所有路由贡献相同的属性。在 Scala 中,同样的模式,即两个独立定义的内部类扩展同一外部类的特质,会触发``冲突基类型''拒绝(附录~\ref{app:scala-multi-target}),因为 DOT 的 self 变量是单值的,无法同时通过多条继承路由解析。表达式问题在 Scala 中仍可使用对象代数或类似模式求解,但组合独立定义的操作本质上会产生上述对共享模式的多目标 $\this$ 解析。
\else
Each operation is an independent mixin that extends a factory's
subtree.
Inheritance recursively merges subtrees that share a label,
so every label in the factory acquires all operations from all
inherited mixins.
Crucially, this applies to return types: $\mathrm{NatPlus}$ returns
values from $\mathrm{NatFactory}$, and $\mathrm{NatEquality}$
independently adds $\mathrm{Equal}$ to the same factory.
Inheriting both causes the values returned by $\mathrm{Plus}$ to
carry $\mathrm{Equal}$ automatically, without modifying either mixin.
In object algebras~\cite{oliveira2012-object-algebras} and finally
tagless interpreters~\cite{carette2009-finally-tagless}, independently defined
operations do not automatically enrich each other's return types;
additional boilerplate or type-class machinery is required.

Multi-target qualified $\this$ (equation~\ref{eq:this}) is essential.
In the test mixin above, $\mathrm{NatPlusSuccessor}$ and
$\mathrm{NatEqualitySuccessor}$ both inherit
$\mathrm{NatDataSuccessor}$.
When the two operations are composed, the resulting
$\mathrm{Successor}$ inherits the $\mathrm{NatDataSuccessor}$ schema
through two independent routes.
The reference
$\qualifiedthis{\mathrm{Successor}}.\mathrm{predecessor}$
inside each operation resolves through both routes, and since
inheritance is idempotent, all routes contribute the same
properties.
In Scala, the same pattern, where two independently defined inner
classes extend the same outer class's trait, triggers a
``conflicting base types'' rejection
(Appendix~\ref{app:scala-multi-target}), because DOT's self variable
is single-valued and cannot resolve through multiple inheritance
routes simultaneously.
The Expression Problem can still be solved in Scala using
object algebras or similar patterns, but composing independently
defined operations inherently produces the multi-target $\this$
resolution to shared schemas just described.
\fi

\paragraph{\bilingual{Church-encoded Nats are tries}{Church 编码的 Nat 是字典树}}
\label{sec:case-study-cartesian}
\ifInheritanceChinese
Church 编码的 Nat 是一个记录,其结构映射了构建它所用的数据构造子路径:$\mathrm{Zero}$ 是一个平坦记录;$\mathrm{S}(\mathrm{Zero})$ 是一个带有 $\mathrm{predecessor}$ 属性、指向一个 $\mathrm{Zero}$ 记录的记录;依此类推。这恰好是字典树的形状。%
\footnote{
  一元 Nat 的字典树可以被视为一个稀疏链表,其深度等于它所代表的数值:在字典树并中,链表中的某些 $\mathrm{Successor}$ 节点指向 $\mathrm{Zero}$(叶节点),因此链表中并非每个位置都被占用。\anon[{} 补充材料]{{} MIXINv2 源码~\cite{mixin2025}}{} 中包含一个字典树深度为对数级的二进制自然数(BinNat)算术。
}
一个同时继承两个 Nat 值的作用域是一个字典树并。下面的测试将 $\mathrm{Plus}$ 应用于 $\{1,2\}$ 的字典树并,被加数为 $\{3,4\}$ 的字典树并,然后用 $\mathrm{Equal}$ 将结果与自身比较:
\else
A Church-encoded Nat is a record whose structure mirrors the
data-constructor path used to build it: $\mathrm{Zero}$ is a flat record;
$\mathrm{S}(\mathrm{Zero})$ is a record with a $\mathrm{predecessor}$
property pointing to a $\mathrm{Zero}$ record; and so on.
This is exactly the shape of a trie.%
\footnote{
  The unary Nat trie can be viewed as a sparse linked list
  whose depth equals the number it represents:
  in a trie union, some $\mathrm{Successor}$ nodes along
  the list point to $\mathrm{Zero}$ (leaf nodes),
  so not every position in the list is occupied.
  The\anon[{} supplementary material]{{} MIXINv2 source~\cite{mixin2025}}{} includes a binary natural number
  (BinNat) arithmetic whose trie depth is logarithmic.
}
A scope that simultaneously inherits two Nat values is a trie union.
The following test applies $\mathrm{Plus}$ to a trie union of $\{1,2\}$
with addend a trie union of $\{3,4\}$, then compares the result to itself
with $\mathrm{Equal}$:
\fi
\[
  \mathrm{CartesianTest} \mapsto \left\{
    \begin{aligned}
      &\mathrm{NatConstants},\; \mathrm{NatPlusZero},\; \mathrm{NatPlusSuccessor},\\
      &\mathrm{NatEqualityZero},\; \mathrm{NatEqualitySuccessor},\\
      &\mathrm{OneOrTwo} \mapsto \left\{\qualifiedthis{\mathrm{CartesianTest}}.\mathrm{One},\; \qualifiedthis{\mathrm{CartesianTest}}.\mathrm{Two}\right\},\\
      &\mathrm{ThreeOrFour} \mapsto \left\{\qualifiedthis{\mathrm{CartesianTest}}.\mathrm{Three},\; \qualifiedthis{\mathrm{CartesianTest}}.\mathrm{Four}\right\},\\
      &\mathrm{Result} \mapsto \left\{\mathrm{OneOrTwo}.\mathrm{Plus},\; \mathrm{addend} \mapsto \{\mathrm{ThreeOrFour}\}\right\},\\
      &\mathrm{Check} \mapsto \left\{\mathrm{Result}.\mathrm{sum}.\mathrm{Equal},\; \mathrm{other} \mapsto \{\mathrm{Result}.\mathrm{sum}\}\right\}
  \end{aligned}\right\}
\]
\ifInheritanceChinese
$\mathrm{Result}.\mathrm{sum}$ 求值为字典树 $\{4,5,6\}$,即所有两两之和的集合。$\mathrm{Equal}$ 随后对从 $\{4,5,6\} \times \{4,5,6\}$ 中取出的每一对进行应用;路径 $\mathrm{CartesianTest}.\mathrm{Check}.\mathrm{equal}$ 同时继承 $\mathrm{True}$ 与 $\mathrm{False}$。笛卡尔积直接来自语义方程。$\mathrm{OneOrTwo}$ 同时继承 $\mathrm{One}$ 与 $\mathrm{Two}$,因此 $\bases(\mathrm{OneOrTwo})$ 包含两条路径。当 $\mathrm{Result}$ 继承 $\mathrm{OneOrTwo}.\mathrm{Plus}$ 时,$\mathrm{Result}$ 的 $\bases$ 包含 $\mathrm{One}$ 与 $\mathrm{Two}$ 各自的 $\mathrm{Plus}$。每个 $\mathrm{Plus}$ 通过 $\this$ 解析其 $\mathrm{addend}$,到达 $\mathrm{ThreeOrFour}$;该值本身是 $\mathrm{Three}$ 与 $\mathrm{Four}$ 的字典树并,因此 $\mathrm{addend}$ 路径的 $\supers$ 包含两个值。由于 $\mathrm{Successor}.\mathrm{Plus}$ 以递增后的被加数递归委托给 $\mathrm{predecessor}.\mathrm{Plus}$,而 $\mathrm{predecessor}$ 也通过字典树并解析,递归在每个字典树节点处分支。因此结果 $\mathrm{sum}$ 收集了通过两个操作数的任意组合所能到达的全部路径,即笛卡尔积。

同样操作单个值的 $\mathrm{Plus}$ 与 $\mathrm{Equal}$ 文件,无需修改,即产生逻辑编程的关系语义:$\overrides$ 与 $\supers$ 收集所有继承路径,每个操作都在其上分布。
\else
$\mathrm{Result}.\mathrm{sum}$ evaluates to the trie $\{4,5,6\}$, the set of all pairwise sums.
$\mathrm{Equal}$ then applies to every pair drawn from $\{4,5,6\} \times \{4,5,6\}$;
the path $\mathrm{CartesianTest}.\mathrm{Check}.\mathrm{equal}$
inherits both $\mathrm{True}$ and $\mathrm{False}$.
The Cartesian product arises directly from the semantic equations.
$\mathrm{OneOrTwo}$ inherits both $\mathrm{One}$ and $\mathrm{Two}$,
so $\bases(\mathrm{OneOrTwo})$ contains two paths.
When $\mathrm{Result}$ inherits
$\mathrm{OneOrTwo}.\mathrm{Plus}$, the $\bases$ of
$\mathrm{Result}$ include the $\mathrm{Plus}$ of both
$\mathrm{One}$ and $\mathrm{Two}$.
Each $\mathrm{Plus}$ resolves its $\mathrm{addend}$ via
$\this$, reaching $\mathrm{ThreeOrFour}$, which is itself a
trie union of $\mathrm{Three}$ and $\mathrm{Four}$:
the $\supers$ of the $\mathrm{addend}$ path contain both values.
Since $\mathrm{Successor}.\mathrm{Plus}$ recursively delegates to
$\mathrm{predecessor}.\mathrm{Plus}$ with an incremented addend,
and $\mathrm{predecessor}$ also resolves through the trie union,
the recursion branches at every trie node.
The result $\mathrm{sum}$ therefore collects all paths reachable
through any combination of the two operands, which is the
Cartesian product.

The same $\mathrm{Plus}$ and $\mathrm{Equal}$ files that operate on
single values thus produce, without modification, the relational
semantics of logic programming:
$\overrides$ and $\supers$ collect all inheritance paths, and
every operation distributes over them.
\fi

\section{\bilingual{Discussion}{讨论}}

\subsection{\bilingual{Emergent Phenomena}{涌现现象}}
\label{sec:emergent-phenomena}

\ifInheritanceChinese
在 MIXINv2 中开发真实程序的过程中，我们发现了三种现象，每一种都是对某种已有设计模式的非预期用法，超出了该模式的原始能力。

\begin{description}
  \item[关系语义]
    第~\ref{sec:case-study}~节案例研究中的笛卡尔积之所以出现，是因为字典树是一种集合结构，因此对字典树的深度合并即为集合并，而对合并后字典树的操作则遍历其操作数的笛卡尔积。
    这并不局限于 $\mathrm{Nat}$：任何 Datalog 元组都可编码为链表，而深度合并在这些元组上计算集合操作，我们认为这正是 Datalog 语义。
    附录~\ref{app:datalog-encoding}~将 Datalog 编码进继承演算，仅用第~\ref{sec:syntax}~节的三个基本构造即可表达变量连接、常量匹配、多体规则和递归规则。
    这种逻辑编程语义是对编码案例研究算术的工厂模式的非预期用法中涌现出来的。

  \item[对不可扩展性的免疫]\label{item:immunity}
    在第~\ref{sec:case-study}~节的案例研究中，$\mathrm{Nat}$ 编码被证明是无类型表达问题~\cite{wadler1998-expression-problem}的一个实例：%
    \footnote{
      表达问题有两个要求：沿操作轴和数据构造子轴均可扩展，以及静态类型安全。无类型演算无法满足类型安全要求，因此我们只考虑双轴可扩展性要求。
    }
    一个新操作和一个新数据类型各自作为独立的 mixin 添加，沿操作轴和情形轴均进行了扩展，且无需修改任何已有 mixin。
    在其他场景中，访问者模式允许添加新操作但不允许添加新情形；而这里两者都允许，因为第~\ref{sec:syntax}~节的每个构造都可通过深度合并（公式~\ref{eq:overrides}）进行扩展，所以不可扩展的那一半根本无法写出；第~\ref{sec:asymmetry}~节（定理~\ref{thm:cartesian-nonexpressibility}）则形式化了使这一结论无条件成立的不对称性。
    在其他地方恢复第二个扩展维度需要专门的机制~\cite{liang1995-monad-transformers, lammel2003-scrap-your-boilerplate, loh2006-open-data-types,
      swierstra2008-data-types-a-la-carte, carette2009-finally-tagless, oliveira2012-object-algebras,
      kiselyov2013-extensible-effects, wu2014-effect-handlers-in-scope, kiselyov2015-freer-monads,
    wang2016-expression-problem-trivially, leijen2017-row-typed-algebraic-effects, poulsen2023-hefty-algebras}。
    第~\ref{sec:semantic-variants}~节分析了为何深度合并是唯一能实现这一点的组合机制。
    这第二个可扩展性维度是对访问者模式的非预期用法中涌现出来的。

\end{description}

\label{sec:practical-function-color}%
更多现象依赖于 FFI，而本文中的示例排除了 FFI。其中一种是函数颜色盲~\cite{nystrom2015-function-color}：同一个 MIXINv2 文件可以在同步和异步运行时下运行，这是通过继承进行依赖注入的非预期用法。补充材料\anon[]{~\cite{mixin2025}}给出了相关示例。
\else
Developing real programs in MIXINv2 surfaced three phenomena, each an
off-label use of an adopted design pattern beyond its original abilities.

\begin{description}
  \item[Relational semantics]
    The Cartesian product in the case study of
    Section~\ref{sec:case-study} arose because a trie is a set
    structure, so deep merge over tries is set union and an operation
    over the merged trie ranges over the product of its operands.
    This is not specific to $\mathrm{Nat}$: any Datalog tuple encodes
    as a linked list, and deep merge computes set operations over such
    tuples, which we believe is Datalog semantics.
    Appendix~\ref{app:datalog-encoding} encodes Datalog into
    inheritance-calculus, expressing variable joins,
    constant matches, multi-body rules, and recursive rules with
    the three primitives of Section~\ref{sec:syntax} alone.
    This logic-programming semantics emerged from an off-label use of
    the factory pattern that encodes the case study's arithmetic.

  \item[Immunity to nonextensibility]\label{item:immunity}
    In the case study of Section~\ref{sec:case-study}, the
    $\mathrm{Nat}$ encoding turned out to be an instance of the
    untyped Expression
    Problem~\cite{wadler1998-expression-problem}:%
    \footnote{
      The Expression Problem has two requirements: extensibility along
      both the operation and the data-constructor axis, and static
      type safety. An untyped calculus cannot meet the type-safety
      requirement, so we consider only the two-axis extensibility
      requirement.
    }
    a new operation and a new data type were each added as a separate
    mixin, extending along both the operation axis and the case axis,
    with no existing mixin modified.
    Elsewhere the visitor pattern admits new operations but not new
    cases; here it admits both, because every construct of
    Section~\ref{sec:syntax} is extensible via deep merge
    (equation~\ref{eq:overrides}), so the nonextensible half cannot be
    written, and Section~\ref{sec:asymmetry}
    (Theorem~\ref{thm:cartesian-nonexpressibility}) formalizes the
    asymmetry that makes this unconditional.
    Recovering that second dimension elsewhere takes dedicated
    mechanisms~\cite{liang1995-monad-transformers, lammel2003-scrap-your-boilerplate, loh2006-open-data-types,
      swierstra2008-data-types-a-la-carte, carette2009-finally-tagless, oliveira2012-object-algebras,
      kiselyov2013-extensible-effects, wu2014-effect-handlers-in-scope, kiselyov2015-freer-monads,
    wang2016-expression-problem-trivially, leijen2017-row-typed-algebraic-effects, poulsen2023-hefty-algebras}.
    Section~\ref{sec:semantic-variants} analyzes why deep merge
    is the only composition mechanism that achieves this.
    This second extensibility dimension emerged from an off-label use
    of the visitor pattern.

\end{description}

\label{sec:practical-function-color}%
Further phenomena rely on the FFI, which the examples in this paper
exclude. One is function color
blindness~\cite{nystrom2015-function-color}: the same MIXINv2 file runs
under both synchronous and asynchronous runtimes, an off-label use of
dependency injection via inheritance. The
supplementary material\anon[]{~\cite{mixin2025}} gives the example.
\fi

\subsection{\bilingual{Semantic Alternatives}{语义替代方案}}
\label{sec:semantic-variants}

\ifInheritanceChinese
上述\hyperref[item:immunity]{对不可扩展性的免疫}条目确立了继承演算对不可扩展性的免疫性。我们将此作为设计目标，并通过调查我们已知的所有候选方案，询问是否有替代设计能共享这一性质。

\paragraph{\bilingual{Degrees of freedom}{自由度}}
以下分析假设 AST 是一棵\emph{命名树}：从父节点到子节点的每条边都携带一个标签，因此每个节点都由从根节点出发的路径唯一标识。对于配置语言，这自然成立：JSON、YAML 或 Nix 属性集中的每个键都是一个标签，而嵌套产生路径。%
\footnote{
  Nix 是一门具有一等函数的函数式语言；仅凭其属性集并不构成命名树 DSL。然而，NixOS 模块系统~\cite{nixosmodules}将用户界面限制为通过递归合并组合的嵌套属性集，而正是这种 DSL 层面的结构自然地构成了命名树。
}
对于主流编程语言，AST 并不直接是命名树，因为匿名表达式位置和隐式作用域缺乏标签。然而，经过 ANF 变换~\cite{flanagan1993-essence-compiling-continuations}后，每个中间结果都绑定到一个名字，而剩余的匿名位置%
\footnote{
  例如，$\mathbf{let}$ 的主体。
}
可以被分配合成的新鲜标签，例如第~\ref{sec:forward-translation}~节中的 $\mathrm{result}$ 和 $\mathrm{tailCall}$，从而得到一棵命名树。

在此前提下，每种程序是带引用的命名树的语言，都可以由两个有限集来刻画：一个由\emph{节点类型}构成的集合 $\mathcal{N}$，每种节点类型决定子树如何组织其子节点；以及一个由\emph{引用类型}构成的集合 $\mathcal{R}$，每种引用类型决定树中一个位置如何引用另一个位置。语法由 $\mathcal{N}$ 和 $\mathcal{R}$ 的选择固定；语义则由每种节点类型如何构造组合以及每种引用类型如何解析其目标来决定。

{\small
  \begin{center}
    \begin{tabular}{lp{0.38\columnwidth}p{0.33\columnwidth}}
      \bilingual{\textbf{Language}}{\textbf{语言}}
      & \bilingual{\textbf{Node types $\mathcal{N}$}}{\textbf{节点类型 $\mathcal{N}$}}
      & \bilingual{\textbf{Reference types $\mathcal{R}$}}{\textbf{引用类型 $\mathcal{R}$}} \\
      \hline
      Java (OOP)
      & \bilingual{class, interface, method,
      block, conditional, loop,
      \texttt{throw}, lambda, \ldots}{类、接口、方法、块、条件、循环、\texttt{throw}、lambda、\ldots}
      & \bilingual{variable, field access, method call,
      \texttt{extends}/\texttt{implements},
      \texttt{import}, \ldots}{变量、字段访问、方法调用、\texttt{extends}/\texttt{implements}、\texttt{import}、\ldots} \\[3pt]
      Haskell (FP)
      & \bilingual{$\lambda$, application, \texttt{let},
      \texttt{case}, \texttt{data},
      type class, instance, \ldots}{$\lambda$、应用、\texttt{let}、\texttt{case}、\texttt{data}、类型类、实例、\ldots}
      & \bilingual{variable, constructor,
      type class dispatch,
      \texttt{import}, \ldots}{变量、构造子、类型类分派、\texttt{import}、\ldots} \\[3pt]
      Scala\footnote{
        \bilingual{Scala is a multi-paradigm language.}{Scala 是一门多范式语言。}
      }
      & \bilingual{class, trait, object, method,
      block, conditional, \texttt{match},
      \texttt{while}, \texttt{throw}, \ldots}{类、trait、对象、方法、块、条件、\texttt{match}、\texttt{while}、\texttt{throw}、\ldots}
      & \bilingual{variable, member access,
      \texttt{extends}, \texttt{with},
      implicits, \texttt{import}, \ldots}{变量、成员访问、\texttt{extends}、\texttt{with}、隐式参数、\texttt{import}、\ldots} \\
      \hline
      $\lambda$-\bilingual{calculus}{演算}
      & \bilingual{abstraction, application}{抽象、应用}
      & \bilingual{variable}{变量} \\
      \bilingual{Inheritance-calculus}{继承演算}
      & \bilingual{record}{记录}
      & \bilingual{inheritance}{继承} \\
    \end{tabular}
  \end{center}
}

\paragraph{\bilingual{Immunity to nonextensibility}{对不可扩展性的免疫}}
程序 $P$ 中的路径 $p$ 是\emph{可扩展的}，若存在程序 $Q$，使得通过引用将 $P$ 与 $Q$ 组合后，在不修改 $P$ 的前提下，$p$ 处能够观察到新的属性。一种语言\emph{对不可扩展性免疫}，若其每个良构程序的每条路径都是可扩展的。

我们现在对设计空间进行分类，以解释为何我们选择第~\ref{sec:mixin-trees}~节的语义。两个维度保持自由：\emph{组合机制}，决定单一引用类型如何合并定义；以及\emph{引用解析}，决定单一引用类型如何找到其目标。

\subsubsection*{\bilingual{The composition mechanism}{组合机制}}

每种组合机制决定了一个抽象的定义如何被合并进另一个抽象。下表按各候选方案是否满足可交换性~(C)、幂等性~(I) 和可结合性~(A)，以及其可扩展性模型对已知候选方案进行分类。

\begin{center}
  \begin{tabular}{lcccll}
    \bilingual{\textbf{Candidate}}{\textbf{候选方案}} & \textbf{C} & \textbf{I} & \textbf{A}
    & \bilingual{\textbf{Extensibility}}{\textbf{可扩展性}} & \bilingual{\textbf{Representative}}{\textbf{代表}} \\
    \hline
    \bilingual{Deep merge}{深度合并}
    & $\checkmark$ & $\checkmark$ & $\checkmark$
    & \bilingual{universal}{通用} & \bilingual{this paper}{本文} \\
    \bilingual{Shallow merge}{浅层合并}
    & $\times$ & $\checkmark$ & $\checkmark$
    & \bilingual{top-level only}{仅顶层} & JS spread \\
    \bilingual{Conflict rejection}{冲突拒绝}
    & $\checkmark$ & $\checkmark$ & $\checkmark$
    & \bilingual{none on overlap}{重叠时无扩展}
    & Harper--Pierce~\cite{harper1991-symmetric-record-concatenation} \\
    \bilingual{Priority merge}{优先级合并}
    & $\times$ & $\times$ & $\checkmark$
    & \bilingual{universal}{通用} & Jsonnet~\texttt{+}, Nix~\texttt{//} \\
    $\beta$-\bilingual{reduction}{归约}
    & $\times$ & $\times$ & $\times$
    & \bilingual{opt-in}{按需选择} & Java, Haskell, Scala, \ldots \\
    \bilingual{Conditional}{条件并入}
    & \multicolumn{3}{c}{\bilingual{causes divergence}{导致发散}}
    & --- & --- \\
  \end{tabular}
\end{center}

\begin{description}
  \item[\bilingual{Deep merge}{深度合并}]
    对同名定义的合并递归进入其子树（$\overrides$，公式~\ref{eq:overrides}）。每个标签都是一个扩展点，无论原作者是否预期了扩展。

  \item[\bilingual{Shallow merge}{浅层合并}]
    嵌套标签无法被独立扩展。

  \item[\bilingual{Conflict rejection}{冲突拒绝}]
    同名定义被作为错误拒绝。三个等式性质均成立，但为空洞成立：这些等式得到满足，是因为用于验证它们的组合被拒绝了。冲突拒绝阻止了对已有标签的扩展~\cite{harper1991-symmetric-record-concatenation}，而这恰恰是表达问题~\cite{wadler1998-expression-problem}所要求的可组合性。

  \item[\bilingual{Priority merge}{优先级合并}]
    一个来源具有优先权~\cite{jsonnet,dolstra2010-nixos}，因此 $A + B \neq B + A$，引入了顺序依赖。

  \item[$\beta$-\bilingual{reduction}{归约}]
    只有作者放置的参数才是扩展点；函数体是封闭的。要使一个新的方面可扩展，原函数必须被重写以接受一个额外参数。任何基于函数应用的组合机制都继承了这一限制，因为 $\beta$-归约使函数体封闭。表达问题~\cite{wadler1998-expression-problem}在此类语言中之所以困难，原因相同：可扩展性是按需选择的。模拟对不可扩展性免疫的框架（访问者模式、对象代数~\cite{oliveira2012-object-algebras}、finally tagless 解释器~\cite{carette2009-finally-tagless}）正是为绕过这一限制而生的机制。

  \item[\bilingual{Conditional incorporation}{条件并入}]
    根据其他定义的存在与否来移除或有条件地包含定义，会产生依赖循环：某定义是否存在，取决于另一定义是否存在，而后者又取决于前者。当添加和移除交替出现时，递归求值（附录~\ref{app:well-definedness}）无法到达不动点，从而导致发散。这与图灵完备性固有的发散不同：图灵完备程序发散是因为计算无界，而条件并入发散是因为直接后果算子 $T_P$ 振荡而非单调递增。
\end{description}

\noindent
在上述候选方案中，深度合并是我们找到的唯一一个可交换、幂等、可结合且对不可扩展性免疫的方案。其他候选方案各自引入了对不可扩展性的脆弱性，或导致发散。

\subsubsection*{\bilingual{The reference resolution}{引用解析}}

一旦将记录上的深度合并确定为组合机制，六个语义方程中的四个便由语法决定：$\properties$~(\ref{eq:properties}) 读取记录的标签，$\bases$~(\ref{eq:bases}) 读取其继承来源，$\supers$~(\ref{eq:supers}) 取传递闭包，$\overrides$~(\ref{eq:overrides}) 实现合并。唯一剩余的自由度是引用的解析方式：$\resolve$~(\ref{eq:resolve}) 和 $\this$~(\ref{eq:this}) 方程。我们考察三个已知候选方案。

\begin{description}
  \item[\bilingual{Early binding}{早期绑定}]
    使用早期绑定时，引用在定义时即解析到一个固定路径，与记录后来如何被继承无关。CUE~\cite{cue2019}是这方面的典型：字段引用被静态解析，嵌套字段无法引用其封闭上下文中由后续组合贡献的属性。这使得 mixin 无法充当 $\lambda$-演算的调用栈：$\beta$-归约要求被调用方能观察到调用点提供的参数，而那是一次后续组合。使用早期绑定时，引用在此组合发生前就已被冻结。在方程中，$\resolve$~(\ref{eq:resolve}) 调用 $\this$~(\ref{eq:this})，后者遍历继承的树结构；迟绑定正是使 $\this$ 能够观察到后续组合的 mixin 所贡献属性的原因。若没有图灵完备的 mixin，函数必须被重新引入为组合机制，从而再次暴露上述按需选择的可扩展性问题。

  \item[\bilingual{Dynamic scope}{动态作用域}]
    使用动态作用域时，引用在使用点而非定义点处解析。Jsonnet~\cite{jsonnet}采用此方式：\texttt{self} 始终指向最终合并后的对象，而非定义点处的对象。联合文件系统在文件系统层面表现出相同的性质：符号链接的相对路径在解引用时相对于挂载上下文解析，而非相对于定义该链接的层。动态作用域破坏了第~\ref{sec:forward-translation}~节的 $\lambda$-演算嵌入。替换引理（引理~\ref{lem:substitution}）依赖于基本情况 $M = y$（$y \neq x$）：$y$ 的 de~Bruijn 指标 $m$ 从定义点路径 $\init(p_{\mathrm{def}})$ 向上走到一个与 $x$ 的绑定器不同的作用域层级，因此 $x$ 作用域处的替换 $\mathrm{argument} \mapsto \mathcal{T}(V)$ 不产生任何效果。动态作用域下，这种分离不再成立。%
    \footnote{
      反例：$(\lambda x.\, \lambda y.\, x)\; V_1\; V_2$。$\lambda y.\, x$ 内部对 $x$ 的引用翻译为 $\dbi{1}.\mathrm{argument}$（de~Bruijn 指标 $1$）。使用迟绑定时，$\this$ 从定义点作用域（内层 $\lambda$）出发，向上走一层到外层 $\lambda$，到达槽 $\mathrm{argument} \mapsto \mathcal{T}(V_1)$。使用动态作用域时，引用从使用点上下文解析，其中最内层封闭作用域是应用 $\{C_1.\mathrm{result},\; \mathrm{argument} \mapsto \mathcal{T}(V_2)\}$；向上走一层到达该作用域的 $\mathrm{argument}$ 槽，返回 $\mathcal{T}(V_2)$ 而非 $\mathcal{T}(V_1)$。这是动态作用域的经典变量捕获失败。
    }
    与早期绑定一样，mixin 无法完全替代函数。

  \item[\bilingual{Single-target $\this$}{单目标 $\this$}]
    单目标 $\this$ 是一个有前途的候选方案：单路径引理~(\ref{lem:single-path})表明，对于翻译后的 $\lambda$-项，$\this$~(\ref{eq:this}) 中的前沿集 $S$ 始终恰好包含一条路径，因此 $\lambda$-演算嵌入不需要多目标解析。然而，单目标 $\this$ 因另一原因而对不可扩展性脆弱。当多条继承路径到达同一作用域时，单目标解析必须拒绝该情况或选择其中一条路径。NixOS 模块系统~\cite{nixosmodules}会引发重复声明错误（附录~\ref{app:scala-multi-target}）。这拒绝了两个独立编写的模块定义同一选项的嵌套扩展，而这种组合在求解表达问题~\cite{wadler1998-expression-problem}时会出现，如第~\ref{sec:expression-problem}~节所示。

    我们对 MIXINv2 的前三个实现使用了 $\this$ 解析的标准技术：遵循 Cook~\cite{cook1989-denotational-inheritance}的基于闭包的不动点，以及使用 de~Bruijn 指标~\cite{debruijn1972-nameless-dummies}的基于栈的环境查找。当继承引入到同一封闭作用域的多条路径时，三种实现都产生了隐晦的错误。困难在于结构上。闭包捕获单一环境，因此不动点组合子只能为单一 $\this$ 目标求解。基于栈的环境是一条线性链，每个作用域层级只有一个父节点。这两种数据结构都内嵌了一个与多重继承自然产生的多目标情况不相容的单值假设（附录~\ref{app:scala-multi-target}）。公式~(\ref{eq:this})是在放弃这一假设后涌现出来的：它跟踪一个\emph{前沿集} $S$，而非单一当前作用域。
\end{description}

\noindent
下表总结了两个层面上所有替代方案的脆弱性。

\begin{center}
  \begin{tabular}{llll}
    & \bilingual{\textbf{Candidate}}{\textbf{候选方案}} & \bilingual{\textbf{Vulnerability}}{\textbf{脆弱性}}
    & \bilingual{\textbf{Representative}}{\textbf{代表}} \\
    \hline
    \multirow{5}{*}{\rotatebox[origin=c]{90}{\scriptsize \bilingual{composition}{组合}}}
    & \bilingual{Shallow merge}{浅层合并}
    & \bilingual{nested extensions not composable}{嵌套扩展不可组合}
    & JS spread \\
    & \bilingual{Conflict rejection}{冲突拒绝}
    & \bilingual{same-label extensions rejected}{同名扩展被拒绝}
    & Harper--Pierce~\cite{harper1991-symmetric-record-concatenation} \\
    & \bilingual{Priority merge}{优先级合并}
    & \bilingual{ordering dependency}{顺序依赖}
    & Jsonnet~\texttt{+}~\cite{jsonnet},
    Nix~\texttt{//}~\cite{dolstra2010-nixos} \\
    & $\beta$-\bilingual{reduction}{归约}
    & \bilingual{opt-in extensibility}{按需选择的可扩展性}
    & Java, Haskell, Scala, \ldots \\
    & \bilingual{Conditional}{条件并入}
    & \bilingual{causes divergence}{导致发散}
    & --- \\
    \hline
    \multirow{3}{*}{\rotatebox[origin=c]{90}{\scriptsize \bilingual{resolution}{解析}}}
    & \bilingual{Early binding}{早期绑定}
    & \bilingual{mixin cannot replace function}{mixin 无法替代函数}
    & CUE~\cite{cue2019} \\
    & \bilingual{Dynamic scope}{动态作用域}
    & \bilingual{mixin cannot replace function}{mixin 无法替代函数}
    & Jsonnet~\cite{jsonnet}, union FS \\
    & \bilingual{Single-target $\this$}{单目标 $\this$}
    & \bilingual{nested extensions rejected}{嵌套扩展被拒绝}
    & NixOS modules~\cite{nixosmodules} \\
  \end{tabular}
\end{center}

\noindent
我们在两个层面考察的所有替代方案，均对不可扩展性不免疫。

我们通过尝试本节中的每个替代方案并发现其失败，从而得到了这一语义：浅层合并无法组合嵌套扩展；优先级合并引入了顺序依赖；早期绑定在组合能够提供引用目标之前就将引用冻结；动态作用域破坏了 $\lambda$-演算嵌入所依赖的替换；单目标 $\this$ 拒绝了表达问题所要求的多目标组合。深度合并加上多目标迟绑定 $\this$，正是最终剩下的选择。
\else
The \hyperref[item:immunity]{immunity-to-nonextensibility item} above established that inheritance-calculus is immune
to nonextensibility.
We adopt this as a design goal and ask whether an alternative
design could share this property
by surveying every candidate known to us.

\paragraph{\bilingual{Degrees of freedom}{自由度}}
The analysis below assumes that the AST is a \emph{named tree}:%
every edge from parent to child carries a label, so that every
node is uniquely identified by the path from the root.
For configuration languages this is naturally the case: every
key in a JSON, YAML, or Nix attribute set is a label, and nesting
produces paths.%
\footnote{
  Nix is a functional language with first-class functions;
  its attribute sets alone do not form a named-tree DSL.
  The NixOS module system~\cite{nixosmodules}, however,
  restricts the user-facing interface to nested attribute sets
  composed by recursive merging, and it is this DSL-level
  structure that naturally forms a named tree.
}
For mainstream programming languages the AST is not directly a
named tree, since anonymous expression positions and implicit
scopes lack labels.
After ANF
transformation~\cite{flanagan1993-essence-compiling-continuations}, however, every intermediate
result is bound to a name, and the remaining anonymous positions%
\footnote{
  For example, the body of a $\mathbf{let}$.
}
can be assigned synthetic
fresh labels such as $\mathrm{result}$ and $\mathrm{tailCall}$
in Section~\ref{sec:forward-translation}, yielding a named
tree.

Under this premise, every language whose programs are such
named trees with references can be characterized by two finite sets:
a set~$\mathcal{N}$ of \emph{node types}, each determining how a
subtree organizes its children, and a set~$\mathcal{R}$ of
\emph{reference types}, each determining how one position in the
tree refers to another.
The syntax is fixed by the choice of $\mathcal{N}$ and
$\mathcal{R}$; the semantics is determined by how each node type
structures composition and how each reference type resolves its
target.

{\small
  \begin{center}
    \begin{tabular}{lp{0.38\columnwidth}p{0.33\columnwidth}}
      \textbf{Language}
      & \textbf{Node types $\mathcal{N}$}
      & \textbf{Reference types $\mathcal{R}$} \\
      \hline
      Java (OOP)
      & class, interface, method,
      block, conditional, loop,
      \texttt{throw}, lambda, \ldots
      & variable, field access, method call,
      \texttt{extends}/\texttt{implements},
      \texttt{import}, \ldots \\[3pt]
      Haskell (FP)
      & $\lambda$, application, \texttt{let},
      \texttt{case}, \texttt{data},
      type class, instance, \ldots
      & variable, constructor,
      type class dispatch,
      \texttt{import}, \ldots \\[3pt]
      Scala\footnote{
        Scala is a multi-paradigm language.
      }
      & class, trait, object, method,
      block, conditional, \texttt{match},
      \texttt{while}, \texttt{throw}, \ldots
      & variable, member access,
      \texttt{extends}, \texttt{with},
      implicits, \texttt{import}, \ldots \\
      \hline
      $\lambda$-calculus
      & abstraction, application
      & variable \\
      Inheritance-calculus
      & record
      & inheritance \\
    \end{tabular}
  \end{center}
}

\paragraph{\bilingual{Immunity to nonextensibility}{对不可扩展性的免疫}}
A path~$p$ in a program~$P$ is \emph{extensible} if there exists
a program~$Q$ such that composing $P$ with~$Q$ through a reference
makes new properties observable at~$p$ without modifying~$P$.
A language is \emph{immune to nonextensibility} if every path in
every well-formed program is extensible.

We now classify the design space to explain why we chose the
semantics of Section~\ref{sec:mixin-trees}.
Two axes remain free:
the \emph{composition mechanism}, which determines how the single reference
type incorporates definitions;
and the \emph{reference resolution}, which determines how the single reference
type finds its target.

\subsubsection*{\bilingual{The composition mechanism}{组合机制}}

Every composition mechanism determines how the definitions of one
abstraction are incorporated into another.
The following table classifies the known candidates by whether
they satisfy commutativity~(C), idempotence~(I), and
associativity~(A), and by their extensibility model.

\begin{center}
  \begin{tabular}{lcccll}
    \bilingual{\textbf{Candidate}}{\textbf{候选方案}} & \textbf{C} & \textbf{I} & \textbf{A}
    & \bilingual{\textbf{Extensibility}}{\textbf{可扩展性}} & \bilingual{\textbf{Representative}}{\textbf{代表}} \\
    \hline
    \bilingual{Deep merge}{深度合并}
    & $\checkmark$ & $\checkmark$ & $\checkmark$
    & \bilingual{universal}{通用} & \bilingual{this paper}{本文} \\
    \bilingual{Shallow merge}{浅层合并}
    & $\times$ & $\checkmark$ & $\checkmark$
    & \bilingual{top-level only}{仅顶层} & JS spread \\
    \bilingual{Conflict rejection}{冲突拒绝}
    & $\checkmark$ & $\checkmark$ & $\checkmark$
    & \bilingual{none on overlap}{重叠时无扩展}
    & Harper--Pierce~\cite{harper1991-symmetric-record-concatenation} \\
    \bilingual{Priority merge}{优先级合并}
    & $\times$ & $\times$ & $\checkmark$
    & \bilingual{universal}{通用} & Jsonnet~\texttt{+}, Nix~\texttt{//} \\
    $\beta$-\bilingual{reduction}{归约}
    & $\times$ & $\times$ & $\times$
    & \bilingual{opt-in}{按需选择} & Java, Haskell, Scala, \ldots \\
    \bilingual{Conditional}{条件并入}
    & \multicolumn{3}{c}{\bilingual{causes divergence}{导致发散}}
    & --- & --- \\
  \end{tabular}
\end{center}

\begin{description}
  \item[Deep merge]
    Merging same-label definitions recurses into their subtrees
    ($\overrides$, equation~\ref{eq:overrides}).
    Every label is an extension point, whether or not the
    original author anticipated extension.

  \item[Shallow merge]
    Nested labels cannot be independently extended.

  \item[Conflict rejection]
    Same-label definitions are rejected as errors.
    All three identities hold, but vacuously:
    the identities are satisfied because the compositions that
    would test them are refused.
    Conflict rejection prevents extending existing
    labels~\cite{harper1991-symmetric-record-concatenation}, precisely the composability that the
    Expression Problem~\cite{wadler1998-expression-problem} requires.

  \item[Priority merge]
    One source takes
    precedence~\cite{jsonnet,dolstra2010-nixos},
    so $A + B \neq B + A$, introducing
    an ordering dependency.

  \item[$\beta$-reduction]
    Only the parameters placed by the author are extension
    points; the function body is closed.
    To make a new aspect extensible, the original function must
    be rewritten to accept an additional parameter.
    Any composition mechanism based on function application
    inherits this limitation, because $\beta$-reduction
    closes the function body.
    The Expression Problem~\cite{wadler1998-expression-problem} is hard
    in such languages for the same reason:
    extensibility is opt-in.
    The frameworks that simulate immunity to nonextensibility
    (visitor pattern,
      object algebras~\cite{oliveira2012-object-algebras},
    finally tagless interpreters~\cite{carette2009-finally-tagless})
    are the machinery needed to work around this limitation.

  \item[Conditional incorporation]
    Removing or conditionally including definitions
    based on the presence of other definitions
    creates a dependency cycle:
    whether a definition exists depends on
    whether another definition exists, which in turn
    depends on the first.
    The recursive evaluation
    (Appendix~\ref{app:well-definedness}) cannot
    reach a fixed point when addition and removal
    alternate, causing divergence.
    This differs from the divergence inherent to
    Turing completeness: Turing-complete programs
    diverge because computation is unbounded, whereas
    conditional incorporation diverges because the
    immediate consequence operator $T_P$ oscillates
    rather than ascending monotonically.
\end{description}

\noindent
Among the candidates above, deep merge is the only one we found
that is commutative, idempotent, associative, and
immune to nonextensibility.
The other candidates each introduce a vulnerability to
nonextensibility or cause divergence.

\subsubsection*{\bilingual{The reference resolution}{引用解析}}

Once deep merge over records is fixed as the composition
mechanism, four of the six semantic equations are determined
by the syntax:
$\properties$~(\ref{eq:properties}) reads the labels of a
record, $\bases$~(\ref{eq:bases}) reads its inheritance
sources, $\supers$~(\ref{eq:supers}) takes the transitive
closure, and $\overrides$~(\ref{eq:overrides}) implements
the merge.
The only remaining degree of freedom is how references
resolve: the $\resolve$~(\ref{eq:resolve}) and
$\this$~(\ref{eq:this}) equations.
We examine the three known candidates.

\begin{description}
  \item[Early binding]
    With early binding, references resolve at definition time to a
    fixed path, independently of how the record is later inherited.
    CUE~\cite{cue2019} exemplifies this: field references are
    statically resolved, and a nested field cannot refer to
    properties of its enclosing context contributed by later
    composition.
    This prevents the mixin from serving as a call stack for the
    $\lambda$-calculus: $\beta$-reduction requires the callee to
    observe the argument supplied at the call site, which is a
    later composition.
    With early binding, references are frozen before this
    composition occurs.
    In the equations,
    $\resolve$~(\ref{eq:resolve}) calls
    $\this$~(\ref{eq:this}), which walks the inherited
    tree structure; late binding is what allows
    $\this$ to see properties contributed by later-composed
    mixins.
    Without a Turing-complete mixin, functions must be
    reintroduced as the composition mechanism, re-exposing
    the opt-in extensibility problem above.

  \item[Dynamic scope]
    With dynamic scope, references resolve at the point of use
    rather than the definition site.
    Jsonnet~\cite{jsonnet} follows this approach: \texttt{self}
    always refers to the final merged object, not the object at
    the definition site.
    Union file systems exhibit the same property at the
    file-system level: a symbolic link's relative path resolves
    against the mount context at dereference time, not against the
    layer that defined the link.
    Dynamic scope breaks the $\lambda$-calculus embedding
    of Section~\ref{sec:forward-translation}.
    The Substitution Lemma (Lemma~\ref{lem:substitution})
    relies on the base case $M = y$ ($y \neq x$):
    the de~Bruijn index $m$ of~$y$ walks up from the
    definition-site path $\init(p_{\mathrm{def}})$ to a scope
    level different from~$x$'s binder, so the substitution
    $\mathrm{argument} \mapsto \mathcal{T}(V)$ at~$x$'s scope
    has no effect.
    With dynamic scope this separation fails.%
    \footnote{
      Counterexample:
      $(\lambda x.\, \lambda y.\, x)\; V_1\; V_2$.
      The reference to~$x$ inside $\lambda y.\, x$ translates to
      $\dbi{1}.\mathrm{argument}$ (de~Bruijn index~$1$).
      With late binding, $\this$ starts from the definition-site
      scope (the inner $\lambda$) and walks up one level to the
      outer $\lambda$, reaching the slot
      $\mathrm{argument} \mapsto \mathcal{T}(V_1)$.
      With dynamic scope, the reference resolves from the
      use-site context, where the innermost enclosing scope is
      the application $\{C_1.\mathrm{result},\;
      \mathrm{argument} \mapsto \mathcal{T}(V_2)\}$;
      walking up one level reaches this scope's
      $\mathrm{argument}$ slot, returning $\mathcal{T}(V_2)$
      instead of $\mathcal{T}(V_1)$.
      This is the classical variable-capture failure of dynamic
      scoping.
    }
    As with early binding, the mixin cannot serve as a complete
    replacement for functions.

  \item[Single-target $\this$]
    Single-target $\this$ is a promising candidate:
    the Single-Path Lemma~(\ref{lem:single-path}) shows that
    for translated $\lambda$-terms, the frontier set~$S$ in
    $\this$~(\ref{eq:this}) always contains exactly one path,
    so the $\lambda$-calculus embedding does not require
    multi-target resolution.
    However, single-target $\this$ is vulnerable to
    nonextensibility for a different reason.
    When multiple inheritance routes reach the same scope,
    single-target resolution must reject the situation or select
    one route.
    The NixOS module system~\cite{nixosmodules} raises a
    duplicate-declaration error
    (Appendix~\ref{app:scala-multi-target}).
    This rejects nested extensions where two independently authored
    modules define the same option, a composition that arises
    when solving the Expression
    Problem~\cite{wadler1998-expression-problem},
    as Section~\ref{sec:expression-problem} demonstrates.

    Our first three implementations of MIXINv2 used
    standard techniques for $\this$ resolution: closure-based
    fixed points following Cook~\cite{cook1989-denotational-inheritance}, and
    stack-based environment lookup using de~Bruijn
    indices~\cite{debruijn1972-nameless-dummies}.
    All three produced subtle bugs when inheritance introduced
    multiple routes to the same enclosing scope.
    The difficulty was structural.
    A closure captures a single environment, so a fixed-point
    combinator solves for a single $\this$ target.
    A stack-based environment is a linear chain where each scope
    level has one parent.
    Both data structures embed a single-valued assumption
    incompatible with the multi-target situation that multiple
    inheritance naturally produces
    (Appendix~\ref{app:scala-multi-target}).
    Equation~(\ref{eq:this}) emerged from abandoning this
    assumption: it tracks a \emph{frontier set}~$S$ rather than a
    single current scope.
\end{description}

\noindent
The following table summarizes the vulnerabilities of all
alternatives at both levels.

\begin{center}
  \begin{tabular}{llll}
    & \bilingual{\textbf{Candidate}}{\textbf{候选方案}} & \bilingual{\textbf{Vulnerability}}{\textbf{脆弱性}}
    & \bilingual{\textbf{Representative}}{\textbf{代表}} \\
    \hline
    \multirow{5}{*}{\rotatebox[origin=c]{90}{\scriptsize \bilingual{composition}{组合}}}
    & \bilingual{Shallow merge}{浅层合并}
    & \bilingual{nested extensions not composable}{嵌套扩展不可组合}
    & JS spread \\
    & \bilingual{Conflict rejection}{冲突拒绝}
    & \bilingual{same-label extensions rejected}{同名扩展被拒绝}
    & Harper--Pierce~\cite{harper1991-symmetric-record-concatenation} \\
    & \bilingual{Priority merge}{优先级合并}
    & \bilingual{ordering dependency}{顺序依赖}
    & Jsonnet~\texttt{+}~\cite{jsonnet},
    Nix~\texttt{//}~\cite{dolstra2010-nixos} \\
    & $\beta$-\bilingual{reduction}{归约}
    & \bilingual{opt-in extensibility}{按需选择的可扩展性}
    & Java, Haskell, Scala, \ldots \\
    & \bilingual{Conditional}{条件并入}
    & \bilingual{causes divergence}{导致发散}
    & --- \\
    \hline
    \multirow{3}{*}{\rotatebox[origin=c]{90}{\scriptsize \bilingual{resolution}{解析}}}
    & \bilingual{Early binding}{早期绑定}
    & \bilingual{mixin cannot replace function}{mixin 无法替代函数}
    & CUE~\cite{cue2019} \\
    & \bilingual{Dynamic scope}{动态作用域}
    & \bilingual{mixin cannot replace function}{mixin 无法替代函数}
    & Jsonnet~\cite{jsonnet}, union FS \\
    & \bilingual{Single-target $\this$}{单目标 $\this$}
    & \bilingual{nested extensions rejected}{嵌套扩展被拒绝}
    & NixOS modules~\cite{nixosmodules} \\
  \end{tabular}
\end{center}

\noindent
No alternative we examined at either level is immune to
nonextensibility.

We arrived at this semantics by trying each alternative in this section
and finding that it failed:
shallow merge could not compose nested extensions;
priority merge introduced ordering dependencies;
early binding froze references before composition could
supply them;
dynamic scope broke the substitution that the
$\lambda$-calculus embedding relies on;
single-target $\this$ rejected the multi-target compositions
that the Expression Problem demands.
Deep merge with multi-target late-binding $\this$ is what
remained.
\fi

\section{\bilingual{Related Work}{相关工作}}
\label{sec:related-work}

\subsection{\bilingual{Computational Models and \texorpdfstring{$\lambda$}{λ}-Calculus Semantics}{计算模型与 \texorpdfstring{$\lambda$}{λ}-演算语义}}

\paragraph{\bilingual{Minimal Turing-complete models}{最小图灵完备模型}}
\ifInheritanceChinese
Sch\"onfinkel~\cite{schonfinkel1924-combinatory-logic} 与 Curry~\cite{curry1958-combinatory-logic} 证明了 $\lambda$-演算可以归约为三个组合子 $S$、$K$、$I$%
\footnote{
  组合子 $I$ 是冗余的,因为 $I = SKK$。
}。将 $\lambda$-项编译为 SKI 的括号抽象是机械的,但会导致项大小的指数级膨胀;在 $\lambda$-演算中已有的计算模式之外,不会涌现出任何新的计算模式。在语义上,两者共享相同的模型:Meyer~\cite{meyer1982-model-of-lambda-calculus} 证明了满足五个附加条件的组合代数恰好就是 $\lambda$-代数。从未有人展示过一个独立的 SKI 语义域,并从中推导出 $\lambda$-演算的语义。方向始终是反过来的:Scott 的 $D_\infty$ 是为 $\lambda$-演算而构建的,而 SKI 随后被解释于其中。Dolan~\cite{dolan2013-mov-turing-complete} 证明了 x86 的 \texttt{mov} 指令单独就已图灵完备,其方法是利用寻址模式作为可变内存上的隐式算术;由此得到的程序是正确的,但不提供任何超越 RAM 机的抽象机制。图灵机本身虽然是通用的,却缺乏随机访问:读取距离为 $d$ 处的格需要 $d$ 次顺序的磁带移动。这些模型中的每一个,都在某个已有模型的语义框架内实现了图灵完备性,且对其中任何一个都从未展示过独立的语义域。
\else
Sch\"onfinkel~\cite{schonfinkel1924-combinatory-logic} and
Curry~\cite{curry1958-combinatory-logic} showed that the
$\lambda$-calculus reduces to three combinators $S$, $K$, $I$
\footnote{
  The $I$ combinator is redundant since $I = SKK$.
}.
The bracket abstraction that compiles $\lambda$-terms to SKI
is mechanical but incurs exponential blow-up in term size;
no new computational patterns emerge beyond those already
present in the $\lambda$-calculus.
Semantically, the two share the same models:
Meyer~\cite{meyer1982-model-of-lambda-calculus} showed that combinatory algebras
satisfying five additional conditions are exactly lambda
algebras.
No independent semantic domain for SKI has ever been
exhibited from which $\lambda$-calculus semantics could be
derived. The direction has always been the reverse:
Scott's $D_\infty$ was built for the $\lambda$-calculus,
and SKI was interpreted inside it afterwards.
Dolan~\cite{dolan2013-mov-turing-complete} showed that the x86 \texttt{mov}
instruction alone is Turing complete by exploiting
addressing modes as implicit arithmetic on mutable memory;
the resulting programs are correct but provide no
abstraction mechanisms beyond those of the RAM Machine.
The Turing Machine itself, while universal, lacks random
access: reading a cell at distance~$d$ requires $d$
sequential tape movements.
Each of these models achieves Turing completeness
within the semantic framework of an existing model,
and no independent semantic domain has been exhibited
for any of them.
\fi

%% === 审稿人备忘：SF-calculus 与 inheritance-calculus 的关系 ===
%%
%% Barry Jay 的 SF-calculus（Intensional Computation with Higher-Order
%% Functions, TCS 2019, doi:10.1016/j.tcs.2019.02.016；前身为
%% Confusion in the Church-Turing Thesis, arXiv:1410.7103, 2014）
%% 和 inheritance-calculus 都试图解决 λ-calculus 的 extensionality 限制，
%% 但方向完全不同。
%%
%% 共同出发点：λ-calculus 是 extensional 的——只能观察函数的输入输出行为，
%% 无法检查程序的内部结构。两个内部结构不同但行为相同的 λ 项不可区分。
%%   - Jay 的诊断：λ-calculus 缺乏 intensional 操作（检查项的结构），
%%     无法定义结构等式判定。
%%   - 本文的诊断：λ-calculus 的 function body 是 closed 的，无法从外部
%%     添加新的可观察投影（open extension）。
%%
%% 核心机制对比：
%%   SF-calculus：保留 S、K（或 λ），添加 F（factorisation）算子，
%%     可以将闭合范式拆解为组成部分。程序仍然是函数，只是多了检查函数
%%     结构的能力。属于 λ-calculus 的 *扩展*。
%%   Inheritance-calculus：用三个不同的原语（record、definition、
%%     inheritance）替代 λ-calculus 的三个原语。β-reduction 是 deep merge
%%     的特例（Section 3 的翻译 T 将函数应用映射为继承）。
%%     属于 λ-calculus 的 *包含*（subsumption）。
%%
%% 关于 F 算子的编码：
%%   F 做两件事：(1) 判断结构形态（atom vs. compound），
%%   (2) 提取子部分。
%%   在 inheritance-calculus 中：
%%     - 提取子部分是平凡的：所有值都是透明 record，通过 qualified this
%%       投影任意 slot 即可。Theorem 6 的 separating context 就在读取
%%       argument slot（↑².argument）。
%%     - 判断结构形态并分支：需要 Acceptance pattern（Section 4）。
%%       NatVisitor 对 Nat 做的事情与 F 对 SF-calculus 项做的事情
%%       完全对应：VisitZero/SuccessorVisitor 对应 F 的 atom/compound
%%       分支。Visitor 可以通过 open extension 添加到已有 term 上，
%%       不修改原始定义。
%%   Acceptance 在精神上就是 F：两者都是对 closed normal form 做
%%   constructor dispatch。区别在于 F 是不可扩展的 built-in primitive
%%   （硬编码了 S 和 F 两个 atom），而 Acceptance 是可扩展的 encoding
%%   （新增 constructor 只需继承新的 Accepted* 分支，即 Expression Problem
%%   immunity）。
%%
%% 精确差异：
%%   SF-calculus 把 intensional observation 做成了 object-level 的
%%   primitive reduction rule（F 是一个 term）；
%%   inheritance-calculus 把它留在了 meta-level（观察者总能通过
%%   properties(p) 查询任何路径），而 term-level 的分支需要通过
%%   Acceptance pattern 显式编码。
%%   换言之：inheritance-calculus 中 F 的 *读取* 能力是免费的
%%   （所有值都是透明 record），但 F 的 *分支* 能力需要 Acceptance。
%%   SF-calculus 的 F 则把读取和分支捆绑在一个不可扩展的 primitive 里。
%%   Acceptance 是可扩展的（新增 constructor 只需继承新分支，
%%   即 Expression Problem immunity），而 F 硬编码了固定的 atom 集合
%%   （S 和 F）。
%%
%% 关于 Jay 所谓的 "trivial solution of the Halting Problem"：
%%   Jay 用非标准编码将每个程序（包括不停机的）映射到闭合范式（静态数据），
%%   然后结构等式判定变得平凡。这并不推翻图灵定理——它改变了问题的
%%   形式化方式（把程序当静态数据而非可执行代码），而非解决了原来的问题。
%%   Jay 的真正论点是：λ-definability 和 Turing computability 不是同一回事，
%%   Church-Turing Thesis 的等价性证明隐含了编码假设。
%%
%% 正文篇幅不够论述 embeds SF-calculus，暂不展开。
%% 如果审稿人提到 SF-calculus，可以基于上述分析撰写回复或补充段落。
%% ================================================================

\paragraph{\bilingual{Semantic descriptions of the $\lambda$-calculus}{$\lambda$-演算的语义描述}}
\ifInheritanceChinese
$\lambda$-演算的语义可以用几种根本上不同的方式来定义。\emph{操作语义}将 $\beta$-归约 $(\lambda x.\,M)\,N \to M[N/x]$ 作为计算的基本概念~\cite{church1936-lambda-calculus,church1941-calculi-of-lambda-conversion,barendregt1984-lambda-calculus}。这要求捕获-回避替换,其微妙之处促使了 de~Bruijn 索引~\cite{debruijn1972-nameless-dummies} 与显式替换演算~\cite{abadi1991-explicit-substitutions} 的出现。它还要求选择求值策略:按名调用与按值调用~\cite{plotkin1975-call-by-name-call-by-value},或 Abramsky 的惰性策略~\cite{abramsky1990-lazy-lambda-calculus}。\emph{指称语义}~\cite{scott1970-mathematical-theory-of-computation,scott1976-data-types-as-lattices} 将项解释于一个数学域中,但赋予 $\beta$-归约意义需要一个自反域 $D \cong [D \to D]$,由 Cantor 定理知该域无集合论解。Scott 的 $D_\infty$ 通过将 $D$ 构造为连续格的逆极限~\cite{scott1972-continuous-lattices} 来解决这一问题。由此得到的语义是适当的,但不是完全抽象的~\cite{plotkin1977-lcf-as-programming-language,milner1977-fully-abstract-models}:它将观察等价所区分的项等同了起来。\emph{博弈语义}~\cite{abramsky2000-full-abstraction-pcf,hyland2000-full-abstraction-pcf,nickau1994-hereditarily-sequential-functionals} 对 PCF 与惰性 $\lambda$-演算~\cite{abramsky1993-full-abstraction-lazy-lambda,abramsky1995-games-full-abstraction-lazy-lambda} 实现了完全抽象,需要竞技场、策略与无害性条件。在这些方法的每一种中,函数空间 $[D \to D]$ 都是核心障碍:操作语义通过替换隐式地需要它,而指称语义则显式地构造它,博弈语义则通过策略加以精化。这些方法中的每一种都涉及高阶结构%
\footnote{
  指称语义中的函数空间,博弈语义中的策略,以及操作语义中的替换。
}
并在其语义域中保留了函数空间。在每种情形下,指称帐和操作帐都是相互独立的形式体系。三者都继承了 $\beta$-归约在递归程序上的发散:例如,有向图上一个朴素的传递闭包在每个框架下都会发散。
\else
The semantics of the $\lambda$-calculus can be defined in
several fundamentally different ways.
\emph{Operational semantics} takes $\beta$-reduction
$(\lambda x.\,M)\,N \to M[N/x]$ as the primitive notion of
computation~\cite{church1936-lambda-calculus,church1941-calculi-of-lambda-conversion,barendregt1984-lambda-calculus}.
This requires capture-avoiding substitution, a mechanism whose
subtleties motivated de~Bruijn
indices~\cite{debruijn1972-nameless-dummies} and explicit substitution
calculi~\cite{abadi1991-explicit-substitutions}.
It also requires a choice of evaluation
strategy: call-by-name versus
call-by-value~\cite{plotkin1975-call-by-name-call-by-value} or the lazy
strategy of Abramsky~\cite{abramsky1990-lazy-lambda-calculus}.
\emph{Denotational semantics}~\cite{scott1970-mathematical-theory-of-computation,scott1976-data-types-as-lattices}
interprets terms in a mathematical domain, but giving meaning
to $\beta$-reduction requires a reflexive domain
$D \cong [D \to D]$, which has no set-theoretic solution by
Cantor's theorem. Scott's $D_\infty$ resolves this by
constructing $D$ as an inverse limit of continuous
lattices~\cite{scott1972-continuous-lattices}.
The resulting semantics is adequate but not fully
abstract~\cite{plotkin1977-lcf-as-programming-language,milner1977-fully-abstract-models}:
it identifies terms that observational equivalence
distinguishes.
\emph{Game semantics}~\cite{abramsky2000-full-abstraction-pcf,hyland2000-full-abstraction-pcf,nickau1994-hereditarily-sequential-functionals}
achieves full abstraction for PCF and for the lazy
$\lambda$-calculus~\cite{abramsky1993-full-abstraction-lazy-lambda,abramsky1995-games-full-abstraction-lazy-lambda},
requiring arenas, strategies, and innocence conditions.
In each of these approaches, the function space $[D \to D]$
is the central obstacle: operational semantics needs it
implicitly via substitution. In contrast, denotational semantics
constructs it explicitly, and game semantics refines it via
strategies.
Each of these approaches involves higher-order structure%
\footnote{
  Function spaces
  in denotational semantics, strategies
  in game semantics, and substitution
  in operational semantics.
}
and retains the function space in its semantic domain.
The denotational and operational accounts remain separate
formalisms in each case.
All three inherit $\beta$-reduction's divergence on
recursive programs: a na\"ive transitive closure on a cyclic
graph, for instance, diverges in every framework.
\fi

\paragraph{\bilingual{Set-theoretic models of the $\lambda$-calculus}{$\lambda$-演算的集合论模型}}
\ifInheritanceChinese
$\lambda$-演算的几个模型族通过在集合论而非域论框架下工作,避免了 Scott 的 $D_\infty$ 的逆极限构造。这些模型消除了 Scott 拓扑,但并未消除函数空间本身:它们仍然是高阶的,通过注入、交叉类型或多重集关系将 $[D \to D]$ 编码于 $D$ 之中。\emph{图模型}~\cite{engeler1981-algebras-and-combinators,plotkin1993-set-theoretic-lambda-models} 通过幂集格上的 web 注入来编码函数;Bucciarelli 与 Salibra~\cite{bucciarelli2008-graph-lambda-theories} 证明了最大的合理图理论等于 Böhm 树理论 $\mathcal{B}$。\emph{过滤器模型}~\cite{barendregt1983-filter-lambda-model,coppo1978-intersection-type-assignment} 将一个项解释为可从中推导出的交叉类型集合;Dezani-Ciancaglini 等人~\cite{dezani1998-intersection-types-bohm-trees} 证明了 BCD 过滤器模型的 $\lambda$-理论恰好与 Böhm 树相等性重合。\emph{关系模型}~\cite{ehrhard2005-finiteness-spaces,bucciarelli2001-phase-semantics-exponentials} 通过有限多重集从线性逻辑的关系语义中涌现;Ehrhard~\cite{ehrhard2005-finiteness-spaces} 证明了在这一范畴中,纯 $\lambda$-演算的标准自反对象不能存在。这三个模型族都是语义模型:它们将 $\lambda$-项解释于一个外部数学结构中,并在该结构中将 $\beta$-归约验证为一个等式,而不是定义一个具有自身原语的独立演算。每一种都编码函数空间而非消除它,且每一种都需要高阶机制才能为惰性 $\lambda$-演算获得完全抽象。
\else
Several families of $\lambda$-calculus models avoid the
inverse-limit construction of Scott's $D_\infty$ by working
in set-theoretic rather than domain-theoretic settings.
These models eliminate Scott topology but not the function
space itself: they remain higher-order, encoding $[D \to D]$
inside $D$ via injections, intersection types, or multiset
relations.
\emph{Graph models}~\cite{engeler1981-algebras-and-combinators,plotkin1993-set-theoretic-lambda-models}
encode functions via a web injection on powerset lattices;
Bucciarelli and Salibra~\cite{bucciarelli2008-graph-lambda-theories}
proved that the greatest sensible graph theory equals the
B\"ohm tree theory~$\mathcal{B}$.
\emph{Filter models}~\cite{barendregt1983-filter-lambda-model,coppo1978-intersection-type-assignment}
interpret a term as the set of intersection types derivable
for it;
Dezani-Ciancaglini et al.~\cite{dezani1998-intersection-types-bohm-trees}
showed that the BCD filter model's $\lambda$-theory
coincides exactly with B\"ohm tree equality.
\emph{Relational models}~\cite{ehrhard2005-finiteness-spaces,bucciarelli2001-phase-semantics-exponentials}
arise from the relational semantics of linear logic via
finite multisets;
Ehrhard~\cite{ehrhard2005-finiteness-spaces} showed that a standard
reflexive object for the pure $\lambda$-calculus cannot
exist in this category.
All three families are semantic models: they interpret
$\lambda$-terms in an external mathematical structure and
validate $\beta$-reduction as an equation in that structure,
rather than defining an independent calculus with its own
primitives.
Each encodes the function space rather than eliminating it,
and each requires higher-order machinery to obtain full
abstraction for the lazy $\lambda$-calculus.
\fi

\paragraph{\bilingual{Recursive equations and fixpoint semantics}{递归方程与不动点语义}}
\ifInheritanceChinese
Van~Emden 与 Kowalski~\cite{vanemden1976-predicate-logic-semantics} 证明了谓词逻辑作为编程语言的语义允许三种等价的刻画:操作性的\footnote{通常是 SLD 归结。}、模型论的\footnote{基于最小 Herbrand 模型。}以及不动点的\footnote{即基原子幂集上即时后承算子 $T_P$ 的最小不动点。}。Aczel~\cite{aczel1977-inductive-definitions} 通过幂集格上的单调算子对归纳定义给出了系统性处理,为 Datalog 风格的语义所依赖的逻辑基础奠定了基础。Leroy 与 Grall~\cite{leroy2006-coinductive-big-step} 通过余归纳定义处理了大步语义中的发散,需要在归纳求值判断之外引入一个单独的余归纳判断来说明非终止。Van~Emden 与 Kowalski 的 $T_P$ 在一阶域上为递归 Datalog 程序计算 $\mathrm{lfp}$。传统的 $\lambda$-演算语义在语义域中需要函数空间,且递归的 $\lambda$-项在 $\beta$-归约下通常会发散;利用一阶域上的 $T_P$ 为 $\lambda$-演算赋予一种改变哪些项收敛的不动点语义,尚未被探索。在近期工作中,Yang~\cite{yang2026-co-lambda} 将 tabling 求值~\cite{tamaki1986-tabled-resolution} 直接引入纯 $\lambda$-演算:首部规范化被读作一个余代数方程,其状态为被驻留的 $\lambda$-项,即具有可判定同一性的一阶对象,而 tabling 求解该方程,在有限时间内将 $\Omega$ 这样的非生产性循环判定为 $\bot$。该处理保留了标准的惰性 (Lévy--Longo) 树语义,仅改变了求值器;它占据了图~\ref{fig:calculi-matrix} 中不动点 $\lambda$ 节点。它在机制上也有所不同。表以整个被驻留的项为键,而非以进入作用域树的路径为键,因此在路径前缀上重叠的查询之间不共享部分结果。偏序不是集合包含下的幂集格,而是树近似序,其中 $\bot$ 叶被精化为子树,且不动点是唯一的(最小与最大因有护性而重合):一个重入调用被以单一的最终 $\bot$ 回答,而后续任何迭代都不会精化该 $\bot$,因此迭代无法增大一个累加器。其比惰性 $\lambda$-演算更确定的部分因此恰好是图~\ref{fig:calculi-matrix} 中非生产性的 $\bot$,即发散被判定为确定 $\bot$ 答案的那些项,别无其他;而幂集序才使得迭代可以积累并产生新的非 $\bot$ 值。
\else
Van~Emden and Kowalski~\cite{vanemden1976-predicate-logic-semantics} showed that
the semantics of predicate logic as a programming language admits
three equivalent characterizations:
operational\footnote{
  Typically SLD resolution.
}, model-theoretic\footnote{
  Based on least Herbrand models.
}, and fixed-point\footnote{
  Namely, the least fixed point of the immediate consequence operator~$T_P$ on the powerset of ground atoms.
}.
Aczel~\cite{aczel1977-inductive-definitions} gave a systematic treatment of
inductive definitions via monotone operators on powerset lattices,
providing the logical foundations that Datalog-style semantics
rests on.
Leroy and Grall~\cite{leroy2006-coinductive-big-step} treated divergence
in big-step semantics via coinductive definitions,
requiring a separate coinductive judgment to account for
non-termination alongside the inductive evaluation judgment.
Van~Emden and Kowalski's $T_P$ computes
$\mathrm{lfp}$ for recursive Datalog programs over a
first-order domain.
Conventional $\lambda$-calculus semantics requires function
spaces in the semantic domain, and recursive $\lambda$-terms
ordinarily diverge under $\beta$-reduction;
using $T_P$ over a first-order domain to give the
$\lambda$-calculus a fixpoint semantics that changes which
terms converge has not been explored.
In recent work, Yang~\cite{yang2026-co-lambda} brings tabled
evaluation~\cite{tamaki1986-tabled-resolution} to the pure
$\lambda$-calculus directly: head normalisation is read as a
coalgebraic equation whose states are interned
$\lambda$-terms, first-order objects with decidable identity,
and tabling solves it, deciding unproductive loops such
as~$\Omega$ as~$\bot$ in finite time.
That treatment preserves the standard lazy
(L\'evy--Longo) tree semantics and changes only the evaluator;
it occupies the fixpoint~$\lambda$ node of
Figure~\ref{fig:calculi-matrix}.
It also differs in mechanism.
The table is keyed by whole interned terms, not by paths
into a scope tree, so partial results are not shared between
queries that overlap on a path prefix.
The order is not a powerset lattice under set inclusion but
the tree approximation order, in which $\bot$ leaves are
refined into subtrees, and the fixpoint is unique (least and
greatest coincide by guardedness): a re-entrant call is
answered with a single final~$\bot$ that no later iteration
refines, so the iteration cannot grow an accumulator.
Its more-definedness over the lazy $\lambda$-calculus is
therefore exactly the unproductive~$\bot$ of
Figure~\ref{fig:calculi-matrix}, terms whose divergence is
decided as a definite~$\bot$ answer, and nothing more;
the powerset order is what lets iteration accumulate and
produce new non-$\bot$ values.
\fi

\subsection{\bilingual{Inheritance, Objects, and Records}{继承、对象与记录}}

\paragraph{\bilingual{Denotational semantics of inheritance}{继承的指称语义}}
\ifInheritanceChinese
Cook~\cite{cook1989-denotational-inheritance} 给出了继承的第一个指称语义,将对象建模为递归记录。继承是\emph{生成器}%
\footnote{
  这些是从 self 到完整对象的函数。
}与\emph{包装器}%
\footnote{
  包装器是修改生成器的函数。
}的复合,由一个不动点构造来求解。Cook 的关键洞见是,继承是一种可适用于任何形式的递归定义的通用机制,而不仅仅是面向对象的方法。这一洞见是本工作的出发点之一。该模型有四个原语(记录、函数、生成器和包装器),并需要一个不动点组合子。复合是不对称的:包装器修改生成器,反之则不然。不动点构造依赖函数域中的 Scott 连续性。由此得到的指称不能直接执行,因此需要一个单独的操作语义来在运行时定义方法分派与对象构造。
\else
Cook~\cite{cook1989-denotational-inheritance} gave the first denotational
semantics of inheritance, modeling objects as recursive records.
Inheritance is composition of \emph{generators}
\footnote{
  These are functions from self to complete object.
} and
\emph{wrappers}\footnote{
  Wrappers are functions that modify generators.
},
solved by a fixed-point construction.
Cook's key insight is that inheritance is a general mechanism
applicable to any form of recursive definition, not only
object-oriented methods. This insight is one of the starting points of the
present work.
The model has four primitives (records, functions, generators, and
wrappers) and requires a fixed-point combinator. Composition is
asymmetric: a wrapper modifies a generator but not vice versa.
The fixed-point construction relies on Scott continuity in a
domain of functions. The resulting denotations are not directly
executable, so a separate operational semantics is needed to
define method dispatch and object construction at runtime.
\fi

\paragraph{\bilingual{Mixins and traits}{Mixin 与 trait}}
\ifInheritanceChinese
Bracha 与 Cook~\cite{bracha1990-mixin-inheritance} 将 mixin 形式化为抽象子类,即从超类参数到子类的函数。由于 mixin 应用是函数复合,它\emph{既不可交换也不幂等};C3 线性化算法~\cite{barrett1996-c3-linearization} 后来被开发出来以施加一个确定性顺序。Schärli 等人~\cite{scharli2003-traits} 引入了 trait,其复合是可交换的,但拒绝同名冲突~\cite{ducasse2006-traits-fine-grained-reuse}。两种机制都在\emph{方法层级}操作:对同名的两个定义进行复合会导致覆盖或冲突,而非嵌套结构的递归合并。两者都预设了 $\lambda$-演算:方法体是函数,而复合机制本身不是图灵完备的。Mixin 与 trait 均不支持深合并。
\else
Bracha and Cook~\cite{bracha1990-mixin-inheritance} formalized mixins as
abstract subclasses, that is, functions from a superclass parameter to a
subclass.
Because mixin application is function composition, it is
\emph{neither commutative nor idempotent}; the C3 linearization
algorithm~\cite{barrett1996-c3-linearization} was later developed to impose a
deterministic order.
Sch\"arli et al.~\cite{scharli2003-traits} introduced traits,
whose composition is commutative but rejects same-name
conflicts~\cite{ducasse2006-traits-fine-grained-reuse}.
Both mechanisms operate at the \emph{method level}: composing two
definitions of the same name results in an override or a conflict,
not a recursive merge of nested structure.
Both presuppose the $\lambda$-calculus: method bodies are functions,
and the composition mechanism itself is not Turing complete.
Neither mixins nor traits are deep-mergeable.
\fi

\paragraph{\bilingual{Object and record calculi}{对象演算与记录演算}}
\ifInheritanceChinese
Abadi 与 Cardelli~\cite{abadi1996-theory-of-objects} 将对象作为基本概念,但方法是以 self 为参数的函数,且对象扩展是不对称的。Boudol~\cite{boudol2004-recursive-record-semantics} 证明了自引用记录需要一个不安全的不动点算子,其适定性依赖于求值策略。Harper 与 Pierce~\cite{harper1991-symmetric-record-concatenation} 给出了一个具有对称连接的记录演算,该演算拒绝同标签的复合;Cardelli~\cite{cardelli1992-extensible-records} 研究了带子类型的可扩展记录;Rémy~\cite{remy1989-row-polymorphism} 为 ML 中的记录和变体引入了行多态。这些演算全都在\emph{平坦}记录上操作:没有嵌套结构的递归合并,没有惰性观察,$\this$ 也需要一个显式的不动点组合子。
\else
Abadi and Cardelli~\cite{abadi1996-theory-of-objects} treat objects as the
primitive notion, but methods are functions that take self as a
parameter and object extension is asymmetric.
Boudol~\cite{boudol2004-recursive-record-semantics} showed that self-referential
records require an unsafe fixed-point operator whose
well-definedness depends on the evaluation strategy.
Harper and Pierce~\cite{harper1991-symmetric-record-concatenation} gave a record
calculus with symmetric concatenation that rejects same-label
composition;
Cardelli~\cite{cardelli1992-extensible-records} studied extensible records
with subtyping;
R\'emy~\cite{remy1989-row-polymorphism} introduced row polymorphism for
records and variants in ML.
All of these calculi operate on \emph{flat} records: there is
no recursive merging of nested structure, no lazy observation,
and $\this$ requires an explicit fixed-point combinator.
\fi

\paragraph{\bilingual{Family polymorphism and DOT}{族多态与 DOT}}
\ifInheritanceChinese
Ernst~\cite{ernst2001-family-polymorphism} 引入了族多态;Ernst、Ostermann 与 Cook~\cite{ernst2006-virtual-class-calculus} 在虚类演算中将其形式化,其中类通过路径表达式 \texttt{this.out.C} 来访问。Amin 等人~\cite{amin2016-dependent-object-types} 在 DOT 演算中形式化了路径依赖类型,这是 Scala 的理论基础。在这两个框架中,每条路径都解析为\emph{单一}封闭对象中的\emph{单一}类。当多条继承路径通向同一作用域时,单值解析必须选择其中一条路径并丢弃其余;DOT 的 self 变量 $\nu(x : T)\, d$ 将 $x$ 绑定到单一对象,而 Scala 在编译时拒绝由此产生的交叉类型(附录~\ref{app:scala-multi-target})。两个框架均未提供集合值解析的机制。
\else
Ernst~\cite{ernst2001-family-polymorphism} introduced family polymorphism;
Ernst, Ostermann, and Cook~\cite{ernst2006-virtual-class-calculus}
formalized it in the virtual class calculus, where classes are
accessed via path expressions \texttt{this.out.C}.
Amin et al.~\cite{amin2016-dependent-object-types} formalized path-dependent types
in the DOT calculus, the theoretical foundation of Scala.
In both frameworks, each path resolves to a
\emph{single} class in a \emph{single} enclosing object.
When multiple inheritance routes lead to the same scope,
single-valued resolution must choose one route and discard the
others; DOT's self variable $\nu(x : T)\, d$ binds $x$ to a
single object, and Scala rejects the resulting intersection
type at compile time
(Appendix~\ref{app:scala-multi-target}).
Neither framework provides a mechanism for set-valued
resolution.
\fi

\subsection{\bilingual{Configuration Languages and Module Systems}{配置语言与模块系统}}

\paragraph{\bilingual{Configuration languages}{配置语言}}
\ifInheritanceChinese
CUE~\cite{cue2019} 在单一格中统一了类型与值,其中以 \texttt{\&} 表示的合一是可交换、可结合且幂等的。CUE 的设计受到 Google 的 Borg 配置语言 (GCL) 的启发,后者使用了图合一。CUE 的格包含标量类型,其冲突语义规定合一不相容标量会得到 $\bot$;标量对于配置语言是否必要,还是仅凭树结构就已足够,是一个有趣的设计问题。Jsonnet~\cite{jsonnet} 通过 \texttt{+} 运算符提供了具有 mixin 语义的对象继承:复合\emph{不是}可交换的,因为在标量冲突上右侧胜出,而深合并需要在每个字段上通过 \texttt{+:} 语法显式选择加入,因此深合并虽可用但并非默认行为。Dhall~\cite{dhall2017} 采取函数式方法:它是一个带记录的类型化 $\lambda$-演算,其中复合是记录合并算子,而非继承机制。
\else
CUE~\cite{cue2019} unifies types and values in a single lattice
where unification, denoted by \texttt{\&}, is commutative, associative, and
idempotent.
CUE's design was motivated by Google's Borg Configuration
Language (GCL), which used graph unification.
CUE's lattice includes scalar types with a conflict
semantics where unifying incompatible scalars yields $\bot$; whether
scalars are necessary for a configuration language, or whether
tree structure alone suffices, is an interesting design question.
Jsonnet~\cite{jsonnet} provides object inheritance via the
\texttt{+} operator with mixin semantics: composition is
\emph{not} commutative since the right-hand side wins on scalar
conflicts, and deep merging requires explicit opt-in via the \texttt{+:}
syntax on a per-field basis, so deep merge is
available but not the default.
Dhall~\cite{dhall2017} takes a functional approach: it is a typed
$\lambda$-calculus with records, where composition is a record merge
operator, not an inheritance mechanism.
\fi

\paragraph{\bilingual{Module systems and deep merge}{模块系统与深合并}}
\ifInheritanceChinese
NixOS 模块系统~\cite{dolstra2010-nixos,nixosmodules} 通过\emph{递归合并嵌套属性集}来组合模块,而非通过线性化或方法层级的覆盖,这与先前的 mixin 和 trait 演算中区分树层级继承与方法层级复合的深合并语义相同。当两个模块定义了相同的嵌套路径时,其子树被合并而非覆盖。NixOS 模块系统集成了深合并、递归的 $\this$、模块的延迟求值,以及其上的丰富类型检查层。同一机制为 NixOS 系统配置、Home Manager~\cite{homemanager}、nix-darwin~\cite{nixdarwin}、flake-parts~\cite{flakeparts} 以及数十个其他生态系统提供支持,共同管理任意复杂度的配置。一个似乎尚未被完全解决的方面是 $\this$ 在多重继承下的语义。当两个各自继承自同一延迟模块的独立模块被组合时,每次导入都会引入对共享选项自己的声明;模块系统的 \texttt{merge\-Option\-Decls} 以静态错误拒绝重复声明(附录~\ref{app:scala-multi-target})。模块系统无法表达两条路径对同一选项的贡献是均等的,即递归合并自然产生的多目标情形,正是这一情形促成了第~\ref{sec:mixin-trees} 节的观测语义。
\else
The NixOS module system~\cite{dolstra2010-nixos,nixosmodules}
composes modules by \emph{recursively merging nested attribute sets},
not by linearization or method-level override, which is the same deep-merge
semantics that distinguishes tree-level inheritance from method-level
composition in prior mixin and trait calculi.
When two modules define the same nested path, their subtrees are
merged rather than overridden.
The NixOS module system incorporates deep merge,
recursive $\this$, deferred evaluation of modules, and a
rich type-checking layer on top.
The same mechanism powers
NixOS system configuration, Home
Manager~\cite{homemanager}, nix-darwin~\cite{nixdarwin},
flake-parts~\cite{flakeparts}, and dozens of other ecosystems,
collectively managing configurations of arbitrary complexity.
The one aspect that appears not to have been fully addressed is the
semantics of $\this$ under multiple inheritance.
When two independent modules that inherit from the same
deferred module are composed, each import introduces its own
declaration of the shared options; the module system's
\texttt{merge\-Option\-Decls} rejects duplicate declarations
with a static error (Appendix~\ref{app:scala-multi-target}).
The module system cannot express that both routes contribute
equally to the same option, that is, the multi-target situation that
recursive merging naturally gives rise to, which motivated the
observational semantics of Section~\ref{sec:mixin-trees}.
\fi

\subsection{\bilingual{Our Companion Work}{我们的配套工作}}

\paragraph{\bilingual{Modular software architectures}{模块化软件架构}}
\label{sec:practical-modular}
\ifInheritanceChinese
插件系统、组件框架与依赖注入框架通过声明式继承实现了软件的可扩展性。第~\ref{sec:mixin-trees} 节所确立的继承的对称性与可结合性解释了这类系统为何有效:组件可以以任意顺序继承而不改变行为。MIXINv2\anon[~(补充材料)]{~\cite{mixin2025}} 本身就是一个依赖注入框架:每个作用域将其依赖声明为具名槽,复合跨越作用域边界按名解析它们,从而使继承演算的模式显式化而非随意拼凑。
\else
Plugin systems, component frameworks, and dependency injection frameworks
enable software extensibility through declarative inheritance. The
symmetric, associative nature of inheritance established in Section~\ref{sec:mixin-trees} explains why such
systems work: components can be inherited in any order without changing
behavior.
MIXINv2\anon[~(supplementary material)]{~\cite{mixin2025}}
is itself a dependency injection framework:
each scope declares its dependencies as named slots, and
composition resolves them by name across scope boundaries,
making the inheritance-calculus patterns explicit rather than ad~hoc.
\fi

\paragraph{\bilingual{Union file systems}{联合文件系统}}
\label{sec:practical-union-fs}
\ifInheritanceChinese
联合文件系统,如 UnionFS、OverlayFS 与 AUFS,将多个目录层次叠加以呈现统一视图。其语义与继承演算相似,表现在后面的层覆盖前面的层且所有层的文件均可访问,但其设计在具体之处有所不同:继承是不对称的,标量文件的冲突解决是不可交换的,且层与层之间不存在晚绑定或动态分派。我们将继承演算直接实现为联合文件系统\anon[~(作为补充材料附上)]{~\cite{ratarmount2025}},仅凭三个原语就获得了文件查询的晚绑定语义与动态分派,无需任何额外机制。在软件包管理器、构建系统与操作系统发行版中,外部工具链将配置转换为文件树;而在这里,文件系统作为自身的配置,转换即是继承:这是一种\emph{文件系统即编译器},其中目录树既是源程序又是编译目标。
\else
Union file systems, such as UnionFS, OverlayFS, and AUFS, layer multiple directory
hierarchies to present a unified view. Their semantics resembles
inheritance-calculus in that later layers override earlier layers and files from all
layers remain accessible, but their designs differ in specific ways:
inheritance is asymmetric, conflict resolution for scalar files
is noncommutative, and there is no late-binding or dynamic
dispatch across layers.
We implemented inheritance-calculus directly as a union file
system\anon[~(included as supplementary material)]{~\cite{ratarmount2025}},
obtaining late-binding semantics and dynamic dispatch for file lookups
with no additional mechanism beyond the three primitives.
In package managers, build systems, and OS distributions,
external toolchains transform configuration into file trees;
here, the file system serves as configuration to itself,
and the transformation is inheritance: a
\emph{file-system-as-compiler} in which the directory tree is
both the source program and the compilation target.
\fi

\section{\bilingual{Future Work}{未来工作}}
\label{sec:future-work}

\ifInheritanceChinese
\anon[辅助实现]{MIXINv2~\cite{mixin2025}}{} 是继承演算的一个可执行实现,它已经超越了本文所呈现的无类型演算:它包含编译期检查,以验证所有引用都解析到 mixin 树中的有效路径,以及一个从宿主语言引入标量值的外部函数接口 (FFI)。以下所描述的未来工作,关注的是这一实现与一门完全实用的语言之间的差距,横跨理论基础与库层面的编码两个方面。
\else
\anon[The supplementary implementation]{MIXINv2~\cite{mixin2025}}{} is an executable implementation of inheritance-calculus
that already goes beyond the untyped calculus presented here:
it includes compile-time checking that all references resolve to
valid paths in the mixin tree, and a foreign-function interface (FFI)
that introduces scalar values from the host language.
The future work described below concerns the distance between this
implementation and a fully practical language, spanning both
theoretical foundations and library-level encodings.
\fi

\subsection{\bilingual{Type System}{类型系统}}
\ifInheritanceChinese
MIXINv2 现有的编译期检查验证每条引用路径都解析到继承结果中存在的属性。将这些检查的可靠性形式化,即证明类型良好的程序在运行时不产生悬空引用,是未来工作的一个方向。这样的形式化与 DOT 演算~\cite{amin2016-dependent-object-types} 的关系,类似于继承演算与 $\lambda$-演算的关系:有类型的继承演算可以看作不含 $\lambda$ 的 DOT,它保留了路径依赖类型,同时以 mixin 继承取代函数。正式的类型系统将适合在 Coq 和 Agda 等证明助手中机械化,以验证类型安全、进展和保持等性质。

在现有检查的可靠性之外,额外的类型系统特性将强化这门语言。一个关键的例子是\emph{完全性检查}:验证一个继承为所有必需槽位提供了实现,而不仅仅是引用能够解析。在无类型的继承演算中,写下 $\mathrm{magnitude} \mapsto \{\}$ 纯属说明性质:它定义了一个结构槽位,但不对填充它的内容施加任何约束。完全性检查将在静态上强制这类声明,拒绝使必需槽位保持未填充状态的继承。
\else
MIXINv2's existing compile-time checks verify that every
reference path resolves to a property that exists in the inherited
result. Formalizing the soundness of these checks, that is, proving that
well-typed programs do not produce dangling references at
runtime, is one direction of future work.
Such a formalization would relate to the DOT
calculus~\cite{amin2016-dependent-object-types} in a manner analogous to how
inheritance-calculus relates to the $\lambda$-calculus: a Typed
inheritance-calculus can be viewed as DOT without $\lambda$, retaining
path-dependent types while replacing functions with
mixin inheritance.
A formal type system would be amenable to mechanization in proof
assistants, such as Coq and Agda, for verifying type safety, progress, and
preservation properties.

Beyond soundness of the existing checks, additional type system
features would strengthen the language.
A key example is \emph{totality checking}: verifying that an
inheritance provides implementations for all required slots, not
merely that references resolve.
In untyped inheritance-calculus, writing $\mathrm{magnitude} \mapsto \{\}$
is purely documentary: it defines a structural slot but imposes no
constraint on what must fill it.
Totality checking would enforce such declarations statically,
rejecting inheritances that leave required slots unfilled.
\fi

\subsection{\bilingual{Standard Library}{标准库}}
\ifInheritanceChinese
其余几项涉及标准库:它们不需要对演算或类型系统进行扩展,但需要在现有框架内进行非平凡的编码。

继承演算\emph{默认是开放的}:继承可以自由合并任意两个 mixin,一个值可以同时具有多个数据构造子(例如,$\{\mathrm{Zero},\; \mathrm{Odd},\; \mathrm{half} \mapsto \{\mathrm{Zero}\}\}$)。这是该演算自然的 trie 语义,也是第~\ref{sec:expression-problem}~节求解表达式问题的基础。然而,许多操作(例如等值测试)假设值是\emph{线性的},即恰好具有一个数据构造子。第~\ref{sec:case-study}~节的观察者模式并不强制这一点:将其应用于一个具有多个数据构造子的值时,所有回调都会触发,其结果被继承在一起。

\emph{模式验证}可以在库层面使用现有原语来实现。观察者可以测试某个特定的数据构造子是否存在,返回一个布尔值;第二次布尔派发随后基于该结果进行分支。通过链接这样的测试,工厂可以验证一个值恰好匹配一个数据构造子并携带所需的属性。这并不限制继承本身(继承仍然不受约束),但提供了一种在不变式违反传播之前检测到它的方式。

同样的技术,即观察者返回布尔值后接布尔派发,可以产生\emph{封闭模式匹配}。第~\ref{sec:case-study}~节的观察者模式本质上是\emph{开放的}:新的数据构造子情况可以通过继承添加。封闭匹配(其中恰好执行一个分支)可以编码为一个\emph{访问者链}:一个由 if-then-else 节点组成的链表,每个节点测试一个数据构造子,在不匹配时向后传递。这类似于 GHC.Generics 的 $\mathrm{(:+:)}$ 和类型表示~\cite{magalhaes2010-generic-deriving},其中每个链节对应一个求和项。一个重要的推论是,这样的链\emph{不}可交换,因为它们具有固定的顺序,因此无法通过开放继承来扩展,而这正是封闭派发的正确语义。
\else
The remaining items concern the standard library: they require
no extensions to the calculus or type system, but involve nontrivial
encodings within the existing framework.

Inheritance-calculus is \emph{open by default}: inheritance
can freely merge any two mixins, and a value may simultaneously
inhabit multiple data constructors (e.g.,
$\{\mathrm{Zero},\; \mathrm{Odd},\; \mathrm{half} \mapsto \{\mathrm{Zero}\}\}$).
This is the natural trie semantics of the
calculus, and the basis for solving the Expression Problem
in Section~\ref{sec:expression-problem}.
However, many operations, such as equality testing, assume that values
are \emph{linear}, inhabiting exactly one data constructor.
The observer pattern of Section~\ref{sec:case-study} does not enforce
this: applied to a multi-data-constructor value, all callbacks fire and
their results are inherited.

\emph{Schema validation} can be implemented at the library level
using existing primitives.
An observer can test whether a particular data constructor is present,
returning a Boolean; a second Boolean dispatch then branches on the
result.
By chaining such tests, a factory can validate that a value matches
exactly one data constructor and carries the required properties.
This does not restrict inheritance itself, which remains
unconstrained, but provides a way to detect invariant violations
before they propagate.

The same technique, an observer returning Boolean followed by Boolean
dispatch, yields \emph{closed pattern matching}.
The observer pattern of Section~\ref{sec:case-study} is inherently
\emph{open}: new data-constructor cases can be added through inheritance.
Closed matching, where exactly one branch is taken, can be encoded as
a \emph{visitor chain}: a linked list of if-then-else nodes, each
testing one data constructor and falling through on mismatch.
This is analogous to GHC.Generics' $\mathrm{(:+:)}$ sum
representation~\cite{magalhaes2010-generic-deriving}, where each link
corresponds to one summand.
An important consequence is that such chains do \emph{not}
commute, as they have a fixed order, and therefore cannot be extended
through open inheritance, which is the correct semantics for closed
dispatch.
\fi

\section{\bilingual{Conclusion}{结论}}

\ifInheritanceChinese
对不可扩展性的免疫指导了继承演算的设计。第~\ref{sec:semantic-variants}~节排除了每一种替代的组合机制和引用解析,使得带有多目标晚绑定限定 $\this$ 的深合并成为唯一的幸存者。mixin 充当了统一模块、类、对象、方法、字段和局部绑定的单一抽象;mixin 树的深合并使继承具有可交换性、幂等性和可结合性,从而多重继承的线性化问题不再出现。第~\ref{sec:mixin-trees}~节的六个一阶方程,具有指称语义且可由 tabling~\cite{tamaki1986-tabled-resolution} 计算,足以定义语义;观察者通过查询一棵惰性构造的树中的路径来驱动计算。这些方程是对标准面向对象语义的最小偏离:唯一的改变是 $\overrides$ 方程~(\ref{eq:overrides}),它以子树的深合并取代了方法层面的遮蔽。面向对象编程传统上被归类为命令式的;去掉可变状态和函数体,便得到纯面向对象编程,%
\footnote{
  ``纯''具有双重含义:不纯粹意味着不纯洁,即值是不可变的且组合没有副作用;而不混合意味着组合完全依赖继承,没有函数式抽象,例如 $\lambda$ 或 $\beta$-归约。
}
这是声明式编程的一个最小计算模型。

由于 $\overrides$ 合并的是子树而非方法体,函数不再是原语。$\lambda$-演算在第~\ref{sec:translation}~节通过六条翻译规则宏表达到继承演算中。第~\ref{sec:bohm-tree}~节的 Böhm 树对应将这一一阶语义扩展到了 $\lambda$-演算:先前的高阶语义框架被证明忠实于继承演算一阶方程的一个子语言的首次迭代近似。

继承演算进一步获得了 $\lambda$-演算无法表达的三种现象;事实上,正如第~\ref{sec:asymmetry}~节所示,在 Felleisen 意义下~\cite{felleisen1991-expressive-power},继承演算严格地比惰性 $\lambda$-演算更具表达力。第~\ref{sec:expression-problem}~节在表达式问题~\cite{wadler1998-expression-problem}意义下建立了对不可扩展性的免疫。第~\ref{sec:case-study}~节的案例研究表明,普通算术产生了逻辑编程的关系语义:循环继承层次结构通过与递归 Datalog 所支配的相同的升链条件收敛,附录~\ref{app:datalog-encoding}给出了一个系统性的编码。第~\ref{sec:practical-function-color}~节展示了函数颜色盲性~\cite{nystrom2015-function-color}。

单个 mixin 同时是数据、程序、配置、函数、关系和对象。它宏表达了 $\lambda$-演算,并恢复了逻辑编程的关系语义,同时对两者均无法逃脱的不可扩展性保持免疫。

引言提出了配置本身是否能成为实用的通用编程的问题。答案是肯定的。我们提出继承演算作为实用通用编程的基础,我们相信它优于任何已知的计算模型。
\else
Immunity to nonextensibility guided the design of
inheritance-calculus.
Section~\ref{sec:semantic-variants} eliminated every
alternative composition mechanism and reference resolution,
leaving deep merge with multi-target late-binding qualified
$\this$ as the unique survivor.
The mixin served as the single abstraction unifying module,
class, object, method, field, and local binding;
deep merge of mixin trees made inheritance commutative, idempotent, and associative,
so the linearization problem of multiple inheritance did not arise.
The six first-order equations of Section~\ref{sec:mixin-trees},
denotational and computable by
tabling~\cite{tamaki1986-tabled-resolution}, suffice to
define the semantics; the observer drives computation by
querying paths in a lazily constructed tree.
These equations are a minimal departure from standard
object-oriented semantics: the only change is the $\overrides$ equation~(\ref{eq:overrides}),
which replaces method-level shadowing with deep merge of subtrees.
Object-oriented programming is traditionally classified as
imperative; removing mutable state and function bodies yields
pure object-oriented programming,%
\footnote{
  ``Pure'' carries a double meaning:
  not impure, since values are immutable and composition has no
  side effects; and not mixed, since composition relies entirely on
  inheritance, without functional abstractions such
  as~$\lambda$ or~$\beta$-reduction.
}
a minimal computational model for declarative programming.

Since $\overrides$ merges subtrees rather than method bodies,
functions are no longer primitive.
The $\lambda$-calculus macro-expressed into
inheritance-calculus in six translation rules in
Section~\ref{sec:translation}.
The B\"ohm tree correspondence of
Section~\ref{sec:bohm-tree} extended this first-order
semantics to the $\lambda$-calculus:
the prior higher-order semantic frameworks proved faithful
to the first-iteration approximation of a sublanguage of
inheritance-calculus's first-order equations.

Inheritance-calculus further gains three phenomena
that the $\lambda$-calculus cannot express;
indeed, inheritance-calculus is strictly more expressive
than the lazy $\lambda$-calculus in Felleisen's
sense~\cite{felleisen1991-expressive-power},
as shown in Section~\ref{sec:asymmetry}.
Section~\ref{sec:expression-problem} established immunity
to nonextensibility in the sense of the
Expression Problem~\cite{wadler1998-expression-problem}.
The case study of Section~\ref{sec:case-study} showed that
ordinary arithmetic yields the relational semantics of logic
programming: cyclic inheritance hierarchies converge by
the same ascending chain condition that governs recursive
Datalog, and Appendix~\ref{app:datalog-encoding} gave a
systematic encoding.
Section~\ref{sec:practical-function-color} exhibited
function color
blindness~\cite{nystrom2015-function-color}.

A single mixin is at once data, program, configuration,
function, relation, and object.
It macro-expresses the $\lambda$-calculus and recovers the
relational semantics of logic programming, while remaining
immune to the nonextensibility that neither escapes.

The introduction asked whether configuration alone can be
practical general-purpose programming.
It can.
We propose inheritance-calculus as a foundation for practical
general-purpose programming, one we believe better than any
known computational model.
\fi

\section{\bilingual{Data Availability Statement}{数据可用性声明}}

\ifInheritanceChinese
可执行实现 (MIXINv2) 以及本文中的所有示例以 MIT 许可证作为开源软件公开发布\anon[{} (匿名审查期间省略 URL)]{~\cite{mixin2025}}。每个以继承演算记法呈现的示例都有对应的可执行 MIXINv2 源码,且一套测试套件复现了论文中的每个示例。案例研究 (第~\ref{sec:case-study}~节) 的布尔逻辑和一元自然数算术、附录~\ref{app:datalog-encoding}~的 Datalog 编码 (包括双跳、常量连接和多变量连接关系程序以及重复访问者自连接回归示例),以及函数颜色盲性段落 (第~\ref{sec:practical-function-color}~节) 中讨论的同步和异步运行时实例,均包含在这些源码中。第~\ref{sec:practical-union-fs}~节所讨论的 ratarmount 补丁\anon[也已公开发布 (匿名审查期间省略 URL)]{作为一个开放的拉取请求~\cite{ratarmount2025}可以获取}。\anon[以上所有内容也作为辅助材料一并附上。]{}
\else
The executable implementation (MIXINv2) and all examples from
this paper are openly available as open-source software under
the MIT license\anon[{} (URL omitted for anonymous review)]{~\cite{mixin2025}}.
Every example rendered in inheritance-calculus notation has a
corresponding executable MIXINv2 source, and a test suite
reproduces every example in the paper.
The boolean logic and unary natural number arithmetic of the case
study (Section~\ref{sec:case-study}), the Datalog encodings of
Appendix~\ref{app:datalog-encoding} (including the two-hop,
constant-join, and multi-variable-join relational programs and the
repeated-visitor self-join regression example), and the synchronous
and asynchronous runtime instances discussed in the
function-color-blindness paragraph
(Section~\ref{sec:practical-function-color}) are all included among
these sources.
The ratarmount patch discussed in
Section~\ref{sec:practical-union-fs} is
\anon[also openly available (URL omitted for anonymous review)]{available as an open pull
request~\cite{ratarmount2025}}.
\anon[All of the above are also included as supplementary material.]{}
\fi

\begin{acks}
  Claude Code, GitHub Copilot, Gemini Code Assist, and Codex were utilized to generate sections of this Work,
  including text, tables, graphs, code, formulas, citations, etc.
\end{acks}
\fi % end \ifInheritanceBody (the main matter)

% References follow the main matter and precede the appendix. This copy fires
% whenever the main matter is present (preprint, submission); the appendix-only
% build emits its copy after the appendix instead.
\ifInheritanceBody
\bibliographystyle{ACM-Reference-Format}
\bibliography{references}
\fi

\ifInheritanceAppendix
\appendix

\section{\bilingual{Well-Definedness of the Semantic Functions}{语义函数的良定义性}}
\label{app:well-definedness}

\ifInheritanceChinese
六个语义方程~(\ref{eq:properties})--(\ref{eq:this}) 是一阶的:其输入与输出均为路径及路径的集合,不涉及函数空间、闭包或高阶构造。
本节通过证明这六个方程定义了一个单调算子来确立良定义性;该算子的最小不动点由 tabling 求值~\cite{tamaki1986-tabled-resolution} 可靠地计算。
\else
The six semantic
equations~(\ref{eq:properties})--(\ref{eq:this}) are
first-order: their inputs and outputs are paths and sets of
paths, with no function spaces, closures, or higher-order
constructs.
This section establishes well-definedness by showing that
the six equations define a monotone operator whose least
fixed point is computed soundly by tabled
evaluation~\cite{tamaki1986-tabled-resolution}.
\fi

\begin{definition}[\bilingual{Dependency graph}{依赖图}]\label{def:dependency-graph}
\ifInheritanceChinese
  固定一棵 AST 及其原始函数 $\defines$ 与 $\inherits$。
  \emph{依赖图}~$G$ 以由该 AST 产生的所有语义函数应用
  ($\properties(p)$、\allowbreak$\supers(p)$、\allowbreak$\overrides(p)$、\allowbreak$\bases(p)$、\allowbreak
  $\resolve(\ldots)$、\allowbreak$\this(\ldots)$) 为顶点。
  当通过方程~(\ref{eq:properties})--(\ref{eq:this}) 计算~$u$
  需要~$v$ 的值时,从顶点~$u$ 到顶点~$v$ 有一条有向边。
\else
  Fix an AST with its primitive functions $\defines$ and
  $\inherits$.
  The \emph{dependency graph}~$G$ has as vertices all
  semantic-function applications
  ($\properties(p)$, $\supers(p)$, $\overrides(p)$,
    $\bases(p)$, $\resolve(\ldots)$,
  $\this(\ldots)$) that arise from the AST\@.
  There is a directed edge from vertex~$u$ to vertex~$v$
  whenever computing~$u$ via
  equations~(\ref{eq:properties})--(\ref{eq:this})
  requires the value of~$v$.
\fi
\end{definition}

\begin{definition}[\bilingual{Immediate consequence operator $T_P$}{直接后果算子 $T_P$}]%
  \label{def:tp-operator}
\ifInheritanceChinese
  设 $\mathcal{Q}$ 为所有语义函数查询(即~$G$ 的顶点)的集合。
  对每个查询 $q \in \mathcal{Q}$,设 $\mathcal{V}_q$ 为形状与~$q$ 相匹配的答案值的集合
  (依定义~$q$ 的方程,为路径、路径的集合或路径对的集合)。
  \emph{查询--答案映射}是逐点乘积
  $\prod_{q \in \mathcal{Q}} \mathcal{P}(\mathcal{V}_q)$ 中的一个元素;
  等价地,它为每个查询~$q$ 指派一个候选答案集合 $M(q) \subseteq \mathcal{V}_q$。
  \emph{直接后果算子}~\cite{vanemden1976-predicate-logic-semantics}
  $T_P$ 将查询--答案映射~$M$ 映射为新的查询--答案映射~$T_P(M)$,方式是:对每个查询~$q$,
  求值定义~$q$ 的那条方程的右侧,并以当前近似~$M(q')$ 解释每个子查询~$q'$。
\else
  Let $\mathcal{Q}$ be the set of all
  semantic-function queries
  (vertices of~$G$).
  For each query $q \in \mathcal{Q}$, let
  $\mathcal{V}_q$ be the set of answer values of the
  appropriate shape for~$q$
  (paths, sets of paths, or sets of pairs of paths,
  depending on the equation defining~$q$).
  A \emph{query--answer map} is an element of the
  pointwise product
  $\prod_{q \in \mathcal{Q}} \mathcal{P}(\mathcal{V}_q)$;
  equivalently, it assigns to each query~$q$ a set
  $M(q) \subseteq \mathcal{V}_q$ of candidate answers.
  The \emph{immediate consequence
  operator}~\cite{vanemden1976-predicate-logic-semantics}
  $T_P$ maps a query--answer map~$M$ to a new
  query--answer map~$T_P(M)$ by evaluating, for each
  query~$q$, the right-hand side of the equation that
  defines~$q$, interpreting every sub-query~$q'$ by
  the current approximation~$M(q')$.
\fi
\end{definition}

\begin{lemma}[\bilingual{Monotonicity}{单调性}]\label{lem:monotonicity}
  \ifInheritanceChinese
  算子 $T_P$ 在完备格
  $\Bigl(\prod_{q \in \mathcal{Q}} \mathcal{P}(\mathcal{V}_q),\; \subseteq\Bigr)$
  上是单调的。
  \else
  The operator $T_P$ is monotone on the complete lattice
  $\Bigl(\prod_{q \in \mathcal{Q}} \mathcal{P}(\mathcal{V}_q),\; \subseteq\Bigr)$.
  \fi
\end{lemma}

\begin{proof}
\ifInheritanceChinese
  每条方程的右侧均由 $\cup$、$\in$、等式测试以及正集合
  概括 $\{\, x \mid x \in S, \ldots \,\}$ 构成。
  六条方程均不使用否定、集合差或补集。
  因此,若 $M \subseteq M'$(逐点),则 $T_P(M) \subseteq T_P(M')$。
\else
  Each equation's right-hand side is built from
  $\cup$, $\in$, equality tests, and positive set
  comprehension $\{\, x \mid x \in S, \ldots \,\}$.
  None of the six equations uses negation, set
  difference, or complement.
  Therefore, if $M \subseteq M'$ (pointwise), then
  $T_P(M) \subseteq T_P(M')$.
\fi
\end{proof}

\begin{theorem}[\bilingual{Existence and uniqueness of
  $\mathrm{lfp}(T_P)$}{$\mathrm{lfp}(T_P)$ 的存在性与唯一性}]\label{thm:lfp-exists}
  \ifInheritanceChinese
  算子 $T_P$ 有唯一的最小不动点 $\mathrm{lfp}(T_P)$。
  \else
  The operator $T_P$ has a unique least fixed point
  $\mathrm{lfp}(T_P)$.
  \fi
\end{theorem}

\begin{proof}
\ifInheritanceChinese
  每个幂集格
  $(\mathcal{P}(\mathcal{V}_q), \subseteq)$
  都是完备的,
  因此逐点乘积格
  $\Bigl(\prod_{q \in \mathcal{Q}} \mathcal{P}(\mathcal{V}_q), \subseteq\Bigr)$
  也是完备的。
  由 Knaster--Tarski 定理~\cite{tarski1955-lattice-fixpoint-theorem},
  完备格上的每个单调函数都有唯一的最小不动点。
\else
  Each powerset lattice
  $(\mathcal{P}(\mathcal{V}_q), \subseteq)$
  is complete,
  so the pointwise product lattice
  $\Bigl(\prod_{q \in \mathcal{Q}} \mathcal{P}(\mathcal{V}_q), \subseteq\Bigr)$
  is also complete.
  By the Knaster--Tarski
  theorem~\cite{tarski1955-lattice-fixpoint-theorem},
  every monotone function on a complete lattice has a
  unique least fixed point.
\fi
\end{proof}

\begin{theorem}[\bilingual{Well-definedness}{良定义性}]\label{thm:well-defined}
  \ifInheritanceChinese
  语义函数 $\properties$、\allowbreak$\supers$、\allowbreak$\overrides$、\allowbreak$\bases$、\allowbreak
  $\resolve$、\allowbreak$\this$ 是良定义的:
  其指称语义为 $\mathrm{lfp}(T_P)$,即直接后果算子的唯一最小不动点。
  \else
  The semantic functions
  $\properties$, $\supers$, $\overrides$, $\bases$,
  $\resolve$, $\this$ are well-defined:
  their denotational semantics is
  $\mathrm{lfp}(T_P)$, the unique least fixed point of
  the immediate consequence operator.
  \fi
\end{theorem}

\begin{theorem}[\bilingual{Tabling soundness}{Tabling 可靠性}]%
  \label{thm:tabling-soundness}
  \ifInheritanceChinese
  即使依赖图~$G$ 是无穷的,方程~(\ref{eq:properties})--(\ref{eq:this}) 的
  tabling 求值也绝不会返回错误答案。具体而言:
  \begin{enumerate}
    \item \textbf{可靠性(无假正例)。}
      tabling 返回的每个事实均由某条方程的右侧在其输入上求值产生。
      对求值轨迹进行归纳,每个返回的事实都属于 $\mathrm{lfp}(T_P)$。

    \item \textbf{单调积累。}
      tabling 累加器在各轮迭代之间只通过集合并增长,任何事实都不会被删除。

    \item \textbf{不动点检测的正确性。}
      若求值过程中没有子查询被重入,则所有子查询均在无近似的情况下被完整计算,结果是精确的。
      若发生了重入但累加器是稳定的(即没有新事实被添加),则累加器等于
      $\mathrm{lfp}(T_P)$ 在可达查询上的限制,因为它是一个从~$\bot$ 上升到达的前不动点。

    \item \textbf{可靠的失败。}
      当 tabling 无法确定答案时(无论是因为迭代预算耗尽还是因为存在无穷非循环链),
      它会抛出错误而非返回错误答案。
  \end{enumerate}
  \else
  Tabled evaluation of
  equations~(\ref{eq:properties})--(\ref{eq:this})
  never returns a wrong answer, even when the
  dependency graph~$G$ is infinite.
  Specifically:
  \begin{enumerate}
    \item \textbf{Soundness (no false positives).}
      Every fact returned by tabling was produced by
      evaluating some equation's right-hand side on its
      inputs.
      By induction on the evaluation trace, every
      returned fact belongs to $\mathrm{lfp}(T_P)$.

    \item \textbf{Monotone accumulation.}
      The tabling accumulator only grows between
      iterations via set union; no fact is ever
      removed.

    \item \textbf{Correct fixpoint detection.}
      If no sub-query was re-entered during evaluation,
      all sub-queries were fully computed with no
      approximation, and the result is exact.
      If re-entry occurred but the accumulator is
      stable, meaning no new facts were added, the
      accumulator equals
      $\mathrm{lfp}(T_P)$ restricted to the reachable
      queries, because it is a pre-fixed point reached
      by ascending from~$\bot$.

    \item \textbf{Sound failure.}
      When tabling cannot determine the answer,
      whether because the iteration budget is exhausted
      or because of an infinite non-cyclic chain,
      it raises an error rather than
      returning a wrong answer.
  \end{enumerate}
  \fi
\end{theorem}

\begin{proof}
\ifInheritanceChinese
  对于~(1),每个返回的事实都是某条方程的右侧作用于先前返回的事实上的输出;
  归纳地,所有输入均属于 $\mathrm{lfp}(T_P)$,因此输出也是(由单调性和不动点性质)。

  对于~(2),实现通过 $A \leftarrow A \cup T_P(A)$ 进行积累,这对~$A$ 是单调的。

  对于~(3),当没有重入时,求值是有限无环可达子图上的普通记忆化递归,
  因此每个子查询都被精确计算,不涉及任何近似。
  当发生重入时,设 $R \subseteq \mathcal{Q}$ 为从初始查询可达的查询集合,
  设 $T_P^R$ 表示 $T_P$ 到~$R$ 上映射的限制。
  若累加器稳定,即 $A = A \cup T_P^R(A)$,则
  $T_P^R(A) \subseteq A$,从而~$A$ 是~$T_P^R$ 的一个前不动点。

  由~(1),曾被添加到累加器中的每个事实均由某次右侧求值产生,
  因此属于 $\mathrm{lfp}(T_P^R)$。
  故 $A \subseteq \mathrm{lfp}(T_P^R)$。

  反之,由 Knaster--Tarski 将最小不动点刻画为最小前不动点,
  $\mathrm{lfp}(T_P^R)$ 包含于~$T_P^R$ 的每个前不动点中,特别地包含于~$A$ 中。
  从而 $\mathrm{lfp}(T_P^R) \subseteq A$。

  因此 $A = \mathrm{lfp}(T_P^R)$,即稳定的累加器等于可达查询上的精确指称。

  对于~(4),预算耗尽与栈溢出均产生异常,而非静默的错误答案。
\else
  For~(1), each returned fact is the output of some
  equation's right-hand side applied to previously
  returned facts; by induction, all inputs are in
  $\mathrm{lfp}(T_P)$, so the output is too (by
  monotonicity and the fixed-point property).

  For~(2), the implementation accumulates via
  $A \leftarrow A \cup T_P(A)$, which is monotone in~$A$.

  For~(3), when no reentry occurs, the evaluation is
  ordinary memoised recursion on a finite acyclic
  reachable subgraph, so every sub-query is computed
  exactly and no approximation is involved.
  When reentry occurs, let $R \subseteq \mathcal{Q}$ be
  the set of queries reachable from the initial query,
  and let $T_P^R$ denote the restriction of $T_P$ to
  maps on~$R$.
  If the accumulator is stable, that is,
  $A = A \cup T_P^R(A)$, then
  $T_P^R(A) \subseteq A$, so $A$ is a pre-fixed point
  of~$T_P^R$.

  By~(1), every fact ever added to the accumulator was
  produced by some right-hand side evaluation and
  therefore belongs to $\mathrm{lfp}(T_P^R)$.
  Hence $A \subseteq \mathrm{lfp}(T_P^R)$.

  Conversely, by the Knaster--Tarski characterization
  of the least fixed point as the least pre-fixed
  point, $\mathrm{lfp}(T_P^R)$ is contained in every
  pre-fixed point of~$T_P^R$, in particular in~$A$.
  Thus $\mathrm{lfp}(T_P^R) \subseteq A$.

  Therefore $A = \mathrm{lfp}(T_P^R)$, that is, the stable
  accumulator equals the exact denotation on the
  reachable queries.

  For~(4), budget exhaustion and stack overflow both
  produce exceptions, never silent wrong answers.
\fi
\end{proof}

% 终止性依赖程序结构:无环有限子图→一趟即止;有环有限逐查询值域→有限格上升链条件收敛;
% 无穷可达子图→固有发散,属图灵完备性,非求值策略的缺陷。
\paragraph{\bilingual{Termination}{终止性}}
\ifInheritanceChinese
求值是否终止取决于程序:
\begin{itemize}
  \item \textbf{无环、有限可达子图。}
    一趟即足;求值终止。
  \item \textbf{有环、逐查询值域有限。}
    多趟迭代在有限格的上升链条件下收敛~\cite{bancilhon1986-recursive-query-strategies}。
    这是 Datalog 的情形(第~\ref{sec:recursive-rules} 节)。
  \item \textbf{无穷可达子图。}
    求值可能发散;这是图灵完备性的固有属性,而非求值策略的缺陷。
\end{itemize}
\else
Whether evaluation terminates depends on the
program:
\begin{itemize}
  \item \textbf{Acyclic, finite reachable subgraph.}
    One pass suffices; evaluation terminates.
  \item \textbf{Cyclic, finite per-query domain.}
    Multiple passes converge by the ascending chain
    condition on a finite
    lattice~\cite{bancilhon1986-recursive-query-strategies}.
    This is the Datalog case
    (Section~\ref{sec:recursive-rules}).
  \item \textbf{Infinite reachable subgraph.}
    Evaluation may diverge; this is inherent to
    Turing completeness and not a defect of the
    evaluation strategy.
\end{itemize}
\fi

\section{B\"ohm Tree Correspondence: Proofs}
\label{app:bohm-tree-proofs}

This appendix proves two results about the sublanguage of
inheritance-calculus corresponding to the $\lambda$-calculus
(Section~\ref{sec:bohm-tree}).
First, the first-iteration approximation
$T_P\!\uparrow\!1$ is fully abstract with respect to
B\"ohm tree equivalence
(Theorems~\ref{thm:adequacy} and~\ref{thm:first-order-semantics}),
establishing that the first iteration coincides with
the lazy $\lambda$-calculus.
Second, the full $\mathrm{lfp}(T_P)$ semantics is strictly
more defined in the powerset sense of
Figure~\ref{fig:calculi-matrix}: translated $\lambda$-terms
that diverge under $\beta$-reduction may converge to definite
sets, as shown in Section~\ref{sec:bohm-tree}.

\subsection{B\"ohm Trees}

We briefly recall the definition of B\"ohm
trees~\cite{barendregt1984-lambda-calculus}.
A \emph{head reduction} $M \to_h M'$ contracts
the outermost $\beta$-redex only: if $M$ has the form
$(\lambda x.\, e)\; v\; M_1 \cdots M_k$, then
$M \to_h e[v/x]\; M_1 \cdots M_k$.
A term $M$ is in \emph{head normal form} (HNF) if it has no
head redex, that is, $M = \lambda x_1 \ldots x_m.\;
y\; M_1 \cdots M_k$ where $y$ is a variable.

The \emph{B\"ohm tree} $\mathrm{BT}(M)$ of a $\lambda$-term $M$
is an infinite labeled tree defined by:
\[
  \mathrm{BT}(M) =
  \begin{cases}
    \bot & \text{if } M \text{ has no HNF} \\[6pt]
    \lambda x_1 \ldots x_m.\;
    y\bigl(\mathrm{BT}(M_1),\; \ldots,\; \mathrm{BT}(M_k)\bigr)
    &
    \begin{aligned}[t]
      &\text{if } M \to^*_h \\[-2pt]
      &\lambda x_1 \ldots x_m.\; y\; M_1 \cdots M_k
    \end{aligned}
  \end{cases}
\]
B\"ohm tree equivalence $\mathrm{BT}(M) = \mathrm{BT}(N)$ is the
standard observational equivalence of the lazy
$\lambda$-calculus~\cite{abramsky1993-full-abstraction-lazy-lambda,abramsky1995-games-full-abstraction-lazy-lambda}.

\subsection{Path Encoding}

We define a correspondence between positions in the B\"ohm tree
and paths in the mixin tree.
When the translation $\mathcal{T}$ is applied to an ANF
$\lambda$-term, the resulting mixin tree has a specific shape:
\begin{itemize}
  \item An abstraction $\lambda x.\, M$ translates to a record
    with own properties $\{\mathrm{argument}, \mathrm{result}\}$,
    which is referred to as the \emph{abstraction shape}.
  \item A $\mathbf{let}$-binding
    $\mathbf{let}\; x = V_1\; V_2 \;\mathbf{in}\; M$ translates
    to a record with own properties $\{x, \mathrm{result}\}$,
    where $x$ holds the encapsulated application
    $\{\mathcal{T}(V_1),\; \mathrm{argument} \mapsto
    \mathcal{T}(V_2)\}$.
  \item A tail call $V_1\; V_2$ translates to a record with own
    properties $\{\mathrm{tailCall}, \mathrm{result}\}$, where
    $\mathrm{tailCall}$ holds the encapsulated application and
    $\mathrm{result}$ projects
    $\mathrm{tailCall}.\mathrm{result}$.
  \item A variable $x$ with de~Bruijn index $n$ translates to
    an indexed reference $\dbi{n}.\mathrm{argument}$.
\end{itemize}

\noindent
Each B\"ohm tree position is a sequence of navigation steps from
the root.
In the B\"ohm tree, the possible steps at a node
$\lambda x_1 \ldots x_m.\; y\; M_1 \cdots M_k$ are:
entering the body by peeling off one abstraction layer, and
entering the $i$-th argument $M_i$ of the head variable.
In the mixin tree, these correspond to following the
$\mathrm{result}$ label (for entering the body) and the
$\mathrm{argument}$ label after $\this$ resolution (for
entering an argument position).

Rather than comparing absolute paths across different mixin
trees, which would be sensitive to internal structure, we
observe only the \emph{convergence} behavior accessible via
the $\mathrm{result}$ projection
(Definition~\ref{def:convergence}).
Each $\mathrm{result}$-following step corresponds to one head
reduction step in the $\lambda$-calculus: a tail call's
$\mathrm{result}$ projects through the encapsulated
application, and a $\mathbf{let}$-binding's $\mathrm{result}$
enters the continuation.
The abstraction shape at the end signals that a head
normal form has been reached.

\subsection{Single-Path Lemma}

The $\this$ function (equation~\ref{eq:this}) is designed for the
general case of multi-target $\this$ resolution.
For translated $\lambda$-terms, we show it degenerates to
single-target resolution.

\begin{lemma}[Single-Path]\label{lem:single-path}
  For any closed ANF $\lambda$-term $M$, at every invocation of
  $\this(S,\; p_{\mathrm{def}},\; n)$ during evaluation of
  $\mathcal{T}(M)$, the frontier set $S$ contains exactly one
  path.
\end{lemma}

\begin{proof}
  By case analysis on the translation rules of~$\mathcal{T}$.
  Since $\mathcal{T}$ is compositional, the structural invariant
  established for each rule propagates to all inherited subtrees:
  every subtree appearing in the mixin tree is itself the image of
  some sub-term under~$\mathcal{T}$, so the invariant holds
  throughout the runtime composition.

  \paragraph{Case $\mathcal{T}(\lambda x.\, M)
    = \{\mathrm{argument} \mapsto \{\},\;
  \mathrm{result} \mapsto \mathcal{T}(M)\}$}
  This is a record literal with
  $\defines = \{\mathrm{argument}, \mathrm{result}\}$ and
  no inheritance sources, so $\inherits = \varnothing$.
  At the root path $p$, we have $\overrides(p) = \{p\}$ (a single
  path) since there are no inheritance sources to introduce
  additional branches.
  Inside $\mathcal{T}(M)$, any reference to $x$ is
  $\dbi{n}.\mathrm{argument}$ where $n$ is the de~Bruijn index
  pointing to this scope level.
  Since $\mathcal{T}(\lambda x.\, M)$ introduces exactly one
  scope level, there is exactly one
  mixin at that level, hence $\this$ finds exactly one matching
  pair $(p_{\mathrm{site}},\; p_{\mathrm{override}})$ in $\supers$.

  \paragraph{Case
    $\mathcal{T}(\mathbf{let}\; x = V_1\; V_2
    \;\mathbf{in}\; M)
    = \{x \mapsto \{\mathcal{T}(V_1),\;
      \mathrm{argument} \mapsto \mathcal{T}(V_2)\},\;
  \mathrm{result} \mapsto \mathcal{T}(M)\}$}
  The outer record has own properties $\{x, \mathrm{result}\}$
  and no inheritance sources at the outer level
  ($\inherits = \varnothing$), so $\overrides(\text{root})
  = \{\text{root}\}$.
  Inside the $x$ subtree, the inheritance
  $\{\mathcal{T}(V_1),\; \mathrm{argument} \mapsto
  \mathcal{T}(V_2)\}$ has exactly one inheritance source
  ($\mathcal{T}(V_1)$), creating exactly one $\bases$ entry.
  By induction on $V_1$, the single-path property holds inside
  $\mathcal{T}(V_1)$.
  Inside $\mathcal{T}(M)$, the induction hypothesis applies
  directly.

  \paragraph{Case
    $\mathcal{T}(\mathbf{let}\; x = V
    \;\mathbf{in}\; M)
    = \{x \mapsto \{\mathrm{result} \mapsto \mathcal{T}(V)\},\;
  \mathrm{result} \mapsto \mathcal{T}(M)\}$}
  The outer record has own properties $\{x, \mathrm{result}\}$
  and no inheritance sources at the outer level, so
  $\overrides(\text{root}) = \{\text{root}\}$.
  The $x$ subtree is itself a record literal with own
  property $\{\mathrm{result}\}$ and no inheritance
  sources.
  Thus neither entering the bound value nor entering the
  body introduces branching in $\supers$.
  By induction, the single-path property holds inside both
  $\mathcal{T}(V)$ and $\mathcal{T}(M)$.

  \paragraph{Case $\mathcal{T}(V_1\; V_2) =
    \{\mathrm{tailCall} \mapsto
      \{\mathcal{T}(V_1),\; \mathrm{argument} \mapsto
      \mathcal{T}(V_2)\},\;
      \mathrm{result} \mapsto
  \{\mathrm{tailCall}.\mathrm{result}\}\}$}
  Identical to the $\mathbf{let}$-binding case: the outer record
  has own properties $\{\mathrm{tailCall}, \mathrm{result}\}$
  with no inheritance sources at the outer level, and the
  $\mathrm{tailCall}$ subtree contains exactly one inheritance
  source.

  \paragraph{Case $\mathcal{T}(x)
  = \dbi{n}.\mathrm{argument}$}
  This is a reference, not a record.
  The $\this$ function walks up $n$ steps from the enclosing scope.
  By the compositionality argument above, every scope level
  encountered during the walk, including inherited subtrees
  reached through reference resolution, is the image of some
  sub-term under~$\mathcal{T}$, so each step of $\this$
  encounters exactly one matching pair in $\supers$.

  \medskip\noindent
  The key structural invariant is the following.
  For every subtree generated by~$\mathcal{T}$,
  its root has either no inheritance source or exactly one
  inheritance source.
  Moreover, the only roots with an inheritance source are the
  application wrappers introduced by the application
  $\mathbf{let}$-binding and tail-call rules, and in both cases
  that source is unique and encapsulated under a distinguished
  child ($x$ or $\mathrm{tailCall}$), never at the enclosing
  scope level.
  Therefore every scope level traversed by $\this$ contributes
  at most one matching pair in $\supers$, so the frontier set
  remains a singleton throughout evaluation.
\end{proof}

\subsection{Substitution Lemma}

The core mechanism of the translation is that inheritance with
$\mathrm{argument} \mapsto \mathcal{T}(V)$ plays the role of
substitution.  We make this precise.

\begin{lemma}[Substitution]\label{lem:substitution}
  Let $M$ be an ANF term in which $x$ may occur free,
  and let $V$ be a closed value.%
  \footnote{
    The proof uses the Single-Path Lemma
    (Lemma~\ref{lem:single-path}), which requires the inherited tree
    to occur within a closed term. This holds whenever
    $\lambda x.\, M$ applied to $V$ is a subterm of a closed ANF
    program.
  }
  Define the \emph{inherited tree}
  $C = \{\mathcal{T}(\lambda x.\, M),\;
  \mathrm{argument} \mapsto \mathcal{T}(V)\}$.
  Then for every path $p$ under $C \snoc \mathrm{result}$ and
  every label $\ell$:
  \[
    \ell \in \properties(p)
    \text{ in } C \snoc \mathrm{result}
    \quad\Longleftrightarrow\quad
    \ell \in \properties(p)
    \text{ in } \mathcal{T}(M[V/x])
  \]
  where $M[V/x]$ is the usual capture-avoiding substitution.
\end{lemma}

\begin{proof}
  By structural induction on $M$.

  \paragraph{Base case: $M = x$}
  Then $\mathcal{T}(x) =
  \dbi{n}.\mathrm{argument}$ (where $n$ is the de~Bruijn index
  of $x$), and
  $C \snoc \mathrm{result}$ contains this reference.
  In the inherited tree $C$, $\resolve$ and $\this$ resolve the reference by
  navigating from the reference's enclosing scope up $n$ steps
  to the binding $\lambda$, where the inheritance
  $\{\mathcal{T}(\lambda x.\, x),\;
  \mathrm{argument} \mapsto \mathcal{T}(V)\}$ inherits into the
  $\mathrm{argument}$ slot, providing $\mathcal{T}(V)$.
  Since $C$ is the $\mathrm{tailCall}$ subtree of
  $\mathcal{T}\bigl((\lambda x.\, M)\; V\bigr)$, which is a
  closed ANF term, Lemma~\ref{lem:single-path} applies and
  $\this$ finds exactly one path,
  so $\resolve$ returns $\mathcal{T}(V)$'s subtree.
  Since $M[V/x] = V$, we have
  $\mathcal{T}(M[V/x]) = \mathcal{T}(V)$, and the properties
  coincide.

  \paragraph{Base case: $M = y$ where $y \neq x$ and $y$ is
  $\lambda$-bound}
  Then $\mathcal{T}(y) = \dbi{m}.\mathrm{argument}$ where $m$
  is the de~Bruijn index of~$y$.
  This walks up to the scope that binds~$y$, which is different
  from the scope that binds~$x$, so the inheritance with
  $\mathrm{argument} \mapsto \mathcal{T}(V)$ at $x$'s binding
  scope has no effect.
  Since $M[V/x] = y$, the properties are identical.

  \paragraph{Base case: $M = y$ where $y \neq x$ and $y$ is
  let-bound}
  Then $\mathcal{T}(y) = y.\mathrm{result}$.
  This is a lexical reference to the sibling property~$y$ in the
  enclosing record, so the inherited tree~$C$ does not alter its
  resolution.
  Since $M[V/x] = y$, the properties are identical.

  \paragraph{Inductive case: $M = \lambda y.\, M'$}
  Then $\mathcal{T}(M) =
  \{\mathrm{argument} \mapsto \{\},\;
  \mathrm{result} \mapsto \mathcal{T}(M')\}$.
  The subtree at $C \snoc \mathrm{result}$ is another
  abstraction-shaped record.
  References to $x$ inside $\mathcal{T}(M')$ resolve through
  $\this$ in exactly the same way (with one additional scope
  level from the inner $\lambda$), and by induction on $M'$, the
  properties under
  $C \snoc \mathrm{result} \snoc \mathrm{result}$
  coincide with those of
  $\mathcal{T}(M'[V/x])$.
  Since
  $(\lambda y.\, M')[V/x] = \lambda y.\, (M'[V/x])$
  (assuming $y$ is fresh), the result follows.

  \paragraph{Inductive case:
    $M = \mathbf{let}\; z = V_1\; V_2$
  $\mathbf{in}\; M'$}
  Then
  \[
    \mathcal{T}(M) =
    \{z \mapsto \{\mathcal{T}(V_1),\;
      \mathrm{argument} \mapsto \mathcal{T}(V_2)\},\;
    \mathrm{result} \mapsto \mathcal{T}(M')\}.
  \]
  The substitution distributes:
  $M[V/x] =
  \mathbf{let}\; z = V_1[V/x]\; V_2[V/x]
  \;\mathbf{in}\; M'[V/x]$.
  By induction on $V_1$, $V_2$, and $M'$, the properties under
  each subtree ($z$ and $\mathrm{result}$) coincide,
  since references to $x$ inside each subtree resolve to
  $\mathcal{T}(V)$ via the same $\this$ mechanism.

  \paragraph{Inductive case:
    $M = \mathbf{let}\; z = V'
  \;\mathbf{in}\; M'$}
  Then
  \[
    \mathcal{T}(M) =
    \{z \mapsto \{\mathrm{result} \mapsto \mathcal{T}(V')\},\;
    \mathrm{result} \mapsto \mathcal{T}(M')\}.
  \]
  The substitution distributes:
  $M[V/x] =
  \mathbf{let}\; z = V'[V/x]
  \;\mathbf{in}\; M'[V/x]$.
  By induction on $V'$ and $M'$, the properties under both
  subtrees, specifically $z$ and $\mathrm{result}$, coincide.
  In particular, occurrences of~$z$ in $M'$ are translated as the
  lexical reference $z.\mathrm{result}$ on both sides, so the
  wrapper introduced by the value $\mathbf{let}$-binding is
  preserved by the substitution.

  \paragraph{Inductive case: $M = V_1\; V_2$ (tail call)}
  Analogous to the $\mathbf{let}$-binding case with
  $\mathrm{tailCall}$ in place of $z$.
\end{proof}

\subsection{Convergence Preservation}

The convergence criterion (Definition~\ref{def:convergence})
follows the $\mathrm{result}$ chain from the root.
Each step in this chain corresponds to one head reduction step
in the $\lambda$-calculus.
We make this precise.

\begin{lemma}[Result-step]\label{lem:result-step}
  Let $(\lambda x.\, M)\; V$ be a tail call, which is a $\beta$-redex.
  Then for every path $p$ and label $\ell$:
  \[
    \ell \in \properties\!\bigl(\,
    \text{root} \snoc \mathrm{result} \snoc p\,\bigr)
    \text{ in } \mathcal{T}\bigl((\lambda x.\, M)\; V\bigr)
    \;\Longleftrightarrow\;
    \ell \in \properties\!\bigl(\,
    \text{root} \snoc p\,\bigr)
    \text{ in } \mathcal{T}(M[V/x])
  \]
  That is, following one $\mathrm{result}$ projection in the
  tail call's mixin tree yields the same properties as the
  root of the reduct's mixin tree.
\end{lemma}

\begin{proof}
  The translation gives:
  \[
    \mathcal{T}\bigl((\lambda x.\, M)\; V\bigr) =
    \{\mathrm{tailCall} \mapsto
      \{\mathcal{T}(\lambda x.\, M),\;
      \mathrm{argument} \mapsto \mathcal{T}(V)\},\;
      \mathrm{result} \mapsto
    \{\mathrm{tailCall}.\mathrm{result}\}\}
  \]
  The $\mathrm{result}$ label at the root is defined as
  $\mathrm{tailCall}.\mathrm{result}$, so
  $\text{root} \snoc \mathrm{result}$ resolves to the
  $\mathrm{result}$ path inside the inheritance
  $\mathrm{tailCall} = \{\mathcal{T}(\lambda x.\, M),\;
  \mathrm{argument} \mapsto \mathcal{T}(V)\}$.
  This inheritance is exactly the inherited tree $C$ of
  Lemma~\ref{lem:substitution}, and
  $\mathrm{tailCall} \snoc \mathrm{result}$ is
  $C \snoc \mathrm{result}$.
  By Lemma~\ref{lem:substitution}, the properties under
  $C \snoc \mathrm{result}$ coincide with those of
  $\mathcal{T}(M[V/x])$.

  \medskip\noindent
  \emph{Let-binding variant.}
  An analogous result holds for
  $\mathbf{let}\; x = V_1\; V_2 \;\mathbf{in}\; M'$
  where $V_1 = \lambda y.\, M''$.
  The translation gives:
  \[
    \mathcal{T}\bigl(\mathbf{let}\; x = V_1\; V_2
    \;\mathbf{in}\; M'\bigr) =
    \{x \mapsto \{\mathcal{T}(V_1),\;
      \mathrm{argument} \mapsto \mathcal{T}(V_2)\},\;
    \mathrm{result} \mapsto \mathcal{T}(M')\}
  \]
  Since $\mathrm{result} \mapsto \mathcal{T}(M')$ is a property
  definition, the subtree at
  $\text{root} \snoc \mathrm{result}$ is $\mathcal{T}(M')$.
  Inside $\mathcal{T}(M')$, each reference to $x$ is
  $x.\mathrm{result}$, which resolves to the $\mathrm{result}$
  path inside the $x$ subtree
  $\{\mathcal{T}(\lambda y.\, M''),\;
  \mathrm{argument} \mapsto \mathcal{T}(V_2)\}$.
  This is the inherited tree $C$ of
  Lemma~\ref{lem:substitution}, so
  $x.\mathrm{result}$ has the same properties as
  $\mathcal{T}(M''[V_2/y])$.
  In the $\lambda$-calculus,
  $M = (\lambda x.\, M')(V_1\; V_2)$, and the head reduct is
  $N = M'[(V_1\; V_2)/x]$.
  In $\mathcal{T}(N)$, each occurrence of $x$ is replaced by
  $(V_1\; V_2)$, whose result, by the tail-call variant above,
  also yields $\mathcal{T}(M''[V_2/y])$.
  Therefore the properties at
  $\text{root} \snoc \mathrm{result} \snoc p$ in
  $\mathcal{T}(M)$ coincide with those at
  $\text{root} \snoc p$ in $\mathcal{T}(N)$ for all $p$.

  \medskip\noindent
  \emph{Value let-binding variant.}
  An analogous result holds for
  $\mathbf{let}\; x = V \;\mathbf{in}\; M'$.
  The translation gives:
  \[
    \mathcal{T}\bigl(\mathbf{let}\; x = V \;\mathbf{in}\; M'\bigr) =
    \{x \mapsto \{\mathrm{result} \mapsto \mathcal{T}(V)\},\;
    \mathrm{result} \mapsto \mathcal{T}(M')\}.
  \]
  Since $\mathrm{result} \mapsto \mathcal{T}(M')$ is a property
  definition, the subtree at
  $\text{root} \snoc \mathrm{result}$ is $\mathcal{T}(M')$.
  Inside $\mathcal{T}(M')$, each occurrence of $x$ is translated
  as the lexical reference $x.\mathrm{result}$, which resolves to
  the subtree $\mathcal{T}(V)$.
  This is exactly the translation of the reduct
  $M'[V/x]$.
  Therefore the properties at
  $\text{root} \snoc \mathrm{result} \snoc p$ in
  $\mathcal{T}(\mathbf{let}\; x = V \;\mathbf{in}\; M')$
  coincide with those at
  $\text{root} \snoc p$ in $\mathcal{T}(M'[V/x])$ for all $p$.
\end{proof}

\begin{theorem}[Convergence preservation]%
  \label{thm:convergence-preservation}
  If $M \to_h N$ (head reduction), then
  $\mathcal{T}(M){\Downarrow}$ if and only if
  $\mathcal{T}(N){\Downarrow}$.
\end{theorem}

\begin{proof}
  A head reduction step in ANF takes one of three forms:
  a tail call $(\lambda x.\, M')\; V \to_h M'[V/x]$,
  an application $\mathbf{let}$-binding
  $\mathbf{let}\; x = V_1\; V_2 \;\mathbf{in}\; M'
  \to_h M'[(V_1\; V_2)/x]$,
  or a value $\mathbf{let}$-binding
  $\mathbf{let}\; x = V \;\mathbf{in}\; M'
  \to_h M'[V/x]$.
  The second and third forms are simply the ANF
  presentation of one $\beta$-step in the underlying
  $\lambda$-term.
  In all three cases, Lemma~\ref{lem:result-step}
  (tail-call, application let-binding, and value
  let-binding variants) shows that
  the properties at $\text{root} \snoc \mathrm{result}^n$ in
  $\mathcal{T}(M)$ coincide with the properties at
  $\text{root} \snoc \mathrm{result}^{n-1}$ in
  $\mathcal{T}(N)$ for all $n \ge 1$, where $N$ is the
  head reduct.
  Hence the abstraction shape appears at depth $n$ in the
  pre-reduct if and only if it appears at depth $n - 1$ in
  the post-reduct.
\end{proof}

\begin{remark}[Iteration levels]
  \label{rem:convergence-lfp}
  Section~\ref{sec:bohm-tree} showed that the B\"ohm tree
  correspondence holds at the first iteration
  $T_P\!\uparrow\!1$ and that $\mathrm{lfp}(T_P)$ is
  strictly more defined in the powerset sense of
  Figure~\ref{fig:calculi-matrix}.
  Here we make the connection precise.
  Tabled evaluation~\cite{tamaki1986-tabled-resolution}
  computes $\mathrm{lfp}(T_P)$ by iterating:
  $T_P\!\uparrow\!0 = \bot$,
  $T_P\!\uparrow\!1 = T_P(\bot)$,
  $T_P\!\uparrow\!2 = T_P(T_P(\bot))$,
  \ldots\ until stabilisation
  (Theorem~\ref{thm:tabling-soundness}).

  \begin{description}
    \item[$T_P\!\uparrow\!1$ corresponds to
      $\beta$-reduction.]
      $T_P\!\uparrow\!1$ evaluates each equation once,
      with re-entrant calls returning~$\bot$
      (the value from $T_P\!\uparrow\!0$).
      The Result-Step Lemma
      (Lemma~\ref{lem:result-step}) establishes a
      structural isomorphism between the equation
      systems of $\mathcal{T}(M)$ and
      $\mathcal{T}(N)$, shifted by one
      $\mathrm{result}$ step.
      This isomorphism is purely structural and holds
      for any number of iterations.
      Theorems~\ref{thm:adequacy}
      and~\ref{thm:first-order-semantics}
      are proved at this level.
  \end{description}

  \paragraph{Observation granularity}%
  \label{rem:set-vs-domain}
  The set-valued semantics at $T_P\!\uparrow\!1$ is
  strictly finer than the domain-theoretic semantics of
  the lazy $\lambda$-calculus.
  In the latter, a non-terminating computation yields a
  single undifferentiated~$\bot$;
  in the former, one can in principle distinguish
  $\varnothing$ (a definite empty set, arrived at in
  finite time) from non-stabilisation of the Kleene
  iteration (no finite approximation suffices).
  The convergence criterion
  (Definition~\ref{def:convergence}) deliberately
  projects away this distinction: it observes only
  whether the $\mathrm{result}$ chain reaches an abstraction
  shape, reducing the observation to a single bit.
  Theorems~\ref{thm:adequacy}
  and~\ref{thm:first-order-semantics}
  are stated at this coarsened observation level.

  \begin{description}
    \item[$\mathrm{lfp}(T_P)$ restores full precision.]
      At $\mathrm{lfp}(T_P)$, the full precision of the
      set-valued semantics is restored: re-entrant calls
      that returned~$\varnothing$ at $T_P\!\uparrow\!1$
      may now converge to non-empty sets, so the semantics
      distinguishes ``definitionally empty'' from
      ``convergent after further iteration,''
      and $M{\Downarrow}_\lambda$ implies
      $\mathcal{T}(M){\Downarrow}_{\mathrm{lfp}}$,
      but not conversely.
  \end{description}
\end{remark}

\subsection{Adequacy}

\begin{definition}[Inheritance-convergence]\label{def:convergence}
  A mixin tree $T$ \emph{converges}, written $T{\Downarrow}$,
  if there exists $n \ge 0$ such that:
  \[
    \{\mathrm{argument},\, \mathrm{result}\}
    \;\subseteq\;
    \properties\!\bigl(\,
      \underbrace{\text{root} \snoc \mathrm{result}
      \snoc \cdots \snoc \mathrm{result}}_{n}
    \,\bigr)
  \]
  in the semantics of
  Section~\ref{sec:mixin-trees}.
\end{definition}

\begin{theorem}[Adequacy]\label{thm:adequacy}
  A closed ANF $\lambda$-term $M$ has a head normal form if and
  only if $\mathcal{T}(M){\Downarrow}$.
\end{theorem}

A closed ANF term at the top level is either an abstraction
(a value) or a computation (an application
  $\mathbf{let}$-binding, a value $\mathbf{let}$-binding,
or a tail call).
An abstraction translates to a record whose root immediately
has the abstraction shape; a computation's root has the form
$\{x, \mathrm{result}\}$ or
$\{\mathrm{tailCall}, \mathrm{result}\}$, and one must follow
the $\mathrm{result}$ chain to find the eventual value.

\begin{proof}[Proof of Theorem~\ref{thm:adequacy}]
  ($\Rightarrow$)
  Suppose $M \to^*_h \lambda x.\, M'$ in $k$ head reduction
  steps.
  We show $\mathcal{T}(M)$ converges at depth $\le k$ by
  induction on $k$.

  \paragraph{Base case ($k = 0$)}
  $M$ is already an abstraction
  $\lambda x.\, M'$.
  Then $\mathcal{T}(M) =
  \{\mathrm{argument} \mapsto \{\},\;
  \mathrm{result} \mapsto \mathcal{T}(M')\}$,
  whose root has
  $\defines = \{\mathrm{argument}, \mathrm{result}\}$.
  These labels are directly in $\defines$ of the root, so
  the recursive evaluation returns them without further
  recursion, and $\mathcal{T}(M)$ converges at depth $n = 0$.

  \paragraph{Inductive step ($k \ge 1$)}
  $M$ is not an abstraction,
  so $M$ is either a tail call, an application
  $\mathbf{let}$-binding, or a value
  $\mathbf{let}$-binding.
  Each form corresponds to one head $\beta$-step in the
  underlying $\lambda$-term.
  Let $N$ be the head reduct, so
  $M \to_h N \to^{k-1}_h \lambda x.\, M'$.
  Note that $N$ is closed: head reduction substitutes a
  closed value into a term whose only free variable is the
  substitution target, preserving closedness.
  By Lemma~\ref{lem:result-step} (tail-call,
    application let-binding, or value let-binding
  variant), the properties at
  $\text{root} \snoc \mathrm{result} \snoc p$
  in $\mathcal{T}(M)$ coincide with those at
  $\text{root} \snoc p$ in $\mathcal{T}(N)$.
  By the induction hypothesis, $\mathcal{T}(N)$
  converges at depth $\le k - 1$, so $\mathcal{T}(M)$
  converges at depth $\le k$.

  \medskip\noindent
  ($\Leftarrow$)
  Suppose $\mathcal{T}(M){\Downarrow}$ at depth $n$.
  We show $M$ has a head normal form by induction on $n$.

  \paragraph{Base case ($n = 0$)}
  The root of $\mathcal{T}(M)$ has
  abstraction shape
  $\{\mathrm{argument}, \mathrm{result}\}$.
  Only the abstraction rule $\mathcal{T}(\lambda x.\, M')$
  produces a root with own property $\mathrm{argument}$.
  (For tail calls and $\mathbf{let}$-bindings, the root has
    $\defines = \{\mathrm{tailCall}, \mathrm{result}\}$
    or $\{x, \mathrm{result}\}$, and $\mathrm{argument}$ cannot
    appear via $\supers$ because
    $\inherits = \varnothing$ at the root, giving
  $\bases(\text{root}) = \varnothing$.)
  Therefore $M$ is an abstraction, hence already in head normal
  form.

  \paragraph{Inductive step ($n \ge 1$)}
  The abstraction shape
  appears at depth $n$ but not at depth $0$.
  By the base case argument, $M$ is not an abstraction, so at the
  head position of this closed top-level ANF term, $M$ is either
  a tail call, an application $\mathbf{let}$-binding, or a value
  $\mathbf{let}$-binding.
  In each case, the ANF grammar requires the function
  position to be a value~$V_1$
  ($V ::= x \mid \lambda x.\, M$).
  Since $M$ is closed, $V_1$ cannot be a free variable,
  so $V_1$ must be a $\lambda$-abstraction;
  the head form is therefore a $\beta$-redex.
  By Lemma~\ref{lem:result-step} (tail-call,
    application let-binding, or value let-binding
  variant), the properties at depth
  $n - 1$ in $\mathcal{T}(N)$ match those at depth $n$
  in $\mathcal{T}(M)$, where $N$ is the head reduct.
  So $\mathcal{T}(N)$ converges at depth $n - 1$.
  Since $M$ is closed and head reduction preserves
  closedness (it substitutes a closed value for the only
  free variable), $N$ is a closed ANF term.
  By the induction hypothesis, $N$ has a head normal
  form, and since $M \to_h N$, so does $M$.
\end{proof}

\subsection{First-Order Semantics of the Lazy \texorpdfstring{$\lambda$}{λ}-Calculus}

Adequacy relates a single term's head normal form to its
inheritance-convergence.
The following theorem lifts this to an equivalence between
terms, establishing that the translation $\mathcal{T}$
gives the lazy $\lambda$-calculus a first-order semantics
that preserves and reflects B\"ohm tree equivalence.

\begin{definition}[$\lambda$-contextual equivalence under $\mathcal{T}$]%
  \label{def:ctx-equiv}
  For closed $\lambda$-terms $M$ and $N$, write
  $\mathcal{T}(M) \approx_\lambda \mathcal{T}(N)$ if for every
  closing $\lambda$-calculus context $C[\cdot]$:
  \[
    \mathcal{T}(C[M]){\Downarrow}
    \;\Longleftrightarrow\;
    \mathcal{T}(C[N]){\Downarrow}
  \]
\end{definition}

\noindent
The subscript $\lambda$ emphasizes that the quantification ranges
over $\lambda$-calculus contexts, not over all inheritance-calculus contexts.
This definition observes only convergence\footnote{
  Convergence provides exactly a single bit of information.
}, but
quantifying over all $\lambda$-contexts makes it a fine-grained
equivalence: different contexts can supply different arguments
and project different results, probing every aspect of the
term's behavior within the $\lambda$-definable fragment.

\begin{theorem}[First-order semantics of the lazy $\lambda$-calculus]%
  \label{thm:first-order-semantics}
  For closed $\lambda$-terms $M$ and $N$:
  \[
    \mathcal{T}(M) \approx_\lambda \mathcal{T}(N)
    \quad\Longleftrightarrow\quad
    \mathrm{BT}(M) = \mathrm{BT}(N)
  \]
\end{theorem}

\begin{proof}[Proof of Theorem~\ref{thm:first-order-semantics}]
  The translation $\mathcal{T}$ is compositional: for every
  $\lambda$-calculus context $C[\cdot]$, the mixin tree
  $\mathcal{T}(C[M])$ is determined by the mixin tree
  $\mathcal{T}(M)$ and the translation of the context.
  This means $\mathcal{T}$ preserves the context structure
  needed for the argument.

  \medskip\noindent
  ($\Leftarrow$)
  Suppose $\mathrm{BT}(M) = \mathrm{BT}(N)$.
  B\"ohm tree equivalence is a congruence~\cite{barendregt1984-lambda-calculus},
  so for every context $C[\cdot]$,
  $\mathrm{BT}(C[M]) = \mathrm{BT}(C[N])$.
  In particular, $C[M]$ has a head normal form iff $C[N]$ does.
  By Theorem~\ref{thm:adequacy},
  $\mathcal{T}(C[M]){\Downarrow}$ iff
  $\mathcal{T}(C[N]){\Downarrow}$.
  Hence $\mathcal{T}(M) \approx_\lambda \mathcal{T}(N)$.

  \medskip\noindent
  ($\Rightarrow$)
  Suppose $\mathrm{BT}(M) \neq \mathrm{BT}(N)$.
  B\"ohm tree equivalence coincides with observational
  equivalence for the lazy
  $\lambda$-calculus~\cite{abramsky1993-full-abstraction-lazy-lambda,abramsky1995-games-full-abstraction-lazy-lambda}:
  there exists a context $C[\cdot]$ such that $C[M]$ converges
  and $C[N]$ diverges, or vice versa.
  By Theorem~\ref{thm:adequacy},
  $\mathcal{T}(C[M]){\Downarrow}$ and
  $\mathcal{T}(C[N]){\Uparrow}$.
  Hence $\mathcal{T}(M) \not\approx_\lambda \mathcal{T}(N)$.
\end{proof}

\noindent
The proof rests on two pillars: Adequacy
(Theorem~\ref{thm:adequacy}), which is our contribution, and
the classical result that B\"ohm tree equivalence equals
observational equivalence for the lazy
$\lambda$-calculus~\cite{abramsky1993-full-abstraction-lazy-lambda,abramsky1995-games-full-abstraction-lazy-lambda}.
Neither pillar requires defining or comparing
absolute paths inside mixin trees.

\section{\bilingual{Expressive Asymmetry: Proofs}{表达力的不对称:证明}}
\label{app:expressiveness}

\begin{definition}[\bilingual{Macro-expressibility, after Felleisen~\cite{felleisen1991-expressive-power}}{宏可表达性,据 Felleisen~\cite{felleisen1991-expressive-power}}]
  \label{def:macro-expressibility}
  \ifInheritanceChinese
  语言 $\mathscr{L}_1$ 的一个构造 $C$ 在语言 $\mathscr{L}_0$ 中是\emph{宏可表达的},若存在一个翻译 $\mathcal{E}$,使得对每一个含有 $C$ 的出现的程序 $P$,翻译 $\mathcal{E}(P)$ 是通过将 $C$ 的每次出现替换为一个仅依赖于 $C$ 的子表达式的 $\mathscr{L}_0$ 表达式而得到的,并保持 $P$ 的所有其他构造不变。
  \else
  A construct $C$ of language $\mathscr{L}_1$ is \emph{macro-expressible}
  in language $\mathscr{L}_0$ if there exists a translation $\mathcal{E}$
  such that for every program $P$ containing occurrences of $C$,
  the translation $\mathcal{E}(P)$ is obtained by replacing each
  occurrence of $C$ with an $\mathscr{L}_0$ expression that depends only
  on $C$'s subexpressions, leaving all other constructs of $P$
  unchanged.
  \fi
\end{definition}

\begin{lemma}[\bilingual{Composition of macro-expressible translations}{宏可表达翻译的复合}]
  \label{lem:macro-composition}
  \ifInheritanceChinese
  若 $\mathcal{E}_1$ 在 $\mathscr{L}_0$ 中宏可表达 $\mathscr{L}_1$ 的构造,且 $\mathcal{E}_2$ 在 $\mathscr{L}_1$ 中宏可表达 $\mathscr{L}_2$ 的构造,则复合 $\mathcal{E}_1 \circ \mathcal{E}_2$ 在 $\mathscr{L}_0$ 中宏可表达 $\mathscr{L}_2$ 的构造。
  \else
  If $\mathcal{E}_1$ macro-expresses the constructs of
  $\mathscr{L}_1$ in $\mathscr{L}_0$ and
  $\mathcal{E}_2$ macro-expresses the constructs of
  $\mathscr{L}_2$ in $\mathscr{L}_1$, then the composition
  $\mathcal{E}_1 \circ \mathcal{E}_2$ macro-expresses the
  constructs of $\mathscr{L}_2$ in $\mathscr{L}_0$.
  \fi
\end{lemma}

\begin{proof}
  \ifInheritanceChinese
  $\mathscr{L}_2$ 的每个构造 $F$ 在 $\mathscr{L}_1$ 上有一个固定的语法抽象 $A_F^{(2)}$。
  出现在 $A_F^{(2)}$ 中的每个 $\mathscr{L}_1$-构造 $G$ 在 $\mathscr{L}_0$ 上有一个固定的语法抽象 $A_G^{(1)}$。
  将 $A_F^{(2)}$ 中的每个 $G$ 替换为 $A_G^{(1)}$,得到 $F$ 在 $\mathscr{L}_0$ 上的一个固定语法抽象,从而满足 E4。
  \else
  Each construct~$F$ of~$\mathscr{L}_2$ has a fixed
  syntactic abstraction~$A_F^{(2)}$
  over~$\mathscr{L}_1$.
  Each $\mathscr{L}_1$-construct~$G$ appearing
  in~$A_F^{(2)}$ has a fixed
  syntactic abstraction~$A_G^{(1)}$
  over~$\mathscr{L}_0$.
  Substituting~$A_G^{(1)}$ for every~$G$
  in~$A_F^{(2)}$ yields a fixed syntactic
  abstraction over~$\mathscr{L}_0$ for~$F$,
  satisfying~E4.
  \fi
\end{proof}

\begin{definition}[\bilingual{Language universe}{语言全集}]
  \label{def:language-universe}
  \ifInheritanceChinese
  设 $\mathscr{U}$ 为一个语言,其构造子是继承演算的构造子、惰性 $\lambda$-演算的构造子(包括 $\lambda$-抽象、应用与变量)以及 $\mathbf{let}$-绑定的并集。
  $\mathscr{U}$ 的语义通过将每个 $\mathscr{U}$-程序翻译为继承演算来给出,分三步:
  \begin{enumerate}
    \item 将每个极大 $\lambda$-演算子表达式(即根节点与所有内部节点仅使用 $\lambda$-演算构造子的每棵子树)转化为 A-范式~\cite{flanagan1993-essence-compiling-continuations}。
    \item 将翻译 $\mathcal{T}$(第~\ref{sec:forward-translation} 节)应用于每个 ANF $\lambda$-演算子表达式。
    \item 按第~\ref{sec:mixin-trees} 节的规则对所得继承演算程序求值。
  \end{enumerate}
  当一个 $\mathscr{U}$-表达式交错使用 $\lambda$-演算与继承演算的构造子(例如,一个 $\mathbf{let}$-绑定,其被绑定的值是一个继承演算表达式)时,翻译从叶到根组合地应用:每个 $\lambda$-演算构造子通过 $\mathcal{T}$ 替换为其继承演算等价物,同时将继承演算子表达式视为不透明的值。
  由于 $\mathcal{T}$ 的每条规则(第~\ref{sec:forward-translation} 节)都是一个固定的语法抽象,仅依赖于被翻译的子表达式,因此这种自底向上的应用是良定义的,并为步骤~(3) 产生一个纯继承演算程序。

  若一个 $\mathscr{U}$-程序 $P$ 的翻译收敛,则 $P$ 收敛,记为 $\mathrm{eval}_{\mathscr{U}}(P)$。
  \else
  Let $\mathscr{U}$ be the language whose constructors are
  the union of the constructors of inheritance-calculus,
  the constructors of the lazy $\lambda$-calculus
  , including $\lambda$-abstraction, application, and variable, and
  $\mathbf{let}$-binding.
  The semantics of $\mathscr{U}$ is given by translating
  every $\mathscr{U}$-program to inheritance-calculus in
  three steps:
  \begin{enumerate}
    \item Convert every maximal $\lambda$-calculus
      subexpression (that is, every subtree whose root
        and all internal nodes use only $\lambda$-calculus
      constructors) to A-normal
      form~\cite{flanagan1993-essence-compiling-continuations}.
    \item Apply the translation $\mathcal{T}$
      (Section~\ref{sec:forward-translation}) to every ANF
      $\lambda$-calculus subexpression.
    \item Evaluate the resulting inheritance-calculus
      program by the rules of
      Section~\ref{sec:mixin-trees}.
  \end{enumerate}
  When a $\mathscr{U}$-expression interleaves
  $\lambda$-calculus and inheritance-calculus constructors
  (e.g.\ a $\mathbf{let}$-binding whose bound value is an
  inheritance-calculus expression), the translation is
  applied compositionally from leaves to root: each
  $\lambda$-calculus construct is replaced by its
  inheritance-calculus equivalent via~$\mathcal{T}$,
  treating inheritance-calculus subexpressions as opaque
  values.
  Since each rule of~$\mathcal{T}$
  (Section~\ref{sec:forward-translation}) is a fixed
  syntactic abstraction depending only on the translated
  subexpressions, this bottom-up application is
  well-defined and yields a pure inheritance-calculus
  program for step~(3).

  A $\mathscr{U}$-program $P$ converges,
  $\mathrm{eval}_{\mathscr{U}}(P)$, iff its translation
  converges.
  \fi
\end{definition}

\begin{remark}\label{rem:universe-properties}
  \ifInheritanceChinese
  惰性 $\lambda$-演算与继承演算是 $\mathscr{U}$ 的保守性限制:各自通过去掉另一方的构造子来得到。
  对于纯继承演算程序,步骤~(1) 与~(2) 是平凡的,因此 $\mathrm{eval}_{\mathscr{U}}$ 与继承演算的语义一致。
  对于纯 $\lambda$-演算程序,步骤~(1)--(3) 的复合由充分性(定理~\ref{thm:adequacy})与惰性求值一致。
  \else
  The lazy $\lambda$-calculus and inheritance-calculus are
  conservative restrictions of $\mathscr{U}$: each is
  obtained by removing the other's constructors.
  For a pure inheritance-calculus program, steps~(1)
  and~(2) are vacuous, so
  $\mathrm{eval}_{\mathscr{U}}$ agrees with the
  inheritance-calculus semantics.
  For a pure $\lambda$-calculus program, the composition
  of steps~(1)--(3) agrees with lazy evaluation by
  Adequacy (Theorem~\ref{thm:adequacy}).
  \fi
\end{remark}

\ifInheritanceChinese
三步翻译是有效的(ANF 转换与 $\mathcal{T}$ 是语法变换),而继承演算的递归求值(第~\ref{sec:mixin-trees} 节)是一个半判定过程,因此 $\mathrm{eval}_{\mathscr{U}}$ 是递归可枚举的,满足 Felleisen 框架~\cite{felleisen1991-expressive-power} 的要求。
\else
The three-step translation is effective (ANF conversion
and~$\mathcal{T}$ are syntactic transformations) and the
recursive evaluation of inheritance-calculus
(Section~\ref{sec:mixin-trees}) is a semi-decision
procedure, so $\mathrm{eval}_{\mathscr{U}}$ is
recursively enumerable, as required by
Felleisen's framework~\cite{felleisen1991-expressive-power}.
\fi

\begin{lemma}[\bilingual{$\mathbf{let}$-binding is macro-expressible
  in the lazy $\lambda$-calculus}{$\mathbf{let}$-绑定在惰性 $\lambda$-演算中是宏可表达的}]
  \label{lem:let-macro}
  \ifInheritanceChinese
  翻译 $\mathbf{let}\; x = e_1 \;\mathbf{in}\; e_2 \;\mapsto\; (\lambda x.\, e_2)\; e_1$ 是 $\mathbf{let}$ 在惰性 $\lambda$-演算中的一个宏可表达的消去。
  \else
  The translation
  $\mathbf{let}\; x = e_1 \;\mathbf{in}\; e_2
  \;\mapsto\; (\lambda x.\, e_2)\; e_1$
  is a macro-expressible elimination of $\mathbf{let}$
  in the lazy $\lambda$-calculus.
  \fi
\end{lemma}

\begin{proof}
  \ifInheritanceChinese
  该翻译将每个 $\mathbf{let}$-绑定替换为一个仅依赖于两个子表达式 $e_1$ 与 $e_2$ 的 $\lambda$-表达式。
  这是 Felleisen~\cite{felleisen1991-expressive-power} 的命题 3.7。
  \else
  The translation replaces each $\mathbf{let}$-binding
  with a $\lambda$-expression that depends only on the
  two subexpressions $e_1$ and~$e_2$.
  This is Proposition~3.7 of
  Felleisen~\cite{felleisen1991-expressive-power}.
  \fi
\end{proof}

\begin{lemma}[\bilingual{Equal expressiveness of the lazy
  $\lambda$-calculus and its ANF fragment}{惰性 $\lambda$-演算与其 ANF 片段的等同表达力}]
  \label{lem:anf-equiv}
  \ifInheritanceChinese
  在 $\mathscr{U}$ 中,惰性 $\lambda$-演算与 ANF 惰性 $\lambda$-演算(第~\ref{sec:forward-translation} 节)具有相同的宏可表达力。
  \else
  In $\mathscr{U}$, the lazy $\lambda$-calculus and the
  ANF lazy $\lambda$-calculus
  (Section~\ref{sec:forward-translation}) are equally
  macro-expressive.
  \fi
\end{lemma}

\begin{proof}
  \ifInheritanceChinese
  每个 ANF 项都是一个 $\lambda$-项,因此 ANF 片段嵌入完整 $\lambda$-演算的映射是恒等变换,平凡地宏可表达。
  对于反方向,ANF 变换~\cite{flanagan1993-essence-compiling-continuations} 由两步构成,每步都是宏可表达的:
  \begin{enumerate}
    \item \emph{\bilingual{Application flattening}{应用展开}。}
      将每个嵌套应用 $(e_1\; e_2)\; e_3$ 替换为 $\mathbf{let}\; x = e_1\; e_2 \;\mathbf{in}\; x\; e_3$。
      其语法抽象为 $A(\cdot_1, \cdot_2, \cdot_3) = \mathbf{let}\; x = \cdot_1\; \cdot_2 \;\mathbf{in}\; x\; \cdot_3$,仅依赖三个子表达式。
    \item \emph{\bilingual{Value lifting}{值提升}。}
      将应用位置上每个自由值 $V$ 替换为 $\mathbf{let}\; x = V \;\mathbf{in}\; x$。
      其语法抽象为 $A(\cdot_1) = \mathbf{let}\; x = \cdot_1 \;\mathbf{in}\; x$。
  \end{enumerate}
  由引理~\ref{lem:let-macro},每个所得 $\mathbf{let}$-绑定本身在 $\lambda$-演算中是宏可表达的。
  由引理~\ref{lem:macro-composition},复合这些宏可表达翻译,得到一个从完整 $\lambda$-演算到其 ANF 片段的宏可表达翻译。
  \else
  Every ANF term is a $\lambda$-term, so the embedding
  of the ANF fragment into the full $\lambda$-calculus is
  the identity, which is trivially macro-expressible.
  For the reverse direction, the ANF
  transformation~\cite{flanagan1993-essence-compiling-continuations} consists of
  two steps, each macro-expressible:
  \begin{enumerate}
    \item \emph{Application flattening.}
      Replace each nested application
      $(e_1\; e_2)\; e_3$ with
      $\mathbf{let}\; x = e_1\; e_2 \;\mathbf{in}\;
      x\; e_3$.
      The syntactic abstraction is
      $A(\cdot_1, \cdot_2, \cdot_3)
      = \mathbf{let}\; x = \cdot_1\; \cdot_2
      \;\mathbf{in}\; x\; \cdot_3$,
      depending only on the three subexpressions.
    \item \emph{Value lifting.}
      Replace each free value~$V$ used in an application
      position with
      $\mathbf{let}\; x = V \;\mathbf{in}\; x$.
      The syntactic abstraction is
      $A(\cdot_1) = \mathbf{let}\; x = \cdot_1
      \;\mathbf{in}\; x$.
  \end{enumerate}
  By Lemma~\ref{lem:let-macro}, each resulting
  $\mathbf{let}$-binding is itself
  macro-expressible in the $\lambda$-calculus.
  By Lemma~\ref{lem:macro-composition}, composing
  these macro-expressible translations
  yields a macro-expressible translation from the full
  $\lambda$-calculus to its ANF fragment.
  \fi
\end{proof}

\begin{theorem}[\bilingual{Forward macro-expressibility}{正向宏可表达性}]
  \label{thm:forward-macro}
  \ifInheritanceChinese
  翻译 $\mathcal{T}$(第~\ref{sec:forward-translation} 节)是 $\lambda$-演算嵌入继承演算的一个宏可表达的嵌入。
  \else
  The translation $\mathcal{T}$
  (Section~\ref{sec:forward-translation}) is a macro-expressible
  embedding of the $\lambda$-calculus into inheritance-calculus.
  \fi
\end{theorem}

\begin{proof}
  \ifInheritanceChinese
  我们验证 Felleisen~\cite[Definition~3.11]{felleisen1991-expressive-power} 的条件 E4:对于每个元数为 $a$ 的 $\lambda$-演算构造子 $F$,必须存在一个固定的 $a$ 元语法抽象 $A_F(\cdot_1,\ldots,\cdot_a)$ 使得 $\mathcal{T}(F(e_1,\ldots,e_a)) = A_F(\mathcal{T}(e_1),\ldots,\mathcal{T}(e_a))$。
  我们逐条检验 $\mathcal{T}$ 的规则(第~\ref{sec:forward-translation} 节):
  \begin{enumerate}
    \item \emph{\bilingual{$\lambda$-bound variable}{$\lambda$-绑定变量}}($x$,de~Bruijn 索引为 $n'$)。
      这是一个零元构造子,因此没有子表达式。
      语法抽象是常量 $A = \dbi{n'}.\mathrm{argument}$。
    \item \emph{\bilingual{Let-bound variable}{Let-绑定变量}}($x$)。
      零元。$A = x.\mathrm{result}$。
    \item \emph{\bilingual{$\lambda$-abstraction}{$\lambda$-抽象}},记为 $\lambda x.\, M$。
      关于体 $M$ 是一元的。
      $A(\cdot_1) = \{\mathrm{argument} \mapsto \{\},\; \mathrm{result} \mapsto \cdot_1\}$。
    \item \emph{\bilingual{Application let-binding}{应用 Let-绑定}}($\mathbf{let}\; x = V_1\; V_2 \;\mathbf{in}\; M$)。
      关于 $V_1$、$V_2$ 与 $M$ 是三元的。
      $A(\cdot_1, \cdot_2, \cdot_3) = \{x \mapsto \{\cdot_1,\; \mathrm{argument} \mapsto \cdot_2\},\; \mathrm{result} \mapsto \cdot_3\}$。
    \item \emph{\bilingual{Value let-binding}{值 Let-绑定}}($\mathbf{let}\; x = V \;\mathbf{in}\; M$)。
      关于 $V$ 与 $M$ 是二元的。
      $A(\cdot_1, \cdot_2) = \{x \mapsto \{\mathrm{result} \mapsto \cdot_1\},\; \mathrm{result} \mapsto \cdot_2\}$。
    \item \emph{\bilingual{Tail call}{尾调用}}($V_1\; V_2$)。
      二元。\newline
      $A(\cdot_1, \cdot_2) = \{\mathrm{tailCall} \mapsto \{\cdot_1,\; \mathrm{argument} \mapsto \cdot_2\},\; \mathrm{result} \mapsto \{\mathrm{tailCall}.\mathrm{result}\}\}$。
  \end{enumerate}
  在每种情况下,语法抽象 $A_F$ 都是一个固定的继承演算上下文,其空位由子表达式的递归翻译填充,从而满足 E4。
  \else
  We verify condition~E4 of
  Felleisen~\cite[Definition~3.11]{felleisen1991-expressive-power}:
  for each $\lambda$-calculus construct~$F$ of arity~$a$,
  there must exist a fixed $a$-ary syntactic abstraction
  $A_F(\cdot_1,\ldots,\cdot_a)$ over
  inheritance-calculus such that
  $\mathcal{T}(F(e_1,\ldots,e_a)) =
  A_F(\mathcal{T}(e_1),\ldots,\mathcal{T}(e_a))$.
  We check each rule of $\mathcal{T}$
  (Section~\ref{sec:forward-translation}):
  \begin{enumerate}
    \item \emph{$\lambda$-bound variable}
      ($x$ with de~Bruijn index $n'$).
      This is a nullary construct and therefore has no subexpressions.
      The syntactic abstraction is the constant
      $A = \dbi{n'}.\mathrm{argument}$.
    \item \emph{Let-bound variable} ($x$).
      Nullary.
      $A = x.\mathrm{result}$.
    \item \emph{$\lambda$-abstraction}, denoted $\lambda x.\, M$.
      Unary in the body~$M$.
      $A(\cdot_1) = \{\mathrm{argument} \mapsto \{\},\;
      \mathrm{result} \mapsto \cdot_1\}$.
    \item \emph{Application let-binding}
      ($\mathbf{let}\; x = V_1\; V_2
      \;\mathbf{in}\; M$).
      Ternary in~$V_1$, $V_2$, and~$M$.
      $A(\cdot_1, \cdot_2, \cdot_3) =
      \{x \mapsto \{\cdot_1,\;
        \mathrm{argument} \mapsto \cdot_2\},\;
      \mathrm{result} \mapsto \cdot_3\}$.
    \item \emph{Value let-binding}
      ($\mathbf{let}\; x = V \;\mathbf{in}\; M$).
      Binary in~$V$ and~$M$.
      $A(\cdot_1, \cdot_2) =
      \{x \mapsto \{\mathrm{result} \mapsto \cdot_1\},\;
      \mathrm{result} \mapsto \cdot_2\}$.
    \item \emph{Tail call} ($V_1\; V_2$).
      Binary.\newline
      $A(\cdot_1, \cdot_2) =
      \{\mathrm{tailCall} \mapsto
        \{\cdot_1,\;
        \mathrm{argument} \mapsto \cdot_2\},\;
        \mathrm{result} \mapsto
      \{\mathrm{tailCall}.\mathrm{result}\}\}$.
  \end{enumerate}
  In each case the syntactic abstraction~$A_F$ is a fixed
  inheritance-calculus context whose holes are filled
  by the recursive translations of the subexpressions,
  satisfying~E4.
  \fi
\end{proof}

\begin{theorem}[\bilingual{Nonexpressibility of inheritance}{继承的不可宏表达性}]
  \label{thm:cartesian-nonexpressibility}
  \ifInheritanceChinese
  惰性 $\lambda$-演算不能宏可表达 $\mathscr{U}$(定义~\ref{def:language-universe})的继承构造子。
  \else
  The lazy $\lambda$-calculus cannot macro-express the
  inheritance constructors of $\mathscr{U}$
  (Definition~\ref{def:language-universe}).
  \fi
\end{theorem}

\begin{proof}
  \ifInheritanceChinese
  我们对 $L_0 = \text{惰性 $\lambda$-演算}$,$L_1 = \mathscr{U}$ 应用 Felleisen~\cite{felleisen1991-expressive-power} 的定理 3.14(i)。
  由定义~\ref{def:language-universe},$\mathscr{U} = L_0 + \{$mixin-树构造子$\}$ 是 $L_0$ 的一个保守性扩展(Felleisen~\cite{felleisen1991-expressive-power} 的定义 3.2)。
  我们验证四个条件:
  \begin{enumerate}
    \item[(i)] \emph{$L_0 \subseteq L_1$:} $\mathscr{U}$ 的语法包含所有 $L_0$-短语。
    \item[(ii)] \emph{不引入新的 $L_0$-短语:} mixin-树构造子不在 $L_0$ 中引入新短语。
    \item[(iii)] \emph{构造子不相交:} mixin-树构造子与 $L_0$-构造子不相交。
    \item[(iv)] \emph{语义保持:} 对纯 $\lambda$-项,定义~\ref{def:language-universe} 的步骤~(1) 与~(2) 是平凡的,步骤~(3) 由充分性(定理~\ref{thm:adequacy})给出惰性求值,因此 $\mathrm{eval}_{\mathscr{U}}$ 在 $L_0$-短语上与 $\mathrm{eval}_{L_0}$ 一致。
  \end{enumerate}
  该定理要求我们给出两个 $L_0$-项 $M$ 与 $N$,满足 $M \cong_0 N$ 但 $M \not\cong_1 N$,其中 $\cong_0$ 是 $L_0$ 中的操作等价,$\cong_1$ 是 $\mathscr{U}$ 中的操作等价。

  \paragraph{\bilingual{Operational equivalence in $L_0$}{$L_0$ 中的操作等价}}
  由 Böhm 树等价~\cite{abramsky1993-full-abstraction-lazy-lambda,abramsky1995-games-full-abstraction-lazy-lambda}(定理~\ref{thm:first-order-semantics}),对于封闭 $\lambda$-项 $M$ 与 $N$:$M \cong_0 N$ 当且仅当 $\mathrm{BT}(M) = \mathrm{BT}(N)$。

  \paragraph{\bilingual{Separating pair}{分离对}}
  定义 Church 布尔值与 Church 布尔相等:$\mathrm{true} = \lambda t.\,\lambda f.\, t$,$\mathrm{false} = \lambda t.\,\lambda f.\, f$,以及 $\mathrm{eq} = \lambda a.\,\lambda b.\, (a\;b)\;(b\;\mathrm{false}\;\mathrm{true})$。
  令
  \[
    M = \mathrm{eq}\;\mathrm{false}\;\mathrm{false}
    \qquad
    N = \mathrm{eq}\;\mathrm{true}\;\mathrm{true}.
  \]
  两者在惰性 $\lambda$-演算中均归约为 Church~$\mathrm{true}$,因此 $\mathrm{BT}(M) = \mathrm{BT}(N)$,从而 $M \cong_0 N$。

  \paragraph{\bilingual{Distinguishing $\mathscr{U}$-context}{区分 $\mathscr{U}$-上下文}}
  我们现在构造一个区分 $M$ 与 $N$ 的 $\mathscr{U}$-上下文。
  $M$ 与 $N$ 是语法封闭的 $\lambda$-项;由定义~\ref{def:language-universe},它们在 $\mathscr{U}$ 中的语义通过 ANF 转换后接 $\mathcal{T}$ 将其翻译为继承演算来给出。
  在扩展的 ANF 约定(第~\ref{sec:forward-translation} 节)下,应用位置上的每个自由值首先被绑定到一个名字:
  \else
  We apply Theorem~3.14(i) of
  Felleisen~\cite{felleisen1991-expressive-power} to the pair
  $L_0 = \text{lazy $\lambda$-calculus}$,
  $L_1 = \mathscr{U}$.
  By Definition~\ref{def:language-universe},
  $\mathscr{U} = L_0 + \{$mixin-tree constructors$\}$
  is a conservative extension of~$L_0$
  (Definition~3.2 of
  Felleisen~\cite{felleisen1991-expressive-power}).
  We verify the four conditions:
  \begin{enumerate}
    \item[(i)] \emph{$L_0 \subseteq L_1$:}
      The syntax of~$\mathscr{U}$ includes all
      $L_0$-phrases.
    \item[(ii)] \emph{No new $L_0$-phrases:}
      The mixin-tree constructors do not introduce new
      phrases in~$L_0$.
    \item[(iii)] \emph{Disjoint constructors:}
      The mixin-tree constructors are disjoint from the
      $L_0$-constructors.
    \item[(iv)] \emph{Semantics preserved:}
      For pure $\lambda$-terms, steps~(1) and~(2) of
      Definition~\ref{def:language-universe} are vacuous
      and step~(3) yields lazy evaluation by Adequacy
      (Theorem~\ref{thm:adequacy}), so
      $\mathrm{eval}_{\mathscr{U}}$ agrees with
      $\mathrm{eval}_{L_0}$ on $L_0$-phrases.
  \end{enumerate}
  The theorem requires us to exhibit two $L_0$-terms
  $M$ and $N$ with $M \cong_0 N$ but
  $M \not\cong_1 N$, where $\cong_0$ is operational
  equivalence in $L_0$ and $\cong_1$ is operational
  equivalence in $\mathscr{U}$.

  \paragraph{Operational equivalence in $L_0$}
  By B\"ohm tree
  equivalence~\cite{abramsky1993-full-abstraction-lazy-lambda,abramsky1995-games-full-abstraction-lazy-lambda}
  (Theorem~\ref{thm:first-order-semantics}), for closed
  $\lambda$-terms $M$ and $N$:
  $M \cong_0 N$ iff $\mathrm{BT}(M) = \mathrm{BT}(N)$.

  \paragraph{Separating pair}
  Define Church booleans and Church boolean equality:
  $\mathrm{true} = \lambda t.\,\lambda f.\, t$,
  $\mathrm{false} = \lambda t.\,\lambda f.\, f$, and
  $\mathrm{eq} = \lambda a.\,\lambda b.\,
  (a\;b)\;(b\;\mathrm{false}\;\mathrm{true})$.
  Let
  \[
    M = \mathrm{eq}\;\mathrm{false}\;\mathrm{false}
    \qquad
    N = \mathrm{eq}\;\mathrm{true}\;\mathrm{true}.
  \]
  Both reduce to Church~$\mathrm{true}$ in the
  lazy $\lambda$-calculus, so
  $\mathrm{BT}(M) = \mathrm{BT}(N)$ and hence
  $M \cong_0 N$.

  \paragraph{Distinguishing $\mathscr{U}$-context}
  We now construct a $\mathscr{U}$-context that
  separates $M$ and $N$.
  $M$ and $N$ are syntactically closed $\lambda$-terms;
  by Definition~\ref{def:language-universe}, their
  semantics in $\mathscr{U}$ is given by translating
  them to inheritance-calculus via ANF conversion followed
  by~$\mathcal{T}$.
  Under the extended ANF convention
  (Section~\ref{sec:forward-translation}), every free
  value used in an application is first bound to a name:
  \fi
  \begin{align*}
    \mathrm{ANF}(M) ={} &
    \mathbf{let}\;\mathrm{eq} =
    \lambda a.\,\lambda b.\,
    \mathbf{let}\; x_1 = a\;b \;\mathbf{in}\; \\
    &\qquad
    \mathbf{let}\; f_1 = \mathrm{false}
    \;\mathbf{in}\;
    \mathbf{let}\; x_2 = b\;f_1 \;\mathbf{in}\; \\
    &\qquad
    \mathbf{let}\; t_1 = \mathrm{true}
    \;\mathbf{in}\;
    \mathbf{let}\; x_3 = x_2\;t_1
    \;\mathbf{in}\;
    x_1\;x_3 \\
    & \mathbf{in}\;\;
    \mathbf{let}\; f_1 = \mathrm{false}
    \;\mathbf{in}\;
    \mathbf{let}\;\mathrm{eqFalse} = \mathrm{eq}\;f_1 \\
    &\quad
    \;\mathbf{in}\;
    \mathbf{let}\; f_2 = \mathrm{false}
    \;\mathbf{in}\;
    \mathrm{eqFalse}\;f_2
  \end{align*}
  \ifInheritanceChinese
  应用 $\mathcal{T}$:
  \else
  Applying~$\mathcal{T}$:
  \fi
  \begin{align*}
    \mathcal{T}(\mathrm{ANF}(M)) \;=\; \{&
      \mathrm{eq} \mapsto \{\mathrm{result} \mapsto
      \mathcal{T}(\lambda a.\,\lambda b.\,\cdots)\}, \\
      & f_1 \mapsto \{\mathrm{result} \mapsto
      \mathcal{T}(\mathrm{false})\}, \\
      & \mathrm{eqFalse} \mapsto
      \{\mathrm{eq}.\mathrm{result},\;
        \mathrm{argument} \mapsto
      \{f_1.\mathrm{result}\}\}, \\
      & f_2 \mapsto \{\mathrm{result} \mapsto
      \mathcal{T}(\mathrm{false})\}, \\
      & \mathrm{tailCall} \mapsto
      \{\mathrm{eqFalse}.\mathrm{result},\;
        \mathrm{argument} \mapsto
      \{f_2.\mathrm{result}\}\}, \\
      & \mathrm{result} \mapsto
      \{\mathrm{tailCall}.\mathrm{result}\}
    \;\}
  \end{align*}
  \ifInheritanceChinese
  对 $N$ 类似,以 $\mathrm{true}$ 替换 $\mathrm{false}$。
  注意 $\mathrm{eq}$ 是 $\mathcal{T}(\mathrm{ANF}(M))$ 的一个具名子项,通过值 let-绑定绑定。
  在 $\mathrm{eq}$ 的翻译内部,$\lambda b$ 作用域包含 $x_1$、$f_1$、$x_2$、$t_1$、$x_3$ 的 let-绑定以及一个尾调用。
  特别地,$\dbi{1}.\mathrm{argument}$ 到达外层 $\lambda a$ 作用域,而 $\dbi{0}.\mathrm{argument}$ 到达 $\lambda b$ 自身。\footnote{
    这些索引使用压缩约定,仅计数 $\lambda$-作用域。
    如翻译脚注(第~\ref{sec:forward-translation} 节)所述,每个 let-绑定与尾调用在继承演算中引入一个额外的作用域层级;形式索引相应地更大,但目标作用域相同。
  }

  $\mathscr{U}$-上下文使用继承在 $\mathrm{eq}$ 的 $\lambda b$ 作用域内添加一个 $\mathrm{repr}$ 记录,该记录投影第一个操作数。
  定义上下文 $C[\cdot]$ 为:
  \else
  and symmetrically for $N$ with $\mathrm{true}$ in place
  of $\mathrm{false}$.
  Note that $\mathrm{eq}$ is a named child
  of $\mathcal{T}(\mathrm{ANF}(M))$, bound via a value
  let-binding.
  Inside the translation of $\mathrm{eq}$, the inner
  $\lambda b$ scope contains the let-bindings
  for $x_1$, $f_1$, $x_2$, $t_1$, $x_3$, and a tail call.
  In particular, $\dbi{1}.\mathrm{argument}$ reaches
  the outer $\lambda a$ scope and
  $\dbi{0}.\mathrm{argument}$ reaches $\lambda b$
  itself.\footnote{
    These indices use the condensed
    convention, counting only $\lambda$-scopes.
    As noted in the translation footnote
    (Section~\ref{sec:forward-translation}),
    each let-binding and tail call introduces an
    additional scope level in inheritance-calculus;
    the formal index is correspondingly larger,
    but the target scope is the same.
  }

  The $\mathscr{U}$-context uses inheritance to add
  a $\mathrm{repr}$ record inside the $\lambda b$ scope
  of $\mathrm{eq}$ that projects the first operand.
  Define the context $C[\cdot]$ as:
  \fi
  \begin{align*}
    C[\cdot] ={} & \mathbf{let}\;
    \mathrm{extended} = \{[\cdot],\;
      \mathrm{eq} \mapsto \{
        \mathrm{result} \mapsto \{ \\
          &\quad
          \mathrm{result} \mapsto \{
            \mathrm{repr} \mapsto \{
              \mathrm{firstOperand}
              \mapsto \{\dbi{2}.\mathrm{argument}\}
    \}\}\}\}\} \\
    & \mathbf{in}\;
    \mathrm{extended}.\mathrm{tailCall}.\mathrm{repr}.\mathrm{firstOperand}
  \end{align*}
  \ifInheritanceChinese
  这里 $[\cdot]$ 是空位。
  注意 mixin 中 $\mathrm{result}$ 的两层嵌套:外层 $\mathrm{result}$ 进入 $\mathrm{eq}$ 的值 let-绑定包装器,而内层 $\mathrm{result}$ 进入 $\lambda a$ 作用域。
  第二个 $\mathrm{result}$ 则包含 $\lambda b$ 作用域,因此 $\mathrm{repr}$ 被放置在 $\lambda b$ 作用域层级。
  mixin $\{\mathrm{eq} \mapsto \{\mathrm{result} \mapsto \{\mathrm{result} \mapsto \{\mathrm{repr} \mapsto \cdots\}\}\}\}$ 是一个继承演算表达式:它仅使用 $\mathscr{U}$ 的 mixin-树构造子。
  外层 $\mathbf{let}$-绑定与投影是 $\lambda$-演算构造(引理~\ref{lem:let-macro})。
  因此 $C[\cdot]$ 是一个合法的 $\mathscr{U}$-上下文。

  $\mathrm{repr}$ 中的 de~Bruijn 索引 $\dbi{2}$ 计数作用域层级(在压缩约定下):$\dbi{0}$ 是 $\mathrm{repr}$ 自身,$\dbi{1}$ 是 $\lambda b$ 作用域,而 $\dbi{2}$ 是 $\lambda a$ 作用域。
  因此 $\mathrm{firstOperand}$ 投影 $\lambda a$ 的参数。

  由于 $\mathrm{eqFalse}$(相应地,$\mathrm{eqTrue}$)继承 $\mathrm{eq}.\mathrm{result}$,而 $\mathrm{tailCall}$ 继承 $\mathrm{eqFalse}.\mathrm{result}$($\lambda b$ 作用域),mixin 在 $\lambda b$ 作用域层级放置的 $\mathrm{repr}$ 记录通过继承链传播到 $\mathrm{tailCall}$。
  在 $C[M]$ 中:$\mathrm{eqFalse}$ 绑定向 $\lambda a$ 作用域提供 $\mathrm{argument} \mapsto \{f_1.\mathrm{result}\}$(即 $\mathrm{false}$),因此 $\mathrm{firstOperand}$ 求值为 $\mathrm{false} = \lambda t.\,\lambda f.\, f$。
  在 $C[N]$ 中:$\mathrm{eqTrue}$ 绑定提供 $\mathrm{argument} \mapsto \mathrm{true}$,因此 $\mathrm{firstOperand}$ 求值为 $\mathrm{true} = \lambda t.\,\lambda f.\, t$。

  定义 $D[\cdot] = (\lambda x.\, x\;\Omega\;I)\; C[\cdot]$。
  则 $\mathrm{eval}_{\mathscr{U}}$ 对 $D[M]$ 成立(因为 $\mathrm{false}\;\Omega\;I \to I$,收敛),但对 $D[N]$ 不成立(因为 $\mathrm{true}\;\Omega\;I \to \Omega$,发散)。
  因此 $M \not\cong_1 N$。

  \paragraph{\bilingual{Conclusion}{结论}}
  我们有 $M \cong_0 N$ 但 $M \not\cong_1 N$,故 ${\cong_0} \neq {(\cong_1|_{L_0})}$。
  由 Felleisen~\cite{felleisen1991-expressive-power} 的定理 3.14(i),惰性 $\lambda$-演算不能宏可表达 $\mathscr{U}$ 的继承构造子。
  \else
  Here $[\cdot]$ is the hole.
  Note the two levels of $\mathrm{result}$ nesting in the
  mixin: the outer $\mathrm{result}$ enters the value
  let-binding wrapper for $\mathrm{eq}$, and the inner
  $\mathrm{result}$ enters the $\lambda a$ scope.
  The second $\mathrm{result}$ then contains the
  $\lambda b$ scope, so $\mathrm{repr}$ is placed at
  the $\lambda b$ scope level.
  The mixin $\{\mathrm{eq} \mapsto \{\mathrm{result}
      \mapsto \{\mathrm{result} \mapsto
  \{\mathrm{repr} \mapsto \cdots\}\}\}\}$ is an
  inheritance-calculus expression: it uses only the
  mixin-tree constructors of $\mathscr{U}$.
  The outer $\mathbf{let}$-binding and projection are
  $\lambda$-calculus constructs (Lemma~\ref{lem:let-macro}).
  Thus $C[\cdot]$ is a valid $\mathscr{U}$-context.

  The de~Bruijn index $\dbi{2}$ in
  $\mathrm{repr}$ counts scope levels
  (in the condensed convention):
  $\dbi{0}$ is $\mathrm{repr}$ itself,
  $\dbi{1}$ is the $\lambda b$ scope, and
  $\dbi{2}$ is the $\lambda a$ scope.
  So $\mathrm{firstOperand}$ projects
  $\lambda a$'s argument.

  Since $\mathrm{eqFalse}$ (respectively $\mathrm{eqTrue}$)
  inherits $\mathrm{eq}.\mathrm{result}$,
  and $\mathrm{tailCall}$ inherits
  $\mathrm{eqFalse}.\mathrm{result}$ (the $\lambda b$
  scope), the $\mathrm{repr}$ record placed by the mixin
  at the $\lambda b$ scope level propagates through
  the inheritance chain to $\mathrm{tailCall}$.
  In $C[M]$: the $\mathrm{eqFalse}$ binding supplies
  $\mathrm{argument} \mapsto \{f_1.\mathrm{result}\}$
  (that is, $\mathrm{false}$) to the $\lambda a$ scope, so
  $\mathrm{firstOperand}$ evaluates to
  $\mathrm{false} = \lambda t.\,\lambda f.\, f$.
  In $C[N]$: the $\mathrm{eqTrue}$ binding supplies
  $\mathrm{argument} \mapsto \mathrm{true}$, so
  $\mathrm{firstOperand}$ evaluates to
  $\mathrm{true} = \lambda t.\,\lambda f.\, t$.

  Define $D[\cdot] = (\lambda x.\, x\;\Omega\;I)\;
  C[\cdot]$.
  Then $\mathrm{eval}_{\mathscr{U}}$ holds for $D[M]$
  (since $\mathrm{false}\;\Omega\;I \to I$, which
  converges)
  but not for $D[N]$
  (since $\mathrm{true}\;\Omega\;I \to \Omega$, which
  diverges).
  Therefore $M \not\cong_1 N$.

  \paragraph{Conclusion}
  We have $M \cong_0 N$ but $M \not\cong_1 N$, so
  ${\cong_0} \neq {(\cong_1|_{L_0})}$.
  By Theorem~3.14(i) of
  Felleisen~\cite{felleisen1991-expressive-power},
  the lazy $\lambda$-calculus cannot macro-express the
  inheritance constructors of~$\mathscr{U}$.
  \fi
\end{proof}

\begin{theorem}[\bilingual{Expressive asymmetry}{表达力的不对称}]
  \label{thm:expressive-asymmetry}
  \ifInheritanceChinese
  在语言全集 $\mathscr{U}$(定义~\ref{def:language-universe})中,继承演算比惰性 $\lambda$-演算具有严格更强的表达力:它能宏可表达 $\lambda$-演算能宏可表达的每个 $\mathscr{U}$-构造,但反之不然(Felleisen~\cite{felleisen1991-expressive-power} 的定义 3.17)。
  \else
  In the language universe $\mathscr{U}$
  (Definition~\ref{def:language-universe}),
  inheritance-calculus is strictly more expressive than
  the lazy $\lambda$-calculus: it can macro-express every
  $\mathscr{U}$-construct that the $\lambda$-calculus
  can, but not conversely
  (Definition~3.17 of
  Felleisen~\cite{felleisen1991-expressive-power}).
  \fi
\end{theorem}

\begin{proof}
  \ifInheritanceChinese
  我们须证明:(a)~继承演算能宏可表达惰性 $\lambda$-演算能宏可表达的每个 $\mathscr{U}$-构造,以及 (b)~反之不成立。

  \paragraph{\bilingual{Part (b): the lazy $\lambda$-calculus $\not\geq$
  inheritance-calculus}{第 (b) 部分:惰性 $\lambda$-演算 $\not\geq$ 继承演算}}
  $\lambda$-演算包含其自身的构造子($\lambda$、应用、变量),并能宏可表达 $\mathbf{let}$-绑定(引理~\ref{lem:let-macro}),但不能宏可表达 $\mathscr{U}$ 的继承构造子(定理~\ref{thm:cartesian-nonexpressibility})。
  继承演算包含那些构造子。
  因此,惰性 $\lambda$-演算对于 $\mathscr{U}$ 而言表达力不及继承演算。

  \paragraph{\bilingual{Part (a): inheritance-calculus $\geq$ the lazy
  $\lambda$-calculus}{第 (a) 部分:继承演算 $\geq$ 惰性 $\lambda$-演算}}
  我们证明继承演算能宏可表达 $\lambda$-演算所包含的或能宏可表达的每个 $\mathscr{U}$-构造。
  $\lambda$-演算包含 $\lambda$-抽象、应用与变量;它能宏可表达 $\mathbf{let}$-绑定(引理~\ref{lem:let-macro})。
  我们须证明继承演算能宏可表达这四个构造。
  翻译分两个阶段进行:
  \begin{enumerate}
    \item \emph{\bilingual{ANF conversion}{ANF 转换}。}
      A-范式变换~\cite{flanagan1993-essence-compiling-continuations} 将嵌套应用替换为 $\mathbf{let}$-绑定;由引理~\ref{lem:anf-equiv},它在 $\lambda$-演算中是宏可表达的。
      在反方向,$\mathbf{let}$-绑定由 Felleisen~\cite{felleisen1991-expressive-power} 的命题 3.7(重述为引理~\ref{lem:let-macro})在 $\lambda$-演算中是宏可表达的。
      因此,ANF 转换是 $\lambda$-演算与其 ANF 片段之间的一个宏可表达的等价。
    \item \emph{\bilingual{The translation $\mathcal{T}$}{翻译 $\mathcal{T}$}。}
      每个 ANF $\lambda$-演算构造子映射到继承演算上的一个固定语法抽象(定理~\ref{thm:forward-macro}),满足 Felleisen~\cite[Definition~3.11]{felleisen1991-expressive-power} 的条件 E4。
  \end{enumerate}
  由引理~\ref{lem:macro-composition},复合两个宏可表达翻译,得到一个从 $\lambda$-演算(含 $\mathbf{let}$)到继承演算的宏可表达翻译。
  由于继承演算平凡地包含其自身的构造子,它能宏可表达 $\lambda$-演算所包含的或能宏可表达的每个 $\mathscr{U}$-构造。
  \else
  We must show:
  (a)~inheritance-calculus can macro-express every
  $\mathscr{U}$-construct that the lazy $\lambda$-calculus
  can macro-express, and
  (b)~the converse fails.

  \paragraph{Part (b): the lazy $\lambda$-calculus $\not\geq$
  inheritance-calculus}
  The $\lambda$-calculus contains its own constructors
  ($\lambda$, application, variable) and can macro-express
  $\mathbf{let}$-binding (Lemma~\ref{lem:let-macro}), but
  cannot macro-express the inheritance constructors of
  $\mathscr{U}$
  (Theorem~\ref{thm:cartesian-nonexpressibility}).
  Inheritance-calculus contains those constructors.
  Therefore the lazy $\lambda$-calculus is not at least
  as expressive as inheritance-calculus with respect
  to~$\mathscr{U}$.

  \paragraph{Part (a): inheritance-calculus $\geq$ the lazy
  $\lambda$-calculus}
  We show that inheritance-calculus can macro-express
  every $\mathscr{U}$-construct that the
  $\lambda$-calculus contains or can macro-express.
  The $\lambda$-calculus contains $\lambda$-abstraction,
  application, and variable; it can macro-express
  $\mathbf{let}$-binding
  (Lemma~\ref{lem:let-macro}).
  We must show that inheritance-calculus can
  macro-express all four of these constructs.
  The translation proceeds in two stages:
  \begin{enumerate}
    \item \emph{ANF conversion.}
      The A-normal form
      transformation~\cite{flanagan1993-essence-compiling-continuations}
      replaces nested applications with
      $\mathbf{let}$-bindings; it is macro-expressible
      in the $\lambda$-calculus by
      Lemma~\ref{lem:anf-equiv}.
      In the reverse direction,
      $\mathbf{let}$-bindings are macro-expressible in
      the $\lambda$-calculus by
      Proposition~3.7 of
      Felleisen~\cite{felleisen1991-expressive-power}
      (restated as Lemma~\ref{lem:let-macro}).
      Thus ANF conversion is a macro-expressible
      equivalence between the $\lambda$-calculus and its
      ANF fragment.
    \item \emph{The translation $\mathcal{T}$.}
      Each ANF $\lambda$-calculus construct maps to a
      fixed syntactic abstraction over
      inheritance-calculus
      (Theorem~\ref{thm:forward-macro}), satisfying
      condition~E4 of
      Felleisen~\cite[Definition~3.11]{felleisen1991-expressive-power}.
  \end{enumerate}
  By Lemma~\ref{lem:macro-composition}, composing the
  two macro-expressible translations gives a
  macro-expressible translation from
  the $\lambda$-calculus (with $\mathbf{let}$) into
  inheritance-calculus.
  Since inheritance-calculus trivially contains its own
  constructors, it can macro-express every
  $\mathscr{U}$-construct that the $\lambda$-calculus
  contains or can macro-express.
  \fi
\end{proof}

\section{\bilingual{Multi-Target \texorpdfstring{$\this$}{this} Resolution in Other Systems}{其他系统中的多目标 \texorpdfstring{$\this$}{this} 解析}}
\label{app:scala-multi-target}

\ifInheritanceChinese
本附录说明两个具有代表性的系统——Scala~3 和 NixOS 模块系统——拒绝了继承演算能够自然处理的多目标 $\this$ 解析模式。
\else
This appendix demonstrates that two representative systems, Scala~3
and the NixOS module system, reject the multi-target $\this$ resolution
pattern that inheritance-calculus handles naturally.
\fi

\subsection*{\bilingual{Scala~3}{Scala~3}}

\ifInheritanceChinese
考虑两个独立扩展同一外部类的对象，每个对象各自提供内部 trait 的一份副本：
\else
Consider two objects that independently extend an outer class,
each providing its own copy of an inner trait:
\fi

\begin{lstlisting}[style=scala]
class MyOuter:
  trait MyInner:
    def outer = MyOuter.this

object Object1 extends MyOuter
object Object2 extends MyOuter
object HasMultipleOuters extends Object1.MyInner
                     with Object2.MyInner
\end{lstlisting}

\noindent
\ifInheritanceChinese
Scala~3 以如下错误拒绝了 \texttt{HasMultipleOuters}：
\else
Scala~3 rejects \texttt{HasMultipleOuters} with the error:
\fi

\begin{lstlisting}
trait MyInner is extended twice
object HasMultipleOuters cannot be instantiated since
  it has conflicting base types
  Object1.MyInner and Object2.MyInner
\end{lstlisting}

\noindent
\ifInheritanceChinese
该拒绝并非表面层面的限制，而是 DOT 类型论基础的必然后果~\cite{amin2016-dependent-object-types}。在 DOT 中，一个对象被构造为 $\nu(x : T)\, d$，其中 self 变量~$x$ 绑定到\emph{单一}对象。\texttt{Object1.MyInner} 和 \texttt{Object2.MyInner} 是不同的路径依赖类型，每个都将外部 $\this$ 约束到不同的对象。它们的交集要求 \texttt{outer} 同时返回 \texttt{Object1} 和 \texttt{Object2}，但 DOT 的 self 变量是单值的，使得这一交集不可实现。
\else
The rejection is not a surface-level restriction but a consequence
of DOT's type-theoretic foundations~\cite{amin2016-dependent-object-types}.
In DOT, an object is constructed as $\nu(x : T)\, d$, where the
self variable~$x$ binds to a \emph{single} object.
\texttt{Object1.MyInner} and \texttt{Object2.MyInner} are distinct
path-dependent types, each constraining the outer $\this$
to a different object.
Their intersection requires \texttt{outer} to simultaneously
return \texttt{Object1} and \texttt{Object2}, but DOT's self
variable is single-valued, making this intersection unrealizable.
\fi

\ifInheritanceChinese
等价的继承演算定义如下：
\else
The equivalent inheritance-calculus definition is:
\fi
\begin{align*}
  \{ \quad \mathrm{MyOuter} &\mapsto \{
      \mathrm{MyInner} \mapsto \{
    \mathrm{outer} \mapsto \{\qualifiedthis{\mathrm{MyOuter}}\}\}\}, \\
    \mathrm{Object1} &\mapsto \{\mathrm{MyOuter}\}, \\
    \mathrm{Object2} &\mapsto \{\mathrm{MyOuter}\}, \\
    \mathrm{HasMultipleOuters} &\mapsto \{
      \mathrm{Object1}.\mathrm{MyInner},\;
    \mathrm{Object2}.\mathrm{MyInner}\}
  \quad \}
\end{align*}
\ifInheritanceChinese
此定义在继承演算中是良定义的。$\mathrm{outer}$ 内部的引用 $\qualifiedthis{\mathrm{MyOuter}}$ 具有 de~Bruijn 索引 $n = 1$：从 $\mathrm{outer}$ 的封闭作用域 $\mathrm{MyInner}$ 出发，经过一步 $\this$ 便到达 $\mathrm{MyOuter}$。当 $\mathrm{HasMultipleOuters}$ 同时继承 $\mathrm{Object1}.\mathrm{MyInner}$ 和 $\mathrm{Object2}.\mathrm{MyInner}$ 时，$\this$ 函数（方程~\ref{eq:this}）通过在 $\supers(\mathrm{HasMultipleOuters})$ 中搜索与 $\mathrm{MyInner}$ 定义位置匹配的覆盖路径来解析 $\qualifiedthis{\mathrm{MyOuter}}$。它找到两条继承位置路径，一条经由 $\mathrm{Object1}$，另一条经由 $\mathrm{Object2}$，并返回两者。两条路径都通向继承自同一 $\mathrm{MyOuter}$ 的记录，因此查询 $\mathrm{HasMultipleOuters}.\mathrm{outer}$ 无论走哪条路由都能得到 $\mathrm{MyOuter}$ 的 $\properties$。由于属性解析通过集合并来聚合结果（这是一种本质上幂等的操作），来自两条路由的相同属性定义无冲突地合并。因此，该演算消除了对任何人为路由优先级或线性化规则的需求。
\else
This is well-defined in inheritance-calculus.
The reference
$\qualifiedthis{\mathrm{MyOuter}}$ inside $\mathrm{outer}$
has de~Bruijn index $n = 1$:
starting from $\mathrm{outer}$'s enclosing scope
$\mathrm{MyInner}$, one $\this$ step reaches
$\mathrm{MyOuter}$.
When $\mathrm{HasMultipleOuters}$ inherits from both
$\mathrm{Object1}.\mathrm{MyInner}$ and
$\mathrm{Object2}.\mathrm{MyInner}$,
the $\this$ function (equation~\ref{eq:this}) resolves
$\qualifiedthis{\mathrm{MyOuter}}$ by searching through
$\supers(\mathrm{HasMultipleOuters})$ for override paths matching
$\mathrm{MyInner}$'s definition site.
It finds two inheritance-site paths, one through
$\mathrm{Object1}$ and one through $\mathrm{Object2}$, and
returns both.
Both paths lead to records that inherit from the same
$\mathrm{MyOuter}$, so querying
$\mathrm{HasMultipleOuters}.\mathrm{outer}$ yields the
$\properties$ of $\mathrm{MyOuter}$ regardless of which
route is taken.
Because attribute resolution aggregates results via set union (which is an intrinsically idempotent operation), identical property definitions from the two routes merge without conflict. Thus, the calculus eliminates the need for any artificial route-priority or linearization rules.
\fi

\subsection*{\bilingual{NixOS Module System}{NixOS 模块系统}}

\ifInheritanceChinese
NixOS 模块系统以静态错误拒绝了相同的模式。该翻译使用了模块系统自身的抽象：\texttt{deferredModule} 对应 trait（即未求值的模块值，被导入到其他不动点中），\texttt{submoduleWith} 对应对象（在其自身的不动点中求值）。
\else
The NixOS module system rejects the same pattern with a static error.
The translation uses the module system's own abstractions:
\texttt{deferredModule} for traits, which are unevaluated module values imported
into other fixpoints, and \texttt{submoduleWith} for objects, which are evaluated
in their own fixpoints.
\fi

\begin{lstlisting}[style=nix]
let
  lib = (import <nixpkgs> { }).lib;
  result = lib.evalModules {
    modules = [
      (toplevel@{ config, ... }: {
        # class MyOuter { trait MyInner {
        #   def outer = MyOuter.this } }
        options.MyOuter = lib.mkOption {
          default = { };
          type = lib.types.deferredModuleWith {
            staticModules = [
              (MyOuter: {
                options.MyInner = lib.mkOption {
                  default = { };
                  type = lib.types.deferredModuleWith {
                    staticModules = [
                      (MyInner: {
                        options.outer = lib.mkOption {
                          default = { };
                          type =
                            lib.types.deferredModuleWith {
                              staticModules =
                                [ toplevel.config.MyOuter ];
                            };
                        };
                      })
                    ];
                  };
                };
              })
            ];
          };
        };
        # object Object1 extends MyOuter
        options.Object1 = lib.mkOption {
          default = { };
          type = lib.types.submoduleWith {
            modules = [ toplevel.config.MyOuter ];
          };
        };
        # object Object2 extends MyOuter
        options.Object2 = lib.mkOption {
          default = { };
          type = lib.types.submoduleWith {
            modules = [ toplevel.config.MyOuter ];
          };
        };
        # object HasMultipleOuters extends
        #   Object1.MyInner with Object2.MyInner
        options.HasMultipleOuters = lib.mkOption {
          default = { };
          type = lib.types.submoduleWith {
            modules = [
            toplevel.config.Object1.MyInner
            toplevel.config.Object2.MyInner
          ];
          };
        };
      })
    ];
  };
in
  builtins.attrNames result.config.HasMultipleOuters.outer
\end{lstlisting}

\noindent
\ifInheritanceChinese
NixOS 模块系统以如下错误拒绝了此代码：
\else
The NixOS module system rejects this with:
\fi

\begin{lstlisting}
error: The option `HasMultipleOuters.outer'
  in `<unknown-file>'
  is already declared
  in `<unknown-file>'.
\end{lstlisting}

\noindent
\ifInheritanceChinese
该拒绝发生的原因是：\texttt{Object1.\allowbreak{}My\-Inner} 和 \texttt{Object2.\allowbreak{}My\-Inner} 各自携带了 \texttt{My\-Inner} 延迟模块的完整 \texttt{static\-Modules}，其中包括声明 \texttt{options.\allowbreak{}outer}。当 \texttt{HasMultipleOuters} 同时导入两者时，模块系统的 \texttt{merge\-Option\-Decls} 遇到了同一选项的两个声明，并产生静态错误；它无法表达两个声明源自同一 trait 且应被统一这一事实。
\else
The rejection occurs because
\texttt{Object1.\allowbreak{}My\-Inner} and
\texttt{Object2.\allowbreak{}My\-Inner} each carry the full
\texttt{static\-Modules} of the \texttt{My\-Inner} deferred
module, including the declaration
\texttt{options.\allowbreak{}outer}.
When \texttt{HasMultipleOuters} imports both, the module system's
\texttt{merge\-Option\-Decls} encounters two declarations of the
same option and raises a static error; it cannot express that
both declarations originate from the same trait and should be
unified.
\fi

\ifInheritanceChinese
该拒绝与上文 Scala 的``冲突基类型''错误类似：两个系统都假设每个选项或类型成员只有单一的声明位置，并拒绝两条继承路由独立引入同一声明的多目标情形。在继承演算中，相同的模式是良定义的，因为 $\overrides$（方程~\ref{eq:overrides}）能够识别两条路由通向同一定义，而 $\supers$（方程~\ref{eq:supers}）则无重复地收集两个继承位置上下文。
\else
The rejection is analogous to Scala's ``conflicting base types''
error above: both systems assume that each option or type member
has a single declaration site, and reject the multi-target situation
where two inheritance routes introduce the same declaration
independently.
In inheritance-calculus, the same pattern is well-defined because
$\overrides$ (equation~\ref{eq:overrides}) recognizes that both
routes lead to the same definition, and $\supers$
(equation~\ref{eq:supers}) collects both inheritance-site
contexts without duplication.
\fi

\section{Encoding Datalog in Inheritance-Calculus}
\label{app:datalog-encoding}

Section~\ref{sec:emergent-phenomena} observed that
inheritance-calculus natively produces the relational semantics
of logic programming.
This appendix makes that observation precise by giving a
systematic encoding of Datalog into inheritance-calculus.
The encoding uses no extensions to the calculus; it is
expressible in the three primitives of Section~\ref{sec:syntax}.
All examples in this appendix have executable MIXINv2
implementations\anon[~(supplementary material)]{~\cite{mixin2025}}.
There are three ingredients to keep separate.
First, relations are stored as tries, so tuple lookup is indexed
by prefixes rather than by scanning a flat set of facts.
Second, the control skeleton is written in the visitor-pattern
style already used in Section~\ref{sec:case-study}; from a
functional-programming viewpoint, this is the same role played by
Scott-encoded case analysis.
Third, inheritance union and tabled evaluation give the resulting
program its relational, least-fixed-point semantics.
The appendix proceeds in that order: data representation first,
then control flow, then recursion.

\subsection{Relations as Tries}

A binary relation over a finite domain $D = \{a, b, c, \ldots\}$
is represented as a \emph{trie}: a mixin tree whose paths
encode tuples.
The trie is built from three base scopes:
$\mathrm{Relation}$ (abstract base),
$\mathrm{Cons}$ (non-empty node), and $\mathrm{Nil}$ (terminator).
Each constant $d \in D$ inherits from $\mathrm{Cons}$ and carries
a $\mathrm{TailMap}$ that provides per-constant sub-tries:
\begin{align*}
  &\mathrm{Relation} \mapsto \{\},\quad
  \mathrm{Cons} \mapsto \{\mathrm{Relation},\;
  \mathrm{Tail} \mapsto \{\mathrm{Relation}\}\},\quad
  \mathrm{Nil} \mapsto \{\mathrm{Relation}\}\\[2pt]
  &d \mapsto \{\mathrm{Cons},\;
    \mathrm{Tail} \mapsto \{\},\;
  \mathrm{TailMap} \mapsto \{d \mapsto \{\mathrm{Tail}\}\}\}
\end{align*}
A fact $R(d_1, d_2)$ is encoded as a trie path of depth~2,
with a $\mathrm{Nil}$ terminator.
The per-constant sub-tries are accessed via
$\mathrm{TailMap}.d$:
\begin{align*}
  R(a, b) \;\leadsto\quad
  R \mapsto \{a,\; \mathrm{TailMap} \mapsto \{a \mapsto
  \{b,\; \mathrm{TailMap} \mapsto \{b \mapsto \{\mathrm{Nil}\}\}\}\}\}
\end{align*}
Multiple facts for the same relation are accumulated by
inheritance (trie union).
Adding $R(b, c)$:
\[
  R \mapsto \left\{
    \begin{aligned}
      &a,\; b,\\
      &\mathrm{TailMap} \mapsto \left\{
        \begin{aligned}
          &a \mapsto \{b,\; \mathrm{TailMap} \mapsto \{b \mapsto \{\mathrm{Nil}\}\}\},\\
          &b \mapsto \{c,\; \mathrm{TailMap} \mapsto \{c \mapsto \{\mathrm{Nil}\}\}\}
      \end{aligned}\right\}
  \end{aligned}\right\}
\]
The point of the trie representation is not merely that it is a
tree-shaped encoding.
It also provides the indexing structure used by the subsequent
joins: once a constant $d$ has been selected, the encoding moves
directly to $\mathrm{TailMap}.d$ and continues from the matching
sub-trie.
This is why the appendix speaks about tries rather than about an
arbitrary encoding of finite sets of tuples.

\subsection{Infrastructure}

Each constant $d$ carries four groups of boilerplate
definitions used by the encoding.
These play a role analogous to derived type-class instances
in Haskell: each constant must provide a fixed set of
properties that the encoding's generic patterns rely on.
Inheritance-calculus has no built-in derivation mechanism,
so the boilerplate is written out manually; however, the
amount of boilerplate per constant is $O(1)$
and does not grow with the domain size $|D|$,
so adding a new constant
does not require modifying any existing scope, preserving
the expression-problem property.

% 我们可以通过引入类似ECMAScript的Symbol机制来消除这些重复定义，但这让语言更复杂了

\begin{description}
  \item[Acceptance]
    The visitor-pattern infrastructure from
    Section~\ref{sec:case-study}.
    Operationally, this is the object-oriented spelling of a
    case split on the current constant.
    Each constant $d$ defines a $\mathrm{VisitorMap}$
    containing a branch keyed by $d$, a $\mathrm{Visitor}$
    with a $\mathrm{Visited}$ slot for per-constant results,
    and an $\mathrm{Accepted}$ property that delegates to the
    matching branch's $\mathrm{Visited}$ value.
    This is the dispatch mechanism: given a scope that may
    contain multiple constants, $\mathrm{Acceptance}$ selects the
    branch corresponding to each constant present and
    collects the per-constant $\mathrm{Visited}$ results.

  \item[WithVisitorTail]
    Maps each constant $d$ to its sub-trie entry
    $\mathrm{TailMap}.d$, labeled as $\mathrm{VisitorTail}$.
    This provides the right-hand side of a join\footnote{
      In the context of the visitor pattern, this acts as the visitor.
    }: when two relations' $\mathrm{VisitorMap}$ entries are
    composed, $\mathrm{VisitorTail}$ carries the joined
    relation's row data for the matched constant.

  \item[WithVisiteeTail]
    Maps each constant $d$ to its sub-trie entry
    $\mathrm{TailMap}.d$, labeled as $\mathrm{VisiteeTail}$.
    This provides the left-hand side\footnote{
      This acts as the visitee.
    }: the
    relation being iterated. Together,
    $\mathrm{WithVisitorTail}$ and $\mathrm{WithVisiteeTail}$
    enable value-level joins where only constants present
    in \emph{both} relations produce output.

  \item[Replacement]
    Given a $\mathrm{NewTail}$ value, produces a
    $\mathrm{Replaced}$ scope that inherits the label $d$
    but substitutes $\mathrm{TailMap}.d$ with
    $\mathrm{NewTail}$.
    This is the output construction mechanism: it
    builds output tuples by substituting new sub-tries into
    constant wrappers.
    Crucially, $\mathrm{Replacement}$ is defined on the
    per-constant $\mathrm{TailMap}$ entry (not on the merged
    scope), ensuring that output construction operates on
    individual constants rather than their union.
\end{description}

As a running example, we use the executable two-hop encoding
$\mathrm{twoHop}(X, Z) \mathrel{{:}\!{-}} \mathrm{edge}(X, Y),\; \mathrm{edge}(Y, Z)$
from
\anon[\nolinkurl{RelationalTwoHopRepeatedVisitor.mixin.yaml}~(supplementary material)]{\nolinkurl{RelationalTwoHopRepeatedVisitor.mixin.yaml}~\cite{mixin2025}}.
Its top-level scopes are named after the stages of the
encoding: \texttt{\_XAcceptance} iterates the first column,
\texttt{\_YAcceptance} performs the shared-variable join,
\texttt{\_NilAcceptance0} confirms the end of the first fact,
\texttt{\_ZAcceptance} iterates the output column, and
\texttt{\_NilAcceptance1} triggers the final
\texttt{Replacement} chain.
The schematic encodings below should be read as abstractions
of that executable MIXINv2 file.
If one temporarily ignores relational accumulation and reads
each dispatch as an ordinary case split, the control flow is:
choose $X$ from the first relation, choose a matching $Y$,
confirm that the first tuple is complete, choose $Z$ from the
second relation's matching row, confirm that the second tuple
is complete, and emit $(X, Z)$.
The rest of the appendix explains how that familiar
single-valued control skeleton is realised by visitor dispatch
over tries and then lifted to relational semantics by
inheritance union.

\subsection{Encoding Rules}

A Datalog rule $H(\bar{X}) \mathrel{{:}\!{-}} B_1(\bar{Y}_1), \ldots, B_n(\bar{Y}_n)$
is encoded by nesting three mechanisms inside that control
skeleton:

\begin{enumerate}
  \item \textbf{Column iteration.}
    For each column position, an $\mathrm{Acceptance}$ dispatch
    iterates over the constants present.
    The $\mathrm{VisitorMap}$ is populated from
    $\mathrm{WithVisiteeTail}$, so each matched constant's
    $\mathrm{Visitor}$ receives $\mathrm{VisiteeTail}$: the
    per-constant sub-trie at that position.
    The result of processing each constant is placed in
    $\mathrm{Visited}$, and $\mathrm{Accepted}$ collects all
    per-constant results.
    (Details in Sections~\ref{sec:encoding-variables},
    \ref{sec:encoding-constants}, and~\ref{sec:encoding-multi}.)

  \item \textbf{Join.}
    For each shared variable between body atoms, the inner
    $\mathrm{Acceptance}$'s $\mathrm{VisitorMap}$ inherits from
    \emph{both} $\mathrm{WithVisitorTail}$ (from the joined
    relation) and $\mathrm{WithVisiteeTail}$ (from the iterated
    relation).
    A constant $d$ produces a $\mathrm{Visitor}$ entry only if
    $d$ appears in both relations, achieving a value-level
    join.
    The $\mathrm{Visitor}$ then receives both
    $\mathrm{VisiteeTail}$ and $\mathrm{VisitorTail}$,
    providing access to both relations' data for the
    matched constant.

  \item \textbf{Output construction.}
    Inside the innermost $\mathrm{Nil}$ confirmation (verifying
    that each body atom's tuple is complete), chain
    $\mathrm{Replacement}$ operations from the innermost
    output column outward, each setting $\mathrm{NewTail}$ to
    the previous $\mathrm{Replaced}$ value, terminated by
    $\mathrm{Nil}$.
    The final $\mathrm{Replaced}$ value is assigned to
    $\mathrm{Visited}$.
\end{enumerate}

\subsection{Variable Join}
\label{sec:encoding-variables}

Consider the simplest join:
$\mathrm{composed}(X, Z) \mathrel{{:}\!{-}} \alpha(X, Y),\; \beta(Y, Z)$.
The executable self-join instance of this pattern is
\anon[\nolinkurl{RelationalTwoHopRepeatedVisitor.mixin.yaml}~(supplementary material)]{\nolinkurl{RelationalTwoHopRepeatedVisitor.mixin.yaml}~\cite{mixin2025}},
whose scope names are shown in parentheses below.
This is the right place to see the division of labour from the
section introduction: visitor dispatch provides the control flow,
the trie provides direct access to the matching tails, and
$\supers$ accumulates the outputs produced by each successful
match.

The shared variable $Y$ requires matching constants
between $\alpha$'s second column and $\beta$'s first column.
The encoding nests five $\mathrm{Acceptance}$ dispatches:
iterate $X$ from $\alpha$, join on $Y$ between $\alpha$
and $\beta$, confirm $\mathrm{Nil}$ after $Y$ in $\alpha$,
iterate $Z$ from $\beta$'s matched row, and confirm
$\mathrm{Nil}$ after $Z$.
We present each level as a named sub-scope.

\paragraph{Level~1: iterate $X$ from $\alpha$}
(Executable scope: \texttt{\_XAcceptance}.)
{\small
  \[
    \mathrm{xAcc} \mapsto \left\{
      \begin{aligned}
        &\alpha.\mathrm{Acceptance},\; \mathrm{VisitorMap} \mapsto \{\alpha.\mathrm{WithVisiteeTail}\},\\
        &\mathrm{Visitor} \mapsto \left\{
          \begin{aligned}
            &\_X \mapsto \{\mathrm{VisiteeTail}\},\; \ldots,\\
            &\mathrm{Visited} \mapsto \{\mathrm{yAcc}.\mathrm{Accepted}\}
        \end{aligned}\right\}
    \end{aligned}\right\}
\]}

\paragraph{Level~2: join $\alpha$'s $Y$-column with $\beta$'s first column}

(Executable scope: \texttt{\_YAcceptance}.)

The $\mathrm{VisitorMap}$ inherits from \emph{both}
$\beta.\mathrm{WithVisitorTail}$ and
$\_X.\mathrm{WithVisiteeTail}$, producing an entry for
constant $d$ only if $d$ appears in both $\alpha$'s
second column (via $\_X$) and $\beta$'s first column.
{\small
  \[
    \mathrm{yAcc} \mapsto \left\{
      \begin{aligned}
        &\_X.\mathrm{Acceptance},\\
        &\mathrm{VisitorMap} \mapsto \{\beta.\mathrm{WithVisitorTail},\; \_X.\mathrm{WithVisiteeTail}\},\\
        &\mathrm{Visitor} \mapsto \left\{
          \begin{aligned}
            &\_Y_0 \mapsto \{\mathrm{VisiteeTail}\},\; \mathrm{VisitorTail} \mapsto \{\beta.\mathrm{Tail}\},\\
            &\ldots,\\
            &\mathrm{Visited} \mapsto \{\mathrm{nilAcc}_0.\mathrm{Accepted}\}
        \end{aligned}\right\}
    \end{aligned}\right\}
\]}

\paragraph{Level~3: confirm $\mathrm{Nil}$ after $Y$ in $\alpha$}

(Executable scope: \texttt{\_NilAcceptance0}.)

This checks that $\alpha(X, Y)$ is a complete two-column
fact (its trie path terminates at $\mathrm{Nil}$ after $Y$).
Inside the $\mathrm{Nil}$ branch, the encoding proceeds to
iterate $\beta$'s second column via $\mathrm{VisitorTail}$:
{\small
  \[
    \mathrm{nilAcc}_0 \mapsto \left\{
      \begin{aligned}
        &\_Y_0.\mathrm{Acceptance},\\
        &\mathrm{VisitorMap} \mapsto \left\{\mathrm{Nil} \mapsto \left\{
            \begin{aligned}
              &\ldots,\\
              &\mathrm{Visited} \mapsto \{\mathrm{zAcc}.\mathrm{Accepted}\}
        \end{aligned}\right\}\right\}
    \end{aligned}\right\}
\]}

\paragraph{Level~4: iterate $Z$ from $\beta$'s matched row}
(Executable scope: \texttt{\_ZAcceptance}.)
{\small
  \[
    \mathrm{zAcc} \mapsto \left\{
      \begin{aligned}
        &\mathrm{VisitorTail}.\mathrm{Acceptance},\; \mathrm{VisitorMap} \mapsto \{\mathrm{VisitorTail}.\mathrm{WithVisiteeTail}\},\\
        &\mathrm{Visitor} \mapsto \left\{
          \begin{aligned}
            &\_Z \mapsto \{\mathrm{VisiteeTail}\},\; \ldots,\\
            &\mathrm{Visited} \mapsto \{\mathrm{nilAcc}_1.\mathrm{Accepted}\}
        \end{aligned}\right\}
    \end{aligned}\right\}
\]}

\paragraph{Level~5: confirm $\mathrm{Nil}$ after $Z$ and construct output}
(Executable scope: \texttt{\_NilAcceptance1}, with
\texttt{\_ZReplacement} and \texttt{\_XReplacement}.)
{\small
  \[
    \mathrm{nilAcc}_1 \mapsto \left\{
      \begin{aligned}
        &\mathrm{VisiteeTail}.\mathrm{Acceptance},\\
        &\mathrm{VisitorMap} \mapsto \left\{\mathrm{Nil} \mapsto \left\{
            \begin{aligned}
              &\mathrm{zR} \mapsto \{\_Z.\mathrm{Replacement},\; \mathrm{NewTail} \mapsto \{\mathrm{Nil}\}\},\\
              &\mathrm{xR} \mapsto \{\_X.\mathrm{Replacement},\; \mathrm{NewTail} \mapsto \{\mathrm{zR}.\mathrm{Replaced}\}\},\\
              &\mathrm{Visited} \mapsto \{\mathrm{xR}.\mathrm{Replaced}\}
        \end{aligned}\right\}\right\}
    \end{aligned}\right\}
\]}

\paragraph{Final result}
\[
  \mathrm{composed} \mapsto \{\mathrm{xAcc}.\mathrm{Accepted}\}
\]

\paragraph{Semantics}
At Level~2, the combined $\mathrm{VisitorMap}$ ensures that a
constant $d$ produces a $\mathrm{Visitor}$ only when $d$
appears in both $\alpha$'s second column \emph{and}
$\beta$'s first column.
The $\mathrm{Visitor}$ for each matched constant $d$
independently receives $\mathrm{VisiteeTail}$ ($\alpha$'s
data for~$d$) and $\mathrm{VisitorTail}$ ($\beta$'s row
for~$d$).
Output construction (Level~5) operates inside this
per-constant $\mathrm{Visitor}$, so $\mathrm{Replacement}$
acts on the \emph{individual} constant's data, not on the
merged scope.
Each matched constant $d$ produces its own output tuples
in $\mathrm{Visited}$, and these are collected by
$\mathrm{Accepted}$ via trie union ($\supers$).
This yields standard Datalog tuple-level semantics: only
constants that actually match between both relations
contribute to the output.
For the concrete instance
$\mathrm{edge} = \{(a,b),\; (b,c),\; (c,a)\}$, as encoded in
\anon[\nolinkurl{RelationalTwoHopRepeatedVisitor.mixin.yaml}~(supplementary material)]{\nolinkurl{RelationalTwoHopRepeatedVisitor.mixin.yaml}~\cite{mixin2025}},
Level~1 selects $a$, Level~2 matches $b$, Level~3 confirms
that $\mathrm{edge}(a,b)$ ends in $\mathrm{Nil}$, Level~4
iterates $\mathrm{edge}.\mathrm{TailMap}.b$ and selects $c$,
and Level~5 confirms $\mathrm{Nil}$ again before the two
$\mathrm{Replacement}$ operations reconstruct $(a,c)$.

\subsection{Constant Match and Navigation}
\label{sec:encoding-constants}

Constants in Datalog atoms use two mechanisms, neither of
which is a join.
They are simpler than variable joins because no value from a
second relation needs to be matched against the currently chosen
constant:

\paragraph{Type existence match}
An atom $R(X, b)$ with a constant in the second position
requires checking that constant $b$ exists in $R$'s
second column at the matched row.
This is encoded as an $\mathrm{Acceptance}$ with a
hand-crafted $\mathrm{VisitorMap}$ containing \emph{only}
the $b$ branch (instead of composing from
  $\mathrm{WithVisitorTail}$, which would create branches for
all constants):
{\small
  \begin{align*}
    \mathrm{VisitorMap} \mapsto \{b \mapsto \{\ldots\}\}
\end{align*}}
If $b$ exists in the dispatching scope, the
$b$ branch fires; otherwise, $\mathrm{Accepted}$
receives no matching branch.

\paragraph{Direct navigation}
An atom $R(b, Z)$ with a constant in the first position
navigates directly to $b$'s sub-trie:
\[
  \mathrm{edgeB}_Z \mapsto \{R.\mathrm{TailMap}.b\}
\]
This bypasses the visitor mechanism entirely, accessing
the per-constant $\mathrm{TailMap}$ entry to restrict the
scope to only $b$'s data.

\paragraph{Example}
$\mathrm{viaB}(X, Z) \mathrel{{:}\!{-}} \mathrm{edge}(X, b),\; \mathrm{edge}(b, Z)$
combines both mechanisms.
The executable MIXINv2 encoding is
\anon[\nolinkurl{RelationalConstant.mixin.yaml}~(supplementary material)]{\nolinkurl{RelationalConstant.mixin.yaml}~\cite{mixin2025}}.
The encoding iterates $\mathrm{edge}$ for $X$ (Level~1),
then checks for constant $b$ in $X$'s row with a
hand-crafted $\mathrm{VisitorMap}$ (Level~2), navigates
$\mathrm{edge}.\mathrm{TailMap}.b$ for $Z$ (Level~3),
confirms $\mathrm{Nil}$ (Level~4), and constructs $(X, Z)$
via $\mathrm{Replacement}$.

\subsection{Multi-Body Rules}
\label{sec:encoding-multi}

Rules with $n > 2$ body atoms chain additional join levels.
Each shared variable introduces one combined
$\mathrm{WithVisitorTail}$/$\mathrm{WithVisiteeTail}$ join
inside the previous level's $\mathrm{Visitor}$,
and the output construction is placed inside the innermost
$\mathrm{Nil}$ confirmation.
So the binary-rule encoding is the whole story structurally:
larger bodies only add more nested copies of the same pattern.

\paragraph{Three-body rule}
$\mathrm{chain}(X, W) \mathrel{{:}\!{-}} \alpha(X, Y),\; \beta(Y, Z),\; \gamma(Z, W)$
requires two joins: on $Y$ (between $\alpha$ and $\beta$)
and on $Z$ (between $\beta$ and $\gamma$).
The structure is: iterate $X$ from $\alpha$, join on $Y$
(getting both $\alpha$'s and $\beta$'s data), confirm
$\mathrm{Nil}$ after $Y$ in $\alpha$, then inside the
$\mathrm{Nil}$ branch join $\beta$'s $Z$-column with
$\gamma$'s first column, iterate $W$ from $\gamma$'s
matched row, confirm $\mathrm{Nil}$, and construct
$(X, W)$.
The $Z$-join is nested inside the $Y$-join's
$\mathrm{Nil}$ confirmation, ensuring that $\beta$'s data
for the matched $Y$ constant (available via
$\mathrm{VisitorTail}$) is used to iterate $\beta$'s
second column.

\paragraph{Mixed constants and variables}
$\mathrm{constJoin}(X, W) \mathrel{{:}\!{-}} \mathrm{edge}(X, Y),\; \mathrm{edge}(Y, b),\; \mathrm{edge}(b, W)$
combines a variable join on $Y$, a constant match on $b$
(hand-crafted $\mathrm{VisitorMap}$ with only the $b$ entry),
and direct navigation via $\mathrm{edge}.\mathrm{TailMap}.b$
for $W$.
The corresponding executable MIXINv2 file is
\anon[\nolinkurl{RelationalConstantJoin.mixin.yaml}~(supplementary material)]{\nolinkurl{RelationalConstantJoin.mixin.yaml}~\cite{mixin2025}}.

\subsection{Multi-Variable Join}

When a rule shares \emph{multiple} variables between two
body atoms, each shared variable requires its own
$\mathrm{Acceptance}$ check.
The checks are chained: the outer $\mathrm{Acceptance}$
joins on one column pair using combined
$\mathrm{WithVisitorTail}$/$\mathrm{WithVisiteeTail}$, and
its $\mathrm{Visitor}$ contains an inner $\mathrm{Acceptance}$
joining on another column pair.

\paragraph{Example}
$\mathrm{multiJoin}(X, Y) \mathrel{{:}\!{-}} r(X, Y),\; s(X, Y)$
shares both $X$ and $Y$.
The executable MIXINv2 encoding is
\anon[\nolinkurl{RelationalMultiVariableJoin.mixin.yaml}~(supplementary material)]{\nolinkurl{RelationalMultiVariableJoin.mixin.yaml}~\cite{mixin2025}}.
The outer $\mathrm{Acceptance}$ joins $r$'s and $s$'s first
columns (variable $X$) using combined
$\mathrm{WithVisitorTail}$/$\mathrm{WithVisiteeTail}$.
Inside the outer $\mathrm{Visitor}$, an inner
$\mathrm{Acceptance}$ joins $r$'s and $s$'s second columns
(variable $Y$) similarly.
Both joins must match for the output to be produced, and
the per-constant $\mathrm{Visited}$ pattern ensures that
only constants present in both relations contribute.

\subsection{Recursive Rules}
\label{sec:recursive-rules}

Recursive Datalog rules, such as transitive closure
$\mathrm{path}(X, Y) \mathrel{{:}\!{-}} \mathrm{edge}(X, Z),\; \mathrm{path}(Z, Y)$,
are encoded identically to non-recursive rules: the
head relation $\mathrm{path}$ appears both as a defined scope
and as a reference in the body (providing
$\mathrm{WithVisitorTail}$ for the $Z$-join).
This should not be read as ordinary $\this$ of a
single-valued lazy program.
Operationally, tabled evaluation
(Section~\ref{sec:mixin-trees},
Appendix~\ref{app:well-definedness})
computes $\mathrm{lfp}(T_P)$: each query into
$\mathrm{path}$ triggers further evaluation through
the recursive body, and the trie union ($\supers$)
accumulates all reachable tuples.
This corresponds exactly to \emph{Naive
evaluation}~\cite{bancilhon1986-recursive-query-strategies} of the
$T_P$
operator~\cite{vanemden1976-predicate-logic-semantics}:
each iteration applies every rule to the current
fact set and adds newly derived facts.
Because Datalog-encoded relations have finitely many
constants per query, the per-query answer domain is
finite, and convergence is guaranteed by the ascending
chain condition on a finite lattice.
\emph{Semi-Naive} evaluation~\cite{bancilhon1986-recursive-query-strategies},
which restricts each iteration to rules that use at
least one newly derived fact, is a natural
optimisation left for future work.

\paragraph{The recursive rule as monadic pseudocode}
The recursive rule above roughly corresponds to the
following pseudocode (Haskell-like syntax with
  \texttt{RebindableSyntax} and \texttt{OverloadedLists}
  semantics, where do-notation desugars to a
\texttt{Set}-based bind and list literals denote sets):
\begin{lstlisting}[style=haskell]
{-# LANGUAGE RebindableSyntax, OverloadedLists #-}
s >>= f = unions [f x | x <- s]
return x = [x]
guard True  = [()]
guard False = []
edge = [("a","b"), ("b","c"), ("c","a")]
path = edge <> do
  (x, y0) <- edge
  (y1, z) <- path
  guard (y0 == y1)
  return (x, z)
\end{lstlisting}
\noindent
Each line of the pseudocode has a direct counterpart in
the inheritance-calculus encoding
(Section~\ref{sec:encoding-variables}):
\begin{description}
  \item[\texttt{edge}] corresponds to the base-fact trie;
  \item[\texttt{path = edge <> do \ldots}]
    corresponds to $\mathrm{path}$ inheriting $\mathrm{edge}$
    (base facts) and the recursive body (derived facts via
    $\supers$);
  \item[\texttt{(x, y0) <- edge}]
    corresponds to Level~1 ($\mathrm{xAcc}$), iterating $X$
    from $\mathrm{edge}$;
  \item[\texttt{(y1, z) <- path}]
    corresponds to Level~2 ($\mathrm{yAcc}$), joining on $Y$
    with $\mathrm{path}$, and Level~4 ($\mathrm{zAcc}$),
    iterating $Z$;
  \item[\texttt{guard (y0 == y1)}]
    corresponds to the $\mathrm{WithVisitorTail} \cap
    \mathrm{WithVisiteeTail}$ intersection in Level~2: a
    constant produces a $\mathrm{Visitor}$ entry only if it
    appears in both relations;
  \item[\texttt{return (x, z)}]
    corresponds to the $\mathrm{Replacement}$ chain in
    Level~5, constructing the output tuple $(X, Z)$.
\end{description}
The only difference is the evaluation strategy:
Haskell evaluates \texttt{path} lazily (diverges),
whereas inheritance-calculus evaluates it via tabled
$T_P$ iteration (converges to $\mathrm{lfp}$).

However, this monadic program does not converge because
\texttt{path} is defined recursively.%
\footnote{
  Enabling the \texttt{RecursiveDo} extension does not help:
  \texttt{mfix} ties a knot at the \emph{element} level via laziness,
  whereas the transitive closure equation requires a
  fixed point at the \emph{container} level, that is, finding a
  set~$S$ such that $S = F(S)$.
}
Inheritance-calculus solves this through tabled evaluation
(Section~\ref{sec:mixin-trees}), which computes the
least fixed point of the $T_P$ operator by iterating
over the lattice of mixin trees, precisely the
container-level fixed point that the pseudocode above cannot express.

\subsection{Semantic Correspondence with Standard Datalog}
\label{sec:datalog-semantics}

The encoding produces standard Datalog tuple-level semantics.
Putting the previous subsections together: tries provide indexed
tuple storage, visitor dispatch supplies the control skeleton,
and per-constant $\mathrm{Replacement}$ keeps output
construction local to each witness.
The key mechanism is the per-constant $\mathrm{Visited}$
pattern: output construction ($\mathrm{Replacement}$)
occurs inside each matched constant's $\mathrm{Visitor}$
scope, operating on the \emph{individual} constant's data
rather than on the merged scope.
The $\mathrm{Accepted}$ property then collects all per-constant
$\mathrm{Visited}$ results via trie union ($\supers$).

For example, given $\mathrm{edge} = \{(a,b),\; (b,c),\; (c,a)\}$
(a~3-cycle), the two-hop self-join
$\mathrm{twoHop}(X, Z) \mathrel{{:}\!{-}} \mathrm{edge}(X, Y),\; \mathrm{edge}(Y, Z)$
produces exactly $3$ pairs:
$\{(a,c),\; (b,a),\; (c,b)\}$, matching standard Datalog.
This is the behavior exhibited by
\anon[\nolinkurl{RelationalTwoHopRepeatedVisitor.mixin.yaml}~(supplementary material)]{\nolinkurl{RelationalTwoHopRepeatedVisitor.mixin.yaml}~\cite{mixin2025}}.
When the $Y$-join matches constant $b$ (present in both
  $\mathrm{edge}$'s second column from $(a,b)$ and
$\mathrm{edge}$'s first column), the $\mathrm{Visitor}$ for
$b$ receives $\mathrm{VisitorTail} = \mathrm{edge}.\mathrm{TailMap}.b$
(containing only~$c$), and $\mathrm{Replacement}$ constructs
a single output tuple $(a, c)$, rather than all of $\{a,b,c\}$.

This value-level precision arises because each constant's
$\mathrm{Replacement}$ operates on $\mathrm{TailMap}.d$
(the per-constant sub-trie) rather than on the merged
trie.
The same point is visible in
the repeated-visitor self-join regression example\anon[~(supplementary material)]{~\cite{mixin2025}},
where repeated visitor dispatch prevents witnesses from
different constants from being merged spuriously.
Although $\bases$ (equation~\ref{eq:bases}) still computes
cross-joins at the primitive level, the encoding routes
computation through per-constant $\mathrm{TailMap}$ entries
before any merging occurs, achieving the same effect as
standard Datalog's tuple-level evaluation.

\section{Evaluation Trace of Multi-Target \texorpdfstring{$\this$}{this} Resolution}
\label{app:evaluation-trace}

This appendix gives a step-by-step evaluation trace of the program of
Appendix~\ref{app:scala-multi-target}, computing
$\properties((\text{``HasMultipleOuters''},\, \text{``outer''}))$ by the six
functions of Section~\ref{sec:definitions}
(equations~\ref{eq:properties} to~\ref{eq:this}).
Every call to $\overrides$, $\bases$, $\resolve$, $\this$, and $\supers$
is shown with its arguments and returned set.
The trace lives in the metalanguage: every argument and result is a
\emph{path}, that is, a sequence of labels
$(\ell_1, \ldots, \ell_n)$ identifying a position in the tree
(Section~\ref{sec:definitions}), not an inheritance-calculus expression.
To keep labels distinct from the metalanguage variables of
Section~\ref{sec:definitions} (the italic $p$, $\ell$, $n$, $S$), we
quote each concrete label, so a path is written as a parenthesized,
comma-separated sequence of quoted labels.
Thus $(\text{``HasMultipleOuters''}, \text{``outer''})$ is the path that
the inheritance-calculus projection
$\mathrm{HasMultipleOuters}.\mathrm{outer}$ identifies, and the root path
is the empty sequence $()$.

\paragraph{Program}
\begin{align*}
  \{ \quad \mathrm{MyOuter} &\mapsto \{
      \mathrm{MyInner} \mapsto \{
    \mathrm{outer} \mapsto \{\qualifiedthis{\mathrm{MyOuter}}\}\}\}, \\
    \mathrm{Object1} &\mapsto \{\mathrm{MyOuter}\}, \\
    \mathrm{Object2} &\mapsto \{\mathrm{MyOuter}\}, \\
    \mathrm{HasMultipleOuters} &\mapsto \{
      \mathrm{Object1}.\mathrm{MyInner},\;
    \mathrm{Object2}.\mathrm{MyInner}\}
  \quad \}
\end{align*}

\paragraph{AST primitives}
Parsing populates $\defines$ and $\inherits$
(Section~\ref{sec:definitions}):
\begin{align*}
  \defines(()) &= \{
    \text{``MyOuter''},\;
    \text{``Object1''},\;
    \text{``Object2''},\;
  \text{``HasMultipleOuters''}\}, \\
  \defines((\text{``MyOuter''})) &= \{\text{``MyInner''}\}, \\
  \defines((\text{``MyOuter''}, \text{``MyInner''})) &= \{\text{``outer''}\}, \\
  \inherits((\text{``MyOuter''}, \text{``MyInner''}, \text{``outer''})) &=
  \{(1,\, ())\}, \\
  \inherits((\text{``Object1''})) = \inherits((\text{``Object2''})) &=
  \{(0,\, (\text{``MyOuter''}))\}, \\
  \inherits((\text{``HasMultipleOuters''})) &=
  \{(0,\, (\text{``Object1''}, \text{``MyInner''})),\;
  (0,\, (\text{``Object2''}, \text{``MyInner''}))\}.
\end{align*}
All other $\defines$ and $\inherits$ are empty.
Each $\inherits$ entry is a reference pair $(n,\, w)$,
where $n$ is the number of upward steps and $w$ the
downward projection list (Section~\ref{sec:definitions}).

\paragraph{Trace}
We present the trace as a numbered list of semantic-function calls in
dependency order: a call appears after the calls it depends on, and each
item states, in English, which earlier items it uses, how it computes, and
what it yields. Each distinct call is listed once, so a single item may be
used by several later items, the reuse that an implementation could realize
by tabling~\cite{tamaki1986-tabled-resolution}.
The arguments and results are metalanguage values: a label is a quoted
string such as $\text{``outer''}$, and a path is a tuple of labels such as
$(\text{``HasMultipleOuters''}, \text{``outer''})$, with the root path
written $()$.
The query is
$\properties((\text{``HasMultipleOuters''}, \text{``outer''}))$.

\begin{enumerate}
  \item \label{trace:overrides-MyOuter-MyInner}Compute $\overrides((\text{``MyOuter''}, \text{``MyInner''}))$, the paths sharing the identity of $(\text{``MyOuter''}, \text{``MyInner''})$. It keeps $(\text{``MyOuter''}, \text{``MyInner''})$ and adds any same-label definition reached through the supers of its enclosing scope. Result: $\{(\text{``MyOuter''}, \text{``MyInner''})\}$.
  \item \label{trace:supers-MyOuter-MyInner}Compute $\supers((\text{``MyOuter''}, \text{``MyInner''}))$ using item~\ref{trace:overrides-MyOuter-MyInner}. It takes the reflexive-transitive closure of $\bases$ from $(\text{``MyOuter''}, \text{``MyInner''})$, pairing each reachable override with the inheritance site through which it is reached. Result: $\{((\text{``MyOuter''}), (\text{``MyOuter''}, \text{``MyInner''}))\}$.
  \item \label{trace:overrides-Has-outer}Compute $\overrides((\text{``HasMultipleOuters''}, \text{``outer''}))$, the paths sharing the identity of $(\text{``HasMultipleOuters''}, \text{``outer''})$ using item~\ref{trace:supers-MyOuter-MyInner}. It keeps $(\text{``HasMultipleOuters''}, \text{``outer''})$ and adds any same-label definition reached through the supers of its enclosing scope. Result: $\{(\text{``HasMultipleOuters''}, \text{``outer''}), (\text{``MyOuter''}, \text{``MyInner''}, \text{``outer''})\}$.
  \item \label{trace:bases-Has-outer}Compute $\bases((\text{``HasMultipleOuters''}, \text{``outer''}))$. $(\text{``HasMultipleOuters''}, \text{``outer''})$ carries no inheritance reference ($\inherits = \varnothing$), so it has no direct base. Result: $\varnothing$.
  \item \label{trace:bases-MyOuter-MyInner-outer}Compute $\bases((\text{``MyOuter''}, \text{``MyInner''}, \text{``outer''}))$. $(\text{``MyOuter''}, \text{``MyInner''}, \text{``outer''})$ carries the reference pair $(1, ())$ (take one step upward, then project the empty projection), whose target is obtained by item~\ref{trace:resolve-Has-MyOuter-MyInner-1}.
  \item \label{trace:this-Has-MyOuter-MyInner-1}Compute one $\this$ step: from the frontier $\{(\text{``HasMultipleOuters''})\}$, take the supers of each frontier path and keep those whose definition site is $(\text{``MyOuter''}, \text{``MyInner''})$, collecting their inheritance sites as the new frontier. Result: $\{(\text{``Object1''}), (\text{``Object2''})\}$.
  \item \label{trace:resolve-Has-MyOuter-MyInner-1}Compute $\resolve$ for the reference pair $(1, ())$ defined at $(\text{``MyOuter''}, \text{``MyInner''})$, reached from inheritance site $(\text{``HasMultipleOuters''})$ using item~\ref{trace:this-Has-MyOuter-MyInner-1}. It takes one step upward through $\this$ and then appends the projection. Result: $\{(\text{``Object1''}), (\text{``Object2''})\}$.
  \item \label{trace:resolve-0-MyOuter}Compute $\resolve$ for the reference pair $(0, (\text{``MyOuter''}))$ defined at $()$, reached from inheritance site $()$. It takes no step upward through $\this$ and then appends the projection. Result: $\{(\text{``MyOuter''})\}$.
  \item \label{trace:supers-Has-outer}Compute $\supers((\text{``HasMultipleOuters''}, \text{``outer''}))$ using items~\ref{trace:overrides-Has-outer}, \ref{trace:bases-Has-outer}, \ref{trace:bases-MyOuter-MyInner-outer}, \ref{trace:resolve-Has-MyOuter-MyInner-1} and~\ref{trace:resolve-0-MyOuter}. It takes the reflexive-transitive closure of $\bases$ from $(\text{``HasMultipleOuters''}, \text{``outer''})$, pairing each reachable override with the inheritance site through which it is reached. Result: $\{((), (\text{``MyOuter''})), ((), (\text{``Object1''})), ((), (\text{``Object2''})), ((\text{``HasMultipleOuters''}), (\text{``HasMultipleOuters''}, \text{``outer''})), ((\text{``HasMultipleOuters''}), (\text{``MyOuter''}, \text{``MyInner''}, \text{``outer''}))\}$.
  \item \label{trace:properties}Compute $\properties((\text{``HasMultipleOuters''}, \text{``outer''}))$ using item~\ref{trace:supers-Has-outer}. It unions the $\defines$ of every override path in that $\supers$ set. Only $(\text{``MyOuter''})$ defines a label, namely $\text{``MyInner''}$, reached through both $(\text{``Object1''})$ and $(\text{``Object2''})$; set union collapses the two routes. Result: $\{\text{``MyInner''}\}$.
\end{enumerate}


\begin{remark}
  The crux is the $\this$ step (item~\ref{trace:this-Has-MyOuter-MyInner-1}).
  The inheritance reference at
  $(\text{``MyOuter''}, \text{``MyInner''}, \text{``outer''})$ takes one step
  upward, so $\this$ moves one step up from
  $(\text{``MyOuter''}, \text{``MyInner''})$ and yields the frontier
  $\{(\text{``Object1''}), (\text{``Object2''})\}$, the two inheritance
  sites that incorporated $\text{``MyInner''}$. This is precisely the
  multi-target situation that Scala~3 and the NixOS module system reject
  (Appendix~\ref{app:scala-multi-target}): there $\this$ would have to
  select a single site. Reaching the properties of $(\text{``MyOuter''})$
  then happens because both $(\text{``Object1''})$ and
  $(\text{``Object2''})$ have $(\text{``MyOuter''})$ as a base, so the final
  $\supers$ contains $(\text{``MyOuter''})$ through both routes; since
  attribute resolution aggregates by set union, the two routes collapse to
  the single answer
  $\properties((\text{``HasMultipleOuters''}, \text{``outer''}))
  = \{\text{``MyInner''}\} = \properties((\text{``MyOuter''}))$.
\end{remark}

\fi % end \ifInheritanceAppendix

% Appendix-only build (supplement.tex): no main matter ran above, so emit the
% bibliography here, after the appendix.
\ifInheritanceBody\else
\bibliographystyle{ACM-Reference-Format}
\bibliography{references}
\fi

% \clearpage before closing CJK* ships every remaining page and float while the
% CJK active characters still mean CJK (the running head and any deferred float
% are typeset in the output routine, after \end{CJK*} would otherwise restore the
% default UTF-8 meaning).
\ifInheritanceChinese\clearpage\end{CJK*}\fi
\end{document}
